% Generated mechanically from frozen input p07/ja/etale-cohomology.tex.
% Do not edit this generated file; edit the bound locale source instead.

% OK, start here.
%

\setcounter{chapter}{58}
\chapter{エタール・コホモロジー}



\phantomsection
\label{etale-cohomology-section-phantom}



\section{序論}
\label{etale-cohomology-section-introduction}

\noindent
本章は、スキームのエタール・コホモロジーを扱う一連の章の最初のものである。
ここでは、エタール位相と、この位相におけるアーベル層のコホモロジーについて、
ごく基礎的な事柄を論じる。扱う話題の多くは初読時には安全に読み飛ばせる。
何を読み飛ばすかを決めるための指針については、次節を参照されたい。

\medskip\noindent
本章の初版は、Johan de Jong が 2009 年秋に Columbia University で講義した
エタール・コホモロジーの講義前半のノートをもとに作られた。
当初のノート作成者は Thibaut Pugin、Zachary Maddock、および Min Lee である。
講義の後半は跡公式を扱う章に収録されている。The Trace Formula の
第 \ref{trace-section-introduction} 節を参照されたい。




\section{初読時に読み飛ばすべき節}
\label{etale-cohomology-section-skip}

\noindent
本章の内容は、代数空間、Deligne--Mumford スタック、代数スタックなどに関する
理論の展開に用いることを意図している。そのため、スキームのエタール・
コホモロジーに関する当初の解説へ、かなり技術的な内容を加えた。
このような内容は、「トポス」という語の出現頻度、集合論に関する議論、あるいは
スキームの射のきわめて一般的な性質を扱う証明によって見分けられる。
これらの議論の一部は初読時には読み飛ばしてよい。

\medskip\noindent
とりわけ、次の各節を読み飛ばすことを勧める。
\begin{enumerate}
\item 大トポスと小トポスの比較、第 \ref{etale-cohomology-section-compare} 節。
\item 射の復元、第 \ref{etale-cohomology-section-morphisms} 節。
\item 押し出しと引き戻し、第 \ref{etale-cohomology-section-monomorphisms} 節。
\item 性質 (A)、第 \ref{etale-cohomology-section-A} 節。
\item 性質 (B)、第 \ref{etale-cohomology-section-B} 節。
\item 性質 (C)、第 \ref{etale-cohomology-section-C} 節。
\item 小エタール・サイトの位相的不変性、
第 \ref{etale-cohomology-section-topological-invariance} 節。
\item 整かつ普遍的単射な射、
第 \ref{etale-cohomology-section-integral-universally-injective} 節。
\item 大サイトと順像、第 \ref{etale-cohomology-section-big} 節。
\item 大下方シュリークの完全性、第 \ref{etale-cohomology-section-exactness-lower-shriek} 節。
\end{enumerate}
これらの節のほかにも読み飛ばしてよい結果が散在しており、読者は上記のキーワードに
よってそれらを見分けることができる。



%9.08.09
\section{プロローグ}
\label{etale-cohomology-section-prologue}

\noindent
この講義では、もう一つのコホモロジー理論を扱う。まず注意すべきことは、
Zariski 位相が完全には満足できるものではないということである。
期待する結果を与えない主な理由の一つは、$X$ が複素多様体で
$\mathcal{F}$ が定数層ならば
$$
H^i(X, \mathcal{F}) = 0, \quad \text{すべての } i > 0.
$$
となることである。理由は次のとおりである。既約スキーム（とくに多様体）では、
任意の二つの空でない開部分集合が交わるため、定数層の制限写像は全射となる。
このような層を {\it flasque} であるという。この場合、すべての高次
{\v C}ech コホモロジー群が消滅し、したがってすべての高次 Zariski
コホモロジー群も消滅する。言い換えると、Zariski 位相にはこの高次
コホモロジーを検出するのに「十分な数の」開集合がない。

\medskip\noindent
一方、$X$ が滑らかな射影的複素多様体ならば
$$
H_{Betti}^{2 \dim X}(X (\mathbf{C}), \Lambda) = \Lambda \quad \text{ただし }
\Lambda = \mathbf{Z}, \ \mathbf{Z}/n\mathbf{Z},
$$
である。ここで $X(\mathbf{C})$ は $X$ の複素点全体を表す。
この性質を代数幾何学、とりわけ正標数の場合にも再現できることが望ましい。




\section{エタール位相}
\label{etale-cohomology-section-etale-topology}

\noindent
Zariski 位相を細分するために、単に開集合を余分に「追加」することは非常に難しい。
位相を定義する一つの効率的な方法は、開集合だけでなく、その上にあるいくつかの
スキームも考えることである。エタール位相を定義するには、エタールな射
$\varphi : U \to X$ をすべて考える。$X$ が $\mathbf{C}$ 上の滑らかな射影多様体ならば、
これは次を意味する。
\begin{enumerate}
\item $U$ は滑らかな多様体の非交和であり、
\item $\varphi$ は（解析的に）局所同型である。
\end{enumerate}
ここで「解析的に」とは、$\mathbf{C}$ 上の通常の（超越的）位相に関するものである。
したがって第二の条件は、$\varphi$ の微分が至る所で最大階数をもつことを意味する
（とくに $U$ のすべての連結成分は $X$ と同じ次元をもつ）。

\medskip\noindent
二重被覆、すなわち大まかには多様体間の次数 $2$ の有限写像、たとえば
$$
\Spec(\mathbf{C}[t])
\longrightarrow
\Spec(\mathbf{C}[t]),
\quad t \longmapsto t^2
$$
は、一点のみからなるファイバーをもつならエタール射ではない。この例では
$t = 0$ のときにそうなる。$\mathbf{C}$ 上の多様体間の有限写像がエタールであるには、
すべてのファイバーが同じ個数の点をもたなければならない。この例の写像の始域から
$t = 0$ の点を除けば、射はエタールになる。しかし始域からほかの点も除くことができ、
その後も射はエタールである。エタール位相を考えるには、このような射をすべて
調べなければならない。Zariski 位相と異なり、これら自身は単なる $X$ の開部分集合である
必要はないが、その像は常に開である。

\begin{definition}
\label{etale-cohomology-definition-etale-covering-initial}
射の族 $\{ \varphi_i : U_i \to X\}_{i \in I}$ は、各 $\varphi_i$ がエタール射であり、
それらの像が $X$ を被覆する、すなわち
$X = \bigcup_{i \in I} \varphi_i(U_i)$ であるとき、{\it エタール被覆}と呼ばれる。
\end{definition}

\noindent
これによってエタール位相が「定義」される。言い換えると、これで層が何であるかを
述べられる。集合（それぞれ、アーベル群、ベクトル空間など）の
{\it エタール層} $\mathcal{F}$ とは、$X$ 上の次のデータである。
\begin{enumerate}
\item 各エタール射 $\varphi : U \to X$ に対する集合（それぞれ、アーベル群、
ベクトル空間など）$\mathcal{F}(U)$、
\item エタール・スキームの各対 $U, \ U'$（$X$ 上）と、各射
$U \to U'$（$X$ 上の射であり、これは自動的にエタールである）に対する制限写像
$\rho^{U'}_U : \mathcal{F}(U') \to \mathcal{F}(U)$。
\end{enumerate}
これらのデータは、$\rho^U_U = \text{id}$ が恒等射 $U \to U$ の場合に成り立ち、
$\rho^{U'}_U \circ \rho^{U''}_{U'} = \rho^{U''}_U$ が、スキームの射
$U \to U' \to U''$ が $X$ 上エタールである場合に成り立つという条件に加え、
次の {\it 層公理}を満たさなければならない。
\begin{itemize}
\item[(*)] 任意のエタール被覆 $\{ \varphi_i : U_i \to U\}_{i \in I}$ に対し、図式
$$
\xymatrix{
\mathcal{F} (U) \ar[r] &
\Pi_{i \in I} \mathcal{F} (U_i) \ar@<1ex>[r] \ar@<-1ex>[r] &
\Pi_{i, j \in I} \mathcal{F} (U_i \times_U U_j)
}
$$
は集合（それぞれ、アーベル群、ベクトル空間など）の圏における等化子図式である。
\end{itemize}

\begin{remark}
\label{etale-cohomology-remark-i-is-j}
直前の主張では $i = j$ の場合を忘れないことが本質的である。これは一般に
（Zariski 位相の場合と異なり）きわめて非自明な条件である。実際、重要な被覆が
一つの要素しかもたないこともしばしばある。
\end{remark}

\noindent
恒等射はエタール射であるから、エタール層の大域切断を計算でき、コホモロジーは
単に対応する右導来関手となる。言い換えると、理論をさらに展開して主張を厳密に
述べさえすれば、コホモロジーを定義するうえで障害はない。




\section{エタール位相の威力}
\label{etale-cohomology-section-feats}

\noindent
自然数 $n \in \mathbf{N} = \{1, 2, 3, 4, \ldots\}$ に対して
$$
H_\etale^2 (\mathbf{P}^1_\mathbf{C}, \mathbf{Z}/n\mathbf{Z}) =
\mathbf{Z}/n\mathbf{Z}.
$$
が成り立つ。より一般に、$X$ が複素多様体ならば、有限体に係数をもつその
エタール Betti 数は $X(\mathbf{C})$ の通常の Betti 数と一致する。すなわち
$$
\dim_{\mathbf{F}_q} H_\etale^{2i} (X, \mathbf{F}_q) =
\dim_{\mathbf{F}_q} H_{Betti}^{2i} (X(\mathbf{C}), \mathbf{F}_q).
$$
である。これはきわめて満足のいく結果である。しかし、これらの等式が成り立つのは
ねじれ係数の場合に限られ、一般には成り立たない。整数係数の場合には
$$
H_\etale^2 (\mathbf{P}^1_\mathbf{C}, \mathbf{Z}) = 0.
$$
となる。これに対して、$H_{Betti}^2(\mathbf{P}^1(\mathbf{C}), \mathbf{Z}) = \mathbf{Z}$ である。
これは、位相空間 $\mathbf{P}^1(\mathbf{C})$ が $2$ 次元球面と同相だからである。
ねじれ係数から極限操作によって非ねじれ係数へ戻る方法があり、これについては
まもなく述べる。




\section{一つの計算}
\label{etale-cohomology-section-computation}

\noindent
$\mathbf{P}^1_\mathbf{C}$ の係数
$\Lambda = \mathbf{Z}/n\mathbf{Z}$ のコホモロジーをどのように計算するのだろうか。
{\v C}ech コホモロジーを用いる。$\mathbf{P}^1_\mathbf{C}$ の被覆は二つの標準開集合
$U_0, U_1$ によって与えられる。これらはいずれも
$\mathbf{A}^1_\mathbf{C}$ と同型であり、その共通部分は
$\mathbf{A}^1_\mathbf{C} \setminus \{0\} = \mathbf{G}_{m, \mathbf{C}}$
と同型である。エタール・コホモロジーでも Mayer--Vietoris 完全系列が成り立つ。
したがって、完全系列
$$
H_\etale^{i-1}(U_0\cap U_1, \Lambda) \to
H_\etale^i(\mathbf{P}^1_\mathbf{C}, \Lambda) \to
H_\etale^i(U_0, \Lambda) \oplus
H_\etale^i(U_1, \Lambda) \to H_\etale^i(U_0\cap U_1,
\Lambda).
$$
を得る。期待される答えを得るには、第三項に現れる直和が消滅することを示す必要がある。
実際、通常の位相の場合と同様に
$$
H_\etale^q (\mathbf{A}^1_\mathbf{C}, \Lambda) = 0
\quad \text{任意の } q \geq 1 \text{ に対して},
$$
であり、また
$$
H_\etale^q (\mathbf{A}^1_\mathbf{C} \setminus \{0\}, \Lambda) = \left\{
\begin{matrix}
\Lambda & \text{（}q = 1\text{ のとき）} \\
0 & \text{（}q \geq 2\text{ のとき）}.
\end{matrix}
\right.
$$
が成り立つ。これらの結果はすでにかなり難しい（初等的な証明はあるだろうか）。
以下では、エタール・コホモロジーの機構を利用できるものとして、この計算をどのように
行うかを説明する。

\medskip\noindent
{\bf 高次コホモロジー。} これは次の一般的な事実によって処理できる。すなわち、
$X$ が $\mathbf{C}$ 上のアフィン曲線ならば
$$
H_\etale^q (X, \mathbf{Z}/n\mathbf{Z}) = 0 \quad \text{任意の } q \geq 2 \text{ に対して}.
$$
これは曲線の一般点を考え、Galois コホモロジーを用いることによって証明される。
したがって、考えるべきものは次数 1 のコホモロジーだけである。

\medskip\noindent
{\bf 次数 1 のコホモロジー。} 次の同一視を用いる。
\begin{eqnarray*}
H_\etale^1 (X, \mathbf{Z}/n\mathbf{Z}) = \left\{
\begin{matrix}
\text{集合の層 }\mathcal{F}\text{ であって、エタール・サイト }X_\etale
\text{ 上で次を備えるもの} \\
\text{作用 }\mathbf{Z}/n\mathbf{Z} \times \mathcal{F} \to \mathcal{F}
\text{ であって }\mathcal{F}\text{ が }\mathbf{Z}/n\mathbf{Z}\text{-トーサーとなるもの}
\end{matrix}
\right\}
\Big/ \cong
\\
 = \left\{
\begin{matrix}
\text{有限エタール射 }Y \to X\text{ と、次を満たす自由な} \\
\mathbf{Z}/n\mathbf{Z}\text{ 作用を併せたもの：}
X = Y/(\mathbf{Z}/n\mathbf{Z}).
\end{matrix}
\right\}
\Big/ \cong.
\end{eqnarray*}
最初の同一視は非常に一般的であり（サイト上の任意のコホモロジー理論について成り立つ）、
エタール位相とは無関係である。第二の同一視は降下理論の帰結である。最後の集合は、
具体的に扱うことのできる幾何学的対象の集まりを記述している。

\medskip\noindent
曲線 $\mathbf{A}^1_\mathbf{C}$ には非自明な有限エタール被覆が存在しない。したがって
$H_\etale^1 (\mathbf{A}^1_\mathbf{C}, \mathbf{Z}/n\mathbf{Z}) = 0$.
これは位相的に見ることもできるし、次の段落の議論を用いて見ることもできる。

\medskip\noindent
有限エタール被覆
$\varphi : Y \to \mathbf{A}^1_\mathbf{C} \setminus \{0\}$.
を記述しよう。$Y$ が連結な場合を考えれば十分なので、以下そう仮定する。
Riemann--Hurwitz の公式を適用して $Y$ がどのようなものになり得るかを求める
（もちろん、これは少し大げさな方法であり、望むなら次の部分は読み飛ばしてよい）。
この射が $n$ 対 1 であるとし、射影的コンパクト化
$$
\xymatrix{
{Y\ } \ar@{^{(}->}[r] \ar[d]^\varphi &
{\bar Y} \ar[d]^{\bar\varphi} \\
{\mathbf{A}^1_\mathbf{C} \setminus \{0\}} \ar@{^{(}->}[r] &
{\mathbf{P}^1_\mathbf{C}}
}
$$
を考える。$\varphi$ はエタールで分岐しないが、$\bar{\varphi}$ は
0 と $\infty$ で分岐し得る。0 の逆像が点 $y_1,
\ldots, y_r$ であり、その分岐指数が $e_1, \ldots e_r$ であるとする。また、
$\infty$ の逆像が点 $y_1', \ldots, y_s'$ であり、その分岐指数が
$d_1, \ldots d_s$ であるとする。とくに $\sum e_i = n = \sum d_j$ である。
Riemann--Hurwitz の公式を適用すると
$$
2 g_Y - 2 = -2n + \sum (e_i - 1) + \sum (d_j - 1)
$$
を得る。したがって $g_Y = 0$、$r = s = 1$、かつ $e_1 = d_1 = n$ である。
ゆえに $Y \cong {\mathbf{A}^1_\mathbf{C} \setminus \{0\}}$ であり、
$\varphi(z) = \lambda z^n$ がある $\lambda \in \mathbf{C}^*$ について成り立つことが容易に分かる。
$Y$ のパラメータを取り直せば $\lambda = 1$ と仮定できる。したがって、この被覆は
$n$ 乗根を $\mathbf{A}^1_{\mathbf{C}} \setminus \{0\}$ 上の座標について取ることによって与えられる。

\medskip\noindent
分類すべきなのは ${\mathbf{A}^1_\mathbf{C} \setminus \{0\}}$ の被覆と、それに対する自由な
$\mathbf{Z}/n\mathbf{Z}$ 作用との組であったことを思い出そう。
この場合、そのような各作用は、$Y$ の自己同型であって $z$ を $\zeta_n z$ に送るものに
対応する。ここで $\zeta_n$ は原始 $n$ 乗根である。そのような作用は
$\phi(n)$ 個ある（ここで $\phi(n)$ は Euler 関数を表す）。したがって、連結有限
エタール被覆はちょうど $\phi(n)$ 個あり、各々が与えられた自由な
$\mathbf{Z}/n\mathbf{Z}$ 作用をもち、原始 $n$ 乗根に対応する。一方、次数 $n$ の非連結有限
エタール被覆（基底は $\mathbf{A}^1_{\mathbf{C}} \setminus \{0\}$）で、与えられた自由な
$\mathbf{Z}/n\mathbf{Z}$ 作用を備えるものは、非原始的な $n$ 乗根（$1$ の根）と
一対一に対応する。この対応の確認は読者に委ねる。言い換えると、この計算は
$$
H_\etale^1 (\mathbf{A}^1_\mathbf{C} \setminus \{0\},
\mathbf{Z}/n\mathbf{Z}) =
\Hom(\mu_n(\mathbf{C}), \mathbf{Z}/n\mathbf{Z}) \cong \mathbf{Z}/n\mathbf{Z}.
$$
最初の同一視は標準的であるが、第二の同一視はそうではない。
注 \ref{etale-cohomology-remark-normalize-H1-Gm} を参照されたい。
Riemann--Hurwitz の公式の証明はコホモロジーの計算を用いないので、消滅などの詳細を
補えば、以上は実際に証明を与えている。




\section{非ねじれ係数}
\label{etale-cohomology-section-nontorsion}

\noindent
非ねじれ係数を調べるために、次のように定義する。
$$
H_\etale^i (X, \mathbf{Q}_\ell) :=
\left( \lim_n H_\etale^i(X, \mathbf{Z}/\ell^n\mathbf{Z}) \right)
\otimes_{\mathbf{Z}_\ell} \mathbf{Q}_\ell.
$$
記号 $\lim_n$ は、コホモロジー群
$H_\etale^i(X, \mathbf{Z}/\ell^n\mathbf{Z})$ の系（添字は $n$）の
{\it 極限}を表す。Categories の第 \ref{categories-section-posets-limits} 節を参照されたい。
したがって、ある種の両立条件を満たす層の系を調べる必要がある。




\section{層理論}
\label{etale-cohomology-section-sheaf-theory}
%9.10.09

\noindent
ここからサイトと層について本格的に論じ始める。
続く数節に収め得る有用な抽象理論は驚くほど多い。その一部は、Sites、
Modules on Sites、Cohomology on Sites など、先行する章ですでに展開されている。
ここでは内容を増やしすぎず、本章後半の議論を理解するために必要な分だけを述べる。




\section{前層}
\label{etale-cohomology-section-presheaves}

\noindent
本節の参考文献は Sites の第 \ref{sites-section-presheaves} 節である。

\begin{definition}
\label{etale-cohomology-definition-presheaf}
$\mathcal{C}$ を圏とする。$\mathcal{C}$ 上の{\it 集合の前層}（それぞれ、
{\it アーベル群の前層}）とは、関手 $\mathcal{C}^{opp} \to
\textit{Sets}$（それぞれ $\textit{Ab}$）のことである。
\end{definition}

\noindent
{\bf 用語。} $U \in \Ob(\mathcal{C})$ のとき、$\mathcal{F}(U)$ の元を
$\mathcal{F}$ の $U$ 上の{\it 切断}と呼ぶ。$\varphi : V \to U$ を
$\mathcal{C}$ における射とすると、写像
$\mathcal{F}(\varphi) : \mathcal{F}(U) \to \mathcal{F}(V)$
を{\it 制限写像}と呼び、しばしば $s \mapsto s|_V$、またときには
$s \mapsto \varphi^*s$ と書く。記法 $s|_V$ は制限写像が
$\varphi$ に依存するため曖昧であるが、これは標準的な記法の濫用である。また、記法
$\Gamma(U, \mathcal{F}) = \mathcal{F}(U)$.
も用いる。

\medskip\noindent
$\mathcal{F}$ が関手であるということは、射 $W \to V \to U$ が
$\mathcal{C}$ に属し、$s \in \Gamma(U, \mathcal{F})$ ならば
$(s|_V)|_W = s |_W$ が成り立つということである。ここでは今述べた記法の濫用を用いた。
さらに、恒等射 $\text{id}_U : U \to U$ に対応する制限写像は恒等写像である。

\medskip\noindent
$\mathcal{C}$ 上の集合の前層（それぞれアーベル群の前層）の圏を
$\textit{PSh} (\mathcal{C})$（それぞれ $\textit{PAb}
(\mathcal{C})$）と書く。これは $\mathcal{C}^{opp}$ から
$\textit{Sets}$（それぞれ $\textit{Ab}$）への関手の圏である。言い換えると、
前層の射とは関手の自然変換である。圏 $\textit{PSh}(\mathcal{C})$ および
$\textit{PAb}(\mathcal{C})$ を考えるのは、圏 $\mathcal{C}$ が小さい場合に限る。
（別段の断りがない限り圏は小さいものとするのが本書の規約であり、小さくない圏は
Categories の注 \ref{categories-remark-big-categories} に列挙されているべきである。）

\begin{example}
\label{etale-cohomology-example-representable-presheaf}
対象 $X \in \Ob(\mathcal{C})$ が与えられたとき、次の対応を考える。
$$
\begin{matrix}
U & \longmapsto & h_X(U) = \Mor_\mathcal{C}(U, X) \\
V \xrightarrow{\varphi} U & \longmapsto &
(\psi \mapsto \psi \circ \varphi) : h_X(U) \to h_X(V).
\end{matrix}
$$
これは関手 $h_X : \mathcal{C}^{opp} \to \textit{Sets}$、したがって前層を定める。
これを {\it $X$ に付随する表現可能前層}と呼ぶ。
任意のサイト上の任意の位相について、表現可能前層が層になるわけではない。
\end{example}

\begin{lemma}[Yoneda]
\label{etale-cohomology-lemma-yoneda}
\begin{slogan}
対象間の射は、それらが表現する関手間の自然変換と一対一に対応する。
\end{slogan}
$\mathcal{C}$ を圏とし、$X, Y \in
\Ob(\mathcal{C})$ とする。このとき自然な全単射
$$
\begin{matrix}
\Mor_\mathcal{C}(X, Y) &
\longrightarrow &
\Mor_{\textit{PSh}(\mathcal{C})} (h_X, h_Y) \\
\psi &
\longmapsto &
h_\psi = \psi \circ - : h_X \to h_Y.
\end{matrix}
$$
が存在する。
\end{lemma}

\begin{proof}
Categories の補題 \ref{categories-lemma-yoneda} を参照されたい。
\end{proof}




\section{サイト}
\label{etale-cohomology-section-sites}


\begin{definition}
\label{etale-cohomology-definition-family-morphisms-fixed-target}
$\mathcal{C}$ を圏とする。{\it 固定された終域をもつ射の族}
$\mathcal{U} = \{\varphi_i : U_i \to U\}_{i\in I}$ とは、次のデータである。
\begin{enumerate}
\item 対象 $U \in \mathcal{C}$、
\item 集合 $I$（空でもよい）、および
\item 各 $i\in I$ に対する射 $\varphi_i : U_i \to U$ で、$\mathcal{C}$ に属し、終域が
$U$ であるもの。
\end{enumerate}
\end{definition}

\noindent
{\it 固定された終域をもつ射の族の射}という概念がある。その特別な場合が
{\it 細分}という概念である。この内容の参考文献は
Sites の第 \ref{sites-section-refinements} 節である。

\begin{definition}
\label{etale-cohomology-definition-site}
{\it サイト}\footnote{ここでサイトと呼ぶものは、
\cite[Expos\'e II, D\'efinition 1.3]{SGA4} では前位相を備えた圏と呼ばれる。
\cite{ArtinTopologies} では Grothendieck 位相をもつ圏と呼ばれる。}とは、圏
$\mathcal{C}$ と、{\it 被覆}と呼ばれる固定された終域をもつ射の族からなる集合
$\text{Cov}(\mathcal{C})$ であって、次を満たすものからなる。
\begin{enumerate}
\item （同型）$\varphi : V \to U$ が $\mathcal{C}$ における同型ならば、
$\{\varphi : V \to U\}$ は被覆である。
\item （局所性）$\{\varphi_i : U_i \to U\}_{i\in I}$ が被覆であり、各
$i \in I$ に対して被覆
$\{\psi_{ij} : U_{ij} \to U_i \}_{j\in I_i}$ が与えられているならば、
$$
\{
\varphi_i \circ \psi_{ij} : U_{ij} \to U
\}_{(i, j)\in \prod_{i\in I} \{i\} \times I_i}
$$
も被覆である。
\item （基底変換）$\{U_i \to U\}_{i\in I}$ が被覆であり、
$V \to U$ が $\mathcal{C}$ における射ならば、次が成り立つ。
\begin{enumerate}
\item 各 $i \in I$ に対し、ファイバー積 $U_i \times_U V$ が
$\mathcal{C}$ に存在する。
\item $\{U_i \times_U V \to V\}_{i\in I}$ は被覆である。
\end{enumerate}
\end{enumerate}
\end{definition}

\noindent
本章では、サイトの基礎にある圏は常に「小さい」、すなわちその対象全体が集合をなし、
サイトの被覆全体も（上の定義のように）集合をなすものとする。本章ではたいてい、
対象や被覆の集まりを集合へ縮小するための議論を省略する。さらに詳しくは
Sites の注 \ref{sites-remark-no-big-sites} を参照されたい。

\begin{example}
\label{etale-cohomology-example-site-topological-space}
$X$ が位相空間ならば、それに付随するサイト $X_{Zar}$ を次のように定める。
$X_{Zar}$ の対象は $X$ の開部分集合、これらの間の射は包含写像、被覆は通常の位相的な
（全射的な）被覆である。$U, V \subset W \subset X$ が開部分集合ならば
$U \times_W V = U \cap V$ が存在することに注意しよう。したがって、この圏はファイバー積を
もつ。すべての確認は自明であり、期待どおりに機能する。
\end{example}




\section{層}
\label{etale-cohomology-section-sheaves}

\begin{definition}
\label{etale-cohomology-definition-sheaf}
集合の前層 $\mathcal{F}$（それぞれアーベル群の前層）がサイト
$\mathcal{C}$ 上のものであるとする。任意の被覆
$\{\varphi_i : U_i \to U\}_{i\in I} \in \text{Cov} (\mathcal{C})$
に対して写像
$$
\mathcal{F}(U) \longrightarrow \prod\nolimits_{i\in I} \mathcal{F}(U_i)
$$
が単射であるとき、これを{\it 分離前層}という。ここで、この写像は
$s \mapsto (s|_{U_i})_{i\in I}$ である。前層 $\mathcal{F}$ が{\it 層}であるとは、
任意の被覆
$\{\varphi_i : U_i \to U\}_{i\in I} \in \text{Cov} (\mathcal{C})$
に対して図式
\begin{equation}
\label{etale-cohomology-equation-sheaf-axiom}
\xymatrix{
\mathcal{F}(U) \ar[r] &
\prod_{i\in I} \mathcal{F}(U_i) \ar@<1ex>[r] \ar@<-1ex>[r] &
\prod_{i, j \in I} \mathcal{F}(U_i \times_U U_j),
}
\end{equation}
が集合の圏（それぞれアーベル群の圏）における等化子図式であることをいう。
ここで最初の写像は $s \mapsto (s|_{U_i})_{i\in I}$ であり、右側の二つの写像は
$(s_i)_{i\in I} \mapsto (s_i |_{U_i \times_U U_j})$ および
$(s_i)_{i\in I} \mapsto (s_j |_{U_i \times_U U_j})$
である。
\end{definition}

\begin{remark}
\label{etale-cohomology-remark-empty-covering}
空被覆（$I = \emptyset$ の場合）について、これは
$\mathcal{F}(\emptyset)$ が空積、すなわち対応する圏の終対象であることを意味する
（$\textit{Sets}$ と $\textit{Ab}$ のいずれについても一元集合である）。
\end{remark}

\begin{example}
\label{etale-cohomology-example-sheaf-site-space}
位相空間に付随するサイト $X_{Zar}$ についてこの定義を書き下すと、通常の層の概念が得られる。
例 \ref{etale-cohomology-example-site-topological-space} を参照されたい。
\end{example}

\begin{definition}
\label{etale-cohomology-definition-category-sheaves}
$\Sh(\mathcal{C})$（それぞれ $\textit{Ab}(\mathcal{C})$）を、
$\textit{PSh}(\mathcal{C})$（それぞれ $\textit{PAb}(\mathcal{C})$）の、
対象が層であるような充満部分圏とする。これは $\mathcal{C}$ 上の
{\it 集合の層の圏}（それぞれ{\it アーベル群の層の圏}）である。
\end{definition}




\section{G-集合の例}
\label{etale-cohomology-section-G-sets}

\noindent
$G$ を群とし、サイト $\mathcal{T}_G$ を次のように定める。その基礎圏は
$G$-集合の圏、すなわち対象が左 $G$-作用を備えた集合で射が同変写像である圏とする。
$\mathcal{T}_G$ の被覆は、族 $\{\varphi_i : U_i \to U\}_{i\in I}$ であって
$U = \bigcup_{i\in I} \varphi_i(U_i)$ を満たすものとする。

\medskip\noindent
サイト $\mathcal{T}_G$ には特別な対象がある。それは左移動による自然な作用を備えた
$G$-集合 $G$ である。これを ${}_G G$ と書く。自然な群同型
$$
\begin{matrix}
\rho : & G^{opp} & \longrightarrow & \text{Aut}_{G\textit{-Sets}}({}_G G) \\
& g & \longmapsto & (h \mapsto hg).
\end{matrix}
$$
が存在することに注意しよう。とくに、任意の前層 $\mathcal{F}$ に対して、集合
$\mathcal{F}({}_G G)$ は $G$-作用を $\rho$ を介して受け継ぐ。
（$\mathcal{F}$ の反変性により、集合 $\mathcal{F}({}_G G)$ は再び左
$G$-集合になることに注意する。）実際、関手
$$
\begin{matrix}
\Sh(\mathcal{T}_G) & \longrightarrow & G\textit{-Sets} \\
\mathcal{F} & \longmapsto & \mathcal{F}({}_G G)
\end{matrix}
$$
は圏同値である。その準逆は関手 $X \mapsto
h_X$ である。完全な証明は Sites の第
\ref{sites-section-example-sheaf-G-sets} 節に譲り、ここではこれが成り立つ理由を説明する。
\begin{enumerate}
\item
$S$ が $G$-集合ならば、それを軌道の非交和 $S = \coprod_{i\in I}
O_i$ に分解できる。被覆 $\{O_i \to S\}_{i\in I}$ に対する層公理は、次の図式が
$$
\xymatrix{
\mathcal{F}(S) \ar[r] &
\prod_{i\in I} \mathcal{F}(O_i) \ar@<1ex>[r] \ar@<-1ex>[r] &
\prod_{i, j \in I} \mathcal{F}(O_i \times_S O_j)
}
$$
等化子であると述べている。$G\textit{-Sets}$ におけるファイバー積が
$\textit{Sets}$ におけるファイバー積から誘導されることに注意し、
$\mathcal{F}(\emptyset)$ が $G$-一元集合であることを用いると、
$$
\prod_{i, j \in I} \mathcal{F}(O_i \times_S O_j) = \prod_{i \in I}
\mathcal{F}(O_i)
$$
を得る。さらに、上の二つの写像は実際には同じである。したがって層公理は単に
$\mathcal{F}(S) = \prod_{i\in I} \mathcal{F}(O_i)$ と述べている。
\item
$S$ が $G$-集合 $S= G/H$ であり、$\mathcal{F}$ が $\mathcal{T}_G$ 上の層ならば、
次が成り立つと主張する。
$$
\mathcal{F}(G/H) = \mathcal{F}({}_G G)^H
$$
とくに $\mathcal{F}(\{*\}) = \mathcal{F}({}_G G)^G$ である。これを見るため、
被覆 $\{ {}_G G \to G/H \}$（$S$ の被覆）に層公理を用いる。このとき
\begin{eqnarray*}
{}_G G \times_{G/H} {}_G G & \cong & G \times H \\
(g_1, g_2) & \longmapsto & (g_1, g_1^{-1} g_2)
\end{eqnarray*}
は ${}_G G$ のコピーの非交和である（$G$-集合として）。したがって層公理は
$$
\xymatrix{
\mathcal{F} (G/H) \ar[r] &
\mathcal{F}({}_G G) \ar@<1ex>[r] \ar@<-1ex>[r] &
\prod_{h\in H} \mathcal{F}({}_G G)
}
$$
となる。ここで右側の二つの写像は $s \mapsto (s)_{h \in H}$ および $s \mapsto
(hs)_{h \in H}$ である。したがって、主張どおり
$\mathcal{F}(G/H) = \mathcal{F}({}_G G)^H$ である。
\end{enumerate}
これは主張した圏同値を完全に証明するものではないが、少なくとも層 $\mathcal{F}$ が
${}_G G$ 上の切断によって完全に決定されることを示している。詳細（および集合論上の注意）は
Sites の第 \ref{sites-section-example-sheaf-G-sets} 節にある。




\section{層化}
\label{etale-cohomology-section-sheafification}

\begin{definition}
\label{etale-cohomology-definition-0-cech}
$\mathcal{F}$ をサイト $\mathcal{C}$ 上の前層とし、
$\mathcal{U} = \{U_i \to U\} \in \text{Cov} (\mathcal{C})$ とする。
$\mathcal{F}$ の $\mathcal{U}$ に関する
{\it 第零 {\v C}ech コホモロジー群}を次で定める。
$$
\check H^0 (\mathcal{U}, \mathcal{F}) =
\left\{
(s_i)_{i\in I} \in \prod\nolimits_{i\in I }\mathcal{F}(U_i)
\text{ であって }
s_i|_{U_i \times_U U_j} = s_j |_{U_i \times_U U_j}
\right\}.
$$
\end{definition}

\noindent
次の標準的な写像が存在する：
$\mathcal{F}(U) \to \check H^0 (\mathcal{U}, \mathcal{F})$,
$s \mapsto (s |_{U_i})_{i\in I}$.
被覆 $\mathcal{V} = \{V_j \to V\}_{j \in J}$ から
$\mathcal{U}$ への{\it 被覆の射}とは、三つ組
$(\chi, \alpha, \chi_j)$ であって、$\chi : V \to U$ が射、
$\alpha : J \to I$ が集合の写像であり、各 $j \in J$ に対して射
$\chi_j$ が可換図式
$$
\xymatrix{
V_j \ar[rr]_{\chi_j} \ar[d] & & U_{\alpha(j)} \ar[d] \\
V \ar[rr]^\chi & & U.
}
$$
に入るものをいう。データ $\chi, \alpha, \{\chi_j\}_{j \in J}$ が与えられたとき、次の写像を定める：
\begin{eqnarray*}
\check H^0(\mathcal{U}, \mathcal{F}) & \longrightarrow &
\check H^0(\mathcal{V}, \mathcal{F}) \\
(s_i)_{i\in I} & \longmapsto &
\left(\chi_j^*\left(s_{\alpha(j)}\right)\right)_{j\in J}.
\end{eqnarray*}
このとき次を主張する。
\begin{enumerate}
\item この写像は定義良好である。
\item この写像は $\chi$ のみに依存し、
$\alpha, \{\chi_j\}_{j \in J}$ の選択には依存しない。
\end{enumerate}
第一の事実の証明は省略する。第二項を見るため、別の三つ組
$(\psi, \beta, \psi_j)$ で $\chi = \psi$ を満たすものを考える。このとき可換図式
$$
\xymatrix{
V_j \ar[rrr]_{(\chi_j, \psi_j)} \ar[dd] & & &
U_{\alpha(j)} \times_U U_{\beta(j)} \ar[dl] \ar[dr] \\
& & U_{\alpha(j)} \ar[dr] & &
U_{\beta(j)} \ar[dl] \\
V \ar[rrr]^{\chi = \psi} & & & U.
}
$$
を得る。切断 $s \in \mathcal{F}(\mathcal{U})$ が与えられたとする。その
$\mathcal{F}(V_j)$ での像は、三つ組
$(\chi, \alpha, \{\chi_j\}_{j \in J})$
が与える写像の下で
$\chi_j^*s_{\alpha(j)}$ であり、三つ組 $(\psi, \beta, \{\psi_j\}_{j \in J})$
が与える写像による像は $\psi_j^*s_{\beta(j)}$ である。仮定により
$s \in \check H^0(\mathcal{U}, \mathcal{F})$ なので、両者は等しい。実際、いずれも共通値
$$
s_{\alpha(j)}|_{U_{\alpha(j)} \times_U U_{\beta(j)}} =
s_{\beta(j)}|_{U_{\alpha(j)} \times_U U_{\beta(j)}}
$$
を図式中の写像 $(\chi_j, \psi_j)$ によって引き戻したものである。

\begin{theorem}
\label{etale-cohomology-theorem-sheafification}
$\mathcal{C}$ をサイトとし、$\mathcal{F}$ を $\mathcal{C}$ 上の前層とする。
\begin{enumerate}
\item 対応
$$
U \mapsto \mathcal{F}^+(U) :=
\colim_{\mathcal{U} \text{ が }U\text{ の被覆}}
\check H^0(\mathcal{U}, \mathcal{F})
$$
は前層を定める。また、この余極限は有向余極限である。
\item 前層の標準的な写像 $\mathcal{F} \to \mathcal{F}^+$ が存在する。
\item $\mathcal{F}$ が分離前層ならば、$\mathcal{F}^+$ は層であり、(2) の写像は単射である。
\item $\mathcal{F}^+$ は分離前層である。
\item $\mathcal{F}^\# = (\mathcal{F}^+)^+$ は層であり、標準的な写像は関手的同型
$$
\Hom_{\textit{PSh}(\mathcal{C})}(\mathcal{F}, \mathcal{G}) =
\Hom_{\Sh(\mathcal{C})}(\mathcal{F}^\#, \mathcal{G})
$$
を任意の $\mathcal{G} \in \Sh(\mathcal{C})$ に対して誘導する。
\end{enumerate}
\end{theorem}

\begin{proof}
Sites の定理 \ref{sites-theorem-plus} を参照されたい。
\end{proof}

\noindent
言い換えると、これは自然な写像
$\mathcal{F} \to \mathcal{F}^\#$ が忘却関手
$\Sh(\mathcal{C}) \to \textit{PSh}(\mathcal{C})$ の左随伴であることを意味する。




\section{コホモロジー}
\label{etale-cohomology-section-cohomology}

\noindent
次の基本結果により、サイト上のアーベル群の層に対してコホモロジーを定義できる。

\begin{theorem}
\label{etale-cohomology-theorem-enough-injectives}
サイト上のアーベル群の層の圏は、十分多くの入射対象をもつアーベル圏である。
\end{theorem}

\begin{proof}
Modules on Sites の補題 \ref{sites-modules-lemma-abelian-abelian} および
Injectives の定理 \ref{injectives-theorem-sheaves-injectives} を参照されたい。
\end{proof}

\noindent
したがって、コホモロジーを切断関手の右導来関手として定義できる。すなわち、
$U \in \Ob(\mathcal{C})$ かつ
$\mathcal{F} \in \textit{Ab}(\mathcal{C})$ ならば、
$$
H^p(U, \mathcal{F}) :=
R^p\Gamma(U, \mathcal{F}) =
H^p(\Gamma(U, \mathcal{I}^\bullet))
$$
と定める。ここで $\mathcal{F} \to \mathcal{I}^\bullet$ は入射分解である。この定義のためには、
関手 $\Gamma(U, -)$ が左完全であることを確認すべきである。これは実際に成り立ち、圏
$\textit{Ab}(\mathcal{C})$ がアーベル圏となる理由の一部でもある。Modules on Sites の補題
\ref{sites-modules-lemma-abelian-abelian} を参照されたい。大域切断関手とその右導来関手を含む、
サイト上のコホモロジーについてのより一般的な議論は、Cohomology on Sites の第
\ref{sites-cohomology-section-cohomology-sheaves} 節を参照されたい。



\section{fpqc トポロジー}
\label{etale-cohomology-section-fpqc}
%9.15.09

\noindent
エタール・コホモロジーに進む前に、準連接層に対してよく振る舞う fpqc トポロジーを
少し調べる。

\begin{definition}
\label{etale-cohomology-definition-fpqc-covering}
$T$ をスキームとする。$T$ の {\it fpqc 被覆}とは、次の条件を満たす族
$\{ \varphi_i : T_i \to T\}_{i \in I}$ のことである。
\begin{enumerate}
\item 各 $\varphi_i$ は平坦射であり、
$\bigcup_{i\in I} \varphi_i(T_i) = T$ である。
\item 各アフィン開集合 $U \subset T$ に対して、有限集合 $K$、写像
$\mathbf{i} : K \to I$、およびアフィン開集合
$U_{\mathbf{i}(k)} \subset T_{\mathbf{i}(k)}$ が存在して、
$U = \bigcup_{k \in K} \varphi_{\mathbf{i}(k)}(U_{\mathbf{i}(k)})$.
となる。
\end{enumerate}
\end{definition}

\begin{remark}
\label{etale-cohomology-remark-fpqc}
第一の条件が fp に対応する。これはフランス語で忠実平坦を意味する
{\it fid\`element plat} の略である。第二の条件は qc、すなわち
{\it quasi-compact} に対応する。第二の条件が成り立つならば、
第一の条件の後半は不要である。
\end{remark}

\begin{example}
\label{etale-cohomology-example-fpqc-coverings}
fpqc 被覆の例を挙げる。
\begin{enumerate}
\item $T$ の任意の Zariski 開被覆は fpqc 被覆である。
\item 族 $\{\Spec(B) \to \Spec(A)\}$ が fpqc 被覆であることと、
$A \to B$ が忠実平坦な環準同型であることは同値である。
\item $f: X \to Y$ が平坦、全射かつ準コンパクトならば、$\{ f: X\to
Y\}$ は fpqc 被覆である。
\item 射
$\varphi :
\coprod_{x \in \mathbf{A}^1_k} \Spec(\mathcal{O}_{\mathbf{A}^1_k, x})
\to \mathbf{A}^1_k$
を考える。ここで $k$ は体である。この射は平坦かつ全射である。しかし準コンパクトではなく、実際、族
$\{\varphi\}$ は fpqc 被覆ではない。
\item $\mathbf{A}^2_k = \Spec(k[x, y])$ と書く。標準開集合の射を
$i_x : D(x) \to \mathbf{A}^2_k$ および
$i_y : D(y) \to \mathbf{A}^2_k$ と記す。このとき族
$\{i_x, i_y, \Spec(k[[x, y]]) \to \mathbf{A}^2_k\}$
および
$\{i_x, i_y, \Spec(\mathcal{O}_{\mathbf{A}^2_k, 0}) \to \mathbf{A}^2_k\}$
は fpqc 被覆である。
\end{enumerate}
\end{example}

\begin{lemma}
\label{etale-cohomology-lemma-site-fpqc}
スキームの圏上の fpqc 被覆全体はサイトの公理を満たす。
\end{lemma}

\begin{proof}
Topologies の補題 \ref{topologies-lemma-fpqc} を参照されたい。
\end{proof}

\noindent
この補題によってスキームの圏の fpqc サイトを定義できそうに見える。しかし、
fpqc トポロジーを考える際には集合論的な問題が生じる。Topologies の第
\ref{topologies-section-fpqc} 節を参照されたい。これは、サイトが「小さい」ことを
要求している一方で、スキームからなるどの小圏も、所与の空でないスキームの
fpqc 被覆の余終系を含み得ないことから生じる。厳密に言えば、このことは「部分的な」
fpqc サイトを定義する妨げにはならないが、そうするのは賢明とは思われない。
回避策は、fpqc トポロジーに関する層という概念を認める一方で（以下を参照）、
すべての fpqc 層からなる圏を考えることを禁ずることである。

\begin{definition}
\label{etale-cohomology-definition-sheaf-property-fpqc}
$S$ をスキームとする。$S$ 上のスキームの圏を $\Sch/S$ と記す。
関手 $\mathcal{F} : (\Sch/S)^{opp} \to \textit{Sets}$、すなわち集合の前層を考える。
$\mathcal{F}$ は、任意の fpqc 被覆 $\{U_i \to U\}_{i \in I}$ で
$S$ 上のスキームの被覆であるものに対して
図式 (\ref{etale-cohomology-equation-sheaf-axiom}) が等化子図式となるとき、
{\it fpqc トポロジーに関する層条件を満たす}という。
\end{definition}

\noindent
同様に、$\mathcal{F}$ は、任意の開被覆 $U = \bigcup_{i \in I} U_i$ に対して図式
(\ref{etale-cohomology-equation-sheaf-axiom}) が等化子図式となるとき、
{\it Zariski トポロジーに関する層条件を満たす}という。Schemes の定義
\ref{schemes-definition-representable-by-open-immersions} を参照されたい。
これは明らかに、任意のスキーム $T$ で $S$ 上のものに対して、$\mathcal{F}$ を
$T$ の開集合に制限したものが（通常の）層であることと同値である。

\begin{lemma}
\label{etale-cohomology-lemma-fpqc-sheaves}
$\mathcal{F}$ を $\Sch/S$ 上の前層とする。このとき、$\mathcal{F}$ が
fpqc トポロジーに関する層条件を満たすための必要十分条件は次の二条件である。
\begin{enumerate}
\item $\mathcal{F}$ は Zariski トポロジーに関する層条件を満たす。
\item 任意の忠実平坦射 $\Spec(B) \to \Spec(A)$ で、
$S$ 上のアフィンスキームの射であるものに対して、被覆
$\{\Spec(B) \to \Spec(A)\}$ に層公理が成り立つ。すなわち、図式
$$
\xymatrix{
\mathcal{F}(\Spec(A)) \ar[r] &
\mathcal{F}(\Spec(B)) \ar@<1ex>[r] \ar@<-1ex>[r] &
\mathcal{F}(\Spec(B \otimes_A B))
}
$$
は等化子図式である。
\end{enumerate}
\end{lemma}

\begin{proof}
Topologies の補題 \ref{topologies-lemma-sheaf-property-fpqc} を参照されたい。
\end{proof}

\noindent
fpqc トポロジーに関する層条件を満たす前層 $\mathcal{F}$ で
$\Sch/S$ 上のものは、別の見方をすれば次のデータである。
\begin{enumerate}
\item 各 $T/S$ に対して、通常の（すなわち Zariski）層
$\mathcal{F}_T$ で $T_{Zar}$ 上のもの。
\item 任意の写像 $f : T' \to T$ で $S$ 上のものに対して、制限写像
$f^{-1}\mathcal{F}_T \to \mathcal{F}_{T'}$
\end{enumerate}
であって、次を満たすもの：
\begin{enumerate}
\item[(a)] 制限写像は関手的である。
\item[(b)] $f : T' \to T$ が開埋め込みならば、制限写像
$f^{-1}\mathcal{F}_T \to \mathcal{F}_{T'}$ は同型である。
\item[(c)] 任意の忠実平坦射
$\Spec(B) \to \Spec(A)$ で $S$ 上のものに対して、図式
$$
\xymatrix{
\mathcal{F}_{\Spec(A)}(\Spec(A)) \ar[r] &
\mathcal{F}_{\Spec(B)}(\Spec(B)) \ar@<1ex>[r] \ar@<-1ex>[r] &
\mathcal{F}_{\Spec(B \otimes_A B)}(\Spec(B \otimes_A B))
}
$$
は等化子である。
\end{enumerate}
データ (1)、(2) と条件 (a)、(b) は、Zariski トポロジーに関する層条件を満たす
$\Sch/S$ 上の前層のデータを与える。補題 \ref{etale-cohomology-lemma-fpqc-sheaves} により、
このとき条件 (c) が fpqc トポロジーに関する層条件を得るために十分である。

\begin{example}
\label{etale-cohomology-example-quasi-coherent}
前層
$$
\begin{matrix}
\mathcal{F} : & (\Sch/S)^{opp} & \longrightarrow & \textit{Ab} \\
& T/S & \longmapsto & \Gamma(T, \Omega_{T/S}).
\end{matrix}
$$
を考える。微分と局所化の両立性から、$\mathcal{F}$ は Zariski サイト上の層である。
しかし、これは fpqc トポロジーに関する層条件を満たさない。実際、
$S = \Spec(\mathbf{F}_p)$ の場合と、射
$$
\varphi :
V = \Spec(\mathbf{F}_p[v])
\to
U = \Spec(\mathbf{F}_p[u])
$$
を考える。この射は $u$ を $v^p$ に写すことで与えられる。族 $\{\varphi\}$ は fpqc 被覆であるが、制限写像
$\mathcal{F}(U) \to \mathcal{F}(V)$
は生成元 $\text{d}u$ を $\text{d}(v^p) = 0$ に写す。したがってこれは零写像であり、図式
$$
\xymatrix{
\mathcal{F}(U) \ar[r]^{0} &
\mathcal{F}(V) \ar@<1ex>[r] \ar@<-1ex>[r] &
\mathcal{F}(V \times_U V)
}
$$
は等化子ではない。後で、$\mathcal{F}$ が実際にはエタール・サイトおよび
滑らかなサイト上の層を与えることを見る。
\end{example}

\begin{lemma}
\label{etale-cohomology-lemma-representable-sheaf-fpqc}
$\Sch/S$ 上の任意の表現可能前層は fpqc トポロジーに関する層条件を満たす。
\end{lemma}

\begin{proof}
Descent の補題
\ref{descent-lemma-fpqc-universal-effective-epimorphisms} を参照されたい。
\end{proof}

\noindent
この事実の証明は、次節で論じる準連接層の降下を用いるので、後で再び取り上げる。
この補題を洗練された言い方で表すと、{\it fpqc トポロジーは標準トポロジーより弱い}、
あるいは fpqc トポロジーは{\it 準標準的}である、となる。サイトの枠組みにおける
この議論については、Sites の第 \ref{sites-section-representable-sheaves} 節を参照されたい。

\begin{remark}
\label{etale-cohomology-remark-fpqc-finest}
fpqc トポロジーは Zariski、エタール、滑らか、syntomic、および fppf
トポロジーのいずれよりも細かい。したがって、fpqc トポロジーに関する層条件を
満たす任意の前層は、Zariski、エタール、滑らか、syntomic、および fppf
サイト上の層となる。特に、例えば表現可能前層はスキームのエタール・サイト上の層となる。
\end{remark}

\begin{example}
\label{etale-cohomology-example-additive-group-sheaf}
$S$ をスキームとする。加法群スキーム
$\mathbf{G}_{a, S} = \mathbf{A}^1_S$ で $S$ 上のものを考える。Groupoids の例
\ref{groupoids-example-additive-group} を参照されたい。
対応する表現可能前層は
$$
h_{\mathbf{G}_{a, S}}(T) =
\Mor_S(T, \mathbf{G}_{a, S}) =
\Gamma(T, \mathcal{O}_T).
$$
によって与えられる。上の議論から、これは fpqc トポロジーに関する層条件を満たす
集合の前層である。一方、明らかに環の前層でもある。したがって、これを関手
$$
\begin{matrix}
\mathcal{O} : &
(\Sch/S)^{opp} &
\longrightarrow &
\textit{Rings} \\
&
T/S &
\longmapsto &
\Gamma(T, \mathcal{O}_T)
\end{matrix}
$$
であって fpqc トポロジーに関する層条件を満たすものと考えられる。
これに対応して $\mathcal{O}$-加群などの概念が得られる。
\end{example}




\section{忠実平坦降下}
\label{etale-cohomology-section-fpqc-descent}

\noindent
本節では準連接加群の忠実平坦降下を論じる。より正確には、準連接加群が
fpqc 被覆に関して有効降下を満たすことを証明する。

\begin{definition}
\label{etale-cohomology-definition-descent-datum}
$\mathcal{U} = \{ t_i : T_i \to T\}_{i \in I}$ を、終域を固定した
スキームの射の族とする。$\mathcal{U}$ に関する準連接層の
{\it 降下データ}とは、次を満たす族
$((\mathcal{F}_i)_{i \in I}, (\varphi_{ij})_{i, j \in I})$ のことである。
\begin{enumerate}
\item $\mathcal{F}_i$ は $T_i$ 上の準連接層である。
\item $\varphi_{ij} : \text{pr}_0^* \mathcal{F}_i \to
\text{pr}_1^* \mathcal{F}_j$ は $T_i \times_T T_j$ 上の加群の同型である。
\end{enumerate}
さらに{\it コサイクル条件}を満たす。すなわち、図式
$$
\xymatrix{
\text{pr}_0^*\mathcal{F}_i \ar[dr]_{\text{pr}_{02}^*\varphi_{ik}}
\ar[rr]^{\text{pr}_{01}^*\varphi_{ij}} & &
\text{pr}_1^*\mathcal{F}_j \ar[dl]^{\text{pr}_{12}^*\varphi_{jk}} \\
& \text{pr}_2^*\mathcal{F}_k
}
$$
は $T_i \times_T T_j \times_T T_k$ 上で可換である。この降下データが
{\it 有効}であるとは、準連接層 $\mathcal{F}$ で $T$ 上のものと
$\mathcal{O}_{T_i}$-加群の同型
$\varphi_i : t_i^* \mathcal{F} \cong \mathcal{F}_i$ が存在し、
写像 $\varphi_{ij}$ と両立すること、すなわち
$$
\varphi_{ij} = \text{pr}_1^* (\varphi_j) \circ \text{pr}_0^* (\varphi_i)^{-1}.
$$
が成り立つことをいう。
\end{definition}

\noindent
本節と次節では、次の定理の証明に用いるいくつかの材料と、それに関連する内容を論じる。

\begin{theorem}
\label{etale-cohomology-theorem-descent-quasi-coherent}
$\mathcal{V} = \{T_i \to T\}_{i\in I}$ が fpqc 被覆ならば、
$\mathcal{V}$ に関する準連接層のすべての降下データは有効である。
\end{theorem}

\begin{proof}
Descent の命題
\ref{descent-proposition-fpqc-descent-quasi-coherent} を参照されたい。
\end{proof}

\noindent
言い換えれば、準連接層のファイバー圏は fpqc サイト上のスタックである。
この定理の証明は二段階からなる。第一段階では、Zariski 被覆に対する主張が、
層の標準的な貼り合わせ（Sheaves の第
\ref{sheaves-section-glueing-sheaves} 節を参照）と準連接性の局所性から
容易に（または周知の結果として）従うことを確認する。第二段階は、
fpqc 被覆が $\{\Spec(B) \to \Spec(A)\}$ という形で、
$A \to B$ が忠実平坦な環準同型である場合である。これは代数の補題であり、以下に述べる。

\medskip\noindent
{\bf 加群の降下。}
$A \to B$ を環準同型とし、複体
$$
(B/A)_\bullet : B \to B \otimes_A B \to B \otimes_A B \otimes_A B \to \ldots
$$
を考える。ここで $B$ は次数 0、$B \otimes_A B$ は次数 1 などに置かれ、写像は
\begin{eqnarray*}
b & \mapsto & 1 \otimes b - b \otimes 1, \\
b_0 \otimes b_1 & \mapsto & 1 \otimes b_0 \otimes b_1 - b_0 \otimes 1 \otimes
b_1 + b_0 \otimes b_1 \otimes 1, \\
& \text{以下同様。}
\end{eqnarray*}
で与えられる。

\begin{lemma}
\label{etale-cohomology-lemma-algebra-descent}
$A \to B$ が忠実平坦ならば、複体 $(B/A)_\bullet$ は正の次数で完全であり、
$H^0((B/A)_\bullet) = A$ である。
\end{lemma}

\begin{proof}
Descent の補題 \ref{descent-lemma-ff-exact} を参照されたい。
\end{proof}

\noindent
Grothendieck はこれを三段階で証明する。第一に、写像 $A
\to B$ が切断をもつと仮定し、次数 0 の $A$ だけを非零項とする複体への
明示的なホモトピーを構成する。第二に、この結果を証明するには、忠実平坦な基底変換
$A \to A'$ の後で証明すれば十分であり、$B$ を $B' = B \otimes_A A'$ に
置き換えればよいことを確認する。第三に、忠実平坦な基底変換
$A \to A' = B$ を適用し、写像
$A' = B \to B' = B \otimes_A B$ が自然な切断をもつことを指摘する。

\medskip\noindent
同じ方針によって次の補題が証明される。

\begin{lemma}
\label{etale-cohomology-lemma-descent-modules}
$A \to B$ が忠実平坦で、$M$ が $A$-加群ならば、複体
$(B/A)_\bullet \otimes_A M$ は正の次数で完全であり、
$H^0((B/A)_\bullet \otimes_A M) = M$ である。
\end{lemma}

\begin{proof}
Descent の補題 \ref{descent-lemma-ff-exact} を参照されたい。
\end{proof}

\begin{definition}
\label{etale-cohomology-definition-descent-datum-modules}
$A \to B$ を環準同型、$N$ を $B$-加群とする。$N$ の
$A \to B$ に関する{\it 降下データ}とは、同型
$\varphi : N \otimes_A B \cong B \otimes_A N$ であって、
$B \otimes_A B$-加群の同型であり、
$B \otimes_A B \otimes_A B$-加群の図式
$$
\xymatrix{
{N \otimes_A B \otimes_A B} \ar[dr]_{\varphi_{02}} \ar[rr]^{\varphi_{01}} & &
{B \otimes_A N \otimes_A B} \ar[dl]^{\varphi_{12}} \\
& {B \otimes_A B \otimes_A N}
}
$$
が可換となるものをいう。ここで
$\varphi_{01} = \varphi \otimes \text{id}_B$ であり、
$\varphi_{12}$ と $\varphi_{02}$ も同様に定める。
\end{definition}

\noindent
$N' = B \otimes_A M$ がある $A$-加群 M によって与えられるならば、
これは写像
$$
\begin{matrix}
\varphi_\text{can}: & N' \otimes_A B & \to & B \otimes_A N' \\
& b_0 \otimes m \otimes b_1 & \mapsto & b_0 \otimes b_1 \otimes m.
\end{matrix}
$$
によって定まる標準的な降下データをもつ。

\begin{definition}
\label{etale-cohomology-definition-effective-modules}
降下データ $(N, \varphi)$ が{\it 有効}であるとは、ある
$A$-加群 $M$ が存在して $(N, \varphi) \cong (B \otimes_A M,
\varphi_\text{can})$ となることをいう。ここでは降下データの同型を明らかな意味で用いる。
\end{definition}

\noindent
定理 \ref{etale-cohomology-theorem-descent-quasi-coherent} は次の結果から従う。

\begin{theorem}
\label{etale-cohomology-theorem-descent-modules}
$A \to B$ が忠実平坦ならば、$A\to B$ に関する降下データは有効である。
\end{theorem}

\begin{proof}
Descent の命題 \ref{descent-proposition-descent-module} を参照されたい。
証明の別の見方については、Descent の注意
\ref{descent-remark-homotopy-equivalent-cosimplicial-algebras} も参照されたい。
\end{proof}

\begin{remarks}
\label{etale-cohomology-remarks-theorem-modules-exactness}
加群の降下に関する結果には、いくつかの応用がある。
\begin{enumerate}
\item 被覆 $\{\Spec(B) \to \Spec(A)\}$ で $A \to B$ が
忠実平坦であるものに対して、{\v C}ech 複体が正の次数で完全であること。
これはコホモロジーのいくつかの消滅を与える。
\item $(N, \varphi)$ が忠実平坦写像 $A \to B$ に関する降下データならば、
対応する $A$-加群は
$$
M = \Ker \left(
\begin{matrix}
N & \longrightarrow & B \otimes_A N \\
n & \longmapsto & 1 \otimes n - \varphi(n \otimes 1)
\end{matrix}
\right).
$$
で与えられる。Descent の命題 \ref{descent-proposition-descent-module} を参照されたい。
\end{enumerate}
\end{remarks}




%9.17.09
\section{準連接層}
\label{etale-cohomology-section-quasi-coherent}

\noindent
加群の降下を準連接層の研究に応用できる。

\begin{proposition}
\label{etale-cohomology-proposition-quasi-coherent-sheaf-fpqc}
任意の準連接層 $\mathcal{F}$ で $S$ 上のものに対して、前層
$$
\begin{matrix}
\mathcal{F}^a : & \Sch/S & \to & \textit{Ab}\\
& (f: T \to S) & \mapsto & \Gamma(T, f^*\mathcal{F})
\end{matrix}
$$
は fpqc トポロジーに関する層条件を満たす $\mathcal{O}$-加群である。
\end{proposition}

\begin{proof}
これは Descent の補題
\ref{descent-lemma-sheaf-condition-holds} で証明されている。
ここでは証明の概略を示す。補題 \ref{etale-cohomology-lemma-fpqc-sheaves} で示したように、
Zariski 被覆とアフィンスキームの忠実平坦射に対して層条件を確認すれば十分である。
Zariski 被覆に対する層条件は標準的なスキーム論である。実際、
$\Gamma(U, i^\ast \mathcal{F}) = \mathcal{F}(U)$ が、
$i : U \hookrightarrow S$ が開埋め込みならば成り立つ。

\medskip\noindent
被覆 $\left\{\Spec(B)\to \Spec(A)\right\}$ で $A\to B$ が忠実平坦であり、
$\mathcal{F}|_{\Spec(A)} = \widetilde{M}$
となる場合、これは
$M = H^0\left((B/A)_\bullet \otimes_A M \right)$、すなわち
\begin{align*}
0 \to M \to B \otimes_A M \to B \otimes_A B \otimes_A M
\end{align*}
が補題 \ref{etale-cohomology-lemma-descent-modules} により完全であるという事実に対応する。
\end{proof}

\noindent
環付きサイト上の準連接層には抽象的な概念がある。ここで簡単に導入する。
詳しくは Modules on Sites の第 \ref{sites-modules-section-local} 節を参照されたい。
$\mathcal{C}$ を圏とし、$U$ を $\mathcal{C}$ の対象とする。このとき
$\mathcal{C}/U$ は $U$ 上の対象の圏を表す。Categories の例
\ref{categories-example-category-over-X} を参照されたい。
$\mathcal{C}$ がサイトならば、$\mathcal{C}/U$ もサイトである。すなわち、
$V/U$ の被覆とは、族 $\{V_i/U \to V/U\}$ で
$\mathcal{C}/U$ の終域を固定した射からなり、$\{V_i \to V\}$ が
$\mathcal{C}$ の被覆となるものをいう。さらに、任意の層
$\mathcal{F}$ で $\mathcal{C}$ 上のものが与えられると、その{\it 制限}
$\mathcal{F}|_{\mathcal{C}/U}$（明らかな仕方で定義する）も層である。
詳しくは Sites の第 \ref{sites-section-localize} 節を参照されたい。

\begin{definition}
\label{etale-cohomology-definition-ringed-site}
$\mathcal{C}$ を{\it 環付きサイト}、すなわち環の層
$\mathcal{O}$ を備えたサイトとする。$\mathcal{O}$-加群の層
$\mathcal{F}$ で $\mathcal{C}$ 上のものが{\it 準連接}であるとは、任意の
$U \in \Ob(\mathcal{C})$ に対して被覆
$\{U_i \to U\}_{i\in I}$ で $\mathcal{C}$ の被覆であるものが存在し、制限
$\mathcal{F}|_{\mathcal{C}/U_i}$ が自由 $\mathcal{O}$-加群の間の
$\mathcal{O}$-線形写像
$$
\bigoplus\nolimits_{k \in K} \mathcal{O}|_{\mathcal{C}/U_i}
\longrightarrow
\bigoplus\nolimits_{l \in L} \mathcal{O}|_{\mathcal{C}/U_i}.
$$
の余核と同型であることをいう。$K$ にわたる直和は前層
$V \mapsto \bigoplus_{k \in K} \mathcal{O}(V)$ に付随する層であり、
もう一方についても同様である。
\end{definition}

\noindent
上のような一般的定義を与えられることは有用だが、この概念は一般にはよく振る舞わない。

\begin{remark}
\label{etale-cohomology-remark-final-object}
$\mathcal{C}$ が終対象、例えば $S$ をもつ場合、上の定義の条件は
$U = S$ に対して確認するだけで十分である。Modules on Sites の補題
\ref{sites-modules-lemma-local-final-object} を参照されたい。
\end{remark}

\begin{theorem}[準連接層に関するメタ定理]
\label{etale-cohomology-theorem-quasi-coherent}
$S$ をスキームとする。$\mathcal{C}$ をサイトとし、次を仮定する。
\begin{enumerate}
\item 基礎にある圏 $\mathcal{C}$ は $\Sch/S$ の充満部分圏である。
\item $T \in \Ob(\mathcal{C})$ の任意の Zariski 被覆は
$\mathcal{C}$ の被覆によって細分される。
\item $S/S$ は $\mathcal{C}$ の対象である。
\item $\mathcal{C}$ のすべての被覆はスキームの fpqc 被覆である。
\end{enumerate}
このとき前層 $\mathcal{O}$ は $\mathcal{C}$ 上の層であり、任意の準連接
$\mathcal{O}$-加群で $(\mathcal{C}, \mathcal{O})$ 上のものは、
$\mathcal{F}^a$ という形をもつ。ここで $\mathcal{F}$ は $S$ 上の準連接層である。
\end{theorem}

\begin{proof}
いくつかの形式的議論を行えば、これはまさに定理
\ref{etale-cohomology-theorem-descent-quasi-coherent} である。詳細は省略する。Descent の命題
\ref{descent-proposition-equivalence-quasi-coherent} では、
大 Zariski、fppf、エタール、滑らか、および syntomic サイトで $S$ のもの、
ならびに小 Zariski サイトと小エタール・サイトで $S$ のものについて、
この定理のより精密な形を証明する。
\end{proof}

\noindent
言い換えれば、スキーム $S$ 上の準連接加群と、準連接 $\mathcal{O}$-加群で
定理に現れるサイト $\mathcal{C}$ 上のものとの間に違いはない。
大および小サイト $(\Sch/S)_{fppf}$、$S_\etale$ などについての
より精密な主張は、Descent の各節
\ref{descent-section-quasi-coherent-sheaves},
\ref{descent-section-quasi-coherent-cohomology}, および
\ref{descent-section-quasi-coherent-sheaves-bis}.
にある。本章では、これらのどの状況でも成り立つ結果を簡潔に述べるために、
「定理 \ref{etale-cohomology-theorem-quasi-coherent} におけるサイト」と呼ぶことがある。






\section{{\v C}ech コホモロジー}
\label{etale-cohomology-section-cech-cohomology}

\noindent
次の目標は、降下理論を用いて等式
$H^i(\mathcal{C}, \mathcal{F}^a) = H_{Zar}^i(S, \mathcal{F})$
が、任意の準連接層 $\mathcal{F}$ で $S$ 上のものと、
定理 \ref{etale-cohomology-theorem-quasi-coherent} に現れる任意のサイト
$\mathcal{C}$ に対して成り立つことを示すことである。そのために、サイト上の
{\v C}ech コホモロジーを導入する。詳しくは \cite{ArtinTopologies} および
Cohomology on Sites の各節 \ref{sites-cohomology-section-cech},
\ref{sites-cohomology-section-cech-functor}
および \ref{sites-cohomology-section-cech-cohomology-cohomology}
を参照されたい。

\begin{definition}
\label{etale-cohomology-definition-cech-complex}
$\mathcal{C}$ を圏とする。$\mathcal{U} = \{U_i \to U\}_{i \in I}$ を、
$\mathcal{C}$ の終域を固定した射の族とする。すべてのファイバー積
$U_{i_0} \times_U \ldots \times_U U_{i_p}$ が $\mathcal{C}$ に存在すると仮定する。
$\mathcal{F} \in \textit{PAb}(\mathcal{C})$ をアーベル群の前層とする。
{\it {\v C}ech 複体} $\check{\mathcal{C}}^\bullet(\mathcal{U}, \mathcal{F})$ を
$$
\prod_{i_0 \in I} \mathcal{F}(U_{i_0}) \to
\prod_{i_0, i_1 \in I} \mathcal{F}(U_{i_0} \times_U U_{i_1}) \to
\prod_{i_0, i_1, i_2 \in I}
\mathcal{F}(U_{i_0} \times_U U_{i_1} \times_U U_{i_2}) \to \ldots
$$
によって定義する。ここで第一項は次数 0 に置かれ、写像は通常のものである。
{\it {\v C}ech コホモロジー群}を
$$
\check{H}^p(\mathcal{U}, \mathcal{F}) =
H^p(\check{\mathcal{C}}^\bullet(\mathcal{U}, \mathcal{F})).
$$
によって定義する。
\end{definition}

\noindent
上の定義では、添字 $(i_0, \ldots, i_p)$ に重複を許すことが本質的である。

\begin{lemma}
\label{etale-cohomology-lemma-cech-presheaves}
記法と仮定は定義 \ref{etale-cohomology-definition-cech-complex} のとおりとする。関手
$\check{\mathcal{C}}^\bullet(\mathcal{U}, -)$ は圏
$\textit{PAb}(\mathcal{C})$ 上で完全である。
\end{lemma}

\noindent
言い換えれば、$0\to \mathcal{F}_1\to \mathcal{F}_2\to \mathcal{F}_3\to 0$
がアーベル群の前層の短完全列ならば、
$$
0 \to \check{\mathcal{C}}^\bullet\left(\mathcal{U}, \mathcal{F}_1\right)
\to\check{\mathcal{C}}^\bullet(\mathcal{U}, \mathcal{F}_2) \to
\check{\mathcal{C}}^\bullet(\mathcal{U}, \mathcal{F}_3)\to 0
$$
は複体の短完全列である。

\begin{proof}
これは前層の短完全列の定義から直ちに従う。実際、アーベル群の前層の圏は、
ある圏から $\textit{Ab}$ への関手の圏なので、自動的にアーベル圏となる。
列 $\mathcal{F}_1\to \mathcal{F}_2\to \mathcal{F}_3$ が
$\textit{PAb}$ で完全であるための必要十分条件は、任意の
$U \in \Ob(\mathcal{C})$ に対して列
$\mathcal{F}_1(U) \to \mathcal{F}_2(U) \to \mathcal{F}_3(U)$ が
$\textit{Ab}$ で完全となることである。したがって上の複体は、各次数で短完全列の
積にすぎない。Cohomology on Sites の補題
\ref{sites-cohomology-lemma-cech-exact-presheaves} も参照されたい。
\end{proof}

\noindent
これにより $\check{H}^\bullet(\mathcal{U}, -)$ は $\delta$-関手である。
次に、これが普遍 $\delta$-関手であることを示す。そのためには、この関手が
{\it 消去可能}であることを示す必要がある。まず Yoneda の補題を思い出す。

\begin{lemma}[Yoneda の補題]
\label{etale-cohomology-lemma-yoneda-presheaf}
任意の前層 $\mathcal{F}$ が圏 $\mathcal{C}$ 上で与えられ、
$U \in \Ob(\mathcal{C})$ をとると、関手的同型
$$
\Hom_{\textit{PSh}(\mathcal{C})}(h_U, \mathcal{F}) =
\mathcal{F}(U).
$$
が存在する。
\end{lemma}

\begin{proof}
Categories の補題 \ref{categories-lemma-yoneda} を参照されたい。
\end{proof}

\noindent
集合 $E$ が与えられたとき、本節では $\mathbf{Z}[E]$ で
$E$ 上の自由アーベル群を表す。式で書けば
$\mathbf{Z}[E] = \bigoplus_{e \in E} \mathbf{Z}$ である。すなわち
$\mathbf{Z}[E]$ は自由 $\mathbf{Z}$-加群であり、その基底は
$E$ の元からなる。この記法を用いて、集合の前層上の
自由アーベル前層を導入する。

\begin{definition}
\label{etale-cohomology-definition-free-abelian-presheaf}
圏 $\mathcal{C}$ と集合の前層 $\mathcal{G}$ をとる。
{\it $\mathcal{G}$ 上の自由アーベル前層}を
$\mathbf{Z}_\mathcal{G}$ と書き、規則
$$
\mathbf{Z}_\mathcal{G}(U)
=
\mathbf{Z}[\mathcal{G}(U)]
$$
によって定義する。ここで $U \in \Ob(\mathcal{C})$ であり、
制限写像は $\mathcal{G}$ の制限写像から誘導される。
特別な場合 $\mathcal{G} = h_U$ には、単に
$\mathbf{Z}_U = \mathbf{Z}_{h_U}$.
と書く。
\end{definition}

\noindent
関手 $\mathcal{G} \mapsto \mathbf{Z}_\mathcal{G}$ は忘却関手
$\textit{PAb}(\mathcal{C}) \to \textit{PSh}(\mathcal{C})$ の左随伴である。
したがって任意の前層 $\mathcal{F}$ に対し、標準同型
$$
\Hom_{\textit{PAb}(\mathcal{C})}(\mathbf{Z}_U, \mathcal{F})
=
\Hom_{\textit{PSh}(\mathcal{C})}(h_U, \mathcal{F})
=
\mathcal{F}(U)
$$
があり、最後の等号は Yoneda の補題による。特に次の結果を得る。

\begin{lemma}
\label{etale-cohomology-lemma-cech-complex-describe}
記法と仮定は定義 \ref{etale-cohomology-definition-cech-complex} のとおりとする。
{\v C}ech 複体 $\check{\mathcal{C}}^\bullet(\mathcal{U}, \mathcal{F})$
は次のように明示的に記述できる：
\begin{eqnarray*}
\check{\mathcal{C}}^\bullet(\mathcal{U}, \mathcal{F})
& = &
\left(
\prod_{i_0 \in I}
\Hom_{\textit{PAb}(\mathcal{C})}(\mathbf{Z}_{U_{i_0}}, \mathcal{F}) \to
\prod_{i_0, i_1 \in I}
\Hom_{\textit{PAb}(\mathcal{C})}(
\mathbf{Z}_{U_{i_0} \times_U U_{i_1}}, \mathcal{F}) \to \ldots
\right) \\
& = &
\Hom_{\textit{PAb}(\mathcal{C})}\left(
\left(
\bigoplus_{i_0 \in I} \mathbf{Z}_{U_{i_0}} \leftarrow
\bigoplus_{i_0, i_1 \in I} \mathbf{Z}_{U_{i_0} \times_U U_{i_1}} \leftarrow
\ldots
\right), \mathcal{F}\right)
\end{eqnarray*}
\end{lemma}

\begin{proof}
これは上の式から従う。Cohomology on Sites の補題
\ref{sites-cohomology-lemma-cech-map-into} も参照されたい。
\end{proof}

\noindent
したがって、最後の $\Hom$ の第一引数にある複体だけを調べればよい。

\begin{lemma}
\label{etale-cohomology-lemma-exact}
記法と仮定は定義 \ref{etale-cohomology-definition-cech-complex} のとおりとする。
アーベル群の前層の複体
\begin{align*}
\mathbf{Z}_\mathcal{U}^\bullet \quad : \quad
\bigoplus_{i_0 \in I} \mathbf{Z}_{U_{i_0}} \leftarrow
\bigoplus_{i_0, i_1 \in I} \mathbf{Z}_{U_{i_0} \times_U U_{i_1}} \leftarrow
\bigoplus_{i_0, i_1, i_2 \in I}
\mathbf{Z}_{U_{i_0} \times_U U_{i_1} \times_U U_{i_2}} \leftarrow
\ldots
\end{align*}
は次数 $0$ を除くすべての次数で $\textit{PAb}(\mathcal{C})$ において完全である。
\end{lemma}

\begin{proof}
任意の $V\in \Ob(\mathcal{C})$ に対し、アーベル群の複体
$\mathbf{Z}_\mathcal{U}^\bullet(V)$ は
$$
\begin{matrix}
\mathbf{Z}\left[
\coprod_{i_0\in I} \Mor_\mathcal{C}(V, U_{i_0})\right]
\leftarrow
\mathbf{Z}\left[
\coprod_{i_0, i_1 \in I}
\Mor_\mathcal{C}(V, U_{i_0} \times_U U_{i_1})\right]
\leftarrow \ldots = \\
\bigoplus_{\varphi : V \to U}
\left(
\mathbf{Z}\left[\coprod_{i_0 \in I} \Mor_\varphi(V, U_{i_0})\right]
\leftarrow
\mathbf{Z}\left[\coprod_{i_0, i_1\in I} \Mor_\varphi(V, U_{i_0}) \times
\Mor_\varphi(V, U_{i_1})\right]
\leftarrow
\ldots
\right)
\end{matrix}
$$
である。ここで
$$
\Mor_{\varphi}(V, U_i)
=
\{ V \to U_i \text{ であって } V \to U_i \to U \text{ が } \varphi \text{ に等しいもの}\}.
$$
$S_\varphi = \coprod_{i\in I} \Mor_\varphi(V, U_i)$ とおく。このとき
$$
\mathbf{Z}_\mathcal{U}^\bullet(V)
=
\bigoplus_{\varphi : V \to U}
\left(
\mathbf{Z}[S_\varphi] \leftarrow
\mathbf{Z}[S_\varphi \times S_\varphi] \leftarrow
\mathbf{Z}[S_\varphi \times S_\varphi \times S_\varphi] \leftarrow
\ldots \right).
$$
したがって、各 $S = S_\varphi$ に対して複体
\begin{align*}
\mathbf{Z}[S] \leftarrow
\mathbf{Z}[S \times S] \leftarrow
\mathbf{Z}[S \times S \times S] \leftarrow \ldots
\end{align*}
が負の次数で完全であることを示せば十分である。このためには、明示的な
ホモトピーを与えればよい。$s\in S$ を固定し、
$K: n_{(s_0, \ldots, s_p)} \mapsto n_{(s, s_0,
\ldots, s_p)}.$ と定める。$K$ が作用素
$$
\delta :
\eta_{(s_0, \ldots, s_p)}
\mapsto
\sum\nolimits_{i = 0}^p (-1)^p \eta_{(s_0, \ldots, \hat s_i, \ldots, s_p)}.
$$
の零ホモトピーであることは容易に確かめられる。詳しくは
Cohomology on Sites の補題 \ref{sites-cohomology-lemma-homology-complex}
を参照されたい。
\end{proof}

\begin{lemma}
\label{etale-cohomology-lemma-hom-injective}
記法と仮定は定義 \ref{etale-cohomology-definition-cech-complex} のとおりとする。
$\mathcal{I}$ が $\textit{PAb}(\mathcal{C})$ の入射対象ならば、
$\check H^p(\mathcal{U}, \mathcal{I}) = 0$ がすべての $p > 0$ に対して成り立つ。
\end{lemma}

\begin{proof}
{\v C}ech 複体は、関手
$\Hom_{\textit{PAb}(\mathcal{C})}(-, \mathcal{I}) $ を複体 $
\mathbf{Z}^\bullet_\mathcal{U} $ に適用して得られる。すなわち、
$$
\check H^p(\mathcal{U}, \mathcal{I}) = H^p
(\Hom_{\textit{PAb}(\mathcal{C})} (\mathbf{Z}^\bullet_\mathcal{U},
\mathcal{I})).
$$
いま見たように $\mathbf{Z}^\bullet_\mathcal{U}$ は負の次数で完全であり、
関手 $\Hom_{\textit{PAb}(\mathcal{C})}(-,
\mathcal{I})$ も完全である。したがって $\Hom_{\textit{PAb}(\mathcal{C})}
(\mathbf{Z}^\bullet_\mathcal{U}, \mathcal{I})$ は正の次数で完全である。
\end{proof}

\begin{theorem}
\label{etale-cohomology-theorem-cech-derived}
記法と仮定は定義 \ref{etale-cohomology-definition-cech-complex} のとおりとする。
$\textit{PAb}(\mathcal{C})$ 上で、関手 $\check{H}^p(\mathcal{U}, -)$ は
$\check{H}^0(\mathcal{U}, -)$ の右導来関手である。
\end{theorem}

\begin{proof}
補題 \ref{etale-cohomology-lemma-hom-injective} により、関手
$\check H^p(\mathcal{U}, -)$ は消去可能なので普遍
$\delta$-関手である。$\check H^0(\mathcal{U}, -)$ の右導来関手も同様である。
次数 $0$ で一致するので、普遍 $\delta$-関手の普遍性により両者は一致する。
詳しくは Cohomology on Sites の補題
\ref{sites-cohomology-lemma-cech-cohomology-derived-presheaves} を参照されたい。
\end{proof}

\begin{remark}
\label{etale-cohomology-remark-presheaves-no-topology}
以上の主張はすべて前層に関するものであり、ここまではまだ位相を用いていないことに注意する。
\end{remark}




\section{{\v C}ech コホモロジーからコホモロジーへのスペクトル系列}
\label{etale-cohomology-section-cech-ss}

\noindent
このスペクトル系列は、層のコホモロジーに関する基礎的結果を証明するうえで
基本的な役割を果たす。

\begin{lemma}
\label{etale-cohomology-lemma-forget-injectives}
忘却関手 $\textit{Ab}(\mathcal{C})\to \textit{PAb}(\mathcal{C})$
は入射対象を入射対象に写す。
\end{lemma}

\begin{proof}
忘却関手が左随伴をもち、その左随伴である層化が完全関手であるという事実から
形式的に従う。詳しくは Cohomology on Sites の補題
\ref{sites-cohomology-lemma-injective-abelian-sheaf-injective-presheaf}
を参照されたい。
\end{proof}

\begin{theorem}
\label{etale-cohomology-theorem-cech-ss}
$\mathcal{C}$ をサイトとする。任意の被覆
$\mathcal{U} = \{U_i \to U\}_{i \in I}$ で $U \in \Ob(\mathcal{C})$
のものと、任意のアーベル群の層 $\mathcal{F}$ で $\mathcal{C}$
上のものに対して、スペクトル系列
$$
E_2^{p, q}
=
\check H^p(\mathcal{U}, \underline{H}^q(\mathcal{F}))
\Rightarrow
H^{p+q}(U, \mathcal{F}),
$$
が存在する。ここで $\underline{H}^q(\mathcal{F})$ はアーベル群の前層
$V \mapsto H^q(V, \mathcal{F})$ である。
\end{theorem}

\begin{proof}
入射分解 $\mathcal{F}\to \mathcal{I}^\bullet$ を
$\textit{Ab}(\mathcal{C})$ において選び、二重複体
$\check{\mathcal{C}}^\bullet(\mathcal{U}, \mathcal{I}^\bullet)$
と写像
$$
\xymatrix{
\Gamma(U, I^\bullet) \ar[r] &
\check{\mathcal{C}}^\bullet(\mathcal{U}, \mathcal{I}^\bullet) \\
& \check{\mathcal{C}}^\bullet(\mathcal{U}, \mathcal{F}) \ar[u]
}
$$
を考える。ここで水平写像は、左端の列への自然写像
$\Gamma(U, I^\bullet) \to
\check{\mathcal{C}}^0(\mathcal{U}, \mathcal{I}^\bullet)$
であり、垂直写像は $\mathcal{F}\to \mathcal{I}^0$ から誘導され、
最下段に入る。仮定により $\mathcal{I}^\bullet$ は
$\textit{Ab}(\mathcal{C})$ における入射対象の複体であるから、補題
\ref{etale-cohomology-lemma-forget-injectives} により
$\textit{PAb}(\mathcal{C})$ における入射対象の複体でもある。
したがって、二重複体の各行は正の次数で完全であり
（補題 \ref{etale-cohomology-lemma-hom-injective}）、写像
$\check{\mathcal{C}}^0(\mathcal{U}, \mathcal{I}^\bullet)
\to \check{\mathcal{C}}^1(\mathcal{U}, \mathcal{I}^\bullet)$
の核は $\Gamma(U, \mathcal{I}^\bullet)$ に等しい。実際、
$\mathcal{I}^\bullet$ は層の複体である。特に、全複体のコホモロジーは
大域切断関手 $H^0(U, \mathcal{F})$ の通常のコホモロジーである。

\medskip\noindent
垂直方向については、第 $q$ コホモロジー群で第 $p$ 列から得られるものは
$$
\prod_{i_0, \ldots, i_p}
H^q(U_{i_0} \times_U \ldots \times_U U_{i_p}, \mathcal{F})
=
\prod_{i_0, \ldots, i_p}
\underline{H}^q(\mathcal{F})(U_{i_0} \times_U \ldots \times_U U_{i_p})
$$
であり、これは項 $E_1^{p, q}$ に現れる。したがってこれは標準的な二重複体の
スペクトル系列であり、その $E_2$-ページは主張どおりである。詳しくは
Cohomology on Sites の補題
\ref{sites-cohomology-lemma-cech-spectral-sequence} を参照されたい。
\end{proof}

\begin{remark}
\label{etale-cohomology-remark-grothendieck-ss}
これは関手の合成
$$
\textit{Ab}(\mathcal{C}) \longrightarrow
\textit{PAb}(\mathcal{C}) \xrightarrow{\check H^0} \textit{Ab}.
$$
に対する Grothendieck スペクトル系列である。
\end{remark}








\section{スキームの大サイトと小サイト}
\label{etale-cohomology-section-big-small}

\noindent
$S$ をスキームとする。$\tau$ を、以下で扱う位相の一つとする。すなわち
\hfil\break $\tau \in \{fppf, syntomic, smooth, \etale, Zariski\}$ である。
もちろんエタール位相だけに関心があるなら、以下を通じて
$\tau = \etale$ と仮定すればよい。さらに、エタール射、エタール被覆、
およびエタール・サイトについては、節 \ref{etale-cohomology-section-etale-site} 以降で
より詳しく論じる。準連接層のコホモロジーを論じるためには、大
$\tau$-サイトと、$\tau \in \{\etale, Zariski\}$ の場合には
小 $\tau$-サイトで $S$ のものとを導入するのが便利である。そのため、まず
$\tau$-被覆の概念を導入する。

\begin{definition}
\label{etale-cohomology-definition-tau-covering}
(Topologies の各定義
\ref{topologies-definition-fppf-covering},
\ref{topologies-definition-syntomic-covering},
\ref{topologies-definition-smooth-covering},
\ref{topologies-definition-etale-covering}, および
\ref{topologies-definition-zariski-covering} を参照されたい。)
$\tau \in \{fppf, syntomic, smooth, \etale, Zariski\}$ とする。
終域を固定したスキームの射の族 $\{f_i : T_i \to T\}_{i \in I}$ が
{\it $\tau$-被覆}であるとは、各 $f_i$ がそれぞれ有限表示平坦、
syntomic、滑らか、エタール、または開埋め込みであり、かつ
$\bigcup f_i(T_i) = T$ が成り立つことをいう。
\end{definition}

\noindent
すべての $\tau$-被覆からなるクラスは、定義 \ref{etale-cohomology-definition-site}
（本章におけるサイトの定義）の公理 (1)、(2)、(3) を満たす。
Topologies の各補題
\ref{topologies-lemma-fppf},
\ref{topologies-lemma-syntomic},
\ref{topologies-lemma-smooth},
\ref{topologies-lemma-etale}, および
\ref{topologies-lemma-zariski} を参照されたい。

\medskip\noindent
以下で用いるサイトを導入しよう。\cite{SGA4} とは異なり、宇宙は選ばない。
その代わりに「部分宇宙」（Sets の節 \ref{sets-section-sets-hierarchy}
にいう十分大きな集合）を選び、その集合に含まれるすべてのスキームを考える。
もちろん、選んだ基礎スキーム $S$ が部分宇宙に含まれるようにする。
台となる圏を選んだ後、それをサイトにする十分大きな $\tau$-被覆の集合を選ぶ。
詳細はスキーム上の位相を扱う章にある。選択には大きな自由度があるが、最終的に
実際の選択は $S$ のエタール（または他の）コホモロジーに影響しない
（\cite{SGA4} でも、最終的には宇宙の実際の選択は問題にならない）。
さらに、本文は強非可到達基数、すなわち宇宙を用いてよいと考える読者が、
それを代わりに用いられるように書かれている。

\begin{definition}
\label{etale-cohomology-definition-tau-site}
$S$ をスキームとする。
$\tau \in \{fppf, syntomic, smooth, \etale, \linebreak[0] Zariski\}$ とする。
\begin{enumerate}
\item {\it 大 $\tau$-サイト（基礎スキーム $S$）}とは、上で説明したように、
また Topologies の各定義
\ref{topologies-definition-big-small-fppf},
\ref{topologies-definition-big-small-syntomic},
\ref{topologies-definition-big-small-smooth},
\ref{topologies-definition-big-small-etale}, および
\ref{topologies-definition-big-small-Zariski} でより詳しく構成される
サイト $(\Sch/S)_\tau$ のいずれかをいう。
\item $\tau \in \{\etale, Zariski\}$ のとき、{\it 小
$\tau$-サイト（基礎スキーム $S$）}とは、充満部分圏 $S_\tau$ で
$(\Sch/S)_\tau$ のものをいい、その対象はスキーム $T$ で $S$ 上にあり、構造射
$T \to S$ がそれぞれエタール射または開埋め込みであるものからなる圏をいう。
$S_\tau$ における被覆とは、被覆 $\{U_i \to U\}$ で
$(\Sch/S)_\tau$ におけるもののうち、$U$ が $S_\tau$ の対象であるものをいう。
\end{enumerate}
\end{definition}

\noindent
サイト $(\Sch/S)_\tau$ の台となる圏は、適切な「閉性」をもつ。
すなわち、その圏のスキーム $T$ をとると、$T$ の任意の局所閉部分スキームは
$(\Sch/S)_\tau$ のある対象と同型である。他の閉性として、スキームの
ファイバー積、可算非交和、および与えられたスキーム上の有限型スキームを
とる操作に関して閉じている。またアフィン・スキーム $\Spec(R)$ が与えられれば、
圏の内部にとどまったまま $R$ を完備化、局所化し、またはイデアルで割ることが
できる、などが挙げられる。他方、例えば $(\Sch/S)_\tau$ のスキームの
任意の非交和をとると、この圏の外に出ることがある。
また、対象 $T$ で $(\Sch/S)_\tau$ のものが与えられたとき、
$\tau$-被覆 $\{T_i \to T\}_{i \in I}$（定義
\ref{etale-cohomology-definition-tau-covering} の意味で）が存在しても、それが
$(\Sch/S)_\tau$ における被覆ではないことがある。例えばスキーム
$T_i$ が圏 $(\Sch/S)_\tau$ の対象でない場合である。しかし、サイト
$(\Sch/S)_\tau$ は、被覆 $\{U_j \to T\}_{j \in J}$ で
$(\Sch/S)_\tau$ におけるものが常に存在し、それが被覆
$\{T_i \to T\}_{i \in I}$ を細分するように選ばれている。Topologies の各補題
\ref{topologies-lemma-fppf-induced},
\ref{topologies-lemma-syntomic-induced},
\ref{topologies-lemma-smooth-induced},
\ref{topologies-lemma-etale-induced}, および
\ref{topologies-lemma-zariski-induced} を参照されたい。
本章では、これらの問題をほとんど意識しないことにする。

\medskip\noindent
$\mathcal{F}$ が $(\Sch/S)_\tau$ または $S_\tau$ 上の層ならば、
$$ 
H^p_\tau(U, \mathcal{F}), \text{ 特に }
H^p_\tau(S, \mathcal{F})
$$
で、$\mathcal{F}$ の、サイトの対象 $U$ 上におけるコホモロジー群を表す。
節 \ref{etale-cohomology-section-cohomology} を参照されたい。したがって
$H^p_{fppf}(S, \mathcal{F})$,
$H^p_{syntomic}(S, \mathcal{F})$,
$H^p_{smooth}(S, \mathcal{F})$,
$H^p_\etale(S, \mathcal{F})$, および
$H^p_{Zar}(S, \mathcal{F})$ が得られる。最後の二つは、大エタール・サイトと
小エタール・サイト、または大 Zariski サイトと小 Zariski サイトのいずれを
指すか曖昧になりうる。しかし、次の補題によりこの曖昧さは無害である。

\begin{lemma}
\label{etale-cohomology-lemma-compare-cohomology-big-small}
$\tau \in \{\etale, Zariski\}$ とする。
$\mathcal{F}$ を $(\Sch/S)_\tau$ 上で定義されたアーベル群の層とする。
$\mathcal{F}$ の $S$ 上のコホモロジー群は、
$\mathcal{F}|_{S_\tau}$ の $S$ 上のコホモロジー群と一致する。
\end{lemma}

\begin{proof}
Topologies の各補題 \ref{topologies-lemma-at-the-bottom} および
\ref{topologies-lemma-at-the-bottom-etale}
により、関手 $S_\tau \to (\Sch/S)_\tau$ は Sites の補題
\ref{sites-lemma-bigger-site} の仮定を満たす。したがって本補題は
Cohomology on Sites の補題
\ref{sites-cohomology-lemma-cohomology-bigger-site} から従う。
\end{proof}

\noindent
$S$ の大または小エタール・サイト上の層の圏は、
$(\Sch/S)_\etale$ または $S_\etale$ のアフィン対象からなる充満部分圏だけで
決まり、それらの間の標準エタール被覆（下で定義する）だけを考えればよい。
こうしてサイト $(\textit{Aff}/S)_\etale$ と $S_{affine, \etale}$ が得られる。
Topologies の定義 \ref{topologies-definition-big-small-etale} を参照されたい。
比較結果は Topologies の各補題
\ref{topologies-lemma-affine-big-site-etale} および
\ref{topologies-lemma-alternative} で証明されている。以下では、本章で扱う
いくつかの位相における標準被覆を定義する。

\begin{definition}
\label{etale-cohomology-definition-standard-tau}
(Topologies の各定義
\ref{topologies-definition-standard-fppf},
\ref{topologies-definition-standard-syntomic},
\ref{topologies-definition-standard-smooth},
\ref{topologies-definition-standard-etale}, および
\ref{topologies-definition-standard-Zariski} を参照されたい。)
$\tau \in \{fppf, syntomic, smooth, \etale, Zariski\}$ とする。
$T$ をアフィン・スキームとする。{\it 標準 $\tau$-被覆（対象 $T$）}とは、
族 $\{f_j : U_j \to T\}_{j = 1, \ldots, m}$ であって、各 $U_j$ がアフィンで、
各 $f_j$ がそれぞれ有限表示平坦、標準 syntomic、標準滑らか、エタール、
または $T$ の標準主開集合の埋め込みであり、かつ
$T = \bigcup f_j(U_j)$ が成り立つものをいう。
\end{definition}

\begin{lemma}
\label{etale-cohomology-lemma-tau-affine}
$\tau \in \{fppf, syntomic, smooth, \etale, Zariski\}$ とする。
アフィン・スキームの任意の $\tau$-被覆は、標準 $\tau$-被覆によって
細分できる。
\end{lemma}

\begin{proof}
Topologies の各補題
\ref{topologies-lemma-fppf-affine},
\ref{topologies-lemma-syntomic-affine},
\ref{topologies-lemma-smooth-affine},
\ref{topologies-lemma-etale-affine}, および
\ref{topologies-lemma-zariski-affine} を参照されたい。
\end{proof}

\noindent
完全を期すため、ここで扱っているコホモロジー群が部分宇宙の選択に依存しない
ことを述べ、証明する。小エタール・トポスの不変性は、後の補題
\ref{etale-cohomology-lemma-etale-topos-independent-partial-universe} で証明する。
この補題で用いる記法と用語については、Topologies の節
\ref{topologies-section-change-alpha} を参照されたい。

\begin{lemma}
\label{etale-cohomology-lemma-cohomology-enlarge-partial-universe}
$\tau \in \{fppf, syntomic, smooth, \etale, Zariski\}$ とする。
$S$ をスキームとする。$(\Sch/S)_\tau$ と $(\Sch'/S)_\tau$ を二つの
大 $\tau$-サイト（基礎スキーム $S$）とし、前者が後者に含まれると仮定する。
このとき
\begin{enumerate}
\item 任意のアーベル群の層 $\mathcal{F}'$ で $(\Sch'/S)_\tau$ 上で
定義されたものと、任意の対象 $U$ で $(\Sch/S)_\tau$ のものに対して
$$
H^p_\tau(U, \mathcal{F}'|_{(\Sch/S)_\tau}) =
H^p_\tau(U, \mathcal{F}')
$$
が成り立つ。言い換えれば、大きい方のサイトで計算した
$\mathcal{F}'$ の $U$ 上のコホモロジーは、小さい方のサイトに制限した
$\mathcal{F}'$ の $U$ 上のコホモロジーと一致する。
\item 任意のアーベル群の層 $\mathcal{F}$ で $(\Sch/S)_\tau$ 上のものに対し、
アーベル群の層 $\mathcal{F}'$ で $(\Sch/S)_\tau'$ 上のもののうち、その
$(\Sch/S)_\tau$ への制限が $\mathcal{F}$ と同型であるものが存在する。
\end{enumerate}
\end{lemma}

\begin{proof}
Topologies の補題 \ref{topologies-lemma-change-alpha} により、包含関手
$(\Sch/S)_\tau \to (\Sch'/S)_\tau$ は Sites の補題
\ref{sites-lemma-bigger-site} の仮定を満たす。これから (2) が従い、(1) は
Cohomology on Sites の補題
\ref{sites-cohomology-lemma-cohomology-bigger-site} から従う。
\end{proof}




\section{エタール・トポス}
\label{etale-cohomology-section-etale-topos}

\noindent
{\it トポス}とは、サイト上の集合の層の圏である。Sites の定義
\ref{sites-definition-topos} を参照されたい。したがって「スキームの
エタール・トポス」という語で、スキームの小エタール・サイト上の層の圏を
表すのが慣例である。形式的な定義は次のとおりである。

\begin{definition}
\label{etale-cohomology-definition-etale-topos}
$S$ をスキームとする。
\begin{enumerate}
\item $S$ の{\it エタール・トポス}、または{\it 小エタール・トポス}とは、
集合の層の圏 $\Sh(S_\etale)$ で、小エタール・サイトで $S$ のもの上の
層からなる圏をいう。
\item $S$ の{\it Zariski トポス}、または{\it 小 Zariski トポス}とは、
集合の層の圏 $\Sh(S_{Zar})$ で、小 Zariski サイトで $S$ のもの上の
層からなる圏をいう。
\item $\tau \in \{fppf, syntomic, smooth, \etale, Zariski\}$ に対し、
{\it 大 $\tau$-トポス}とは、大 $\tau$-トポスで $S$ のものの上の
集合の層の圏をいう。
\end{enumerate}
\end{definition}

\noindent
$S$ の小 Zariski トポスは、単に $S$ の台位相空間上の集合の層の圏である。
Topologies の補題 \ref{topologies-lemma-Zariski-usual} を参照されたい。
小エタール・トポスは小エタール・サイトの構成における選択に依存しないが、
一般に大トポスはそれらの選択に依存する。

\medskip\noindent
大または小エタール・トポスは、$(\Sch/S)_\etale$ または
$S_\etale$ のアフィン対象からなる充満部分圏だけで決まる。
Topologies の各補題 \ref{topologies-lemma-affine-big-site-etale} および
\ref{topologies-lemma-alternative} を参照されたい。例えば次の補題の証明で
このことを用いる。

\begin{lemma}
\label{etale-cohomology-lemma-etale-topos-independent-partial-universe}
$S$ をスキームとする。$S$ のエタール・トポスは、定義
\ref{etale-cohomology-definition-tau-site} における小エタール・サイトの構成に
（標準同値を除いて）依存しない。
\end{lemma}

\begin{proof}
二つの大エタール・サイト $\Sch_\etale$ と $\Sch_\etale'$ が
$S$ を含むとき、明らかな記法のもとで
$\Sh(S_\etale) \cong \Sh(S_\etale')$ となることを示せばよい。
Topologies の補題 \ref{topologies-lemma-contained-in} により
$\Sch_\etale \subset \Sch_\etale'$ と仮定してよい。Sets の補題
\ref{sets-lemma-what-is-in-it} により、$S$ 上エタールな任意の
アフィン・スキームは $\Sch_\etale$ と $\Sch_\etale'$ の双方の対象と
同型である。したがって誘導関手
$S_{affine, \etale} \to S_{affine, \etale}'$
は同値である。さらに、この関手とその擬逆写像がともに標準エタール被覆を
標準エタール被覆に写すことは明らかである。ゆえに結果は Topologies の補題
\ref{topologies-lemma-alternative} から従う。
\end{proof}





\section{準連接層のコホモロジー}
\label{etale-cohomology-section-cohomology-quasi-coherent}
%9.22.09

\noindent
まず簡単な補題から始める（実際には、述べたものより一般に成り立つ）。
これは、標準被覆の {\v C}ech 複体が、形
$\{\Spec(B) \to \Spec(A)\}$ で $A \to B$ が忠実平坦である fpqc 被覆の
{\v C}ech 複体に等しいことを述べる。

\begin{lemma}
\label{etale-cohomology-lemma-cech-complex}
位相 $\tau \in \{fppf, syntomic, smooth, \etale, Zariski\}$ をとる。
$S$ をスキームとする。$\mathcal{F}$ を $(\Sch/S)_\tau$ 上の
アーベル群の層、または $S_\tau$ 上のアーベル群の層
（$\tau = \etale$ の場合）とし、
$\mathcal{U} = \{U_i \to U\}_{i \in I}$
をこのサイトの標準 $\tau$-被覆とする。
$V = \coprod_{i \in I} U_i$ とおく。このとき
\begin{enumerate}
\item $V$ はアフィン・スキームである。
\item $\mathcal{V} = \{V \to U\}$ は fpqc 被覆であり、
さらに $\tau$-被覆でもある。ただし $\tau = Zariski$ の場合を除く。
\item 二つの {\v C}ech 複体
$\check{\mathcal{C}}^\bullet (\mathcal{U}, \mathcal{F})$ と
$\check{\mathcal{C}}^\bullet (\mathcal{V}, \mathcal{F})$ は一致する。
\end{enumerate}
\end{lemma}

\begin{proof}
標準 $\tau$-被覆の定義は Topologies の各定義
\ref{topologies-definition-standard-Zariski},
\ref{topologies-definition-standard-etale},
\ref{topologies-definition-standard-smooth},
\ref{topologies-definition-standard-syntomic}, および
\ref{topologies-definition-standard-fppf} にある。
定義により各スキーム $U_i$ はアフィンで、$I$ は有限集合である。
したがって $V$ はアフィン・スキームである。$V \to U$ が平坦かつ全射で
あることは明らかなので、$\mathcal{V}$ は fpqc 被覆である。例
\ref{etale-cohomology-example-fpqc-coverings} を参照されたい。Zariski の場合を除けば、
被覆 $\mathcal{V}$ は $\tau$-被覆でもある。Topologies の各定義
\ref{topologies-definition-etale-covering},
\ref{topologies-definition-smooth-covering},
\ref{topologies-definition-syntomic-covering}, および
\ref{topologies-definition-fppf-covering} を参照されたい。

\medskip\noindent
被覆 $\mathcal{U}$ は $\mathcal{V}$ の細分であることに注意する。
したがって {\v C}ech 複体の写像
$\check{\mathcal{C}}^\bullet (\mathcal{V}, \mathcal{F}) \to
\check{\mathcal{C}}^\bullet (\mathcal{U}, \mathcal{F})$ が存在する。
Cohomology on Sites の式
(\ref{sites-cohomology-equation-map-cech-complexes}) を参照されたい。
次に、$T = \coprod_{j \in J} T_j$ が、$\mathcal{F}$ の定義されたサイト内の
スキームの非交和ならば、終域を固定した射の族
$\{T_j \to T\}_{j \in J}$ は Zariski 被覆である。したがって
\begin{equation}
\label{etale-cohomology-equation-sheaf-coprod}
\mathcal{F}(T) =
\mathcal{F}(\coprod\nolimits_{j \in J} T_j) =
\prod\nolimits_{j \in J} \mathcal{F}(T_j)
\end{equation}
が成り立つ。これは $\mathcal{F}$ の層条件による。ゆえに上の
{\v C}ech 複体の写像は各次数で同型である。実際、
$$
V \times_U \ldots \times_U V
=
\coprod\nolimits_{i_0, \ldots i_p} U_{i_0} \times_U \ldots \times_U U_{i_p}
$$
がスキームとして成り立つ。
\end{proof}

\noindent
等式 (\ref{etale-cohomology-equation-sheaf-coprod}) は一般の前層に対しては偽であることに
注意する。層に対してさえ任意のサイトで成り立つわけではない。余積が被覆を
与えるとは限らず、また非交とは限らないからである。しかし、通常用いるサイト
（少なくとも本章で調べるすべてのサイト）では成り立つ。

\begin{remark}
\label{etale-cohomology-remark-refinement}
補題 \ref{etale-cohomology-lemma-cech-complex} の主張では、被覆 $\mathcal{U}$ は
$\mathcal{V}$ の細分であるが、逆は成り立たない。形 $\{V \to U\}$ の被覆は、
$U$ のすべての被覆からなる圏の始部分圏をなさない。それでも、{\v C}ech
コホモロジー $\check H^n(U, \mathcal{F})$（これは、被覆
$\mathcal{U}$ で $U$ のものからなる圏の反対圏上で、$\mathcal{F}$ の
$\mathcal{U}$ に関する {\v C}ech コホモロジー群の余極限として定義される）は、
被覆 $\{V \to U\}$ を用いて計算できる。必要になれば、精密な補題
（層に対してのみ成り立つ）を定式化してここに加えることにする。
\end{remark}

\begin{lemma}[コホモロジーの局所性]
\label{etale-cohomology-lemma-locality-cohomology}
$\mathcal{C}$ をサイト、$\mathcal{F}$ を $\mathcal{C}$ 上のアーベル群の層、
$U$ を $\mathcal{C}$ の対象、$p > 0$ を整数とし、$\xi \in
H^p(U, \mathcal{F})$ とする。このとき、被覆
$\mathcal{U} = \{U_i \to U\}_{i \in I}$ で、$U$ のもので
$\mathcal{C}$ におけるもののうち、$\xi |_{U_i} = 0$ がすべての
$i \in I$ に対して成り立つものが存在する。
\end{lemma}

\begin{proof}
入射分解 $\mathcal{F} \to \mathcal{I}^\bullet$ を選ぶ。このとき
$\xi$ はコサイクル $\tilde{\xi} \in \mathcal{I}^p(U)$ で
$d^p(\tilde{\xi}) = 0$ を満たすものによって表される。仮定により、列
$\mathcal{I}^{p - 1} \to \mathcal{I}^p \to \mathcal{I}^{p + 1}$ は
$\textit{Ab}(\mathcal{C})$ で完全である。これは、被覆
$\mathcal{U} = \{U_i \to U\}_{i \in I}$ で、
$\tilde{\xi}|_{U_i} = d^{p - 1}(\xi_i)$ が、ある
$\xi_i \in \mathcal{I}^{p-1}(U_i)$ に対して成り立つものが存在することを意味する。
コホモロジー類 $\xi|_{U_i}$ はコサイクル $\tilde{\xi}|_{U_i}$ によって
表され、後者はコバウンダリなので、この類は消える。詳しくは
Cohomology on Sites の補題
\ref{sites-cohomology-lemma-kill-cohomology-class-on-covering} を参照されたい。
\end{proof}

\begin{theorem}
\label{etale-cohomology-theorem-zariski-fpqc-quasi-coherent}
$S$ をスキームとし、$\mathcal{F}$ を準連接 $\mathcal{O}_S$-加群とする。
$\mathcal{C}$ を $(\Sch/S)_\tau$ で
$\tau \in \{fppf, syntomic, smooth, \etale, Zariski\}$ に対応するもの、
または $S_\etale$ のいずれかとする。このとき
$$
H^p(S, \mathcal{F}) = H^p_\tau(S, \mathcal{F}^a)
$$
がすべての $p \geq 0$ に対して成り立つ。ここで
\begin{enumerate}
\item 左辺は、層 $\mathcal{F}$ の、スキーム $S$ の台位相空間上における
通常のコホモロジーを表す。
\item 右辺は、アーベル群の層 $\mathcal{F}^a$（命題
\ref{etale-cohomology-proposition-quasi-coherent-sheaf-fpqc} を参照）のサイト
$\mathcal{C}$ 上におけるコホモロジーを表す。
\end{enumerate}
\end{theorem}

\begin{proof}
次を示す：
$H^p(U, f^*\mathcal{F}) = H^p_\tau(U, \mathcal{F}^a)$
これは、任意の対象 $f : U \to S$ でサイト $\mathcal{C}$ のものに対して成り立つ。
層の性質により、$p = 0$ の場合には主張は正しい。

\medskip\noindent
$U$ がアフィンであると仮定する。このとき
$H^p_\tau(U, \mathcal{F}^a) = 0$ がすべての $p > 0$ に対して
成り立つことを示したい。$p$ に関する帰納法を用いる。
\begin{enumerate}
\item[p = 1]
$\xi \in H^1_\tau(U, \mathcal{F}^a)$ をとる。補題
\ref{etale-cohomology-lemma-locality-cohomology} により、fpqc 被覆
$\mathcal{U} = \{U_i \to U\}_{i \in I}$ で、
$\xi|_{U_i} = 0$ がすべての $i \in I$ に対して成り立つものが存在する。
$\mathcal{U}$ を細分することにより、
$\mathcal{U}$ が標準 $\tau$-被覆であると仮定してよい。定理
\ref{etale-cohomology-theorem-cech-ss} のスペクトル系列を適用すると、$\xi$ は
コホモロジー類 $\check \xi \in \check H^1(\mathcal{U}, \mathcal{F}^a)$
から来ることが分かる。被覆
$\mathcal{V} = \{\coprod_{i\in I} U_i \to U\}$ を考える。補題
\ref{etale-cohomology-lemma-cech-complex} により、
$\check H^\bullet(\mathcal{U}, \mathcal{F}^a) =
\check H^\bullet(\mathcal{V}, \mathcal{F}^a)$.
他方、$\mathcal{V}$ は形 $\{\Spec(B) \to \Spec(A)\}$ の被覆であり、
$f^*\mathcal{F} = \widetilde{M}$ がある $A$-加群 $M$ に対して
成り立つので、{\v C}ech 複体
$\check{\mathcal{C}}^\bullet(\mathcal{V}, \mathcal{F})$
は複体 $(B/A)_\bullet \otimes_A M$ にほかならない。補題
\ref{etale-cohomology-lemma-descent-modules} により
$H^p((B/A)_\bullet \otimes_A M) = 0$ が $p > 0$ に対して成り立つ。
したがって $\check \xi = 0$ であり、ゆえに $\xi = 0$ である。
\item[p > 1]
$\xi \in H^p_\tau(U, \mathcal{F}^a)$ をとる。補題
\ref{etale-cohomology-lemma-locality-cohomology} により、fpqc 被覆
$\mathcal{U} = \{U_i \to U\}_{i \in I}$ で、
$\xi|_{U_i} = 0$ がすべての $i \in I$ に対して成り立つものが存在する。
$\mathcal{U}$ を細分することにより、$\mathcal{U}$ が標準
$\tau$-被覆であると仮定してよい。定理
\ref{etale-cohomology-theorem-cech-ss} のスペクトル系列を適用する。
交わり $U_{i_0} \times_U \ldots \times_U U_{i_p}$ はアフィンであるから、
帰納法の仮定によりコホモロジー群
$$
E_2^{p, q} = \check H^p(\mathcal{U}, \underline{H}^q(\mathcal{F}^a))
$$
はすべての $0 < q < p$ に対して消える。したがって $\xi$ は
$\check \xi \in \check H^p(\mathcal{U}, \mathcal{F}^a)$ から来なければならない。
$\mathcal{U}$ を、ただ一つの射を含む被覆 $\mathcal{V}$ で置き換え、再び補題
\ref{etale-cohomology-lemma-descent-modules} を用いると、{\v C}ech コホモロジー類
$\check \xi$ は零でなければならず、したがって $\xi = 0$ であることが分かる。
\end{enumerate}
次に、$U$ が分離的であると仮定する。アフィン開被覆
$U = \bigcup_{i \in I} U_i$ を $U$ に対して選ぶ。このとき族
$\mathcal{U} = \{U_i \to U\}_{i \in I}$ は fpqc 被覆である。
さらに、すべての交わり
$U_{i_0} \times_U \ldots \times_U U_{i_p}$ は、$U$ が分離的なのでアフィンである。したがって、定理
\ref{etale-cohomology-theorem-cech-ss} のスペクトル系列では、第零行を除くすべての行が零である。ゆえに
$$
H^p_\tau(U, \mathcal{F}^a) =
\check H^p(\mathcal{U}, \mathcal{F}^a) =
\check H^p(\mathcal{U}, \mathcal{F}) = H^p(U, \mathcal{F})
$$
である。最後の等号は標準的なスキーム論から従う。Cohomology of Schemes の補題
\ref{coherent-lemma-cech-cohomology-quasi-coherent} を参照されたい。

\medskip\noindent
一般の場合は技術的であり、ここに与えた証明を拡張するにはスペクトル系列の射を
論じる必要があるので、ここでは扱わない。主張は Descent の命題
\ref{descent-proposition-same-cohomology-quasi-coherent}
（その証明では少し異なる方法をとる）と Cohomology on Sites の補題
\ref{sites-cohomology-lemma-cohomology-of-open} を組み合わせることから従う。
\end{proof}

\begin{remark}
\label{etale-cohomology-remark-right-derived-global-sections}
定理 \ref{etale-cohomology-theorem-zariski-fpqc-quasi-coherent} について補足する。
$S$ は圏 $\mathcal{C}$ の終対象なので、右辺のコホモロジー群は単に
大域切断関手の右導来関手である。実際、証明は
$H^p(U, f^*\mathcal{F}) = H^p_\tau(U, \mathcal{F}^a)$
が、任意の対象 $f : U \to S$ でサイト $\mathcal{C}$ のものに対して成り立つことを示している。
\end{remark}





\section{層の例}
\label{etale-cohomology-section-examples-sheaves}

\noindent
$S$ と $\tau$ は節 \ref{etale-cohomology-section-big-small} と同様のものとする。
任意の表現可能な前層が $(\Sch/S)_\tau$ または $S_\tau$ 上の層であることは
すでに見た。補題 \ref{etale-cohomology-lemma-representable-sheaf-fpqc} および注意
\ref{etale-cohomology-remark-fpqc-finest} を参照されたい。以下にいくつかの特別な場合を挙げる。

\begin{definition}
\label{etale-cohomology-definition-additive-sheaf}
節 \ref{etale-cohomology-section-big-small} のサイト $(\Sch/S)_\tau$ または $S_\tau$ の
いずれにおいても、次のように定める。
\begin{enumerate}
\item 層 $T \mapsto \Gamma(T, \mathcal{O}_T)$ を $\mathcal{O}_S$ または
$\mathbf{G}_a$ と書き、基礎スキームを明示したいときには
$\mathbf{G}_{a, S}$ とも書く。
\item 同様に、層 $T \mapsto \Gamma(T, \mathcal{O}^*_T)$ を
$\mathcal{O}_S^*$ または $\mathbf{G}_m$ と書き、基礎スキームを明示したいときには
$\mathbf{G}_{m, S}$ とも書く。
\item 任意のサイト上の{\it 定数層} $\underline{\mathbf{Z}/n\mathbf{Z}}$ とは、
定数前層 $U \mapsto \mathbf{Z}/n\mathbf{Z}$ の層化である。
\end{enumerate}
\end{definition}

\noindent
例えば第一のものが層であることは定理 \ref{etale-cohomology-theorem-quasi-coherent} から従う。
第二のものは第一のものの部分前層であり、それ自身が層であることは容易に分かる。
第三のものは定義により層である。これらの層はいずれも表現可能であることに注意する。
第一と第二は、それぞれスキーム $\mathbf{G}_{a, S}$ と
$\mathbf{G}_{m, S}$ によって表現される。Groupoids の節
\ref{groupoids-section-group-schemes} を参照されたい。第三は有限エタール群スキーム
$\mathbf{Z}/n\mathbf{Z}_S$ によって表現され、これは
$(\mathbf{Z}/n\mathbf{Z})_S$ と書かれることもある。このスキームは、$n$ 個の
$S$ のコピーに $S$ 上の明らかな群スキーム構造を備えたものにすぎない。Groupoids の例
\ref{groupoids-example-constant-group} および次の注意を参照されたい。

\begin{remark}
\label{etale-cohomology-remark-constant-locally-constant-maps}
$G$ を抽象群とする。節 \ref{etale-cohomology-section-big-small} のサイト
$(\Sch/S)_\tau$ または $S_\tau$ のいずれにおいても、層化
$\underline{G}$、すなわち $G$ に付随する定数前層をそのサイトの
{\it Zariski 位相}について層化したものは、すでに
$$
\Gamma(U, \underline{G}) =
\{\text{Zariski 局所定数写像 }U \to G\}
$$
を与える。この Zariski 層は、Groupoids の例
\ref{groupoids-example-constant-group} により群スキーム $G_S$ で表現される。
補題 \ref{etale-cohomology-lemma-representable-sheaf-fpqc} により、任意の表現可能な前層は
$\tau$-位相についても層条件を満たす。したがって、上の Zariski 層化
$\underline{G}$ は $\tau$-層化でもあると結論できる。
\end{remark}

\begin{definition}
\label{etale-cohomology-definition-structure-sheaf}
$S$ をスキームとする。$S$ の{\it 構造層}とは、環の層
$\mathcal{O}_S$ で、上で論じたサイト $S_{Zar}$、$S_\etale$、または
$(\Sch/S)_\tau$ のいずれかの上のものをいう。
\end{definition}

\noindent
どのサイト上で考えているかに混同の余地がある場合には、添字を用いて明示する。
例えば、$\mathcal{O}_{S_\etale}$ と書くことにより、$S$ の小エタール・サイト上で
考えていることを強調する場合がある。

\begin{remark}
\label{etale-cohomology-remark-special-case-fpqc-cohomology-quasi-coherent}
上で導入した用語では、定理
\ref{etale-cohomology-theorem-zariski-fpqc-quasi-coherent} の特別な場合として
$$
H_{fppf}^p(X, \mathbf{G}_a) =
H_\etale^p(X, \mathbf{G}_a) =
H_{Zar}^p(X, \mathbf{G}_a) =
H^p(X, \mathcal{O}_X)
$$
がすべての $p \geq 0$ に対して成り立つ。さらに、記法
$H^p_{fppf}(X, \mathcal{O}_X)$ を、$X$ の大 fppf サイト上の
構造層のコホモロジーを表すために用いることもできる。
\end{remark}




\section{Picard 群}
\label{etale-cohomology-section-picard-groups}

\noindent
次の定理は Hilbert の定理 90 と呼ばれることもある。

\begin{theorem}
\label{etale-cohomology-theorem-picard-group}
任意のスキーム $X$ に対して標準的な同一視
\begin{align*}
H_{fppf}^1(X, \mathbf{G}_m) & = H^1_{syntomic}(X, \mathbf{G}_m) \\
& = H^1_{smooth}(X, \mathbf{G}_m) \\
& = H_\etale^1(X, \mathbf{G}_m) \\
& = H^1_{Zar}(X, \mathbf{G}_m) \\
& = \Pic(X) \\
& = H^1(X, \mathcal{O}_X^*)
\end{align*}
\end{theorem}

\begin{proof}
$\tau$ を節 \ref{etale-cohomology-section-big-small} で考えた位相の一つとする。
Cohomology on Sites の補題
\ref{sites-cohomology-lemma-h1-invertible} により
$H^1_\tau(X, \mathbf{G}_m) =
H^1_\tau(X, \mathcal{O}_\tau^*) =
\Pic(\mathcal{O}_\tau)$
であることが分かる。ここで $\mathcal{O}_\tau$ はサイト
$(\Sch/X)_\tau$ の構造層である。可逆 $\mathcal{O}_\tau$-加群は
準連接 $\mathcal{O}_\tau$-加群である。定理 \ref{etale-cohomology-theorem-quasi-coherent}、
またはより精密な Descent の命題
\ref{descent-proposition-equivalence-quasi-coherent} により
$\Pic(\mathcal{O}_\tau) = \Pic(X)$ であることが分かる。
最後の等号も同じ方法で証明される。
\end{proof}






\section{エタール・サイト}
\label{etale-cohomology-section-etale-site}

\noindent
ここからは、スキームのエタール・サイトをさらに詳しく調べる。
第一歩として、エタール射の概念を少し論じる。





\section{エタール射}
\label{etale-cohomology-section-etale-morphism}

\noindent
詳しくは、形式的な定義について Morphisms の節
\ref{morphisms-section-etale} を、エタール射の興味深い性質の概観について
\'Etale Morphisms の各節
\ref{etale-section-etale-morphisms},
\ref{etale-section-structure-etale-map},
\ref{etale-section-etale-smooth},
\ref{etale-section-topological-etale},
\ref{etale-section-functorial-etale}, および
\ref{etale-section-properties-permanence} を参照されたい。

\medskip\noindent
代数 $A$ が代数閉体 $k$ 上有限型で、その微分加群
$\Omega_{A/k}$ が次元に等しい階数の有限局所自由加群であるとき、これを
{\it 滑らか}であるということを思い出そう。スキーム $X$ が $k$ 上
{\it 滑らか}であるとは、それが $k$ 上局所有限型で、各アフィン開集合が
滑らかな $k$-代数のスペクトルであることをいう。$k$ が代数閉でない場合、
$k$-代数 $A$ が滑らかな $k$-代数であるとは、$A \otimes_k \overline{k}$ が
滑らかな $\overline{k}$-代数であることをいう。環準同型 $A \to B$ が滑らかで
あるとは、それが平坦かつ有限表示で、任意の素イデアル $\mathfrak p \subset A$
に対してファイバー環 $\kappa(\mathfrak p) \otimes_A B$ が剰余体
$\kappa(\mathfrak p)$ 上滑らかであることをいう。より一般に、スキームの射が
{\it 滑らか}であるとは、それが平坦かつ局所有限表示で、幾何学的ファイバーが
滑らかであることをいう。

\medskip\noindent
これらの事実については Morphisms の節
\ref{morphisms-section-smooth} を参照されたい。これを用いてエタール射を
次のように定義できる。

\begin{definition}
\label{etale-cohomology-definition-etale-morphism}
スキームの射が{\it エタール}であるとは、相対次元 0 の滑らかな射であることをいう。
\end{definition}

\noindent
特に、スキームの射 $X \to S$ が滑らかで
$\Omega_{X/S} = 0$ ならば、それはエタールである。

\begin{proposition}
\label{etale-cohomology-proposition-etale-morphisms}
エタール射について次が成り立つ。
\begin{enumerate}
\item $k$ を体とする。スキームの射 $U \to \Spec(k)$ がエタールであるための
必要十分条件は、$U \cong \coprod_{i \in I} \Spec(k_i)$ と書けて、各
$i \in I$ に対して環 $k_i$ が体であり、$k$ の有限分離拡大となることである。
\item $\varphi : U \to S$ をスキームの射とする。次の条件は同値である：
\begin{enumerate}
\item $\varphi$ はエタールである。
\item $\varphi$ は局所有限表示かつ平坦であり、そのすべてのファイバーは
エタールである。
\item $\varphi$ は平坦、非分岐、かつ局所有限表示である。
\end{enumerate}
\item 環準同型 $A \to B$ がエタールであるための必要十分条件は、
$B \cong A[x_1, \ldots, x_n]/(f_1, \ldots, f_n)$
と書けて、$\Delta = \det \left( \frac{\partial f_i}{\partial x_j} \right)$ が
$B$ で可逆となることである。
\item エタール射の基底変換はエタールである。
\item エタール射の合成はエタールである。
\item エタール射のファイバー積および積はエタールである。
\item エタール射の相対次元は 0 である。
\item $Y \to X$ をエタール射とする。$X$ が被約（それぞれ正則）ならば、
$Y$ もそうである。
\item エタール射は開写像である。
\item $X \to S$ と $Y \to S$ がエタールならば、任意の $S$-射
$X \to Y$ もエタールである。
\end{enumerate}
\end{proposition}

\begin{proof}
これらの事実（およびさらに多くの事実）は前章までに証明した。
参照箇所は次のとおりである：
(1) Morphisms, 補題 \ref{morphisms-lemma-etale-over-field}。
(2) Morphisms, 補題 \ref{morphisms-lemma-etale-flat-etale-fibres} および
\ref{morphisms-lemma-flat-unramified-etale}。
(3) Algebra, 補題 \ref{algebra-lemma-etale-standard-smooth}。
(4) Morphisms, 補題 \ref{morphisms-lemma-base-change-etale}。
(5) Morphisms, 補題 \ref{morphisms-lemma-composition-etale}。
(6) (4) と (5) から形式的に従う。
(7) Morphisms, 補題 \ref{morphisms-lemma-etale-locally-quasi-finite} および
\ref{morphisms-lemma-locally-quasi-finite-rel-dimension-0}。
(8) Algebra, 補題 \ref{algebra-lemma-reduced-goes-up} および
\ref{algebra-lemma-Rk-goes-up} を参照されたい。同種の結果については
\'Etale Morphisms の節 \ref{etale-section-properties-permanence} も参照されたい。
(9) Morphisms, 補題 \ref{morphisms-lemma-fppf-open} および
\ref{morphisms-lemma-etale-flat} を参照されたい。
(10) Morphisms, 補題 \ref{morphisms-lemma-etale-permanence} を参照されたい。
\end{proof}

\begin{definition}
\label{etale-cohomology-definition-standard-etale}
環準同型 $A \to B$ が{\it 標準エタール}であるとは、
$B \cong \left(A[t]/(f)\right)_g$ と書けて、$f, g \in A[t]$、$f$ はモニック、
かつ $\text{d}f/\text{d}t$ が $B$ で可逆となることをいう。
\end{definition}

\noindent
標準エタールな環準同型はエタールである。実際、
$B = \left(A[t]/(f)\right)_g$、$f, g \in A[t]$ で、$f$ はモニック、かつ
$\text{d}f/\text{d}t$ が $B$ で可逆であると仮定する。このとき $A[t]/(f)$ は
有限自由 $A$-加群であり、その階数はモニック多項式 $f$ の次数に等しい。
したがって $B$ は、この自由代数の局所化として $A$ 上有限表示かつ平坦である。
$B$ がエタールであることの証明を終えるには、ファイバー環
$$
\kappa(\mathfrak p) \otimes_A B
\cong
\kappa(\mathfrak p) \otimes_A (A[t]/(f))_g
\cong
\kappa(\mathfrak p)[t, 1/\overline{g}]/(\overline{f})
$$
が有限分離体拡大の有限積であることを示せば十分である。ここで
$\overline{f}, \overline{g} \in \kappa(\mathfrak p)[t]$ は $f$ と $g$ の像である。
互いに異なる既約モニック因子の冪への分解を
$$
\overline{f} = \overline{f}_1 \ldots \overline{f}_a
\overline{f}_{a + 1}^{e_1} \ldots \overline{f}_{a + b}^{e_b}
$$
$\overline{f}$ に対してとったものとする。ここで各因子は
$\overline{f}_i$ で、$e_1, \ldots, e_b > 1$ とする。仮定により
$\text{d}\overline{f}/\text{d}t$ は
$\kappa(\mathfrak p)[t, 1/\overline{g}]$ で可逆である。したがって少なくとも
すべての $\overline{f}_i$ で $i > a$ のものは可逆である。ゆえに
$$
\kappa(\mathfrak p)[t, 1/\overline{g}]/(\overline{f})
\cong
\prod\nolimits_{i \in I} \kappa(\mathfrak p)[t]/(\overline{f}_i)
$$
を得る。ここで $I \subset \{1, \ldots, a\}$ は添字 $i$ のうち、
$\overline{f}_i$ が $\overline{g}$ を割らないものからなる部分集合である。さらに、
$\text{d}\overline{f}/\text{d}t$ の因子
$\kappa(\mathfrak p)[t]/(\overline{f}_i)$ における像は、明らかに単元と
$\text{d}\overline{f}_i/\text{d}t$ の積に等しい。したがって
$\kappa_i = \kappa(\mathfrak p)[t]/(\overline{f}_i)$ は
$\kappa(\mathfrak p)$ の有限体拡大であり、その最小多項式が分離的な一元で
生成される。すなわち、望んだとおり体拡大
$\kappa_i/\kappa(\mathfrak p)$ は有限分離拡大である。

\medskip\noindent
任意のエタールな環準同型は局所的に標準エタールであることが分かる。
これを定式化するため、次の用語を導入する。環準同型 $A \to B$ が、素イデアル
$\mathfrak q$ で $B$ のものにおいて{\it エタール}であるとは、ある
$h \in B$ で $h \not \in \mathfrak q$ を満たし、$A \to B_h$ がエタールと
なるものが存在することをいう。結果は次のとおりである。

\begin{theorem}
\label{etale-cohomology-theorem-standard-etale}
環準同型 $A \to B$ が素イデアル $\mathfrak q$ においてエタールであるための
必要十分条件は、ある $g \in B$ で $g \not \in \mathfrak q$ を満たし、
$B_g$ が $A$ 上標準エタールとなるものが存在することである。
\end{theorem}

\begin{proof}
Algebra の命題 \ref{algebra-proposition-etale-locally-standard} を参照されたい。
\end{proof}





\section{エタール被覆}
\label{etale-cohomology-section-etale-covering}

\noindent
定義を思い出そう。

\begin{definition}
\label{etale-cohomology-definition-etale-covering}
スキーム $U$ の{\it エタール被覆}とは、スキームの射の族
$\{\varphi_i : U_i \to U\}_{i \in I}$ で、次を満たすものをいう：
\begin{enumerate}
\item 各 $\varphi_i$ はエタール射である。
\item $U_i$ は $U$ を被覆する。すなわち
$U = \bigcup_{i\in I}\varphi_i(U_i)$ である。
\end{enumerate}
\end{definition}

\begin{lemma}
\label{etale-cohomology-lemma-etale-fpqc}
任意のエタール被覆は fpqc 被覆である。
\end{lemma}

\begin{proof}
(Topologies の補題
\ref{topologies-lemma-zariski-etale-smooth-syntomic-fppf-fpqc} も参照されたい。)
$\{\varphi_i : U_i \to U\}_{i \in I}$ をエタール被覆とする。
エタール射は平坦であり、被覆の元はその終域を被覆するので、性質 fp
（忠実平坦）が満たされる。性質 qc（準コンパクト）を確かめるため、
$V \subset U$ をアフィン開集合とし、
$\varphi_i^{-1}(V) = \bigcup_{j \in J_i} V_{ij}$ と書く。ここで
$V_{ij} \subset U_i$ はアフィン開集合である。$\varphi_i$ は開写像
（エタール射は開写像だから）なので、
$V = \bigcup_{i\in I} \bigcup_{j \in J_i} \varphi_i(V_{ij})$
が $V$ の開被覆であることが分かる。さらに、$V$ は準コンパクトなので、
この被覆は有限な細分をもつ。
\end{proof}

\noindent
したがって、fpqc 被覆について正しい任意の主張は、エタール被覆については
{\it なおさら}正しい。例えば、エタール・サイトは準標準的である。

\begin{definition}
\label{etale-cohomology-definition-big-etale-site}
(詳しくは節 \ref{etale-cohomology-section-big-small} または Topologies の節
\ref{topologies-section-etale} を参照されたい。)
$S$ をスキームとする。{\it $S$ 上の大エタール・サイト}とはサイト
$(\Sch/S)_\etale$ のことである。定義 \ref{etale-cohomology-definition-tau-site} を参照されたい。
{\it $S$ 上の小エタール・サイト}とはサイト $S_\etale$ のことである。
定義 \ref{etale-cohomology-definition-tau-site} を参照されたい。同様に、$S$ 上の
{\it 大 Zariski サイト}および{\it 小 Zariski サイト}を定義し、それぞれ
$(\Sch/S)_{Zar}$ および $S_{Zar}$ と書く。
\end{definition}

\noindent
大まかにいえば、$S$ の大エタール・サイトは $S$ 上のスキームからなり、
その被覆はエタール被覆である。$S$ の小エタール・サイトは $S$ 上エタールな
スキームからなり、その被覆はエタール被覆である。実際、命題
\ref{etale-cohomology-proposition-etale-morphisms} により、$S_\etale$ の対象間の任意の射は
エタールである。したがって、$\{U_i \to U\}_{i \in I}$ が
$S_\etale$ における被覆であることを確かめるには、
$\coprod U_i \to U$ が全射であることを確かめれば十分である。

\medskip\noindent
小エタール・サイトは大エタール・サイトより対象が少なく、$S$ 上の
エタール位相の「開集合」だけを含む。これは大エタール・サイトの充満部分圏で、
その位相は大サイトの位相から誘導される。したがって、大エタール・サイトから
小エタール・サイトへの制限関手は完全で、入射対象を入射対象へ写す。
このことから次が従う。

\begin{proposition}
\label{etale-cohomology-proposition-cohomology-restrict-small-site}
$S$ をスキーム、$\mathcal{F}$ を $(\Sch/S)_\etale$ 上のアーベル群の層とする。
このとき $\mathcal{F}|_{S_\etale}$ は $S_\etale$ 上の層であり、
$$
H^p_\etale(S, \mathcal{F}|_{S_\etale}) =
H^p_\etale(S, \mathcal{F})
$$
がすべての $p \geq 0$ に対して成り立つ。
\end{proposition}

\begin{proof}
これは補題 \ref{etale-cohomology-lemma-compare-cohomology-big-small} の特別な場合である。
\end{proof}

\noindent
節 \ref{etale-cohomology-section-big-small} で導入した一般的な記法に従い、上のコホモロジー群を
$H_\etale^p(S, \mathcal{F})$ と書く。





%9.24.09
\section{クンマー理論}
\label{etale-cohomology-section-kummer}

\noindent
$n \in \mathbf{N}$ とし、次で定義される関手 $\mu_n$ を考える：
$$
\begin{matrix}
\Sch^{opp} & \longrightarrow & \textit{Ab} \\
S & \longmapsto &
\mu_n(S)
=
\{t \in \Gamma(S, \mathcal{O}_S^*) \mid t^n = 1 \}.
\end{matrix}
$$
Groupoids の例 \ref{groupoids-example-roots-of-unity} により、これは
表現可能な関手であり、それを表現するスキームも $\mu_n$ と書く。補題
\ref{etale-cohomology-lemma-representable-sheaf-fpqc} により、この関手は fpqc 位相に関する
層条件を満たす（特に、エタール位相、Zariski 位相などに関する層条件も満たす）。

\begin{lemma}
\label{etale-cohomology-lemma-kummer-sequence}
$n\in \mathcal{O}_S^*$ ならば、
$$
0 \to
\mu_{n, S} \to
\mathbf{G}_{m, S} \xrightarrow{(\cdot)^n}
\mathbf{G}_{m, S} \to 0
$$
は $S$ の小エタール・サイトと大エタール・サイトの両方の上で、層の短完全列である。
\end{lemma}

\begin{proof}
定義により、層 $\mu_{n, S}$ は写像 $(\cdot)^n$ の核である。したがって、
最後の写像が全射であることを示せば十分である。$U$ を $S$ 上のスキームとし、
$f \in \mathbf{G}_m(U) = \Gamma(U, \mathcal{O}_U^*)$ とする。
$U$ のエタール被覆で、その各元の上で $f$ の制限が $n$ 乗となるものが
見つかることを示す必要がある。次のようにおく：
$$
U' =
\underline{\Spec}_U(\mathcal{O}_U[T]/(T^n-f))
\xrightarrow{\pi}
U.
$$
(相対スペクトルについては Constructions の節
\ref{constructions-section-spec-via-glueing} または
\ref{constructions-section-spec} を参照されたい。)
$\Spec(A) \subset U$ をアフィン開集合とし、$f|_{\Spec(A)}$ が単元
$a \in A^*$ に対応するとする。このとき $\pi^{-1}(\Spec(A)) = \Spec(B)$ であり、
$B = A[T]/(T^n - a)$ である。環準同型 $A \to B$ は階数 $n$ の有限自由射なので
忠実平坦である。したがって $\Spec(B) \to \Spec(A)$ は全射である。これは
$U$ の任意のアフィン開集合について成り立つので、$\pi$ は全射である。
さらに、$n$ と $T^{n - 1}$ は $B$ で可逆であるから、
$nT^{n-1} \in B^*$ であり、環準同型 $A \to B$ は標準エタール、特に
エタールである。これは $U$ の任意のアフィン開集合について成り立つので、
$\pi$ はエタールである。ゆえに
$\mathcal{U} = \{\pi : U' \to U\}$ はエタール被覆である。さらに、
$f|_{U'} = (f')^n$ であり、ここで $f'$ は $T$ の
$\Gamma(U', \mathcal{O}_{U'}^*)$ における類である。したがって
$\mathcal{U}$ は所望の性質をもつ。
\end{proof}

\begin{remark}
\label{etale-cohomology-remark-no-kummer-sequence-zariski}
補題 \ref{etale-cohomology-lemma-kummer-sequence} は「エタール」を「Zariski」に置き換えると
偽である。エタール位相は滑らかな位相より粗いので、Topologies の補題
\ref{topologies-lemma-zariski-etale-smooth} を参照すれば、この列は滑らかな
位相においても完全であることが従う。
\end{remark}

\noindent
定理 \ref{etale-cohomology-theorem-picard-group}、補題 \ref{etale-cohomology-lemma-kummer-sequence}、および
コホモロジーの一般的性質により、コホモロジーの長完全列
$$
\xymatrix{
0 \ar[r] &
H_\etale^0(S, \mu_{n, S}) \ar[r] &
\Gamma(S, \mathcal{O}_S^*) \ar[r]^{(\cdot)^n} &
\Gamma(S, \mathcal{O}_S^*) \ar@(rd, ul)[rdllllr]
\\
& H_\etale^1(S, \mu_{n, S}) \ar[r] &
\Pic(S) \ar[r]^{(\cdot)^n} &
\Pic(S) \ar@(rd, ul)[rdllllr] \\
& H_\etale^2(S, \mu_{n, S}) \ar[r] &
\ldots
}
$$
を得る。少なくとも、これは $n$ が $S$ 上可逆ならば成り立つ。$n$ が $S$ 上
可逆でないときには、次の補題を適用できる。

\begin{lemma}
\label{etale-cohomology-lemma-kummer-sequence-syntomic}
任意の $n \in \mathbf{N}$ に対して、列
$$
0 \to
\mu_{n, S} \to
\mathbf{G}_{m, S} \xrightarrow{(\cdot)^n}
\mathbf{G}_{m, S} \to 0
$$
はサイト $(\Sch/S)_{fppf}$ および $(\Sch/S)_{syntomic}$ 上の層の短完全列である。
\end{lemma}

\begin{proof}
定義により、層 $\mu_{n, S}$ は写像 $(\cdot)^n$ の核である。したがって、
最後の写像が全射であることを示せば十分である。syntomic 位相は fppf 位相より
弱いので、Topologies の補題
\ref{topologies-lemma-zariski-etale-smooth-syntomic-fppf} を参照すれば、
syntomic 位相について証明すれば十分である。$U$ を $S$ 上のスキームとし、
$f \in \mathbf{G}_m(U) = \Gamma(U, \mathcal{O}_U^*)$ とする。
$U$ の syntomic 被覆で、その各元の上で $f$ の制限が $n$ 乗となるものが
見つかることを示す必要がある。次のようにおく：
$$
U' =
\underline{\Spec}_U(\mathcal{O}_U[T]/(T^n-f))
\xrightarrow{\pi}
U.
$$
(相対スペクトルについては Constructions の節
\ref{constructions-section-spec-via-glueing} または
\ref{constructions-section-spec} を参照されたい。)
$\Spec(A) \subset U$ をアフィン開集合とし、$f|_{\Spec(A)}$ が単元
$a \in A^*$ に対応するとする。このとき $\pi^{-1}(\Spec(A)) = \Spec(B)$ であり、
$B = A[T]/(T^n - a)$ である。環準同型 $A \to B$ は階数 $n$ の有限自由射なので
忠実平坦である。したがって $\Spec(B) \to \Spec(A)$ は全射である。これは
$U$ の任意のアフィン開集合について成り立つので、$\pi$ は全射である。
さらに、$B$ は $A$ 上の大域的相対完全交叉なので、環準同型 $A \to B$ は
標準 syntomic、特に syntomic である。これは $U$ の任意のアフィン開集合に
ついて成り立つので、$\pi$ は syntomic である。ゆえに
$\mathcal{U} = \{\pi : U' \to U\}$ は syntomic 被覆である。さらに、
$f|_{U'} = (f')^n$ であり、ここで $f'$ は $T$ の
$\Gamma(U', \mathcal{O}_{U'}^*)$ における類である。したがって
$\mathcal{U}$ は所望の性質をもつ。
\end{proof}

\begin{remark}
\label{etale-cohomology-remark-no-kummer-sequence-smooth-etale-zariski}
補題 \ref{etale-cohomology-lemma-kummer-sequence-syntomic} は、滑らかな位相、エタール位相、
または Zariski 位相に対しては偽である。
\end{remark}

\noindent
定理 \ref{etale-cohomology-theorem-picard-group}、補題
\ref{etale-cohomology-lemma-kummer-sequence-syntomic}、およびコホモロジーの一般的性質により、
コホモロジーの長完全列
$$
\xymatrix{
0 \ar[r] &
H_{fppf}^0(S, \mu_{n, S}) \ar[r] &
\Gamma(S, \mathcal{O}_S^*) \ar[r]^{(\cdot)^n} &
\Gamma(S, \mathcal{O}_S^*) \ar@(rd, ul)[rdllllr]
\\
& H_{fppf}^1(S, \mu_{n, S}) \ar[r] &
\Pic(S) \ar[r]^{(\cdot)^n} &
\Pic(S) \ar@(rd, ul)[rdllllr] \\
& H_{fppf}^2(S, \mu_{n, S}) \ar[r] &
\ldots
}
$$
を任意のスキーム $S$ と任意の整数 $n$ に対して得る。もちろん、syntomic
コホモロジーについても同様の列がある。

\medskip\noindent
$n \in \mathbf{N}$ とし、$S$ を任意のスキームとする。$\mu_n$ に値をとる
第一コホモロジー群を記述する、より直接的な別の方法がある。対
$(\mathcal{L}, \alpha)$ を考える。ここで $\mathcal{L}$ は $S$ 上の可逆層で、
$\alpha : \mathcal{L}^{\otimes n} \to \mathcal{O}_S$ は $n$ 乗テンソル冪の
自明化である（このテンソル冪は可逆層 $\mathcal{L}$ のものである）。
$(\mathcal{L}', \alpha')$ をもう一つの
このような対とする。同型
$\varphi : (\mathcal{L}, \alpha) \to (\mathcal{L}', \alpha')$ とは、図式
$\varphi : \mathcal{L} \to \mathcal{L}'$ が可逆層の同型であって、次の図式が
$$
\xymatrix{
\mathcal{L}^{\otimes n} \ar[d]_{\varphi^{\otimes n}} \ar[r]_\alpha &
\mathcal{O}_S \ar[d]^1 \\
(\mathcal{L}')^{\otimes n} \ar[r]^{\alpha'} &
\mathcal{O}_S \\
}
$$
可換となるものをいう。したがって
\begin{equation}
\label{etale-cohomology-equation-isomorphisms-pairs}
\mathit{Isom}_S((\mathcal{L}, \alpha), (\mathcal{L}', \alpha'))
=
\left\{
\begin{matrix}
\emptyset & \text{もし} & \text{これらが同型でないならば} \\
H^0(S, \mu_{n, S})\cdot \varphi & \text{もし} &
\varphi \text{ が対の同型ならば}
\end{matrix}
\right.
\end{equation}
を得る。さらに、二つの対 $(\mathcal{L}, \alpha)$、$(\mathcal{L}', \alpha')$ が
与えられたとき、そのテンソル積
$$
(\mathcal{L}, \alpha) \otimes (\mathcal{L}', \alpha')
=
(\mathcal{L} \otimes \mathcal{L}', \alpha \otimes \alpha')
$$
もまた対である。対 $(\mathcal{O}_S, 1)$ はこのテンソル積演算の単位元であり、
逆元は
$$
(\mathcal{L}, \alpha)^{-1} = (\mathcal{L}^{\otimes -1}, \alpha^{\otimes -1}).
$$
で与えられる。したがって、対の同型類の集合はアーベル群をなす。さらに、
$$
(\mathcal{L}, \alpha)^{\otimes n}
=
(\mathcal{L}^{\otimes n}, \alpha^{\otimes n})
\xrightarrow{\alpha}
(\mathcal{O}_S, 1)
$$
は同型である。したがって、この群の各元の位数は $n$ を割る。この群は一般には
$n$-捩れ、すなわち $\Pic(S)$ における捩れ部分{\bf ではない}ことに注意されたい。

\begin{lemma}
\label{etale-cohomology-lemma-describe-h1-mun}
$S$ をスキームとする。標準的な同一視
$$
H_\etale^1(S, \mu_n) =
\text{上記の対 }(\mathcal{L}, \alpha)\text{ の同型類からなる群}
$$
が $n$ が $S$ 上可逆ならば存在する。一般には
$$
H_{fppf}^1(S, \mu_n) =
\text{上記の対 }(\mathcal{L}, \alpha)\text{ の同型類からなる群}.
$$
同じ結果は fppf を syntomic で置き換えても成り立つ。
\end{lemma}

\begin{proof}
まず第二の同型を証明する。$(\mathcal{L}, \alpha)$ を上記の対とする。
アフィン開被覆 $S = \bigcup U_i$ で
$\mathcal{L}|_{U_i} \cong \mathcal{O}_{U_i}$ となるものを選ぶ。
$s_i \in \mathcal{L}(U_i)$ を生成元とする。このとき
$\alpha(s_i^{\otimes n}) = f_i \in \mathcal{O}_S^*(U_i)$ である。
$U_i = \Spec(A_i)$ と書けば、大域的相対完全交叉
$A_i \to B_i = A_i[T]/(T^n - f_i)$ で、$f_i$ の像が $n$ 乗として
$B_i$ に現れるものが存在することが分かる。言い換えると、
$V_i = \Spec(B_i)$ とおけば、syntomic 被覆
$\mathcal{V} = \{V_i \to S\}_{i \in I}$ と自明化
$\varphi_i : (\mathcal{L}, \alpha)|_{V_i} \to (\mathcal{O}_{V_i}, 1)$ を得る。

\medskip\noindent
この結果（被覆 $\mathcal{V}$ の存在）を用いて、この対にコホモロジー類
$H^1_{syntomic}(S, \mu_{n, S})$ の元を対応させる。同値な二つの構成を与える。

\medskip\noindent
第一の構成：{\v C}ech コホモロジーを用いる。二重交叉
$V_i \times_S V_j$ 上で、対の同型
$$
(\mathcal{O}_{V_i \times_S V_j}, 1)
\xrightarrow{\text{pr}_0^*\varphi_i^{-1}}
(\mathcal{L}|_{V_i \times_S V_j}, \alpha|_{V_i \times_S V_j})
\xrightarrow{\text{pr}_1^*\varphi_j}
(\mathcal{O}_{V_i \times_S V_j}, 1)
$$
を得る。式 (\ref{etale-cohomology-equation-isomorphisms-pairs}) により、これは元
$\zeta_{ij} \in \mu_n(V_i \times_S V_j)$ によって与えられる。これらの
$\zeta_{ij}$ が $1$-コサイクル、すなわち
$(\zeta_{i_0i_1}) \in \check C(\mathcal{V}, \mu_n)$ で
$d(\zeta_{i_0i_1}) = 0$ を満たす元を与えることの確認は省略する。したがって、
その類は $\check H^1(\mathcal{V}, \mu_n)$ の元であり、定理
\ref{etale-cohomology-theorem-cech-ss} によりコホモロジー類
$H^1_{syntomic}(S, \mu_{n, S})$ の元へ写る。

\medskip\noindent
第二の構成：トーサーを用いる。前層
$$
\mu_n(\mathcal{L}, \alpha) :
U
\longmapsto
\mathit{Isom}_U((\mathcal{O}_U, 1), (\mathcal{L}, \alpha)|_U)
$$
を $(\Sch/S)_{syntomic}$ 上で考える。これは
$\SheafHom_\mathcal{O}(\mathcal{O}, \mathcal{L})$ の部分前層とみなせる
（内部 Hom 層については Modules on Sites の節
\ref{sites-modules-section-internal-hom} を参照されたい）。この部分前層を
定義する条件は局所的なので、これは層である。式
(\ref{etale-cohomology-equation-isomorphisms-pairs}) により、この層には層 $\mu_{n, S}$ が
自由に作用する。したがって確認すべきことは、局所的に切断をもつことだけである。
これは自明化被覆 $\mathcal{V}$ が存在するので正しい。ゆえに
$\mu_n(\mathcal{L}, \alpha)$ は $\mu_{n, S}$-トーサーであり、Cohomology on Sites
の補題 \ref{sites-cohomology-lemma-torsors-h1} により、対応する元
$H^1_{syntomic}(S, \mu_{n, S})$ の元を得る。

\medskip\noindent
残るのは次を示すことである：
\begin{enumerate}
\item 二つの構成は同じコホモロジー類を与える。
\item 同型な対は同じコホモロジー類を生じる。
\item
$(\mathcal{L}, \alpha) \otimes (\mathcal{L}', \alpha')$
のコホモロジー類は、$(\mathcal{L}, \alpha)$ と
$(\mathcal{L}', \alpha')$ のコホモロジー類の和である。
\item コホモロジー類が自明ならば、その対は自明である。
\item $H^1_{syntomic}(S, \mu_{n, S})$ の任意の元は、ある対の
コホモロジー類である。
\end{enumerate}
(1) の証明は省略する。(2) は第二の構成から明らかである。同型なトーサーは
同じコホモロジー類を与えるからである。(3) は第一の構成から明らかである。
得られる {\v C}ech 類が加法的だからである。(4) は第二の構成から明らかである。
トーサーが自明であるための必要十分条件は大域切断をもつことである。
Cohomology on Sites の補題 \ref{sites-cohomology-lemma-trivial-torsor} を参照されたい。

\medskip\noindent
(5) は次のように分かる（直接的な証明の方が望ましいが）。
$\xi \in H^1_{syntomic}(S, \mu_{n, S})$ と仮定する。このとき $\xi$ は元
$\overline{\xi} \in H^1_{syntomic}(S, \mathbf{G}_{m, S})$ で
$n \overline{\xi} = 0$ を満たすものへ写る。定理
\ref{etale-cohomology-theorem-picard-group} により、$\overline{\xi}$ は可逆層 $\mathcal{L}$ で、
その $n$ 乗テンソル冪が $\mathcal{O}_S$ と同型となるものに対応する。したがって対
$(\mathcal{L}, \alpha')$ で、そのコホモロジー類 $\xi'$ の像
$\overline{\xi'}$ が $H^1_{syntomic}(S, \mathbf{G}_{m, S})$ において同じとなる
ものが存在する。ゆえに $\xi - \xi'$ が対の類であることを示せば十分である。
構成と上のコホモロジー長完全列により、
$\xi - \xi' = \partial(f)$ が、ある $f \in H^0(S, \mathcal{O}_S^*)$ に対して
成り立つことが分かる。対 $(\mathcal{O}_S, f)$ を
考える。この対のコホモロジー類が $\partial(f)$ であることの確認は省略する。
これで第一の同一視の証明が終わる（fppf を syntomic で置き換えたものについて）。

\medskip\noindent
第一の同一視を見るには、$n$ が $S$ 上可逆ならば、証明の最初の部分で構成した
被覆 $\mathcal{V}$ が実際にエタール被覆であることに注意すればよい
（補題 \ref{etale-cohomology-lemma-kummer-sequence} の証明と比較されたい）。証明の残りは位相に
依存しない。ただし最後の議論だけはクンマー列が完全であること、すなわち補題
\ref{etale-cohomology-lemma-kummer-sequence} を用いる。
\end{proof}






\section{近傍、茎、および点}
\label{etale-cohomology-section-stalks}

\noindent
$S$ の任意の幾何点には、完全な茎関手を対応させることができる。
$S_\etale$ 上の層の射が同型であるための必要十分条件は、これらすべての茎上で
同型であることである。アーベル層の複体が完全であるための必要十分条件は、
すべての幾何点において茎の複体が完全であることである。まとめると、これは
スキーム $S$ の小エタール・サイトが十分な点をもつことを意味する。また、
$S$ の小エタール・トポスの任意の点（これは抽象的な概念である）は幾何点から
与えられることも分かる。したがって、ある意味で $S$ の小エタール・トポスは、
幾何点と近傍を用いて理解できる。

\begin{definition}
\label{etale-cohomology-definition-geometric-point}
$S$ をスキームとする。
\begin{enumerate}
\item $S$ の {\it 幾何点}とは、射 $\Spec(k) \to S$ であって、
$k$ が代数閉体であるもののことである。このような点は通常、上線を付けた小文字
$\overline{s}$ で表す。$\overline{s}$ は射だけでなくスキーム $\Spec(k)$ を
表すためにも用い、$\kappa(\overline{s})$ で $k$ を表す。
\item $\overline{s}$ が $s$ の {\it 上にある}とは、
$s \in S$ が $\overline{s}$ の像であることをいう。
\item 幾何点 $\overline{s}$（$S$ のもの）の {\it エタール近傍}とは、可換図式
$$
\xymatrix{
& U \ar[d]^\varphi \\
{\overline{s}} \ar[r]^{\overline{s}} \ar[ur]^{\bar u} & S
}
$$
であって、$\varphi$ がスキームのエタール射であるもののことである。
これを $(U, \overline{u}) \to (S, \overline{s})$ と書く。
\item {\it エタール近傍の射}
$(U, \overline{u}) \to (U', \overline{u}')$
とは、$S$-射 $h: U \to U'$ であって
$\overline{u}' = h \circ \overline{u}$ を満たすもののことである。
\end{enumerate}
\end{definition}

\begin{remark}
\label{etale-cohomology-remark-etale-between-etale}
$U$ と $U'$ は $S$ 上エタールであるから、それらの間の任意の
$S$-射もエタールである。命題 \ref{etale-cohomology-proposition-etale-morphisms} を参照。
とくに、エタール近傍の射はすべてエタールである。
\end{remark}

\begin{remark}
\label{etale-cohomology-remark-etale-neighbourhoods}
$S$ をスキーム、$s \in S$ を点とする。More on Morphisms,
定義 \ref{more-morphisms-definition-etale-neighbourhood} では、
エタール近傍 $(U, u) \to (S, s)$、すなわち $(S, s)$ の近傍という概念を定義した。
$\overline{s}$ が $S$ の幾何点で $s$ の上にあるなら、任意のエタール近傍
$(U, \overline{u}) \to (S, \overline{s})$ から $(U, u)$ を
$(S, s)$ のエタール近傍として得る。ここで $u \in U$ は、$U$ の一意な点であって
$\overline{u}$ が $u$ の上にあるものとする。逆に、エタール近傍 $(U, u)$、
すなわち $(S, s)$ の近傍が与えられると、
剰余体の拡大 $\kappa(u)/\kappa(s)$ は有限分離的であり（命題
\ref{etale-cohomology-proposition-etale-morphisms} を参照）、したがって埋め込み
$\kappa(u) \subset \kappa(\overline{s})$（$\kappa(s)$ 上のもの）を見いだせる。
言い換えれば、幾何点 $\overline{u}$（$U$ のもの）で $u$ の上にあり、
$(U, \overline{u})$ が
$(S, \overline{s})$ のエタール近傍となるものを見いだせる。以下ではこれらの観察を用いて、
二種類のエタール近傍の間を行き来する。
\end{remark}

\begin{lemma}
\label{etale-cohomology-lemma-cofinal-etale}
$S$ をスキーム、$\overline{s}$ を $S$ の幾何点とする。
エタール近傍の圏は余フィルター圏である。より正確には、次が成り立つ。
\begin{enumerate}
\item $(U_i, \overline{u}_i)_{i = 1, 2}$ を $\overline{s}$ の
二つのエタール近傍（$S$ におけるもの）とする。このとき、第三のエタール近傍
$(U, \overline{u})$ と射
$(U, \overline{u}) \to (U_i, \overline{u}_i)$, $i = 1, 2$.
が存在する。
\item $h_1, h_2: (U, \overline{u}) \to (U', \overline{u}')$ を
$\overline{s}$ のエタール近傍の間の二つの射とする。このとき、エタール近傍
$(U'', \overline{u}'')$ と射
$h : (U'', \overline{u}'') \to (U, \overline{u})$
であって $h_1$ と $h_2$ を等化するもの、すなわち
$h_1 \circ h = h_2 \circ h$ を満たすものが存在する。
\end{enumerate}
\end{lemma}

\begin{proof}
(1) について、ファイバー積 $U = U_1 \times_S U_2$ を考える。
エタール射は基底変換で保たれるので、これは $U_1$ と $U_2$ のいずれの上でも
エタールである。命題 \ref{etale-cohomology-proposition-etale-morphisms} を参照。
写像 $\overline{s} \to U$（$(\overline{u}_1, \overline{u}_2)$ により定まるもの）は、
これに $U_1$ と $U_2$ の双方へ写る
エタール近傍の構造を与える。(2) については、$U''$ をファイバー積
$$
\xymatrix{
U'' \ar[r] \ar[d] & U \ar[d]^{(h_1, h_2)} \\
U' \ar[r]^-\Delta & U' \times_S U'.
}
$$
として定義する。$\overline{u}$ と $\overline{u}'$ は $S$ 上で
$\overline{s}$ と一致するから、
$\overline{u}'' = (\overline{u}, \overline{u}')$ は $U''$ の幾何点である。
とくに $U'' \not = \emptyset$ である。さらに、$U'$ は $S$ 上エタールなので、
ファイバー積 $U'\times_S U'$ もそうである（命題
\ref{etale-cohomology-proposition-etale-morphisms} を参照）。したがって、上の注意
\ref{etale-cohomology-remark-etale-between-etale} により縦の矢印 $(h_1, h_2)$ はエタールである。
ゆえに基底変換により $U''$ は $U'$ 上エタールであり、したがって $S$ 上でも
エタールである（エタール射の合成はエタールだからである）。
よって $(U'', \overline{u}'')$ が求める解である。
\end{proof}

\begin{lemma}
\label{etale-cohomology-lemma-geometric-lift-to-cover}
$S$ をスキームとし、$\overline{s}$ を $S$ の幾何点とする。
$(U, \overline{u})$ を $\overline{s}$ のエタール近傍とし、
$\mathcal{U} = \{\varphi_i : U_i \to U \}_{i\in I}$ をエタール被覆とする。
このとき $i \in I$ と $\overline{u}_i : \overline{s} \to U_i$ であって、
$\varphi_i : (U_i, \overline{u}_i) \to (U, \overline{u})$ が
エタール近傍の射となるものが存在する。
\end{lemma}

\begin{proof}
$U = \bigcup_{i\in I} \varphi_i(U_i)$ なので、ファイバー積
$\overline{s} \times_{\overline{u}, U, \varphi_i} U_i$ は、ある $i$ に対して
空でない。このとき、カルテジアン図式
$$
\xymatrix{
\overline{s} \times_{\overline{u}, U, \varphi_i} U_i
\ar[d]^{\text{pr}_1} \ar[r]_-{\text{pr}_2} & U_i
\ar[d]^{\varphi_i} \\
\Spec(k) = \overline{s} \ar@/^1pc/[u]^\sigma
\ar[r]^-{\overline{u}} & U
}
$$
を見る。射影 $\text{pr}_1$ はエタール射の基底変換なのでエタールである。
命題 \ref{etale-cohomology-proposition-etale-morphisms} を参照。したがって、命題
\ref{etale-cohomology-proposition-etale-morphisms} により
$\overline{s} \times_{\overline{u}, U, \varphi_i} U_i$ は $k$ の有限分離拡大の
非交和である。ここで $\overline{s} = \Spec(k)$ である。しかし $k$ は代数閉体なので、
これらの拡大はすべて自明であり、切断 $\sigma$（$\text{pr}_1$ のもの）が存在する。
合成 $\text{pr}_2 \circ \sigma$ が $\overline{u}$ と整合する写像を与える。
\end{proof}

\begin{definition}
\label{etale-cohomology-definition-stalk}
$S$ をスキーム、$\mathcal{F}$ を $S_\etale$ 上の前層、
$\overline{s}$ を $S$ の幾何点とする。$\mathcal{F}$ の {\it 茎}
（$\overline{s}$ におけるもの）とは、
$$
\mathcal{F}_{\overline{s}}
=
\colim_{(U, \overline{u})} \mathcal{F}(U)
$$
である。ここで $(U, \overline{u})$ は $\overline{s}$ のすべての
エタール近傍（$S$ におけるもの）を走る。
\end{definition}

\noindent
補題 \ref{etale-cohomology-lemma-cofinal-etale} により、この余極限はフィルター圏を
添字とする。すなわち、エタール近傍の圏の反対圏を添字とする。
言い換えれば、$\mathcal{F}_{\overline{s}}$ の元は
三つ組 $(U, \overline{u}, \sigma)$（ここで
$\sigma \in \mathcal{F}(U)$）と考えられる。二つの三つ組
$(U, \overline{u}, \sigma)$, $(U', \overline{u}', \sigma')$
が茎の同じ元を定めるための必要十分条件は、第三のエタール近傍
$(U'', \overline{u}'')$ とエタール近傍の射
$h : (U'', \overline{u}'') \to (U, \overline{u})$,
$h' : (U'', \overline{u}'') \to (U', \overline{u}')$ であって、
$h^*\sigma = (h')^*\sigma'$ が $\mathcal{F}(U'')$ において成り立つものが
存在することである。
Categories, 節 \ref{categories-section-directed-colimits} を参照。

\begin{lemma}
\label{etale-cohomology-lemma-stalk-gives-point}
$S$ をスキームとし、$\overline{s}$ を $S$ の幾何点とする。関手
\begin{align*}
u : S_\etale & \longrightarrow \textit{Sets}, \\
U & \longmapsto
|U_{\overline{s}}|
=
\{\overline{u} \text{ であって }(U, \overline{u})
\text{ が }\overline{s}\text{ のエタール近傍であるもの}\}.
\end{align*}
を考える。ここで $|U_{\overline{s}}|$ は幾何学的ファイバーの台集合を表す。
このとき $u$ は点 $p$（サイト $S_\etale$ のもの）を定め
（Sites, 定義 \ref{sites-definition-point}）、それに付随する茎関手
$\mathcal{F} \mapsto \mathcal{F}_p$
（Sites, 式 \ref{sites-equation-stalk}）は、上で定義した関手
$\mathcal{F} \mapsto \mathcal{F}_{\overline{s}}$ である。
\end{lemma}

\begin{proof}
補題 \ref{etale-cohomology-lemma-geometric-lift-to-cover} の証明で、スキーム
$U_{\overline{s}}$ が $\overline{s}$ と同型なスキームの非交和であることを見た。
したがって $|U_{\overline{s}}|$ は、$U$ の幾何点で $\overline{s}$ の上にあるものの集合、
すなわち定義 \ref{etale-cohomology-definition-geometric-point} の図式に収まる射
$\overline{u} : \overline{s} \to U$ の集まりとも考えられる。
これより $u(S)$ は一元集合であり、さらに
$u(U \times_V W) = u(U) \times_{u(V)} u(W)$
が、$U \to V$ と $W \to V$ が $S_\etale$ の射であるたびに成り立つ。
また、被覆 $\{U_i \to U\}_{i \in I}$（$S_\etale$ におけるもの）が与えられると、
補題 \ref{etale-cohomology-lemma-geometric-lift-to-cover} により
$\coprod u(U_i) \to u(U)$ は全射である。したがって Sites, 命題
\ref{sites-proposition-point-limits} を適用でき、$p$ は
$S_\etale$ の点である。最後に、関手
$\mathcal{F} \mapsto \mathcal{F}_{\overline{s}}$ は、Sites, 式
\ref{sites-equation-stalk} の関手
$\mathcal{F} \mapsto \mathcal{F}_p$（$p$ に付随するもの）と
まったく同じ余極限で与えられる。
これで最後の主張が従う。
\end{proof}

\begin{remark}
\label{etale-cohomology-remark-map-stalks}
$S$ をスキームとし、$\overline{s} : \Spec(k) \to S$ および
$\overline{s}' : \Spec(k') \to S$ を $S$ の二つの幾何点とする。
{\it 幾何点の射 $a : \overline{s} \to \overline{s}'$}とは、単に
射 $a : \Spec(k) \to \Spec(k')$ であって
$\overline{s}' \circ a = \overline{s}$ を満たすもののことである。
このような射が与えられると、$\overline{s}'$ のエタール近傍の圏から
$\overline{s}$ のエタール近傍の圏への関手を、規則
$(U, \overline{u}') \mapsto (U, \overline{u}' \circ a)$ により得る。
したがって、Categories, 補題
\ref{categories-lemma-functorial-colimit} から標準写像
$$
\mathcal{F}_{\overline{s}'}
=
\colim_{(U, \overline{u}')} \mathcal{F}(U)
\longrightarrow
\colim_{(U, \overline{u})} \mathcal{F}(U)
=
\mathcal{F}_{\overline{s}}
$$
を得る。茎の元を三つ組として記述すると、これは
$\mathcal{F}_{\overline{s}'}$ の元で、三つ組
$(U, \overline{u}', \sigma)$ により表されるものを
$\mathcal{F}_{\overline{s}}$ の元で、三つ組
$(U, \overline{u}' \circ a, \sigma)$ により表されるものへ写す。
上の関手は明らかに圏同値なので、
この標準写像は茎関手の同型である。

\medskip\noindent
$a$ に対応する茎の写像の向きが正しいことを確認しておこう。
Sites, 定義 \ref{sites-definition-morphism-points} によれば、上のことは、
$a$ が射 $a : p \to p'$ を定めることを意味する。この射は点 $p, p'$ の間の射であり、
これらはサイト $S_\etale$ の点で、補題 \ref{etale-cohomology-lemma-stalk-gives-point} により
$\overline{s}, \overline{s}'$ に付随する。点のより一般的な射
（$S$ の点の特殊化に対応するもの）も存在する。これらは後で記述し、
同型ではない。節 \ref{etale-cohomology-section-specialization} を参照。
\end{remark}

\begin{lemma}
\label{etale-cohomology-lemma-stalk-exact}
$S$ をスキームとし、$\overline{s}$ を $S$ の幾何点とする。
\begin{enumerate}
\item 茎関手
$\textit{PAb}(S_\etale) \to \textit{Ab}$,
$\mathcal{F}  \mapsto  \mathcal{F}_{\overline{s}}$
は完全である。
\item $(\mathcal{F}^\#)_{\overline{s}} = \mathcal{F}_{\overline{s}}$ が、
任意の集合の前層 $\mathcal{F}$（$S_\etale$ 上のもの）に対して成り立つ。
\item 関手
$\textit{Ab}(S_\etale) \to \textit{Ab}$,
$\mathcal{F} \mapsto \mathcal{F}_{\overline{s}}$ は完全である。
\item 同様に、茎関手
$\textit{PSh}(S_\etale) \to \textit{Sets}$ および
$\Sh(S_\etale) \to \textit{Sets}$、
$\mathcal{F} \mapsto \mathcal{F}_{\overline{s}}$ は完全であり
（Categories, 定義 \ref{categories-definition-exact} を参照）、
任意の余極限と可換である。
\end{enumerate}
\end{lemma}

\begin{proof}
直接的な議論による証明を示す前に、この結果は Modules on Sites, 節
\ref{sites-modules-section-stalks} の一般論から従うことに注意する。
実際、補題 \ref{etale-cohomology-lemma-stalk-gives-point} により
$\mathcal{F} \mapsto \mathcal{F}_{\overline{s}}$ は $S$ の小エタール・サイトの
一点から生じる。(1), (2), (3) の直接証明だけを与え、(4) の直接証明は省略する。

\medskip\noindent
$\textit{PAb}(S_\etale)$ 上の関手としての完全性は、有向余極限がすべての余極限および
有限極限と可換であることから形式的に従う。(2) の茎の同一視は、写像
$$
\kappa :
\mathcal{F}_{\overline{s}}
\longrightarrow
(\mathcal{F}^\#)_{\overline{s}}
$$
を介して得られる。これは自然射 $\mathcal{F}\to \mathcal{F}^\#$ により
誘導される。定理 \ref{etale-cohomology-theorem-sheafification} を参照。
この写像がアーベル群の同型であると主張する。単射性を示し、全射性の証明は省略する。

\medskip\noindent
$\sigma\in \mathcal{F}_{\overline{s}}$ とする。あるエタール近傍
$(U, \overline{u})\to (S, \overline{s})$ が存在し、$\sigma$ はある切断
$s \in \mathcal{F}(U)$ の像である。$\kappa(\sigma) = 0$ が
$(\mathcal{F}^\#)_{\overline{s}}$ において成り立つなら、エタール近傍の射
$(U', \overline{u}')\to (U, \overline{u})$ であって、
$s|_{U'}$ が $\mathcal{F}^\#(U')$ において零となるものが存在する。
したがって、エタール被覆 $\{U_i'\to U'\}_{i\in I}$ であって、
$s|_{U_i'}=0$ が $\mathcal{F}(U_i')$ においてすべての $i$ に対し
成り立つものが存在する。
補題 \ref{etale-cohomology-lemma-geometric-lift-to-cover} により、ある $i \in I$ と射
$\overline{u}_i': \overline{s} \to U_i'$ が存在して、
$(U_i', \overline{u}_i') \to (U', \overline{u}')\to (U, \overline{u})$
はいずれもエタール近傍の射となる。したがって $\sigma = 0$ である。
実際、$(U_i', \overline{u}_i') \to (U, \overline{u})$ は
$s|_{U'_i}=0$ を満たすエタール近傍の射である。
これで $\kappa$ が単射であることが示された。

\medskip\noindent
$\textit{Ab}(S_\etale) \to \textit{Ab}$ が完全な関手であることを示すため、
$\textit{Ab}(S_\etale)$ の任意の短完全列
$
0\to \mathcal{F}\to \mathcal{G}\to \mathcal H \to 0.
$
を考える。これは前層の完全列
$$
0 \to \mathcal{F} \to \mathcal{G} \to \mathcal H \to
\mathcal H/^p\mathcal{G} \to 0,
$$
を与える。ここで $/^p$ は $\textit{PAb}(S_\etale)$ における商を表す。
$\overline{s}$ における茎をとると、
$(\mathcal H /^p\mathcal{G})_{\bar{s}} =
(\mathcal H /\mathcal{G})_{\bar{s}} = 0$ であることが分かる。実際、
$\mathcal H/^p\mathcal{G}$ の層化は $0$ である。したがって、
$$
0\to \mathcal{F}_{\overline{s}} \to \mathcal{G}_{\overline{s}} \to
\mathcal{H}_{\overline{s}} \to 0 = (\mathcal H/^p\mathcal{G})_{\overline{s}}
$$
は完全である。実際、茎をとる操作は前層からの関手として完全である。
\end{proof}

\begin{theorem}
\label{etale-cohomology-theorem-exactness-stalks}
$S$ をスキームとする。集合の層の射
$a : \mathcal{F} \to \mathcal{G}$ が単射（それぞれ全射）であるための
必要十分条件は、茎上の写像
$a_{\overline{s}} : \mathcal{F}_{\overline{s}} \to \mathcal{G}_{\overline{s}}$
が $S$ のすべての幾何点に対して単射（それぞれ全射）であることである。
$S_\etale$ 上のアーベル層の列が
完全であるための必要十分条件は、$S$ の幾何点におけるすべての茎上で
完全であることである。
\end{theorem}

\begin{proof}
茎上の完全性が必要であることは補題 \ref{etale-cohomology-lemma-stalk-exact} から従う。
逆を示すには、層の射が全射（それぞれ単射）であるための必要十分条件が、
すべての茎上で全射（それぞれ単射）であることを示せばよい。
全射の場合を証明し、単射の場合の証明は省略する。

\medskip\noindent
$\alpha : \mathcal{F} \to \mathcal{G}$ を、すべての幾何点に対して
$\mathcal{F}_{\overline{s}} \to \mathcal{G}_{\overline{s}}$ が
全射であるような層の射とする。
$U \in \Ob(S_\etale)$
と $s \in \mathcal{G}(U)$ を固定する。各 $u \in U$ に対し、
$\overline{u} \to U$ で $u$ の上にあるものと、エタール近傍
$(V_u , \overline{v}_u) \to (U, \overline{u})$ を選ぶ。ただし
$s|_{V_u} = \alpha(s_{V_u})$ が、ある
$s_{V_u} \in \mathcal{F}(V_u)$ に対して成り立つようにする。
$\alpha$ は茎上で全射なので、これは可能である。このとき
$\{V_u \to U\}_{u \in U}$ はエタール被覆であり、その上で $s$ の制限は
写像 $\alpha$ の像に属する。したがって $\alpha$ は全射である。
Sites, 節 \ref{sites-section-sheaves-injective} を参照。
\end{proof}

\begin{remarks}
\label{etale-cohomology-remarks-enough-points}
幾何学的サイトの点について述べる。
\begin{enumerate}
\item 定理 \ref{etale-cohomology-theorem-exactness-stalks} は、$S_\etale$ の点の族で
$S$ の幾何点によって与えられるもの（補題 \ref{etale-cohomology-lemma-stalk-gives-point}）が保守的であると
述べている。Sites, 定義 \ref{sites-definition-enough-points} を参照。
とくに $S_\etale$ は十分な点をもつ。
\item $\mathcal{F}$ を大エタール・サイト
\label{etale-cohomology-item-stalks-big}
$S$ の層とする。$T \to S$ を大エタール・サイト（$S$ のもの）の対象とし、
$\overline{t}$ を $T$ の幾何点とする。このとき
$\mathcal{F}_{\overline{t}}$ を制限 $\mathcal{F}|_{T_\etale}$ の茎として定義する。
これは $\mathcal{F}$ の $T$ の小エタール・サイトへの制限である。
言い換えれば、$\mathcal{F}$ の茎を、任意のスキーム
$T/S \in \Ob((\Sch/S)_\etale)$ の任意の幾何点において定義できる。
\item このサイトのすべての対象のすべての幾何点を考えることにより、
$S$ の大エタール・サイトも十分な点をもつ（\ref{etale-cohomology-item-stalks-big} を参照）。
\end{enumerate}
\end{remarks}

\noindent
次の補題は初読時には読み飛ばすべきである。

\begin{lemma}
\label{etale-cohomology-lemma-points-small-etale-site}
$S$ をスキームとする。
\begin{enumerate}
\item $p$ を小エタール・サイト $S_\etale$（$S$ のもの）の点で、
関手 $u : S_\etale \to \textit{Sets}$ によって与えられるものとする。
このとき、幾何点 $\overline{s}$（$S$ のもの）が存在し、$p$ は
$S_\etale$ の点で、補題 \ref{etale-cohomology-lemma-stalk-gives-point} において
$\overline{s}$ に付随するものと同型である。
\item $p : \Sh(pt) \to \Sh(S_\etale)$ を $S$ の小エタール・トポスの点とする。
このとき $p$ は $S$ の幾何点から生じる。すなわち、茎関手
$\mathcal{F} \mapsto \mathcal{F}_p$ は、定義 \ref{etale-cohomology-definition-stalk} で
定義した茎関手と同型である。
\end{enumerate}
\end{lemma}

\begin{proof}
Sites, 補題 \ref{sites-lemma-point-site-topos} により、サイトの点と
付随するトポスの点の間には一対一対応があるので、(1) を証明すれば十分である。
Sites, 命題 \ref{sites-proposition-point-limits} により、関手 $u$ は次の性質をもつ。
(a) $u(S) = \{*\}$、(b) $u(U \times_V W) = u(U) \times_{u(V)} u(W)$、および
(c) $\{U_i \to U\}$ がエタール被覆なら
$\coprod u(U_i) \to u(U)$ は全射である。
とくに $U' \subset U$ が開部分スキームなら $u(U') \subset u(U)$ である。
さらに Sites, 補題 \ref{sites-lemma-point-site-topos} により
$u(U) = p^{-1}(h_U^\#)$ と書ける。言い換えれば、$u(U)$ は表現可能層
$h_U$ の茎である。$U = V \amalg W$ なら
$h_U = (h_V \amalg h_W)^\#$ であり、
$u(U) = u(V) \amalg u(W)$ を得る。これは $p^{-1}$ が完全だからである。

\medskip\noindent
$u$ の $S_{Zar}$ への制限を考える。Sites, 例
\ref{sites-example-point-topological} および
\ref{sites-example-point-topology} により、一意な点 $s \in S$ が存在し、
開集合 $S' \subset S$ に対して、$u(S') = \{*\}$ となるのは
$s \in S'$ の場合であり、$u(S') = \emptyset$ となるのは
$s \not \in S'$ の場合である。$\varphi : U \to S$ が $S_\etale$ の対象なら、
$\varphi(U) \subset S$ は開であり（命題 \ref{etale-cohomology-proposition-etale-morphisms} を参照）、
$\{U \to \varphi(U)\}$ はエタール被覆であることに注意する。したがって
$u(U) = \emptyset \Leftrightarrow s \not \in \varphi(U)$ である。

\medskip\noindent
幾何点 $\overline{s} : \overline{s} \to S$ で $s$ の上にあるものを選ぶ。
慣例的な記号の濫用については定義 \ref{etale-cohomology-definition-geometric-point} を参照。
$\varphi : U \to S$ を $S_\etale$ の対象で、$U$ がアフィンであるものとする。
$\varphi$ は分離的であり、ファイバー $U_s$（$\varphi$ の $s$ 上のもの）は
$\Spec(\kappa(s))$ 上のアフィン・スキームであって、有限分離拡大 $k_i$
（$\kappa(s)$ のもの）の有限積のスペクトルである。したがって \'Etale Morphisms, 補題
\ref{etale-lemma-etale-etale-local-technical} を適用すると、
エタール近傍 $(V, \overline{v})$（$(S, \overline{s})$ のもの）であって
$$
U \times_S V = U_1 \amalg \ldots \amalg U_n \amalg W
$$
となるものを得る。ここで $U_i \to V$ は同型であり、$W$ は
$\overline{v}$ の上に点をもたない。したがって
$$
u(U) \times u(V) =
u(U \times_S V) =
u(U_1) \amalg \ldots \amalg u(U_n) \amalg u(W)
$$
であり、もちろん $u(U_i) = u(V)$ でもある。$V$ を少し縮小すれば、
$V$ は $s$ の上にちょうど一点をもつと仮定でき、その結果 $W$ は
$s$ の上に点をもたない。上の議論により $u(W) = \emptyset$ となる。
したがって
$$
u(U) \times u(V) =
u(U_1) \amalg \ldots \amalg u(U_n) =
\coprod\nolimits_{i = 1, \ldots, n} u(V)
$$
を得る。ここで等号は集合 $u(V)$ への写像と両立する。
$u(V) \not = \emptyset$ であることに注意する。実際、$s$ は $V \to S$ の像に属する。
とくに、この状況では $u(U)$ は $n$ 個の元をもつ有限集合である。
実際、左辺はファイバー $u(U)$ を $u(V)$ の任意の元上にもち、
右辺は濃度 $n$ のファイバーをもつ。

\medskip\noindent
極限
$$
\lim_{(V, \overline{v})} u(V)
$$
を、エタール近傍 $(V, \overline{v})$（$\overline{s}$ のもの）の圏上で考える。
$(V, \overline{v})$ で $V$ がアフィンであるものの部分圏上で極限をとっても、
同じ値が得られることは明らかである。前段落を $V$ と $U$ の役割を
入れ替えて適用すると、この場合 $u(V)$ は常に空でない有限集合である。
さらに、補題 \ref{etale-cohomology-lemma-cofinal-etale} によりこの極限は余フィルター的である。
したがって Categories, 節 \ref{categories-section-codirected-limits} により、
この極限は空でない。この極限から元 $x$ を選ぶ。これは、
$x_{V, \overline{v}} \in u(V)$ が各エタール近傍
$(V, \overline{v})$（$(S, \overline{s})$ のもの）に対して得られ、
エタール近傍の任意の射
$\varphi : (V', \overline{v}') \to (V, \overline{v})$ に対して
$u(\varphi)(x_{V', \overline{v}'}) = x_{V, \overline{v}}$ が成り立つことを意味する。

\medskip\noindent
$x$ の選択を用いて、関手的な全単射
$$
c : |U_{\overline{s}}| \longrightarrow u(U)
$$
を $U \in \Ob(S_\etale)$ に対して構成すれば証明は完了する。
$|U_{\overline{s}}|$ の記述については、
補題 \ref{etale-cohomology-lemma-stalk-gives-point} とその証明を参照。まず、
$U$ がアフィンの場合に写像を構成すれば十分であると主張する。
この主張の証明は省略する。$U \to S$ が $S_\etale$ に属し、
$U$ がアフィンであると仮定し、
$\overline{u} : \overline{s} \to U$ を $|U_{\overline{s}}|$ の元とする。
$(V, \overline{v})$ を、$U \times_S V$ が証明の第三段落のように
分解するように選ぶ。このとき対 $(\overline{u}, \overline{v})$ は
$U \times_S V$ の幾何点で $\overline{v}$ の上にあるものを与え、
成分 $U_i$（$U \times_S V$ のもの）の一つを定める。より正確には、切断
$\sigma : V \to U \times_S V$（射影 $\text{pr}_U$ のもの）であって
$(\overline{u}, \overline{v}) = \sigma \circ \overline{v}$ を満たすものが存在する。
$c(\overline{u}) = u(\text{pr}_U)(u(\sigma)(x_{V, \overline{v}})) \in u(U)$
と置く。これが $(V, \overline{v})$ の選択に依存しないことを確かめなければならない。
補題 \ref{etale-cohomology-lemma-cofinal-etale} により、エタール近傍の圏は余フィルター圏である。
したがって、エタール近傍の射
$\varphi : (V', \overline{v}') \to (V, \overline{v})$ と、射影の切断
$\sigma' : V' \to U \times_S V'$ であって
$(\overline{u}, \overline{v'}) = \sigma' \circ \overline{v}'$ を満たすものが与えられたとき、
$u(\sigma')(x_{V', \overline{v}'}) = u(\sigma)(x_{V, \overline{v}})$
となることを示せば十分である。図式
$$
\xymatrix{
V' \ar[d]^{\sigma'} \ar[r]_\varphi & V \ar[d]^\sigma \\
U \times_S V' \ar[r]^{1 \times \varphi} &
U \times_S V
}
$$
を考える。
この図式は可換とは限らない。その理由は、スキーム $V'$ と $V$ が連結とは限らず、
したがって上で $\sigma'$ と $\sigma$ を構成するために用いた分解が
一意とは限らないことである。しかし、構成により
$\sigma \circ \varphi \circ \overline{v}' =
(1 \times \varphi) \circ \sigma' \circ \overline{v}'$
である。したがって、$U \times_S V$ は $S$ 上エタールなので、
開近傍 $V'' \subset V'$（$\overline{v'}$ のもの）が存在し、$V''$ に制限すると
図式は可換になる。Morphisms, 補題 \ref{morphisms-lemma-value-at-one-point} を参照。
つまり、上の図式を
$$
\xymatrix{
V'' \ar[r] \ar[d]^{\sigma'|_{V''}} &
V' \ar[d]^{\sigma'} \ar[r]_\varphi &
V \ar[d]^\sigma \\
U \times_S V'' \ar[r] &
U \times_S V' \ar[r]^{1 \times \varphi} &
U \times_S V
}
$$
へ拡張して、左の正方形と外側の長方形を可換にできる。
$u$ は関手なので、図式のどちらの経路を通っても
$x_{V'', \overline{v}'}$ は $u(U \times_S V)$ の同じ元へ写る。
一方、この元は $x_{V', \overline{v}'}$ と $x_{V, \overline{v}}$、
すなわちそれぞれ $u(V')$ と $u(V)$ の元へ写る。
したがって、求める等式
$u(\sigma')(x_{V', \overline{v}'}) = u(\sigma)(x_{V, \overline{v}})$.
が従う。

\medskip\noindent
同様にして、構成 $c : |U_{\overline{s}}| \to u(U)$ が $U$ に関して
関手的であることが示される。詳細は省略する。最後に、第三段落の結果により
写像 $c$ が全単射であることは明らかである。これで補題の証明が完了する。
\end{proof}







\section{他の位相における点}
\label{etale-cohomology-section-points-topologies}

\noindent
この節では、スキームのエタール・サイト以外のいくつかのサイトについて、
点の存在を簡単に論じる。この節で用いる用語については Sites, 節
\ref{sites-section-sites-enough-points} および Topologies, 節
\ref{topologies-section-procedure} 以下を参照。幾何学的サイトはすべて十分な点をもつ。

\begin{lemma}
\label{etale-cohomology-lemma-points-fppf}
$S$ をスキームとする。次のサイトはすべて十分な点をもつ。
$S_{affine, Zar}$, $S_{Zar}$, $S_{affine, \etale}$, $S_\etale$,
$(\Sch/S)_{Zar}$, $(\textit{Aff}/S)_{Zar}$,
$(\Sch/S)_\etale$, $(\textit{Aff}/S)_\etale$,
$(\Sch/S)_{smooth}$, $(\textit{Aff}/S)_{smooth}$,
$(\Sch/S)_{syntomic}$, $(\textit{Aff}/S)_{syntomic}$,
$(\Sch/S)_{fppf}$ および $(\textit{Aff}/S)_{fppf}$。
\end{lemma}

\begin{proof}
各大サイトに付随するトポスは、サイト $(\textit{Aff}/S)_\tau$ により
定義されるトポスと同値である。Topologies, 補題
\ref{topologies-lemma-affine-big-site-Zariski},
\ref{topologies-lemma-affine-big-site-etale},
\ref{topologies-lemma-affine-big-site-smooth},
\ref{topologies-lemma-affine-big-site-syntomic} および
\ref{topologies-lemma-affine-big-site-fppf} を参照。
サイト $(\textit{Aff}/S)_\tau$ に対する結果は、Deligne の結果
Sites, 補題 \ref{sites-lemma-criterion-points} から直ちに従う。

\medskip\noindent
$S_{Zar}$ に対する結果は明らかである。$S_{affine, Zar}$ に対する結果は
Deligne の結果から従う。$S_\etale$ に対する結果は、定理
\ref{etale-cohomology-theorem-exactness-stalks}（の証明）から従うか、または
Topologies, 補題 \ref{topologies-lemma-alternative} と
$S_{affine, \etale}$ に適用した Deligne の結果から従う。
\end{proof}

\noindent
上の補題は点の存在を保証するが、それらの点がどのようなものかは教えない。
{\it いくつかの}点は次のように明示的に構成できる。
$\overline{s} : \Spec(k) \to S$ を幾何点で、$k$ が代数閉体であるものとする。
関手
$$
u : (\Sch/S)_{fppf} \longrightarrow \textit{Sets},
\quad
u(U) = U(k) = \Mor_S(\Spec(k), U).
$$
$U \mapsto U(k)$ は有限極限と可換することに注意する。実際、
$S(k) = \{\overline{s}\}$ かつ
$(U_1 \times_U U_2)(k) = U_1(k) \times_{U(k)} U_2(k)$ である。さらに、
$\{U_i \to U\}$ が fppf 被覆なら
$\coprod U_i(k) \to U(k)$ は全射である。Sites, 命題
\ref{sites-proposition-point-limits} により、$u$ は点 $p$
（$(\Sch/S)_{fppf}$ のもの）を定め、その茎は
$$
\mathcal{F}_p = \colim_{(U, x)} \mathcal{F}(U)
$$
である。ここで余極限は、通常どおり組 $U \to S$, $x \in U(k)$ 上でとる。
しかし、この圏には始対象 $(\Spec(k), \text{id})$ がある。したがって
$$
\mathcal{F}_p = \mathcal{F}(\Spec(k))
$$
であり、これはそれほど興味深くない。実際、一般にはこれらの点は保守的な
点の族をなさない。次の注意では、より興味深い種類の点を記述する。

\begin{remark}
\label{etale-cohomology-remark-points-fppf-site}
\begin{reference}
これは \cite{Schroeer} で論じられている。
\end{reference}
$S = \Spec(A)$ をアフィン・スキームとする。$(p, u)$ をサイト
$(\textit{Aff}/S)_{fppf}$ の点とする。Sites, 節
\ref{sites-section-points} および \ref{sites-section-construct-points} を参照。
$B = \mathcal{O}_p$ を点 $p$ における構造層の茎とする。次を思い出そう。
$$
B = \colim_{(U, x)} \mathcal{O}(U) =
\colim_{(\Spec(C), x_C)} C
$$
ここで $x_C \in u(\Spec(C))$ である。$\Spec(B)$ が
$(\textit{Aff}/S)_{fppf}$ の対象であり、かつ元
$x_B \in u(\Spec(B))$ で整合系 $x_C$ へ写るものが存在することもある。
この場合、近傍の系は始対象をもち、
$\mathcal{F}_p = \mathcal{F}(\Spec(B))$ が任意の層 $\mathcal{F}$
（$(\textit{Aff}/S)_{fppf}$ 上のもの）に対して従う。関手
$\mathcal{F} \mapsto \mathcal{F}(\Spec(B))$ が
$\Sh((\textit{Aff}/S)_{fppf})$ の点を定めるなら、
$B$ は局所 $A$-代数でなければならず、任意の忠実平坦かつ有限表示な環準同型
$B \to B'$ に対して切断 $B' \to B$ が存在することは容易に分かる。
逆に、このような任意の $A$-代数 $B$ に対し、関手
$\mathcal{F} \mapsto \mathcal{F}(\Spec(B))$ は一点の茎関手である。
詳細は省略する。サイト $(\textit{Aff}/S)_{fppf}$ の一般の点が
どのようなものかは明らかでない。
\end{remark}











\section{アーベル層の台}
\label{etale-cohomology-section-support}

\noindent
まず局所切断の台について述べる。

\begin{lemma}
\label{etale-cohomology-lemma-support-subsheaf-final}
$S$ をスキームとする。$\mathcal{F}$ を $S$ のエタール・トポスの
終対象の部分層とする（Sites, 例 \ref{sites-example-singleton-sheaf} を参照）。
このとき一意な開集合 $W \subset S$ で
$\mathcal{F} = h_W$ を満たすものが存在する。
\end{lemma}

\begin{proof}
この条件は、$\mathcal{F}(U)$ が一元集合または空集合であることを、
すべての $\varphi : U \to S$（$\Ob(S_\etale)$ に属するもの）に対して意味する。
とくに局所切断は常に貼り合わさる。$\mathcal{F}(U) \not = \emptyset$ なら、
$\mathcal{F}(\varphi(U)) \not = \emptyset$ である。実際、
$\{\varphi : U \to \varphi(U)\}$ は被覆である。したがって
$W = \bigcup_{\varphi : U \to S, \mathcal{F}(U) \not = \emptyset} \varphi(U)$.
ととればよい。
\end{proof}

\begin{lemma}
\label{etale-cohomology-lemma-zero-over-image}
$S$ をスキームとする。$\mathcal{F}$ を $S_\etale$ 上のアーベル層とし、
$\sigma \in \mathcal{F}(U)$ を局所切断とする。次を満たす開部分集合
$W \subset U$ が存在する。
\begin{enumerate}
\item $W \subset U$ は $U$ の Zariski 開部分集合で
$\sigma|_W = 0$ を満たすもののうち最大である。
\item すべての $\varphi : V \to U$（$S_\etale$ に属するもの）に対して
$$
\sigma|_V = 0 \Leftrightarrow \varphi(V) \subset W,
$$
が成り立つ。
\item すべての幾何点 $\overline{u}$（$U$ のもの）に対して
$$
(U, \overline{u}, \sigma) = 0\text{（}\mathcal{F}_{\overline{s}}\text{ において）}
\Leftrightarrow
\overline{u} \in W
$$
が成り立つ。ここで $\overline{s} = (U \to S) \circ \overline{u}$ である。
\end{enumerate}
\end{lemma}

\begin{proof}
$\mathcal{F}$ はエタール位相に関する層なので、$\mathcal{F}$ の $U_{Zar}$ への
制限は $U$ 上の Zariski 位相に関する層である。したがって性質 (1) をもつ
Zariski 開集合 $W$ が存在する。Modules, 補題
\ref{modules-lemma-support-section-closed} を参照。
$\varphi : V \to U$ を $S_\etale$ の矢印とする。
$\varphi(V) \subset U$ は開部分集合であり、
$\{V \to \varphi(V)\}$ はエタール被覆であることに注意する。したがって
$\sigma|_V = 0$ なら、$\mathcal{F}$ の層条件により
$\sigma|_{\varphi(V)} = 0$ である。これで (2) が示された。
(3) を示すには、$(U, \overline{u}, \sigma)$ が
$\mathcal{F}_{\overline{s}}$ の零元を定めるなら
$\overline{u} \in W$ であることを示さなければならない。
実際、この仮定は、エタール近傍の射
$(V, \overline{v}) \to (U, \overline{u})$ であって
$\sigma|_V = 0$ を満たすものが存在することを意味する。したがって (2) により
$V \to U$ は $W$ の中へ写り、ゆえに $\overline{u} \in W$ である。
\end{proof}

\noindent
$S$ をスキームとし、$s \in S$ とする。$\mathcal{F}$ を
$S_\etale$ 上の層とする。注意 \ref{etale-cohomology-remark-map-stalks} により、
層 $\mathcal{F}$ の茎の同型類は、$s$ の上にある幾何点においてよく定義される。

\begin{definition}
\label{etale-cohomology-definition-support}
$S$ をスキームとし、$\mathcal{F}$ を $S_\etale$ 上のアーベル層とする。
\begin{enumerate}
\item {\it $\mathcal{F}$ の台}とは、点 $s \in S$ であって、
$\mathcal{F}_{\overline{s}} \not = 0$ が任意の（ある）幾何点
$\overline{s}$（$s$ の上にあるもの）に対して成り立つもの全体の集合である。
\item $\sigma \in \mathcal{F}(U)$ を切断とする。{\it $\sigma$ の台}とは、
閉部分集合 $U \setminus W$ のことである。ここで $W \subset U$ は
$U$ の最大の開部分集合で、その上で $\sigma$ の制限が零となるものである
（補題 \ref{etale-cohomology-lemma-zero-over-image} を参照）。
\end{enumerate}
\end{definition}

\noindent
一般に、アーベル層の台は閉とは限らない。例えば
$S = \mathbf{A}^1_{\mathbf{C}}$ と仮定する。
$i_t : \Spec(\mathbf{C}) \to S$ を点 $t \in \mathbf{C}$ の包含とする。
後で $\mathcal{F}_t = i_{t, *}(\mathbf{Z}/2\mathbf{Z})$ が、
台がちょうど $\{t\}$ であるアーベル層だと分かる。節
\ref{etale-cohomology-section-closed-immersions} を参照。このとき
$$
\bigoplus\nolimits_{n \in \mathbf{N}} \mathcal{F}_n
$$
は台 $\{1, 2, 3, \ldots\} \subset S$ をもつアーベル層である。
これは茎をとる操作が余極限と可換だからである。補題
\ref{etale-cohomology-lemma-stalk-exact} を参照。したがって、これは台が閉でない
アーベル層の例である。層と切断の台に関する基本的な事実をいくつか挙げる。

\begin{lemma}
\label{etale-cohomology-lemma-support-section-closed}
$S$ をスキームとする。$\mathcal{F}$ を $S_\etale$ 上のアーベル層とし、
$U \in \Ob(S_\etale)$ および $\sigma \in \mathcal{F}(U)$ とする。
\begin{enumerate}
\item $\sigma$ の台は $U$ において閉である。
\item $\sigma + \sigma'$ の台は
$\sigma, \sigma' \in \mathcal{F}(U)$ の台の和集合に含まれる。
\item $\varphi : \mathcal{F} \to \mathcal{G}$ が
$S_\etale$ 上のアーベル層の射なら、$\varphi(\sigma)$ の台は
$\sigma \in \mathcal{F}(U)$ の台に含まれる。
\item $\mathcal{F}$ の台は、$\mathcal{F}$ のすべての局所切断の台の像の和集合である。
\item $\mathcal{F} \to \mathcal{G}$ が全射なら、$\mathcal{G}$ の台は
$\mathcal{F}$ の台の部分集合である。
\item $\mathcal{F} \to \mathcal{G}$ が単射なら、$\mathcal{F}$ の台は
$\mathcal{G}$ の台の部分集合である。
\end{enumerate}
\end{lemma}

\begin{proof}
(1) は定義により成り立つ。(2) と (3) は、$\mathcal{F}$ と
$\mathcal{G}$ の $U_{Zar}$ への制限について成り立つため従う。
Modules, 補題 \ref{modules-lemma-support-section-closed} を参照。
(4) は補題 \ref{etale-cohomology-lemma-zero-over-image} の (3) から直ちに従う。
(5) と (6) は他の各部分から従う。
\end{proof}

\begin{lemma}
\label{etale-cohomology-lemma-support-sheaf-rings-closed}
$S_\etale$ 上の環の層の台は閉である。
\end{lemma}

\begin{proof}
これは、われわれの規約では環が $0$ であるための必要十分条件が
$1 = 0$ であり、したがって環の層の台が単位切断の台だからである。
\end{proof}




\section{Hensel 環}
\label{etale-cohomology-section-henselian-ring}

\noindent
まず Stacks プロジェクトですでに何度も用いられてきた定理を述べる。
この結果には多くの版があるが、ここでは代数的な版だけを述べる。

\begin{theorem}
\label{etale-cohomology-theorem-quasi-finite-etale-locally}
$A\to B$ を有限型の環準同型、$\mathfrak p \subset A$ を素イデアルとする。
このとき、エタールな環準同型 $A \to A'$ と、素イデアル
$\mathfrak p' \subset A'$ で $\mathfrak p$ の上にあり、次を満たすものが存在する。
\begin{enumerate}
\item
$\kappa(\mathfrak p) = \kappa(\mathfrak p')$,
\item
$ B \otimes_A A' = B_1\times \ldots \times B_r \times C$,
\item
$ A'\to B_i$ は有限であり、一意な素イデアル $q_i\subset B_i$ で
$\mathfrak p'$ の上にあるものが存在する。
\item ファイバー $\Spec(C \otimes_{A'} \kappa(\mathfrak p'))$
（$C$ の $\mathfrak p'$ 上のもの）の既約成分はすべて次元が少なくとも 1 である。
\end{enumerate}
\end{theorem}

\begin{proof}
Algebra, 補題 \ref{algebra-lemma-etale-makes-quasi-finite-finite} または
\cite[Th\'eor\`eme 18.12.1]{EGA4} を参照。スキームの射によるさまざまな版については
More on Morphisms, 節 \ref{more-morphisms-section-etale-localization} を参照。
\end{proof}

\noindent
Hensel の補題を思い出そう。この補題には多くの版がある。二つを挙げる。
\begin{enumerate}
\item[(f)] $f\in \mathbf{Z}_p[T]$ がモニックで、
$f \bmod p = g_0 h_0$ かつ $gcd(g_0, h_0) = 1$ なら、
$f$ は $f = gh$ と分解し、$\bar g = g_0$ および $\bar h = h_0$ を満たす。
\item[(r)] $f \in \mathbf{Z}_p[T]$ とモニックな $a_0 \in \mathbf{F}_p$ が
$\bar f(a_0) =0$ かつ $\bar f'(a_0) \neq 0$ を満たすなら、
$a \in \mathbf{Z}_p$ で $f(a) = 0$ かつ $\bar a = a_0$ を満たすものが存在する。
\end{enumerate}
いずれの版も正しい（後でこれを見る）。第一の版は互いに素な因子への分解の
持ち上げを求め、第二の版は極大イデアルを法とする単根の持ち上げを求める。
一般の局所環に対してこれらの条件を課すことは互いに同値であり、さらに多くの
他の条件とも同値であることが分かる。根の持ち上げ性はしばしば最も確認しやすいので、
これを Hensel 局所環の定義として用いる。

%10.01.09
\begin{definition}
\label{etale-cohomology-definition-henselian}
（Algebra, 定義 \ref{algebra-definition-henselian} を参照。）
局所環 $(R, \mathfrak m, \kappa)$ が {\it Hensel 的}であるとは、
任意のモニックな $f \in R[T]$ と任意の $a_0 \in \kappa$ で
$\bar f(a_0) = 0$ および $\bar f'(a_0) \neq 0$ を満たすものに対して、
$a \in R$ が存在し、$f(a) = 0$ および
$a \bmod \mathfrak m = a_0$ が成り立つことをいう。
\end{definition}

\noindent
Hensel 局所環の代表例として完備局所環を念頭に置くとよい。
完備局所環とは、局所環 $(R, \mathfrak m)$ で
$R \cong \lim_n R/\mathfrak m^n$ を満たすもの、すなわち
$\mathfrak m$-進位相に関して完備かつ分離的なものだった
（Algebra, 定義 \ref{algebra-definition-complete-local-ring}）。

\begin{theorem}
\label{etale-cohomology-theorem-hensel}
完備局所環は Hensel 的である。
\end{theorem}

\begin{proof}
ニュートン法による。Algebra, 補題 \ref{algebra-lemma-complete-henselian} を参照。
\end{proof}

\begin{theorem}
\label{etale-cohomology-theorem-henselian}
$(R, \mathfrak m, \kappa)$ を局所環とする。次は同値である。
\begin{enumerate}
\item $R$ は Hensel 的である。
\item 任意の $f\in R[T]$ と、任意の分解 $\bar f = g_0 h_0$
（$\kappa[T]$ におけるもので $\gcd(g_0, h_0)=1$ を満たすもの）に対して、
分解 $f = gh$（$R[T]$ におけるもの）で
$\bar g = g_0$ および $\bar h = h_0$ を満たすものが存在する。
\item 任意の有限 $R$-代数 $S$ は、$R$ 上有限な局所環の有限積と同型である。
\item 任意の有限型 $R$-代数 $A$ は積 $A \cong A' \times C$ と同型である。
ここで $A' \cong A_1 \times \ldots \times A_r$ は有限局所 $R$-代数の積であり、
$C \otimes_R \kappa$ のすべての既約成分は次元が少なくとも 1 である。
\item $A$ がエタールな $R$-代数で、$\mathfrak n$ が $A$ の極大イデアルで
$\mathfrak m$ の上にあり、$\kappa \cong A/\mathfrak n$ を満たすなら、
同型 $\varphi : A \cong R \times A'$ であって
$\varphi(\mathfrak n) = \mathfrak m \times A' \subset R \times A'$ を満たすものが存在する。
\end{enumerate}
\end{theorem}

\begin{proof}
これは Algebra, 補題 \ref{algebra-lemma-characterize-henselian} の結果の一部に
すぎない。上の (5) は Algebra, 補題 \ref{algebra-lemma-characterize-henselian} の
(8) に対応するが、少し異なる形で述べている。
\end{proof}

\begin{lemma}
\label{etale-cohomology-lemma-finite-over-henselian}
$R$ が Hensel 的で、$A$ が有限 $R$-代数なら、$A$ は Hensel 局所環の有限積である。
\end{lemma}

\begin{proof}
Algebra, 補題 \ref{algebra-lemma-finite-over-henselian} を参照。
\end{proof}

\begin{definition}
\label{etale-cohomology-definition-strictly-henselian}
局所環 $R$ が {\it 強 Hensel 的}であるとは、それが Hensel 的であり、
その剰余体が分離閉であることをいう。
\end{definition}

\begin{example}
\label{etale-cohomology-example-powerseries}
$R = \mathbf{C}[[t]]$ の場合、エタールな $R$-代数は、自明な拡大
$R \to R$ と拡大 $R \to R[X, X^{-1}]/(X^n-t)$ の有限積である。
後者は開集合 $D(t) \subset \Spec(R)$ を経由するので、任意のエタール被覆は
被覆 $\{\text{id} : \Spec(R) \to \Spec(R)\}$ により細分される。
以下で、これは Hensel 環のスペクトルのエタール被覆に関するかなり一般的な
事実であることを見る。このことから、強 Hensel 環のスペクトルの高次
エタール・コホモロジーが零であることが分かる。
\end{example}

\begin{theorem}
\label{etale-cohomology-theorem-henselization}
$(R, \mathfrak m, \kappa)$ を局所環とし、
$\kappa\subset\kappa^{sep}$ を分離代数閉包とする。標準的な平坦局所環準同型
$R \to R^h \to R^{sh}$ であって、次を満たすものが存在する。
\begin{enumerate}
\item $R^h$, $R^{sh}$ はエタールな $R$-代数のフィルター余極限である。
\item $R^h$ は Hensel 的であり、$R^{sh}$ は強 Hensel 的である。
\item $\mathfrak m R^h$（それぞれ $\mathfrak m R^{sh}$）は
$R^h$（それぞれ $R^{sh}$）の極大イデアルである。
\item $\kappa = R^h/\mathfrak m R^h$ であり、
$\kappa^{sep} = R^{sh}/\mathfrak m R^{sh}$ は $\kappa$ の拡大として成り立つ。
\end{enumerate}
\end{theorem}

\begin{proof}
$R^h$ と $R^{sh}$ の構造は Algebra, 補題
\ref{algebra-lemma-henselization} および
\ref{algebra-lemma-strict-henselization} で記述されている。
\end{proof}

\noindent
定理 \ref{etale-cohomology-theorem-henselization} で構成された環を、それぞれ局所環 $R$ の
{\it Hensel 化}および {\it 強 Hensel 化}という。Algebra, 定義
\ref{algebra-definition-henselization} を参照。$R$ の多くの性質はその（強）
Hensel 化に反映される。More on Algebra, 節
\ref{more-algebra-section-permanence-henselization} を参照。




\section{構造層の茎}
\label{etale-cohomology-section-stalks-structure-sheaf}

\noindent
本節では、幾何点における構造層の茎を、対応する「通常の」点における
局所環の強 Hensel 化と同一視する。

\begin{lemma}
\label{etale-cohomology-lemma-describe-etale-local-ring}
\begin{slogan}
スキームのエタール位相における構造層の茎は強 Hensel 化である。
\end{slogan}
$S$ をスキームとする。$\overline{s}$ を $S$ の幾何点で
$s \in S$ 上にあるものとする。$\kappa = \kappa(s)$ とおき、
$\kappa \subset \kappa^{sep} \subset \kappa(\overline{s})$ により
$\kappa$ の $\kappa(\overline{s})$ における分離代数閉包を表す。
このとき標準的な同一視
$$
(\mathcal{O}_{S, s})^{sh}
\cong
(\mathcal{O}_S)_{\overline{s}}
$$
がある。ここで左辺は定理 \ref{etale-cohomology-theorem-henselization} で記述した
局所環 $\mathcal{O}_{S, s}$ の強 Hensel 化であり、右辺は
構造層 $\mathcal{O}_S$ の $S_\etale$ 上の幾何点
$\overline{s}$ における茎である。
\end{lemma}

\begin{proof}
$\Spec(A) \subset S$ を $s$ のアフィン近傍とし、
$\mathfrak p \subset A$ を $s$ に対応する素イデアルとする。
これらの選択のもとで、標準同型
$\mathcal{O}_{S, s} = A_{\mathfrak p}$ および
$\kappa(s) = \kappa(\mathfrak p)$ がある。したがって
$\kappa(\mathfrak p) \subset \kappa^{sep} \subset \kappa(\overline{s})$.
次を思い出そう。
$$
(\mathcal{O}_S)_{\overline{s}} =
\colim_{(U, \overline{u})} \mathcal{O}(U)
$$
ここで余極限は $(S, \overline{s})$ のエタール近傍全体にわたる。
エタール近傍 $(U, \overline{u})$ のうち、$U$ がアフィンであり、
$U \to S$ が $\Spec(A)$ を経由するものは余終系をなす。言い換えると、
$$
(\mathcal{O}_S)_{\overline{s}} = \colim_{(B, \mathfrak q, \phi)} B
$$
であることが分かる。ここで余極限は、エタール $A$-代数 $B$ であって、
素イデアル $\mathfrak q$（$\mathfrak p$ 上にあるもの）と
$\kappa(\mathfrak p)$-代数写像
$\phi : \kappa(\mathfrak q) \to \kappa(\overline{s})$ を備えたもの全体にわたる。
$\kappa(\mathfrak q)$ は $\kappa(\mathfrak p)$ 上有限分離的なので、
$\phi$ の像は $\kappa^{sep}$ に含まれることに注意する。
以上の翻訳により、この補題の主張は Algebra, 補題
\ref{algebra-lemma-strict-henselization-different} の主張と同値である。
\end{proof}

\begin{definition}
\label{etale-cohomology-definition-etale-local-rings}
$S$ をスキームとし、$\overline{s}$ を $S$ の幾何点で
点 $s \in S$ 上にあるものとする。
\begin{enumerate}
\item $S$ の $\overline{s}$ における {\it エタール局所環}とは、
構造層 $\mathcal{O}_S$ の $S_\etale$ 上の $\overline{s}$ における
茎である。これを {\it $\mathcal{O}_{S, s}$ の強 Hensel 化}
（幾何点 $\overline{s}$ に関するもの）と呼ぶこともある。
記法は $\mathcal{O}_{S, \overline{s}}^{sh}$ とする。
\item {\it $\mathcal{O}_{S, s}$ の Hensel 化}とは、
$S$ の $s$ における局所環の Hensel 化である。
Algebra, 定義 \ref{algebra-definition-henselization} および
定理 \ref{etale-cohomology-theorem-henselization} を参照せよ。
記法は $\mathcal{O}_{S, s}^h$ とする。
\item $S$ の $\overline{s}$ における {\it 強 Hensel 化}とは、
スキーム $\Spec(\mathcal{O}_{S, \overline{s}}^{sh})$ のことである。
\item $S$ の $s$ における {\it Hensel 化}とは、
スキーム $\Spec(\mathcal{O}_{S, s}^h)$ のことである。
\end{enumerate}
\end{definition}

\noindent
$f : T \to S$ をスキームの射とする。$\overline{t}$ を $T$ の
幾何点とし、その像 $\overline{s}$ は $S$ にあるとする。
$t \in T$ および $s \in S$ をそれぞれの像とする。
このとき、標準的な可換図式
$$
\xymatrix{
\Spec(\mathcal{O}^{sh}_{T, \overline{t}}) \ar[r] \ar[d] &
\Spec(\mathcal{O}^h_{T, t}) \ar[r] \ar[d] &
T \ar[d]^f \\
\Spec(\mathcal{O}^{sh}_{S, \overline{s}}) \ar[r] &
\Spec(\mathcal{O}^h_{S, s}) \ar[r] &
S
}
$$
を得る。これは $T$ と $S$ の Hensel 化および強 Hensel 化からなる。
$t$ と $s$ のアフィン近傍を選び、
Algebra, 補題 \ref{algebra-lemma-henselian-functorial-improve} および
\ref{algebra-lemma-strictly-henselian-functorial-improve} で与えられる
（強）Hensel 化の関手性を用いれば証明できる。

\begin{lemma}
\label{etale-cohomology-lemma-describe-henselization}
$S$ をスキームとし、$s \in S$ とする。このとき
$$
\mathcal{O}_{S, s}^h =
\colim_{(U, u)} \mathcal{O}(U)
$$
である。ここで余極限は、エタール近傍 $(U, u)$、すなわち
$(S, s)$ の近傍で $\kappa(s) = \kappa(u)$ を満たすものからなる
フィルター圏にわたる。
\end{lemma}

\begin{proof}
この補題は More on Morphisms, 補題
\ref{more-morphisms-lemma-describe-henselization} の複製である。
\end{proof}

\begin{remark}
\label{etale-cohomology-remark-henselization-Noetherian}
$S$ をスキームとし、$s \in S$ とする。$S$ が局所 Noether
スキームならば、$\mathcal{O}_{S, s}^h$ も Noether 環であり、
両者の完備化は同じである：
$$
\widehat{\mathcal{O}_{S, s}} \cong \widehat{\mathcal{O}_{S, s}^h}.
$$
とくに
$\mathcal{O}_{S, s} \subset
\mathcal{O}_{S, s}^h \subset
\widehat{\mathcal{O}_{S, s}}$.
$\mathcal{O}_{S, s}$ の Hensel 化は一般にその完備化よりはるかに小さく、
しかも完備化の多くの性質を受け継ぐ。たとえば、
$\mathcal{O}_{S, s}$ が被約ならば $\mathcal{O}_{S, s}^h$ も被約であるが、
完備化については一般にこれは成り立たない。
今後ここに参照を挿入する。
\end{remark}

\begin{lemma}
\label{etale-cohomology-lemma-etale-site-locally-ringed}
$S$ をスキームとする。小エタールサイト $S_\etale$ は、
その構造層 $\mathcal{O}_S$ を備えると局所環付きサイトである。
Modules on Sites, 定義 \ref{sites-modules-definition-locally-ringed} を参照せよ。
\end{lemma}

\begin{proof}
これは、茎
$(\mathcal{O}_S)_{\overline{s}} = \mathcal{O}^{sh}_{S, \overline{s}}$
が局所環であり、かつ $S_\etale$ が十分多くの点をもつことから従う。
補題 \ref{etale-cohomology-lemma-describe-etale-local-ring}、
定理 \ref{etale-cohomology-theorem-exactness-stalks}、および
注意 \ref{etale-cohomology-remarks-enough-points} を参照せよ。
これにより小エタールサイトが局所環付きとなることについては、
Modules on Sites, 補題 \ref{sites-modules-lemma-locally-ringed-stalk} および
\ref{sites-modules-lemma-ringed-stalk-not-zero} を参照せよ。
\end{proof}






\section{小エタール・トポスの関手性}
\label{etale-cohomology-section-functoriality}

\noindent
これまではエタールサイトの関手性、すなわちスキームの射が与えられたときに
何が起こるかを論じてこなかった。厳密で形式的な議論は
Topologies, 節 \ref{topologies-section-etale} にある。
本節および以下の数節では、この内容を小エタールサイトの状況に限って
簡潔に論じる。

\medskip\noindent
$f : X \to Y$ をスキームの射とする。このとき関手
\begin{equation}
\label{etale-cohomology-equation-functorial}
u : Y_\etale \longrightarrow X_\etale, \quad
V/Y \longmapsto X \times_Y V/X.
\end{equation}
を得る。この関手は次の重要な性質をもつ。
\begin{enumerate}
\item $u(\text{終対象}) = \text{終対象}$,
\item $u$ はファイバー積を保つ,
\item $\{V_j \to V\}$ が $Y_\etale$ の被覆ならば、
$\{u(V_j) \to u(V)\}$ は $X_\etale$ の被覆である。
\end{enumerate}
これらはいずれも容易に確認できる（省略する）。その帰結として、
{\it サイトの射}と呼ばれるもの
$$
f_{small} : X_\etale \longrightarrow Y_\etale,
$$
を得る。Sites, 定義 \ref{sites-definition-morphism-sites} および
Sites, 命題 \ref{sites-proposition-get-morphism} を参照せよ。
エタール層やエタール・コホモロジーを扱うために、この抽象概念を
詳しく知る必要はない。通常は、エタール層上に関手
$f_{small, *}$（順像）および $f_{small}^{-1}$（逆像）が存在することと、
それらのいくつかの簡単な性質を知れば十分である。
以下の数節でこれらの性質を論じるが、証明にはより抽象的な内容を
参照することがある。そのような抽象的状況で証明するのが自然なことが
多いからである。


\section{順像}
\label{etale-cohomology-section-direct-image}

\noindent
前層の順像を定義しよう。

\begin{definition}
\label{etale-cohomology-definition-direct-image-presheaf}
$f: X\to Y$ をスキームの射とする。
$\mathcal{F} $ を $X_\etale$ 上の集合の前層とする。
$\mathcal{F}$ の（$f$ による）{\it 順像}は
$$
f_*\mathcal{F} : Y_\etale^{opp} \longrightarrow \textit{Sets}, \quad
(V/Y) \longmapsto \mathcal{F}(X \times_Y V/X).
$$
である。$f_* = f_{small, *}$ と書くこともあるが、これはほかの
順像関手（通常の Zariski 順像や $f_{big, *}$ など）と区別するためである。
\end{definition}

\noindent
エタール射の基底変換は再びエタールなので、これはよく定義された
エタール前層である。より圏論的に言えば、$f_*\mathcal{F}$ は関手の合成
$\mathcal{F} \circ u$ であり、ここで $u$ は式
(\ref{etale-cohomology-equation-functorial}) のものである。これにより、この構成が前層
$\mathcal{F}$ に関して関手的であることは明らかであり、したがって関手
$$
f_* = f_{small, *} :
\textit{PSh}(X_\etale)
\longrightarrow
\textit{PSh}(Y_\etale)
$$
を得る。$\mathcal{F}$ がアーベル群の前層ならば、
$f_*\mathcal{F}$ もアーベル群の前層であり、
$$
f_* = f_{small, *} :
\textit{PAb}(X_\etale)
\longrightarrow
\textit{PAb}(Y_\etale)
$$
を得ることに注意する。これは上と同じく、まったく同じ規則で定義される。

\begin{remark}
\label{etale-cohomology-remark-direct-image-sheaf}
層の順像は層であると主張する。実際、
$\{V_j \to V\}$ が $Y_\etale$ のエタール被覆ならば、
$\{X \times_Y V_j \to X \times_Y V\}$ は $X_\etale$ の
エタール被覆である。したがって、$\mathcal{F}$ に対する
$\{X \times_Y V_i \to X \times_Y V\}$ に関する層条件は、
$f_*\mathcal{F}$ に対する $\{V_i \to V\}$ に関する層条件と同値である。
ゆえに $\mathcal{F}$ が層ならば、$f_*\mathcal{F}$ も層である。
\end{remark}

\begin{definition}
\label{etale-cohomology-definition-direct-image-sheaf}
$f: X\to Y$ をスキームの射とする。
$\mathcal{F} $ を $X_\etale$ 上の集合の層とする。
$\mathcal{F}$ の（$f$ による）{\it 順像}は
$$
f_*\mathcal{F} : Y_\etale^{opp} \longrightarrow \textit{Sets}, \quad
(V/Y) \longmapsto \mathcal{F}(X \times_Y V/X)
$$
である。これは注意 \ref{etale-cohomology-remark-direct-image-sheaf} により層である。
$f_* = f_{small, *}$ と書くこともあるが、これはほかの順像関手
（通常の Zariski 順像や $f_{big, *}$ など）と区別するためである。
\end{definition}

\noindent
上とまったく同じ議論が成り立ち、関手
$$
f_* = f_{small, *} :
\Sh(X_\etale)
\longrightarrow
\Sh(Y_\etale)
$$
および
$$
f_* = f_{small, *} :
\textit{Ab}(X_\etale)
\longrightarrow
\textit{Ab}(Y_\etale)
$$
を得る。これらも{\it 順像}と呼ばれる。

\medskip\noindent
アーベル層上の関手 $f_*$ は左完全である。（アーベル圏間の関手が
左完全であることの意味については Homology, 節
\ref{homology-section-functors} を参照せよ。）実際、
$0 \to \mathcal{F}_1 \to \mathcal{F}_2 \to \mathcal{F}_3$
が $X_\etale$ 上で完全ならば、任意の
$U/X \in \Ob(X_\etale)$
に対してアーベル群の列
$0 \to \mathcal{F}_1(U) \to \mathcal{F}_2(U) \to \mathcal{F}_3(U)$
は完全である。したがって、任意の $V/Y \in \Ob(Y_\etale)$ に対して
アーベル群の列
$0 \to f_*\mathcal{F}_1(V) \to f_*\mathcal{F}_2(V) \to f_*\mathcal{F}_3(V)$
は完全である。これは上の列で $U = X \times_Y V$ としたものだからである。

\begin{definition}
\label{etale-cohomology-definition-higher-direct-images}
$f: X \to Y$ をスキームの射とする。
右導来関手 $\{R^pf_*\}_{p \geq 1}$、すなわち
$f_* : \textit{Ab}(X_\etale) \to \textit{Ab}(Y_\etale)$
の右導来関手を{\it 高次順像}と呼ぶ。
\end{definition}

\noindent
高次順像とその導来圏版については、
（将来ここに参照を挿入）でさらに詳しく論じる。



\section{逆像}
\label{etale-cohomology-section-inverse-image}

\noindent
本節では、小エタールサイト上の層の引き戻しについて簡潔に論じる。
その厳密な構成は Topologies, 節 \ref{topologies-section-etale} にある。

\begin{definition}
\label{etale-cohomology-definition-inverse-image}
$f: X\to Y$ をスキームの射とする。{\it 逆像}、または
{\it 引き戻し}\footnote{集合の層またはアーベル群の層の引き戻しには
$f^{-1}$ という記号を用い、環付きサイト／トポスの射に沿う加群の層の
引き戻しには $f^*$ を用いる。}関手とは、関手
$$
f^{-1} = f_{small}^{-1} :
\Sh(Y_\etale)
\longrightarrow
\Sh(X_\etale)
$$
および
$$
f^{-1} = f_{small}^{-1} :
\textit{Ab}(Y_\etale)
\longrightarrow
\textit{Ab}(X_\etale)
$$
であり、これらは $f_* = f_{small, *}$ の左随伴である。したがって
$f^{-1}$ は次の性質によって特徴づけられる。
$$
\Hom_{{\Sh(X_\etale)}} (f^{-1}\mathcal{G}, \mathcal{F})
=
\Hom_{\Sh(Y_\etale)} (\mathcal{G}, f_*\mathcal{F})
$$
この等式は、任意の $\mathcal{F} \in \Sh(X_\etale)$ および
$\mathcal{G} \in \Sh(Y_\etale)$ に対して関手的に成り立つ。同様に
$$
\Hom_{{\textit{Ab}(X_\etale)}} (f^{-1}\mathcal{G}, \mathcal{F})
=
\Hom_{\textit{Ab}(Y_\etale)} (\mathcal{G}, f_*\mathcal{F})
$$
を得る。ここで $\mathcal{F} \in \textit{Ab}(X_\etale)$ および
$\mathcal{G} \in \textit{Ab}(Y_\etale)$ である。
\end{definition}

\noindent
このような随伴が存在することは自明ではない。一方、それはかなり一般的な
状況で存在する。下の注意 \ref{etale-cohomology-remark-functoriality-general} を参照せよ。
一般論によれば、$f^{-1}\mathcal{G}$ は前層
\begin{equation}
\label{etale-cohomology-equation-pullback}
U/X
\longmapsto
\colim_{U \to X \times_Y V} \mathcal{G}(V/Y)
\end{equation}
に付随する層である。ここで余極限は、対
$(V/Y, \varphi : U/X \to X \times_Y V/X)$ の圏にわたって取る。
これを見るには、Sites, 命題 \ref{sites-proposition-get-morphism} を
式 (\ref{etale-cohomology-equation-functorial}) の関手 $u$ に適用し、Sites, 節
\ref{sites-section-continuous-functors} および
\ref{sites-section-functoriality-PSh} にある
$u_s = (u_p\ )^\#$ の記述を用いればよい。引き戻しの基本的性質の
いくつかを証明するため、ときどきこの公式を用いる。

\begin{lemma}
\label{etale-cohomology-lemma-stalk-pullback}
$f : X \to Y$ をスキームの射とする。
\begin{enumerate}
\item 関手
$f^{-1} : \textit{Ab}(Y_\etale) \to \textit{Ab}(X_\etale)$
は完全である。
\item 関手
$f^{-1} : \Sh(Y_\etale) \to \Sh(X_\etale)$
は完全である。すなわち、有限極限および有限余極限と可換である。
Categories, 定義 \ref{categories-definition-exact} を参照せよ。
\item $\overline{x} \to X$ を幾何学的点とし、
$\mathcal{G}$ を $Y_\etale$ 上の層とする。このとき標準的な同一視
$$
(f^{-1}\mathcal{G})_{\overline{x}} = \mathcal{G}_{\overline{y}}.
$$
がある。ここで $\overline{y} = f \circ \overline{x}$ である。
\item 任意のエタール射 $V \to Y$ に対して
$f^{-1}h_V = h_{X \times_Y V}$ である。
\end{enumerate}
\end{lemma}

\begin{proof}
集合の層における $f^{-1}$ の完全性は、Sites, 命題
\ref{sites-proposition-get-morphism} を式
(\ref{etale-cohomology-equation-functorial}) の関手 $u$ に適用した帰結である。
実際、引き戻しの完全性はトポス（または、そう呼びたければサイト）の
射の定義の一部である。したがって (2) が成り立つ。(1) はこれから従う。
実際、アーベル層 $\mathcal{G}$ が $Y_\etale$ 上に与えられたとき、
$f^{-1}\mathcal{F}$ の台となる集合の層は、$f^{-1}$ を $\mathcal{F}$ の
台となる集合の層に施したものと同じだからである。Sites, 節
\ref{sites-section-sheaves-algebraic-structures} を参照せよ。
Modules on Sites, 補題 \ref{sites-modules-lemma-flat-pullback-exact}
も参照せよ。文献では、(1) と (2) を (3) から定理
\ref{etale-cohomology-theorem-exactness-stalks} によって導くこともある。

\medskip\noindent
(3) は引き戻しの茎に関する一般的事実である。Sites, 補題
\ref{sites-lemma-point-morphism-sites} を参照せよ。
ここでは (3) を次のように直接証明する。補題
\ref{etale-cohomology-lemma-stalk-exact} により、茎を取ることは層化と可換であることに
注意する。ここで $f^{-1}\mathcal{G}$ は前層
$$
U \longrightarrow \colim_{U \to X \times_Y V} \mathcal{G}(V),
$$
に付随する層であったことを思い出そう。式
(\ref{etale-cohomology-equation-pullback}) を参照せよ。したがって
\begin{align*}
(f^{-1}\mathcal{G})_{\overline{x}}
& = \colim_{(U, \overline{u})} f^{-1}\mathcal{G}(U) \\
& = \colim_{(U, \overline{u})}
\colim_{a : U \to X \times_Y V} \mathcal{G}(V) \\
& = \colim_{(V, \overline{v})} \mathcal{G}(V) \\
& = \mathcal{G}_{\overline{y}}
\end{align*}
を得る。第三の等号では、対 $(U, \overline{u})$ と射
$a : U \to X \times_Y V$ が対 $(V, a \circ \overline{u})$ に対応する。

\medskip\noindent
(4) も、$f^{-1}h_V$ を定義する余極限を同一視することにより、同様に
証明できる。あるいは、米田の補題（Categories, 補題
\ref{categories-lemma-yoneda}）と関手的な等式
$$
\Mor_{\Sh(X_\etale)}(f^{-1}h_V, \mathcal{F}) =
\Mor_{\Sh(Y_\etale)}(h_V, f_*\mathcal{F}) =
f_*\mathcal{F}(V) = \mathcal{F}(X \times_Y V)
$$
を、表現可能前層が層であるという事実と併せて用いてもよい。
完全に一般的な結果については Sites, 補題
\ref{sites-lemma-pullback-representable-sheaf} も参照せよ。
\end{proof}

\noindent
関手の対 $(f_*, f^{-1})$ は小エタール・トポスの射
$$
f_{small} :
\Sh(X_\etale)
\longrightarrow
\Sh(Y_\etale)
$$
を定める。層のコホモロジーに関する一般論の多くは、トポスおよび
トポスの射に対して成り立つ。以下では、結果が一般的なものか、
エタール・トポスに固有のものかを、できるかぎり明記する。

\begin{remark}
\label{etale-cohomology-remark-functoriality-general}
より一般に、$\mathcal{C}_1, \mathcal{C}_2$ を終対象とファイバー積を
もつサイトとする。$u: \mathcal{C}_2 \to \mathcal{C}_1$ を次を満たす
関手とする。
\begin{enumerate}
\item $\{V_i \to V\}$ が $\mathcal{C}_2$ の被覆ならば、
$\{u(V_i) \to u(V)\}$ は $\mathcal{C}_1$ の被覆である
（このとき $u$ は {\it 連続}であるという）。
\item $u$ は有限極限と可換である（すなわち $u$ は左完全であり、
言い換えれば $u$ はファイバー積と終対象を保つ）。
\end{enumerate}
このとき
$f_*: \Sh(\mathcal{C}_1) \to \Sh(\mathcal{C}_2)$
を $ f_* \mathcal{F}(V) = \mathcal{F}(u(V))$ により定義できる。
さらに、完全関手 $f^{-1}$ が存在し、これは $f_*$ の左随伴である。
Sites, 定義 \ref{sites-definition-morphism-sites} および命題
\ref{sites-proposition-get-morphism} を参照せよ。注意：トポスの射を得るには、
$u$ が連続でファイバー積と可換であることだけを仮定しても十分ではない。
\end{remark}




\section{大トポスの関手性}
\label{etale-cohomology-section-functoriality-big-topoi}

\noindent
$f : X \to Y$ をスキームの射とすると、$f$ に付随するトポスの射は
多数ある。その一覧については Topologies, 節
\ref{topologies-section-change-topologies} を参照せよ。おそらく最もよく
用いられるものは、トポスの射
$$
f_{big} = f_{big, \tau} :
\Sh((\Sch/X)_\tau)
\longrightarrow
\Sh((\Sch/Y)_\tau)
$$
である。ここで $\tau \in \{Zariski, \etale, smooth, syntomic, fppf\}$
である。それぞれは連続関手
$$
(\Sch/Y)_\tau \longrightarrow (\Sch/X)_\tau, \quad
V/Y \longmapsto X \times_Y V/X
$$
に対応する。この関手は終対象、ファイバー積、および被覆を保つので、
サイトの射
$$
f_{big} : (\Sch/X)_\tau \longrightarrow (\Sch/Y)_\tau.
$$
を定める。Topologies, 節 \ref{topologies-section-zariski},
\ref{topologies-section-etale},
\ref{topologies-section-smooth},
\ref{topologies-section-syntomic}, および
\ref{topologies-section-fppf} を参照せよ。特に、$f_{big}$ に沿う順像は
規則
$$
(f_{big, *}\mathcal{F})(V/Y) = \mathcal{F}(X \times_Y V/X)
$$
で与えられる。これらのトポスの射は、きわめて簡単に記述できる逆像関手
$f_{big}^{-1}$ をもつ。すなわち
$$
(f_{big}^{-1}\mathcal{G})(U/X) = \mathcal{G}(U/Y)
$$
である。ここで $U/Y$ の構造射は、構造射 $U \to X$ と $f$ の合成である。
Topologies, 補題 \ref{topologies-lemma-morphism-big},
\ref{topologies-lemma-morphism-big-etale},
\ref{topologies-lemma-morphism-big-smooth},
\ref{topologies-lemma-morphism-big-syntomic}, および
\ref{topologies-lemma-morphism-big-fppf} を参照せよ。







\section{関手性と加群の層}
\label{etale-cohomology-section-morphisms-modules}

\noindent
本節では、Descent, 節 \ref{descent-section-quasi-coherent-sheaves},
\ref{descent-section-quasi-coherent-cohomology}, および
\ref{descent-section-quasi-coherent-sheaves-bis}
で説明した内容の一部を、エタール位相の設定で定式化し直す。
$f : X \to Y$ をスキームの射とする。上の節
\ref{etale-cohomology-section-functoriality}, \ref{etale-cohomology-section-direct-image}, および
\ref{etale-cohomology-section-inverse-image}
で、これが小エタールサイトの射 $f_{small}$ を誘導することを見た。
Descent, 注意 \ref{descent-remark-change-topologies-ringed} では、$f$ がさらに
自然な写像
$$
f_{small}^\sharp :
\mathcal{O}_{Y_\etale}
\longrightarrow
f_{small, *}\mathcal{O}_{X_\etale}
$$
を誘導することも見た。これは $Y_\etale$ 上の環の層の写像であり、
$(f_{small}, f_{small}^\sharp)$ は環付きサイトの射となる。
環付きサイトの射の定義については Modules on Sites, 定義
\ref{sites-modules-definition-ringed-site} を参照せよ。ここでは、
$f_{small}^\sharp$ が両立する写像系
$$
\text{pr}_V^\sharp : \mathcal{O}(V) \longrightarrow \mathcal{O}(X \times_Y V)
$$
によって定義されることだけを思い出しておこう。ここで $V$ は
$Y_\etale$ の対象全体を動く。

\medskip\noindent
この構成がスキームの射の合成と両立することは明らかである。より正確には、
$f : X \to Y$ および $g : Y \to Z$ がスキームの射ならば
$$
(g_{small}, g_{small}^\sharp)
\circ
(f_{small}, f_{small}^\sharp)
=
((g \circ f)_{small}, (g \circ f)_{small}^\sharp)
$$
が環付きトポスの射として成り立つ。さらに Modules on Sites, 定義
\ref{sites-modules-definition-pushforward} により、スキームの射
$f : X \to Y$ が与えられると、よく定義された引き戻し関手と順像関手
\begin{align*}
f_{small}^* :
\textit{Mod}(\mathcal{O}_{Y_\etale})
\longrightarrow
\textit{Mod}(\mathcal{O}_{X_\etale}), \\
f_{small, *} :
\textit{Mod}(\mathcal{O}_{X_\etale})
\longrightarrow
\textit{Mod}(\mathcal{O}_{Y_\etale})
\end{align*}
を得る。これらは通常の仕方で随伴する。$g : Y \to Z$ が別のスキームの
射ならば、合成について上で述べたことから
$(g \circ f)_{small}^* = f_{small}^* \circ g_{small}^*$
および $(g \circ f)_{small, *} = g_{small, *} \circ f_{small, *}$
が成り立つ。

\medskip\noindent
すべての $\mathcal{O}_X$ 加群からなる $X$ 上の圏と、すべての
$\mathcal{O}_{X_\etale}$ 加群からなる $X_\etale$ 上の圏との間には、
かなり大きな違いがある。
しかし Descent, 節 \ref{descent-section-quasi-coherent-sheaves},
\ref{descent-section-quasi-coherent-cohomology}, および
\ref{descent-section-quasi-coherent-sheaves-bis}
の結果によれば、$S$ 上の準連接加群と $S_\etale$ 上の準連接加群を
考えることには大きな違いがない。（たとえば、これはすでに定理
\ref{etale-cohomology-theorem-quasi-coherent} で見た。）特に、$f : X \to Y$ が任意の
スキームの射ならば、準連接層に対して引き戻し関手 $f_{small}^*$ と
$f^*$ は一致する。Descent, 命題
\ref{descent-proposition-equivalence-quasi-coherent-functorial} を参照せよ。
さらに、$f$ が準コンパクトかつ準分離的ならば、順像についても同じことが
成り立つ。Descent, 補題
\ref{descent-lemma-higher-direct-images-small-etale} を参照せよ。

\medskip\noindent
最後に、大サイト上の構造層の関手性について少し述べる。
$f : X \to Y$ をスキームの射とする。位相
$\tau \in \{Zariski,\linebreak[0]
\etale,\linebreak[0] smooth,\linebreak[0] syntomic,\linebreak[0]
fppf\}$ のいずれかを選ぶ。このとき射
$f_{big} : (\Sch/X)_\tau \to (\Sch/Y)_\tau$
は写像
$$
f_{big}^\sharp : \mathcal{O}_Y \longrightarrow f_{big, *}\mathcal{O}_X
$$
によって環付きサイトの射となる。Descent, 注意
\ref{descent-remark-change-topologies-ringed} を参照せよ。実際、これは
上で説明した小サイトの場合と同じ構成によって与えられる。












\section{位相の比較}
\label{etale-cohomology-section-compare-topologies}

\noindent
本節では、異なる位相に関する層を比較すると何が起こるかを調べ始める。

\begin{lemma}
\label{etale-cohomology-lemma-where-sections-are-equal}
$S$ をスキームとし、$\mathcal{F}$ を $S_\etale$ 上の集合の層とする。
$s, t \in \mathcal{F}(S)$ とする。このとき、次の性質によって特徴づけられる
開集合 $W \subset S$ が存在する：射 $f : T \to S$ が $W$ を経由することと、
$s|_T = t|_T$ であることは同値である（制限とは $f_{small}$ による
引き戻しである）。
\end{lemma}

\begin{proof}
$U \in \Ob(S_\etale)$ に対し、$s|_U \not = t|_U$ なら空集合を、そうで
なければ一点集合を対応させる前層を考える。これが $\Sh(S_\etale)$ の
終対象の部分層であることは明らかである。補題
\ref{etale-cohomology-lemma-support-subsheaf-final} により、この前層を表現する開集合
$W \subset S$ が得られる。幾何学的点 $\overline{x}$（$S$ の点）に対し、
$\overline{x} \in W$ であることと、$s$ と $t$ の $\overline{x}$ における
茎が一致することは同値である。補題 \ref{etale-cohomology-lemma-stalk-pullback} による
引き戻しの茎の記述から、$W$ が所望の性質をもつことが分かる。
\end{proof}

\begin{lemma}
\label{etale-cohomology-lemma-describe-pullback}
$S$ をスキームとし、$\tau \in \{Zariski, \etale\}$ とする。射
$$
\pi_S : (\Sch/S)_\tau \longrightarrow S_\tau
$$
を考える。これは Topologies, 補題 \ref{topologies-lemma-at-the-bottom}
または \ref{topologies-lemma-at-the-bottom-etale} の射である。
$\mathcal{F}$ を $S_\tau$ 上の層とする。このとき
$\pi_S^{-1}\mathcal{F}$ は規則
$$
(\pi_S^{-1}\mathcal{F})(T) = \Gamma(T_\tau, f_{small}^{-1}\mathcal{F})
$$
で与えられる。ここで $f : T \to S$ である。さらに、
$\pi_S^{-1}\mathcal{F}$ は fpqc 被覆に関する層条件を満たす。
\end{lemma}

\begin{proof}
射 $i_f : \Sh(T_\tau) \to \Sh(\Sch/S)_\tau)$ があり、
$\pi_S \circ i_f = f_{small}$ となることに注意する。これは射
$T_\tau \to S_\tau$ としての等式である。Topologies, 補題
\ref{topologies-lemma-put-in-T},
\ref{topologies-lemma-morphism-big-small},
\ref{topologies-lemma-put-in-T-etale}, および
\ref{topologies-lemma-morphism-big-small-etale} を参照せよ。
引き戻しの推移性から、所望のとおり
$i_f^{-1} \pi_S^{-1}\mathcal{F} = f_{small}^{-1}\mathcal{F}$ を得る。

\medskip\noindent
$\{g_i : T_i \to T\}_{i \in I}$ を fpqc 被覆とする。最後の主張の意味は
次のとおりである：層 $\mathcal{G}$ が $T_\tau$ 上にあり、切断
$s_i \in \Gamma(T_i, g_{i, small}^{-1}\mathcal{G})$ の
$T_i \times_T T_j$ への引き戻しが一致するとする。このとき一意な切断
$s$ が $\mathcal{G}$ の $T$ 上の切断として存在し、$T_i$ へのその
引き戻しは $s_i$ と一致する。

\medskip\noindent
$V \to T$ を $T_\tau$ の対象とし、$t \in \mathcal{G}(V)$ とする。
各 $i$ に対し、最大の開集合 $W_i \subset T_i \times_T V$ が存在し、
$s_i$ と $t$ の引き戻しは、$\mathcal{G}$ の
$W_i \subset T_i \times_T V$ への引き戻しの切断として一致する。補題
\ref{etale-cohomology-lemma-where-sections-are-equal} を参照せよ。$s_i$ と $s_j$ は
$T_i \times_T T_j$ 上で一致するので、$W_i$ と $W_j$ は
$T_i \times_T T_j \times_T V$ 上の同じ開集合へ引き戻される。
Descent, 補題 \ref{descent-lemma-open-fpqc-covering} により、開集合
$W \subset V$ が得られ、その $T_i \times_T V$ への逆像が $W_i$ となる。

\medskip\noindent
$g_{i, small}^{-1}\mathcal{G}$ の構成により、$\tau$-被覆
$\{T_{ij} \to T_i\}_{j \in J_i}$、各 $j$ に対する開埋め込みまたは
エタール射 $V_{ij} \to T$、切断 $t_{ij} \in \mathcal{G}(V_{ij})$、および
可換図式
$$
\xymatrix{
T_{ij} \ar[r] \ar[d] & V_{ij} \ar[d] \\
T_i \ar[r] & T
}
$$
が存在して、$s_i|_{T_{ij}}$ は $t_{ij}$ の引き戻しとなる。言い換えれば、
被覆 $\{T_i \to T\}$ を $\{T_{ij} \to T\}$ で置き換えることにより、
分解 $T_i \to V_i \to T$（ここで $V_i \in \Ob(T_\tau)$）と切断
$t_i \in \mathcal{G}(V_i)$ が存在し、それが $s_i$ へ、$T_i$ 上で
引き戻されると仮定して
よい。前段落の結果により開集合 $W_i \subset V_i$ が得られ、
$t_i|_{W_i}$ はすべての $s_j$ と $T_j \times_T W_i$ 上で「一致する」。
$T_i \to V_i$ は $W_i$ を経由することに注意する。ゆえに
$\{W_i \to T\}$ は $\tau$-被覆であり、補題が証明された。
\end{proof}

\begin{lemma}
\label{etale-cohomology-lemma-sections-upstairs}
$S$ をスキームとし、$f : T \to S$ を次を満たす射とする。
\begin{enumerate}
\item $f$ は平坦かつ準コンパクトである。
\item $f$ の幾何学的ファイバーは連結である。
\end{enumerate}
$\mathcal{F}$ を $S_\etale$ 上の層とする。このとき
$\Gamma(S, \mathcal{F}) = \Gamma(T, f^{-1}_{small}\mathcal{F})$ である。
\end{lemma}

\begin{proof}
標準写像
$\Gamma(S, \mathcal{F}) \to \Gamma(T, f_{small}^{-1}\mathcal{F})$.
がある。$f$ は全射である（そのファイバーが連結だからである）ので、
この写像は単射である。

\medskip\noindent
写像が全射であることを示すため、
$\alpha \in \Gamma(T, f_{small}^{-1}\mathcal{F})$ とする。
$\{T \to S\}$ は fpqc 被覆なので、補題
\ref{etale-cohomology-lemma-describe-pullback} により、二つの射影による $\alpha$ の
$T \times_S T$ への引き戻しが同じ切断であることを示せば十分である。
$\overline{s} \to S$ を幾何学的点とする。一致が
$(T \times_S T)_{\overline{s}}$ 上で成り立つことを示せば十分である。
実際、$T \times_S T$ の各幾何学的点はこれらの幾何学的ファイバーの
一つに含まれる。言い換えれば、$\alpha|_{T_{\overline{s}}}$ が
$$
(T \times_S T)_{\overline{s}} =
T_{\overline{s}} \times_{\overline{s}} T_{\overline{s}}
$$
上の同じ切断へ、二つの $T_{\overline{s}}$ への射影によって引き戻される
ことを示そうとしている。しかし
$\mathcal{F}|_{T_{\overline{s}}}$ は $\mathcal{F}|_{\overline{s}}$ の
引き戻しなので、値 $\mathcal{F}_{\overline{s}}$ をもつ定値層である。
仮定により $T_{\overline{s}}$ は連結だから、定値層の任意の切断は定値である。
したがって $\alpha|_{T_{\overline{s}}}$ は
$\mathcal{F}_{\overline{s}}$ の元に対応する。ゆえに
$(T \times_S T)_{\overline{s}}$ への二つの引き戻しはいずれもこの同じ元に
対応し、結論を得る。
\end{proof}

\noindent
次は、射が平坦であることを仮定しない補題
\ref{etale-cohomology-lemma-sections-upstairs} の一版である。

\begin{lemma}
\label{etale-cohomology-lemma-sections-upstairs-submersive}
$S$ をスキームとし、$f : X \to S$ を次を満たす射とする。
\begin{enumerate}
\item $f$ は商写像である。
\item $f$ の幾何学的ファイバーは連結である。
\end{enumerate}
$\mathcal{F}$ を $S_\etale$ 上の層とする。このとき
$\Gamma(S, \mathcal{F}) = \Gamma(X, f^{-1}_{small}\mathcal{F})$ である。
\end{lemma}

\begin{proof}
標準写像
$\Gamma(S, \mathcal{F}) \to \Gamma(X, f_{small}^{-1}\mathcal{F})$.
がある。$f$ は全射である（そのファイバーが連結だからである）ので、
この写像は単射である。

\medskip\noindent
写像が全射であることを示すため、
$\tau \in \Gamma(X, f_{small}^{-1}\mathcal{F})$ とする。
$\{U_i \to S\}$ をエタール被覆、$\sigma_i \in \mathcal{F}(U_i)$ を
切断として、$\sigma_i$ が $\tau|_{X \times_S U_i}$ へ引き戻されるものを
十分である。実際、上で示した単射性により、$\sigma_i$ と $\sigma_j$ は
$\mathcal{F}$ の $U_i \times_S U_j$ 上の同じ切断に制限される。
したがって、一意な切断 $\sigma \in \mathcal{F}(S)$ を得て、その制限は
$\sigma_i$（$U_i$ 上）である。このとき $\sigma$ の $X$ への
引き戻しは $\tau$ である。これは局所的に成り立つからである。

\medskip\noindent
$\overline{x}$ を $X$ の幾何学的点とし、その像 $\overline{s}$ を
$S$ における点とする。茎における $\tau$ の像
$$
(f_{small}^{-1}\mathcal{F})_{\overline{x}} = \mathcal{F}_{\overline{s}}
$$
を考える。補題 \ref{etale-cohomology-lemma-stalk-pullback} を参照せよ。
$U \to S$ を $\overline{s}$ のエタール近傍とし、茎でこの像へ写る切断
$\sigma \in \mathcal{F}(U)$ を見つけることができる。したがって $S$ を
$U$ で、$X$ を $X \times_S U$ で置き換えることにより、切断
$\sigma$、すなわち $\mathcal{F}$ の $S$ 上の切断が存在し、その
$(f_{small}^{-1}\mathcal{F})_{\overline{x}}$ における像は $\tau$ と同じで
あると仮定してよい。

\medskip\noindent
補題 \ref{etale-cohomology-lemma-where-sections-are-equal} により、最大の開集合
$W \subset X$ が存在して、$f_{small}^{-1}\sigma$ と $\tau$ は $W$ 上で
一致し、$W$ の形成はさらなる引き戻しと可換である。
$\mathcal{F}$ の幾何学的ファイバー $X_{\overline{s}}$ への引き戻しは、
$\mathcal{F}_{\overline{s}}$ を $\overline{s}$ 上の層と見なしたものの、
$X_{\overline{s}} \to \overline{s}$ による引き戻しであることに注意する。
したがって $\tau$ と $\sigma$ は、値 $\mathcal{F}_{\overline{s}}$ をもつ
定値層の $X_{\overline{s}}$ 上の切断を与え、それらは一点で
一致する。仮定により $X_{\overline{s}}$ は連結なので、$W$ は $X_s$ を
含むと結論できる。異なる幾何学的ファイバーについて同じ議論を行うと、
$W$ はそれが交わるすべてのファイバーを含むことが分かる。$f$ は商写像
なので、$W$ は $s$ のある開近傍（$S$ における近傍）の逆像である。これで証明が
完了する。
\end{proof}

\begin{lemma}
\label{etale-cohomology-lemma-sections-base-field-extension}
$K/k$ を体の拡大とし、$k$ は可分代数閉であるとする。$S$ を $k$ 上の
スキームとする。射影を
$p : S_K = S \times_{\Spec(k)} \Spec(K) \to S$ と書く。
$\mathcal{F}$ を $S_\etale$ 上の層とする。このとき
$\Gamma(S, \mathcal{F}) = \Gamma(S_K, p^{-1}_{small}\mathcal{F})$ である。
\end{lemma}

\begin{proof}
補題 \ref{etale-cohomology-lemma-sections-upstairs} から従う。実際、$p$ は
$\Spec(K) \to \Spec(k)$ の基底変換なので、平坦かつ準コンパクトである
ことは明らかである。一方、$\overline{s} : \Spec(L) \to S$ が幾何学的点
ならば、$p$ の $\overline{s}$ 上のファイバーは $K \otimes_k L$ の
スペクトルであり、Algebra, 補題
\ref{algebra-lemma-separably-closed-irreducible} により既約、したがって
連結である。
\end{proof}









\section{射の復元}
\label{etale-cohomology-section-morphisms}

\noindent
本節では、スキームにその局所環付き小エタール・トポスを対応させる規則が、
適切な意味で忠実充満であることを証明する。定理
\ref{etale-cohomology-theorem-fully-faithful} を参照せよ。

\begin{lemma}
\label{etale-cohomology-lemma-morphism-locally-ringed}
$f : X \to Y$ をスキームの射とする。環付きサイトの射
$(f_{small}, f_{small}^\sharp)$ は $f$ に付随するものであり、局所環付きサイトの射である。
Modules on Sites,
定義 \ref{sites-modules-definition-morphism-locally-ringed-topoi}
を参照せよ。
\end{lemma}

\begin{proof}
補題 \ref{etale-cohomology-lemma-etale-site-locally-ringed} で見たように、
$(X_\etale, \mathcal{O}_{X_\etale})$ と
$(Y_\etale, \mathcal{O}_{Y_\etale})$ は局所環付きサイトなので、
この主張には意味がある。さらに、定理
\ref{etale-cohomology-theorem-exactness-stalks} および注意
\ref{etale-cohomology-remarks-enough-points} により、$X_\etale$ は十分な点をもつ。
したがって、$(f_{small}, f_{small}^\sharp)$ が
Modules on Sites,
補題 \ref{sites-modules-lemma-locally-ringed-morphism}
の条件 (3) を満たすことを示せば十分である。このために点 $p$ を取り、これを
$X_\etale$ の点とする。補題 \ref{etale-cohomology-lemma-points-small-etale-site} により、$p$ は
幾何学的点 $\overline{x}$ に対応し、これは $X$ の点である。補題
\ref{etale-cohomology-lemma-stalk-pullback} により、点 $q = f_{small} \circ p$ は
幾何学的点 $\overline{y} = f \circ \overline{x}$ に対応し、これは $Y$ の点である。
よって示すべきことは、誘導される茎の写像
$$
(\mathcal{O}_Y)_{\overline{y}} \longrightarrow (\mathcal{O}_X)_{\overline{x}}
$$
が局所環準同型であることである。$a \in (\mathcal{O}_Y)_{\overline{y}}$ を、
右辺の極大イデアルの元に写る左辺の元とする。$a$ は三つ組
$(V, \overline{v}, a)$ の同値類であるとする。ここで $V \to Y$ はエタール、
$\overline{v} : \overline{x} \to V$ は $Y$ 上の射、$a \in \mathcal{O}(V)$ である。
これは $(X \times_Y V, \overline{x} \times \overline{v}, \text{pr}_V^\sharp(a))$
の同値類に写り、後者は局所環 $(\mathcal{O}_X)_{\overline{x}}$ の元である。
ところで、これが極大イデアルに属するということは、
$\text{pr}_V^\sharp(a)$ を $\kappa(\overline{x})$ の元へ引き戻すと零になることを
意味する。したがって $a$ を $\kappa(\overline{x})$ へ引き戻しても零となる。
これは $a$ が $(\mathcal{O}_Y)_{\overline{y}}$ の極大イデアルに属することを意味する。
\end{proof}

\begin{lemma}
\label{etale-cohomology-lemma-2-morphism}
$X$, $Y$ をスキームとし、$f : X \to Y$ をスキームの射とする。
$t$ を $2$-射であって $(f_{small}, f_{small}^\sharp)$ からそれ自身へのものとする
（Modules on Sites,
定義 \ref{sites-modules-definition-2-morphism-ringed-topoi} 参照）。
このとき $t = \text{id}$ である。
\end{lemma}

\begin{proof}
これは、$t : f^{-1}_{small} \to f^{-1}_{small}$ が関手の自然変換であり、図式
$$
\xymatrix{
f_{small}^{-1}\mathcal{O}_Y
\ar[rd]_{f_{small}^\sharp}  & &
f_{small}^{-1}\mathcal{O}_Y \ar[ll]^t \ar[ld]^{f_{small}^\sharp} \\
& \mathcal{O}_X
}
$$
が可換であることを意味する。$V \to Y$ がエタールで、$V$ がアフィンであるとする。
Morphisms, 補題 \ref{morphisms-lemma-quasi-affine-finite-type-over-S}
により、はめ込み $i : V \to \mathbf{A}^n_Y$ を選べる。これは $Y$ 上の射である。
層の言葉では、これは $i$ が層の単射
$h_i : h_V \to \prod_{j = 1, \ldots, n} \mathcal{O}_Y$ を誘導することを意味する。
$i'$ を $i$ の $X$ への基底変換とする。これは Schemes,
補題 \ref{schemes-lemma-base-change-immersion} により、はめ込みである。
したがって $i' : X \times_Y V \to \mathbf{A}^n_X$ ははめ込みであり、これはさらに
$h_{i'} : h_{X \times_Y V} \to \prod_{j = 1, \ldots, n} \mathcal{O}_X$
が層の単射であることを意味する。補題 \ref{etale-cohomology-lemma-stalk-pullback} の同一視
$f_{small}^{-1}h_V = h_{X \times_Y V}$ のもとで、写像 $h_{i'}$ は
$$
\xymatrix{
f_{small}^{-1}h_V \ar[r]^-{f^{-1}h_i} &
\prod_{j = 1, \ldots, n} f_{small}^{-1}\mathcal{O}_Y
\ar[r]^{\prod f^\sharp} &
\prod_{j = 1, \ldots, n} \mathcal{O}_X
}
$$
に等しい（検証は省略する）。したがって、写像
$t : f_{small}^{-1}h_V \to f_{small}^{-1}h_V$ は可換図式
$$
\xymatrix{
f_{small}^{-1}h_V \ar[r]^-{f^{-1}h_i} \ar[d]^t &
\prod_{j = 1, \ldots, n} f_{small}^{-1}\mathcal{O}_Y
\ar[r]^-{\prod f^\sharp} \ar[d]^{\prod t} &
\prod_{j = 1, \ldots, n} \mathcal{O}_X \ar[d]^{\text{id}}\\
f_{small}^{-1}h_V \ar[r]^-{f^{-1}h_i} &
\prod_{j = 1, \ldots, n} f_{small}^{-1}\mathcal{O}_Y
\ar[r]^-{\prod f^\sharp} &
\prod_{j = 1, \ldots, n} \mathcal{O}_X
}
$$
に組み込まれる。右側の正方形の可換性は、上で説明した $t$ に関する仮定から従う。
上の議論により横向きの射の合成は単射なので、左の縦向きの射も恒等写像である。
$Y_\etale$ 上の任意の集合の層には、$h_V$ という形の層（ここで $V$ はアフィン）の
（巨大な）余積からの全射が存在する（Topologies,
補題 \ref{topologies-lemma-alternative} と Sites,
補題 \ref{sites-lemma-sheaf-coequalizer-representable} を組み合わせよ）。
よって $t : f_{small}^{-1} \to f_{small}^{-1}$ は、望むとおり恒等自然変換である。
\end{proof}

\begin{lemma}
\label{etale-cohomology-lemma-faithful}
$X$, $Y$ をスキームとする。任意の二つのスキームの射 $a, b : X \to Y$ について、$2$-同型
$(a_{small}, a_{small}^\sharp) \cong (b_{small}, b_{small}^\sharp)$
が環付きトポスの $2$-圏で存在するなら、これらの射は等しい。
\end{lemma}

\begin{proof}
多少紛らわしいので、注意深く論じよう。
$t : a_{small}^{-1} \to b_{small}^{-1}$ をこの $2$-同型とする。
任意の開集合 $V \subset Y$ を考える。$h_V$ は終対象の層 $*$ の部分層である。
したがって $a_{small}^{-1}h_V = h_{a^{-1}(V)}$ と
$b_{small}^{-1}h_V = h_{b^{-1}(V)}$ はともに終対象の層の部分層である。
よって同型
$$
t : a_{small}^{-1}h_V = h_{a^{-1}(V)} \to b_{small}^{-1}h_V = h_{b^{-1}(V)}
$$
は恒等射でなければならず、$a^{-1}(V) = b^{-1}(V)$ である。
したがって $a$ と $b$ は台位相空間上で等しい。次に切断
$f \in \mathcal{O}_Y(V)$ を取る。これは集合の層の射
$f : h_V \to \mathcal{O}_Y$ を定め、またそのような射によって定まる。
これを引き戻し、$t$ を適用すると、可換図式
$$
\xymatrix{
h_{b^{-1}(V)} \ar@{=}[r] &
b_{small}^{-1}h_V \ar[d]^{b_{small}^{-1}(f)} & &
a_{small}^{-1}h_V \ar[d]^{a_{small}^{-1}(f)} \ar[ll]^t &
h_{a^{-1}(V)} \ar@{=}[l]
\\
& b_{small}^{-1}\mathcal{O}_Y
\ar[rd]_{b^\sharp}  & &
a_{small}^{-1}\mathcal{O}_Y \ar[ll]^t \ar[ld]^{a^\sharp} \\
& & \mathcal{O}_X
}
$$
を得る。ここで三角形が可換であることは、Modules on Sites,
節 \ref{sites-modules-section-2-category} における $2$-同型の定義による。
上で見たように、上段の横向きの射の合成は恒等式
$a^{-1}(V) = b^{-1}(V)$ から得られる。したがって図式の可換性より、
$a_{small}^\sharp(f) = b_{small}^\sharp(f)$ が
$\mathcal{O}_X(a^{-1}(V)) = \mathcal{O}_X(b^{-1}(V))$ において成り立つ。
これは任意の開集合 $V$ と
任意の $f \in \mathcal{O}_Y(V)$ に対して成り立つので、スキームの射として
$a = b$ である。
\end{proof}

\begin{lemma}
\label{etale-cohomology-lemma-morphism-ringed-etale-topoi-affines}
$X$, $Y$ をアフィン・スキームとする。局所環付きトポスの射
$$
(g, g^\#) :
(\Sh(X_\etale), \mathcal{O}_X)
\longrightarrow
(\Sh(Y_\etale), \mathcal{O}_Y)
$$
を与える。このとき、一意なスキームの射 $f : X \to Y$ が存在して、
$(g, g^\#)$ は $2$-同型を介して $(f_{small}, f_{small}^\sharp)$ と同一視される。
Modules on Sites,
定義 \ref{sites-modules-definition-2-morphism-ringed-topoi}
を参照せよ。
\end{lemma}

\begin{proof}
この証明では、$\mathcal{O}_X$ と書くのは小エタール・サイト
$X_\etale$ の構造層であり、$\mathcal{O}_Y$ についても同様とする。
$Y = \Spec(B)$ および $X = \Spec(A)$ と書こう。
$B = \Gamma(Y_\etale, \mathcal{O}_Y)$,
$A = \Gamma(X_\etale, \mathcal{O}_X)$
なので、$g^\sharp$ は環準同型 $\varphi : B \to A$ を誘導する。
$f = \Spec(\varphi) : X \to Y$ を対応するアフィン・スキームの射とする。
この $f$ が所要の性質を満たすことを示す。

\medskip\noindent
$V \to Y$ を $Y$ 上エタールなアフィン・スキームとする。したがって
$V = \Spec(C)$ と書ける。ここで $C$ はエタール $B$-代数である。さらに
$$
C = B[x_1, \ldots, x_n]/(P_1, \ldots, P_n)
$$
と書ける。ここで $P_i$ は、$\Delta = \det(\partial P_i/ \partial x_j)$ が
$C$ で可逆となるような多項式である。例えば Algebra,
補題 \ref{algebra-lemma-etale-standard-smooth} を参照せよ。
$T$ が $Y$ 上のスキームなら、$T$-値点としての $V$ の点は、$n$ 個の
$\Gamma(T, \mathcal{O}_T)$ の切断であって多項式方程式
$P_1 = 0, \ldots, P_n = 0$ を満たすものによって与えられる。言い換えれば、層
$h_V$ は $Y_\etale$ 上で二つの写像
$$
\xymatrix{
\prod\nolimits_{i = 1, \ldots, n} \mathcal{O}_Y
\ar@<1ex>[r]^a \ar@<-1ex>[r]_b &
\prod\nolimits_{j = 1, \ldots, n} \mathcal{O}_Y
}
$$
の等化子である。ここで $b(h_1, \ldots, h_n) = 0$ であり、
$a(h_1, \ldots, h_n) =
(P_1(h_1, \ldots, h_n), \ldots, P_n(h_1, \ldots, h_n))$ である。
$g^{-1}$ は完全なので、次の実線部分からなる可換図式の上段も等化子図式である：
$$
\xymatrix{
g^{-1}h_V \ar[r] \ar@{..>}[d] &
\prod\nolimits_{i = 1, \ldots, n} g^{-1}\mathcal{O}_Y
\ar@<1ex>[r]^{g^{-1}a} \ar@<-1ex>[r]_{g^{-1}b} \ar[d]^{\prod g^\sharp} &
\prod\nolimits_{j = 1, \ldots, n} g^{-1}\mathcal{O}_Y \ar[d]^{\prod g^\sharp}\\
h_{X \times_Y V} \ar[r] &
\prod\nolimits_{i = 1, \ldots, n} \mathcal{O}_X
\ar@<1ex>[r]^{a'} \ar@<-1ex>[r]_{b'} &
\prod\nolimits_{j = 1, \ldots, n} \mathcal{O}_X  \\
}
$$
ここで $b'$ は零写像であり、$a'$ は像
$P'_i = \varphi(P_i) \in A[x_1, \ldots, x_n]$ を用いて、
$a'(h_1, \ldots, h_n) =
(P'_1(h_1, \ldots, h_n), \ldots, P'_n(h_1, \ldots, h_n))$.
という、$a$ の定義と同じ規則で定められる。図式の可換性は、大域切断上で
$\varphi = g^\sharp$ であることから従う。さらに $X \times_Y V$ はアフィン・スキーム
$\Spec(A \otimes_B C)$ であり、
$A \otimes_B C = A[x_1, \ldots, x_n]/(P'_1, \ldots, P'_n)$.
であるから、先とまったく同じ議論により下段も等化子図式である。
したがって、この図式に組み込まれる一意な点線の射
$g^{-1}h_V \to h_{X \times_Y V}$ を得る。

\medskip\noindent
層の写像 $g^{-1}h_V \to h_{X \times_Y V}$ が同型であると主張する。
$X$ の小エタール・サイトは十分な点をもつので
（定理 \ref{etale-cohomology-theorem-exactness-stalks}）、茎について示せば十分である。
そこで $\overline{x}$ を $X$ の幾何学的点とし、$p$ を $X$ の小エタール・トポスの
対応する点とする。$q = g \circ p$ と置く。これは $Y$ の小エタール・トポスの点である。
補題 \ref{etale-cohomology-lemma-points-small-etale-site} により、$q$ は幾何学的点
$\overline{y}$ に対応し、これは $Y$ の点である。茎の写像
$$
(g^\sharp)_p :
(\mathcal{O}_Y)_{\overline{y}} =
\mathcal{O}_{Y, q} =
(g^{-1}\mathcal{O}_Y)_p
\longrightarrow
\mathcal{O}_{X, p} =
(\mathcal{O}_X)_{\overline{x}}
$$
を考える。$(g, g^\sharp)$ は {\it 局所}環付きトポスの射なので、
$(g^\sharp)_p$ は強 Hensel 局所環の間の局所環準同型である。
上の大きな可換図式に局所化を適用し、Algebra,
補題 \ref{algebra-lemma-strictly-henselian-solutions} を用いると、望むとおり
$(g^{-1}h_V)_p \to (h_{X \times_Y V})_p$ が同型であることが分かる。

\medskip\noindent
同型 $g^{-1}h_V \to h_{X \times_Y V}$ は関手的であると主張する。
実際、$V_1 \to V_2$ が $Y$ 上エタールなアフィン・スキームの射であるとする。
$V_i = \Spec(C_i)$ と書き、
$$
C_i = B[x_{i, 1}, \ldots, x_{i, n_i}]/(P_{i, 1}, \ldots, P_{i, n_i})
$$
とする。射 $V_1 \to V_2$ は $B$-代数準同型 $C_2 \to C_1$ によって与えられ、
後者はある多項式 $Q_j \in B[x_{1, 1}, \ldots, x_{1, n_1}]$
（$j = 1, \ldots, n_2$）によって与えられる。
このとき、層の図式
$$
\xymatrix{
h_{V_1} \ar[d] \ar[r] & \prod_{i = 1, \ldots, n_1} \mathcal{O}_Y
\ar[d]^{Q_1, \ldots, Q_{n_2}}\\
h_{V_2} \ar[r] & \prod_{i = 1, \ldots, n_2} \mathcal{O}_Y
}
$$
が可換であることは容易に分かる。これを $X_\etale$ へ引き戻すと、実線部分が可換な図式
$$
\xymatrix{
g^{-1}h_{V_1} \ar@{..>}[dd] \ar[rrd] \ar[r] &
\prod_{i = 1, \ldots, n_1} g^{-1}\mathcal{O}_Y
\ar[dd]^{g^\sharp}
\ar[rrd]^{Q_1, \ldots, Q_{n_2}} \\
& & g^{-1}h_{V_2} \ar@{..>}[dd] \ar[r] &
\prod_{i = 1, \ldots, n_2} g^{-1}\mathcal{O}_Y
\ar[dd]^{g^\sharp} \\
h_{X \times_Y V_1} \ar[r] \ar[rrd] &
\prod\nolimits_{i = 1, \ldots, n_1} \mathcal{O}_X
\ar[rrd]^{Q'_1, \ldots, Q'_{n_2}} \\
& & h_{X \times_Y V_2} \ar[r] &
\prod\nolimits_{i = 1, \ldots, n_2} \mathcal{O}_X
}
$$
を得る。ここで $Q'_j \in A[x_{1, 1}, \ldots, x_{1, n_1}]$ は
$Q_j$ の、$\varphi$ による像である。点線の射は存在して二つの正方形を可換にし、
横向きの射は単射なので、図式全体が可換である。これで関手性が示された。
（また、厳密には不要だが、$g^{-1}h_V \to h_{X \times_Y V}$ の構成が
表示の選択に依存しないことも示される。）

\medskip\noindent
ここで $f_{small, *} \cong g_*$ を示せる。実際、$\mathcal{F}$ を
$X_\etale$ 上の層とする。任意のアフィンな $V \in \Ob(X_\etale)$ に対して
\begin{align*}
(g_*\mathcal{F})(V)
& =
\Mor_{\Sh(Y_\etale)}(h_V, g_*\mathcal{F}) \\
& =
\Mor_{\Sh(X_\etale)}(g^{-1}h_V, \mathcal{F}) \\
& =
\Mor_{\Sh(X_\etale)}(h_{X \times_Y V}, \mathcal{F}) \\
& =
\mathcal{F}(X \times_Y V) \\
& =
f_{small, *}\mathcal{F}(V)
\end{align*}
である。第三の等号では、上で構成した同型
$g^{-1}h_V \cong h_{X \times_Y V}$ を用いた。この同型が関手的なので、
これらの同型は明らかに $\mathcal{F}$ に関しても $V$ に関しても関手的である。
$g^{-1}h_V \cong h_{X \times_Y V}$ の関手性を用いた。ところで $Y_\etale$ 上の
任意の層は、アフィン・スキームからなる部分圏への制限によって決まる
（Topologies, 補題 \ref{topologies-lemma-alternative}）。
したがって、望む関手の同型 $f_{small, *} \cong g_*$ を得る。

\medskip\noindent
最後に、上の同型 $f_{small, *} \cong g_*$ を介して写像
$f_{small}^\sharp$ と $g^\sharp$ が一致することを確かめなければならない。
構成により、これは $\mathcal{O}_Y$ の大域切断、すなわち $B$ の元については
すでに成り立つ。再び Topologies,
補題 \ref{topologies-lemma-alternative} により、アフィン $V$ 上の切断についてだけ
確かめればよい。このアフィン・スキームは $Y$ 上エタールである。先と同様に
$V = \Spec(C)$, $C = B[x_i]/(P_j)$ と書くと、座標関数 $x_i$ が
$\mathcal{O}_X$ の同じ切断へ写ることを確認すれば十分である。これらは $X \times_Y V$
上の切断である。
そしてこれはまさに図式
$$
\xymatrix{
g^{-1}h_V \ar[r] \ar@{..>}[d] &
\prod\nolimits_{i = 1, \ldots, n} g^{-1}\mathcal{O}_Y
\ar[d]^{\prod g^\sharp} \\
h_{X \times_Y V} \ar[r] &
\prod\nolimits_{i = 1, \ldots, n} \mathcal{O}_X
}
$$
が可換であることを意味する。よって補題が証明された。
\end{proof}

\noindent
一般のスキームに対する形は次のとおりである。

\begin{theorem}
\label{etale-cohomology-theorem-fully-faithful}
$X$, $Y$ をスキームとし、局所環付きトポスの射
$$
(g, g^\#) :
(\Sh(X_\etale), \mathcal{O}_X)
\longrightarrow
(\Sh(Y_\etale), \mathcal{O}_Y)
$$
を与える。このとき、一意なスキームの射 $f : X \to Y$ が存在して、
$(g, g^\#)$ は $(f_{small}, f_{small}^\sharp)$ と同型になる。
言い換えれば、構成
$$
\Sch \longrightarrow \textit{局所環付きトポス},
\quad
X \longrightarrow (X_\etale, \mathcal{O}_X)
$$
は忠実充満である（右辺では射を $2$-同型を除いて考える）。
\end{theorem}

\begin{proof}
補題 \ref{etale-cohomology-lemma-morphism-ringed-etale-topoi-affines} の証明の議論を
大域的な状況に注意深く合わせれば、この定理を証明できる。しかしここでは、補題
\ref{etale-cohomology-lemma-2-morphism} の結果が与える剛性を用いて、先の補題の結果を貼り合わせ、
大域的な射を得る方法を示したい。残念ながら、議論はいささか煩雑である。

\medskip\noindent
まず $Y$ がアフィンの場合に存在を証明する。この場合、アフィン開被覆
$X = \bigcup U_i$ を選ぶ。各 $i$ に対して、包含射 $j_i : U_i \to X$ は、補題
\ref{etale-cohomology-lemma-morphism-locally-ringed} により局所環付きトポスの射
$(j_{i, small}, j_{i, small}^\sharp) :
(\Sh(U_{i, \etale}), \mathcal{O}_{U_i})
\to
(\Sh(X_\etale), \mathcal{O}_X)$
を誘導する。これを $(g, g^\sharp)$ と合成して、局所環付きトポスの射
$$
(g, g^\sharp) \circ (j_{i, small}, j_{i, small}^\sharp) :
(\Sh(U_{i, \etale}), \mathcal{O}_{U_i})
\to
(\Sh(Y_\etale), \mathcal{O}_Y)
$$
を得る。Modules on Sites,
補題 \ref{sites-modules-lemma-composition-morphisms-locally-ringed-topoi}
を参照せよ。補題 \ref{etale-cohomology-lemma-morphism-ringed-etale-topoi-affines} により、
一意なスキームの射 $f_i : U_i \to Y$ と $2$-同型
$$
t_i :
(f_{i, small}, f_{i, small}^\sharp)
\longrightarrow
(g, g^\sharp) \circ (j_{i, small}, j_{i, small}^\sharp).
$$
$U_{i, i'} = U_i \cap U_{i'}$ と置き、包含射を
$j_{i, i'} : U_{i, i'} \to U_i$ と書く。等式
$j_i \circ j_{i, i'} = j_{i'} \circ j_{i', i}$
が成り立つので、
\begin{align*}
(g, g^\sharp) \circ
(j_{i, small}, j_{i, small}^\sharp) \circ
(j_{i, i', small}, j_{i, i', small}^\sharp)
= \\
(g, g^\sharp) \circ
(j_{i', small}, j_{i', small}^\sharp) \circ
(j_{i', i, small}, j_{i', i, small}^\sharp)
\end{align*}
である。一意性（補題 \ref{etale-cohomology-lemma-faithful} 参照）により、
$f_i \circ j_{i, i'} = f_{i'} \circ j_{i', i}$ を得る。言い換えれば、
スキームの射 $f_i = f \circ j_i$ は、大域的なスキームの射
$f : X \to Y$ の制限である。次の $2$-同型の図式を考える
（記法を簡単にするため、成分 ${}^\sharp$ を省略する）。
$$
\xymatrix{
g \circ j_{i, small} \circ j_{i, i', small}
\ar[rr]^{t_i \star \text{id}_{j_{i, i', small}}}
\ar@{=}[d] & &
f_{small} \circ j_{i, small} \circ j_{i, i', small} \ar@{=}[d] \\
g \circ j_{i', small} \circ j_{i', i, small}
\ar[rr]^{t_{i'} \star \text{id}_{j_{i', i, small}}} & &
f_{small} \circ j_{i', small} \circ j_{i', i, small}
}
$$
記号 $\star$ は水平合成を表す。一般には Categories,
定義 \ref{categories-definition-2-category} を、本件については Sites,
節 \ref{sites-section-2-category} を参照せよ。補題
\ref{etale-cohomology-lemma-2-morphism} により、この図式は可換である。したがって、任意の層
$\mathcal{G}$ を取り、それが $Y_\etale$ 上の層であるとすると、同型
$t_i : f_{small}^{-1}\mathcal{G}|_{U_i} \to g^{-1}\mathcal{G}|_{U_i}$ は
$U_{i, i'}$ 上で一致し、大域的な同型
$t : f_{small}^{-1}\mathcal{G} \to g^{-1}\mathcal{G}$ を得る。
この同型が $\mathcal{G}$ に関して関手的であり、写像 $f_{small}^\sharp$ および
$g^\sharp$ と両立することは明らかである（局所的にこれらの写像と両立するからである）。
これで $Y$ がアフィンの場合に定理が証明された。

\medskip\noindent
一般の場合、$V \subset Y$ をアフィン開集合とする。このとき $h_V$ は終対象の層
$*$ の $Y_\etale$ 上の部分層である。$g$ は完全なので、$g^{-1}h_V$ は
$X_\etale$ 上の終対象の層の部分層である。よって補題
\ref{etale-cohomology-lemma-support-subsheaf-final} により、開部分スキーム $W \subset X$ が存在して
$g^{-1}h_V = h_W$ となる。Modules on Sites,
補題 \ref{sites-modules-lemma-localize-morphism-locally-ringed-topoi} により、
局所環付きトポスの射からなる可換図式
$$
\xymatrix{
(\Sh(W_\etale), \mathcal{O}_W) \ar[r] \ar[d]_{g'} &
(\Sh(X_\etale), \mathcal{O}_X) \ar[d]^g \\
(\Sh(V_\etale), \mathcal{O}_V) \ar[r] &
(\Sh(Y_\etale), \mathcal{O}_Y)
}
$$
が存在する。ここで横向きの射は局所化射であり（包含射 $V \to Y$ および
$W \to X$ によって誘導される）、$g'$ は $g$ から誘導される。
前段落の結果により、スキームの射 $f' : W \to V$ と $2$-同型
$t : (f'_{small}, (f'_{small})^\sharp) \to (g', (g')^\sharp)$ を得る。
先とまったく同様に、これらの射 $f'$ は（アフィン開集合 $V \subset Y$ を動かすとき）
一意性により重なり上で一致するので、射 $f : X \to Y$ を得る。さらに、補題
\ref{etale-cohomology-lemma-2-morphism} により $2$-同型 $t$ も重なり上で両立し、大域的な $2$-同型
$(f_{small}, (f_{small})^\sharp) \to (g, (g)^\sharp)$ を得る。
これが望む結論である。いくつかの詳細は省略する。
\end{proof}










\section{押し出しと引き戻し}
\label{etale-cohomology-section-monomorphisms}

\noindent
スキームの射 $f : X \to Y$ をとる。
以下では、次の条件を考える。
\begin{enumerate}
\item[(A)] 任意のエタール射 $U \to X$ と $u \in U$ に対し、
エタール射 $V \to Y$、非交和分解
$X \times_Y V = W \amalg W'$、および射 $h : W \to U$ が存在し、これは $X$ 上の射であって、
$u$ は $h$ の像に含まれる。
\item[(B)] 任意のエタール射 $V \to Y$ と任意のエタール被覆
$\{U_i \to X \times_Y V\}$ に対し、エタール被覆
$\{V_j \to V\}$ が存在して、各 $j$ について
$X \times_Y V_j = \coprod W_{ij}$ となり、$W_{ij} \to X \times_Y V$ は
$U_i \to X \times_Y V$ を経由する（ある $i$ について）。
\item[(C)] 任意のエタール射 $U \to X$ に対し、エタール射 $V \to Y$
と全射 $X \times_Y V \to U$ が存在し、後者は $X$ 上の射である。
\end{enumerate}
これらの各性質は、関手 $f_{small, *}$ の振る舞いによって意味づけられることが分かる。
続く数節でこれを明らかにする。



\section{性質 (A)}
\label{etale-cohomology-section-A}

\noindent
性質 (A) の定義については、節 \ref{etale-cohomology-section-monomorphisms} を参照せよ。

\begin{lemma}
\label{etale-cohomology-lemma-property-A-implies}
スキームの射 $f : X \to Y$ をとり、(A) を仮定する。
\begin{enumerate}
\item
$f_{small, *} :
\textit{Ab}(X_\etale)
\to
\textit{Ab}(Y_\etale)$
は単射と全射を反映する。
\item $f_{small}^{-1}f_{small, *}\mathcal{F} \to \mathcal{F}$
は任意のアーベル層 $\mathcal{F}$ に対して $X_\etale$ 上で全射である。
\item
$f_{small, *} :
\textit{Ab}(X_\etale)
\to
\textit{Ab}(Y_\etale)$
は忠実である。
\end{enumerate}
\end{lemma}

\begin{proof}
$\mathcal{F}$ を $X_\etale$ 上のアーベル層とする。
$U$ を $X_\etale$ の対象とする。仮定により、被覆
$\{W_i \to U\}$ を $X_\etale$ において選べる。ただし各 $W_i$ は
$X \times_Y V_i$ の開かつ閉な部分スキームとなるものを選べる。ここで
$V_i$ は $Y_\etale$ のある対象である。層の条件から
$$
\mathcal{F}(U) \subset \prod \mathcal{F}(W_i)
$$
が成り立ち、さらに $\mathcal{F}(W_i)$ は
$\mathcal{F}(X \times_Y V_i) = f_{small, *}\mathcal{F}(V_i)$ の直和因子である。
したがって $f_{small, *}$ が単射を反映することは明らかである。

\medskip\noindent
次に、$a : \mathcal{G} \to \mathcal{F}$ をアーベル層の射で、
$f_{small, *}a$ が全射であるものとする。$s \in \mathcal{F}(U)$ をとる。
ここで $U$ は上のようなものとする。上の $W_i$、$V_i$ を用いると、
$s|_{W_i}$ がエタール局所的に $\mathcal{G}$ の切断の $a$ による像であることを
示せば十分である。$\mathcal{F}(W_i)$ は
$\mathcal{F}(X \times_Y V_i)$ の直和因子だから、任意の
$V \in \Ob(Y_\etale)$ と任意の元 $s \in \mathcal{F}(X \times_Y V)$ に対し、
$X \times_Y V$ 上エタール局所的に、これが $\mathcal{G}$ の切断の $a$ による
像であることを示せば十分である。実際、
$\mathcal{F}(X \times_Y V) = f_{small, *}\mathcal{F}(V)$
であり、仮定により被覆 $\{V_j \to V\}$ が存在して、$s$ は
$s_j \in f_{small, *}\mathcal{G}(V_j) = \mathcal{G}(X \times_Y V_j)$ の像となる。
これで $f_{small, *}$ が全射を反映することが示された。

\medskip\noindent
(2)、(3) は (1) から形式的に従う。Modules on Sites,
補題 \ref{sites-modules-lemma-reflect-surjections} を参照せよ。
\end{proof}

\begin{lemma}
\label{etale-cohomology-lemma-locally-quasi-finite-A}
$f : X \to Y$ をスキームの分離かつ局所準有限な射とする。
このとき上の性質 (A) が成り立つ。
\end{lemma}

\begin{proof}
$U \to X$ をエタール射とし、$u \in U$ とする。
幾何学的主張 (A) は、$U$ と $Y$ がアフィンスキームである場合に直ちに帰着する。
点 $x \in X$、$y \in Y$ を $u$ の像とする。$X \to Y$ は局所準有限であり、
$U \to X$ も局所準有限なので（Morphisms, 補題
\ref{morphisms-lemma-etale-locally-quasi-finite} を参照）、
$U \to Y$ は局所準有限である（Morphisms, 補題
\ref{morphisms-lemma-composition-quasi-finite} を参照）。
さらに $X \to Y$ と $U \to Y$ はいずれも分離である。したがって
More on Morphisms, 補題
\ref{more-morphisms-lemma-etale-splits-off-quasi-finite-part-technical-variant}
を両方の射に適用できる。よって、次を満たすエタール近傍
$(V, v) \to (Y, y)$ を選べる。
$$
X \times_Y V = W \amalg R, \quad
U \times_Y V = W' \amalg R'
$$
さらに点 $w \in W$、$w' \in W'$ が存在して、次が成り立つ。
\begin{enumerate}
\item $W$、$R$ は $X \times_Y V$ の開かつ閉な部分スキームである。
\item $W'$、$R'$ は $U \times_Y V$ の開かつ閉な部分スキームである。
\item $W \to V$ と $W' \to V$ は有限である。
\item $w$、$w'$ は $v$ に写る。
\item $\kappa(v) \subset \kappa(w)$ と $\kappa(v) \subset \kappa(w')$
は純非分離である。
\item $W$ または $W'$ の他の点は $v$ に写らない。
\end{enumerate}
次の可換図式を得る。
$$
\xymatrix{
U \ar[d] & U \times_Y V \ar[l] \ar[d] & W' \amalg R' \ar[d] \ar[l] \\
X \ar[d] & X \times_Y V \ar[l] \ar[d] & W \amalg R \ar[l] \\
Y & V \ar[l]
}
$$
$V$ を縮小することにより、$W'$ が $W$ に写ると仮定できる。
実際、$R$ の $W'$ における逆像の像を取り除けばよい。これは
$W' \to V$ が有限であるため閉集合であり、$v$ を含まない。
このとき、$W' \to W$ は有限である。なぜなら $W \to V$ と $W' \to V$ がともに有限だからである。
したがって $W' \to W$ は有限エタールであり、$w$ 上のファイバーには
$\kappa(w) = \kappa(w')$ を満たす点がちょうど一つある。ゆえに $W' \to W$ は、
ある開近傍 $W^\circ$（$w$ を含む）上で同型である。\'Etale Morphisms, 補題
\ref{etale-lemma-finite-etale-one-point} を参照せよ。
$W \to V$ は有限なので、$W \setminus W^\circ$ の像は閉部分集合
$T$ であり、これは $V$ の部分集合で $v$ を含まない。したがって $V$ を
$V \setminus T$ で置き換えると、$W' \to W$ が同型であると仮定できる。
これで分解 $X \times_Y V = W \amalg R$ と射 $W \to U$ は所望の性質をもち、結論を得る。
\end{proof}

\begin{lemma}
\label{etale-cohomology-lemma-integral-A}
$f : X \to Y$ をスキームの整射とする。このとき性質 (A) が成り立つ。
\end{lemma}

\begin{proof}
$U \to X$ をエタール射とし、$u \in U$ を点とする。
エタール射 $V \to Y$、非交和分解
$X \times_Y V = W \amalg W'$、および $X$-射 $W \to U$ であって、
$u$ がその像に含まれるものを見つければよい。
$U$ と $Y$ を縮小して、$U$ と $Y$ がアフィンであると仮定してよい。
この場合、整射は定義によりアフィンなので、$X$ もアフィンである。
$Y = \Spec(A)$、$X = \Spec(B)$、$U = \Spec(C)$ と書く。このとき
$A \to B$ は整環準同型であり、$B \to C$ はエタール環準同型である。
Algebra, 補題 \ref{algebra-lemma-etale} により、有限 $A$-部分代数
$B' \subset B$ とエタール環準同型 $B' \to C'$ であって
$C = B \otimes_{B'} C'$ を満たすものが存在する。したがって問題は、エタール射
$U' = \Spec(C') \to X' = \Spec(B')$
と有限射 $X' \to Y$ の場合に帰着する。この場合、結論は
補題 \ref{etale-cohomology-lemma-locally-quasi-finite-A} から従う。
\end{proof}

\begin{lemma}
\label{etale-cohomology-lemma-when-push-pull-surjective}
スキームの射 $f : X \to Y$ に付随する小エタールトポスの射を
$f_{small} :
\Sh(X_\etale)
\to
\Sh(Y_\etale)$
と書く。次の少なくとも一方を仮定する。
\begin{enumerate}
\item $f$ は整射である。
\item $f$ は分離かつ局所準有限である。
\end{enumerate}
このとき関手
$f_{small, *} :
\textit{Ab}(X_\etale)
\to
\textit{Ab}(Y_\etale)$
は次の性質をもつ。
\begin{enumerate}
\item 射
$f_{small}^{-1}f_{small, *}\mathcal{F} \to \mathcal{F}$
は常に全射である。
\item $f_{small, *}$ は忠実である。
\item $f_{small, *}$ は単射と全射を反映する。
\end{enumerate}
\end{lemma}

\begin{proof}
補題 \ref{etale-cohomology-lemma-locally-quasi-finite-A}、
\ref{etale-cohomology-lemma-integral-A}、および
\ref{etale-cohomology-lemma-property-A-implies} を組み合わせればよい。
\end{proof}



\section{性質 (B)}
\label{etale-cohomology-section-B}

\noindent
性質 (B) の定義については、節 \ref{etale-cohomology-section-monomorphisms} を参照せよ。

\begin{lemma}
\label{etale-cohomology-lemma-property-B-implies}
スキームの射 $f : X \to Y$ に対して (B) が成り立つと仮定する。このとき関手
$f_{small, *} :
\Sh(X_\etale)
\to
\Sh(Y_\etale)$
は全射を全射に送る。
\end{lemma}

\begin{proof}
これは Sites, 補題 \ref{sites-lemma-weaker} から従う。
\end{proof}

\begin{lemma}
\label{etale-cohomology-lemma-simplify-B}
スキームの射 $f : X \to Y$ をとり、次を仮定する。
\begin{enumerate}
\item $V \to Y$ はスキームのエタール射である。
\item $\{U_i \to X \times_Y V\}$ はエタール被覆である。
\item $v \in V$ は点である。
\end{enumerate}
このような任意のデータに対し、エタール近傍
$(V', v') \to (V, v)$、非交和分解
$X \times_Y V' = \coprod W'_i$、および射 $W'_i \to U_i$ であって
$X \times_Y V$ 上のものが存在すると仮定する。このとき性質 (B) が成り立つ。
\end{lemma}

\begin{proof}
省略する。
\end{proof}

\begin{lemma}
\label{etale-cohomology-lemma-finite-B}
$f : X \to Y$ をスキームの有限射とする。このとき性質 (B) が成り立つ。
\end{lemma}

\begin{proof}
エタール射 $V \to Y$、エタール被覆 $\{U_i \to X \times_Y V\}$、および
点 $v \in V$ を考える。補題 \ref{etale-cohomology-lemma-simplify-B} にあるような
$V' \to V$、分解、および射を見つけなければならない。
$V$ と $Y$ を縮小できるので、$V$ と $Y$ がアフィンであると仮定してよい。
$X$ は $Y$ 上有限だから、$X$ もアフィンである。
証明の途中で（有限回）、$(V, v)$ をエタール近傍
$(V', v')$ で置き換え、それに応じて被覆
$\{U_i \to X \times_Y V\}$ を $\{V' \times_V U_i \to X \times_Y V'\}$ で
置き換えてよい。

\medskip\noindent
$X \times_Y V \to V$ は有限なので、相異なる点
$x_1, \ldots, x_n \in X \times_Y V$ は有限個であり、これらは $v$ に写る。
More on Morphisms, 補題
\ref{more-morphisms-lemma-etale-splits-off-quasi-finite-part-technical-variant}
を $X \times_Y V \to V$ と、点 $x_1, \ldots, x_n$（$v$ 上にある）に適用すると、
エタール近傍 $(V', v') \to (V, v)$ であって次を満たすものが得られる。
$$
X \times_Y V' = R \amalg \coprod T_a
$$
ここで $T_a \to V'$ は有限で、ちょうど一つの点 $p_a$ をもち、これは $v'$ 上にあり、
$\kappa(v') \subset \kappa(p_a)$ は純非分離であり、さらに
$R \to V'$ の $v'$ 上のファイバーは空である。
$X \to Y$ は有限なので、$R \to V'$ も有限である。したがって $V'$ を縮小して
$R = \emptyset$ と仮定できる。ゆえに
$X \times_Y V = X_1 \amalg \ldots \amalg X_n$ と仮定でき、点
$x_l \in X_l$ は $v$ 上にちょうど一つで、さらに
$\kappa(v) \subset \kappa(x_l)$ は純非分離である。この性質はエタール近傍
$(V, v)$ を細分しても保たれることに注意する。

\medskip\noindent
各 $l$ に対して $i_l$ と点 $u_l \in U_{i_l}$ を選び、後者が $x_l$ に写るようにする。
次に、有限射 $X \times_Y V \to V$、エタール射
$U_{i_l} \to X \times_Y V$、および点 $u_l$ に性質 (A) を適用する。
これは補題 \ref{etale-cohomology-lemma-integral-A} により許される。
その結果、エタール近傍 $(V', v') \to (V, v)$ と分解
$$
X \times_Y V' = W_l \amalg R_l
$$
および $X$-射 $a_l : W_l \to U_{i_l}$ であって、その像が $u_{i_l}$ を含むものを得る。
図式で表すと次のようになる。
$$
\xymatrix{
& & & U_{i_l} \ar[d] & \\
W_l \ar[rrru] \ar[r] & W_l \amalg R_l \ar@{=}[r] &
X \times_Y V' \ar[r] \ar[d] &
X \times_Y V \ar[r] \ar[d] & X \ar[d] \\
& & V' \ar[r] & V \ar[r] & Y
}
$$
$(V, v)$ を $(V', v')$ で置き換えると、各 $x_l$ は開かつ閉な近傍
$W_l$ に含まれ、包含射 $W_l \to X \times_Y V$ は
$U_i \to X \times_Y V$ を経由する（ある $i$ について）。$W_l$ を
$W_l \cap X_l$ で置き換えると、
これらの開かつ閉な集合は互いに素であり、さらに
$\{x_1, \ldots, x_n\} \subset W_1 \cup \ldots \cup W_n$ となることが分かる。
$X \times_Y V \to V$ は有限なので、$V$ を縮小して、所望のとおり
$X \times_Y V = W_1 \amalg \ldots \amalg W_n$ と仮定できる。
\end{proof}

\begin{lemma}
\label{etale-cohomology-lemma-integral-B}
$f : X \to Y$ をスキームの整射とする。このとき性質 (B) が成り立つ。
\end{lemma}

\begin{proof}
エタール射 $V \to Y$、エタール被覆 $\{U_i \to X \times_Y V\}$、および
点 $v \in V$ を考える。補題 \ref{etale-cohomology-lemma-simplify-B} にあるような
$V' \to V$、分解、および射を見つけなければならない。
$V$ と $Y$ を縮小できるので、$V$ と $Y$ がアフィンであると仮定してよい。
$X$ は $Y$ 上整なので、$X$ と $X \times_Y V$ もアフィンである。
被覆 $\{U_i \to X \times_Y V\}$ を細分できるので、
$\{U_i \to X \times_Y V\}_{i = 1, \ldots, n}$ が標準エタール被覆であると
仮定してよい。$Y = \Spec(A)$、$X = \Spec(B)$、
$V = \Spec(C)$、$U_i = \Spec(B_i)$ と書く。このとき $A \to B$ は
整環準同型であり、$B \otimes_A C \to B_i$ はエタール環準同型である。
Algebra, 補題 \ref{algebra-lemma-etale} により、有限 $A$-部分代数
$B' \subset B$ とエタール環準同型 $B' \otimes_A C \to B'_i$
（各 $i = 1, \ldots, n$ に対するもの）であって、$B_i = B \otimes_{B'} B'_i$ を
満たすものが存在する。したがって問題はエタール被覆
$\{\Spec(B'_i) \to X' \times_Y V\}_{i = 1, \ldots, n}$
の場合に帰着する。ここで $X' = \Spec(B')$ は $Y$ 上有限である。
この場合、結論は補題 \ref{etale-cohomology-lemma-finite-B} から従う。
\end{proof}

\begin{lemma}
\label{etale-cohomology-lemma-what-integral}
スキームの射 $f : X \to Y$ をとる。$f$ は整射（たとえば有限射）であると
仮定する。このとき次が成り立つ。
\begin{enumerate}
\item $f_{small, *}$ は全射を全射に送る（集合の層とアーベル層のいずれについても）。
\item $f_{small}^{-1}f_{small, *}\mathcal{F} \to \mathcal{F}$
は任意のアーベル層 $\mathcal{F}$ に対して $X_\etale$ 上で全射である。
\item
$f_{small, *} :
\textit{Ab}(X_\etale)
\to
\textit{Ab}(Y_\etale)$
は忠実であり、単射と全射を反映する。
\item
$f_{small, *} :
\textit{Ab}(X_\etale)
\to
\textit{Ab}(Y_\etale)$
は完全である。
\end{enumerate}
\end{lemma}

\begin{proof}
(2)、(3) は補題 \ref{etale-cohomology-lemma-when-push-pull-surjective} ですでに示した。
(1) は補題 \ref{etale-cohomology-lemma-integral-B} と \ref{etale-cohomology-lemma-property-B-implies} から従う。
(4) は (1) の帰結である。Modules on Sites,
補題 \ref{sites-modules-lemma-exactness} を参照せよ。
\end{proof}








\section{性質 (C)}
\label{etale-cohomology-section-C}

\noindent
性質 (C) の定義については、節 \ref{etale-cohomology-section-monomorphisms} を参照せよ。

\begin{lemma}
\label{etale-cohomology-lemma-property-C-implies}
スキームの射 $f : X \to Y$ に対して (C) が成り立つと仮定する。このとき関手
$f_{small, *} :
\Sh(X_\etale)
\to
\Sh(Y_\etale)$
は単射と全射を反映する。
\end{lemma}

\begin{proof}
Sites, 補題 \ref{sites-lemma-cover-from-below} から従う。
性質 (C) により関手
$Y_\etale \to X_\etale$、$V \mapsto X \times_Y V$ が
Sites, 補題 \ref{sites-lemma-cover-from-below} の仮定を満たすことの確認は省略する。
\end{proof}

\begin{remark}
\label{etale-cohomology-remark-property-C-strong}
$f : X \to Y$ が開埋め込みならば性質 (C) が成り立つ。実際、
$U \in \Ob(X_\etale)$ ならば $U$ を $Y_\etale$ の対象ともみなすことができ、
$U \times_Y X = U$ である。したがって性質 (C) から
$f_{small, *}$ の完全性は従わない。一般に開埋め込みについてこれは成り立たないからである。
\end{remark}

\begin{lemma}
\label{etale-cohomology-lemma-property-C-closed-implies}
スキームの射 $f : X \to Y$ をとる。任意のエタール射 $V \to Y$ に対し、
次が成り立つと仮定する。
\begin{enumerate}
\item $X \times_Y V \to V$ は性質 (C) をもつ。
\item $X \times_Y V \to V$ は閉射である。
\end{enumerate}
このとき関手
$Y_\etale \to X_\etale$, $V \mapsto X \times_Y V$
はほぼ余連続である。Sites, 定義
\ref{sites-definition-almost-cocontinuous} を参照せよ。
\end{lemma}

\begin{proof}
$V \to Y$ を $Y_\etale$ の対象とし、
$\{U_i \to X \times_Y V\}_{i \in I}$ を $X_\etale$ の被覆とする。
仮定 (1) により、各 $i$ に対してエタール射 $h_i : V_i \to V$ と全射
$X \times_Y V_i \to U_i$ であって $X \times_Y V$ 上のものを選べる。
$\bigcup h_i(V_i) \subset V$ は、閉集合
$Z = \Im(X \times_Y V \to V)$ を含む開集合であることに注意する。
$h_0 : V_0 = V \setminus Z \to V$ を開埋め込みとする。
$\{V_i \to V\}_{i \in I \cup \{0\}}$ がエタール被覆であって、各
$i \in I \cup \{0\}$ に対して次のいずれかが成り立つことは明らかである。
$V_i \times_Y X = \emptyset$（すなわち $i = 0$）であるか、または
$V_i \times_Y X \to V \times_Y X$ が $U_i \to X \times_Y V$ を経由する
（$i \not = 0$）。したがって関手 $Y_\etale \to X_\etale$ はほぼ余連続である。
\end{proof}

\begin{lemma}
\label{etale-cohomology-lemma-integral-homeo-onto-image-C}
$f : X \to Y$ をスキームの整射で、$X$ から $Y$ の閉部分集合への
同相写像を定めるものとする。このとき性質 (C) が成り立つ。
\end{lemma}

\begin{proof}
$g : U \to X$ をエタール射とする。対象 $V \to Y$（$Y_\etale$ のもの）と、
全射 $X \times_Y V \to U$ であって $X$ 上のものを見つけなければならない。
任意の $u \in U$ に対し、対象 $V_u \to Y$（$Y_\etale$ のもの）と射
$h_u : X \times_Y V_u \to U$ であって、$X$ 上のもので
$u \in \Im(h_u)$ を満たすものが存在すると仮定する。このとき
$V = \coprod V_u$ および $h = \coprod h_u$ ととればよく、結論を得る。
したがって、点 $u \in U$ が与えられたとき、上のような組
$(V_u, h_u)$ を見つければよい。そのために $U$ を縮小し、$U$ が
アフィンであると仮定してよい。この場合 $g : U \to X$ は局所準有限である。
$g^{-1}(g(\{u\})) = \{u, u_2, \ldots, u_n\}$ とおく。特殊化
$u_i \leadsto u$ は存在しないので、$U$ をアフィン近傍で置き換えて
$g^{-1}(g(\{u\})) = \{u\}$ とできる。

\medskip\noindent
像 $g(U) \subset X$ は開であるから、$f(g(U))$ は $Y$ の局所閉部分集合である。
開集合 $V \subset Y$ を $f(g(U)) = f(X) \cap V$ となるように選ぶ。このとき $g$ は
$X \times_Y V$ を経由し、得られる $\{U \to X \times_Y V\}$ はエタール被覆である。
$f$ は性質 (B) をもつので（補題 \ref{etale-cohomology-lemma-integral-B} を参照）、
エタール被覆 $\{V_j \to V\}$ であって、射
$X \times_Y V_j \to X \times_Y V$ が $U$ を経由するものが存在する。
したがって $V' = \coprod V_j$ は $Y$ 上エタールであり、射
$h : X \times_Y V' \to U$ であって、その像が $g(U)$ 上へ全射となるものが存在する。
$u$ はそのファイバーにおける唯一の点なので、$h$ の像に含まれ、結論を得る。
\end{proof}

\noindent
次の補題を、整かつ普遍的単射な射に対する $f_{small, *}$ についての
究極の真理へ向かう旅の中継地点\footnote{中継地点とは、長い旅の途中で
食事と休息のために人々が立ち寄る場所である。}と考えてほしい。

\begin{lemma}
\label{etale-cohomology-lemma-integral-universally-injective}
スキームの射 $f : X \to Y$ をとる。$f$ は普遍的単射かつ整である
（たとえば閉埋め込みである）と仮定する。このとき次が成り立つ。
\begin{enumerate}
\item
$f_{small, *} :
\Sh(X_\etale)
\to
\Sh(Y_\etale)$
は単射と全射を反映する。
\item
$f_{small, *} :
\Sh(X_\etale)
\to
\Sh(Y_\etale)$
は押し出しおよび余等化子（より一般には有限連結余極限）と可換である。
\item $f_{small, *}$ は全射を全射に送る（集合の層とアーベル層のいずれについても）。
\item 射
$f_{small}^{-1}f_{small, *}\mathcal{F} \to \mathcal{F}$
は任意の層（集合の層またはアーベル群の層）$\mathcal{F}$ に対して
$X_\etale$ 上で全射である。
\item 関手 $f_{small, *}$ は忠実である（集合の層とアーベル層のいずれについても）。
\item
$f_{small, *} :
\textit{Ab}(X_\etale)
\to
\textit{Ab}(Y_\etale)$
は完全である。
\item 関手 $Y_\etale \to X_\etale$、$V \mapsto X \times_Y V$ は
ほぼ余連続である。
\end{enumerate}
\end{lemma}

\begin{proof}
補題 \ref{etale-cohomology-lemma-integral-A}、
\ref{etale-cohomology-lemma-integral-B}、および
\ref{etale-cohomology-lemma-integral-homeo-onto-image-C}
により、射 $f$ は性質 (A)、(B)、(C) をもつ。さらに補題
\ref{etale-cohomology-lemma-property-C-closed-implies} により、関手
$Y_\etale \to X_\etale$ はほぼ余連続である。ここで次が成り立つ。
\begin{enumerate}
\item 性質 (C) から (1) が従う。補題 \ref{etale-cohomology-lemma-property-C-implies} を参照せよ。
\item ほぼ余連続であることから (2) が従う。Sites, 補題
\ref{sites-lemma-morphism-of-sites-almost-cocontinuous} を参照せよ。
\item 性質 (B) から (3) が従う。補題 \ref{etale-cohomology-lemma-property-B-implies} を参照せよ。
\end{enumerate}
性質 (4)、(5)、(6) は最初の三つから形式的に従う。Sites, 補題
\ref{sites-lemma-exactness-properties} および Modules on Sites, 補題
\ref{sites-modules-lemma-exactness} を参照せよ。性質 (7) は上で示した。
\end{proof}




\section{小エタールサイトの位相不変性}
\label{etale-cohomology-section-topological-invariance}

\noindent
次の定理では、小エタールサイトが次の意味で位相不変であることを示す。
スキームの射 $f : X \to Y$ が普遍同相写像ならば、サイトとして
$X_\etale \cong Y_\etale$ である。これは \'Etale Morphisms, 定理
\ref{etale-theorem-remarkable-equivalence} の結果を強めるものである。
まず射について結果を証明し、その後で圏について結果を述べる。

\begin{theorem}
\label{etale-cohomology-theorem-etale-topological}
$X$ と $Y$ を基礎スキーム $S$ 上の二つのスキームとする。
$S' \to S$ を普遍同相写像とし、$X'$（それぞれ $Y'$）を
$S'$ への基底変換と書く。$X$ が $S$ 上エタールならば、写像
$$
\Mor_S(Y, X) \longrightarrow \Mor_{S'}(Y', X')
$$
は全単射である。
\end{theorem}

\begin{proof}
$Y \to S$ による基底変換の後、$Y = S$ と仮定してよい。
したがって $X$ と $Y$ はともに $S$ 上エタールであると仮定する。
言い換えれば、定理は、$S$ 上エタールなスキームの圏から
$S'$ 上エタールなスキームの圏への基底変換関手が忠実充満であると主張している。

\medskip\noindent
忘却関手
\begin{equation}
\label{etale-cohomology-equation-descent-etale-forget}
\begin{matrix}
\text{降下データ }(X', \varphi')\text{（}S'/S\text{ に関する）} \\
\text{ただし }X'\text{ は }S'\text{ 上エタール}
\end{matrix}
\longrightarrow
\text{スキーム }X'\text{（}S'\text{ 上エタール）}
\end{equation}
を考える。この関手は圏同値であると主張する。一方、関手
\begin{equation}
\label{etale-cohomology-equation-descent-etale}
\text{スキーム }X\text{（}S\text{ 上エタール）} \longrightarrow
\begin{matrix}
\text{降下データ }(X', \varphi')\text{（}S'/S\text{ に関する）} \\
\text{ただし }X'\text{ は }S'\text{ 上エタール}
\end{matrix}
\end{equation}
は \'Etale Morphisms, 補題 \ref{etale-lemma-fully-faithful-cases} により
忠実充満である。したがって主張から定理が従う。

\medskip\noindent
主張を証明する。普遍同相写像は、整かつ普遍的単射かつ全射な射と
同じものであることを思い出そう。Morphisms, 補題
\ref{morphisms-lemma-universal-homeomorphism} を参照せよ。
とくに、Morphisms, 補題 \ref{morphisms-lemma-universally-injective} により
対角射 $\Delta : S' \to S' \times_S S'$ は肥厚である。
したがって \'Etale Morphisms, 定理 \ref{etale-theorem-etale-topological} により、
エタール射 $X' \to S'$ が与えられると、一意な同型
$$
\varphi' : X' \times_S S' \to S' \times_S X'
$$
が存在する。これは $S' \times_S S'$ 上エタールなスキームの同型で、
$\Delta$ により引き戻すと $\text{id} : X' \to X'$ となり、これは $S'$ 上の射である。
$S' \to S' \times_S S' \times_S S'$ も肥厚である
（全単射かつ閉埋め込みである）から、$(X', \varphi')$ は
$S'/S$ に関する降下データである。$\varphi'$ の構成が標準的であることから、
これは $S'$ 上エタールなスキーム間の射と両立する。
言い換えれば、関手 (\ref{etale-cohomology-equation-descent-etale-forget}) の準逆
$X' \mapsto (X', \varphi')$ を得る。これで主張、したがって定理が証明された。
\end{proof}

\begin{theorem}
\label{etale-cohomology-theorem-topological-invariance}
\begin{reference}
\cite[IV Theorem 18.1.2]{EGA}
\end{reference}
スキームの射 $f : X \to Y$ をとる。$f$ は整かつ普遍的単射かつ全射であると
仮定する（すなわち $f$ は普遍同相写像である。Morphisms, 補題
\ref{morphisms-lemma-universal-homeomorphism} を参照せよ）。関手
$$
V \longmapsto V_X = X \times_Y V
$$
は次の圏同値を定める。
$$
\{
\text{スキーム }V\text{（}Y\text{ 上エタール）}
\}
\leftrightarrow
\{
\text{スキーム }U\text{（}X\text{ 上エタール）}
\}
$$
\end{theorem}

\noindent
二つの証明を与える。第一の証明は、全射整射に関する準コンパクト・分離・
エタール射の降下の有効性を用いる。第二の証明は、本章で先に論じた
性質 (A)、(B)、(C) に関する結果を用いる。

\begin{proof}[第一の証明]
定理 \ref{etale-cohomology-theorem-etale-topological} により、この関手は忠実充満である。
あとは本質的全射であることを示せばよい。
$U \to X$ をスキームのエタール射とする。

\medskip\noindent
$U$ と $Y$ がアフィンならば結果が成り立つと仮定する。この場合、アフィン開被覆
$U = \bigcup U_i$ を選び、各 $U_i$ が $Y$ のあるアフィン開集合へ写るようにする。
アフィンの場合の仮定により、エタール射 $V_i \to Y$ であって
$X \times_Y V_i \cong U_i$ が $X$ 上のスキームとして成り立つものを選べる。
$V_{i, i'} \subset V_i$ を、その台位相空間が $U_i \cap U_{i'}$ に対応する
開部分スキームとする。同型
$$
X \times_Y V_{i, i'} \cong U_i \cap U_{i'} \cong X \times_Y V_{i', i}
$$
が $X$ 上のスキームとして存在するので、忠実充満性により同型
$\theta_{i, i'} : V_{i, i'} \to V_{i', i}$ を得る。これは $Y$ 上の
スキームの同型である。これらの同型が
Schemes, 節 \ref{schemes-section-glueing-schemes} のコサイクル条件を満たすことの
確認は省略する。Schemes, 補題 \ref{schemes-lemma-glue-schemes} を適用し、
スキーム $V \to Y$ を、スキーム $V_i$ を同一視 $\theta_{i, i'}$ に沿って
貼り合わせることで得る。構成により $V \to Y$ がエタールで、
$X \times_Y V \cong U$ となることは明らかである。

\medskip\noindent
したがって $U$ と $Y$ がアフィンの場合に補題を示せば十分である。
定理 \ref{etale-cohomology-theorem-etale-topological} の証明で、$U$ には一意な降下データ
$(U, \varphi)$（$X/Y$ に関するもの）が付随することを示した。\'Etale Morphisms, 命題
\ref{etale-proposition-effective} を適用できる
（アフィンの場合への帰着により $U \to X$ はエタールであるだけでなく
準コンパクトかつ分離でもある）。したがってエタール射 $V \to Y$ が存在し、
$X \times_Y V \cong U$ となる。これで証明は完了した。
\end{proof}

\begin{proof}[第二の証明]
定理 \ref{etale-cohomology-theorem-etale-topological} により、この関手は忠実充満である。
あとは本質的全射であることを示せばよい。
$U \to X$ をスキームのエタール射とする。

\medskip\noindent
$U$ と $Y$ がアフィンならば結果が成り立つと仮定する。この場合、アフィン開被覆
$U = \bigcup U_i$ を選び、各 $U_i$ が $Y$ のあるアフィン開集合へ写るようにする。
アフィンの場合の仮定により、エタール射 $V_i \to Y$ であって
$X \times_Y V_i \cong U_i$ が $X$ 上のスキームとして成り立つものを選べる。
$V_{i, i'} \subset V_i$ を、その台位相空間が $U_i \cap U_{i'}$ に対応する
開部分スキームとする。同型
$$
X \times_Y V_{i, i'} \cong U_i \cap U_{i'} \cong X \times_Y V_{i', i}
$$
が $X$ 上のスキームとして存在するので、忠実充満性により同型
$\theta_{i, i'} : V_{i, i'} \to V_{i', i}$ を得る。これは $Y$ 上の
スキームの同型である。これらの同型が Schemes, 節
\ref{schemes-section-glueing-schemes} のコサイクル条件を満たすことの確認は省略する。
Schemes, 補題 \ref{schemes-lemma-glue-schemes} を適用し、スキーム
$V \to Y$ を、スキーム $V_i$ を同一視 $\theta_{i, i'}$ に沿って貼り合わせることで得る。
構成により $V \to Y$ がエタールで、$X \times_Y V \cong U$ となることは明らかである。

\medskip\noindent
したがって、関手
\begin{equation}
\label{etale-cohomology-equation-affine-etale}
\{
\text{アフィンスキーム }V\text{（}Y\text{ 上エタール）}
\}
\leftrightarrow
\{
\text{アフィンスキーム }U\text{（}X\text{ 上エタール）}
\}
\end{equation}
が、$X$ と $Y$ がアフィンの場合に本質的全射であることを示せば十分である。

\medskip\noindent
$U \to X$ を $X$ 上エタールなアフィンスキームとする。
エタール（かつアフィン）射 $V \to Y$ であって、
$X \times_Y V$ が $U$ と $X$ 上で同型になるものを見つけなければならない。
アフィン間のエタール射のファイバーは普遍的に有界であることに注意する。
Morphisms, 補題 \ref{morphisms-lemma-etale-locally-quasi-finite} および
\ref{morphisms-lemma-locally-quasi-finite-qc-source-universally-bounded} を参照せよ。
したがって、整数 $n$（$U \to X$ のファイバーの次数を抑えるもの）に関して帰納法を使える。
エタール射の場合のこの整数の記述については Morphisms, 補題
\ref{morphisms-lemma-etale-universally-bounded} を参照せよ。
$n = 1$ ならば $U \to X$ は開埋め込みであり（\'Etale Morphisms, 定理
\ref{etale-theorem-etale-radicial-open} を参照）、結論は明らかである。
$n > 1$ と仮定する。

\medskip\noindent
補題 \ref{etale-cohomology-lemma-integral-homeo-onto-image-C} により、スキームのエタール射
$W \to Y$ と全射 $W_X \to U$ であって $X$ 上のものが存在する。
$U$ は準コンパクトなので、$W$ を $W$ の有限個のアフィン開集合の
非交和で置き換え、$W$ もアフィンであると仮定してよい。次の図式を得る。
$$
\xymatrix{
U \ar[d] & U \times_Y W \ar[l] \ar[d] & W_X \amalg R \ar@{=}[l]\\
X \ar[d] & W_X \ar[l] \ar[d] \\
Y & W \ar[l]
}
$$
この非交和分解が生じるのは、構成によりアフィンスキームのエタール射
$U \times_Y W \to W_X$ が切断をもつからである。ここで射
$R \to X \times_Y W$ はアフィンスキームのエタール射であり、そのファイバーの
次数は $n - 1$ で普遍的に抑えられる。したがって帰納法の仮定により、
エタールなスキーム $V' \to W$ であって
$R \cong W_X \times_W V'$ を満たすものが存在する。
$V'' = W \amalg V'$ とおくと、スキーム $V''$（$W$ 上エタール）であって、
$W_X$ への基底変換が $U \times_Y W$ と $X \times_Y W$ 上で同型になるものを得る。

\medskip\noindent
ここで降下を用いると、$V$ を $Y$ 上に得ることができ、その $X$ への基底変換は
$U$ と $X$ 上で同型になる。実際、普遍同相写像
$X \times_Y (W \times_Y W) \to (W \times_Y W)$ に対応する関手
(\ref{etale-cohomology-equation-affine-etale}) の忠実充満性により、一意な同型
$\varphi : V'' \times_Y W \to W \times_Y V''$ が存在する。その
$X \times_Y (W \times_Y W)$ への基底変換は、$U \times_Y W$ に対する
標準的な降下データ（$X \times_Y W$ 上のもの）である。とくに $\varphi$ は
コサイクル条件を満たす。したがって Descent, 補題
\ref{descent-lemma-affine} により $\varphi$ は有効である
（上のすべてのスキームがアフィンであることを思い出そう）。
よって $V \to Y$ と同型 $V'' \cong W \times_Y V$ を得て、
$W \times_Y V/W/Y$ 上の標準的な降下データは $\varphi$ と一致する。
Descent, 補題 \ref{descent-lemma-descending-property-etale} により
$V \to Y$ はエタールであることに注意する。さらに、同型
$V_X \cong U$ が存在し、これは次の同型を降下して得られる。
$$
V_X \times_X W_X =
X \times_Y V \times_Y W =
(X \times_Y W) \times_W (W \times_Y V) \cong
W_X  \times_W V'' \cong U \times_Y W
$$
これは構成により存在する。いくつかの詳細は省略する。
\end{proof}

\begin{remark}
\label{etale-cohomology-remark-affine-inside-equivalence}
定理 \ref{etale-cohomology-theorem-topological-invariance} の状況では、$V \mapsto V_X$ が、
エタール射 $V \to Y$（$V$ がアフィンであるもの）と、エタール射
$U \to X$（$U$ がアフィンであるもの）の間の圏同値を誘導することも成り立つ。これはたとえば
Limits, 命題 \ref{limits-proposition-affine} から従う。
\end{remark}

\begin{proposition}[エタールコホモロジーの位相不変性]
\label{etale-cohomology-proposition-topological-invariance}
$X_0 \to X$ をスキームの普遍同相写像とする
（たとえば冪零イデアル層が定める閉埋め込み）。このとき次が成り立つ。
\begin{enumerate}
\item エタールサイト $X_\etale$ と $(X_0)_\etale$ は同型である。
\item エタールトポス $\Sh(X_\etale)$ と $\Sh((X_0)_\etale)$ は同値である。
\item $H^q_\etale(X, \mathcal{F}) = H^q_\etale(X_0, \mathcal{F}|_{X_0})$
がすべての $q$ と任意のアーベル層 $\mathcal{F}$ に対して $X_\etale$ 上で成り立つ。
\end{enumerate}
\end{proposition}

\begin{proof}
圏同値 $X_\etale \to (X_0)_\etale$ は定理
\ref{etale-cohomology-theorem-topological-invariance} によって与えられる。この同値のもとで
エタール被覆が対応することの証明は省略する。したがって (1) が成り立ち、
(2) と (3) は (1) から形式的に従う。
\end{proof}






\section{閉埋め込みと順像}
\label{etale-cohomology-section-closed-immersions}

\noindent
命題 \ref{etale-cohomology-proposition-closed-immersion-pushforward} を正しい一般性のもとで
述べて証明する前に、まず閉埋め込みの場合について簡潔に述べ、証明する。
実際、これまでの議論のいくつかは閉埋め込みの場合にはかなり追いやすくなるので、
ここでは簡略化した形でそれらを繰り返す。

\medskip\noindent
この節の以下では、$i : Z \to X$ を閉埋め込みとする。関手
$$
\Sch/X \longrightarrow \Sch/Z, \quad
U \longmapsto U_Z = Z \times_X U
$$
を、表示のとおり $U \mapsto U_Z$ と書く。閉埋め込みであることは任意の基底変換で
保たれるため、スキーム $U_Z$ は $U$ の閉部分スキームである。

\begin{lemma}
\label{etale-cohomology-lemma-closed-immersion-almost-full}
$i : Z \to X$ をスキームの閉埋め込みとする。
$U, U'$ を $X$ 上エタールなスキームとし、$h : U_Z \to U'_Z$ を
$Z$ 上の射とする。このとき図式
$$
\xymatrix{
U & W \ar[l]_a \ar[r]^b & U'
}
$$
が $X_\etale$ に存在し、$a_Z : W_Z \to U_Z$ は同型であり、
$h = b_Z \circ (a_Z)^{-1}$ である。
\end{lemma}

\begin{proof}
スキーム $M = U \times_X U'$ を考える。部分集合
$\Gamma_h \subset M_Z$ は $h$ のグラフであり、開である。たとえば、$\Gamma_h$ はエタール射
$\text{pr}_{1, Z} : M_Z \to U_Z$ の切断の像だからである。
エタール射、命題 \ref{etale-proposition-properties-sections} を参照されたい。
したがって、開部分スキーム $W \subset M$ が存在して、閉部分集合
$M_Z$ との交わりは $\Gamma_h$ である。$a = \text{pr}_1|_W$ および
$b = \text{pr}_2|_W$ とおく。
\end{proof}

\begin{lemma}
\label{etale-cohomology-lemma-closed-immersion-almost-essentially-surjective}
$i : Z \to X$ をスキームの閉埋め込みとし、
$V \to Z$ をスキームのエタール射とする。このとき、エタール射
$U_i \to X$ と射 $U_{i, Z} \to V$ が存在して、
$\{U_{i, Z} \to V\}$ は $V$ のザリスキー被覆となる。
\end{lemma}

\begin{proof}
$V$ のザリスキー被覆であって、$U_Z$ の形のスキームからなるものを
見つければよい（ここで $U$ は $X$ 上エタール）ので、$X$ と $V$ 上で
ザリスキー局所化してよい。
したがって、$X$ と $V$ はアフィンであると仮定できる。アフィンの場合、これは
代数、補題 \ref{algebra-lemma-lift-etale} である。
\end{proof}

\noindent
$\overline{x} : \Spec(k) \to X$ が $X$ の幾何学的点ならば、
$\overline{x}$ は閉部分スキーム $Z$ を（ただ一通りに）経由するか、または
$Z_{\overline{x}} = \emptyset$ である。$\overline{x}$ が $Z$ を経由するとき、
$\overline{x}$ は実際に $Z$ の幾何学的点なので、そのように呼び、これを表すために
記法「$\overline{x} \in Z$」を用いる。

\begin{lemma}
\label{etale-cohomology-lemma-stalk-pushforward-closed-immersion}
$i : Z \to X$ をスキームの閉埋め込みとする。
$\mathcal{G}$ を $Z_\etale$ 上の集合の層とし、
$\overline{x}$ を $X$ の幾何学的点とする。このとき
$$
(i_{small, *}\mathcal{G})_{\overline{x}} =
\left\{
\begin{matrix}
* & \text{if} & \overline{x} \not \in Z \\
\mathcal{G}_{\overline{x}} & \text{if} & \overline{x} \in Z
\end{matrix}
\right.
$$
である。ここで $*$ は一元集合を表す。
\end{lemma}

\begin{proof}
$U = X \setminus Z$ とおく。
$i_{small, *}\mathcal{G}|_{U_\etale} = *$ は $U$ 上のエタール層の圏における
終対象、すなわち $U$ 上エタールな各スキームに一元集合を対応させる層であることに
注意する。これは $(i_{small, *}\mathcal{G})_{\overline{x}}$ の値を、
$\overline{x} \not \in Z$ の場合について説明する。

\medskip\noindent
次に、$\overline{x} \in Z$ と仮定する。このとき
$$
(i_{small, *}\mathcal{G})_{\overline{x}}
=
\colim_{(U, \overline{u})} \mathcal{G}(U_Z)
$$
であり、その一方で
$$
\mathcal{G}_{\overline{x}}
=
\colim_{(V, \overline{v})} \mathcal{G}(V).
$$
$\mathcal{C}_1 = \{(U, \overline{u})\}^{opp}$ を $\overline{x}$ の
$X$ におけるエタール近傍の圏の反対圏とし、
$\mathcal{C}_2 = \{(V, \overline{v})\}^{opp}$ を $\overline{x}$ の
$Z$ におけるエタール近傍の圏の反対圏とする。標準射
$$
\mathcal{G}_{\overline{x}}
\longrightarrow
(i_{small, *}\mathcal{G})_{\overline{x}}
$$
は関手 $F : \mathcal{C}_1 \to \mathcal{C}_2$、
$F(U, \overline{u}) = (U_Z, \overline{x})$ に対応する。ここで補題
\ref{etale-cohomology-lemma-closed-immersion-almost-essentially-surjective} および
\ref{etale-cohomology-lemma-closed-immersion-almost-full}
から、$\mathcal{C}_1$ は $\mathcal{C}_2$ において共終である。
圏、定義 \ref{categories-definition-cofinal} を参照されたい。
したがって、表示された矢印は同型である。
圏、補題 \ref{categories-lemma-cofinal} を参照されたい。
\end{proof}

\begin{proposition}
\label{etale-cohomology-proposition-closed-immersion-pushforward}
$i : Z \to X$ をスキームの閉埋め込みとする。
\begin{enumerate}
\item 関手
$$
i_{small, *} :
\Sh(Z_\etale)
\longrightarrow
\Sh(X_\etale)
$$
は忠実充満であり、その本質的像は、集合の層 $\mathcal{F}$ であって
$X_\etale$ 上にあり、$X \setminus Z$ への制限が $*$ と同型であるものからなる。
\item 関手
$$
i_{small, *} :
\textit{Ab}(Z_\etale)
\longrightarrow
\textit{Ab}(X_\etale)
$$
は忠実充満であり、その本質的像は、$X_\etale$ 上のアーベル層であって、
台が $Z$ に含まれるものからなる。
\end{enumerate}
いずれの場合にも、$i_{small}^{-1}$ は関手 $i_{small, *}$ の左逆である。
\end{proposition}

\begin{proof}
集合の層の場合を論じる。任意の層 $\mathcal{G}$ で $Z$ 上のものに対して、射
$i_{small}^{-1}i_{small, *}\mathcal{G} \to \mathcal{G}$
は補題 \ref{etale-cohomology-lemma-stalk-pushforward-closed-immersion}
（および定理 \ref{etale-cohomology-theorem-exactness-stalks}）により同型である。
これから形式的に $i_{small, *}$ は忠実充満であることが従う。
サイト、補題 \ref{sites-lemma-exactness-properties} を参照されたい。
$i_{small, *}\mathcal{G}|_{U_\etale} \cong *$ であることは明らかである。
ここで $U = X \setminus Z$ とおいた。逆に、$\mathcal{F}$ を $X$ 上の集合の層で、
$\mathcal{F}|_{U_\etale} \cong *$ を満たすものと仮定する。随伴射
$$
\mathcal{F} \longrightarrow i_{small, *}i_{small}^{-1}\mathcal{F}
$$
を考える。補題 \ref{etale-cohomology-lemma-stalk-pushforward-closed-immersion} と
\ref{etale-cohomology-lemma-stalk-pullback}
をあわせると、これは同型であることがわかる。これで (1) の証明が完了する。
(2) の証明も同一である。
\end{proof}





\section{整かつ普遍的単射な射}
\label{etale-cohomology-section-integral-universally-injective}

\noindent
ここで命題 \ref{etale-cohomology-proposition-closed-immersion-pushforward} の一般形を与える。

\begin{proposition}
\label{etale-cohomology-proposition-integral-universally-injective-pushforward}
$f : X \to Y$ を整かつ普遍的単射なスキームの射とする。
\begin{enumerate}
\item 関手
$$
f_{small, *} :
\Sh(X_\etale)
\longrightarrow
\Sh(Y_\etale)
$$
は忠実充満であり、その本質的像は、集合の層 $\mathcal{F}$ であって
$Y_\etale$ 上にあり、$Y \setminus f(X)$ への制限が $*$ と同型であるものからなる。
\item 関手
$$
f_{small, *} :
\textit{Ab}(X_\etale)
\longrightarrow
\textit{Ab}(Y_\etale)
$$
は忠実充満であり、その本質的像は、$Y_\etale$ 上のアーベル層であって、
台が $f(X)$ に含まれるものからなる。
\end{enumerate}
いずれの場合にも、$f_{small}^{-1}$ は関手 $f_{small, *}$ の左逆である。
\end{proposition}

\begin{proof}
$f$ を
$$
\xymatrix{
X \ar[r]^h & Z \ar[r]^i & Y
}
$$
と分解してよい。ここで $h$ は整、普遍的単射かつ全射であり、
$i : Z \to Y$ は閉埋め込みである。
命題 \ref{etale-cohomology-proposition-closed-immersion-pushforward} を $i$ に適用し、
定理 \ref{etale-cohomology-theorem-topological-invariance} を $h$ に適用する。
\end{proof}









\section{大サイトと順像}
\label{etale-cohomology-section-big}

\noindent
この節では、特定の種類のスキームの射に対する $f_{big, *}$ について、
いくつかの技術的結果を証明する。

\begin{lemma}
\label{etale-cohomology-lemma-monomorphism-big-push-pull}
$\tau \in \{Zariski, \etale, smooth, syntomic, fppf\}$ とする。
$f : X \to Y$ をスキームの単射とする。このとき標準射
$f_{big}^{-1}f_{big, *}\mathcal{F} \to \mathcal{F}$
は、任意の層 $\mathcal{F}$ で $(\Sch/X)_\tau$ 上のものに対して同型である。
\end{lemma}

\begin{proof}
この場合、関手 $(\Sch/X)_\tau \to (\Sch/Y)_\tau$ は連続、余連続かつ
忠実充満である。したがって、結果はサイト、補題
\ref{sites-lemma-back-and-forth} から従う。
\end{proof}

\begin{remark}
\label{etale-cohomology-remark-push-pull-shriek}
補題 \ref{etale-cohomology-lemma-monomorphism-big-push-pull} の状況では、標準射
$\mathcal{F} \to f_{big}^{-1}f_{big!}\mathcal{F}$
は、任意の集合の層 $\mathcal{F}$ で $(\Sch/X)_\tau$ 上のものに対して同型である。
証明は同じである。これはアーベル群の層についても成り立つ。ただし、
アーベル群の層に対する関手 $f_{big!}$ はサイト上の加群、節
\ref{sites-modules-section-exactness-lower-shriek} で定義され、一般には
集合の層に対する $f_{big!}$ と異なることに注意する。アーベル群の層に対する結果は、
サイト上の加群、補題 \ref{sites-modules-lemma-back-and-forth} から従う。
\end{remark}

\begin{lemma}
\label{etale-cohomology-lemma-closed-immersion-cover-from-below}
$f : X \to Y$ をスキームの閉埋め込みとする。
$U \to X$ をシントミック（それぞれ滑らか、それぞれエタール）な射とする。
このとき、シントミック（それぞれ滑らか、それぞれエタール）な射
$V_i \to Y$ と射 $V_i \times_Y X \to U$ が存在して、
$\{V_i \times_Y X \to U\}$ は $U$ のザリスキー被覆となる。
\end{lemma}

\begin{proof}
補題を $\tau = syntomic$ の場合に証明する。
問題は $U$ 上局所的である。したがって、$U$ はアフィンスキームであり、
$Y$ のあるアフィン開部分へ写ると仮定してよい。
よって、次の場合を証明することに帰着する：
$Y = \Spec(A)$、$X = \Spec(A/I)$、および
$U = \Spec(\overline{B})$ とする。ここで
$A/I \to \overline{B}$ はシントミックな環準同型である。
代数、補題 \ref{algebra-lemma-lift-syntomic} により、元
$\overline{g}_i \in \overline{B}$ であって
$\overline{B}_{\overline{g}_i} = A_i/IA_i$ が、あるシントミックな環準同型
$A \to A_i$ に対して成り立つものを見つけられる。ここで後者はシントミックな
環準同型である。これでシントミックの場合の補題が証明された。
滑らかな場合の証明も同じであるが、代数、補題
\ref{algebra-lemma-lift-smooth} を用いる。エタールの場合には代数、補題
\ref{algebra-lemma-lift-etale} を用いる。
\end{proof}

\begin{lemma}
\label{etale-cohomology-lemma-prepare-closed-immersion-almost-cocontinuous}
$f : X \to Y$ をスキームの閉埋め込みとする。
$\{U_i \to X\}$ をシントミック（それぞれ滑らか、それぞれエタール）な被覆とする。
このとき、シントミック（それぞれ滑らか、それぞれエタール）な被覆
$\{V_j \to Y\}$ が存在して、各 $j$ に対し、$V_j \times_Y X = \emptyset$ であるか、
または射 $V_j \times_Y X \to X$ が $U_i$ を、ある $i$ に対して経由する。
\end{lemma}

\begin{proof}
各 $i$ に対して、シントミック（それぞれ滑らか、それぞれエタール）な射
$g_{ij} : V_{ij} \to Y$ と射 $V_{ij} \times_Y X \to U_i$（$X$ 上）を、
$\{V_{ij} \times_Y X \to U_i\}$ がザリスキー被覆となるように選べる。
補題 \ref{etale-cohomology-lemma-closed-immersion-cover-from-below} を参照されたい。
これは特に、$\bigcup_{ij} g_{ij}(V_{ij})$ が閉部分集合 $f(X)$ を含むことを意味する。
したがって、シントミック（それぞれ滑らか、それぞれエタール）な写像 $g_{ij}$ の族は、
開埋め込み $Y \setminus f(X) \to Y$ とあわせて、所望の $Y$ の
シントミック（それぞれ滑らか、それぞれエタール）な被覆をなす。
\end{proof}

\begin{lemma}
\label{etale-cohomology-lemma-closed-immersion-almost-cocontinuous}
$f : X \to Y$ をスキームの閉埋め込みとする。
$\tau \in \{syntomic, smooth, \etale\}$ とする。関手 $V \mapsto X \times_Y V$ は、
ほぼ余連続な関手 $(\Sch/Y)_\tau \to (\Sch/X)_\tau$ を、大 $\tau$ サイトの間に定める
（サイト、定義
\ref{sites-definition-almost-cocontinuous} を参照されたい）。
\end{lemma}

\begin{proof}
次を示さなければならない：射 $V \to Y$ と任意のシントミック
（それぞれ滑らか、それぞれエタール）な被覆
$\{U_i \to X \times_Y V\}$ が与えられたとき、滑らか（それぞれ滑らか、
それぞれエタール）な被覆 $\{V_j \to V\}$ が存在して、各 $j$ に対し、
$X \times_Y V_j$ が空であるか、または $X \times_Y V_j \to Z \times_Y V$ が
$U_i$ の一つを経由する。このことは、上の補題
\ref{etale-cohomology-lemma-prepare-closed-immersion-almost-cocontinuous} を閉埋め込み
$X \times_Y V \to V$ に適用すれば従う。
\end{proof}

\begin{lemma}
\label{etale-cohomology-lemma-closed-immersion-pushforward-exact}
$f : X \to Y$ をスキームの閉埋め込みとする。
$\tau \in \{syntomic, smooth, \etale\}$ とする。
\begin{enumerate}
\item 順像
$f_{big, *} :
\Sh((\Sch/X)_\tau)
\to
\Sh((\Sch/Y)_\tau)$
は余等化子および押し出しと可換である。
\item 順像
$f_{big, *} :
\textit{Ab}((\Sch/X)_\tau)
\to
\textit{Ab}((\Sch/Y)_\tau)$
は完全である。
\end{enumerate}
\end{lemma}

\begin{proof}
これはサイト、補題 \ref{sites-lemma-morphism-of-sites-almost-cocontinuous}、
サイト上の加群、補題
\ref{sites-modules-lemma-morphism-ringed-sites-almost-cocontinuous}、および上の補題
\ref{etale-cohomology-lemma-closed-immersion-almost-cocontinuous} から従う。
\end{proof}

\begin{remark}
\label{etale-cohomology-remark-fppf-closed-immersion-not-closed}
補題 \ref{etale-cohomology-lemma-closed-immersion-pushforward-exact} には $\tau = fppf$ の場合がない。
その理由は、環 $A$、イデアル $I$、および忠実平坦で有限表示な環準同型
$A/I \to \overline{B}$ が与えられても、平坦かつ有限表示な環準同型
$A \to B$ で、$B/IB \not = 0$ かつ $A/I \to B/IB$ が
$\overline{B}$ を経由するものを {\it 一つでも} 見つけられると考える理由はないからである。
したがって、補題 \ref{etale-cohomology-lemma-closed-immersion-almost-cocontinuous} の証明は
fppf 位相には使えない。実際、
$f_{big, *} : \textit{Ab}((\Sch/X)_{fppf})
\to \textit{Ab}((\Sch/Y)_{fppf})$
が $f$ が閉埋め込みのとき完全であるという主張は、おそらく偽である。
反例をご存じなら、
\href{mailto:stacks.project@gmail.com}{stacks.project@gmail.com}
まで電子メールで知らせていただきたい。
\end{remark}












\section{大サイトの下付き感嘆符関手の完全性}
\label{etale-cohomology-section-exactness-lower-shriek}

\noindent
これは単に次の技術的結果である。一般には、関手 $f_{big!}$ は
コンパクト台コホモロジーとはまったく関係がないことに注意する。

\begin{lemma}
\label{etale-cohomology-lemma-exactness-lower-shriek}
$\tau \in \{Zariski, \etale, smooth, syntomic, fppf\}$ とする。
$f : X \to Y$ をスキームの射とする。
$$
f_{big} :
\Sh((\Sch/X)_\tau)
\longrightarrow
\Sh((\Sch/Y)_\tau)
$$
を、位相、補題
\ref{topologies-lemma-morphism-big},
\ref{topologies-lemma-morphism-big-etale},
\ref{topologies-lemma-morphism-big-smooth},
\ref{topologies-lemma-morphism-big-syntomic}、または
\ref{topologies-lemma-morphism-big-fppf} における対応するトポスの射とする。
\begin{enumerate}
\item 関手
$f_{big}^{-1} : \textit{Ab}((\Sch/Y)_\tau) \to \textit{Ab}((\Sch/X)_\tau)$
は左随伴
$$
f_{big!} : \textit{Ab}((\Sch/X)_\tau) \to \textit{Ab}((\Sch/Y)_\tau)
$$
をもち、これは完全である。
\item 関手
$f_{big}^* :
\textit{Mod}((\Sch/Y)_\tau, \mathcal{O})
\to
\textit{Mod}((\Sch/X)_\tau, \mathcal{O})$
は左随伴
$$
f_{big!} :
\textit{Mod}((\Sch/X)_\tau, \mathcal{O})
\to
\textit{Mod}((\Sch/Y)_\tau, \mathcal{O})
$$
をもち、これは完全である。
\end{enumerate}
さらに、二つの関手 $f_{big!}$ は基礎にあるアーベル群の層上で一致する。
\end{lemma}

\begin{proof}
$f_{big}$ は、連続かつ余連続な関手
$u : (\Sch/X)_\tau \to (\Sch/Y)_\tau$、$U/X \mapsto U/Y$ に付随する
トポスの射であったことを思い出す。さらに、
$f_{big}^{-1}\mathcal{O} = \mathcal{O}$ である。したがって、$f_{big!}$ の存在は、
それぞれサイト上の加群、補題 \ref{sites-modules-lemma-g-shriek-adjoint}、および
サイト上の加群、補題 \ref{sites-modules-lemma-lower-shriek-modules} から従う。
$U$ が $(\Sch/X)_\tau$ の対象ならば、関手 $u$ は圏同値
$$
u' :
(\Sch/X)_\tau/U
\longrightarrow
(\Sch/Y)_\tau/U
$$
を誘導する。矢印の両側が $(\Sch/U)_\tau$ に等しいからである。
したがって、基礎にあるアーベル層上で $f_{big!}$ が一致することは、
サイト上の加群、注意 \ref{sites-modules-remark-when-shriek-equal} の議論から従う。
$f_{big!}$ の完全性はサイト上の加群、補題
\ref{sites-modules-lemma-exactness-lower-shriek} から従う。上の関手 $u$ は
ファイバー積および等化子と可換だからである。
\end{proof}

\noindent
次に、代数スタックに付随する異なるサイト上の加群の層を比較する際に
後で有用となる技術的補題を証明する。

\begin{lemma}
\label{etale-cohomology-lemma-compare-structure-sheaves}
$X$ をスキームとし、
$\tau \in \{Zariski, \etale, smooth, syntomic, fppf\}$ とする。
$\mathcal{C}_1 \subset \mathcal{C}_2 \subset (\Sch/X)_\tau$ を、次の性質をもつ
充満部分圏とする：
\begin{enumerate}
\item 対象 $U/X$ で $\mathcal{C}_t$ のものに対して、
\begin{enumerate}
\item $\{U_i \to U\}$ が $(\Sch/X)_\tau$ の被覆ならば、
$U_i/X$ は $\mathcal{C}_t$ の対象である。
\item $U \times \mathbf{A}^1/X$ は $\mathcal{C}_t$ の対象である。
\end{enumerate}
\item $X/X$ は $\mathcal{C}_t$ の対象である。
\end{enumerate}
$\mathcal{C}_t$ に、被覆がちょうど被覆 $\{U_i \to U\}$ で
$(\Sch/X)_\tau$ のものであり、$U \in \Ob(\mathcal{C}_t)$ を満たすものとなる
サイトの構造を入れる。
このとき
\begin{enumerate}
\item[(a)] 関手 $\mathcal{C}_1 \to \mathcal{C}_2$ は忠実充満、連続かつ余連続である。
\end{enumerate}
対応するトポスの射を $g : \Sh(\mathcal{C}_1) \to \Sh(\mathcal{C}_2)$ と書く。
$\mathcal{O}_t$ を $\mathcal{O}$ の $\mathcal{C}_t$ への制限と書く。
$g_!$ をサイト上の加群、定義 \ref{sites-modules-definition-g-shriek} の関手と書く。
\begin{enumerate}
\item[(b)] 標準射 $g_!\mathcal{O}_1 \to \mathcal{O}_2$ は同型である。
\end{enumerate}
\end{lemma}

\begin{proof}
主張 (a) は定義から直ちに従う。この証明では、すべてのスキームは $X$ 上の
スキームであり、スキームのすべての射は $X$ 上のスキームの射である。
$g^{-1}$ は制限で与えられるので、対象 $U$ で $\mathcal{C}_1$ のものに対して
$\mathcal{O}_1(U) = \mathcal{O}_2(U) = \mathcal{O}(U)$ となることに注意する。
$g_!\mathcal{O}_1$ は前層 $g_{p!}\mathcal{O}_1$ に付随する層であり、これは
$V$ で $\mathcal{C}_2$ に属するものに群
$$
\colim_{V \to U} \mathcal{O}(U)
$$
を対応させることを思い出す。ここで $U$ は $\mathcal{C}_1$ の対象を走り、
余極限はアーベル群の圏で取られる。以下では、
$$
V \to U \to U'
$$
が $U, U' \in \Ob(\mathcal{C}_1)$ を満たす射であり、
$f' \in \mathcal{O}(U')$ の制限が $f \in \mathcal{O}(U)$ ならば、
$(V \to U, f)$ と $(V \to U', f')$ は余極限の同じ元を定めることを頻繁に用いる。
また、$g_!\mathcal{O}_1 \to \mathcal{O}_2$ は元 $(V \to U, f)$ を単に
$f$ の $V$ への引き戻しに写す。

\medskip\noindent
全射性。$V$ をスキームとし、$h \in \mathcal{O}(V)$ とする。
射 $V \to X \times \mathbf{A}^1$ を得る。これは $h$ および構造射
$V \to X$ によって誘導される。
$\mathbf{A}^1 = \Spec(\mathbf{Z}[x])$ と書けば、元
$x \in \mathcal{O}(X \times \mathbf{A}^1)$ は $h$ に引き戻される。
仮定 (1)(b) と (2) により $X \times \mathbf{A}^1$ は $\mathcal{C}_1$ の対象なので、
所望の全射性を得る。

\medskip\noindent
単射性。$V$ をスキームとする。
$s = \sum_{i = 1, \ldots, n} (V \to U_i, f_i)$ を上に表示した余極限の元とする。
任意の $i$ に対して射 $f_i : U_i \to X \times \mathbf{A}^1$ を用いると、
$(V \to U_i, f_i)$ は $(f_i : V \to X \times \mathbf{A}^1, x)$ と
余極限の同じ元を定めることがわかる。そこで
$$
f_1 \times \ldots \times f_n : V \to X \times \mathbf{A}^n
$$
を考えると、$s$ は余極限において
$$
\sum\nolimits_{i = 1, \ldots, n}
(f_1 \times \ldots \times f_n : V \to X \times \mathbf{A}^n, x_i) =
(f_1 \times \ldots \times f_n : V \to X \times \mathbf{A}^n,
x_1 + \ldots + x_n)
$$
と同値であることがわかる。ここで $x_1 + \ldots + x_n$ の $V$ への制限が零ならば、
$f_1 \times \ldots \times f_n$ は
$X \times \mathbf{A}^{n - 1} = V(x_1 + \ldots + x_n)$ を経由することがわかる。
したがって、$s$ は余極限において零と同値である。
\end{proof}











%9.29.09
\section{エタールコホモロジー}
\label{etale-cohomology-section-etale-cohomology}

\noindent
以下の諸節では、エタールコホモロジーに関するいくつかの基本的結果を証明する。
位相空間のコホモロジーについて知られており、エタールコホモロジーについても
成り立つ事柄の一例を次に挙げる。

\begin{lemma}[エタールコホモロジーに対する Mayer--Vietoris]
\label{etale-cohomology-lemma-mayer-vietoris}
$X$ をスキームとする。$X = U \cup V$ は二つの開部分の合併であると仮定する。
任意のアーベル層 $\mathcal{F}$ で $X_\etale$ 上のものに対して、コホモロジー長完全列
$$
\begin{matrix}
0 \to
H^0_\etale(X, \mathcal{F}) \to
H^0_\etale(U, \mathcal{F}) \oplus H^0_\etale(V, \mathcal{F}) \to
H^0_\etale(U \cap V, \mathcal{F}) \phantom{\to \ldots} \\
\phantom{0} \to H^1_\etale(X, \mathcal{F}) \to
H^1_\etale(U, \mathcal{F}) \oplus H^1_\etale(V, \mathcal{F}) \to
H^1_\etale(U \cap V, \mathcal{F}) \to \ldots
\end{matrix}
$$
が存在する。この長完全列は $\mathcal{F}$ に関して関手的である。
\end{lemma}

\begin{proof}
$\mathcal{I}$ が単射的アーベル層ならば、
$$
0 \to \mathcal{I}(X) \to \mathcal{I}(U) \oplus \mathcal{I}(V) \to
\mathcal{I}(U \cap V) \to 0
$$
が完全であることに注意する。最初と中央の箇所については、$\mathcal{I}$ が層なので
成り立つ。右端については、サイト上のコホモロジー、補題
\ref{sites-cohomology-lemma-restriction-along-monomorphism-surjective}
により $\mathcal{I}(U) \to \mathcal{I}(U \cap V)$ が全射なので成り立つ。
別の証明方法として、写像 $\mathcal{I}(U) \oplus \mathcal{I}(V) \to
\mathcal{I}(U \cap V)$ の余核が、$\mathcal{I}$ の第一チェックコホモロジー群で、
被覆 $X = U \cup V$ に関するものであり、これが補題
\ref{etale-cohomology-lemma-hom-injective} および \ref{etale-cohomology-lemma-forget-injectives} により消えることを示してもよい。
したがって、$\mathcal{F} \to \mathcal{I}^\bullet$ が単射分解ならば、
$$
0 \to \mathcal{I}^\bullet(X) \to
\mathcal{I}^\bullet(U) \oplus \mathcal{I}^\bullet(V) \to
\mathcal{I}^\bullet(U \cap V) \to 0
$$
は複体の短完全列であり、それに付随するコホモロジー長完全列が補題の主張にある列である。
\end{proof}

\begin{lemma}[相対 Mayer--Vietoris]
\label{etale-cohomology-lemma-relative-mayer-vietoris}
$f : X \to Y$ をスキームの射とする。$X = U \cup V$ は二つの開部分スキームの
合併であると仮定する。
$a = f|_U : U \to Y$、$b = f|_V : V \to Y$、および
$c = f|_{U \cap V} : U \cap V \to Y$ と書く。
任意のアーベル層 $\mathcal{F}$ で $X_\etale$ 上のものに対して、長完全列
$$
0 \to
f_*\mathcal{F} \to
a_*(\mathcal{F}|_U) \oplus b_*(\mathcal{F}|_V) \to
c_*(\mathcal{F}|_{U \cap V}) \to
R^1f_*\mathcal{F} \to \ldots
$$
が $Y_\etale$ 上に存在する。この長完全列は $\mathcal{F}$ に関して関手的である。
\end{lemma}

\begin{proof}
$\mathcal{F} \to \mathcal{I}^\bullet$ を、$\mathcal{F}$ の、$X_\etale$ 上の
単射分解とする。複体の短完全列
$$
0 \to
f_*\mathcal{I}^\bullet \to
a_*\mathcal{I}^\bullet|_U \oplus b_*\mathcal{I}^\bullet|_V \to
c_*\mathcal{I}^\bullet|_{U \cap V} \to
0.
$$
が得られると主張する。実際、任意の $W$ で $Y_\etale$ に属するものと、任意の $n \geq 0$ に対して、
$W$ 上の切断の群の対応する列
$$
0 \to
\mathcal{I}^n(W \times_Y X) \to
\mathcal{I}^n(W \times_Y U)
\oplus \mathcal{I}^n(W \times_Y V) \to
\mathcal{I}^n(W \times_Y (U \cap V)) \to
0
$$
は補題 \ref{etale-cohomology-lemma-mayer-vietoris} の証明で短完全であることが示された。
コホモロジー層を取り、$\mathcal{I}^\bullet|_U$ が $\mathcal{F}|_U$ の単射分解であり、
$\mathcal{I}^\bullet|_V$、$\mathcal{I}^\bullet|_{U \cap V}$ についても同様であることを
用いれば、補題が従う。
\end{proof}







\section{余極限}
\label{etale-cohomology-section-colimit}

\noindent
$(\mathcal{F}_i, \varphi_{ii'})$ がサイト $\mathcal{C}$ 上の層の図式ならば、
その（層の圏における）余極限は前層 $U \mapsto \colim_i \mathcal{F}_i(U)$ の
層化であることを思い出す。サイト、補題 \ref{sites-lemma-colimit-sheaves} を参照されたい。
系が有向であり、$U$ が $\mathcal{C}$ の準コンパクト対象で、準コンパクト対象による
被覆の共終系をもつならば、$\mathcal{F}(U) = \colim \mathcal{F}_i(U)$ である。
サイト、補題 \ref{sites-lemma-directed-colimits-sections} を参照されたい。
アーベル層の余極限の高次コホモロジー群を扱う結果については、サイト上のコホモロジー、補題
\ref{sites-cohomology-lemma-colim-works-over-collection}
を参照されたい。

\medskip\noindent
サイト上のコホモロジー、補題 \ref{sites-cohomology-lemma-colimit} では、この結果を
サイトの逆系上の層の系へ一般化する。ここでは、スキームの逆系上にある
エタール層の系の場合に対応する概念を与える。

\begin{definition}
\label{etale-cohomology-definition-inverse-system-sheaves}
$I$ を前順序集合とする。$(X_i, f_{i'i})$ を $I$ 上のスキームの逆系とする。
{\it 層の系 $(\mathcal{F}_i, \varphi_{i'i})$ で $(X_i, f_{i'i})$ 上のもの}とは、次のデータである：
\begin{enumerate}
\item 層 $\mathcal{F}_i$ で $(X_i)_\etale$ 上にあり、各 $i \in I$ に対するもの、
\item $i' \geq i$ に対する写像
$\varphi_{i'i} : f_{i'i}^{-1}\mathcal{F}_i \to \mathcal{F}_{i'}$
で、$(X_{i'})_\etale$ 上の層の写像であるもの。
\end{enumerate}
ただし、$\varphi_{i''i} = \varphi_{i''i'} \circ f_{i'' i'}^{-1}\varphi_{i'i}$ とする。
これは $i'' \geq i' \geq i$ のときは常に課す。
\end{definition}

\noindent
定義 \ref{etale-cohomology-definition-inverse-system-sheaves} の状況で、$I$ は有向集合であり、
遷移射 $f_{i'i}$ はアフィンであると仮定する。
$X = \lim X_i$ をスキームの圏における極限とする。極限、節
\ref{limits-section-limits} を参照されたい。射影を $f_i : X \to X_i$ と書き、写像
$$
f_i^{-1}\mathcal{F}_i = f_{i'}^{-1}f_{i'i}^{-1}\mathcal{F}_i
\xrightarrow{f_{i'}^{-1}\varphi_{i'i}}
f_{i'}^{-1}\mathcal{F}_{i'}
$$
を考える。これにより $f_i^{-1}\mathcal{F}_i$ は $X_\etale$ 上の層の $I$ 上の系となる
（これを確かめるのはよい演習である）。しばしば、同型
$$
H^q_\etale(X, \colim f_i^{-1}\mathcal{F}_i) =
\colim H^q_\etale(X_i, \mathcal{F}_i)
$$
が存在するかを知りたい。$X_i$ が準コンパクトかつ準分離であることが各 $i$ に対して成り立てば、
これは真であることがわかる。定理 \ref{etale-cohomology-theorem-colimit} を参照されたい。

\begin{lemma}
\label{etale-cohomology-lemma-colimit-affine-sites}
$I$ を有向集合とする。$(X_i, f_{i'i})$ を $I$ 上のスキームの逆系で、
遷移射がアフィンであるものとする。$X = \lim_{i \in I} X_i$ とおく。
位相、補題 \ref{topologies-lemma-alternative} の記法により、
$$
X_{affine, \etale} = \colim (X_i)_{affine, \etale}
$$
がサイト、補題 \ref{sites-lemma-colimit-sites} の意味でサイトとして成り立つ。
\end{lemma}

\begin{proof}
まず、$X$ および $X_i$ が、すべての $i$ に対して準コンパクトかつ準分離である場合を
証明する（関心のあるすべての場合でこれは成り立つ）。この場合、
$X_{affine, \etale}$、それぞれ $(X_i)_{affine, \etale}$ の任意の対象は
$X$ 上有限表示である。さらに、$X$ 上有限表示なスキームの圏は、
$X_i$ 上有限表示なスキームの圏の余極限である。極限、補題
\ref{limits-lemma-descend-finite-presentation} を参照されたい。
$X$ 上エタールなアフィン対象の部分圏についても、極限、補題
\ref{limits-lemma-limit-affine} および \ref{limits-lemma-descend-etale} により同じことが成り立つ。
最後に、$\{U^j \to U\}$ が $X_{affine, \etale}$ の被覆であり、
$U_i^j \to U_i$ が $X_i$ 上エタールなアフィンスキームの射で、その $X$ への基底変換が
$U^j \to U$ ならば、$\{U^j_i \to U_i\}$ のある $X_{i'}$ への基底変換は、
十分大きい $i'$ に対して被覆となることがわかる。極限、補題
\ref{limits-lemma-descend-surjective} を参照されたい。

\medskip\noindent
一般の場合、$U$ を $X_{affine, \etale}$ の対象とする。このとき $U \to X$ は
エタールかつ分離的である（$U$ が分離的だからである）が、一般には準コンパクトでない。
それでも $U \to X$ は局所有限表示であるため、極限、補題
\ref{limits-lemma-descend-finite-presentation-variant} により、ある $i$、準コンパクトかつ
準分離なスキーム $U_i$、および局所有限表示な射 $U_i \to X_i$ が存在して、
その $X$ への基底変換は $U \to X$ となる。このとき
$U = \lim_{i' \geq i} U_{i'}$ であり、ここで
$U_{i'} = U_i \times_{X_i} X_{i'}$ とする。
$i$ を大きくすれば、$U_i$ はアフィンであると仮定してよい。極限、補題
\ref{limits-lemma-limit-affine} を参照されたい。
$U_i \to X_i$ が十分大きい $i$ に対してエタールであることを確かめるため、有限アフィン開被覆
$U_i = U_{i, 1} \cup \ldots \cup U_{i, m}$ であって、
$U_{i, j} \to U_i \to X_i$ がアフィン開部分 $W_{i, j} \subset X_i$ へ写るものを選ぶ。
すると極限、補題 \ref{limits-lemma-descend-etale} を適用でき、
$U_{i, j} \to W_{i, j}$ がエタールであることがわかる。ここで必要なら $i$ を大きくする。
このようにして関手
$\colim (X_i)_{affine, \etale} \to X_{affine, \etale}$ が本質的全射であることがわかる。
忠実充満性は、すでに用いた極限、補題
\ref{limits-lemma-descend-finite-presentation-variant} から直接従う。
被覆に関する主張は、証明の第一段落とまったく同じ方法で証明される。
\end{proof}

\noindent
以上を用いると、余極限とコホモロジーに関する次の一般的結果を得る。

\begin{theorem}
\label{etale-cohomology-theorem-colimit}
$X = \lim_{i \in I} X_i$ を、アフィンな遷移射
$f_{i'i} : X_{i'} \to X_i$ をもつスキームの有向系の極限とする。
$X_i$ はすべての $i \in I$ に対して準コンパクトかつ準分離であると仮定する。
$(\mathcal{F}_i, \varphi_{i'i})$ を $(X_i, f_{i'i})$ 上のアーベル層の系とする。
射影を $f_i : X \to X_i$ と書き、$\mathcal{F} = \colim f_i^{-1}\mathcal{F}_i$ とおく。
このとき
$$
\colim_{i\in I} H_\etale^p(X_i, \mathcal{F}_i) = H_\etale^p(X, \mathcal{F}).
$$
これはすべての $p \geq 0$ に対して成り立つ。
\end{theorem}

\begin{proof}
位相、補題 \ref{topologies-lemma-alternative} により、$\mathcal{F}$ のコホモロジーは
$X_{affine, \etale}$ 上で計算できる。したがって、結果は補題
\ref{etale-cohomology-lemma-colimit-affine-sites} とサイト上のコホモロジー、補題
\ref{sites-cohomology-lemma-colimit} をあわせれば従う。
\end{proof}

\noindent
次の二つの結果は上の定理の特別な場合である。

\begin{lemma}
\label{etale-cohomology-lemma-colimit}
$X$ を準コンパクトかつ準分離なスキームとする。$I$ を有向集合とする。
$(\mathcal{F}_i, \varphi_{ij})$ を $X_\etale$ 上の、$I$ 上のアーベル層の系とする。
このとき
$$
\colim_{i\in I} H_\etale^p(X, \mathcal{F}_i) = H_\etale^p(X,
\colim_{i\in I} \mathcal{F}_i).
$$
\end{lemma}

\begin{proof}
これは定理 \ref{etale-cohomology-theorem-colimit} の特別な場合である。直接の証明も概説する。
すべての $X$ について同時に、$p$ に関する帰納法で証明する。
\begin{enumerate}
\item
$X$ を任意の準コンパクトかつ準分離なスキーム、$\mathcal{U}$ を $X$ の任意の
エタール被覆とする。細分 $\mathcal{V} = \{V_j \to X\}_{j\in J}$ で、$J$ が有限であり、
各 $V_j$ が準コンパクトかつ準分離であり、すべての
$V_{j_0} \times_X \ldots \times_X V_{j_p}$ も準コンパクトかつ準分離となるものが
存在することを示す。
\item
前段階および層の圏における余極限の定義を用いて、定理が $p = 0$ およびすべての
$X$ に対して成り立つことを示す。
\item
コホモロジーの局所性（補題 \ref{etale-cohomology-lemma-locality-cohomology}）、
チェックからコホモロジーへのスペクトル系列（定理 \ref{etale-cohomology-theorem-cech-ss}）、および
上の状況におけるすべての $V_{j_0} \times_X \ldots \times_X V_{j_p}$ に
帰納法の仮定を適用できることを用いて、帰納段階 $p \to p + 1$ を証明する。
\end{enumerate}
\end{proof}

\begin{lemma}
\label{etale-cohomology-lemma-directed-colimit-cohomology}
$A$ を環、$(I, \leq)$ を有向集合、$(B_i, \varphi_{ij})$ を $A$-代数の系とする。
$B = \colim_{i\in I} B_i$ とおく。$X \to \Spec(A)$ を準コンパクトかつ準分離な
スキームの射とする。$\mathcal{F}$ を $X_\etale$ 上のアーベル層とする。
$Y_i = X \times_{\Spec(A)} \Spec(B_i)$、
$Y = X \times_{\Spec(A)} \Spec(B)$、
$\mathcal{G}_i = (Y_i \to X)^{-1}\mathcal{F}$、および
$\mathcal{G} = (Y \to X)^{-1}\mathcal{F}$ と書く。このとき
$$
H_\etale^p(Y, \mathcal{G}) =
\colim_{i\in I} H_\etale^p (Y_i, \mathcal{G}_i).
$$
\end{lemma}

\begin{proof}
これは定理 \ref{etale-cohomology-theorem-colimit} の特別な場合である。次のような直接の証明も概説する。
\begin{enumerate}
\item $V \to Y$ がエタールで、$V$ が準コンパクトかつ準分離であるとする。
このとき、$i\in I$ と $V_i \to Y_i$ が存在して、
$V = V_i \times_{Y_i} Y$ となる。考えているすべてのスキームがアフィンならば、
これは次の代数的主張に対応する：$B = \colim B_i$ で $B \to C$ がエタールならば、
$i \in I$ とエタールな $B_i \to C_i$ が存在して、
$C \cong B \otimes_{B_i} C_i$ となる。これは代数、補題
\ref{algebra-lemma-etale} で証明されている。
\item (1) の状況で、
$\mathcal{G}(V) = \colim_{i' \geq i} \mathcal{G}_{i'}(V_{i'})$
を示す。ここで $V_{i'}$ は $V_i$ の $Y_{i'}$ への基底変換である。
\item (1) により、任意のエタール被覆
$\mathcal{V} = \{V_j \to Y\}_{j\in J}$ で、$J$ が有限で $V_j$ が準コンパクトかつ準分離であるものに対し、$i \in I$ およびエタール被覆
$\mathcal{V}_i = \{V_{ij} \to Y_i\}_{j \in J}$ が存在して、
$\mathcal{V} \cong \mathcal{V}_i \times_{Y_i} Y$ となることがわかる。
\item (2) と (3) から
$$
\check H^*(\mathcal{V}, \mathcal{G})=
\colim_{i\in I} \check H^*(\mathcal{V}_i, \mathcal{G}_i).
$$
が従うことを示す。
\item チェックからコホモロジーへのスペクトル系列
（定理 \ref{etale-cohomology-theorem-cech-ss}）を巧妙に用いる。
\end{enumerate}
\end{proof}

\begin{lemma}
\label{etale-cohomology-lemma-higher-direct-images}
$f: X\to Y$ をスキームの射とし、$\mathcal{F}\in
\textit{Ab}(X_\etale)$ とする。このとき $R^pf_*\mathcal{F}$ は前層
$$
(V \to Y) \longmapsto H_\etale^p(X \times_Y V, \mathcal{F}|_{X \times_Y V}).
$$
に付随する層である。より一般に、$K \in D(X_\etale)$ に対して、$R^pf_*K$ は前層
$$
(V \to Y) \longmapsto H_\etale^p(X \times_Y V, K|_{X \times_Y V}).
$$
に付随する層である。
\end{lemma}

\begin{proof}
この補題は位相空間についても有効であり、その場合の証明も同じである。
層の場合についてはサイト上のコホモロジー、補題
\ref{sites-cohomology-lemma-higher-direct-images}
を、アーベル層の複体の場合についてはサイト上のコホモロジー、補題
\ref{sites-cohomology-lemma-sheafification-cohomology}
を参照されたい。
\end{proof}

\begin{lemma}
\label{etale-cohomology-lemma-relative-colimit}
$S$ をスキームとする。$X = \lim_{i \in I} X_i$ を $S$ 上のスキームの有向系の極限で、
アフィンな遷移射 $f_{i'i} : X_{i'} \to X_i$ をもつものとする。構造射
$g_i : X_i \to S$ および $g : X \to S$ は準コンパクトかつ準分離であると仮定する。
$(\mathcal{F}_i, \varphi_{i'i})$ を $(X_i, f_{i'i})$ 上のアーベル層の系とする。
射影を $f_i : X \to X_i$ と書き、$\mathcal{F} = \colim f_i^{-1}\mathcal{F}_i$ とおく。
このとき
$$
\colim_{i\in I} R^p g_{i, *} \mathcal{F}_i = R^p g_* \mathcal{F}
$$
がすべての $p \geq 0$ に対して成り立つ。
\end{lemma}

\begin{proof}
補題 \ref{etale-cohomology-lemma-higher-direct-images} により、$R^p g_{i, *} \mathcal{F}_i$ は前層
$U \mapsto H^p_\etale(U \times_S X_i, \mathcal{F}_i)$ に付随する層であり、
$R^pg_*\mathcal{F}$ についても同様であることを思い出す。さらに、層の系の余極限は
前層の段階での余極限の層化である。$S_\etale$ の各対象は、準コンパクトかつ
準分離な対象（たとえばアフィンスキーム）による被覆をもつことに注意する。
さらに、$U$ が準コンパクトかつ準分離な対象ならば、
$$
\colim H^p_\etale(U \times_S X_i, \mathcal{F}_i) =
H^p_\etale(U \times_S X, \mathcal{F})
$$
が定理 \ref{etale-cohomology-theorem-colimit} により成り立つ。したがって補題が従う。
\end{proof}

\begin{lemma}
\label{etale-cohomology-lemma-relative-colimit-general}
$I$ を有向集合とする。$g_i : X_i \to S_i$ を $I$ 上のスキームの射の逆系とする。
$g_i$ は準コンパクトかつ準分離であり、$i' \geq i$ に対する遷移射
$f_{i'i} : X_{i'} \to X_i$ および $h_{i'i} : S_{i'} \to S_i$ はアフィンであると仮定する。
$g : X \to S$ を射 $g_i$ の極限とする。極限、節 \ref{limits-section-limits} を参照されたい。
射影を $f_i : X \to X_i$ および $h_i : S \to S_i$ と書く。
$(\mathcal{F}_i, \varphi_{i'i})$ を $(X_i, f_{i'i})$ 上の層の系とする。
$\mathcal{F} = \colim f_i^{-1}\mathcal{F}_i$ とおく。このとき
$$
R^p g_* \mathcal{F} =
\colim_{i \in I} h_i^{-1}R^p g_{i, *} \mathcal{F}_i
$$
がすべての $p \geq 0$ に対して成り立つ。
\end{lemma}

\begin{proof}
補題の写像はどのように構成されるのだろうか。
$i' \geq i$ に対して、可換図式
$$
\xymatrix{
X \ar[r]_{f_{i'}} \ar[d]_g &
X_{i'} \ar[r]_{f_{i'i}} \ar[d]_{g_{i'}} &
X_i \ar[d]^{g_i} \\
S \ar[r]^{h_{i'}} &
S_{i'} \ar[r]^{h_{i'i}} &
S_i
}
$$
をもつ。基底変換写像
$h_{i'i}^{-1}Rg_{i, *}\mathcal{F}_i \to Rg_{i', *}f_{i'i}^{-1}\mathcal{F}_i$
（サイト上のコホモロジー、補題
\ref{sites-cohomology-lemma-base-change-map-flat-case} または
注意 \ref{sites-cohomology-remark-base-change}）を写像
$Rg_{i', *}\varphi_{i'i}$ とあわせると、
$\psi_{i'i} : h_{i' i}^{-1} R^p g_{i, *} \mathcal{F}_i \to
R^pg_{i', *} \mathcal{F}_{i'}$ を得る。同様に、図式の左の正方形を用いて写像
$\psi_i : h_i^{-1}R^pg_{i, *}\mathcal{F}_i \to R^pg_*\mathcal{F}$.
を得る。写像 $h_{i'}^{-1}\psi_{i'i}$ および $\psi_i$ が補題の主張で用いられる写像である。
これが意味をなすには、
$\psi_{i''i} = \psi_{i''i'} \circ h_{i''i'}^{-1}\psi_{i'i}$ および
$\psi_{i'} \circ h_{i'}^{-1}\psi_{i'i} = \psi_i$ を確かめなければならない。これは
サイト上のコホモロジー、注意
\ref{sites-cohomology-remark-compose-base-change-horizontal} から従う。

\medskip\noindent
等式の証明。まず、次元移動を用いる証明を与える\footnote{この方法を用いて
補題の写像を構成することもできる。}。任意の $U$ で $X$ 上アフィンかつエタールなものに対して、
定理 \ref{etale-cohomology-theorem-colimit} により
$$
g_*\mathcal{F}(U) =
H^0(U \times_S X, \mathcal{F}) =
\colim H^0(U_i \times_{S_i} X_i, \mathcal{F}_i) =
\colim g_{i, *}\mathcal{F}_i(U_i)
$$
となる。ここで余極限は、十分大きい $i$ 上で取る。すなわち、ある $i$ と
$U_i$ が存在し、これは $S_i$ 上アフィンかつエタールで、その基底変換が
$U$ であり、これは $S$ 上にあるものとする（補題 \ref{etale-cohomology-lemma-colimit-affine-sites} を参照されたい）。
右辺はサイト、補題 \ref{sites-lemma-colimit} により
$(\colim h_i^{-1}g_{i, *}\mathcal{F}_i)(U)$ に等しい。
これで $p = 0$ に対して補題が証明された。
$(\mathcal{G}_i, \varphi_{i'i})$ が
$\mathcal{G} = \colim f_i^{-1}\mathcal{G}_i$
を満たす系で、$\mathcal{G}_i$ が $X_i$ 上の単射的アーベル層であることが
すべての $i$ に対して成り立つならば、任意の $U$ で $X$ 上アフィンかつエタールなものに対して、
定理 \ref{etale-cohomology-theorem-colimit} により
$$
H^p(U \times_S X, \mathcal{G}) =
\colim H^p(U_i \times_{S_i} X_i, \mathcal{G}_i) = 0
$$
が $p > 0$ に対して成り立つ（余極限は先ほどと同じである）。したがって
$R^pg_*\mathcal{G} = 0$ であり、このような系に対して $p > 0$ の結果を得る。
一般には、系の短完全列
$$
0 \to (\mathcal{F}_i, \varphi_{i'i}) \to
(\mathcal{G}_i, \varphi_{i'i}) \to
(\mathcal{Q}_i, \varphi_{i'i}) \to 0
$$
で、$(\mathcal{G}_i, \varphi_{i'i})$ が上のとおりであるものを選べる。
サイト上のコホモロジー、補題
\ref{sites-cohomology-lemma-colim-sites-injective} を参照されたい。
帰納法により補題は $p - 1$ に対して成り立ち、上の議論により $p$ と
$(\mathcal{G}_i, \varphi_{i'i})$ に対する消滅が成り立つ。
したがって、コホモロジー長完全列により、$p$ と
$(\mathcal{F}_i, \varphi_{i'i})$ に対する結果が従う。

\medskip\noindent
第二の証明。
$S_{affine, \etale} = \colim (S_i)_{affine, \etale}$ であったことを思い出す。
補題 \ref{etale-cohomology-lemma-colimit-affine-sites} を参照されたい。したがって、$U$ が
$S_{affine, \etale}$ の対象ならば、$U = U_i \times_{S_i} S$ と、ある $i$ および
ある $U_i$ で $(S_i)_{affine, \etale}$ に属するものに対して書け、
$$
(\colim_{i \in I} h_i^{-1}R^p g_{i, *} \mathcal{F}_i)(U) =
\colim_{i' \geq i} (R^p g_{i', *}\mathcal{F}_{i'})(U_i \times_{S_i} S_{i'})
$$
となる。これはサイト、補題 \ref{sites-lemma-colimit} と、上で述べた系の遷移写像の
構成による。$R^pg_{i', *}\mathcal{F}_{i'}$ は前層
$U_{i'} \mapsto H^p(U_{i'} \times_{S_{i'}} X_{i'}, \mathcal{F}_{i'})$
に付随する層であり、$R^pg_*\mathcal{F}$ は前層
$U \mapsto H^p(U \times_S X, \mathcal{F})$
に付随する層なので（補題 \ref{etale-cohomology-lemma-higher-direct-images}）、標準可換図式
$$
\xymatrix{
\colim_{i' \geq i}
H^p(U_i \times_{S_i} X_{i'}, \mathcal{F}_{i'}) \ar[r] \ar[d] &
\colim_{i' \geq i}
(R^p g_{i', *}\mathcal{F}_{i'})(U_i \times_{S_i} S_{i'}) \ar[d] \\
H^p(U \times_S X, \mathcal{F}) \ar[r] &
R^pg_*\mathcal{F}(U)
}
$$
を得る。左の垂直矢印は定理 \ref{etale-cohomology-theorem-colimit} により同型であることに注意する。
示したいのは右の垂直矢印が同型であることである。しかし、この矢印の始域と終域が
$S_{affine, \etale}$ 上の層であることはすでにわかっている。したがって、次を示せば十分である：
(1) 終域の元は局所的に始域の元から来る；(2) 始域の元で終域の零に写るものは
局所的に消える。(1) は、上の議論と、下の水平矢印が層化後に同型となる前層の写像から
来るという事実から直ちに従う。
(2) について、$\xi \in  \colim_{i' \geq i}
(R^p g_{i', *}\mathcal{F}_{i'})(U_i \times_{S_i} S_{i'})$
が核に属するとする。$i' \geq i$ と
$\xi_{i'} \in (R^p g_{i', *}\mathcal{F}_{i'})(U_i \times_{S_i} S_{i'})$
で $\xi$ を表すものを選ぶ。標準エタール被覆
$\{U_{i', k} \to U_i \times_{S_i} S_{i'}\}_{k = 1, \ldots, m}$
で、$\xi_{i'}|_{U_{i', k}}$ が
$\xi_{i', k} \in H^p(U_{i', k} \times_{S_{i'}} X_{i'}, \mathcal{F}_{i'})$.
から来るものを選ぶ。$\xi$ が局所的に消えることを証明すれば十分なので、
$U$ をエタール被覆 $\{U_{i', k} \times_{S_{i'}} S \to U = U_i \times_{S_i} S\}$ の
各元で置き換えてよい。この置換後、$\xi$ は群の元 $\xi'$ の像であることがわかる：
$\colim_{i' \geq i} H^p(U_i \times_{S_i} X_{i'}, \mathcal{F}_{i'})$
これは上の図式における像である。$\xi'$ は $R^pg_*\mathcal{F}(U)$ の零に写るので、
もう一度置換を行い、$\xi'$ が $H^p(U \times_S X, \mathcal{F})$ の零に写ると
仮定してよい。しかし左の垂直矢印は同型なので、$\xi' = 0$、したがって所望のとおり
$\xi = 0$ と結論する。
\end{proof}

\begin{lemma}
\label{etale-cohomology-lemma-linus-hamann}
$X = \lim_{i \in I} X_i$ を、アフィンな遷移射 $f_{i'i}$ と射影
$f_i : X \to X_i$ をもつスキームの有向極限とする。
$\mathcal{F}$ を $X_\etale$ 上の層とする。このとき
\begin{enumerate}
\item 標準写像
$\varphi_{i'i} : f_{i'i}^{-1}f_{i, *}\mathcal{F} \to f_{i', *}\mathcal{F}$
が存在して、$(f_{i, *}\mathcal{F}, \varphi_{i'i})$ は定義
\ref{etale-cohomology-definition-inverse-system-sheaves} における $(X_i, f_{i'i})$ 上の層の系となる。
\item $\mathcal{F} = \colim f_i^{-1}f_{i, *}\mathcal{F}$ である。
\end{enumerate}
\end{lemma}

\begin{proof}
位相、補題 \ref{topologies-lemma-alternative} および補題
\ref{etale-cohomology-lemma-colimit-affine-sites} を介すれば、これはサイト、補題
\ref{sites-lemma-colimit-push-pull} の特別な場合である。
\end{proof}

\begin{lemma}
\label{etale-cohomology-lemma-compute-strangely}
$I$ を有向集合とする。$g_i : X_i \to S_i$ を $I$ 上のスキームの射の逆系とする。
$g_i$ は準コンパクトかつ準分離であり、$i' \geq i$ に対する遷移射
$X_{i'} \to X_i$ および $S_{i'} \to S_i$ はアフィンであると仮定する。
$g : X \to S$ を射 $g_i$ の極限とする。極限、節 \ref{limits-section-limits} を参照されたい。
射影を $f_i : X \to X_i$ および $h_i : S \to S_i$ と書く。
$\mathcal{F}$ を $X$ 上のアーベル層とする。このとき
$$
R^pg_*\mathcal{F} = \colim_{i \in I} h_i^{-1}R^pg_{i, *}(f_{i, *}\mathcal{F})
$$
となる。
\end{lemma}

\begin{proof}
補題 \ref{etale-cohomology-lemma-relative-colimit-general} と \ref{etale-cohomology-lemma-linus-hamann} を
形式的に組み合わせればよい。
\end{proof}











\section{余極限と複体}
\label{etale-cohomology-section-colimit-complexes}

\noindent
この節では、節 \ref{etale-cohomology-section-colimit} で層について始めた議論を続け、
さまざまな状況で複体の系のコホモロジーを取ることを論じる。
絶対に必要でない限り、この節を読まないよう読者に強く勧める。

\begin{lemma}
\label{etale-cohomology-lemma-colimit-variant-complexes}
$X = \lim_{i \in I} X_i$ を、アフィンな遷移射
$f_{i'i} : X_{i'} \to X_i$ をもつスキームの有向系の極限とする。
$X_i$ はすべての $i \in I$ に対して準コンパクトかつ準分離であると仮定する。
$\mathcal{F}_i^\bullet$ を $X_{i, \etale}$ 上のアーベル層の複体とする。
$\varphi_{i'i} : f_{i'i}^{-1}\mathcal{F}_i^\bullet \to
\mathcal{F}_{i'}^\bullet$ を $X_{i, \etale}$ 上の複体の写像とし、
$\varphi_{i''i} = \varphi_{i''i'} \circ f_{i'' i'}^{-1}\varphi_{i'i}$
が $i'' \geq i' \geq i$ のときは常に成り立つとする。
ある整数 $a$ が存在して、$\mathcal{F}_i^n = 0$ が
$n < a$ およびすべての $i \in I$ に対して成り立つと仮定する。このとき
$$
H^p_\etale(X, \colim f_i^{-1}\mathcal{F}_i^\bullet) =
\colim H^p_\etale(X_i, \mathcal{F}^\bullet_i)
$$
である。ここで $f_i : X \to X_i$ は射影である。
\end{lemma}

\begin{proof}
これは定理 \ref{etale-cohomology-theorem-colimit} の帰結である。
$\mathcal{F}^\bullet = \colim f_i^{-1}\mathcal{F}_i^\bullet$ とおく。
この定理より
$$
\colim_{i \in I} H_\etale^p(X_i, \mathcal{F}_i^n) =
H_\etale^p(X, \mathcal{F}^n)
$$
がすべての $n, p \in \mathbf{Z}$ に対して成り立つ。
導来圏、補題 \ref{derived-lemma-two-ss-complex-functor} のスペクトル系列
$$
E_{1, i}^{s, t} = H_\etale^t(X_i, \mathcal{F}_i^s) \Rightarrow
H_\etale^{s + t}(X_i, \mathcal{F}_i^\bullet)
$$
および
$$
E_1^{s, t} = H_\etale^t(X, \mathcal{F}^s) \Rightarrow
H_\etale^{s + t}(X, \mathcal{F}^\bullet)
$$
を用いる。$\mathcal{F}_i^n = 0$ が $n < a$ に対して成り立ち
（$a$ は $i$ に依存しない）、$E_{1, i}^{s, t}$（$i$ に依存しない）
および $E_1^{s, t}$ のうち、$H^q_\etale(X_i, \mathcal{F}_i^\bullet)$ と
$H^q_\etale(X, \mathcal{F}^\bullet)$ に寄与する項は、
固定された有限個だけである。また
$E_1^{s, t} = \colim E_{i, i}^{s, t}$ である。
これから所望の主張が従う。細部は省略する。
（スペクトル系列を用いず、複体の「愚直な」切断を用いる別の論証もある。）
\end{proof}

\begin{lemma}
\label{etale-cohomology-lemma-direct-sum-bounded-below-cohomology}
$X$ を準コンパクトかつ準分離なスキームとする。
$K_i \in D(X_\etale)$、$i \in I$ を対象の族とする。
$a \in \mathbf{Z}$ が与えられ、$H^n(K_i) = 0$ が
$n < a$ と $i \in I$ に対して成り立つと仮定する。このとき $R\Gamma(X, \bigoplus_i K_i) =
\bigoplus_i R\Gamma(X, K_i)$.
\end{lemma}

\begin{proof}
$H^p(X, \bigoplus_i K_i) =
\bigoplus_i H^p(X, K_i)$ がすべての $p \in \mathbf{Z}$ に対して
成り立つことを示せばよい。複体 $\mathcal{F}_i^\bullet$ で $K_i$ を表し、
$\mathcal{F}_i^n = 0$ が $n < a$ に対して成り立つものを選ぶ。
単射的対象、補題 \ref{injectives-lemma-derived-products} により、
複体 $\mathcal{F}_i^\bullet$ の直和は対象 $\bigoplus K_i$ を表す。
$\bigoplus \mathcal{F}^\bullet$ は有限直和のフィルター余極限なので、
結果は補題 \ref{etale-cohomology-lemma-colimit-variant-complexes} から従う。
\end{proof}

\begin{lemma}
\label{etale-cohomology-lemma-colimit-complexes}
$S$ をスキームとする。$X = \lim_{i \in I} X_i$ を、$S$ 上のスキームの
有向系で、アフィンな遷移射 $f_{i'i} : X_{i'} \to X_i$ をもつものの極限とする。
$X_i$ はすべての $i \in I$ に対して準コンパクトかつ準分離であると仮定する。
$K \in D^+(S_\etale)$ とする。このとき
$$
\colim_{i \in I} H_\etale^p(X_i, K|_{X_i}) = H_\etale^p(X, K|_X).
$$
この等式はすべての $p \in \mathbf{Z}$ に対して成り立つ。ここで
$K|_{X_i}$ および $K|_X$ は $K$ の $X_i$ および $X$ への引き戻しである。
\end{lemma}

\begin{proof}
$K$ は下に有界な複体 $\mathcal{G}^\bullet$ で、$S_\etale$ 上の
アーベル層からなるものによって表してよい。$\mathcal{G}^n = 0$ が
$n < a$ に対して成り立つとする。$\mathcal{F}^\bullet_i$ と
$\mathcal{F}^\bullet$ で、この複体の $X_i$ と $X$ への引き戻しを表す。
これらの複体は対象 $K|_{X_i}$ と $K|_X$ を表し、項ごとに
$\mathcal{F}^\bullet = \colim f_i^{-1}\mathcal{F}_i^\bullet$
である。したがって補題は補題
\ref{etale-cohomology-lemma-colimit-variant-complexes} から従う。
\end{proof}

\begin{lemma}
\label{etale-cohomology-lemma-relative-colimit-general-complexes}
$I$、$g_i : X_i \to S_i$、$g : X \to S$、$f_i$、$g_i$、$h_i$ を
補題 \ref{etale-cohomology-lemma-relative-colimit-general} と同じものとする。
$0 \in I$ および $K_0 \in D^+(X_{0, \etale})$ とする。
$i \geq 0$ に対し、$K_i$ で $K_0$ の $X_i$ への引き戻しを表す。
$K$ で $K_0$ の $X$ への引き戻しを表す。このとき
$$
R^pg_*K = \colim_{i \geq 0} h_i^{-1}R^pg_{i, *}K_i
$$
がすべての $p \in \mathbf{Z}$ に対して成り立つ。
\end{lemma}

\begin{proof}
整数 $p_0 \in \mathbf{Z}$ を固定する。
整数 $a$ で、$H^j(K_0) = 0$ が $j < a$ に対して成り立つものをとる。
この公式がすべての $p \leq p_0$ に対して成り立つことを、
$a$ に関する下降帰納法により証明する。$a > p_0$ ならば、自明な消滅により、
公式の左辺と右辺は $p \leq p_0$ に対して零である。導来圏、補題
\ref{derived-lemma-negative-vanishing} を参照されたい。
$a \leq p_0$ と仮定する。識別三角
$$
H^a(K_0)[-a] \to K_0 \to \tau_{\geq a + 1}K_0
$$
を考える。この識別三角を $X_i$ と $X$ へ引き戻すと、
$K_i$ と $K$ に対する両立する識別三角を得る。
$p \leq p_0$ に対して、可換図式
$$
\xymatrix{
\colim_{i \geq 0} h_i^{-1}R^{p - 1}g_{i, *}(\tau_{\geq a + 1}K_i)
\ar[r]_-\alpha \ar[d] &
R^{p - 1}g_*(\tau_{\geq a + 1}K) \ar[d] \\
\colim_{i \geq 0} h_i^{-1}R^pg_{i, *}(H^a(K_i)[-a])
\ar[r]_-\beta \ar[d] &
R^pg_*(H^a(K)[-a]) \ar[d] \\
\colim_{i \geq 0} h_i^{-1}R^pg_{i, *}K_i
\ar[r]_-\gamma \ar[d] &
R^pg_*K \ar[d] \\
\colim_{i \geq 0} R^pg_{i, *}\tau_{\geq a + 1}K_i
\ar[r]_-\delta \ar[d] &
R^pg_*\tau_{\geq a + 1}K \ar[d] \\
\colim_{i \geq 0} R^{p + 1}g_{i, *}(H^a(K_i)[-a])
\ar[r]^-\epsilon &
R^{p + 1}g_*(H^a(K)[-a])
}
$$
を考える。その各列は完全である。矢印 $\beta$ と $\epsilon$ は補題
\ref{etale-cohomology-lemma-relative-colimit-general} により同型である。
矢印 $\alpha$ と $\delta$ は帰納法の仮定により同型である。
したがって、所望のとおり $\gamma$ は同型である。
\end{proof}

\begin{lemma}
\label{etale-cohomology-lemma-relative-colimit-general-really-complexes}
$I$、$g_i : X_i \to S_i$、$g : X \to S$、$f_{ii'}$、$f_i$、$g_i$、$h_i$
を補題 \ref{etale-cohomology-lemma-relative-colimit-general} と同じものとする。
$\mathcal{F}_i^\bullet$ を $X_{i, \etale}$ 上のアーベル層の複体とする。
$\varphi_{i'i} : f_{i'i}^{-1}\mathcal{F}_i^\bullet \to
\mathcal{F}_{i'}^\bullet$ を $X_{i, \etale}$ 上の複体の写像とし、
$\varphi_{i''i} = \varphi_{i''i'} \circ f_{i'' i'}^{-1}\varphi_{i'i}$
が $i'' \geq i' \geq i$ のときは常に成り立つとする。
ある整数 $a$ が存在して、$\mathcal{F}_i^n = 0$ が
$n < a$ およびすべての $i \in I$ に対して成り立つと仮定する。このとき
$$
R^pg_*(\colim f_i^{-1}\mathcal{F}_i^\bullet) =
\colim_{i \geq 0} h_i^{-1}R^pg_{i, *}\mathcal{F}_i^\bullet
$$
がすべての $p \in \mathbf{Z}$ に対して成り立つ。
\end{lemma}

\begin{proof}
これは補題 \ref{etale-cohomology-lemma-relative-colimit-general} の帰結である。
$\mathcal{F}^\bullet = \colim f_i^{-1}\mathcal{F}_i^\bullet$ とおく。
この補題より
$$
\colim_{i \in I} h_i^{-1}R^pg_{i, *}\mathcal{F}_i^n = R^pg_*\mathcal{F}^n
$$
がすべての $n, p \in \mathbf{Z}$ に対して成り立つ。
導来圏、補題 \ref{derived-lemma-two-ss-complex-functor} のスペクトル系列
$$
E_{1, i}^{s, t} = R^tg_{i, *}\mathcal{F}_i^s \Rightarrow
R^{s + t}g_{i, *}\mathcal{F}_i^\bullet
$$
および
$$
E_1^{s, t} = R^tg_*\mathcal{F}^s \Rightarrow
R^{s + t}g_*\mathcal{F}^\bullet
$$
を用いる。$\mathcal{F}_i^n = 0$ が $n < a$ に対して成り立ち
（$a$ は $i$ に依存しない）、寄与する
$E_{1, i}^{s, t}$（$i$ に依存しない）および $E_1^{s, t}$ の項は、
固定された有限個だけである。また
$E_1^{s, t} = \colim E_{i, i}^{s, t}$ である。
これから所望の主張が従う。細部は省略する。
（スペクトル系列を用いず、複体の「愚直な」切断を用いる別の論証もある。）
\end{proof}

\begin{lemma}
\label{etale-cohomology-lemma-direct-sum-bounded-below-pushforward}
$f : X \to Y$ をスキームの準コンパクトかつ準分離な射とする。
$K_i \in D(X_\etale)$、$i \in I$ を対象の族とする。
$a \in \mathbf{Z}$ が与えられ、$H^n(K_i) = 0$ が
$n < a$ と $i \in I$ に対して成り立つと仮定する。このとき
$Rf_*(\bigoplus_i K_i) = \bigoplus_i Rf_*K_i$.
\end{lemma}

\begin{proof}
$R^pf_*(\bigoplus_i K_i) =
\bigoplus_i R^pf_*K_i$ がすべての $p \in \mathbf{Z}$ に対して
成り立つことを示せばよい。複体 $\mathcal{F}_i^\bullet$ で $K_i$ を表し、
$\mathcal{F}_i^n = 0$ が $n < a$ に対して成り立つものを選ぶ。
単射的対象、補題 \ref{injectives-lemma-derived-products} により、
複体 $\mathcal{F}_i^\bullet$ の直和は対象 $\bigoplus K_i$ を表す。
$\bigoplus \mathcal{F}^\bullet$ は有限直和のフィルター余極限なので、
結果は補題 \ref{etale-cohomology-lemma-relative-colimit-general-really-complexes} から従う。
\end{proof}














\section{高次順像の茎}
\label{etale-cohomology-section-stalks-direct-image}

\noindent
高次順像の茎は、しばしば次のように計算できる。

\begin{theorem}
\label{etale-cohomology-theorem-higher-direct-images}
$f: X \to S$ をスキームの準コンパクトかつ準分離な射、
$\mathcal{F}$ を $X_\etale$ 上のアーベル層、$\overline{s}$ を
$S$ の幾何点で $s \in S$ 上にあるものとする。このとき
$$
\left(R^nf_* \mathcal{F}\right)_{\overline{s}} =
H_\etale^n( X \times_S \Spec(\mathcal{O}_{S, \overline{s}}^{sh}),
p^{-1}\mathcal{F})
$$
である。ここで $p : X \times_S \Spec(\mathcal{O}_{S, \overline{s}}^{sh}) \to X$
は射影である。$K \in D^+(X_\etale)$ および $n \in \mathbf{Z}$ に対して
$$
\left(R^nf_*K\right)_{\overline{s}} =
H_\etale^n(X \times_S \Spec(\mathcal{O}_{S, \overline{s}}^{sh}), p^{-1}K)
$$
である。実際、
$$
\left(Rf_*K\right)_{\overline{s}}
=
R\Gamma_\etale(X \times_S \Spec(\mathcal{O}_{S, \overline{s}}^{sh}), p^{-1}K)
$$
が $D^+(\textit{Ab})$ において成り立つ。
\end{theorem}

\begin{proof}
$\mathcal{I}$ を $\overline{s}$ の $S$ 上のエタール近傍の圏とする。
補題 \ref{etale-cohomology-lemma-higher-direct-images} により
$$
(R^nf_*\mathcal{F})_{\overline{s}} =
\colim_{(V, \overline{v}) \in \mathcal{I}^{opp}}
H_\etale^n(X \times_S V, \mathcal{F}|_{X \times_S V}).
$$
である。$\mathcal{I}$ を、$\overline{s}$ のアフィンなエタール近傍からなる
始部分圏で置き換えてよい。補題 \ref{etale-cohomology-lemma-describe-etale-local-ring} と
極限、補題
\ref{limits-lemma-directed-inverse-system-affine-schemes-has-limit} により
$$
\Spec(\mathcal{O}_{S, \overline{s}}^{sh}) =
\lim_{(V, \overline{v}) \in \mathcal{I}} V
$$
であることに注意する。ファイバー積は極限と可換なので、
$$
X \times_S \Spec(\mathcal{O}_{S, \overline{s}}^{sh}) =
\lim_{(V, \overline{v}) \in \mathcal{I}} X \times_S V
$$
も得る。補題 \ref{etale-cohomology-lemma-directed-colimit-cohomology} により結論が従う。
第二の変形には、補題 \ref{etale-cohomology-lemma-colimit-complexes} を
補題 \ref{etale-cohomology-lemma-directed-colimit-cohomology} の代わりに用いて同じ論証を行う。

\medskip\noindent
最後の主張を示すには、写像
$\left(Rf_*K\right)_{\overline{s}} \to
R\Gamma_\etale(X \times_S \Spec(\mathcal{O}_{S, \overline{s}}^{sh}), p^{-1}K)$
を $D^+(\textit{Ab})$ において構成し、それがすべての $n$ に対して
上の次数 $n$ のコホモロジー群の同型を実現することを示せば十分である。
そのため、下に有界な複体 $\mathcal{J}^\bullet$ で、$X_\etale$ 上の
単射的アーベル層からなり、$K$ を表すものを選ぶ。複体
$f_*\mathcal{J}^\bullet$ は $Rf_*K$ を表す。したがって複体
$$
(f_*\mathcal{J}^\bullet)_{\overline{s}} =
\colim_{(V, \overline{v}) \in \mathcal{I}^{opp}}
(f_*\mathcal{J}^\bullet)(V)
$$
は $(Rf_*K)_{\overline{s}}$ を表す。各 $V$ に対して写像
$$
(f_*\mathcal{J}^\bullet)(V) =
\Gamma(X \times_S V, \mathcal{J}^\bullet)
\longrightarrow
\Gamma(X \times_S \Spec(\mathcal{O}_{S, \overline{s}}^{sh}),
p^{-1}\mathcal{J}^\bullet)
$$
があり、その終域の複体は
$R\Gamma_\etale(X \times_S \Spec(\mathcal{O}_{S, \overline{s}}^{sh}), p^{-1}K)$
を $D^+(\textit{Ab})$ において表す。
これらの写像の余極限を取れば結果を得る。
\end{proof}

\begin{remark}
\label{etale-cohomology-remark-stalk-fibre}
$f : X \to S$ をスキームの射とする。
$K \in D(X_\etale)$ とする。
$\overline{s}$ を $S$ の幾何点とする。常に標準写像
$$
(Rf_*K)_{\overline{s}}
\longrightarrow
R\Gamma(X \times_S \Spec(\mathcal{O}_{S, \overline{s}}^{sh}), p^{-1}K)
\longrightarrow
R\Gamma(X_{\overline{s}}, K|_{X_{\overline{s}}})
$$
が存在する。ここで $p : X \times_S \Spec(\mathcal{O}_{S, \overline{s}}^{sh}) \to X$
は射影である。実際、可換図式
$$
\xymatrix{
X_{\overline{s}} \ar[r] \ar[d]^{f_{\overline{s}}} &
X \times_S \Spec(\mathcal{O}_{S, \overline{s}}^{sh}) \ar[r]_-p \ar[d]^{f'} &
X \ar[d]^f \\
\overline{s} \ar[r]^-i &
\Spec(\mathcal{O}_{S, \overline{s}}^{sh}) \ar[r]^-j &
S
}
$$
を考える。この図式の二つの正方形に対して、基底変換写像
$$
i^{-1}Rf'_*(p^{-1}K) \to Rf_{\overline{s}, *}(K|_{X_{\overline{s}}})
\quad\text{および}\quad
j^{-1}Rf_*K \to Rf'_*(p^{-1}K)
$$
がある（サイト上のコホモロジー、注意
\ref{sites-cohomology-remark-base-change}）。
大域切断を取ると所望の写像を得る。サイト上のコホモロジー、注意
\ref{sites-cohomology-remark-compose-base-change-horizontal}
により、これら二つの写像の合成は通常の（基底変換）写像
$(Rf_*K)_{\overline{s}} \to R\Gamma(X_{\overline{s}}, K|_{X_{\overline{s}}})$.
である。
\end{remark}






\section{ルレイ・スペクトル系列}
\label{etale-cohomology-section-leray}

\begin{lemma}
\label{etale-cohomology-lemma-prepare-leray}
$f: X \to Y$ を射とし、$\mathcal{I}$ を
$\textit{Ab}(X_\etale)$ の単射的対象とする。$V \in \Ob(Y_\etale)$ とする。このとき
\begin{enumerate}
\item 任意の被覆 $\mathcal{V} = \{V_j\to V\}_{j \in J}$ に対して
$\check H^p(\mathcal{V}, f_*\mathcal{I}) = 0$ がすべての $p > 0$ について成り立つ。
\item $f_*\mathcal{I}$ は関手 $\Gamma(V, -)$ に関して非輪状である。
\item $g : Y \to Z$ ならば、$f_*\mathcal{I}$ は $g_*$ に関して非輪状である。
\end{enumerate}
\end{lemma}

\begin{proof}
等式 $\check{\mathcal{C}}^\bullet(\mathcal{V}, f_*\mathcal{I}) =
\check{\mathcal{C}}^\bullet(\mathcal{V} \times_Y X, \mathcal{I})$
に注意する。この複体の高次コホモロジー群は補題
\ref{etale-cohomology-lemma-hom-injective} により消滅する。これで (1) が示された。
任意の被覆に関する高次チェック・コホモロジー群が消滅する層は
高次コホモロジー群も消滅するので、第二の主張が従う。
これはチェック・コホモロジーからコホモロジーへのスペクトル系列を使う
見事な演習であるが、詳細およびより精密で一般的な主張については、
サイト上のコホモロジー、補題
\ref{sites-cohomology-lemma-cech-vanish-collection} を参照されたい。
(3) は (2) と、補題 \ref{etale-cohomology-lemma-higher-direct-images} における
$R^pg_*$ の記述から従う。
\end{proof}

\noindent
グロタンディーク・スペクトル系列の形式論を用いると、次を得る。

\begin{proposition}[ルレイ・スペクトル系列]
\label{etale-cohomology-proposition-leray}
$f: X \to Y$ をスキームの射とし、$\mathcal{F}$ を $X$ 上のエタール層とする。
このときスペクトル系列
$$
E_2^{p, q} = H_\etale^p(Y, R^qf_*\mathcal{F}) \Rightarrow
H_\etale^{p+q}(X, \mathcal{F}).
$$
が存在する。
\end{proposition}

\begin{proof}
補題 \ref{etale-cohomology-lemma-prepare-leray} および導来圏、節
\ref{derived-section-composition-right-derived-functors} を参照されたい。
\end{proof}









\section{有限射の高次順像の消滅}
\label{etale-cohomology-section-vanishing-finite-morphism}

\noindent
次の目標は、スキームの有限射の高次順像が消滅することを証明することである。

\begin{lemma}
\label{etale-cohomology-lemma-vanishing-etale-cohomology-strictly-henselian}
$R$ を狭義ヘンゼル局所環とする。$S = \Spec(R)$ とおき、
$\overline{s}$ をその閉点とする。このとき大域切断関手
$\Gamma(S, -) : \textit{Ab}(S_\etale) \to \textit{Ab}$
は完全である。実際、$\Gamma(S, \mathcal{F}) = \mathcal{F}_{\overline{s}}$
が任意の集合の層 $\mathcal{F}$ に対して成り立つ。とくに
$$
\forall p\geq 1, \quad H_\etale^p(S, \mathcal{F})=0
$$
がすべての $\mathcal{F}\in \textit{Ab}(S_\etale)$ に対して成り立つ。
\end{lemma}

\begin{proof}
$\Gamma(S, \mathcal{F}) = \mathcal{F}_{\overline{s}}$ を示せば、
茎関手が完全なので $\Gamma(S, -)$ も完全となる。
$(U, \overline{u})$ を $\overline{s}$ のエタール近傍とする。
アフィン開近傍 $\Spec(A)$ で、$\overline{u}$ の $U$ におけるものを選ぶ。
このとき $R \to A$ はエタールであり、
$\kappa(\overline{s}) = \kappa(\overline{u})$ である。
定理 \ref{etale-cohomology-theorem-henselian} により、$A \cong R \times A'$ が、
$R$-代数として、$\kappa(\overline{s}) = \kappa(\overline{u})$
への写像と両立して成り立つことがわかる。
したがって切断
$$
\xymatrix{
\Spec(A) \ar[r] & U \ar[d]\\
& S \ar[ul]
}
$$
を得る。したがって、$\overline{s}$ のエタール近傍の系において恒等写像
$(S, \overline{s}) \to (S, \overline{s})$ は共終である。
ゆえに $\Gamma(S, \mathcal{F}) = \mathcal{F}_{\overline{s}}$ である。
完全関手の高次導来関手は零なので、補題の最後の主張が従う。
導来圏、補題 \ref{derived-lemma-right-derived-exact-functor} を参照されたい。
\end{proof}

\begin{proposition}
\label{etale-cohomology-proposition-finite-higher-direct-image-zero}
$f : X \to Y$ をスキームの有限射とする。
\begin{enumerate}
\item 任意の幾何点 $\overline{y} : \Spec(k) \to Y$ に対して
$$
(f_*\mathcal{F})_{\overline{y}} =
\prod\nolimits_{\overline{x} : \Spec(k) \to X,\ f(\overline{x}) =
\overline{y}} \mathcal{F}_{\overline{x}}.
$$
が $\mathcal{F}$ が $\Sh(X_\etale)$ の対象であるとき成り立ち、また
$$
(f_*\mathcal{F})_{\overline{y}} =
\bigoplus\nolimits_{\overline{x} : \Spec(k) \to X,\ f(\overline{x}) =
\overline{y}} \mathcal{F}_{\overline{x}}.
$$
が $\mathcal{F}$ が $\textit{Ab}(X_\etale)$ の対象であるとき成り立つ。
\item 任意の $q \geq 1$ に対して $R^q f_*\mathcal{F} = 0$ が、
$\mathcal{F}$ が $\textit{Ab}(X_\etale)$ の対象であるとき成り立つ。
\end{enumerate}
\end{proposition}

\begin{proof}
$X_{\overline{y}}^{sh}$ でファイバー積
$X \times_Y \Spec(\mathcal{O}_{Y, \overline{y}}^{sh})$ を表す。
定理 \ref{etale-cohomology-theorem-higher-direct-images} により、
$R^qf_*\mathcal{F}$ の $\overline{y}$ における茎は
$H_\etale^q(X_{\overline{y}}^{sh}, \mathcal{F})$ で計算される。
$f$ は有限なので、$X_{\bar y}^{sh}$ は
$\Spec(\mathcal{O}_{Y, \overline{y}}^{sh})$ 上有限である。したがって、
$X_{\bar y}^{sh} = \Spec(A)$ と、ある環 $A$ で
$\mathcal{O}_{Y, \bar y}^{sh}$ 上有限なものに対して書ける。
後者は狭義ヘンゼルなので、補題 \ref{etale-cohomology-lemma-finite-over-henselian} により、
$A$ はヘンゼル局所環の有限積
$A = A_1 \times \ldots \times A_r$ である。
$\mathcal{O}_{Y, \overline{y}}^{sh}$ の剰余体は分離閉であるから、
各 $A_i$ についても同じである。ゆえに $A_i$ は狭義ヘンゼルである。
したがって $X_{\overline{y}}^{sh} = \coprod_{i = 1}^r \Spec(A_i)$ である。
補題 \ref{etale-cohomology-lemma-vanishing-etale-cohomology-strictly-henselian} の消滅から
$(R^qf_*\mathcal{F})_{\overline{y}} = 0$ が $q > 0$ に対して従い、
定理 \ref{etale-cohomology-theorem-exactness-stalks} により (2) が従う。
(1) は補題 \ref{etale-cohomology-lemma-vanishing-etale-cohomology-strictly-henselian} の
対応する主張から従う。
\end{proof}

\begin{lemma}
\label{etale-cohomology-lemma-finite-pushforward-commutes-with-base-change}
カルテシアン正方形
$$
\xymatrix{
X' \ar[r]_{g'} \ar[d]_{f'} & X \ar[d]^f \\
Y' \ar[r]^g & Y
}
$$
でスキームからなり、$f$ が有限射であるものを考える。
任意の集合の層 $\mathcal{F}$ で $X_\etale$ 上にあるものに対して
$f'_*(g')^{-1}\mathcal{F} = g^{-1}f_*\mathcal{F}$.
である。
\end{lemma}

\begin{proof}
非常に一般的な状況で引き戻し写像
$g^{-1}f_*\mathcal{F} \to f'_*(g')^{-1}\mathcal{F}$ が存在する。
サイト、節 \ref{sites-section-pullback} を参照されたい。
茎で確かめれば十分である（定理 \ref{etale-cohomology-theorem-exactness-stalks}）。
$\overline{y}' : \Spec(k) \to Y'$ を幾何点とする。このとき
\begin{align*}
(f'_*(g')^{-1}\mathcal{F})_{\overline{y}'}
& =
\prod\nolimits_{\overline{x}' : \Spec(k) \to X',\ f' \circ \overline{x}' =
\overline{y}'}
((g')^{-1}\mathcal{F})_{\overline{x}'} \\
& =
\prod\nolimits_{\overline{x}' : \Spec(k) \to X',\ f' \circ \overline{x}' =
\overline{y}'} \mathcal{F}_{g' \circ \overline{x}'} \\
& =
\prod\nolimits_{\overline{x} : \Spec(k) \to X,\ f \circ \overline{x} =
g \circ \overline{y}'} \mathcal{F}_{\overline{x}} \\
& =
(f_*\mathcal{F})_{g \circ \overline{y}'} \\
& =
(g^{-1}f_*\mathcal{F})_{\overline{y}'}
\end{align*}
である。第一の等式は命題
\ref{etale-cohomology-proposition-finite-higher-direct-image-zero} から従う。
第二の等式は補題 \ref{etale-cohomology-lemma-stalk-pullback} から従う。
第三の等式が成り立つのは、図式がカルテシアン正方形であり、したがって写像
$$
\{\overline{x}' : \Spec(k) \to X',\ f' \circ \overline{x}' =
\overline{y}'\}
\longrightarrow
\{\overline{x} : \Spec(k) \to X,\ f \circ \overline{x} =
g \circ \overline{y}'\}
$$
で $\overline{x}'$ を $g' \circ \overline{x}'$ へ送るものが全単射だからである。
第四の等式は命題
\ref{etale-cohomology-proposition-finite-higher-direct-image-zero} から従う。
第五の等式は補題 \ref{etale-cohomology-lemma-stalk-pullback} から従う。
\end{proof}

\begin{lemma}
\label{etale-cohomology-lemma-integral-pushforward-commutes-with-base-change}
カルテシアン正方形
$$
\xymatrix{
X' \ar[r]_{g'} \ar[d]_{f'} & X \ar[d]^f \\
Y' \ar[r]^g & Y
}
$$
でスキームからなり、$f$ が整射であるものを考える。
任意の集合の層 $\mathcal{F}$ で $X_\etale$ 上にあるものに対して
$f'_*(g')^{-1}\mathcal{F} = g^{-1}f_*\mathcal{F}$.
である。
\end{lemma}

\begin{proof}
問題は $Y$ 上局所的なので、$Y$ はアフィンであると仮定してよい。
このとき $X = \lim X_i$ と書け、$f_i : X_i \to Y$ は有限である
（アフィンの場合にはこれは容易であるが、参照として極限、補題
\ref{limits-lemma-integral-limit-finite-and-finite-presentation} を参照されたい）。
遷移射を $p_{i'i} : X_{i'} \to X_i$、射影を
$p_i : X \to X_i$ と書く。$\mathcal{F}_i = p_{i, *}\mathcal{F}$ とおくと、
補題 \ref{etale-cohomology-lemma-linus-hamann} により、系
$(\mathcal{F}_i, \varphi_{i'i})$ で
$\mathcal{F} = \colim p_i^{-1}\mathcal{F}_i$ を満たすものを得る。
$f_*\mathcal{F} = \colim f_{i, *}\mathcal{F}_i$ は補題
\ref{etale-cohomology-lemma-relative-colimit} から得られる。
$X'_i = Y' \times_Y X_i$ とおき、射影を $f'_i$ および $g'_i$ とする。
極限と極限が可換なので $X' = \lim X'_i$ である。
射影を $p'_i : X' \to X'_i$ と書く。このとき
\begin{align*}
g^{-1}f_*\mathcal{F}
& =
g^{-1} \colim f_{i, *}\mathcal{F}_i \\
& =
\colim g^{-1}f_{i, *}\mathcal{F}_i \\
& =
\colim f'_{i, *}(g'_i)^{-1}\mathcal{F}_i \\
& =
f'_*(\colim (p'_i)^{-1}(g'_i)^{-1}\mathcal{F}_i) \\
& =
f'_*(\colim (g')^{-1}p_i^{-1}\mathcal{F}_i) \\
& =
f'_*(g')^{-1} \colim p_i^{-1}\mathcal{F}_i \\
& =
f'_*(g')^{-1}\mathcal{F}
\end{align*}
となり、所望の結果を得る。第一の等式については上を参照されたい。
第二の等式には、引き戻しが余極限と可換であることを用いる。
第三の等式には有限の場合、すなわち補題
\ref{etale-cohomology-lemma-finite-pushforward-commutes-with-base-change} を用いる。
第四の等式には補題 \ref{etale-cohomology-lemma-relative-colimit} を用いる。
第五の等式には $g'_i \circ p'_i = p_i \circ g'$ を用いる。
第六の等式には、引き戻しが余極限と可換であることを用いる。
第七の等式には $\mathcal{F} = \colim p_i^{-1}\mathcal{F}_i$ を用いる。
\end{proof}

\noindent
次の補題は、エタール層と有限全射を扱うコホモロジー的降下の一例である。
固有基底変換定理を証明した後、この結果を大幅に一般化する。

\begin{lemma}
\label{etale-cohomology-lemma-cohomological-descent-finite}
$f : X \to Y$ をスキームの有限全射とする。
$f_n : X_n \to Y$ を、$(n + 1)$ 重ファイバー積で、
$X$ を $Y$ 上で取ったものに等しいものとする。
$\mathcal{F} \in \textit{Ab}(Y_\etale)$ に対し
$\mathcal{F}_n = f_{n, *}f_n^{-1}\mathcal{F}$ とおく。完全列
$$
0 \to \mathcal{F} \to \mathcal{F}_0 \to \mathcal{F}_1 \to
\mathcal{F}_2 \to \ldots
$$
が $Y_\etale$ 上に存在する。さらに、スペクトル系列
$$
E_1^{p, q} = H^q_\etale(X_p, f_p^{-1}\mathcal{F})
$$
が存在し、$H^{p + q}(Y_\etale, \mathcal{F})$ に収束する。
このスペクトル系列は $\mathcal{F}$ に関して関手的である。
\end{lemma}

\begin{proof}
補題の第一の主張を証明すれば、
$E_1^{p, q} = H^q_\etale(Y, \mathcal{F}_p)$ をもち
$H^{p + q}(Y_\etale, \mathcal{F})$ に収束するスペクトル系列を得る。
導来圏、補題 \ref{derived-lemma-two-ss-complex-functor} を参照されたい。
一方、$R^if_{p, *}f_p^{-1}\mathcal{F} = 0$ が $i > 0$ に対して成り立つので
（命題 \ref{etale-cohomology-proposition-finite-higher-direct-image-zero}）、
$$
H^q_\etale(X_p, f_p^{-1}\mathcal{F}) =
H^q_\etale(Y, f_{p, *}f_p^{-1} \mathcal{F}) =
H^q_\etale(Y, \mathcal{F}_p)
$$
を命題 \ref{etale-cohomology-proposition-leray} により得る。したがって補題のスペクトル系列を得る。

\medskip\noindent
補題の第一の主張を証明するため、$X_n$ が $Y$ 上の単体スキームをなすことに注意する。
単体的方法、例 \ref{simplicial-example-fibre-products-simplicial-object} を参照されたい。
さらに、各射影 $d_j : X_{n + 1} \to X_n$ に対して写像
$d_j^{-1} f_n^{-1}\mathcal{F} \to f_{n + 1}^{-1}\mathcal{F}$.
が存在する。これらの写像は写像
$$
\delta_j : \mathcal{F}_n \to \mathcal{F}_{n + 1}
$$
を $j = 0, \ldots, n + 1$ に対して誘導する。これらの写像の交代和により、
微分 $\mathcal{F}_n \to \mathcal{F}_{n + 1}$ を定める。
同様に、標準的な増大写像 $\mathcal{F} \to \mathcal{F}_0$ が存在する。
これはまさに標準写像 $\mathcal{F} \to f_*f^{-1}\mathcal{F}$ である。
この層の列が完全な複体であることを確かめるには、幾何点における茎で
確かめれば十分である（定理 \ref{etale-cohomology-theorem-exactness-stalks}）。
そこで $\overline{y} : \Spec(k) \to Y$ を幾何点とする。
$E = \{\overline{x} : \Spec(k) \to X \mid f(\overline{x}) = \overline{y}\}$ とおく。
$E$ は有限非空集合であり、
$$
(\mathcal{F}_n)_{\overline{y}} =
\bigoplus\nolimits_{e \in E^{n + 1}} \mathcal{F}_{\overline{y}}
$$
が命題 \ref{etale-cohomology-proposition-finite-higher-direct-image-zero} と補題
\ref{etale-cohomology-lemma-stalk-pullback} により成り立つ。したがって、アーベル群 $M$ が
与えられたとき、列
$$
0 \to M \to \bigoplus\nolimits_{e \in E} M \to
\bigoplus\nolimits_{e \in E^2} M \to \ldots
$$
が完全であることを示せばよい。ここで第一の写像は対角写像であり、写像
$\bigoplus_{e \in E^{n + 1}} M  \to \bigoplus_{e \in E^{n + 2}} M$
は、$(n + 2)$ 個の射影 $E^{n + 2} \to E^{n + 1}$ が誘導する写像の
交代和である。これは直接示すこともできるし、単体的方法、補題
\ref{simplicial-lemma-fibre-products-simplicial-object-w-section} を
写像 $E \to \{*\}$ に適用して導くこともできる。
\end{proof}

\begin{remark}
\label{etale-cohomology-remark-cohomological-descent-finite}
補題 \ref{etale-cohomology-lemma-cohomological-descent-finite} の状況で、
$\mathcal{G}$ が $Y_\etale$ 上の集合の層ならば
$$
\Gamma(Y, \mathcal{G}) =
\text{等化子}(
\xymatrix{
\Gamma(X_0, f_0^{-1}\mathcal{G})
\ar@<1ex>[r] \ar@<-1ex>[r] &
\Gamma(X_1, f_1^{-1}\mathcal{G})
}
)
$$
である。これはまったく同じ方法で、層 $\mathcal{G}$ が二つの写像
$f_{0, *}f_0^{-1}\mathcal{G} \to f_{1, *}f_1^{-1}\mathcal{G}$.
の等化子であることを示して証明される。
\end{remark}









%10.06.09
\section{茎上のガロア作用}
\label{etale-cohomology-section-galois-action-stalks}

\noindent
この節では、点 $s$ で $S$ のものの剰余体の絶対ガロア群が、
$s$ 上にある任意の幾何点での茎関手に作用することを定義する。

\medskip\noindent
茎上のガロア作用。
$S$ をスキームとする。
$\overline{s}$ を $S$ の幾何点とする。
$\sigma \in \text{Aut}(\kappa(\overline{s})/\kappa(s))$ とする。
$\sigma$ の茎 $\mathcal{F}_{\overline{s}}$ への作用を、層 $\mathcal{F}$ に対して
次のように定義する：
\begin{equation}
\label{etale-cohomology-equation-galois-action}
\begin{matrix}
\mathcal{F}_{\overline{s}} &
\longrightarrow &
\mathcal{F}_{\overline{s}} \\
(U, \overline{u}, t) &
\longmapsto &
(U, \overline{u} \circ \Spec(\sigma), t).
\end{matrix}
\end{equation}
ここでは、定義 \ref{etale-cohomology-definition-stalk} に続く議論にあるように、
茎の元を三つ組で記述している。これは左作用である。実際、
$\sigma_i \in \text{Aut}(\kappa(\overline{s})/\kappa(s))$ ならば
\begin{align*}
\sigma_1 \cdot (\sigma_2 \cdot (U, \overline{u}, t))
& =
\sigma_1 \cdot (U, \overline{u} \circ \Spec(\sigma_2), t) \\
& =
(U, \overline{u} \circ \Spec(\sigma_2) \circ \Spec(\sigma_1), t) \\
& =
(U, \overline{u} \circ \Spec(\sigma_1 \circ \sigma_2), t) \\
& =
(\sigma_1 \circ \sigma_2) \cdot (U, \overline{u}, t)
\end{align*}
である。この作用が層 $\mathcal{F}$ に関して関手的であることは明らかである。
また、注意 \ref{etale-cohomology-remark-map-stalks} を直接参照してこの作用を定義することもできた。

\begin{definition}
\label{etale-cohomology-definition-algebraic-geometric-point}
$S$ をスキームとする。
$\overline{s}$ を点 $s$ で $S$ のものの上にある幾何点とする。
$\kappa(s) \subset \kappa(s)^{sep} \subset \kappa(\overline{s})$ において、
$\kappa(s)$ の分離代数閉包を、代数閉体 $\kappa(\overline{s})$ の中で取ったものとする。
\begin{enumerate}
\item この状況で $\kappa(s)$ の{\it 絶対ガロア群}とは
$\text{Gal}(\kappa(s)^{sep}/\kappa(s))$ のことである。
これは $\text{Gal}_{\kappa(s)}$ と書かれることもある。
\item 幾何点 $\overline{s}$ は、$\kappa(s) \subset \kappa(\overline{s})$ が
$\kappa(s)$ の代数閉包であるとき{\it 代数的}であるという。
\end{enumerate}
\end{definition}

\begin{example}
\label{etale-cohomology-example-stupid}
幾何点
$\Spec(\mathbf{C}) \to \Spec(\mathbf{Q})$
は代数的でない。
\end{example}

\noindent
$\kappa(s) \subset \kappa(s)^{sep} \subset \kappa(\overline{s})$ を
定義のとおりとする。$\kappa(\overline{s})$ は代数閉なので、写像
$$
\text{Aut}(\kappa(\overline{s})/\kappa(s))
\longrightarrow
\text{Gal}(\kappa(s)^{sep}/\kappa(s)) = \text{Gal}_{\kappa(s)}
$$
は全射であることに注意する。$(U, \overline{u})$ を
$\overline{s}$ のエタール近傍とし、$\overline{u}$ は
点 $u$ で $U$ のものの上にあるとする。$U \to S$ はエタールなので、
剰余体拡大 $\kappa(u)/\kappa(s)$ は有限分離的である。
これから次が従う。
\begin{enumerate}
\item $\sigma \in \text{Aut}(\kappa(\overline{s})/\kappa(s)^{sep})$ ならば、
$\sigma$ は $\mathcal{F}_{\overline{s}}$ に自明に作用する。
\item より正確には、$\text{Aut}(\kappa(\overline{s})/\kappa(s))$ の作用は、
絶対ガロア群 $\text{Gal}_{\kappa(s)}$ の $\mathcal{F}_{\overline{s}}$ への作用を
定め、またそれによって定められる。
\item $(U, \overline{u}, t)$ が元 $\xi$ で $\mathcal{F}_{\overline{s}}$ に属するものを表すとき、
$\text{Gal}(\kappa(s)^{sep}/K)$ の任意の元は自明に作用する。ここで
$\kappa(s) \subset K \subset \kappa(s)^{sep}$ は
$\overline{u}^\sharp : \kappa(u) \to \kappa(\overline{s})$ の像である。
\end{enumerate}
以上より、$\mathcal{F}_{\overline{s}}$ は $\text{Gal}_{\kappa(s)}$-集合となる
（基本群、定義 \ref{pione-definition-G-set-continuous} を参照されたい）。
したがって、茎関手を関手
$$
\Sh(S_\etale) \longrightarrow
\text{Gal}_{\kappa(s)}\textit{-Sets},
\quad
\mathcal{F} \longmapsto \mathcal{F}_{\overline{s}}
$$
とみなしてよく、以後は通常このように茎関手を考える。

\begin{theorem}
\label{etale-cohomology-theorem-equivalence-sheaves-point}
$S = \Spec(K)$ とし、$K$ は体であるとする。
$\overline{s}$ を $S$ の幾何点とする。
$G = \text{Gal}_{\kappa(s)}$ で絶対ガロア群を表す。
茎を取ることは圏同値
$$
\Sh(S_\etale) \longrightarrow G\textit{-Sets},
\quad
\mathcal{F} \longmapsto \mathcal{F}_{\overline{s}}.
$$
を誘導する。
\end{theorem}

\begin{proof}
この関手の逆を構成しよう。基本群、補題
\ref{pione-lemma-sheaves-point} で、$G$-集合 $M$ が与えられるとエタール射
$X \to \Spec(K)$
で、$\Mor_K(\Spec(K^{sep}), X)$ が $M$ と $G$-集合として同型となるものが
存在することを見た。層 $\mathcal{F}$ で、$\Spec(K)_\etale$ 上にあり、
規則 $U \mapsto \Mor_K(U, X)$ により定義されるものを考える。
エタール位相は準標準的なので、これは層である。このとき
$\mathcal{F}_{\overline{s}} = \Mor_K(\Spec(K^{sep}), X) = M$
が $G$-集合として成り立つ（細部は省略する）。これで関手の逆を得て、
証明が終わる。
\end{proof}

\begin{remark}
\label{etale-cohomology-remark-every-sheaf-representable}
定理 \ref{etale-cohomology-theorem-equivalence-sheaves-point} と基本群、補題
\ref{pione-lemma-sheaves-point} の結論は、$\Spec(K)_\etale$ 上の任意の層が、
スキーム $X$ で $\Spec(K)$ 上エタールなものによって表現可能である、と言い換えられる。
これは、任意の層がサイト、定義
\ref{sites-definition-representable-sheaf} の意味で表現可能だということではない。
その理由は、$\Spec(K)_\etale$ の構成において、$\Spec(K)$ 上エタールなスキームの
十分大きい集合を選んだのに対し、$\Spec(K)_\etale$ 上の層は真の類をなすからである。
\end{remark}

\begin{lemma}
\label{etale-cohomology-lemma-global-sections-point}
仮定と記法を定理 \ref{etale-cohomology-theorem-equivalence-sheaves-point} と同じものとする。
関手的な全単射
$$
\Gamma(S, \mathcal{F}) = (\mathcal{F}_{\overline{s}})^G
$$
が存在する。
\end{lemma}

\begin{proof}
これは形式的な議論と定理 \ref{etale-cohomology-theorem-equivalence-sheaves-point} の結果を用いて、
次のように証明できる。層 $\mathcal{F}$ で、
$G$-集合 $M = \mathcal{F}_{\overline{s}}$ に対応するものが与えられたとき、
\begin{eqnarray*}
\Gamma(S, \mathcal{F}) & = &
\Mor_{\Sh(S_\etale)}(h_{\Spec(K)},  \mathcal{F})
\\
& = & \Mor_{G\textit{-Sets}}(\{*\}, M) \\
& = & M^G
\end{eqnarray*}
である。第一の同一視はサイト、節 \ref{sites-section-presheaves} および
\ref{sites-section-representable-sheaves} で説明されている。
第二の同一視は定理 \ref{etale-cohomology-theorem-equivalence-sheaves-point} から従い、
第三の同一視は明らかである。直接証明も与える\footnote{疑い深い読者のために。}。

\medskip\noindent
$t \in \Gamma(S, \mathcal{F})$ を大域切断とする。
すると三つ組 $(S, \overline{s}, t)$ は $\mathcal{F}_{\overline{s}}$ の元を定め、
これは明らかに $G$ の作用で不変である。逆に、$(U, \overline{u}, t)$ が
$\mathcal{F}_{\overline{s}}$ の不変な元を定めるとする。
$U$ を縮小し、$U = \Spec(L)$ で、$K$ のある有限分離体拡大に対して
成り立つと仮定してよい。命題 \ref{etale-cohomology-proposition-etale-morphisms} を参照されたい。
この場合、写像 $\mathcal{F}(U) \to \mathcal{F}_{\overline{s}}$ は単射である。
実際、エタール近傍の任意の射
$(U', \overline{u}') \to (U, \overline{u})$ に対して、制限写像
$\mathcal{F}(U) \to \mathcal{F}(U')$ は単射である。これは $U' \to U$ が
$S_\etale$ の被覆だからである。
$L$ を少し大きくすれば、$K \subset L$ は有限ガロア拡大であると仮定してよい。
ここで
$$
\Spec(L) \times_{\Spec(K)} \Spec(L)
=
\coprod\nolimits_{\sigma \in \text{Gal}(L/K)} \Spec(L)
$$
を用いる。ここで写像 $\Spec(L) \to \Spec(L \otimes_K L)$ は環準同型
$a \otimes b \mapsto a\sigma(b)$ から来る。したがって、
$(U, \overline{u}, t)$ が $G$ のすべての元で不変であるという条件から、
$t \in \mathcal{F}(\Spec(L))$ は二つの射影のいずれによる制限でも
$\mathcal{F}(\Spec(L) \times_{\Spec(K)} \Spec(L))$
の同じ元へ写ることがわかる（ここでは上で述べた単射性を用いる。細部は省略する）。
したがって、層 $\mathcal{F}$ の、エタール被覆
$\{\Spec(L) \to \Spec(K)\}$ に対する条件を用いると、$t$ は $\mathcal{F}$ の
$\Spec(K)$ 上の一意な切断から来ると結論できる。
\end{proof}

\begin{remark}
\label{etale-cohomology-remark-stalk-pullback}
$S$ をスキームとし、$\overline{s} : \Spec(k) \to S$ を
$S$ の幾何点とする。定義により、これは $k$ が代数閉であることを意味する。
とくに $k$ の絶対ガロア群は自明である。したがって定理
\ref{etale-cohomology-theorem-equivalence-sheaves-point} により、$\Spec(k)_\etale$ 上の層の圏は
集合の圏と同値である。この同値は $\Spec(k)$ 上の切断を取ることによって与えられる。
これにより、茎関手の別の定義が得られる。すなわち関手
$$
\Sh(S_\etale) \longrightarrow \textit{Sets}, \quad
\mathcal{F} \longmapsto \mathcal{F}_{\overline{s}}
$$
は関手
$$
\Sh(S_\etale)
\longrightarrow
\Sh(\Spec(k)_\etale) = \textit{Sets},
\quad
\mathcal{F} \longmapsto \overline{s}^*\mathcal{F}
$$
と同型である。これを厳密に証明するには、補題
\ref{etale-cohomology-lemma-stalk-pullback} の (3) を $f = \overline{s}$ として用いればよい。
さらに、これがわかれば、補題 \ref{etale-cohomology-lemma-stalk-pullback} の (3) の一般の場合は
引き戻しの関手性から従う。
\end{remark}






\section{群コホモロジー}
\label{etale-cohomology-section-group-cohomology}

\noindent
以下で $H^i(G, M)$ と書くとき、$G$ は位相群であり、$M$ は連続な
$G$-作用をもつ離散 $G$-加群であり、$H^i(G, -)$ は第 $i$ 右導来関手を表す。
その定義域は圏 $\text{Mod}_G$、すなわちそのような $G$-加群の圏である。
定義 \ref{etale-cohomology-definition-G-module-continuous} および
\ref{etale-cohomology-definition-galois-cohomology} を参照されたい。これには抽象群
$G$ の場合も含まれ、その場合は単に $G$ を離散位相をもつ位相群とみなす。

\medskip\noindent
加群が非離散位相をもつ場合には、\cite{Tate} で導入された連続コホモロジー群を
$H^i_{cont}(G, M)$ と書く。節 \ref{etale-cohomology-section-continuous-group-cohomology} を参照されたい。

\begin{definition}
\label{etale-cohomology-definition-G-module-continuous}
$G$ を位相群とする。
\begin{enumerate}
\item {\it $G$-加群}（{\it 離散 $G$-加群}ともいう）とは、アーベル群 $M$ と、
群準同型による左作用 $a : G \times M \to M$ で、$a$ が、$M$ に離散位相を入れたとき
連続となるものとの組である。
\item {\it $G$-加群の射} $f : M \to N$ とは、$G$-同変準同型で、
$M$ から $N$ へのものである。
\item $G$-加群の圏を $\text{Mod}_G$ と書く。
\end{enumerate}
$R$ を環とする。
\begin{enumerate}
\item {\it $R\text{-}G$-加群}とは、$R$-加群 $M$ と、
左作用 $a : G \times M \to M$ で、$R$-線形写像によるものであり、
$a$ が、$M$ に離散位相を入れたとき連続となるものとの組である。
\item {\it $R\text{-}G$-加群の射} $f : M \to N$ とは、$G$-同変な
$R$-加群準同型で、$M$ から $N$ へのものである。
\item $R\text{-}G$-加群の圏を $\text{Mod}_{R, G}$ と書く。
\end{enumerate}
\end{definition}

\noindent
$a : G \times M \to M$ が連続であるという条件は、任意の $x \in M$ の
固定部分群が $G$ の開部分群であるという条件と同値である。
$G$ が抽象群ならば、$G$ に離散位相を入れることで、これは
$G$-作用をもつアーベル群という概念に一致する。
$\text{Mod}_{\mathbf{Z}, G} = \text{Mod}_G$ であることに注意する。

\medskip\noindent
圏 $\text{Mod}_G$ は十分な単射的対象をもつ。単射的対象、補題
\ref{injectives-lemma-G-modules} を参照されたい。左完全関手
$$
\text{Mod}_G \longrightarrow \textit{Ab},
\quad
M \longmapsto M^G =
\{x \in M \mid g \cdot x = x\ \forall g \in G\}
$$
を考える。$M^G = H^0(G, M)$ と書くこともあれば、
$M^G = \Gamma_G(M)$ と書くこともある。この関手は全右導来関手
$R\Gamma_G(M)$ および第 $i$ 右導来関手
$R^i\Gamma_G(M) = H^i(G, M)$ をもち、後者は任意の $i \geq 0$ に対して定義される。

\medskip\noindent
同じ構成は $H^0(G, -) : \text{Mod}_{R, G} \to \text{Mod}_R$ にも適用できる。
補題 \ref{etale-cohomology-lemma-modules-abelian} で、これが台となる $G$-加群の
コホモロジーと一致することを見る。

\begin{definition}
\label{etale-cohomology-definition-galois-cohomology}
$G$ を位相群とする。$M$ を、連続な $G$-作用をもつ離散 $G$-加群とする。
言い換えると、$M$ は定義 \ref{etale-cohomology-definition-G-module-continuous} で導入した
圏 $\text{Mod}_G$ の対象である。
\begin{enumerate}
\item 右導来関手 $H^i(G, M)$、すなわち $H^0(G, M)$ の
圏 $\text{Mod}_G$ 上での右導来関手を、$M$ の{\it 連続群コホモロジー群}という。
\item $G$ が離散位相を入れた抽象群ならば、$H^i(G, M)$ を
$M$ の{\it 群コホモロジー群}という。
\item $G$ がガロア群ならば、群 $H^i(G, M)$ を
$M$ の{\it ガロアコホモロジー群}という。
\item $G$ が体 $K$ の絶対ガロア群ならば、群 $H^i(G, M)$ を
{\it $K$ のガロアコホモロジー群で $M$ を係数とするもの}ということがある。
この場合、$H^i(K, M)$ と書くことがあり、これは $H^i(G, M)$ の代わりである。
\end{enumerate}
\end{definition}

\begin{lemma}
\label{etale-cohomology-lemma-modules-abelian}
$G$ を位相群とする。$R$ を環とする。
すべての $i \geq 0$ に対して、図式
$$
\xymatrix{
\text{Mod}_{R, G} \ar[rr]_{H^i(G, -)} \ar[d] & &
\text{Mod}_R \ar[d] \\
\text{Mod}_G \ar[rr]^{H^i(G, -)} & &
\textit{Ab}
}
$$
は可換である。ここで縦の矢印は忘却関手である。
\end{lemma}

\begin{proof}
忘却関手を $F : \text{Mod}_{R, G} \to \text{Mod}_G$ と書く。
$F$ は左随伴 $H : \text{Mod}_G \to \text{Mod}_{R, G}$ をもち、
$H(M) = M \otimes_\mathbf{Z} R$ によって与えられる。
$\text{Mod}_G$ の任意の対象は、$\mathbf{Z}[G/U]$ という形の加群の
直和の商である。ここで $U \subset G$ は開部分群である。
$\mathbf{Z}[G/U]$ は $G$-加群を表す。その元は有限 $\mathbf{Z}$-線形結合であり、
右 $U$ 合同類を $G$ の中で取ったものからなる。この加群には左 $G$-作用を入れる。
したがって、$\text{Mod}_G$ の任意の上に有界な複体は、
$\text{Mod}_G$ の上に有界な複体で、その各項が平坦な $\mathbf{Z}$-加群であるものと
擬同型である（導来圏、補題 \ref{derived-lemma-subcategory-left-resolution}）。
ゆえに $LH$ が $D^-(\text{Mod}_G)$ 上に存在し、$H$ を各項が平坦な
$\mathbf{Z}$-加群である任意の複体に適用して計算されることは明らかである。
これは導来圏、補題 \ref{derived-lemma-subcategory-left-acyclics} および命題
\ref{derived-proposition-enough-acyclics} から従う。
導来圏、補題 \ref{derived-lemma-pre-derived-adjoint-functors} から
$$
\text{Ext}^i(\mathbf{Z}, F(M)) = \text{Ext}^i(R, M)
$$
が、$M$ が $\text{Mod}_{R, G}$ に属するとき成り立つことがわかる。
$H^0(G, -) = \Hom(\mathbf{Z}, -)$ が $\text{Mod}_G$ 上で成り立つことに注意する。
ここで $\mathbf{Z}$ は自明な作用をもつ $G$-加群を表す。したがって
$H^i(G, -) = \text{Ext}^i(\mathbf{Z}, -)$ が $\text{Mod}_G$ 上で成り立つ。
同様に、$H^i(G, -) = \text{Ext}^i(R, -)$ が
$\text{Mod}_{R, G}$ 上で成り立つ。以上をあわせれば補題が従う。
\end{proof}

\begin{lemma}
\label{etale-cohomology-lemma-ext-modules-hom}
$G$ を位相群とする。$R$ を環とする。
$M$、$N$ を $R\text{-}G$-加群とする。$M$ が $R$-加群として有限射影的ならば
$\text{Ext}^i(M, N) = H^i(G, M^\vee \otimes_R N)$ である
（記法については証明を参照されたい）。
\end{lemma}

\begin{proof}
加群 $M^\vee = \Hom_R(M, R)$ には $G$ の反傾作用を入れる。
すなわち $(g \cdot \lambda)(m) = \lambda(g^{-1} \cdot m)$ が
$g \in G$、$\lambda \in M^\vee$、$m \in M$ に対して成り立つ。
$G$ の $M^\vee \otimes_R N$ への作用は対角作用、すなわち
$g \cdot (\lambda \otimes n) = g \cdot \lambda \otimes g \cdot n$.
で与えられる。第三の $R\text{-}G$-加群 $E$ に対して
$\Hom(E, M^\vee \otimes_R N) = \Hom(M \otimes_R E, N)$.
である。実際、これは $R$-加群の段階では代数、補題
\ref{algebra-lemma-hom-from-tensor-product} および
\ref{algebra-lemma-evaluation-map-iso-finite-projective} により成り立ち、
$G$-作用の定義は、$R\text{-}G$-加群についても成り立つように選ばれている。
したがって、$M^\vee \otimes_R N$ は単射的 $R\text{-}G$-加群である。
これは $N$ が単射的 $R\text{-}G$-加群である場合に成り立つ。
ゆえに $N \to N^\bullet$ が単射分解ならば、
$M^\vee \otimes_R N \to M^\vee \otimes_R N^\bullet$
も単射分解である。このとき
$$
\Hom(M, N^\bullet) = \Hom(R, M^\vee \otimes_R N^\bullet) =
(M^\vee \otimes_R N^\bullet)^G
$$
左辺は $\text{Ext}^i(M, N)$ を、右辺は
$H^i(G, M^\vee \otimes_R N)$ を計算するので、証明が完了する。
\end{proof}

\begin{lemma}
\label{etale-cohomology-lemma-finite-dim-group-cohomology}
$G$ を位相群とする。$k$ を体とする。
$V$ を $k\text{-}G$-加群とする。
$G$ が位相的に有限生成で $\dim_k(V) < \infty$ ならば、
$\dim_k H^1(G, V) < \infty$ である。
\end{lemma}

\begin{proof}
$g_1, \ldots, g_r \in G$ を $G$ を位相的に生成する元とする。
すなわち、$g_1, \ldots, g_r$ が生成する部分群は稠密である。
補題 \ref{etale-cohomology-lemma-ext-modules-hom} により、$H^1(G, V)$ は $k$-ベクトル空間で、
その元は拡大
$$
0 \to V \to E \to k \to 0
$$
で $k\text{-}G$-加群のものによって与えられる。
$e \in E$ で $1 \in k$ へ写るものを選ぶ。
$$
g_i \cdot e = v_i + e
$$
と、ある $v_i \in V$ に対して書く。これは $g_i \cdot 1 = 1$ だから可能である。
元の列 $v_1, \ldots, v_r \in V$ が拡大 $E$ の同型類を定めると主張する。
これを証明すれば補題が従う。実際、これは考えている拡大類のベクトル空間が
$k$-ベクトル空間 $V^{\oplus r}$ の部分商と同型であることを意味する。
細部は省略する。
$E$ は定義 \ref{etale-cohomology-definition-G-module-continuous} で定めた圏の対象なので、
開部分群 $U$ が存在して、$u \cdot e = e$ がすべての $u \in U$ に対して成り立つ。
任意の $g \in G$ を選ぶ。すると $gU$ は、語 $w$ で、
元 $g_1, \ldots, g_r$ からなるものを含む。
$gu = w$ とする。元 $w \cdot e$ は $v_1, \ldots, v_r$ によって定まるので、
$g \cdot e = (gu) \cdot e = w \cdot e$ もそうである。
\end{proof}

\begin{lemma}
\label{etale-cohomology-lemma-profinite-group-cohomology-torsion}
$G$ を副有限位相群とする。このとき
\begin{enumerate}
\item $H^i(G, M)$ は $i > 0$ および任意の $G$-加群 $M$ に対して捩れ群である。
\item $H^i(G, M) = 0$ は、$M$ が $\mathbf{Q}$-ベクトル空間ならば成り立つ。
\end{enumerate}
\end{lemma}

\begin{proof}
(1) の証明。次元移動により、$H^1(G, M)$ が任意の $G$-加群 $M$ に対して
捩れ群であることを示せば十分である。完全列
$0 \to M \to I \to N \to 0$ で、$I$ が $G$-加群の圏の単射的対象であるものを選ぶ。
すると $H^1(G, M)$ の任意の元は、ある元 $y \in N^G$ の像である。
$x \in I$ で $y$ へ写るものを選ぶ。安定化部分群 $U \subset G$ で
$x$ のものは開であり、
したがって有限指数 $r$ をもつ。$g_1, \ldots, g_r \in G$ を
$G/U$ の代表元系とする。このとき $\sum g_i(x)$ は $I$ の不変元で、
$ry$ へ写る。したがって $r$ は、最初に取った $H^1(G, M)$ の元を零化する。
(2) は、このとき $H^i(G, M)$ が $\mathbf{Q}$-ベクトル空間であると同時に
捩れ群でもあることから従う。
\end{proof}





\section{テイトによる連続コホモロジー}
\label{etale-cohomology-section-continuous-group-cohomology}

\noindent
テイトによる連続コホモロジー（\cite{Tate}）は、連続な非斉次余鎖の複体によって
定義される。これは、$M$ が連続な $G$-作用を備えた任意の位相アーベル群であるときに
定義できる。すなわち、複体
$$
C^\bullet_{cont}(G, M) :
M \to \text{写像}_{cont}(G, M) \to
\text{写像}_{cont}(G \times G, M) \to \ldots
$$
を考える。ここで余微分は $n \geq 1$ に対して次式により定義される。
\begin{align*}
\text{d}(f)(g_1, \ldots, g_{n + 1})
& = g_1(f(g_2, \ldots, g_{n + 1})) \\
&
+ \sum\nolimits_{j = 1, \ldots, n}
(-1)^jf(g_1, \ldots, g_jg_{j + 1}, \ldots, g_{n + 1}) \\
&
+ (-1)^{n + 1}f(g_1, \ldots, g_n)
\end{align*}
また $n = 0$ のとき、$m \in M$ を写像 $g \mapsto g(m) - m$ へ送るものとする。次で定義する。
$$
H^i_{cont}(G, M) = H^i(C^\bullet_{cont}(G, M))
$$
この複体の各項には、$G$ および $G$ の自己積から位相加群 $M$ への連続写像が現れる。
したがって、位相加群の短完全列がコホモロジーの長完全列へ移るかどうかは明らかでない。
もう一つの困難は、位相アーベル群の圏がアーベル圏ではないことである。

\medskip\noindent
しかし、離散 $G$-加群の短完全列からは、連続余鎖複体の短完全列が実際に得られ、
したがって連続コホモロジー群 $H^i_{cont}(G, -)$ の長完全列が得られる。
それゆえ、定義 \ref{etale-cohomology-definition-G-module-continuous} の圏 $\text{Mod}_G$ 上で、関手
$H^i_{cont}(G, M)$ は、ホモロジー、節
\ref{homology-section-cohomological-delta-functor} で定義されるコホモロジー的
$\delta$-関手をなす。定義 \ref{etale-cohomology-definition-galois-cohomology} のコホモロジー
$H^i(G, M)$ は普遍 $\delta$-関手であるから（導来圏、補題
\ref{derived-lemma-right-derived-delta-functor}）、標準写像
$$
H^i(G, M) \longrightarrow H^i_{cont}(G, M)
$$
を $M \in \text{Mod}_G$ に対して得る。$G$ が抽象群である場合（すなわち、$G$ が
離散位相をもつ場合）、または $G$ が副有限群である場合には、これらの写像が同型で
あることが知られている（将来ここに参照を挿入）。位相群 $G$ と
$M \in \Ob(\text{Mod}_G)$ に対してこの写像が同型でないことを示す例をご存じならば、
\href{mailto:stacks.project@gmail.com}{stacks.project@gmail.com} まで電子メールを送られたい。




\section{点のコホモロジー}
\label{etale-cohomology-section-cohomology-point}

\noindent
前節までの議論の帰結として、体のスペクトルのエタール・コホモロジーと
ガロアコホモロジーとの同値を得る。

\begin{lemma}
\label{etale-cohomology-lemma-equivalence-abelian-sheaves-point}
$S = \Spec(K)$ とし、$K$ は体であるとする。
$\overline{s}$ を $S$ の幾何点とする。
$G = \text{Gal}_{\kappa(s)}$ で絶対ガロア群を表す。
茎関手は圏の同値
$$
\textit{Ab}(S_\etale) \longrightarrow \text{Mod}_G,
\quad
\mathcal{F} \longmapsto \mathcal{F}_{\overline{s}}.
$$
を誘導する。
\end{lemma}

\begin{proof}
定理 \ref{etale-cohomology-theorem-equivalence-sheaves-point} で、集合の層と $G$-集合との
同値を見た。アーベル群の層とは可換群構造を備えた集合の層にほかならず、
$G$-加群とは可換群構造を備えた $G$-集合にほかならないので、本補題は
この同値から形式的に従う。
\end{proof}

\begin{lemma}
\label{etale-cohomology-lemma-compare-cohomology-point}
補題 \ref{etale-cohomology-lemma-equivalence-abelian-sheaves-point} の記法と仮定を用いる。
$\mathcal{F}$ を $\Spec(K)_\etale$ 上のアーベル群の層で、
$G$-加群 $M$ に対応するものとする。このとき
\begin{enumerate}
\item $D(\textit{Ab})$ において標準同型
$R\Gamma(S, \mathcal{F}) = R\Gamma_G(M)$,
が成り立つ。
\item $H_\etale^0(S, \mathcal{F}) = M^G$ である。
\item $H_\etale^q(S, \mathcal{F}) = H^q(G, M)$ である。
\end{enumerate}
\end{lemma}

\begin{proof}
補題 \ref{etale-cohomology-lemma-equivalence-abelian-sheaves-point} と
補題 \ref{etale-cohomology-lemma-global-sections-point} を組み合わせればよい。
\end{proof}

\begin{example}
\label{etale-cohomology-example-sheaves-point}
$\Spec(K)_\etale$ 上の層を考える。
$G = \text{Gal}(K^{sep}/K)$ を $K$ の絶対ガロア群とする。
\begin{enumerate}
\item 定数層 $\underline{\mathbf{Z}/n\mathbf{Z}}$ は、加群
$\mathbf{Z}/n\mathbf{Z}$ で自明な $G$-作用をもつものに対応する。
\item 層 $\mathbf{G}_m|_{\Spec(K)_\etale}$ は $(K^{sep})^*$ に、
その $G$-作用を備えたものとして対応する。
\item 層 $\mathbf{G}_a|_{\Spec(K^{sep})}$ は $(K^{sep}, +)$ に、
その $G$-作用を備えたものとして対応する。
\item 層 $\mu_n|_{\Spec(K^{sep})}$ は $\mu_n(K^{sep})$ に、
その $G$-作用を備えたものとして対応する。
\end{enumerate}
注意 \ref{etale-cohomology-remark-special-case-fpqc-cohomology-quasi-coherent} および
定理 \ref{etale-cohomology-theorem-picard-group} により、コホモロジー群について次の同一視を得る。
\begin{align*}
H_\etale^0(S_\etale, \mathbf{G}_m) & =
\Gamma(S, \mathcal{O}_S^*) \\
H_\etale^1(S_\etale, \mathbf{G}_m) & =
H_{Zar}^1(S, \mathcal{O}_S^*) = \Pic(S) \\
H_\etale^i(S_\etale, \mathbf{G}_a) & =
H_{Zar}^i(S, \mathcal{O}_S)
\end{align*}
また、任意の準連接層 $\mathcal{F}$ で $S_\etale$ 上のものに対して
$$
H^i(S_\etale, \mathcal{F}) = H_{Zar}^i(S, \mathcal{F}),
$$
が成り立つ。定理 \ref{etale-cohomology-theorem-zariski-fpqc-quasi-coherent} を参照されたい。
とくに、次の等式列を得る。
$$
0 =
\Pic(\Spec(K)) =
H_\etale^1(\Spec(K)_\etale, \mathbf{G}_m) =
H^1(G, (K^{sep})^*)
$$
これはヒルベルトの定理 90 にほかならない。同様に、$i \geq 1$ に対して
$$
0 = H^i(\Spec(K), \mathcal{O})
= H_\etale^i(\Spec(K)_\etale, \mathbf{G}_a)
= H^i(G, K^{sep})
$$
である。ここで $K^{sep}$ は、$K^{sep}$ を加法を群法とするガロア加群として
見たものを表す。このように、これまでの作業は、ガロアコホモロジー群を計算する
複雑な方法であったとみなすこともできる。
\end{example}

\noindent
次の結果は珍しい事実であり、初読時には読み飛ばすのがよい。

\begin{lemma}
\label{etale-cohomology-lemma-all-modules-quasi-coherent}
$R$ を次元 $0$ の局所環とする。$S = \Spec(R)$ とする。
このとき、任意の $\mathcal{O}_S$-加群で $S_\etale$ 上のものは準連接である。
\end{lemma}

\begin{proof}
$\mathcal{F}$ を $\mathcal{O}_S$-加群で $S_\etale$ 上のものとする。
$\mathcal{F}$ が $R$-加群 $M = \Gamma(S, \mathcal{F})$ によって定まることを示せばよい。
より正確には、$\pi : X \to S$ がエタールならば、
$\Gamma(X, \mathcal{F}) = \Gamma(X, \pi^*\widetilde{M})$ を示せばよい。

\medskip\noindent
$\mathfrak m \subset R$ を極大イデアルとし、$\kappa$ を剰余体とする。
代数、補題 \ref{algebra-lemma-local-dimension-zero-henselian} により、局所環 $R$ は
Hensel である。$X \to S$ がエタールならば、射の性質、補題
\ref{morphisms-lemma-etale-over-field} により、$X$ の台となる位相空間は離散である。
したがって $X$ は、各々が一点をもつアフィンスキームの非交和である。
さらに、$X = \Spec(A)$ が一点をもつアフィンスキームならば、代数、補題
\ref{algebra-lemma-mop-up} により、$R \to A$ は有限エタールである。
この場合、$\Gamma(X, \mathcal{F}) = M \otimes_R A$ を示せばよい。

\medskip\noindent
代数、補題 \ref{algebra-lemma-henselian-cat-finite-etale} により、関手
$A \mapsto A/\mathfrak m A$ は、有限エタール $R$-代数の圏と有限分離
$\kappa$-代数の圏との同値を定める。まず、$A/\mathfrak m A$ が
$\kappa$ のガロア拡大で、ガロア群が $G$ である場合を考える。
各 $\sigma \in G$ に対し、$\sigma : A \to A$ で、対応する $A$ の自己同型で
$R$ 上のものを表す。$N = \Gamma(X, \mathcal{F})$ とする。
このとき $\Spec(\sigma) : X \to X$ は $S$ 上の自己同型であり、したがって
これによる引き戻しは写像 $\sigma : N \to N$ を定める。これは
$\sigma$-線形写像、すなわち $\sigma(an) = \sigma(a) \sigma(n)$ を
$a \in A$ および $n \in N$ に対して満たす写像である。
準連接加群 $\widetilde{N}$ で $X$ 上のものに、$\sigma$ の $N$ への作用から得られる
同型を備え、ガロア降下を適用する。降下、補題
\ref{descent-lemma-galois-descent-more-general} を参照されたい。
この補題から同型 $N = N^G \otimes_R A$ が存在することがわかる。
一方、$N^G = M$ であることは、$\mathcal{F}$ の層条件により明らかである。
したがって所望の同型が成り立つ。

\medskip\noindent
一般の場合（$A$ が局所的で、$R$ 上有限エタールである場合）は、次のように
ガロアの場合から導かれる。有限エタールな $A \to B$ を、$B$ が局所的で、
その剰余体が $\kappa$ のガロア拡大となるように選ぶ。
$G = \text{Aut}(B/R) = \text{Gal}(\kappa_B/\kappa)$ とする。
$H \subset G$ をガロア拡大 $\kappa_B/\kappa_A$ に対応するガロア群とする。
すると上と同様に、$\Gamma(X, \mathcal{F}) = \Gamma(\Spec(B), \mathcal{F})^H$ が示される。
ガロア拡大についての結果を二度用いると、
$$
\Gamma(X, \mathcal{F}) = (M \otimes_R B)^H = M \otimes_R A
$$
を得る。これが所望の等式である。
\end{proof}








\section{曲線のコホモロジー}
\label{etale-cohomology-section-cohomology-curves}

\noindent
次の課題は、代数閉体上の滑らかな曲線のエタール・コホモロジーを捩れ係数で計算し、
とくに次数が少なくとも 3 なら消滅することを示すことである。そのために、一般点での
コホモロジーを計算するが、これはある種のガロアコホモロジーの計算に帰着する。





\section{Brauer 群}
\label{etale-cohomology-section-brauer-groups}

\noindent
体の Brauer 群は有限次元中心単純代数を用いて定義される。この節では Brauer 群に
関係する事実を復習する。その大部分は Brauer 群、節
\ref{brauer-section-introduction} で論じられている。他の文献としては
\cite{SerreCorpsLocaux}、\cite{SerreGaloisCohomology}、または \cite{Weil} を参照されたい。

\begin{theorem}
\label{etale-cohomology-theorem-central-simple-algebra}
$K$ を体とする。単位元をもち結合的な（必ずしも可換でない）$K$-代数 $A$ に
ついて、次の条件は同値である。
\begin{enumerate}
\item $A$ は有限次元中心単純 $K$-代数である。
\item $A$ は有限次元 $K$-ベクトル空間であり、$K$ は $A$ の中心であり、
$A$ は非自明な両側イデアルをもたない。
\item $d \geq 1$ が存在して
$A \otimes_K \bar K \cong \text{Mat}(d \times d, \bar K)$,
が成り立つ。
\item $d \geq 1$ が存在して
$A \otimes_K K^{sep} \cong \text{Mat}(d \times d, K^{sep})$,
が成り立つ。
\item $d \geq 1$ と有限ガロア拡大 $K'/K$ が存在して
$A \otimes_K K' \cong \text{Mat}(d \times d, K')$,
が成り立つ。
\item $n \geq 1$ と、有限次元中心斜体 $D$ であって $K$ 上のものが存在して
$A \cong \text{Mat}(n \times n, D)$ が成り立つ。
\end{enumerate}
整数 $d$ を $A$ の{\it 次数}という。
\end{theorem}

\begin{proof}
これは Brauer 群、補題 \ref{brauer-lemma-finite-central-simple-algebra} の再掲である。
\end{proof}

\begin{lemma}
\label{etale-cohomology-lemma-brauer-inverse}
$A$ を $K$ 上の有限次元中心単純代数とする。このとき
$$
\begin{matrix}
A \otimes_K A^{opp} & \longrightarrow & \text{End}_K(A) \\
\ a \otimes a' & \longmapsto & (x \mapsto a x a')
\end{matrix}
$$
は $K$ 上の代数の同型である。
\end{lemma}

\begin{proof}
Brauer 群、補題 \ref{brauer-lemma-inverse} を参照されたい。
\end{proof}

\begin{definition}
\label{etale-cohomology-definition-brauer-equivalent}
二つの有限次元中心単純代数 $A_1$ と $A_2$ が $K$ 上で{\it 類似}、または
{\it 同値}であるとは、$m, n \geq 1$ が存在して
$\text{Mat}(n \times n, A_1) \cong \text{Mat}(m \times m, A_2)$
となることをいう。$A_1 \sim A_2$ と書く。
\end{definition}

\noindent
Brauer 群、補題 \ref{brauer-lemma-similar} により、これは同値関係である。

\begin{definition}
\label{etale-cohomology-definition-brauer-group}
$K$ を体とする。$K$ の {\it Brauer 群}とは、集合 $\text{Br} (K)$、すなわち $K$ 上の
有限次元中心単純代数の類似類全体に、$K$ 上のテンソル積が誘導する群法を
入れたものである。$A$ の $\text{Br}(K)$ における類を $[A]$ と書く。単位元は、
$[K] = [\text{Mat}(d \times d, K)]$ であり、これは任意の $d \geq 1$ に対して成り立つ。
\end{definition}

\noindent
前の補題から逆元の存在と $-[A] = [A^{opp}]$ が従う。体の Brauer 群は常に
捩れ群である。実際、$[A]$ の位数が $\deg(A)$ を割り切ることを、任意の有限次元
中心単純代数 $A$ について見る（補題 \ref{etale-cohomology-lemma-annihilated-by-degree} を参照）。
一般に Brauer 群は有限生成でない。たとえば、非アルキメデス局所体の Brauer 群は
$\mathbf{Q}/\mathbf{Z}$ である。$\mathbf{C}(x, y)$ の Brauer 群は非可算である。

\begin{lemma}
\label{etale-cohomology-lemma-central-simple-algebra-pgln}
\begin{slogan}
中心単純代数は射影一般線形群のガロアコホモロジーによって分類される。
\end{slogan}
$K$ を体とし、$K^{sep}$ を分離代数閉包とする。このとき、次数 $d$ の中心単純代数で
$K$ 上のものの同型類全体は、非可換コホモロジー
$H^1(\text{Gal}(K^{sep}/K), \text{PGL}_d(K^{sep}))$
と全単射に対応する。
\end{lemma}

\begin{proof}[証明の概略]
スコーレム--ネーターの定理（Brauer 群、定理
\ref{brauer-theorem-skolem-noether} を参照）から、任意の体 $L$ に対して群
$\text{Aut}_{L\text{-代数}}(\text{Mat}_d(L))$ は
$\text{PGL}_d(L)$ に等しいことが従う。定理
\ref{etale-cohomology-theorem-central-simple-algebra} により、次数 $d$ の中心単純代数は
$K$-代数 $\text{Mat}_d(K)$ の形式に対応する。以上を組み合わせると、
次数 $d$ の中心単純代数の同型類が
$H^1(\text{Gal}(K^{sep}/K), \text{PGL}_d(K^{sep}))$
の元に対応することがわかる。捩りの詳細については、たとえば
\cite{SilvermanEllipticCurves} を参照されたい。
\end{proof}

\noindent
$A$ が次数 $d$ の有限次元中心単純代数で、体 $K$ 上のものならば、対応する
コホモロジー類を $\xi_A$ と書く。これは
$H^1(\text{Gal}(K^{sep}/K), \text{PGL}_d(K^{sep}))$
に属する。短完全列
$$
1 \to (K^{sep})^* \to \text{GL}_d(K^{sep}) \to \text{PGL}_d(K^{sep}) \to 1,
$$
を考える。これは次数 2 までのコホモロジー長完全列を与え、その余境界写像は
$$
\delta_d :
H ^1(\text{Gal}(K^{sep}/K), \text{PGL}_d(K^{sep}))
\longrightarrow
H^2(\text{Gal}(K^{sep}/K), (K^{sep})^*).
$$
である。具体的には次のように与えられる。$\xi$ が 1-コサイクル
$(g_\sigma)$ により表されるコホモロジー類ならば、$\delta_d(\xi)$ は 2-コサイクル
\begin{equation}
\label{etale-cohomology-equation-two-cocycle}
(\sigma, \tau)
\longmapsto
\tilde g_\sigma^{-1} \tilde g_{\sigma \tau} \sigma(\tilde g_\tau^{-1})
\in (K^{sep})^*
\end{equation}
の類である。ここで $\tilde g_\sigma \in \text{GL}_d(K^{sep})$ は $g_\sigma$ の持ち上げである。
これを用いると、写像
$$
\delta : \text{Br}(K) \longrightarrow H^2(\text{Gal}(K^{sep}/K), (K^{sep})^*),
\quad
[A] \longmapsto \delta_{\deg A} (\xi_A)
$$
を次のように具体化できる。$A$ が次数 $d$ をもち、$K$ 上であると仮定する。同型
$\varphi : \text{Mat}_d(K^{sep}) \to A \otimes_K K^{sep}$ を選ぶ。
$\sigma \in \text{Gal}(K^{sep}/K)$ に対し、元
$\tilde g_\sigma \in \text{GL}_d(K^{sep})$ を、
$\varphi^{-1} \circ \sigma(\varphi)$ が写像
$x \mapsto \tilde g_\sigma x \tilde g_\sigma^{-1}$ に等しくなるように選ぶ。
$H^2$ の類は 2-コサイクル (\ref{etale-cohomology-equation-two-cocycle}) によって定義される。

\begin{theorem}
\label{etale-cohomology-theorem-brauer-delta}
$K$ を分離代数閉包 $K^{sep}$ をもつ体とする。上で定義した写像
$\delta : \text{Br}(K) \to H^2(\text{Gal}(K^{sep}/K), (K^{sep})^*)$
は群の同型である。
\end{theorem}

\begin{proof}[証明の概略]
$\delta$ が群準同型を定めること、すなわち
$\delta(A \otimes_K B) = \delta(A) + \delta(B)$ を示すには、
コサイクルを用いて直接計算すればよい。

\medskip\noindent
$\delta$ の単射性。可換の場合（$d = 1$）には同一視
$$
H^1(\text{Gal}(K^{sep}/K), \text{GL}_d(K^{sep})) =
H_\etale^1(\Spec(K), \text{GL}_d(\mathcal{O}))
$$
があり、後者は fpqc 降下により自明である。これが非可換の場合にも成り立つならば、
$\delta$ の単射性が直ちに従う（\cite{SGA4.5} を参照）。実際に証明するには、
$\delta([A])$ を、$K$-ベクトル空間 $V$ で左 $A$-加群構造をもち
$\dim_K V = \deg A$ を満たすものの存在に対する障害として読み替えることができる。
$V$ が存在する場合には $A \cong \text{End}_K(V)$ である。

\medskip\noindent
全射性については、コホモロジー類
$\xi \in H^2(\text{Gal}(K^{sep}/K), (K^{sep})^*)$ を選ぶ。このとき有限ガロア拡大
$K^{sep}/K'/K$ が存在し、$\xi$ はある
$\xi' \in H^2(\text{Gal}(K'|K), (K')^*)$ の像である。そこで $K$ 上の中心単純代数を、
データ $K', \xi'$ を用いて具体的に書き下す。
\end{proof}










\section{スキームの Brauer 群}
\label{etale-cohomology-section-brauer-scheme}

\noindent
$S$ をスキームとする。$\mathcal{O}_S$-代数
$\mathcal{A}$ が {\it Azumaya} であるとは、エタール局所的に行列代数であること、
すなわち、あるエタール被覆
$\mathcal{U} = \{ \varphi_i : U_i \to S\}_{i \in I}$ が存在して
$\varphi_i^*\mathcal{A} \cong \text{Mat}_{d_i}(\mathcal{O}_{U_i})$
がある $d_i \geq 1$ に対して成り立つことをいう。このような
$\mathcal{A}$ と $\mathcal{B}$ が {\it 同値}であるとは、有限局所自由
$\mathcal{O}_S$-加群 $\mathcal{F}$ と $\mathcal{G}$ で、すべての $s \in S$ における
階数が正であり、かつ
$$
\mathcal{A} \otimes_{\mathcal{O}_S}
\SheafHom_{\mathcal{O}_S}(\mathcal{F}, \mathcal{F})
\cong
\mathcal{B} \otimes_{\mathcal{O}_S}
\SheafHom_{\mathcal{O}_S}(\mathcal{G}, \mathcal{G})
$$
が $\mathcal{O}_S$-代数として同型となるものが存在することをいう。$S$ の
{\it Brauer 群}とは、集合 $\text{Br}(S)$、すなわち Azumaya
$\mathcal{O}_S$-代数の同値類全体に、テンソル積（$\mathcal{O}_S$ 上）が誘導する
演算を入れたものである。

\begin{lemma}
\label{etale-cohomology-lemma-end-unique-up-to-invertible}
$S$ をスキームとする。$\mathcal{F}$ と $\mathcal{G}$ を、正の階数をもつ有限局所自由な
$\mathcal{O}_S$-加群の層とする。同型
$\SheafHom_{\mathcal{O}_S}(\mathcal{F}, \mathcal{F}) \cong
\SheafHom_{\mathcal{O}_S}(\mathcal{G}, \mathcal{G})$
が $\mathcal{O}_S$-代数の同型として存在するならば、可逆層
$\mathcal{L}$ が $S$ 上に存在して
$\mathcal{F} \otimes_{\mathcal{O}_S} \mathcal{L} \cong \mathcal{G}$
となり、しかもこの同型は、与えられた自己準同型代数の同型を誘導する。
\end{lemma}

\begin{proof}
同型
$\SheafHom_{\mathcal{O}_S}(\mathcal{F}, \mathcal{F}) \to
\SheafHom_{\mathcal{O}_S}(\mathcal{G}, \mathcal{G})$
を固定する。層 $\mathcal{L} \subset \SheafHom(\mathcal{F}, \mathcal{G})$ を、
$\mathcal{O}_S$-加群として、局所同型 $\varphi : \mathcal{F} \to \mathcal{G}$ であって、
$\varphi$ による共役が与えられた自己準同型代数の同型となるものから生成される層とする。
局所計算（$\mathcal{F}$ と $\mathcal{G}$ が有限自由で $S$ がアフィンである場合に帰着する）
により、$\mathcal{L}$ は可逆である。別の局所計算により、評価写像
$$
\mathcal{F} \otimes_{\mathcal{O}_S} \mathcal{L} \longrightarrow \mathcal{G}
$$
は同型である。
\end{proof}

\noindent
次の補題の証明で与える議論は \cite{Saltman-torsion} に見いだせる。

\begin{lemma}
\label{etale-cohomology-lemma-annihilated-by-degree}
\begin{reference}
議論は \cite{Saltman-torsion} から取った。
\end{reference}
$S$ をスキームとする。$\mathcal{A}$ を、階数 $d^2$ の局所自由加群であって $S$ 上の
Azumaya 代数とする。このとき、$\mathcal{A}$ の類は $S$ の Brauer 群において
$d$ で消去される。
\end{lemma}

\begin{proof}
エタール被覆 $\{U_i \to S\}$ と同型
$\mathcal{A}|_{U_i} \to \SheafHom(\mathcal{F}_i, \mathcal{F}_i)$
を、ある局所自由 $\mathcal{O}_{U_i}$-加群 $\mathcal{F}_i$ に対して選ぶ。その階数は
$d$ である。（$\mathcal{F}_i$ は自由であると仮定してよい。）合成
$$
p_i : \mathcal{F}_i^{\otimes d} \to
\wedge^d(\mathcal{F}_i) \to \mathcal{F}_i^{\otimes d}
$$
を考える。第一の射は通常の射影であり、第二の射は、$\mathcal{F}_i$ の最高外積を、
$\mathcal{F}_i^{\otimes d}$ の切断のうち $d$ 文字上の対称群の作用のもとで符号指標に
従って変換するものからなる部分加群と同一視する同型である。このとき
$p_i^2 = d! p_i$ であり、$p_i$ の階数は $1$ である。与えられた同型
$\mathcal{A}|_{U_i} \to \SheafHom(\mathcal{F}_i, \mathcal{F}_i)$
と標準同型
$$
\SheafHom(\mathcal{F}_i, \mathcal{F}_i)^{\otimes d} =
\SheafHom(\mathcal{F}_i^{\otimes d}, \mathcal{F}_i^{\otimes d})
$$
を用い、$p_i$ を $\mathcal{A}^{\otimes d}$ の $U_i$ 上の切断とみなせる。
$p_i|_{U_i \times_S U_j} = p_j|_{U_i \times_S U_j}$ が
$\mathcal{A}^{\otimes d}$ の切断として成り立つと主張する。実際、補題
\ref{etale-cohomology-lemma-end-unique-up-to-invertible} を適用すると、可逆層 $\mathcal{L}_{ij}$ と標準同型
$$
\mathcal{F}_i|_{U_i \times_S U_j} \otimes \mathcal{L}_{ij}
\longrightarrow
\mathcal{F}_j|_{U_i \times_S U_j}.
$$
を得る。この同型により $p_i$ は $p_j$ に写ることがわかる。
$\mathcal{A}^{\otimes d}$ は $S_\etale$ 上の層なので（命題
\ref{etale-cohomology-proposition-quasi-coherent-sheaf-fpqc}）、標準的な大域切断
$p \in \Gamma(S, \mathcal{A}^{\otimes d})$ を得る。局所計算により
$$
\mathcal{H} =
\Im(\mathcal{A}^{\otimes d} \to \mathcal{A}^{\otimes d}, f \mapsto fp)
$$
は階数 $d^d$ の局所自由加群であり、$\mathcal{A}^{\otimes d}$ による左乗法は同型
$\mathcal{A}^{\otimes d} \to \SheafHom(\mathcal{H}, \mathcal{H})$ を誘導する。
言い換えれば、$\mathcal{A}^{\otimes d}$ は $S$ の Brauer 群の自明な元であり、望む結論を得る。
\end{proof}

\noindent
この状況では、定理 \ref{etale-cohomology-theorem-brauer-delta} の同型 $\delta$ の類似物は写像
$$
\delta_S: \text{Br}(S) \to H_\etale^2(S, \mathbf{G}_m).
$$
$\delta_S$ は単射である。$S$ が準コンパクトまたは連結ならば、$\text{Br}(S)$ は捩れ群である。
したがって、この場合 $\delta_S$ の像は $S$ の {\it コホモロジー的 Brauer 群}
$$
\text{Br}'(S) := H_\etale^2(S, \mathbf{G}_m)_\text{torsion}.
$$
に含まれる。ゆえに $S$ が準コンパクトまたは連結ならば、包含 $\text{Br}(S)
\subset \text{Br}'(S)$ がある。これは常に等号であるとは限らない。実際、非分離特異曲面
$S$ であって、$\text{Br}(S) \subset
\text{Br}'(S)$ が真の包含となるものが存在する。$S$ が準射影的ならば
$\text{Br}(S) = \text{Br}'(S)$ である。しかし、たとえば $\mathbf{C}$ 上の滑らかな固有多様体で
これが成り立つかどうかは知られていない。




\section{Artin--Schreier 列}
\label{etale-cohomology-section-artin-schreier}

\noindent
$p$ を素数とする。$S$ を標数 $p$ のスキームとする。{\it Artin--Schreier} 列とは短完全列
$$
0 \longrightarrow \underline{\mathbf{Z}/p\mathbf{Z}}_S \longrightarrow
\mathbf{G}_{a, S} \xrightarrow{F-1} \mathbf{G}_{a, S} \longrightarrow 0
$$
のことである。ここで $F - 1$ は写像 $x \mapsto x^p - x$ である。

\begin{lemma}
\label{etale-cohomology-lemma-vanishing-affine-char-p-p}
$p$ を素数とする。$S$ を標数 $p$ のスキームとする。
\begin{enumerate}
\item $S$ がアフィンならば、$H_\etale^q(S, \underline{\mathbf{Z}/p\mathbf{Z}}) = 0$ が
すべての $q \geq 2$ に対して成り立つ。
\item $S$ が次元 $d$ の準コンパクトかつ準分離的なスキームならば、すべての
$H_\etale^q(S, \underline{\mathbf{Z}/p\mathbf{Z}}) = 0$ が、$q \geq 2 + d$ に対して
成り立つ。
\end{enumerate}
\end{lemma}

\begin{proof}
構造層のエタールコホモロジーは、台となる位相空間上のコホモロジーに等しいことを
思い出そう（定理 \ref{etale-cohomology-theorem-zariski-fpqc-quasi-coherent}）。第一の主張は
Artin--Schreier 完全列と、アフィンスキーム上の構造層のコホモロジーの消滅から従う
（スキームのコホモロジー、補題
\ref{coherent-lemma-quasi-coherent-affine-cohomology-zero}）。第二の主張も同じ議論により、
コホモロジー、命題
\ref{cohomology-proposition-cohomological-dimension-spectral} の消滅と、$S$ がスペクトル空間で
あることから従う（性質、補題
\ref{properties-lemma-quasi-compact-quasi-separated-spectral}）。
\end{proof}

\begin{lemma}
\label{etale-cohomology-lemma-F-1}
$k$ を標数 $p > 0$ の代数閉体とする。$V$ を有限次元 $k$-ベクトル空間とする。
$F : V \to V$ を Frobenius 線形写像、すなわち、$F(\lambda v) = \lambda^p F(v)$ が
すべての $\lambda \in k$ と $v \in V$ に対して成り立つ加法的写像とする。このとき
$F - 1 : V \to V$ は全射であり、その核は有限次元 $\mathbf{F}_p$-ベクトル空間で、
その次元は $\leq \dim_k(V)$ である。
\end{lemma}

\begin{proof}
$F = 0$ ならば主張は成り立つ。$V$ に $F$-安定な部分ベクトル空間によるフィルトレーションが
あり、各次数付き部分について主張が成り立つならば、$(V, F)$ についても成り立つ。
この二つの注意を組み合わせ、$F$ の核は零であると仮定してよい。

\medskip\noindent
$v_1, \ldots, v_n$ を $V$ の基底として選び、$F(v_i) = \sum a_{ij} v_j$ と書く。
$v = \sum \lambda_i v_i$ が核に属するための必要十分条件は
$\sum \lambda_i^p a_{ij} v_j = 0$ である。$k$ は代数閉なので、このことから行列
$(a_{ij})$ は可逆である。その逆行列を $(b_{ij})$ とする。$F - 1$ が全射であることを見るため、
$w = \sum \mu_i v_i \in V$ を取り、方程式
$$
(F - 1)(\sum \lambda_iv_i) =
\sum \lambda_i^p a_{ij} v_j - \sum \lambda_j v_j = \sum \mu_j v_j
$$
を解こうとする。これは
$$
\sum \lambda_j^p v_j - \sum b_{ij} \lambda_i v_j = \sum b_{ij} \mu_i v_j
$$
と同値であり、言い換えれば
$$
\lambda_j^p - \sum b_{ij} \lambda_i = \sum b_{ij} \mu_i,
\quad j = 1, \ldots, \dim(V).
$$
である。代数
$$
A = k[x_1, \ldots, x_n]/
(x_j^p - \sum b_{ij} x_i - \sum b_{ij} \mu_i)
$$
は $k$ 上標準滑らかである（代数、定義 \ref{algebra-definition-standard-smooth}）。実際、行列
$(b_{ij})$ は可逆で、$x_j^p$ の偏微分は零である。$A$ の $k$ 上の基底は、
単項式 $x_1^{e_1} \ldots x_n^{e_n}$ で $e_i < p$ を満たすもの全体なので、
$\dim_k(A) = p^n$ である。$k$ は代数閉だから、$\Spec(A)$ はちょうど $p^n$ 個の点をもつ。
したがって $F - 1$ は全射であり、各ファイバーは $p^n$ 個の点をもつ。すなわち
$F - 1$ の核は $p^n$ 個の元からなる群である。
\end{proof}

\begin{lemma}
\label{etale-cohomology-lemma-F-2}
補題 \ref{etale-cohomology-lemma-F-1} の状況で、$k'/k$ を代数閉体の拡大とする。
$V' = V \otimes_k k'$ とおき、$F' : V' \to V'$ を $F$ が誘導する Frobenius 線形写像とする。
このとき $\Ker(F - 1)$ は $\Ker(F' - 1)$ に同型に写る。
\end{lemma}

\begin{proof}
$v_1, \ldots, v_n$ を $V$ の基底とし、$F(v_i) = \sum a_{ij}v_j$ とする。
$v = \sum \lambda_i v_i$ が $\Ker(F - 1)$ に属するための必要十分条件は、
$\sum \lambda_i^pa_{ij} = \lambda_j$ が $j = 1, \ldots, n$ に対して成り立つことである。
補題 \ref{etale-cohomology-lemma-F-1} により、これらの方程式は閉部分スキーム $Z$ を
$\mathbf{A}^n_k = \Spec(k[\lambda_1, \ldots, \lambda_n])$
の中に切り出し、この部分スキームは $\Spec(k)$ 上有限である。すると $Z(k) = Z(k')$ であり、
結論を得る。
\end{proof}

\begin{lemma}
\label{etale-cohomology-lemma-top-cohomology-coherent}
$X$ を体 $k$ 上分離的有限型なスキームとする。$\mathcal{F}$ を $\mathcal{O}_X$-加群の
連接層とする。このとき、$\dim_k H^d(X, \mathcal{F}) < \infty$ である。ここで
$d = \dim(X)$ とおいた。
\end{lemma}

\begin{proof}
$d$ に関する帰納法で証明する。$d = 0$ の場合、$X$ は有限次元 $k$-代数 $A$ のスペクトルであり
（多様体、補題 \ref{varieties-lemma-algebraic-scheme-dim-0}）、各連接層 $\mathcal{F}$ は有限
$A$-加群 $M = H^0(X, \mathcal{F})$ で $\dim_k M < \infty$ を満たすものに対応するので、
主張は成り立つ。

\medskip\noindent
$d > 0$ とし、次元 $< d$ の分離的有限型スキームについて結果が示されていると仮定する。
スキーム $X$ は Noetherian である。連接層の性質 $\mathcal{P}$ を $X$ 上で
$$
\mathcal{P}(\mathcal{F}) \Leftrightarrow
\dim_k H^d(X, \mathcal{F}) < \infty
$$
で定義する。スキームのコホモロジー、補題 \ref{coherent-lemma-property-initial} の結果を用いて、
$\mathcal{P}$ が $X$ 上のすべての連接層について成り立つことを証明する。

\medskip\noindent
連接層の短完全列
$$
0 \to \mathcal{F}_1 \to \mathcal{F} \to \mathcal{F}_2 \to 0
$$
を $X$ 上で取る。コホモロジー長完全列
$$
H^d(X, \mathcal{F}_1) \to
H^d(X, \mathcal{F}) \to
H^d(X, \mathcal{F}_2)
$$
を考える。したがって、$\mathcal{P}$ が $\mathcal{F}_1$ と $\mathcal{F}_2$ について成り立つならば、
$\mathcal{F}$ についても成り立つ。

\medskip\noindent
$Z \subset X$ を整閉部分スキームとする。$\mathcal{I}$ を $Z$ 上のイデアルの連接層とする。
証明を終えるには、$H^d(X, i_*\mathcal{I}) = H^d(Z, \mathcal{I})$ が有限次元であることを
示せばよい。$\dim(Z) < d$ ならば、このコホモロジー群は零なので結果は成り立つ
（コホモロジー、命題 \ref{cohomology-proposition-vanishing-Noetherian}）。
こうして次の段落で論じる状況に帰着する。

\medskip\noindent
$X$ は次元 $d$ の多様体であり、$\mathcal{F} = \mathcal{I}$ はイデアルの連接層であると
仮定する。この場合、短完全列
$$
0 \to \mathcal{I} \to \mathcal{O}_X \to i_*\mathcal{O}_Z \to 0
$$
がある。ここで $i : Z \to X$ は $\mathcal{I}$ が定める閉部分スキームである。
帰納法の仮定により、$H^{d - 1}(Z, \mathcal{O}_Z) = H^{d - 1}(X, i_*\mathcal{O}_Z)$ は
有限次元である。したがって、構造層について結果を証明すれば十分である。

\medskip\noindent
射 $X \to \Spec(k)$ に Chow の補題（スキームのコホモロジー、補題
\ref{coherent-lemma-chow-Noetherian}）を適用できる。これにより図式
$$
\xymatrix{
X \ar[rd]_g & X' \ar[d]^-{g'} \ar[l]^\pi \ar[r]_i & \mathbf{P}^n_k \ar[dl] \\
& \Spec(k) &
}
$$
を得る。これは Chow の補題の主張にある図式である。また $U \subset X$ を、
$\pi^{-1}(U) \to U$ が同型となるような稠密開部分スキームとする。
$X'$ も多様体であると仮定してよい（スキームのコホモロジー、注意
\ref{coherent-remark-chow-Noetherian}）。射 $i' = (i, \pi) : X' \to \mathbf{P}^n_X$ は
閉埋入である（同所）。したがって
$$
\mathcal{L} = i^*\mathcal{O}_{\mathbf{P}^n_k}(1) \cong
(i')^*\mathcal{O}_{\mathbf{P}^n_X}(1)
$$
は $\pi$-相対的豊富である（たとえば射、補題
\ref{morphisms-lemma-characterize-ample-on-finite-type} による）。したがって、スキームの
コホモロジー、補題 \ref{coherent-lemma-kill-by-twisting} により、ある $n \geq 0$ が存在して、
$R^p\pi_*\mathcal{L}^{\otimes n} = 0$ がすべての $p > 0$ に対して成り立つ。
$\mathcal{G} = \pi_*\mathcal{L}^{\otimes n}$ とおく。零でない大域切断 $s$ を、
$\mathcal{L}^{\otimes n}$ の切断として任意に選ぶ。$\mathcal{G} = \pi_*\mathcal{L}^{\otimes n}$ だから、切断 $s$ は
$\mathcal{G}$ の切断、すなわち写像 $\mathcal{O}_X \to \mathcal{G}$ に対応する。
$s|_U \not = 0$ である。実際、$X'$ は多様体で $\mathcal{L}$ は可逆であり、したがって
$\mathcal{O}_X|_U \to \mathcal{G}|_U$ は零でない。また
$\mathcal{G}|_U = \mathcal{L}^{\otimes n}|_{\pi^{-1}(U)}$ は可逆なので、短完全列
$$
0 \to \mathcal{O}_X \to \mathcal{G} \to \mathcal{Q} \to 0
$$
を得る。ここで $\mathcal{Q}$ は連接で、$X$ の真の閉部分スキーム上に台をもつ。
前と同様に帰納法の仮定を用いて議論すれば、$\dim H^d(X, \mathcal{G}) < \infty$ を
示せば十分である。

\medskip\noindent
Leray スペクトル系列（コホモロジー、補題 \ref{cohomology-lemma-apply-Leray}）により
$H^d(X, \mathcal{G}) = H^d(X', \mathcal{L}^{\otimes n})$ である。
$\overline{X}' \subset \mathbf{P}^n_k$ を $X'$ の閉包とする。このとき $\overline{X}'$ は
次元 $d$ の射影多様体で、$k$ 上のものであり、$X' \subset \overline{X}'$ は稠密開部分である。
可逆層 $\mathcal{L}^{\otimes n}$ は $\mathcal{O}_{\overline{X}'}(n)$ の $X'$ への制限である。
コホモロジー、命題 \ref{cohomology-proposition-cohomological-dimension-spectral} により、写像
$$
H^d(\overline{X}', \mathcal{O}_{\overline{X}'}(n))
\longrightarrow
H^d(X', \mathcal{L}^{\otimes n})
$$
は全射である。左側のコホモロジー群はスキームのコホモロジー、補題
\ref{coherent-lemma-coherent-projective} により有限次元だから、証明は完了する。
\end{proof}

\begin{lemma}
\label{etale-cohomology-lemma-vanishing-variety-char-p-p}
$X$ を代数閉体 $k$ で標数 $p > 0$ のものの上で分離的有限型なスキームとする。このとき
$H_\etale^q(X, \underline{\mathbf{Z}/p\mathbf{Z}}) = 0$ が
$q \geq dim(X) + 1$ に対して成り立つ。
\end{lemma}

\begin{proof}
$d = \dim(X)$ とおく。補題 \ref{etale-cohomology-lemma-vanishing-affine-char-p-p} で示した消滅により、
$H_\etale^{d + 1}(X, \underline{\mathbf{Z}/p\mathbf{Z}}) = 0$ を示せば十分である。
補題 \ref{etale-cohomology-lemma-top-cohomology-coherent} により、$H^d(X, \mathcal{O}_X)$ は有限次元
$k$-ベクトル空間である。したがって、Artin--Schreier 列に付随するコホモロジー長完全列は
$$
H^d(X, \mathcal{O}_X) \xrightarrow{F - 1}
H^d(X, \mathcal{O}_X) \to H^{d + 1}_\etale(X, \mathbf{Z}/p\mathbf{Z}) \to 0
$$
で終わる。補題 \ref{etale-cohomology-lemma-F-1} により、この列の写像 $F - 1$ は全射である。
これで補題が証明された。
\end{proof}

\begin{lemma}
\label{etale-cohomology-lemma-finiteness-proper-variety-char-p-p}
$X$ を代数閉体 $k$ で標数 $p > 0$ のものの上の固有スキームとする。このとき
\begin{enumerate}
\item $H_\etale^q(X, \underline{\mathbf{Z}/p\mathbf{Z}})$ は有限
$\mathbf{Z}/p\mathbf{Z}$-加群であり、これはすべての $q$ に対して成り立つ。また
\item 写像 $H^q_\etale(X, \underline{\mathbf{Z}/p\mathbf{Z}}) \to
H^q_\etale(X_{k'}, \underline{\mathbf{Z}/p\mathbf{Z}})$
は、$k'/k$ が代数閉体の拡大ならば同型である。
\end{enumerate}
\end{lemma}

\begin{proof}
スキームのコホモロジー、補題
\ref{coherent-lemma-proper-over-affine-cohomology-finite} と定理
\ref{etale-cohomology-theorem-zariski-fpqc-quasi-coherent} のコホモロジー比較により、コホモロジー群
$H^q_\etale(X, \mathbf{G}_a) = H^q(X, \mathcal{O}_X)$ は有限次元 $k$-ベクトル空間である。
したがって補題 \ref{etale-cohomology-lemma-F-1} により、Artin--Schreier 列に付随するコホモロジー長完全列は
短完全列
$$
0 \to H_\etale^q(X, \underline{\mathbf{Z}/p\mathbf{Z}}) \to
H^q(X, \mathcal{O}_X) \xrightarrow{F - 1} H^q(X, \mathcal{O}_X) \to 0
$$
に分裂する。さらに、$\mathbf{F}_p$-次元で測ったコホモロジー群
$H_\etale^q(X, \underline{\mathbf{Z}/p\mathbf{Z}})$ の次元は、$k$-次元で測ったベクトル空間
$H^q(X, \mathcal{O}_X)$ の次元以下である。
これで第一の主張が証明された。第二の主張は補題 \ref{etale-cohomology-lemma-F-2} から従う。実際、平坦基底変換により
$H^q(X, \mathcal{O}_X) \otimes_k k' \to H^q(X_{k'}, \mathcal{O}_{X_{k'}})$
は同型である（スキームのコホモロジー、補題
\ref{coherent-lemma-flat-base-change-cohomology}）。
\end{proof}








\section{局所定数層}
\label{etale-cohomology-section-locally-constant}

\noindent
この節は、サイト上の加群、節 \ref{sites-modules-section-locally-constant} のエタールサイトに
対する類似物である。

\begin{definition}
\label{etale-cohomology-definition-finite-locally-constant}
$X$ をスキームとする。$\mathcal{F}$ を $X_\etale$ 上の集合の層とする。
\begin{enumerate}
\item $E$ を集合とする。$\mathcal{F}$ が{\it 値 $E$ の定数層}であるとは、
$\mathcal{F}$ が前層 $U \mapsto E$ の層化であることをいう。記法：$\underline{E}_X$ または
$\underline{E}$。
\item $\mathcal{F}$ が{\it 定数層}であるとは、(1) の層と同型であることをいう。
\item $\mathcal{F}$ が{\it 局所定数}であるとは、被覆 $\{U_i \to X\}$ が存在して
$\mathcal{F}|_{U_i}$ が定数層となることをいう。
\item $\mathcal{F}$ が{\it 有限局所定数}であるとは、局所定数であって、その定数値が
有限集合であることをいう。
\end{enumerate}
$\mathcal{F}$ を $X_\etale$ 上のアーベル群の層とする。
\begin{enumerate}
\item $A$ をアーベル群とする。$\mathcal{F}$ が{\it 値 $A$ の定数層}であるとは、
$\mathcal{F}$ が前層 $U \mapsto A$ の層化であることをいう。記法：$\underline{A}_X$ または
$\underline{A}$。
\item $\mathcal{F}$ が{\it 定数層}であるとは、アーベル層として (1) の層と同型であることをいう。
\item $\mathcal{F}$ が{\it 局所定数}であるとは、被覆 $\{U_i \to X\}$ が存在して
$\mathcal{F}|_{U_i}$ が定数層となることをいう。
\item $\mathcal{F}$ が{\it 有限局所定数}であるとは、局所定数であって、その定数値が
有限アーベル群であることをいう。
\end{enumerate}
$\Lambda$ を環とする。$\mathcal{F}$ を $\Lambda$-加群の層で $X_\etale$ 上のものとする。
\begin{enumerate}
\item $M$ を $\Lambda$-加群とする。$\mathcal{F}$ が{\it 値 $M$ の定数層}であるとは、
$\mathcal{F}$ が前層 $U \mapsto M$ の層化であることをいう。記法：$\underline{M}_X$ または
$\underline{M}$。
\item $\mathcal{F}$ が{\it 定数層}であるとは、$\Lambda$-加群の層として (1) の層と
同型であることをいう。
\item $\mathcal{F}$ が{\it 局所定数}であるとは、被覆 $\{U_i \to X\}$ が存在して
$\mathcal{F}|_{U_i}$ が定数層となることをいう。
\end{enumerate}
\end{definition}

\begin{lemma}
\label{etale-cohomology-lemma-pullback-locally-constant}
$f : X \to Y$ をスキームの射とする。$\mathcal{G}$ が集合、アーベル群、または $\Lambda$-加群の
局所定数層で $Y_\etale$ 上のものならば、$f^{-1}\mathcal{G}$ も $X_\etale$ 上で同じ性質をもつ。
\end{lemma}

\begin{proof}
任意のトポスの射について成り立つ。サイト上の加群、補題
\ref{sites-modules-lemma-pullback-locally-constant} を参照されたい。
\end{proof}

\begin{lemma}
\label{etale-cohomology-lemma-pushforward-locally-constant}
$f : X \to Y$ をスキームの有限エタール射とする。$\mathcal{F}$ が集合の（有限）局所定数層、
アーベル群の（有限）局所定数層、または $\Lambda$-加群の（有限型）局所定数層で
$X_\etale$ 上のものならば、$f_*\mathcal{F}$ も $Y_\etale$ 上で同じ性質をもつ。
\end{lemma}

\begin{proof}
$f_*$ の構成はエタール局所化と可換である。有限エタール射は局所的に同型射の非交和と
同型である。エタール射、補題 \ref{etale-lemma-finite-etale-etale-local} を参照されたい。
したがって補題の主張は、$\mathcal{F}_i$（$i = 1, \ldots, n$）が集合の（有限）局所定数層ならば、
$\prod_{i = 1, \ldots, n} \mathcal{F}_i$ もそうである、という主張に帰着する。
これは明らかである。アーベル群と加群の層についても同様である。
\end{proof}

\begin{lemma}
\label{etale-cohomology-lemma-characterize-finite-locally-constant}
$X$ をスキームとし、$\mathcal{F}$ を $X_\etale$ 上の集合の層とする。このとき次は同値である。
\begin{enumerate}
\item $\mathcal{F}$ は有限局所定数である。
\item $\mathcal{F} = h_U$ が、ある有限エタール射 $U \to X$ に対して成り立つ。
\end{enumerate}
\end{lemma}

\begin{proof}
有限エタール射は局所的に同型射の非交和と同型である。エタール射、補題
\ref{etale-lemma-finite-etale-etale-local} を参照されたい。したがって (2) は (1) を含意する。
逆に、$\mathcal{F}$ が有限局所定数ならば、エタール被覆 $\{X_i \to X\}$ が存在し、
$\mathcal{F}|_{X_i}$ は有限エタール射 $U_i \to X_i$ によって表現される。
降下、補題 \ref{descent-lemma-descent-data-sheaves} の証明とまったく同様に議論すると、
スキームの降下データ $(U_i, \varphi_{ij})$ を $\{X_i \to X\}$ に関して得る（詳細は省略する）。
この降下データは、たとえば降下、補題 \ref{descent-lemma-affine} により有効であり、
その結果得られるスキームの射 $U \to X$ は、降下、補題
\ref{descent-lemma-descending-property-finite} および
\ref{descent-lemma-descending-property-etale} により有限エタールである。
\end{proof}

\begin{lemma}
\label{etale-cohomology-lemma-morphism-locally-constant}
$X$ をスキームとする。
\begin{enumerate}
\item $\varphi : \mathcal{F} \to \mathcal{G}$ を $X_\etale$ 上の集合の局所定数層の射とする。
$\mathcal{F}$ が有限局所定数ならば、エタール被覆 $\{U_i \to X\}$ が存在して、
$\varphi|_{U_i}$ は集合の射に付随する定数層の射となる。
\item $\varphi : \mathcal{F} \to \mathcal{G}$ を $X_\etale$ 上のアーベル群の局所定数層の射とする。
$\mathcal{F}$ が有限局所定数ならば、エタール被覆 $\{U_i \to X\}$ が存在して、
$\varphi|_{U_i}$ はアーベル群の射に付随する定数アーベル層の射となる。
\item $\Lambda$ を環とする。$\varphi : \mathcal{F} \to \mathcal{G}$ を $\Lambda$-加群の
局所定数層で $X_\etale$ 上のものの射とする。$\mathcal{F}$ が有限型ならば、エタール被覆
$\{U_i \to X\}$ が存在して、$\varphi|_{U_i}$ は $\Lambda$-加群の射に付随する
$\Lambda$-加群の定数層の射となる。
\end{enumerate}
\end{lemma}

\begin{proof}
任意のサイト上で成り立つ。サイト上の加群、補題
\ref{sites-modules-lemma-morphism-locally-constant} を参照されたい。
\end{proof}

\begin{lemma}
\label{etale-cohomology-lemma-kernel-finite-locally-constant}
$X$ をスキームとする。
\begin{enumerate}
\item 集合の有限局所定数層の圏は、$\Sh(X_\etale)$ の中で有限極限と有限余極限について閉じている。
\item 有限局所定数アーベル層の圏は $\textit{Ab}(X_\etale)$ の弱 Serre 部分圏である。
\item $\Lambda$ を Noetherian 環とする。$\Lambda$-加群の有限型局所定数層で $X_\etale$ 上のものの圏は
$\textit{Mod}(X_\etale, \Lambda)$ の弱 Serre 部分圏である。
\end{enumerate}
\end{lemma}

\begin{proof}
任意のサイト上で成り立つ。サイト上の加群、補題
\ref{sites-modules-lemma-kernel-finite-locally-constant} を参照されたい。
\end{proof}

\begin{lemma}
\label{etale-cohomology-lemma-tensor-product-locally-constant}
$X$ をスキームとし、$\Lambda$ を環とする。二つの $\Lambda$-加群の局所定数層で
$X_\etale$ 上のもののテンソル積は、$\Lambda$-加群の局所定数層である。
\end{lemma}

\begin{proof}
任意のサイト上で成り立つ。サイト上の加群、補題
\ref{sites-modules-lemma-tensor-product-locally-constant} を参照されたい。
\end{proof}

\begin{lemma}
\label{etale-cohomology-lemma-connected-locally-constant}
$X$ を連結スキームとする。$\Lambda$ を環とし、$\mathcal{F}$ を $\Lambda$-加群の局所定数層とする。
このとき、$\Lambda$-加群 $M$ とエタール被覆 $\{U_i \to X\}$ が存在して、
$\mathcal{F}|_{U_i} \cong \underline{M}|_{U_i}$ となる。
\end{lemma}

\begin{proof}
エタール被覆 $\{U_i \to X\}$ を、$\mathcal{F}|_{U_i}$ が定数となるように選ぶ。たとえば
$\mathcal{F}|_{U_i} \cong \underline{M_i}_{U_i}$ と書ける。
$U_i \times_X U_j$ は、$M_i$ が $M_j$ と同型でなければ空であることに注意する。
各 $\Lambda$-加群 $M$ に対し $I_M = \{i \in I \mid M_i \cong M\}$ とおく。
エタール射は開写像なので、
$U_M = \bigcup_{i \in I_M} \Im(U_i \to X)$
は $X$ の開部分集合である。すると $X = \coprod U_M$ は $X$ の互いに素な開被覆である。
$X$ は連結だから、空でない $U_M$ は一つだけであり、補題が従う。
\end{proof}






\section{局所定数層と基本群}
\label{etale-cohomology-section-pione}

\noindent
いくつかの場合には、局所定数層をスキームの基本群と関係づけることができる。

\begin{lemma}
\label{etale-cohomology-lemma-locally-constant-on-connected}
$X$ を連結スキームとする。$\overline{x}$ を $X$ の幾何学的点とする。
\begin{enumerate}
\item 圏同値
$$
\left\{
\begin{matrix}
\text{有限局所定数な}\\
\text{集合の層で }X_\etale\text{ 上のもの}
\end{matrix}
\right\}
\longleftrightarrow
\left\{
\begin{matrix}
\text{有限 }\pi_1(X, \overline{x})\text{-集合}
\end{matrix}
\right\}
$$
がある。
\item 圏同値
$$
\left\{
\begin{matrix}
\text{有限局所定数な}\\
\text{アーベル群の層で }X_\etale\text{ 上のもの}
\end{matrix}
\right\}
\longleftrightarrow
\left\{
\begin{matrix}
\text{有限 }\pi_1(X, \overline{x})\text{-加群}
\end{matrix}
\right\}
$$
がある。
\item $\Lambda$ を有限環とする。圏同値
$$
\left\{
\begin{matrix}
\text{有限型局所定数な}\\
\Lambda\text{-加群の層で }X_\etale\text{ 上のもの}
\end{matrix}
\right\}
\longleftrightarrow
\left\{
\begin{matrix}
\text{有限 }\pi_1(X, \overline{x})\text{-加群で、}\\
\text{可換する }\Lambda\text{-加群構造を備えたもの}
\end{matrix}
\right\}
$$
がある。
\end{enumerate}
\end{lemma}

\begin{proof}
$\pi_1(X, \overline{x})$ は副有限位相群であることに注意する。基本群、定義
\ref{pione-definition-fundamental-group} を参照されたい。
左辺の圏は節 \ref{etale-cohomology-section-locally-constant} で定義されている。
右辺の圏で用いた記法は、集合については基本群、定義
\ref{pione-definition-G-set-continuous} から、アーベル群については定義
\ref{etale-cohomology-definition-G-module-continuous} から取った。これで記法の意味が明らかになる。

\medskip\noindent
主張 (1) は補題 \ref{etale-cohomology-lemma-characterize-finite-locally-constant} と
基本群、定理 \ref{pione-theorem-fundamental-group} から従う。
(2) と (3) は、台となる（層の）集合に追加の構造を入れることにより、(1) から直ちに従う。
たとえば、$X_\etale$ 上の有限局所定数アーベル群層を与えることは、有限局所定数な
集合の層 $\mathcal{F}$ と、通常の公理を満たす写像
$+ : \mathcal{F} \times \mathcal{F} \to \mathcal{F}$
を与えることに等しい。(1) の同値は積を積に写すので、$+$ を対応する有限
$\pi_1(X, \overline{x})$-集合上の加法に写す。
$\pi_1(X, \overline{x})$-加群を与えることは、両立するアーベル群構造をもつ
$\pi_1(X, \overline{x})$-集合を与えることに等しいから、(2) を得る。(3) もまったく同様に証明される。
\end{proof}

\begin{lemma}
\label{etale-cohomology-lemma-locally-constant-on-connected-geom-unibranch}
$X$ を既約で幾何学的に単枝なスキームとする。
$\overline{x}$ を $X$ の幾何学的点とする。
$\Lambda$ を環とする。圏同値
$$
\left\{
\begin{matrix}
\text{有限型局所定数な}\\
\Lambda\text{-加群の層で }X_\etale\text{ 上のもの}
\end{matrix}
\right\}
\longleftrightarrow
\left\{
\begin{matrix}
\text{有限生成 }\Lambda\text{-加群 }M\text{ で、連続な}\\
\pi_1(X, \overline{x})\text{-作用を備えたもの}
\end{matrix}
\right\}
$$
がある。
\end{lemma}

\begin{proof}
補題 \ref{etale-cohomology-lemma-locally-constant-on-connected} の証明はこの場合には使えない。
実際、有限生成 $\Lambda$-加群 $M$ の台集合は有限とは限らない。

\medskip\noindent
$\nu : X^\nu \to X$ を正規化射とする。射、補題
\ref{morphisms-lemma-normalization-geom-unibranch-univ-homeo}
により、これは普遍同相写像である。基本群、命題
\ref{pione-proposition-universal-homeomorphism} により、これは同型
$\pi_1(X^\nu, \overline{x}) \to \pi_1(X, \overline{x})$
を誘導し、定理 \ref{etale-cohomology-theorem-topological-invariance} により、有限型局所定数
$\Lambda$-加群について、$X_\etale$ 上のものと $X^\nu_\etale$ 上のものとの圏同値を得る。
こうして $X$ が整正規スキームである場合に帰着する。

\medskip\noindent
$X$ を整正規スキームと仮定する。$\eta \in X$ を生成点とする。
$\overline{\eta}$ を $\eta$ 上にある幾何学的点とする。
基本群、命題 \ref{pione-proposition-normal} により、連続全射
$$
\text{Gal}(\kappa(\eta)^{sep}/\kappa(\eta)) = \pi_1(\eta, \overline{\eta})
\longrightarrow
\pi_1(X, \overline{\eta})
$$
があり、その核は基本群、補題
\ref{pione-lemma-normal-pione-quotient-inertia} で記述されている。
$\mathcal{F}$ を有限型局所定数 $\Lambda$-加群の層であって $X_\etale$ 上のものとする。
$M = \mathcal{F}_{\overline{\eta}}$ を $\mathcal{F}$ の $\overline{\eta}$ における茎とする。
節 \ref{etale-cohomology-section-galois-action-stalks} により、
$\text{Gal}(\kappa(\eta)^{sep}/\kappa(\eta))$ の $M$ 上の連続作用を得る。
この作用が上の全射を経由することを示すのが目標である。
$\mathcal{F}$ は有限型なので、$M$ は有限生成 $\Lambda$-加群である。
$\mathcal{F}$ は局所定数なので、すべての $x \in X$ に対し、
$\mathcal{F}$ の $\Spec(\mathcal{O}_{X, x}^{sh})$ への制限は定数である。
したがって、基本群、補題
\ref{pione-lemma-normal-pione-quotient-inertia} の記法のもとで、
$\text{Gal}(K^{sep}/K_x^{sh})$ の $M$ 上の作用は自明である。
これにより望みどおりの分解を得る。

\medskip\noindent
逆に、有限生成 $\Lambda$-加群 $M$ と、その上の
$\pi_1(X, \overline{\eta})$ の連続作用が与えられたとする。
$\mathcal{F}$ を構成して、$M \cong \mathcal{F}_{\overline{\eta}}$ が
$\Lambda[\pi_1(X, \overline{\eta})]$-加群として成り立つようにする。
生成元 $m_1, \ldots, m_r \in M$ を選ぶ。
$\pi_1(X, \overline{\eta})$ の $M$ 上の作用は連続だから、各 $i$ に対して、開部分群
$N_i$ が副有限群 $\pi_1(X, \overline{\eta})$ に存在し、すべての
$\gamma \in H_i$ は $m_i$ を固定する。したがって、開部分群
$H = \bigcap_{i = 1, \ldots, r} H_i$ のすべての元は $M$ のすべての元を固定する。
$H$ を小さく取り直すことにより、$H$ は $\pi_1(X, \overline{\eta})$ の
開正規部分群であると仮定してよい。
$G = \pi_1(X, \overline{\eta})/H$ とおく。$f : Y \to X$ を対応する
Galois 有限エタール $G$-被覆とする。$f_*\underline{\mathbf{Z}}$ は
$\mathbf{Z}[G]$-加群の層であって $X_\etale$ 上のものとみなせる。そこで単に
$$
\mathcal{F} =
f_*\underline{\mathbf{Z}} \otimes_{\underline{\mathbf{Z}[G]}} \underline{M}
$$
と取ればよい。$\mathcal{F}_{\overline{\eta}}$ の計算は読者に委ねる。
この構成が前段落の構成の逆であることの確認も省略する。
\end{proof}

\begin{remark}
\label{etale-cohomology-remark-functorial-locally-constant-on-connected}
補題 \ref{etale-cohomology-lemma-locally-constant-on-connected} と
\ref{etale-cohomology-lemma-locally-constant-on-connected-geom-unibranch} の同値は引き戻しと両立する。
たとえば $f : Y \to X$ を連結スキームの射とする。$\overline{y}$ を $Y$ の
幾何学的点とし、$\overline{x} = f(\overline{y})$ とおく。このとき図式
$$
\xymatrix{
\text{集合の有限局所定数層で }Y_\etale\text{ 上のもの}
\ar[r] &
\text{有限 }\pi_1(Y, \overline{y})\text{-集合} \\
\text{集合の有限局所定数層で }X_\etale\text{ 上のもの}
\ar[r] \ar[u]_{f^{-1}} &
\text{有限 }\pi_1(X, \overline{x})\text{-集合} \ar[u]
}
$$
は可換である。ここで右側の垂直射は、連続準同型
$\pi_1(Y, \overline{y}) \to \pi_1(X, \overline{x})$、すなわち $f$ が誘導するもの
から得られる。これは基本群、定理 \ref{pione-theorem-fundamental-group} の
可換図式から直ちに従う。他の場合にも同様の結果が成り立つ。
\end{remark}









\section{トレース法}
\label{etale-cohomology-section-trace-method}

\noindent
この節の参考文献は \cite[Expos\'e IX, \S 5]{SGA4} である。
ここで扱う内容は、後の補題 \ref{etale-cohomology-lemma-vanishing-easier} の証明で用いる。

\medskip\noindent
$f : Y \to X$ をスキームのエタール射とする。このとき
$$
f_!, f^{-1}, f_*
$$
は $\textit{Ab}(X_\etale)$ と $\textit{Ab}(Y_\etale)$ の間の随伴関手の列である。
関手 $f_!$ は節 \ref{etale-cohomology-section-extension-by-zero} で論じる。
随伴写像 $\text{id} \to f_* f^{-1}$ を {\it 制限写像} と呼ぶ。
随伴写像 $f_! f^{-1} \to \text{id}$ はしばしば {\it トレース写像} と呼ばれる。
$f$ が有限エタールならば、$f_* = f_!$（補題
\ref{etale-cohomology-lemma-shriek-equals-star-finite-etale}）であり、これを写像
$f_*f^{-1} \to \text{id}$ とみなせる。

\begin{definition}
\label{etale-cohomology-definition-trace-map}
$f : Y \to X$ をスキームの有限エタール射とする。
上で述べ、以下で明示的に記述する写像 $f_* f^{-1} \to \text{id}$ を
{\it トレース} と呼ぶ。
\end{definition}

\noindent
$f : Y \to X$ をスキームの有限エタール射とする。トレース写像は次の二つの性質で特徴づけられる。
\begin{enumerate}
\item $X$ 上のエタール局所化と可換である。
\item $Y = \coprod_{i = 1}^d X$ ならば、トレース写像は和写像
$f_*f^{-1} \mathcal{F} = \mathcal{F}^{\oplus d} \to \mathcal{F}$ である。
\end{enumerate}
エタール射、補題 \ref{etale-lemma-finite-etale-etale-local} により、任意の有限エタール射
$f : Y \to X$ は $X$ 上エタール局所的には、ある整数 $d \geq 0$ に対して (2) の形をもつ。
したがって、上の特徴づけを用いてトレース写像を定義できる。特に、トレース写像を構成するために
$f_!$ の存在や $f_!$ と $f_*$ の一致を知る必要はない。
この記述から、$f$ が定次数 $d$ をもつならば、合成
$$
\mathcal{F} \xrightarrow{res}
f_* f^{-1} \mathcal{F} \xrightarrow{trace}
\mathcal{F}
$$
は $d$ 倍写像である。「トレース法」とは次の観察をいう。$\mathcal{F}$ が
$X_\etale$ 上のアーベル群の層であって、$d$ 倍写像が $\mathcal{F}$ 上で同型ならば、写像
$$
H^n_\etale(X, \mathcal{F}) \longrightarrow H^n_\etale(Y, f^{-1}\mathcal{F})
$$
は単射である。実際、
$$
H^n_\etale(Y, f^{-1}\mathcal{F}) = H^n_\etale(X, f_*f^{-1}\mathcal{F})
$$
が成り立つ。これは高次順像の消滅（命題
\ref{etale-cohomology-proposition-finite-higher-direct-image-zero}）と Leray スペクトル系列（命題
\ref{etale-cohomology-proposition-leray}）による。
したがって写像
$$
H^n_\etale(X, \mathcal{F}) \to
H^n_\etale(Y, f^{-1}\mathcal{F})= H^n_\etale(X, f_*f^{-1}\mathcal{F})
\xrightarrow{trace}
H^n_\etale(X, \mathcal{F})
$$
を考えることができ、その合成は（$\mathcal{F}$ と $f$ に関する上の仮定のもとで）同型である。
特に、$H_\etale^q(Y, f^{-1}\mathcal{F}) = 0$ ならば
$H_\etale^q(X, \mathcal{F}) = 0$ でもある。
実際、$d$ 倍写像は $H_\etale^q(X, \mathcal{F})$ 上の同型を誘導し、これは
$H_\etale^q(Y, f^{-1}\mathcal{F})= 0$ を経由する。

\medskip\noindent
これはしばしば次の結果と組み合わせて用いられる。

\begin{lemma}
\label{etale-cohomology-lemma-pullback-filtered}
$S$ を連結スキームとする。$\ell$ を素数とする。$\mathcal{F}$ を
有限型局所定数 $\mathbf{F}_\ell$-ベクトル空間の層であって、$S_\etale$ 上のものとする。
このとき、有限エタール射 $f : T \to S$ であって、その次数が $\ell$ と互いに素なものが存在し、
$f^{-1}\mathcal{F}$ は、逐次商が $\underline{\mathbf{Z}/\ell\mathbf{Z}}_T$ である
有限フィルトレーションをもつ。
\end{lemma}

\begin{proof}
幾何学的点 $\overline{s}$ を $S$ 上に選ぶ。
補題 \ref{etale-cohomology-lemma-locally-constant-on-connected} の同値により、層 $\mathcal{F}$ は、
有限次元 $\mathbf{F}_\ell$-ベクトル空間 $V$ であって、連続な
$\pi_1(S, \overline{s})$-作用を備えたものに対応する。
$G \subset \text{Aut}(V)$ を、この作用を与える準同型
$\rho : \pi_1(S, \overline{s}) \to \text{Aut}(V)$ の像とする。$G$ は有限であることに注意する。
全射な連続準同型
$\overline{\rho} : \pi_1(S, \overline{s}) \to G$
は、Galois 対象 $Y \to S$ であって $\textit{F\'Et}_S$ に属し、自己同型群が
$G = \text{Aut}(Y/S)$ であるものに対応する。基本群、節
\ref{pione-section-finite-etale-under-galois} を参照せよ。
$H \subset G$ を $\ell$-Sylow 部分群とする。
$T = Y/H \to S$ が求めるものだと主張する。実際、$\overline{t} \in T$ を
$\overline{s}$ 上の幾何学的点とする。基本群の関手性により、
$\pi_1(T, \overline{t}) \to \pi_1(S, \overline{s})$
の像は $(\overline{\rho})^{-1}(H)$ である。したがって、$\pi_1(T, \overline{t})$ の
$V$ 上の作用、すなわち $f^{-1}\mathcal{F}$ に対応する作用は、写像
$\pi_1(T, \overline{t}) \to H$ を経由する。注意
\ref{etale-cohomology-remark-functorial-locally-constant-on-connected} を参照せよ。
$H$ は有限 $\ell$-群であるから、表現 $\rho|_{\pi_1(T, \overline{t})}$ の
既約組成因子はすべて階数 $1$ の自明表現である
（これは有限群の表現論に関する簡単な補題である。ここに将来の参照を挿入する）。
補題 \ref{etale-cohomology-lemma-locally-constant-on-connected} の同値により、これは
$f^{-1}\mathcal{F}$ が、値 $\underline{\mathbf{Z}/\ell\mathbf{Z}}_T$ をもつ定数層の
逐次拡大であることを意味する。さらに、$T = Y/H \to S$ の次数は
$\ell$ と互いに素である。実際、それは $H$ の $G$ における指数に等しい。
\end{proof}

\begin{lemma}
\label{etale-cohomology-lemma-action-l-group}
$\Lambda$ を Noether 環とする。$\ell$ を素数とし、$n \geq 1$ とする。
$H$ を有限 $\ell$-群とする。$M$ を有限生成 $\Lambda[H]$-加群であって
$\ell^n$ で零化されるものとする。このとき有限フィルトレーション
$0 = M_0 \subset M_1 \subset \ldots \subset M_t = M$
であって、$\Lambda[H]$-部分加群からなり、$H$ が $M_{i + 1}/M_i$ に自明に作用するという
性質をすべての $i = 0, \ldots, t - 1$ に対してもつものが存在する。
\end{lemma}

\begin{proof}
省略する。ヒント：増大イデアル $\mathfrak m$、すなわち非可換環
$\mathbf{Z}/\ell^n\mathbf{Z}[H]$ のそれが冪零であることを示せ。
\end{proof}

\begin{lemma}
\label{etale-cohomology-lemma-pullback-filtered-modules}
$S$ を既約で幾何学的に単枝なスキームとする。$\ell$ を素数とし、$n \geq 1$ とする。
$\Lambda$ を Noether 環とする。$\mathcal{F}$ を有限型局所定数 $\Lambda$-加群の層であって、
$S_\etale$ 上にあり、$\ell^n$ で零化されるものとする。このとき、有限エタール射
$f : T \to S$ であって、その次数が $\ell$ と互いに素なものが存在し、
$f^{-1}\mathcal{F}$ は、逐次商が $\underline{M}_T$ の形をもつ有限フィルトレーションをもつ。
ここで $\Lambda$-加群 $M$ は有限生成である。
\end{lemma}

\begin{proof}
幾何学的点 $\overline{s}$ を $S$ 上に選ぶ。
補題 \ref{etale-cohomology-lemma-locally-constant-on-connected-geom-unibranch} の同値により、層
$\mathcal{F}$ は、有限生成 $\Lambda$-加群 $M$ であって、連続な
$\pi_1(S, \overline{s})$-作用を備えたものに対応する。
$G \subset \text{Aut}(V)$ を、この作用を与える準同型
$\rho : \pi_1(S, \overline{s}) \to \text{Aut}(M)$ の像とする。
$G$ は有限である。これは $M$ が有限生成 $\Lambda$-加群だからである
（補題 \ref{etale-cohomology-lemma-locally-constant-on-connected-geom-unibranch} の証明を参照せよ）。
全射な連続準同型
$\overline{\rho} : \pi_1(S, \overline{s}) \to G$
は、Galois 対象 $Y \to S$ であって $\textit{F\'Et}_S$ に属し、自己同型群が
$G = \text{Aut}(Y/S)$ であるものに対応する。基本群、節
\ref{pione-section-finite-etale-under-galois} を参照せよ。
$H \subset G$ を $\ell$-Sylow 部分群とする。
$T = Y/H \to S$ が求めるものだと主張する。実際、$\overline{t} \in T$ を
$\overline{s}$ 上の幾何学的点とする。基本群の関手性により、
$\pi_1(T, \overline{t}) \to \pi_1(S, \overline{s})$
の像は $(\overline{\rho})^{-1}(H)$ である。したがって、$\pi_1(T, \overline{t})$ の
$M$ 上の作用、すなわち $f^{-1}\mathcal{F}$ に対応する作用は、写像
$\pi_1(T, \overline{t}) \to H$ を経由する。注意
\ref{etale-cohomology-remark-functorial-locally-constant-on-connected} を参照せよ。
補題 \ref{etale-cohomology-lemma-action-l-group} のように
$0 = M_0 \subset M_1 \subset \ldots \subset M_t = M$ を取る。
これはフィルトレーション
$0 = \mathcal{F}_0 \subset \mathcal{F}_1 \subset \ldots
\subset \mathcal{F}_t = f^{-1}\mathcal{F}$
であって、逐次商が値 $M_{i + 1}/M_i$ をもつ定数層であるものを誘導する。
最後に、$T = Y/H \to S$ の次数は $\ell$ と互いに素である。
実際、それは $H$ の $G$ における指数に等しい。
\end{proof}






\section{ガロアコホモロジー}
\label{etale-cohomology-section-galois-cohomology}

\noindent
この節では、ガロアコホモロジーに関する結果（命題 \ref{etale-cohomology-proposition-serre-galois}）を、
エタールコホモロジーと節 \ref{etale-cohomology-section-trace-method} の技巧を用いて証明する。
これにより、体のスペクトル上の高次エタールコホモロジー群の消滅を証明できるようになる。

\begin{lemma}
\label{etale-cohomology-lemma-nonvanishing-inherited}
$\ell$ を素数、$n$ を整数 $> 0$ とする。
$S$ を準コンパクトかつ準分離なスキームとする。
$X = \lim_{i \in I} X_i$ を $S$-スキームの有向系の極限とし、各
$X_i \to S$ は、$\ell$ と互いに素な定次数をもつ有限エタール射であるとする。
次は同値である。
\begin{enumerate}
\item $\ell$-冪捩れ層 $\mathcal{G}$ が $S$ 上に存在して
$H_\etale^n(S, \mathcal{G}) \neq 0$ である。
\item $\ell$-冪捩れ層 $\mathcal{F}$ が $X$ 上に存在して
$H_\etale^n(X, \mathcal{F}) \neq 0$ である。
\end{enumerate}
実際、$\mathcal{G}$ が与えられたときは $\mathcal{F} = g^{-1}\mathcal{F}$ と取ることができ、
$\mathcal{F}$ が与えられたときは $\mathcal{G} = g_*\mathcal{F}$ と取ることができる。
\end{lemma}

\begin{proof}
$g : X \to S$ と $g_i : X_i \to S$ で構造射を表す。
$\ell$-冪捩れ層 $\mathcal{G}$ で $S$ 上にあり、
$H^n_\etale(S, \mathcal{G}) \not = 0$ となるものを固定する。
$\mathcal{G}_i = g_i^{-1}\mathcal{G}$ で与えられる系は、余極限層が
$g^{-1}\mathcal{G}$ である定理 \ref{etale-cohomology-theorem-colimit} の条件を満たす。したがって
$$
\colim_{i\in I} H^n_\etale(X_i, g_i^{-1}\mathcal{G}) =
H^n_\etale(X, \mathcal{G})
$$
である。$g_i$ は、次数が $\ell$ と互いに素な有限エタール射なので、
「トレース法」を適用でき、写像
$$
H^n_\etale(S, \mathcal{G}) \to H^n_\etale(X_i, g_i^{-1}\mathcal{G})
$$
はすべて単射である（また遷移写像と両立する）ことが分かる。
節 \ref{etale-cohomology-section-trace-method} を参照せよ。したがって余極限は零でない。すなわち
$H^n(X,g^{-1}\mathcal{G}) \neq 0$ であり、$\mathcal{F} = g^{-1}\mathcal{G}$ として
所望の結果を得る。

\medskip\noindent
逆に、$\ell$-冪捩れ層 $\mathcal{F}$ で $X$ 上にあり、
$H^n_\etale(X, \mathcal{F}) \not = 0$ となるものが与えられたとする。
$g_i$ は有限射なので、高次順像は消滅する（命題
\ref{etale-cohomology-proposition-finite-higher-direct-image-zero}）ことに注意する。
すると補題 \ref{etale-cohomology-lemma-relative-colimit} を適用することにより、$g$ についても同じことが従う。
高次順像の消滅と Leray（命題 \ref{etale-cohomology-proposition-leray}）から
$H^n_\etale(X, \mathcal{F}) = H^n(S, g_*\mathcal{F}) \neq 0$ であり、
$\mathcal{G} = g_*\mathcal{F}$ として所望の結果を得る。
\end{proof}

\begin{lemma}
\label{etale-cohomology-lemma-reduce-to-l-group}
$\ell$ を素数、$n$ を整数 $> 0$ とする。$K$ を
$G = Gal(K^{sep}/K)$ をもつ体とし、$H \subset G$ を極大 pro-$\ell$ 部分群、
$L/K$ をそれに対応する体拡大とする。このとき、
$H^n_\etale(\Spec(K), \mathcal{F}) = 0$ がすべての $\ell$-冪捩れ層
$\mathcal{F}$ に対して成り立つことと、
$H^n_\etale(\Spec(L), \underline{\mathbf{Z}/\ell\mathbf{Z}}) = 0$ となることは同値である。
\end{lemma}

\begin{proof}
$L = \bigcup L_i$ と、$K$ 上有限な部分拡大の合併として書く。
$H$ の選び方により、$[L_i : K]$ は $\ell$ と互いに素である。
したがって、補題 \ref{etale-cohomology-lemma-nonvanishing-inherited} と同様に
$\Spec(L) = \lim_{i \in I} \Spec(L_i)$ である。
ゆえに $K$ を $L$ で置き換え、絶対ガロア群 $G$、すなわち $K$ のそれが
副有限 pro-$\ell$ 群であると仮定してよい。

\medskip\noindent
$H^n(\Spec(K), \underline{\mathbf{Z}/\ell\mathbf{Z}}) = 0$ と仮定する。
$\mathcal{F}$ を $\ell$-冪捩れ層であって $\Spec(K)_\etale$ 上のものとする。
$H^n_\etale(\Spec(K), \mathcal{F}) = 0$ を示す。
補題 \ref{etale-cohomology-lemma-equivalence-abelian-sheaves-point} で述べた対応により、層
$\mathcal{F}$ は $\ell$-冪捩れ $G$-加群 $M$ に対応する。
元の任意の有限集合 $x_1, \ldots, x_m \in M$ は、連続性によりある開部分群 $U$ で固定される。
$M'$ を $x_1, \ldots, x_m$ の軌道で生成される加群とする。
$\ell$ のある冪は各 $x_i$ を零化し、軌道は有限なので、これは有限アーベル $\ell$-群である。
$M$ はこれらの部分加群 $M'$ のフィルター付き余極限だから、$\mathcal{F}$ は対応する部分層
$\mathcal{F}' \subset \mathcal{F}$ のフィルター付き余極限である。
この余極限に定理 \ref{etale-cohomology-theorem-colimit} を適用し、$\mathcal{F}$ が有限局所定数層である場合に帰着する。

\medskip\noindent
$M$ を、有限アーベル $\ell$-群であって、副有限 pro-$\ell$ 群 $G$ の連続作用を備えたものとする。
このとき $G$-不変なフィルトレーション
$$
0 = M_0 \subset M_1 \subset \ldots \subset M_r = M
$$
であって、$M_{i + 1}/M_i \cong \mathbf{Z}/\ell \mathbf{Z}$ 上の
$G$-作用が自明となるものがある（これは有限群の表現論に関する簡単な補題である。
ここに将来の参照を挿入する）。したがって、対応する層 $\mathcal{F}$ はフィルトレーション
$$
0 = \mathcal{F}_0 \subset \mathcal{F}_1 \subset \ldots \subset
\mathcal{F}_r = \mathcal{F}
$$
をもち、その逐次商は $\underline{\mathbf{Z}/\ell \mathbf{Z}}$ と同型である。
ゆえに帰納法とコホモロジー長完全列により結論を得る。
\end{proof}

\begin{lemma}
\label{etale-cohomology-lemma-reduce-to-l-group-higher}
$\ell$ を素数、$n$ を整数 $> 0$ とする。$K$ を
$G = Gal(K^{sep}/K)$ をもつ体とし、$H \subset G$ を極大 pro-$\ell$ 部分群、
$L/K$ をそれに対応する体拡大とする。このとき、
$H^q_\etale(\Spec(K),\mathcal{F}) = 0$ が $q \geq n$ とすべての
$\ell$-捩れ層 $\mathcal{F}$ に対して成り立つことと、
$H^n_\etale(\Spec(L), \underline{\mathbf{Z}/\ell\mathbf{Z}}) = 0$ となることは同値である。
\end{lemma}

\begin{proof}
順方向は自明なので、逆方向だけを証明すればよい。$q$ に関する帰納法で進む。
$q = n$ の場合は補題 \ref{etale-cohomology-lemma-reduce-to-l-group} である。
いま $\mathcal{F}$ を $\ell$-冪捩れ層であって $\Spec(K)$ 上のものとする。
$f : \Spec(K^{sep}) \rightarrow \Spec(K)$ を幾何学的点の包含とする。
このとき完全列
$$
0 \rightarrow \mathcal{F} \xrightarrow{res}
f_* f^{-1} \mathcal{F} \rightarrow f_* f^{-1} \mathcal{F}/\mathcal{F} 
\rightarrow 0
$$
を考える。$K^{sep}$ は有限分離拡大のフィルター付き余極限として書けることに注意する。
したがって $f$ は有限エタール射の有向系の極限である。
補題 \ref{etale-cohomology-lemma-nonvanishing-inherited} の証明で見たように、$f$ の高次順像は消滅すると結論できる。
ゆえに $f_* f^{-1} \mathcal{F}$ の高次コホモロジーは幾何学的点上の高次コホモロジーとして表せ、
これは明らかに消滅する。したがって、ここで現れるものはすべて依然として $\ell$-捩れなので、
帰納法の仮定とコホモロジー長完全列を併用して $q + 1$ に対する結果を得る。
\end{proof}

\begin{proposition}
\label{etale-cohomology-proposition-serre-galois}
\begin{reference}
\cite[Chapter II, Section 3, Proposition 5]{SerreGaloisCohomology}
\end{reference}
$K$ を分離代数閉包 $K^{sep}$ をもつ体とする。任意の有限拡大 $K'$ で $K$ のものに対して
$\text{Br}(K') = 0$ であると仮定する。このとき
\begin{enumerate}
\item $H^q(\text{Gal}(K^{sep}/K), (K^{sep})^*) = 0$ が
すべての $q \geq 1$ に対して成り立つ。
\item $H^q(\text{Gal}(K^{sep}/K), M) = 0$ が、任意の捩れ
$\text{Gal}(K^{sep}/K)$-加群 $M$ と任意の $q \geq 2$ に対して成り立つ。
\end{enumerate}
\end{proposition}

\begin{proof}
$p = \text{char}(K)$ とおく。補題 \ref{etale-cohomology-lemma-compare-cohomology-point}、
定理 \ref{etale-cohomology-theorem-brauer-delta}、および例 \ref{etale-cohomology-example-sheaves-point} により、
この命題は次を示すことと同値である。すなわち、
$H^2(\Spec(K'),\mathbf{G}_m|_{\Spec(K')_\etale}) = 0$ が
すべての有限拡大 $K'/K$ に対して成り立つならば、
\begin{itemize}
\item $H^q(\Spec(K),\mathbf{G}_m|_{\Spec(K)_\etale}) = 0$ が
すべての $q \geq 1$ に対して成り立つ。
\item $H^q(\Spec(K),\mathcal{F}) = 0$ が、任意の捩れ層
$\mathcal{F}$ と任意の $q \geq 2$ に対して成り立つ。
\end{itemize}
まず第二の主張を証明する。$\mathcal{F}$ は捩れ層なので、$\ell$-準素分解と、
コホモロジーが余極限（すなわち直和、定理 \ref{etale-cohomology-theorem-colimit} を参照）と両立することを用いて、
各素数について、$H^q(\Spec(K),\mathcal{F}) = 0$、$q \geq 2$ を、すべての
$\ell$-冪捩れ層に対して示すことへ帰着できる。ここで $\ell$ はその素数である。
これにより各素数を個別に解析できる。

\medskip\noindent
$\ell \neq p$ と仮定する。任意の拡大 $K'/K$ に対し、Kummer 完全列（補題
\ref{etale-cohomology-lemma-kummer-sequence}）
$$
0 \to
\mu_{\ell, \Spec{K'}} \to
\mathbf{G}_{m, \Spec{K'}} \xrightarrow{(\cdot)^{\ell}}
\mathbf{G}_{m, \Spec{K'}} \to 0
$$
を考える。$H^q(\Spec{K'},\mathbf{G}_m|_{\Spec(K')_\etale}) = 0$ は、$q = 2$ では仮定により、
$q = 1$ では定理 \ref{etale-cohomology-theorem-picard-group} と $\Pic(K) = (0)$ により成り立つ。
したがってコホモロジー長完全列から、$H^2(\Spec{K'}, \mu_\ell) = 0$ と、
任意の分離拡大 $K'/K$ に対して結論できる。
いま $H$ を極大 pro-$\ell$ 部分群であって $K$ の絶対ガロア群に属するものとし、
$L$ を対応する拡大とする。
$L$ を有限拡大の余極限として書けるから、この余極限に定理 \ref{etale-cohomology-theorem-colimit} を適用して
$H^2(\Spec(L), \mu_\ell) = 0$ を得る。
ここで $\mu_\ell$ は定数層でなければならない。そうでなければ、
$\ell$ と互いに素な次数をもつ $L$ の Galois 拡大、すなわち $L$ に $\ell$ 乗根を添加して得られる拡大が
存在することになるが、これは $L$ の定義に反する。
したがって補題 \ref{etale-cohomology-lemma-reduce-to-l-group-higher} により、$\ell \neq p$ に対する結果を得る。

\medskip\noindent
次に $\ell = p$ と仮定する。Artin-Schrier 完全列（節 \ref{etale-cohomology-section-artin-schreier}）
$$
0 \longrightarrow \underline{\mathbf{Z}/p\mathbf{Z}}_{\Spec{K}} \longrightarrow
\mathbf{G}_{a, \Spec{K}} \xrightarrow{F-1} \mathbf{G}_{a, \Spec{K}} 
\longrightarrow 0
$$
を考える。ここで $F - 1$ は写像 $x \mapsto x^p - x$ である。
注意 \ref{etale-cohomology-remark-special-case-fpqc-cohomology-quasi-coherent} と、アフィンスキームの構造層の
高次コホモロジーの消滅（スキームのコホモロジー、補題
\ref{coherent-lemma-quasi-coherent-affine-cohomology-zero}）により、
$\mathbf{G}_{a, \Spec{K}}$ の高次コホモロジーは消滅することに注意する。
これは標数 $p$ の任意の体に適用できる。特に、体拡大 $L$、すなわち極大 pro-$p$ 部分群
$H$ が定めるものに適用できる。これにより、
$H^n(\Spec{L}, \underline{\mathbf{Z}/p\mathbf{Z}}_{\Spec{L}}) = 0$ と
$n \geq 2$ に対して結論でき、
補題 \ref{etale-cohomology-lemma-reduce-to-l-group-higher} により $\ell = p$ に対する結果が従う。

\medskip\noindent
証明を終えるには、$H^q(\text{Gal}(K^{sep}/K), (K^{sep})^*) = 0$ が
すべての $q \geq 1$ に対して成り立つことを示す必要がある。
$G = \text{Gal}(K^{sep}/K)$ とおき、$M = (K^{sep})^*$ を $G$-加群とみなす。
上で既に $H^1(G, M) = 0$ および $H^2(G, M) = 0$ を示した。
完全列
$$
0 \to A \to M \to M \otimes \mathbf{Q} \to B \to 0
$$
を $G$-加群について考える。上の結果から、$H^i(G, A) = 0$ および
$H^i(G, B) = 0$ が $i > 1$ に対して成り立つ。実際、$A$ と $B$ は捩れ $G$-加群である。
補題 \ref{etale-cohomology-lemma-profinite-group-cohomology-torsion} により、
$H^i(G, M \otimes \mathbf{Q}) = 0$ が $i > 0$ に対して成り立つ。
これから $H^i(G, M) = 0$ が $i \geq 3$ に対しても従うことを確かめるのは楽しい演習である。
\end{proof}

%10.08.09

\begin{definition}
\label{etale-cohomology-definition-Cr}
体 $K$ が {\it $C_r$} であるとは、すべての $0 < d^r < n$ と、任意の斉次多項式
$f \in K[T_1, \ldots, T_n]$ で次数 $d$ のものに対して、$\alpha = (\alpha_1, \ldots, \alpha_n)$、
$\alpha_i \in K$ で、すべてが零ではなく、$f(\alpha) = 0$ となるものが存在することをいう。
このような $\alpha$ を $f$ の {\it 非自明解} と呼ぶ。
\end{definition}

\begin{example}
\label{etale-cohomology-example-algebraically-closed-field-Cr}
代数閉体は $C_r$ である。
\end{example}

\noindent
実際、次の簡単な補題がある。

\begin{lemma}
\label{etale-cohomology-lemma-algebraically-closed-find-solutions}
$k$ を代数閉体とする。$f_1, \ldots, f_s \in k[T_1, \ldots, T_n]$ を、
次数が $d_1, \ldots, d_s$ で $d_i > 0$ となる斉次多項式とする。
$s < n$ ならば、$f_1 = \ldots = f_s = 0$ は共通の非自明解をもつ。
\end{lemma}

\begin{proof}
これは次元論から従う。たとえば、多様体、補題
\ref{varieties-lemma-intersection-in-affine-space} を $s - 1$ 回適用すればよい。
\end{proof}

\noindent
次の結果は $C_1$ 体の Brauer 群を計算する。

\begin{theorem}
\label{etale-cohomology-theorem-C1-brauer-group-zero}
$K$ を $C_1$ 体とする。このとき $\text{Br}(K) = 0$ である。
\end{theorem}

\begin{proof}
$D$ を、中心が $K$ である $K$ 上の有限次元可除代数とする。既に
$$
D \otimes_K K^{sep} \cong \text{Mat}_d(K^{sep})
$$
が内部同型を除いて一意に成り立つことを見た。したがって行列式
$\det : \text{Mat}_d(K^{sep}) \to K^{sep}$ は Galois 不変であり、次数 $d$ の斉次写像
$$
\det = N_\text{red} : D \longrightarrow K
$$
へ降下する。これを {\it 被約ノルム} と呼ぶ。$K$ は $C_1$ なので、$d > 1$ ならば、
零でない $x \in D$ で $N_\text{red}(x) = 0$ を満たすものが存在する。
これは明らかに $x$ が可逆でないことを意味し、矛盾である。したがって $\text{Br}(K) = 0$ である。
\end{proof}

\begin{definition}
\label{etale-cohomology-definition-variety}
$k$ を体とする。{\it 多様体} とは、$k$ 上有限型の分離整スキームをいう。
{\it 曲線} とは次元 $1$ の多様体をいう。
\end{definition}

\begin{theorem}[Tsen の定理]
\label{etale-cohomology-theorem-tsen}
次元 $r$ の多様体で代数閉体 $k$ 上のものの関数体は $C_r$ である。
\end{theorem}

\begin{proof}
射影空間については、体 $k(x_1, \ldots, x_r)$ が $C_r$ であることを直接示せる（演習）。

\medskip\noindent
一般の場合。一般性を失わず、$X$ は射影的であると仮定してよい。
$f \in k(X)[T_1, \ldots, T_n]_d$ かつ $0 < d^r < n$ とする。
$f$ の係数は $\Gamma(X, \mathcal{O}_X(H))$ に属するとする。ただし $H \subset X$ は豊富である。
$\mathbf{\alpha} = (\alpha_1, \ldots, \alpha_n)$ かつ
$\alpha_i \in \Gamma(X, \mathcal{O}_X(eH))$ とする。このとき
$f(\mathbf{\alpha}) \in \Gamma(X, \mathcal{O}_X((de + 1)H))$ である。
方程式系 $f(\mathbf{\alpha}) =0$ を考える。すると漸近的 Riemann--Roch
（多様体、命題 \ref{varieties-proposition-asymptotic-riemann-roch}）により、ある $c > 0$ が存在して
\begin{itemize}
\item 変数の個数は $n\dim_k \Gamma(X, \mathcal{O}_X(eH)) \sim n e^r c$ であり、
\item 方程式の個数は
$\dim_k \Gamma(X, \mathcal{O}_X((de + 1)H)) \sim (de + 1)^r c$ である。
\end{itemize}
$n > d^r$ なので、変数の個数は方程式の個数より多い。方程式は斉次なので、補題
\ref{etale-cohomology-lemma-algebraically-closed-find-solutions} により解が存在する。
\end{proof}

\begin{lemma}
\label{etale-cohomology-lemma-curve-brauer-zero}
$C$ を代数閉体 $k$ 上の曲線とする。このとき $C$ の関数体の Brauer 群は零である。すなわち
$\text{Br}(k(C)) = 0$ である。
\end{lemma}

\begin{proof}
これは Tsen の定理、定理 \ref{etale-cohomology-theorem-tsen}、および定理
\ref{etale-cohomology-theorem-C1-brauer-group-zero} から明らかである。
\end{proof}

\begin{lemma}
\label{etale-cohomology-lemma-cohomology-Gm-function-field-curve}
$k$ を代数閉体、$K/k$ を超越次数 1 の体拡大とする。このとき、すべての
$q \geq 1$ に対して $H_\etale^q(\Spec(K), \mathbf{G}_m) = 0$ である。
\end{lemma}

\begin{proof}
補題 \ref{etale-cohomology-lemma-compare-cohomology-point} により
$H_\etale^q(\Spec(K), \mathbf{G}_m) = H^q(\text{Gal}(K^{sep}/K), (K^{sep})^*)$
であったことを思い出そう。したがって命題 \ref{etale-cohomology-proposition-serre-galois} により、
$K'/K$ が有限体拡大ならば $\text{Br}(K') = 0$ であることを示せば十分である。
ここで $K' = \colim K''$ であることに注意する。ただし $K''$ は、$k$ の有限生成部分拡大で
$K'$ に含まれ、超越次数 $1$ のもの全体を走る。
$\text{Br}(K') = \colim \text{Br}(K'')$ であることに注意すると、有限生成体拡大
$K''/k$ で超越次数 $1$ のものの場合に帰着する。このような体は $k$ 上の曲線の関数体なので、
補題 \ref{etale-cohomology-lemma-curve-brauer-zero} により Brauer 群は自明である。
\end{proof}






\section{乗法群の高次消滅}
\label{etale-cohomology-section-higher-Gm}

\noindent
この節では、代数閉体 $k$ と滑らかな曲線 $X$ を固定し、後者は $k$ 上とする。
$i_x : x \hookrightarrow X$ を $X$ の閉点の包含と書き、
$j : \eta \hookrightarrow X$ を生成点の包含と書く。また、$X_0$ を $X$ の閉点全体の集合と書く。

\begin{theorem}[基本完全列]
\label{etale-cohomology-theorem-fundamental-exact-sequence}
$X$ 上のエタール層の短完全列
$$
0 \longrightarrow
\mathbf{G}_{m, X} \longrightarrow
j_* \mathbf{G}_{m, \eta} \longrightarrow
\bigoplus\nolimits_{x \in X_0} {i_x}_* \underline{\mathbf{Z}}
\longrightarrow 0.
$$
が存在する。
\end{theorem}

\begin{proof}
$\varphi : U \to X$ をエタール射とする。エタール射の性質
（命題 \ref{etale-cohomology-proposition-etale-morphisms}）により、
$U = \coprod_i U_i$ と書け、各 $U_i$ は $X$ への射をもつ滑らかな曲線である。
$U$ に対する上の列は、各 $U_i$ に対する対応する列の積なので、
$U$ が連結、したがって既約である場合を扱えば十分である。この場合、よく知られた完全列
$$
1 \longrightarrow
\Gamma(U, \mathcal{O}_U^*) \longrightarrow
k(U)^* \longrightarrow \bigoplus\nolimits_{y \in U_0} \mathbf{Z}_y.
$$
がある。これは次の列を与える：
$$
0 \longrightarrow \Gamma(U, \mathcal{O}_U^*) \longrightarrow
\Gamma(\eta \times_X U, \mathcal{O}_{\eta \times_X U}^*) \longrightarrow
\bigoplus\nolimits_{x \in X_0}
\Gamma(x \times_X U, \underline{\mathbf{Z}})
$$
定義を展開すれば、これはほかならぬ次の列である：
$$
0 \longrightarrow \mathbf{G}_m(U) \longrightarrow
j_* \mathbf{G}_{m, \eta}(U) \longrightarrow
\left(\bigoplus\nolimits_{x \in X_0} {i_x}_* \underline{\mathbf{Z}}\right)(U).
$$
これにより基本完全列の射が定まり、最後の箇所を除いて完全であることが分かる。
全射性を見るため、$U$ が非特異曲線で $D$ が $U$ 上の因子ならば、
Zariski 開被覆 $\{U_j \to U\}$ が $U$ に存在し、
$D|_{U_j} = \text{div}(f_j)$ がある $f_j \in k(U)^*$ について成り立つことを思い出そう。
\end{proof}

\begin{lemma}
\label{etale-cohomology-lemma-higher-direct-jstar-Gm}
任意の $q \geq 1$ に対して $R^q j_*\mathbf{G}_{m, \eta} = 0$ である。
\end{lemma}

\begin{proof}
$(R^q j_*\mathbf{G}_{m, \eta})_{\bar x} = 0$ であることを、
任意の幾何学的点 $\bar x$（$X$ の点）について示せばよい。

\medskip\noindent
$\bar x$ が閉点 $x$（$X$ の点）の上にあると仮定する。$\Spec(A)$ を
$x$ の $X$ におけるアフィン開近傍とし、$K$ を $A$ の分数体とする。このとき
$$
\Spec(\mathcal{O}^{sh}_{X, \bar x}) \times_X \eta =
\Spec(\mathcal{O}^{sh}_{X, \bar x} \otimes_A K).
$$
環 $\mathcal{O}^{sh}_{X, \bar x} \otimes_A K$ は離散付値環
$\mathcal{O}^{sh}_{X, \bar x}$ の局所化であるから、再び
$\mathcal{O}^{sh}_{X, \bar x}$ であるか、その分数体
$K^{sh}_{\bar x}$ である。しかし、ある局所一意化元が可逆になるので、後者でなければならない。したがって
$$
(R^q j_*\mathbf{G}_{m, \eta})_{(X, \bar x)} = H_\etale^q(\Spec
K^{sh}_{\bar x}, \mathbf{G}_m).
$$
ここで
$\mathcal{O}^{sh}_{X, \bar x} =
\colim_{(U, \bar u) \to \bar x} \mathcal{O}(U) = \colim_{A \subset B} B$
であり、ここで $A \to B$ はエタールであることを思い出す。したがって
$K^{sh}_{\bar x}$ は $K = k(X)$ の代数拡大であり、
補題 \ref{etale-cohomology-lemma-cohomology-Gm-function-field-curve}
を適用すると消滅が得られる。

\medskip\noindent
$\bar x = \bar \eta$ が生成点 $\eta$（$X$ の点）の上にあると仮定する
（実際、この場合は不要である）。このとき
$\mathcal{O}^{sh}_{X, \bar \eta} = \kappa(\eta)^{sep}$ であり、したがって
\begin{eqnarray*}
(R^q j_*\mathbf{G}_{m, \eta})_{\bar \eta}
& = &
H_\etale^q(\Spec(\kappa(\eta)^{sep}) \times_X \eta, \mathbf{G}_m) \\
& = & H_\etale^q (\Spec(\kappa(\eta)^{sep}), \mathbf{G}_m) \\
& = & 0 \ \ \text{ただし } q \geq 1
\end{eqnarray*}
となる。これは対応する Galois 群が自明だからである。
\end{proof}

\begin{lemma}
\label{etale-cohomology-lemma-cohomology-jstar-Gm}
すべての $p \geq 1$ に対して $H_\etale^p(X, j_*\mathbf{G}_{m, \eta}) = 0$ である。
\end{lemma}

\begin{proof}
Leray スペクトル系列は
$$
E_2^{p, q} = H_\etale^p(X, R^qj_*\mathbf{G}_{m, \eta}) \Rightarrow
H_\etale^{p+q}(\eta, \mathbf{G}_{m, \eta}),
$$
と書ける。補題 \ref{etale-cohomology-lemma-cohomology-Gm-function-field-curve} により、
右辺は $p + q \geq 1$ のとき消滅する。
補題 \ref{etale-cohomology-lemma-higher-direct-jstar-Gm} により、
左辺は $q \geq 1$ のとき消滅する。よって所望の消滅が得られる。
\end{proof}

\begin{lemma}
\label{etale-cohomology-lemma-cohomology-istar-Z}
すべての $q \geq 1$ に対して $H_\etale^q(X, \bigoplus_{x \in X_0} {i_x}_*
\underline{\mathbf{Z}}) = 0$ である。
\end{lemma}

\begin{proof}
$X$ が準コンパクトかつ準分離なら、コホモロジーは余極限と可換するので、
$H_\etale^q(X, {i_x}_* \underline{\mathbf{Z}})$ の消滅を示せば十分である。
ところが閉点の包含 $i_x$ は有限であるから、
$R^p {i_x}_* \underline{\mathbf{Z}} = 0$ がすべての $p \geq 1$ に対して成り立つことが、
命題 \ref{etale-cohomology-proposition-finite-higher-direct-image-zero} から従う。
Leray スペクトル系列を適用すると、
$H_\etale^q(X, {i_x}_* \underline{\mathbf{Z}}) =
H_\etale^q(x, \underline{\mathbf{Z}})$ であることが分かる。
最後に、$x$ は代数閉体のスペクトルなので、$x$ 上のすべての高次コホモロジーは消滅する。
\end{proof}

\noindent
以上の一連の補題から、次の結果を得る。

\begin{theorem}
\label{etale-cohomology-theorem-vanishing-cohomology-Gm-curve}
$X$ を代数閉体上の滑らかな曲線とする。このとき
$$
H_\etale^q(X, \mathbf{G}_m) = 0 \ \ \text{任意の } q \geq 2.
$$
\end{theorem}

\begin{proof}
上の議論を参照せよ。
\end{proof}

\noindent
さらに、コホモロジー長完全列
$$
0 \to
H_\etale^0(X, \mathbf{G}_m) \to
H_\etale^0(X, j_*\mathbf{G}_{m\eta}) \to
H_\etale^0(X, \bigoplus {i_x}_*\underline{\mathbf{Z}}) \to
H_\etale^1(X, \mathbf{G}_m) \to 0
$$
も得られるが、これは周知の列
$$
0 \to H_{Zar}^0(X, \mathcal{O}_X^*) \to k(X)^* \to \text{Div}(X)
\to \Pic(X) \to 0.
$$
にほかならない。





\section{曲線の Picard 群}
\label{etale-cohomology-section-pic-curves}

\noindent
次の段階では Kummer 完全列を用いて、有限係数をもつ曲線のコホモロジー群について
いくつかの情報を導く。長完全列で消滅を得るため、Picard 群に関するいくつかの事実を復習する。

\medskip\noindent
$X$ を代数閉体 $k$ 上の滑らかな射影曲線とする。
$g = \dim_k H^1(X, \mathcal{O}_X)$ を $X$ の種数とする。短完全列
$$
0 \to \Pic^0(X) \to \Pic(X) \xrightarrow{\deg} \mathbf{Z} \to 0.
$$
が存在する。アーベル群 $\Pic^0(X)$ は
$\Pic^0(X) = \underline{\Picardfunctor}^0_{X/k}(k)$ と同一視できる。すなわち、
これは $k$-値点として得られるもので、対応するアーベル多様体は
$\underline{\Picardfunctor}^0_{X/k}$ であり、$k$ 上の次元 $g$ をもつ。
したがって、$n \in k^*$ ならば
$\Pic^0(X)[n] \cong (\mathbf{Z}/n\mathbf{Z})^{2g}$
がアーベル群として成り立つ。Picard Schemes of Curves, 節
\ref{pic-section-picard-curve} および Groupoids, 節
\ref{groupoids-section-abelian-varieties} を参照せよ。
この重要な事実、すなわち代数閉体上の滑らかな射影曲線の Picard 群における
捩れ部分の記述には、初等的な証明が存在しないようである。

\begin{lemma}
\label{etale-cohomology-lemma-cohomology-smooth-projective-curve}
$X$ を種数 $g$ の滑らかな射影曲線とし、これは代数閉体 $k$ 上にあるとする。
$n \geq 1$ は $k$ で可逆とする。このとき標準的な同一視
$$
H_\etale^q(X, \mu_n) =
\left\{
\begin{matrix}
\mu_n(k) & \text{ただし }q = 0, \\
\Pic^0(X)[n] & \text{ただし }q = 1, \\
\mathbf{Z}/n\mathbf{Z} & \text{ただし }q = 2, \\
0 & \text{ただし }q \geq 3.
\end{matrix}
\right.
$$
がある。$\mu_n \cong \underline{\mathbf{Z}/n\mathbf{Z}}$ なので、これにより
（非標準的な）同一視
$$
H_\etale^q(X, \underline{\mathbf{Z}/n\mathbf{Z}}) \cong
\left\{
\begin{matrix}
\mathbf{Z}/n\mathbf{Z} & \text{ただし }q = 0, \\
(\mathbf{Z}/n\mathbf{Z})^{2g} & \text{ただし }q = 1, \\
\mathbf{Z}/n\mathbf{Z} & \text{ただし }q = 2, \\
0 & \text{ただし }q \geq 3.
\end{matrix}
\right.
$$
が得られる。
\end{lemma}

\begin{proof}
定理 \ref{etale-cohomology-theorem-picard-group} と
\ref{etale-cohomology-theorem-vanishing-cohomology-Gm-curve} は、$\mathbf{G}_m$ の $X$ 上の
エタールコホモロジーを $X$ の Picard 群によって決定する。
Kummer 完全列 $0\to \mu_{n, X} \to \mathbf{G}_{m, X}
\to \mathbf{G}_{m, X}\to 0$（補題 \ref{etale-cohomology-lemma-kummer-sequence}）
から、コホモロジー長完全列
$$
\xymatrix{
0 \ar[r] & \mu_n(k) \ar[r] &
k^* \ar[r]^{(\cdot)^n} &
k^* \ar@(rd, ul)[rdllllr] \\
& H_\etale^1(X, \mu_n) \ar[r] &
\Pic(X) \ar[r]^{(\cdot)^n} &
\Pic(X) \ar@(rd, ul)[rdllllr] \\
& H_\etale^2(X, \mu_n) \ar[r] & 0 \ar[r] & 0 \ldots
}
$$
を得る。$n$ 乗写像 $k^* \to k^*$ は、$k$ が代数閉であるため全射である。
したがって、写像 $n : \Pic(X) \to \Pic(X)$ の核と余核を計算する必要がある。
行が完全である可換図式
$$
\xymatrix{
0 \ar[r] &
\Pic^0(X) \ar[r] \ar@{>>}[d]^{(\cdot)^n} &
\Pic(X) \ar[r]_-\deg \ar[d]^{(\cdot)^n} &
\mathbf{Z} \ar[r] \ar@{^{(}->}[d]^n & 0 \\
0 \ar[r] &
\Pic^0(X) \ar[r] &
\Pic(X) \ar[r]^-\deg &
\mathbf{Z} \ar[r] & 0
}
$$
を考える。群 $\Pic^0(X)$ は $k$-値点からなる群で、その群スキームは
$\underline{\Picardfunctor}^0_{X/k}$ である。Picard Schemes of Curves, 補題
\ref{pic-lemma-picard-pieces} を参照せよ。同じ補題によれば、
$\underline{\Picardfunctor}^0_{X/k}$ は次元 $g$ のアーベル多様体であり、$k$ 上にある。
その意味は Groupoids, 定義 \ref{groupoids-definition-abelian-variety} による。
したがって、Groupoids, 命題 \ref{groupoids-proposition-review-abelian-varieties} により、
左の垂直写像は全射である。蛇の補題を適用すると、補題の主張どおりの標準的同一視が得られる。

\medskip\noindent
補題の非標準的同一視を得るには、$n : \Pic^0(X) \to \Pic^0(X)$ の核が
$(\mathbf{Z}/n\mathbf{Z})^{\oplus 2g}$ と同型であることを示す必要がある。
これも Groupoids, 命題 \ref{groupoids-proposition-review-abelian-varieties} の一部である。
\end{proof}

\begin{lemma}
\label{etale-cohomology-lemma-pullback-on-h2-curve}
$\pi : X \to Y$ を代数閉体 $k$ 上の滑らかな射影曲線の間の非定値射とし、
$n \geq 1$ は $k$ で可逆とする。写像
$$
\pi^* : H^2_\etale(Y, \mu_n) \longrightarrow H^2_\etale(X, \mu_n)
$$
は $\pi$ の次数による乗法で与えられる。
\end{lemma}

\begin{proof}
補題 \ref{etale-cohomology-lemma-cohomology-smooth-projective-curve} で二つのコホモロジー群
$ H^2_\etale(Y, \mu_n)$ と $H^2_\etale(X, \mu_n)$ を
$\mathbf{Z}/n\mathbf{Z}$ と同一視したので、この主張には意味がある。
実際、$\mathcal{L}$ が次数 $1$ の直線束で、$Y$ 上にあり、類
$[\mathcal{L}] \in H^1_\etale(Y, \mathbf{G}_m)$ をもつならば、
$[\mathcal{L}]$ の境界は $H^2_\etale(Y, \mu_n)$ の生成元である。
ここでいう境界は Kummer 完全列に付随するコホモロジー長完全列の境界写像である。
したがって補題の結果は、直線束 $\pi^*\mathcal{L}$ の $X$ 上での次数が
$\deg(\pi)$ であることから従う。いくつかの細部は省略する。
\end{proof}

\begin{lemma}
\label{etale-cohomology-lemma-vanishing-cohomology-mu-smooth-curve}
$X$ を代数閉体 $k$ 上の滑らかなアフィン曲線とし、$n \in k^*$ とする。
$X \subset \overline{X}$ を滑らかな射影コンパクト化とする
（Varieties, 注 \ref{varieties-remark-smooth-projective-compactification}）。
$g$ を $\overline{X}$ の種数、$r$ を $\overline{X} \setminus X$ の点の個数とする。このとき
\begin{enumerate}
\item $H_\etale^0(X, \mu_n) = \mu_n(k)$;
\item $H_\etale^1(X, \mu_n) \cong (\mathbf{Z}/n\mathbf{Z})^{2g+r-1}$、および
\item $H_\etale^q(X, \mu_n) = 0$ がすべての $q \geq 2$ に対して成り立つ。
\end{enumerate}
\end{lemma}

\begin{proof}
$X = \overline{X} - \{ x_1, \ldots, x_r\}$ と書く。このとき
$\Pic(X) = \Pic(\overline{X})/ R$ であり、ここで $R$ は
$\mathcal{O}_{\overline{X}}(x_i)$（$1 \leq i \leq r$）で生成される部分群である。
$r \geq 1$ なので、$\Pic^0(\overline{X}) \to \Pic(X)$ は全射である。
したがって $\Pic(X)$ は可除である（補題 \ref{etale-cohomology-lemma-cohomology-smooth-projective-curve} の
証明における議論を参照）。Kummer 完全列を適用すると (1) と (3) を得る。(2) には
\begin{align*}
H_\etale^1(X, \mu_n)
& =
\{(\mathcal L, \alpha)\}/\cong \\
& =
(\{(\bar{\mathcal L}, D, \bar \alpha)\}/\cong)/\tilde{R}
\end{align*}
を用いる。最初の等号と記法については補題 \ref{etale-cohomology-lemma-describe-h1-mun} を参照せよ。
第二行の三つ組 $(\bar{\mathcal L}, D, \bar \alpha)$ は、可逆加群
$\bar{\mathcal L}$（$\overline{X}$ 上で次数 $0$）、因子 $D$（$\overline{X}$ 上で次数 $0$）で
$\left\{x_1, \ldots, x_r\right\}$ に台をもち、さらに同型
$\bar{\alpha}: \bar{\mathcal L}^{\otimes n} \cong \mathcal{O}_{\bar{X}}(D)$ をもつものからなる。
三つ組の同型は、補題 \ref{etale-cohomology-lemma-describe-h1-mun} に先行する議論と同様に定義する。
さらに $\tilde{R}$ は、次の形の三つ組
$(\mathcal{O}_{\overline{X}}(D'), n D', 1^{\otimes n})$
で生成される部分群であり、ここで $D'$ は $\left\{x_1, \ldots, x_r\right\}$ に台をもち、
次数 0 である。上の議論より、三つ組 $(\bar{\mathcal L},\ D,\ \bar \alpha)$ を対
$(\bar{\mathcal L}|_X, \bar \alpha|_X)$ に送る写像によって第二の等号を得る。
したがって完全列
$$
0 \longrightarrow
H_\etale^1(\overline{X}, \mu_n) \longrightarrow
H_\etale^1(X, \mu_n) \longrightarrow
\bigoplus_{i = 1}^r \mathbf{Z}/n\mathbf{Z}
\xrightarrow{\ \sum\ }
\mathbf{Z}/n\mathbf{Z} \longrightarrow 0
$$
を得る。ここで中央の写像は、三つ組 $(\bar{ \mathcal L}, D, \bar \alpha)$ の類で、
$D = \sum_{i = 1}^r a_i (x_i)$ を満たすものを $r$-組
$(a_i)_{i = 1}^r$ に送る。あとは補題
\ref{etale-cohomology-lemma-cohomology-smooth-projective-curve} を用いて階数を数えればよい。
\end{proof}

\begin{remark}
\label{etale-cohomology-remark-natural-proof}
前の系を証明する「自然な」方法は、$X$ を $\bar X$ から切除することである。
これは可能だが、その理論をまだ展開していない。
\end{remark}

\begin{remark}
\label{etale-cohomology-remark-normalize-H1-Gm}
$k$ を代数閉体とする。$n$ を $k$ の標数と素な整数とする。次を思い出そう：
$$
\mathbf{G}_{m, k} = \mathbf{A}^1_k \setminus \{0\} =
\mathbf{P}^1_k \setminus \{0, \infty\}
$$
標準的同型
$$
H^1_\etale(\mathbf{G}_{m, k}, \mu_n) = \mathbf{Z}/n\mathbf{Z}
$$
があると主張する。これは何を意味するだろうか。それは、元 $1_k$ が
$H^1_\etale(\mathbf{G}_{m, k}, \mu_n)$ に存在し、任意の射
$\Spec(k') \to \Spec(k)$ に対して、写像
$\mathbf{G}_{m, k'} \to \mathbf{G}_{m, k}$ のエタールコホモロジー上の引き戻し写像が
$1_k$ を $1_{k'}$ に送る、という意味である（特に、この元は $k$ のすべての自己同型で固定される）。
これを見るため、$\mu_{n, \mathbf{Z}}$-トーサー
$\mathbf{G}_{m, \mathbf{Z}} \to \mathbf{G}_{m, \mathbf{Z}}$,
$x \mapsto x^n$ を考える。トーサーと第一コホモロジーとの同一視により、
これを引き戻すと標準的な元 $1_k$ が得られる。捩り戻すことにより、標準的同一視
$$
H^1_\etale(\mathbf{G}_{m, k}, \mathbf{Z}/n\mathbf{Z}) =
\Hom(\mu_n(k), \mathbf{Z}/n\mathbf{Z}),
$$
があることが分かる。すなわち、これらの同型は代数閉体の間の写像と両立し、
特に $k$ の自己同型と両立する。
\end{remark}










\section{零延長}
\label{etale-cohomology-section-extension-by-zero}

\noindent
Modules on Sites, 節 \ref{sites-modules-section-localize} の一般論により、
次の定義を行うことができる。

\begin{definition}
\label{etale-cohomology-definition-extension-zero}
$j : U \to X$ をスキームのエタール射とする。
\begin{enumerate}
\item 制限関手
$j^{-1} : \Sh(X_\etale) \to \Sh(U_\etale)$
は左随伴
$j_!^{Sh} : \Sh(U_\etale) \to \Sh(X_\etale)$.
をもつ。
\item 制限関手
$j^{-1} : \textit{Ab}(X_\etale) \to \textit{Ab}(U_\etale)$
は左随伴をもち、これを
$j_! : \textit{Ab}(U_\etale) \to \textit{Ab}(X_\etale)$
と書き、{\it 零延長}と呼ぶ。
\item $\Lambda$ を環とする。制限関手
$j^{-1} : \textit{Mod}(X_\etale, \Lambda) \to
\textit{Mod}(U_\etale, \Lambda)$
は左随伴をもち、これを
$j_! : \textit{Mod}(U_\etale, \Lambda) \to
\textit{Mod}(X_\etale, \Lambda)$
と書き、{\it 零延長}と呼ぶ。
\end{enumerate}
\end{definition}

\noindent
$\mathcal{F}$ が $X_\etale$ 上のアーベル層ならば、一般には
$j_!\mathcal{F} \not = j_!^{Sh}\mathcal{F}$ である。一方、
$j_!$ は $\Lambda$-加群の層に対して、台となるアーベル層上の $j_!$ と一致する
（Modules on Sites, 注 \ref{sites-modules-remark-localize-shriek-equal}）。
関手 $j_!$ は関手的同型
$$
\Hom_X(j_!\mathcal{F}, \mathcal{G}) = \Hom_U(\mathcal{F}, j^{-1}\mathcal{G})
$$
によって特徴づけられる。これはすべての $\mathcal{F} \in \textit{Ab}(U_\etale)$ と
$\mathcal{G} \in \textit{Ab}(X_\etale)$ に対して成り立つ。
$\Lambda$-加群の層についても同様である。

\medskip\noindent
定義 \ref{etale-cohomology-definition-extension-zero} の関手をさらに明示的に記述するため、
$j^{-1}$ は関手 $U_\etale \to X_\etale$ による制限にほかならないことを思い出そう。
言い換えると、$j^{-1}\mathcal{G}(U') = \mathcal{G}(U')$ が、
$U'$ が $U$ 上エタールであるとき成り立つ。一方、
$\mathcal{F} \in \textit{Ab}(U_\etale)$ に対して前層
\begin{equation}
\label{etale-cohomology-equation-j-p-shriek}
j_{p!}\mathcal{F} : X_\etale \longrightarrow \textit{Ab},
\quad
V \longmapsto \bigoplus\nolimits_{V \to U} \mathcal{F}(V \to U)
\end{equation}
を考える。このとき $j_!\mathcal{F}$ は $j_{p!}\mathcal{F}$ の層化である。
これは Modules on Sites, 補題 \ref{sites-modules-lemma-extension-by-zero} で証明されている。
より一般には Modules on Sites, 節 \ref{sites-modules-section-localize} と
\ref{sites-modules-section-exactness-lower-shriek} の議論を参照せよ。

\begin{exercise}
\label{etale-cohomology-exercise-jshriek-direct}
関手 $j_!$ が、式 (\ref{etale-cohomology-equation-j-p-shriek}) で与えられる関手
$j_{p!}$ の層化として定義され、$j^{-1}$ の左随伴であることを直接証明せよ。
\end{exercise}

\begin{proposition}
\label{etale-cohomology-proposition-describe-jshriek}
$j : U \to X$ をスキームのエタール射とする。
$\mathcal{F}$ を $\textit{Ab}(U_\etale)$ の対象とする。
$\overline{x} : \Spec(k) \to X$ が $X$ の幾何学的点ならば、
$$
(j_!\mathcal{F})_{\overline{x}} =
\bigoplus\nolimits_{\overline{u} : \Spec(k) \to U,\ j(\overline{u}) =
\overline{x}} \mathcal{F}_{\bar{u}}.
$$
特に、$j_!$ は完全関手である。
\end{proposition}

\begin{proof}
$j_!$ の完全性は非常に一般的である。Modules on Sites, 補題
\ref{sites-modules-lemma-extension-by-zero-exact} を参照せよ。
もちろん、これは茎の記述からも従う。茎の公式は Modules on Sites, 補題
\ref{sites-modules-lemma-stalk-j-shriek} と、幾何学的点による小エタールサイトの点の記述
（補題 \ref{etale-cohomology-lemma-points-small-etale-site}）から従う。

\medskip\noindent
後で使うため、同型
\begin{align*}
(j_!\mathcal{F})_{\overline{x}}
& =
(j_{p!}\mathcal{F})_{\overline{x}} \\
& =
\colim_{(V, \overline{v})} j_{p!}\mathcal{F}(V) \\
& =
\colim_{(V, \overline{v})}
\bigoplus\nolimits_{\varphi : V \to U}
\mathcal{F}(V \xrightarrow{\varphi} U) \\
& \to
\bigoplus\nolimits_{\overline{u} : \Spec(k) \to U,\ j(\overline{u}) =
\overline{x}} \mathcal{F}_{\bar{u}}.
\end{align*}
は Modules on Sites, 補題 \ref{sites-modules-lemma-stalk-j-shriek} で構成され、
$(V, \overline{v}, \varphi, s)$ を $s$ の類に送り、この類は $\mathcal{F}$ の
$\overline{u} = \varphi(\overline{v})$ における茎に属することに注意する。
\end{proof}

\begin{lemma}
\label{etale-cohomology-lemma-jshriek-open}
$j : U \to X$ をスキームの開埋め込みとする。任意のアーベル層
$\mathcal{F}$ を $U_\etale$ 上にとると、随伴写像
$j^{-1}j_*\mathcal{F} \to \mathcal{F}$ および
$\mathcal{F} \to j^{-1}j_!\mathcal{F}$ は同型である。
実際、$j_!\mathcal{F}$ は $X_\etale$ 上の一意なアーベル層であり、
$U$ への制限が $\mathcal{F}$ で、$X \setminus U$ の幾何学的点における茎が零である。
\end{lemma}

\begin{proof}
第一の主張は定義をたどって証明することを読者に勧めるが、ここでは非常に一般的な
Modules on Sites, 補題 \ref{sites-modules-lemma-restrict-back} の特別な場合であることを用いる。
第二の主張について、$\mathcal{G}$ が $X_\etale$ 上のアーベル層で、その $U$ への制限が
$\mathcal{F}$ ならば、随伴性から写像 $j_!\mathcal{F} \to \mathcal{G}$ を得る。
命題 \ref{etale-cohomology-proposition-describe-jshriek} により、この写像は $U$ の幾何学的点の茎で同型である。
したがって、$\mathcal{G}$ が $X \setminus U$ の幾何学的点で零の茎をもつならば、
$j_!\mathcal{F} \to \mathcal{G}$ は定理 \ref{etale-cohomology-theorem-exactness-stalks} により同型である。
\end{proof}

\begin{lemma}[零延長は基底変換と可換する]
\label{etale-cohomology-lemma-shriek-base-change}
$f: Y \to X$ をスキームの射とし、$j: V \to X$ をエタール射とする。
ファイバー積
$$
\xymatrix{
V' = Y \times_X V \ar[d]_{f'} \ar[r]_-{j'} & Y \ar[d]^f \\
V \ar[r]^j & X
}
$$
を考える。このとき、アーベル層および加群の層上で
$j'_! f'^{-1} = f^{-1} j_!$ が成り立つ。
\end{lemma}

\begin{proof}
$j'_! f'^{-1}$ は $f'_* (j')^{-1}$ の左随伴であり、$f^{-1} j_!$ は
$j^{-1}f_*$ の左随伴なので、これは成り立つ。さらに、
$f'_* (j')^{-1} = j^{-1}f_*$ である。なぜなら $f_*$ は構成によりエタール局所化と
可換するからである。実際、この補題はサイトの射の設定で
非常に一般的に成り立つ。Modules on Sites, 補題
\ref{sites-modules-lemma-localize-morphism-ringed-sites} を参照せよ。
\end{proof}

\begin{lemma}
\label{etale-cohomology-lemma-shriek-into-star-separated-etale}
$j : U \to X$ を分離エタール射とする。このとき、関手的な単射
$j_!\mathcal{F} \to j_*\mathcal{F}$ がアーベル層および $\Lambda$-加群の層上に存在する。
\end{lemma}

\begin{proof}
アーベル層の場合に証明する。標準的写像
$$
j_{p!}\mathcal{F} \to j_*\mathcal{F}
$$
を構成しよう。これは $X_\etale$ 上のアーベル前層の写像であり、任意のアーベル層
$\mathcal{F}$ を $U_\etale$ 上にとる。ここで $j_{p!}$ は式
(\ref{etale-cohomology-equation-j-p-shriek}) のとおりである。この写像を層化すれば、所望の写像
$j_!\mathcal{F} \to j_*\mathcal{F}$ を得る。両辺をエタール射 $V \to X$ 上で評価すると
$$
j_{p!}\mathcal{F}(V) =
\bigoplus\nolimits_{\varphi : V \to U} \mathcal{F}(V \xrightarrow{\varphi} U)
\quad\text{かつ}\quad
j_*\mathcal{F}(V) = \mathcal{F}(V \times_X U)
$$
を得る。各 $\varphi$ に対して開かつ閉な埋め込み
$$
\Gamma_\varphi = (1, \varphi) : V \longrightarrow V \times_X U
$$
が $U$ 上にある。これは $U$ 上エタールなスキーム間の射なので開であり、
$V$ 上分離なスキームの切断なので閉である
（Schemes, 補題 \ref{schemes-lemma-section-immersion}）。したがって切断
$s_\varphi \in \mathcal{F}(V \xrightarrow{\varphi} U)$ に対し、
一意な切断 $s'_\varphi$ が $\mathcal{F}(V \times_X U)$ に存在し、
$s_\varphi$ を $\Gamma_\varphi$ による引き戻しとしてもち、
$\Gamma_\varphi$ の像の補集合への制限は零である。

\medskip\noindent
この写像が単射であることを示すため、上の公式の元
$\sum_{i = 1, \ldots, n} s_{\varphi_i}$ が $j_{p!}\mathcal{F}(V)$ に属し、
$j_*\mathcal{F}(V)$ で零に写ると仮定する。示すべきことは、
$\sum_{i = 1, \ldots, n} s_{\varphi_i}$ が $V$ のあるエタール被覆の各要素上で
零に制限されることである。すべての二者間の等化子
（$V$ の開かつ閉な部分）を射 $\varphi_i : V \to U$ について見て、$V$ 上局所的に作業すれば、射の像
$\Gamma_{\varphi_1}, \ldots, \Gamma_{\varphi_n}$ が互いに素であると仮定できる。
仮定は $\sum_{i = 1, \ldots, n} s'_{\varphi_i} = 0$ なので、
$s'_{\varphi_i} = 0$ が各 $i$ について成り立つことが直ちに従う
（これらの切断の台が互いに素であるため）。したがって、所望どおり、
$s_{\varphi_i} = 0$ がすべての $i$ について成り立つ。
\end{proof}

\begin{lemma}
\label{etale-cohomology-lemma-shriek-equals-star-finite-etale}
$j : U \to X$ を有限エタール射とする。このとき、写像
$j_! \to j_*$（補題 \ref{etale-cohomology-lemma-shriek-into-star-separated-etale}）
はアーベル層および $\Lambda$-加群の層上で同型である。
\end{lemma}

\begin{proof}
$j_!\mathcal{F} \to j_*\mathcal{F}$ が $X$ 上エタール局所的に同型であることを
確かめれば十分である。したがって $U \to X$ は有限個の同型の非交和であると仮定できる。
\'Etale Morphisms, 補題 \ref{etale-lemma-finite-etale-etale-local} を参照せよ。
この場合の証明は省略する。
\end{proof}

\begin{lemma}
\label{etale-cohomology-lemma-ses-associated-to-open}
$X$ をスキームとする。$Z \subset X$ を閉部分スキームとし、$U \subset X$ をその補集合とする。
$i : Z \to X$ と $j : U \to X$ を包含射と書く。任意のアーベル層
$\mathcal{F}$ を $X_\etale$ 上にとると、標準的短完全列
$$
0 \to j_!j^{-1}\mathcal{F} \to \mathcal{F} \to i_*i^{-1}\mathcal{F} \to 0
$$
が $X_\etale$ 上に存在する。
\end{lemma}

\begin{proof}
関係する関手の随伴性から射を得る。幾何学的点 $\overline{x}$（$X$ の点）について、
$\overline{x} \in U$ ならば左側の写像は茎上で同型であり、
$i_*i^{-1}\mathcal{F}$ の茎は零である。または $\overline{x} \in Z$ ならば右側の写像は
茎上で同型であり、$j_!j^{-1}\mathcal{F}$ の茎は零である。ここでは補題
\ref{etale-cohomology-lemma-stalk-pushforward-closed-immersion} と命題
\ref{etale-cohomology-proposition-describe-jshriek} の茎の記述を用いた。
\end{proof}

\begin{lemma}
\label{etale-cohomology-lemma-compatible-shriek-push-finite}
スキームのファイバー積図式
$$
\xymatrix{
U \ar[d]_g \ar[r]_{j'} & X \ar[d]^f \\
V \ar[r]^j & Y
}
$$
を考える。ここで $f$ は有限で、$j$ は開埋め込みである。このとき関手として
$f_* \circ j'_! = j_! \circ g_*$ が
$\textit{Ab}(U_\etale) \to \textit{Ab}(Y_\etale)$ 上で成り立つ。
\end{lemma}

\begin{proof}
$\mathcal{F}$ を $\textit{Ab}(U_\etale)$ の対象とする。
$\overline{y}$ を $Y$ の幾何学的点とし、開部分 $V$ に含まれないものとする。このとき
$$
(f_*j'_!\mathcal{F})_{\overline{y}} =
\bigoplus\nolimits_{\overline{x},\ f(\overline{x}) = \overline{y}}
(j'_!\mathcal{F})_{\overline{x}} = 0
$$
となる。これは命題 \ref{etale-cohomology-proposition-finite-higher-direct-image-zero} と、
$j'_!\mathcal{F}$ の $\overline{x} \not \in U$ における茎が補題
\ref{etale-cohomology-lemma-jshriek-open} により零であることによる。一方、
$$
j^{-1}f_*j'_!\mathcal{F} =
g_*(j')^{-1}j'_!\mathcal{F} =
g_*\mathcal{F}
$$
が補題 \ref{etale-cohomology-lemma-finite-pushforward-commutes-with-base-change} と
\ref{etale-cohomology-lemma-jshriek-open} により成り立つ。したがって、補題
\ref{etale-cohomology-lemma-jshriek-open} による $j_!$ の特徴づけから
$f_*j'_!\mathcal{F} = j_!g_*\mathcal{F}$ を得る。この同一視が
$\mathcal{F}$ に関して関手的であることの確認は省略する。
\end{proof}







\section{構成可能層}
\label{etale-cohomology-section-constructible}

\noindent
スキーム $X$ をとる。$X$ の {\it 構成可能局所閉部分スキーム}とは、
局所閉部分スキーム $T \subset X$ であって、$T$ の基礎位相空間が
$X$ の構成可能部分集合であるものをいう。$T, T' \subset X$ が同じ基礎位相空間をもつ
局所閉部分スキームならば、エタールサイトの位相不変性（定理
\ref{etale-cohomology-theorem-topological-invariance}）により $T_\etale \cong T'_\etale$ である。
したがって、以下の定義では局所閉部分スキームを被約であると仮定してよい。

\begin{definition}
\label{etale-cohomology-definition-constructible}
スキーム $X$ をとる。
\begin{enumerate}
\item $X_\etale$ 上の集合の層が {\it 構成可能}であるとは、任意のアフィン開集合
$U \subset X$ に対して、$U$ の構成可能局所閉部分スキームへの有限分解
$U = \coprod_i U_i$ で、$\mathcal{F}|_{U_i}$ がすべての $i$ について
有限局所定数となるものが存在することをいう。
\item $X_\etale$ 上のアーベル群の層が {\it 構成可能}であるとは、任意のアフィン開集合
$U \subset X$ に対して、$U$ の構成可能局所閉部分スキームへの有限分解
$U = \coprod_i U_i$ で、$\mathcal{F}|_{U_i}$ がすべての $i$ について
有限局所定数となるものが存在することをいう。
\item $\Lambda$ を Noetherian 環とする。$\Lambda$-加群の層で $X_\etale$ 上のものが
{\it 構成可能}であるとは、任意のアフィン開集合 $U \subset X$ に対して、
$U$ の構成可能局所閉部分スキームへの有限分解 $U = \coprod_i U_i$ で、
$\mathcal{F}|_{U_i}$ がすべての $i$ について有限型かつ局所定数となるものが
存在することをいう。
\end{enumerate}
\end{definition}

\noindent
これは標準的に受け入れられている定義と思われる。スキームのエタールサイトを
越える一般化には、次の別の定義の方が適しているかもしれない。すなわち、
$h_U$, $j_{U!}\mathbf{Z}/n\mathbf{Z}$, $j_{U!}\underline{\Lambda}$ から出発する。
ここで $U$ は $X_\etale$ の準コンパクトかつ準分離な対象をすべて動く。次いで、
これらを含み有限極限と有限余極限について閉じている
$\Sh(X_\etale)$, $\textit{Ab}(X_\etale)$, および
$\textit{Mod}(X_\etale, \underline{\Lambda})$ の最小の充満部分圏をとることにより、
構成可能層を定義する方法である。補題 \ref{etale-cohomology-lemma-constructible-abelian} と補題
\ref{etale-cohomology-lemma-category-constructible-sets},
\ref{etale-cohomology-lemma-category-constructible-abelian},
\ref{etale-cohomology-lemma-category-constructible-modules} から、$X$ が準コンパクトかつ準分離ならば、
この方法で同じ圏が得られることが従う。しかし一般には同じ圏は得られない。

\medskip\noindent
スキームの局所閉部分スキームへの非交和分解 $U = \coprod U_i$ を
$U$ の {\it 分割}と呼ぶ（Topology, 節 \ref{topology-section-stratifications} と比較せよ）。

\begin{lemma}
\label{etale-cohomology-lemma-constructible-quasi-compact-quasi-separated}
準コンパクトかつ準分離なスキーム $X$ をとり、$\mathcal{F}$ を
$X_\etale$ 上の集合の層とする。次の条件は同値である。
\begin{enumerate}
\item $\mathcal{F}$ は構成可能である。
\item 開被覆 $X = \bigcup U_i$ で、$\mathcal{F}|_{U_i}$ が構成可能となるものが存在する。
\item 構成可能局所閉部分スキームによる分割 $X = \bigcup X_i$ で、
$\mathcal{F}|_{X_i}$ が有限局所定数となるものが存在する。
\end{enumerate}
$\Lambda$ が Noetherian のとき、アーベル層と $\Lambda$-加群の層についても
同様の主張が成り立つ。
\end{lemma}

\begin{proof}
(1) が (2) を導くことは明らかである。

\medskip\noindent
(2) を仮定する。任意の $x \in X$ に対して、ある $i$ と $V_x \subset U_i$ があり、
後者は $x$ のアフィン開近傍である。したがって有限アフィン開被覆 $X = \bigcup V_j$ で、
各 $j$ に対して局所閉構成可能部分集合への有限分解
$V_j = \coprod V_{j, k}$ が存在し、$\mathcal{F}|_{V_{j, k}}$ が有限局所定数と
なるものをとれる。Topology, 補題
\ref{topology-lemma-quasi-compact-open-immersion-constructible-image} により、
各 $V_{j, k}$ は $X$ の部分集合として構成可能である。Topology, 補題
\ref{topology-lemma-constructible-partition-refined-by-stratification} により、
構成可能局所閉な層からなる有限層状化 $X = \coprod X_l$ で、各 $V_{j, k}$ が
$X_l$ の合併となるものをとれる。よって (3) が成り立つ。

\medskip\noindent
(3) を仮定する。$U \subset X$ をアフィン開集合とする。このとき $U \cap X_i$ は
$U$ の構成可能局所閉部分集合であり（たとえば Properties, 補題
\ref{properties-lemma-locally-constructible} による）、$U = \coprod U \cap X_i$ は
定義 \ref{etale-cohomology-definition-constructible} の意味で $U$ の分割である。よって (1) が成り立つ。
\end{proof}

\begin{lemma}
\label{etale-cohomology-lemma-constructible-constructible}
準コンパクトかつ準分離なスキーム $X$ をとる。$\mathcal{F}$ を、集合、アーベル群、
または $\Lambda$-加群（$\Lambda$ は Noetherian）の層で $X_\etale$ 上のものとする。
構成可能局所閉部分スキーム $T_i \subset X$ で、(a) $X = \bigcup T_j$ かつ
(b) $\mathcal{F}|_{T_j}$ が構成可能となるものが存在するならば、
$\mathcal{F}$ は構成可能である。
\end{lemma}

\begin{proof}
まず、$X$ はスペクトル位相について準コンパクトなので、被覆は有限であると
仮定してよい（Topology, 補題
\ref{topology-lemma-constructible-hausdorff-quasi-compact} および Properties, 補題
\ref{properties-lemma-quasi-compact-quasi-separated-spectral}）。各 $T_i$ はそれ自体で
準コンパクトかつ準分離なスキームであることに注意する（$X$ 内で構成可能だからである。
詳細は省略する）。したがって、$T_i = \coprod T_{i, j}$ という $T_i$ の局所閉構成可能
部分への有限分割で、$\mathcal{F}|_{T_{i, j}}$ が有限局所定数となるものを
とれる（補題 \ref{etale-cohomology-lemma-constructible-quasi-compact-quasi-separated}）。Topology, 補題
\ref{topology-lemma-constructible-in-constructible} により、$T_{i, j}$ は $X$ の
構成可能局所閉部分スキームである。そこで $X = \bigcup T_{i, j}$ に Topology, 補題
\ref{topology-lemma-constructible-partition-refined-by-stratification} を適用すれば、
$X$ の所望の分割が得られる。
\end{proof}

\begin{lemma}
\label{etale-cohomology-lemma-constructible-local}
スキーム $X$ をとる。集合、アーベル群、または $\Lambda$-加群
（$\Lambda$ は Noetherian）の層の構成可能性は、$X$ 上 Zariski 局所的に
判定できる。
\end{lemma}

\begin{proof}
主張は、$X = \bigcup U_i$ が開被覆で、$\mathcal{F}|_{U_i}$ が構成可能ならば、
$\mathcal{F}$ も構成可能であるという意味である。$U \subset X$ がアフィン開集合ならば、
$U = \bigcup U \cap U_i$ であり、$\mathcal{F}|_{U \cap U_i}$ は構成可能である
（構成可能層の開集合への制限が構成可能であることは明らかである）。補題
\ref{etale-cohomology-lemma-constructible-quasi-compact-quasi-separated} により、
$\mathcal{F}|_U$ は構成可能である。すなわち、$U$ の適切な分割が存在する。
\end{proof}

\begin{lemma}
\label{etale-cohomology-lemma-pullback-constructible}
スキームの射 $f : X \to Y$ をとる。$\mathcal{F}$ が集合、アーベル群、または
$\Lambda$-加群（$\Lambda$ は Noetherian）の構成可能層で $Y_\etale$ 上のものならば、
$f^{-1}\mathcal{F}$ についても $X_\etale$ 上で同じことが成り立つ。
\end{lemma}

\begin{proof}
補題 \ref{etale-cohomology-lemma-constructible-local} により、$X$ と $Y$ がアフィンである場合に
帰着する。補題 \ref{etale-cohomology-lemma-constructible-quasi-compact-quasi-separated} により、
$X$ の構成可能局所閉部分スキームへの有限分割で、各部分上で
$f^{-1}\mathcal{F}$ が有限局所定数となるものを見つければ十分である。
これを得るには、$Y$ の、$\mathcal{F}$ に適合する分割を引き戻し、補題
\ref{etale-cohomology-lemma-pullback-locally-constant} を用いればよい。
\end{proof}

\begin{lemma}
\label{etale-cohomology-lemma-constructible-abelian}
スキーム $X$ をとる。
\begin{enumerate}
\item 集合の構成可能層の圏は、$\Sh(X_\etale)$ の中で有限極限と有限余極限に
ついて閉じている。
\item 構成可能アーベル層の圏は、$\textit{Ab}(X_\etale)$ の弱 Serre 部分圏である。
\item $\Lambda$ を Noetherian 環とする。構成可能な $\Lambda$-加群層で
$X_\etale$ 上のものの圏は、$\textit{Mod}(X_\etale, \Lambda)$ の弱 Serre 部分圏である。
\end{enumerate}
\end{lemma}

\begin{proof}
(3) を証明する。Homology, 補題
\ref{homology-lemma-characterize-weak-serre-subcategory} の判定法を用いる。
$\varphi : \mathcal{F} \to \mathcal{G}$ を構成可能な $\Lambda$-加群層の射とする。
$\mathcal{K} = \Ker(\varphi)$ と $\mathcal{Q} = \Coker(\varphi)$ が構成可能であることを
示さなければならない。同様に、$0 \to \mathcal{F} \to \mathcal{E} \to \mathcal{G} \to 0$
を $\Lambda$-加群層の短完全列で、$\mathcal{F}$, $\mathcal{G}$ が構成可能なものとする。
$\mathcal{E}$ が構成可能であることを示さなければならない。いずれの場合も、$X$ を
アフィン開被覆の各要素で置き換えられるので、$X$ はアフィンであると仮定してよい。
さらに $X$ を、$X$ の構成可能局所閉部分スキームによる有限分割の各部分で置き換えてよい。
その各部分上で $\mathcal{F}$ と $\mathcal{G}$ は有限型かつ局所定数である。したがって補題
\ref{etale-cohomology-lemma-kernel-finite-locally-constant} を適用して結論を得る。

\medskip\noindent
(1) と (2) の証明もほぼ同じなので省略する。
\end{proof}

\begin{lemma}
\label{etale-cohomology-lemma-support-constructible}
準コンパクトかつ準分離なスキーム $X$ をとる。
\begin{enumerate}
\item $\mathcal{F} \to \mathcal{G}$ を $X_\etale$ 上の集合の構成可能層の射とする。
点 $x \in X$ で、$\mathcal{F}_{\overline{x}} \to \mathcal{G}_{\overline{x}}$ が全射、
または単射、または与えられた集合の写像と同型になるものの集合は、$X$ 内で構成可能である。
\item $\mathcal{F}$ を $X_\etale$ 上の構成可能アーベル層とする。
$\mathcal{F}$ の台は構成可能である。
\item $\Lambda$ を Noetherian 環とする。$\mathcal{F}$ を構成可能な $\Lambda$-加群層で
$X_\etale$ 上のものとする。$\mathcal{F}$ の台は構成可能である。
\end{enumerate}
\end{lemma}

\begin{proof}
(1) を証明する。$X = \coprod X_i$ を、局所閉構成可能部分スキームによる $X$ の分割で、
各部分上で $\mathcal{F}$ と $\mathcal{G}$ がともに有限局所定数となるものとする
（$\mathcal{F}$ と $\mathcal{G}$ のそれぞれに補題
\ref{etale-cohomology-lemma-constructible-quasi-compact-quasi-separated} を用い、共通細分をとる）。
各部分への射の制限に補題 \ref{etale-cohomology-lemma-morphism-locally-constant} を適用すればよい。

\medskip\noindent
(2) と (3) の証明は省略する。
\end{proof}

\noindent
次の補題は後で非常に有用になる。大まかには、構成可能層の圏が一種の弱い
「Noether 性」をもつことを述べている。

\begin{lemma}
\label{etale-cohomology-lemma-colimit-constructible}
準コンパクトかつ準分離なスキーム $X$ をとる。
$\mathcal{F} = \colim_{i \in I} \mathcal{F}_i$ を、集合の層、アーベル層、または
加群の層のフィルター余極限とする。
\begin{enumerate}
\item $\mathcal{F}$ と $\mathcal{F}_i$ が集合の構成可能層ならば、ind-対象
$\mathcal{F}_i$ は値 $\mathcal{F}$ をもつ本質的定数対象である。
\item $\mathcal{F}$ と $\mathcal{F}_i$ がアーベル群の構成可能層ならば、ind-対象
$\mathcal{F}_i$ は値 $\mathcal{F}$ をもつ本質的定数対象である。
\item $\Lambda$ を Noetherian 環とする。$\mathcal{F}$ と $\mathcal{F}_i$ が
$\Lambda$-加群の構成可能層ならば、ind-対象 $\mathcal{F}_i$ は値
$\mathcal{F}$ をもつ本質的定数対象である。
\end{enumerate}
\end{lemma}

\begin{proof}
(1) を証明する。構成可能層の有限極限と有限余極限が構成可能であること
（補題 \ref{etale-cohomology-lemma-kernel-finite-locally-constant}）を、以下断りなく用いる。
各 $i$ に対し、$T_i \subset X$ を、点 $x \in X$ で
$\mathcal{F}_{i, \overline{x}} \to \mathcal{F}_{\overline{x}}$ が全射でないものの
集合とする。$\mathcal{F}_i$ と $\mathcal{F}$ は構成可能なので、
$T_i$ は $X$ の構成可能部分集合である（補題 \ref{etale-cohomology-lemma-support-constructible}）。
$\mathcal{F}$ の茎は有限であり、$\mathcal{F} = \colim_{i \in I} \mathcal{F}_i$ なので、
すべての $x \in X$ に対し $x \not \in T_i$ となるのは $i$ が十分大きいときである。
Properties, 補題 \ref{properties-lemma-quasi-compact-quasi-separated-spectral} により
$X$ はスペクトル空間であり、Topology, 補題
\ref{topology-lemma-constructible-hausdorff-quasi-compact} により $X$ の構成可能位相は
準コンパクトである。したがって $T_i = \emptyset$ となるのは $i$ が十分大きいときである。
よって $\mathcal{F}_i \to \mathcal{F}$ は $i$ が十分大きいとき全射である。
ここで、$\mathcal{F}_i \to \mathcal{F}$ がすべての $i$ について全射であると仮定する。
$i \in I$ を選ぶ。$i' \geq i$ に対し、$S_{i'} \subset X$ を、点 $x$ で
$\Im(\mathcal{F}_{i, \overline{x}} \to \mathcal{F}_{\overline{x}})$ の元の個数と
$\Im(\mathcal{F}_{i, \overline{x}} \to \mathcal{F}_{i', \overline{x}})$ の元の個数が
等しくないものの集合とする。$\mathcal{F}_i$, $\mathcal{F}_{i'}$, $\mathcal{F}$ は
構成可能なので、$S_{i'}$ は $X$ の構成可能部分集合である
（詳細は省略する。ヒント：補題 \ref{etale-cohomology-lemma-support-constructible} を用いる）。
$\mathcal{F}_i$ と $\mathcal{F}$ の茎は有限であり、
$\mathcal{F} = \colim_{i' \geq i} \mathcal{F}_{i'}$ なので、すべての $x \in X$ に対し
$x \not \in S_{i'}$ となるのは $i'$ が十分大きいときである。上と同じ議論により、
十分大きい $i'$ で $S_{i'} = \emptyset$ となるものをとれる。したがって、所望のとおり
$\mathcal{F}_i \to \mathcal{F}_{i'}$ は $\mathcal{F}$ を経由する。

\medskip\noindent
(2) を証明する。構成可能アーベル層は集合の構成可能層でもあることに注意すれば、
(2) は (1) から従う。

\medskip\noindent
(3) を証明する。構成可能な $\Lambda$-加群層の圏がアーベル圏であること
（補題 \ref{etale-cohomology-lemma-kernel-finite-locally-constant}）を、以下断りなく用いる。
各 $i$ に対し、$\mathcal{Q}_i$ を射 $\mathcal{F}_i \to \mathcal{F}$ の余核とする。
その台 $T_i$、すなわち $\mathcal{Q}_i$ の台は $X$ の構成可能部分集合である。
実際、$\mathcal{Q}_i$ は構成可能である（補題 \ref{etale-cohomology-lemma-support-constructible}）。
$\mathcal{F}$ の茎は有限
$\Lambda$-加群であり、$\mathcal{F} = \colim_{i \in I} \mathcal{F}_i$ なので、すべての
$x \in X$ に対し $x \not \in T_i$ となるのは $i$ が十分大きいときである。Properties, 補題
\ref{properties-lemma-quasi-compact-quasi-separated-spectral} により $X$ はスペクトル空間であり、
Topology, 補題 \ref{topology-lemma-constructible-hausdorff-quasi-compact} により
$X$ の構成可能位相は準コンパクトである。したがって $T_i = \emptyset$ となるのは
$i$ が十分大きいときである。これで第一の主張が証明された。第二の主張については、
$\mathcal{F}_i \to \mathcal{F}$ がすべての $i$ について全射であると仮定する。
$i \in I$ を選ぶ。$i' \geq i$ に対し、$\mathcal{K}_{i'}$ を
$\Ker(\mathcal{F}_i \to \mathcal{F})$ の $\mathcal{F}_{i'}$ における像とする。
その台 $S_{i'}$ は、$\mathcal{K}_{i'}$ が構成可能なので $X$ の構成可能部分集合である。
$\mathcal{K}_{i'}$ は $\Ker(\mathcal{F}_i \to \mathcal{F})$ から来る。その茎は有限
$\Lambda$-加群であり、
$\mathcal{F} = \colim_{i' \geq i} \mathcal{F}_{i'}$ なので、すべての $x \in X$ に対し
$x \not \in S_{i'}$ となるのは $i'$ が十分大きいときである。上と同じ議論により、
十分大きい $i'$ で $S_{i'} = \emptyset$ となるものをとれる。したがって、所望のとおり
$\mathcal{F}_i \to \mathcal{F}_{i'}$ は $\mathcal{F}$ を経由する。
\end{proof}

\begin{lemma}
\label{etale-cohomology-lemma-tensor-product-constructible}
スキーム $X$ と Noetherian 環 $\Lambda$ をとる。構成可能な $\Lambda$-加群の二つの層で
$X_\etale$ 上のもののテンソル積は、構成可能な $\Lambda$-加群層である。
\end{lemma}

\begin{proof}
問題は直ちに $X$ がアフィンである場合に帰着する。$X$ の構成可能局所閉な層による
任意の二つの分割は同じ型の共通細分をもち、引き戻しはテンソル積と可換なので、
補題 \ref{etale-cohomology-lemma-tensor-product-locally-constant} に帰着する。
\end{proof}

\begin{lemma}
\label{etale-cohomology-lemma-tensor-constructible}
$\Lambda \to \Lambda'$ を Noetherian 環の準同型とする。スキーム $X$ をとり、
$\mathcal{F}$ を構成可能な $\Lambda$-加群層で $X_\etale$ 上のものとする。このとき
$\mathcal{F} \otimes_{\underline{\Lambda}} \underline{\Lambda'}$ は構成可能な
$\Lambda'$-加群層である。
\end{lemma}

\begin{proof}
省略する。ヒント：アフィン局所的には同じ層状化を用いればよい。
\end{proof}










\section{射に関する補助補題}
\label{etale-cohomology-section-stratify-morphisms}

\noindent
構成可能層の関手性に関する性質を証明する際に有用ないくつかの補題を述べる。

\begin{lemma}
\label{etale-cohomology-lemma-etale-stratified-finite}
$U \to X$ を準コンパクトかつ準分離なスキームのエタール射とする
（たとえば Noetherian スキームのエタール射）。このとき、構成可能局所閉部分スキーム
による分割 $X = \coprod_i X_i$ で、$X_i \times_X U \to X_i$ がすべての $i$ について
有限エタールとなるものが存在する。
\end{lemma}

\begin{proof}
$U \to X$ が分離的ならば、これは More on Morphisms, 補題
\ref{more-morphisms-lemma-stratify-flat-fp-qf} である。一般の場合、$X$ はアフィンであると
仮定してよい。有限アフィン開被覆 $U = \bigcup U_j$ を選ぶ。すべての射
$U_j \to X$ と $U_j \cap U_{j'} \to X$ に前の場合を適用し、得られる分割の共通細分
$X = \coprod X_i$ を選ぶ。さらに分割を細分して、$X_i$ もアフィンであると仮定してよい。
$i$ を固定し、$V = U \times_X X_i$ とおく。射
$V_j = U_j \times_X X_i \to X_i$ と
$V_{jj'} = (U_j \cap U_{j'}) \times_X X_i \to X_i$ は有限エタールである。
したがって $V_j$ と $V_{jj'}$ はアフィンスキームであり、$V_{jj'} \subset V_j$ は
開かつ閉である（$V_{jj'} \to X_i$ は固有なので、Morphisms, 補題
\ref{morphisms-lemma-image-proper-scheme-closed} を適用できる）。さらに
$V = \bigcup V_j$ は分離的である。実際、
$\mathcal{O}(V_j) \to \mathcal{O}(V_{jj'})$ は全射である。Schemes, 補題
\ref{schemes-lemma-characterize-separated} を参照せよ。よって前の場合を $V \to X_i$ に
適用でき、必要なら分割をさらに細分できる（実際には不要だが、そのことは用いない）。
\end{proof}

\noindent
Noetherian の場合には、Noether 帰納法と次の興味深い補題を用いて前の補題を証明できる。

\begin{lemma}
\label{etale-cohomology-lemma-generically-finite}
$f: X \to Y$ を準コンパクト、準分離、かつ局所有限型なスキームの射とする。
$\eta$ が $Y$ のある既約成分の生成点で、$f^{-1}(\eta)$ が有限ならば、
開集合 $V \subset Y$ で、$\eta$ を含み $f^{-1}(V) \to V$ が有限となるものが存在する。
\end{lemma}

\begin{proof}
これは Morphisms, 補題 \ref{morphisms-lemma-generically-finite} である。
\end{proof}

\noindent
次の補題の主張は少し強めることができる。

\begin{lemma}
\label{etale-cohomology-lemma-decompose-quasi-finite-morphism}
$f : Y \to X$ をアフィンスキームの準有限かつ有限表示な射とする。
\begin{enumerate}
\item アフィンスキームの全射 $X' \to X$ と閉部分スキーム
$Z' \subset Y' = X' \times_X Y$ で、次を満たすものが存在する。
\begin{enumerate}
\item $Z' \subset Y'$ は肥厚化である。
\item $Z' \to X'$ は有限エタール射である。
\end{enumerate}
\item 局所閉、構成可能、かつアフィンな層による有限分割 $X = \coprod X_i$ と、
全射な有限局所自由射 $X'_i \to X_i$ で、
$Y'_i = X'_i \times_X Y \to X'_i$ の被約化が
$\coprod_{j = 1}^{n_i} (X'_i)_{red} \to (X'_i)_{red}$ と、ある $n_i$ に対して
同型になるものが存在する。
\end{enumerate}
\end{lemma}

\begin{proof}
$X' = \coprod X'_i$ とおけば (2) が (1) を導くことがわかる。
$X = \Spec(A)$, $Y = \Spec(B)$ と書く。$A$ を有限型 $\mathbf{Z}$-代数 $A_i$ の
フィルター余極限として書く。$B$ は有限表示な $A$-代数なので、ある $0 \in I$ と
有限型な環準同型 $A_0 \to B_0$ が存在して、$B = \colim B_i$ かつ
$B_i = A_i \otimes_{A_0} B_0$ となる。Algebra, 補題
\ref{algebra-lemma-colimit-category-fp-algebras} を参照せよ。十分大きい $i$ に対して
$A_i \to B_i$ は準有限である。Limits, 補題
\ref{limits-lemma-descend-quasi-finite} を参照せよ。したがって $\mathbf{Z}$ 上有限型な
代数の場合、特に Noetherian の場合に帰着する（詳細は省略する）。

\medskip\noindent
$X$ と $Y$ が Noetherian であると仮定する。この場合、$X$ の任意の局所閉部分集合は
構成可能である。補題 \ref{etale-cohomology-lemma-generically-finite} と Noether 帰納法により、有限分割
$X = \coprod X_i$、すなわち $X$ の局所閉な層による分割で、$Y \times_X X_i \to X_i$ が有限となる
ものが存在する。この分割を細分してアフィンな層を得られる。したがって $X$ を
$X' = \coprod X_i$ で置き換えれば、$Y \to X$ は有限であると仮定してよい。

\medskip\noindent
$X$ と $Y$ は Noetherian であり、$Y \to X$ は有限であると仮定する。
全射、平坦、準有限な射 $U \to X$ による基底変換後に (2) を証明できるとする。
すると分割 $U = \coprod U_i$ と有限局所自由射 $U'_i \to U_i$ があって、
$U'_i \times_X Y \to U'_i$ は
$\coprod_{j = 1}^{n_i} (U'_i)_{red} \to (U'_i)_{red}$ と、ある $n_i$ に対して同型である。
前段落の議論により、局所閉アフィンな層による分割 $X = \coprod X_j$ で、
$X_j \times_X U_i \to X_j$ がすべての $i, j$ について有限となるものをとれる。
Morphisms, 補題 \ref{morphisms-lemma-finite-flat} により、各
$X_j \times_X U_i \to X_j$ は有限局所自由である。したがって
$X_j \times_X U'_i \to X_j$ も有限局所自由である（Morphisms, 補題
\ref{morphisms-lemma-composition-finite-locally-free}）。よって
$X = \coprod X_j$ と $X_j' = \coprod_i X_j \times_X U'_i$ は $Y \to X$ に対する解を
与える。したがって（Noetherian の場合には）全射平坦準有限な基底変換後に
結果を証明すれば十分である。

\medskip\noindent
Morphisms, 補題 \ref{morphisms-lemma-massage-finite} を適用すると、$Y$ はアフィンスキーム
$Z$ の閉部分スキームであり、$Z = \bigcup_{i \in I} Z_i$ という、$X$ へ同型に写る
閉部分スキームの有限合併に集合論的に等しいと仮定してよい。この場合には、機能する
有限分割 $X = \coprod X_j$ をアフィン局所閉な層によって見つける
（言い換えれば $X'_j = X_j$）。$T_i = Y \cap Z_i$ とおく。これは $X$ の閉部分スキームである。
$X$ は Noetherian なので、アフィン局所閉部分スキームによる有限分割
$X = \coprod X_j$ で、各 $X_j \times_X T_i$ が集合論的に層
$X_j \times_X Z_i$ の合併となるものをとれる。$X$ を $X_j$ で置き換えることにより、
$I = I_1 \amalg I_2$ であり、$Z_i \subset Y$ が $i \in I_1$ に対して成り立ち、
$Z_i \cap Y = \emptyset$ が $i \in I_2$ に対して成り立つと仮定してよい。
$Z$ を $\bigcup_{i \in I_1} Z_i$ で置き換えることにより、$Y = Z$ と仮定してよい。
最後に、再び $X$ を上のような分割の各部分で置き換えて、任意の
$i, i' \subset I$ に対し交わり $Z_i \cap Z_{i'}$ が空であるか、または集合論的に
$Z_i$ および $Z_{i'}$ と等しいと仮定してよい。これは明らかに、$Y$ が
$Z_i$ の非交和に集合論的に等しいことを意味し、示すべきことが従う。
\end{proof}








\section{構成可能層についてさらに}
\label{etale-cohomology-section-more-constructible}

\noindent
$\Lambda$ を Noether 環とし、$X$ をスキームとする。
しばしば $X_\etale$ を、環の層 $\underline{\Lambda}$ を備えた
環付きサイトとみなす。アーベル層を扱う場合には、
$\Lambda = \mathbf{Z}/n\mathbf{Z}$ と取ることが多い。ここで $n$ は適当な整数である。

\begin{lemma}
\label{etale-cohomology-lemma-jshriek-constructible}
準コンパクトかつ準分離なスキームのエタール射
$j : U \to X$ を取る。
\begin{enumerate}
\item 層 $h_U$ は集合の構成可能層である。
\item 層 $j_!\underline{M}$ は、有限アーベル群 $M$ に対して
構成可能アーベル層である。
\item $\Lambda$ が Noether 環で $M$ が有限 $\Lambda$-加群ならば、
$j_!\underline{M}$ は $\Lambda$-加群の構成可能層であり、$X_\etale$ 上にある。
\end{enumerate}
\end{lemma}

\begin{proof}
補題 \ref{etale-cohomology-lemma-etale-stratified-finite} により、分割 $\coprod_i X_i$ が存在し、
$\pi_i : j^{-1}(X_i) \to X_i$ は有限エタールとなる。$h_U$ の $X_i$ への制限は
$h_{j^{-1}(X_i)}$ であり、補題
\ref{etale-cohomology-lemma-characterize-finite-locally-constant} により有限局所定数である。
(2) と (3) については、補題 \ref{etale-cohomology-lemma-shriek-base-change} および
\ref{etale-cohomology-lemma-shriek-equals-star-finite-etale} により
$$
j_!(\underline{M})|_{X_i} =
\pi_{i!}(\underline{M}) =
\pi_{i*}(\underline{M})
$$
であることに注意する。したがって、有限エタール射
$\pi : Y \to X$ について補題を示せば十分である。
これは補題 \ref{etale-cohomology-lemma-pushforward-locally-constant} である。
\end{proof}

\begin{lemma}
\label{etale-cohomology-lemma-torsion-colimit-constructible}
$X$ を準コンパクトかつ準分離なスキームとする。
\begin{enumerate}
\item $\mathcal{F}$ を $X_\etale$ 上の集合の層とする。
このとき $\mathcal{F}$ は集合の構成可能層のフィルター余極限である。
\item $\mathcal{F}$ を $X_\etale$ 上のねじれアーベル層とする。
このとき $\mathcal{F}$ は構成可能アーベル層のフィルター余極限である。
\item $\Lambda$ を Noether 環とし、$\mathcal{F}$ を $\Lambda$-加群の層であって
$X_\etale$ 上にあるものとする。このとき $\mathcal{F}$ は
$\Lambda$-加群の構成可能層のフィルター余極限である。
\end{enumerate}
\end{lemma}

\begin{proof}
$\mathcal{B}$ を $X_\etale$ の準コンパクトかつ準分離な対象全体とする。
Modules on Sites の補題
\ref{sites-modules-lemma-module-filtered-colimit-constructibles} により、
任意の集合の層は次の形の層のフィルター余極限である：
$$
\text{余等化子}\left(
\xymatrix{
\coprod\nolimits_{j = 1, \ldots, m} h_{V_j}
\ar@<1ex>[r] \ar@<-1ex>[r] &
\coprod\nolimits_{i = 1, \ldots, n} h_{U_i}
}
\right)
$$
ここで $V_j$ と $U_i$ は $X_\etale$ の準コンパクトかつ準分離な対象である。
補題 \ref{etale-cohomology-lemma-jshriek-constructible} および
\ref{etale-cohomology-lemma-constructible-abelian} により、これらの余等化子は構成可能である。
これで (1) が示された。

\medskip\noindent
$\Lambda$ を Noether 環とする。Modules on Sites の補題
\ref{sites-modules-lemma-module-filtered-colimit-constructibles} により、
$\Lambda$-加群の層 $\mathcal{F}$ は次の形の加群のフィルター余極限である：
$$
\Coker\left(
\bigoplus\nolimits_{j = 1, \ldots, m} j_{V_j!}\underline{\Lambda}_{V_j}
\longrightarrow
\bigoplus\nolimits_{i = 1, \ldots, n} j_{U_i!}\underline{\Lambda}_{U_i}
\right)
$$
ここで $V_j$ と $U_i$ は $X_\etale$ の準コンパクトかつ準分離な対象である。
補題 \ref{etale-cohomology-lemma-jshriek-constructible} および
\ref{etale-cohomology-lemma-constructible-abelian} により、これらの余核は構成可能である。
これで (3) が示された。

\medskip\noindent
(2) を証明する。まず $\mathcal{F} = \bigcup \mathcal{F}[n]$ と書く。
ここで $\mathcal{F}[n]$ は $n$-ねじれ部分層である。このとき
$\mathcal{F}[n]$ を $\mathbf{Z}/n\mathbf{Z}$-加群の層とみなし、
(3) を適用できる。
\end{proof}

\begin{lemma}
\label{etale-cohomology-lemma-check-constructible}
準コンパクトかつ準分離なスキームの全射
$f : X \to Y$ を取る。
\begin{enumerate}
\item $\mathcal{F}$ を $Y_\etale$ 上の集合の層とする。このとき
$\mathcal{F}$ が構成可能であることと $f^{-1}\mathcal{F}$ が構成可能であることは同値である。
\item $\mathcal{F}$ を $Y_\etale$ 上のアーベル層とする。このとき
$\mathcal{F}$ が構成可能であることと $f^{-1}\mathcal{F}$ が構成可能であることは同値である。
\item $\Lambda$ を Noether 環とし、$\mathcal{F}$ を $\Lambda$-加群の層であって
$Y_\etale$ 上にあるものとする。このとき $\mathcal{F}$ が構成可能であることと
$f^{-1}\mathcal{F}$ が構成可能であることは同値である。
\end{enumerate}
\end{lemma}

\begin{proof}
一方の含意は補題 \ref{etale-cohomology-lemma-pullback-constructible} から従う。
逆に、$f^{-1}\mathcal{F}$ が構成可能であると仮定する。
$\mathcal{F} = \colim \mathcal{F}_i$ と書く。これは補題
\ref{etale-cohomology-lemma-torsion-colimit-constructible} による、集合、アーベル群、または加群の
構成可能層のフィルター余極限である。$f^{-1}$ は左随伴なので余極限と可換し
（Categories の補題 \ref{categories-lemma-adjoint-exact}）、したがって
$f^{-1}\mathcal{F} = \colim f^{-1}\mathcal{F}_i$ である。補題
\ref{etale-cohomology-lemma-colimit-constructible} により、写像
$f^{-1}\mathcal{F}_i \to f^{-1}\mathcal{F}$ は全射であり、これは十分大きい
すべての $i$ に対して成り立つ。
$f$ は全射なので、補題 \ref{etale-cohomology-lemma-stalk-pullback} と定理
\ref{etale-cohomology-theorem-exactness-stalks} を用いて茎を見れば、写像
$\mathcal{F}_i \to \mathcal{F}$ が全射であり、これは十分大きいすべての $i$ に
対して成り立つと結論できる。
したがって $\mathcal{F}$ はある構成可能層 $\mathcal{G}$ の商である。
同じ議論を $\mathcal{G} \times_\mathcal{F} \mathcal{G}$、または
$\mathcal{G} \to \mathcal{F}$ の核にもう一度適用する。
$f^{-1}$ が完全であること、および集合・アーベル群・加群の構成可能層の圏が、
$\Sh(Y_\etale)$、$\Sh(X_\etale)$、$\textit{Ab}(Y_\etale)$、
$\textit{Ab}(X_\etale)$、$\textit{Mod}(Y_\etale, \Lambda)$、および
$\textit{Mod}(X_\etale, \Lambda)$ の内部で有限（余）極限または（余）核を取る操作により
保たれることから、結論を得る。補題 \ref{etale-cohomology-lemma-constructible-abelian} を参照されたい。
\end{proof}

\begin{lemma}
\label{etale-cohomology-lemma-pushforward-constructible}
$f : X \to Y$ をスキームの有限エタール射とし、$\Lambda$ を Noether 環とする。
$\mathcal{F}$ が集合の構成可能層、アーベル群の構成可能層、または
$\Lambda$-加群の構成可能層であって $X_\etale$ 上にあるならば、
$f_*\mathcal{F}$ についても同じことが成り立ち、これは $Y_\etale$ 上にある。
\end{lemma}

\begin{proof}
補題 \ref{etale-cohomology-lemma-constructible-local} により、これは $Y$ 上 Zariski 局所的に
確認すれば十分であり、補題 \ref{etale-cohomology-lemma-check-constructible} により $Y$ を
エタール被覆で置き換えてよい（$f_*$ の構成はエタール局所化と可換する）。
有限エタール射はエタール局所的に同型射の非交和と同型である。
\'Etale Morphisms の補題 \ref{etale-lemma-finite-etale-etale-local} を参照されたい。
したがって集合の層の場合、本補題は、$\mathcal{F}_i$（$i = 1, \ldots, n$）が
集合の構成可能層ならば $\prod_{i = 1, \ldots, n} \mathcal{F}_i$ もそうである、
という主張になる。これは明らかである。アーベル群の層と加群の層についても同様である。
\end{proof}

\begin{lemma}
\label{etale-cohomology-lemma-category-constructible-sets}
$X$ を準コンパクトかつ準分離なスキームとする。集合の構成可能層の圏は
$\Sh(X_\etale)$ の充満部分圏であり、その対象は次の余等化子となる層
$\mathcal{F}$ である：
$$
\xymatrix{
\mathcal{F}_1
\ar@<1ex>[r] \ar@<-1ex>[r]
&
\mathcal{F}_0 \ar[r]
&
\mathcal{F}}
$$
ここで $\mathcal{F}_i$（$i = 0, 1$）は、形 $h_U$ の層の有限余積であり、
$U$ は $X_\etale$ の準コンパクトかつ準分離な対象である。
\end{lemma}

\begin{proof}
補題 \ref{etale-cohomology-lemma-torsion-colimit-constructible} の証明で、この形の層が
構成可能であることを見た。逆を示すには、集合の任意の構成可能層
$\mathcal{F}$ に対して全射 $\mathcal{F}_0 \to \mathcal{F}$ を取れ、その始域
$\mathcal{F}_0$ が補題に述べた形であると仮定すればよい。このとき、補題
\ref{etale-cohomology-lemma-constructible-abelian} により後者も構成可能なので、全射
$\mathcal{F}_1 \to \mathcal{F}_0 \times_\mathcal{F} \mathcal{F}_0$ が得られる。

\medskip\noindent
Topology の補題
\ref{topology-lemma-constructible-partition-refined-by-stratification} により、
$X = \coprod_{i \in I} X_i$ という有限階層化を選べ、$\mathcal{F}$ は各階層上
有限局所定数となる。$I$ の濃度に関する帰納法で結果を示す。
$i \in I$ を $I$ の半順序における極小元とする。このとき
$X_i \subset X$ は閉である。帰納法により、準コンパクトかつ準分離な対象
$U_\alpha$ が $(X \setminus X_i)_\etale$ に有限個存在し、全射
$\coprod h_{U_\alpha} \to \mathcal{F}|_{X \setminus X_i}$ がある。
これらは写像
$$
\coprod h_{U_\alpha} \to \mathcal{F}
$$
を定め、これは $X \setminus X_i$ に制限すると全射である。補題
\ref{etale-cohomology-lemma-characterize-finite-locally-constant} により、
$\mathcal{F}|_{X_i} = h_V$ となるスキーム $V$ があり、これは $X_i$ 上有限エタールである。
$\overline{v}$ を $V$ の幾何点で、$\overline{x} \in X_i$ 上にあるものとする。
$\overline{v}$ は茎 $\mathcal{F}_{\overline{x}} = V_{\overline{x}}$ の元とみなせる。
したがって、エタール近傍 $(U, \overline{u})$ を $\overline{x}$ に対して取れ、さらに
切断 $s \in \mathcal{F}(U)$ を、その $\overline{x}$ での茎が $\overline{v}$ を
与えるように取れる。$s$ を写像 $s : h_U \to \mathcal{F}$ とみなし、$X_i$ に
制限すると、射 $s|_{X_i} : U \times_X X_i \to V$ を得る。これは $X_i$ 上の射であり、
$\overline{u}$ を $\overline{v}$ に写す。$V$ は準コンパクトである
（$X_i$ 上有限であり、その基底は準コンパクトスキーム $X$ の閉部分スキームである）から、
これらの切断を有限個 $s^{(1)}, \ldots, s^{(m)}$ 取れる。これらは
$\mathcal{F}$ の $U^{(1)}, \ldots, U^{(m)}$ 上の切断であり、共同全射
$$
\coprod s^{(j)}|_{X_i} : \coprod U^{(j)} \times_X X_i \longrightarrow V
$$
が定まる。すると所望の全射
$$
\coprod h_{U_\alpha} \amalg \coprod h_{U^{(j)}} \to \mathcal{F}
$$
を得る。
\end{proof}

\begin{lemma}
\label{etale-cohomology-lemma-category-constructible-modules}
$X$ を準コンパクトかつ準分離なスキームとし、$\Lambda$ を Noether 環とする。
$\Lambda$-加群の構成可能層の圏は、ちょうど次の形の加群の圏である：
$$
\Coker\left(
\bigoplus\nolimits_{j = 1, \ldots, m} j_{V_j!}\underline{\Lambda}_{V_j}
\longrightarrow
\bigoplus\nolimits_{i = 1, \ldots, n} j_{U_i!}\underline{\Lambda}_{U_i}
\right)
$$
ここで $V_j$ と $U_i$ は $X_\etale$ の準コンパクトかつ準分離な対象である。
実際、$U_i$ と $V_j$ はアフィンであると仮定してもよい。
\end{lemma}

\begin{proof}
補題 \ref{etale-cohomology-lemma-torsion-colimit-constructible} の証明で、この形の加群が
構成可能であることを見た。構成可能加群の圏はアーベル圏なので
（補題 \ref{etale-cohomology-lemma-constructible-abelian}）、構成可能加群 $\mathcal{F}$ が
与えられたとき、全射
$$
\bigoplus\nolimits_{i = 1, \ldots, n} j_{U_i!}\underline{\Lambda}_{U_i}
\longrightarrow \mathcal{F}
$$
が存在することを示せば十分である。ここで $U_i$ は $X_\etale$ のアフィン対象である。
Modules on Sites の補題
\ref{sites-modules-lemma-module-filtered-colimit-constructibles} により、全射
$$
\Psi :
\bigoplus\nolimits_{i \in I} j_{U_i!}\underline{\Lambda}_{U_i}
\longrightarrow
\mathcal{F}
$$
が存在する。ここで $U_i$ はアフィンであり、直和は無限かもしれない添字集合
$I$ にわたる。各有限部分集合 $I' \subset I$ に対して
$$
T_{I'} = \text{Supp}(\Coker(
\bigoplus\nolimits_{i \in I'} j_{U_i!}\underline{\Lambda}_{U_i}
\longrightarrow \mathcal{F}))
$$
と置く。構成可能層の定義そのものから、集合 $T_{I'}$ は $X$ の構成可能部分集合である。
$T_{I'} = \emptyset$ となる $I'$ が存在することを示したい。
各茎 $\mathcal{F}_{\overline{x}}$ は有限型 $\Lambda$-加群であり、$\Psi$ は全射なので、
各 $x \in X$ に対してある $I'$ が存在し、$x \not \in T_{I'}$ となる。すなわち
$\emptyset = \bigcap_{I' \subset I\text{ finite}} T_{I'}$ である。
Properties の補題
\ref{properties-lemma-quasi-compact-quasi-separated-spectral} により $X$ は
スペクトル空間であるから、Topology の補題
\ref{topology-lemma-constructible-hausdorff-quasi-compact} により $X$ の
構成可能位相は準コンパクトである。したがって、所望どおり
$T_{I'} = \emptyset$ となる有限部分集合 $I' \subset I$ が存在する。
\end{proof}

\begin{lemma}
\label{etale-cohomology-lemma-category-constructible-abelian}
$X$ を準コンパクトかつ準分離なスキームとする。構成可能アーベル層の圏は、
ちょうど次の形のアーベル層の圏である：
$$
\Coker\left(
\bigoplus\nolimits_{j = 1, \ldots, m}
j_{V_j!}\underline{\mathbf{Z}/m_j\mathbf{Z}}_{V_j}
\longrightarrow
\bigoplus\nolimits_{i = 1, \ldots, n}
j_{U_i!}\underline{\mathbf{Z}/n_i\mathbf{Z}}_{U_i}
\right)
$$
ここで $V_j$ と $U_i$ は $X_\etale$ の準コンパクトかつ準分離な対象であり、
$m_j$、$n_i$ は正整数である。実際、$U_i$ と $V_j$ はアフィンであると
仮定してもよい。
\end{lemma}

\begin{proof}
補題 \ref{etale-cohomology-lemma-category-constructible-modules} を
$\Lambda = \mathbf{Z}/n\mathbf{Z}$ として適用すればよい。ここで $X$ が
準コンパクトなので、任意の構成可能アーベル層はある正整数 $n$ で零化されることを
用いた（詳細は省略する）。
\end{proof}

\begin{lemma}
\label{etale-cohomology-lemma-constructible-is-compact}
$X$ を準コンパクトかつ準分離なスキームとし、$\Lambda$ を Noether 環とする。
$\mathcal{F}$ を、集合、アーベル群、または $\Lambda$-加群の構成可能層であって
$X_\etale$ 上にあるものとする。また $\mathcal{G} = \colim \mathcal{G}_i$ を、
集合、アーベル群、または $\Lambda$-加群の層のフィルター余極限とする。このとき
$$
\Mor(\mathcal{F}, \mathcal{G}) = \colim \Mor(\mathcal{F}, \mathcal{G}_i)
$$
この等式は、$\Lambda$-加群の場合も含め、$X_\etale$ 上の集合、アーベル群、または
加群の層の圏で成り立つ。
\end{lemma}

\begin{proof}
集合の層の場合を考える。補題 \ref{etale-cohomology-lemma-category-constructible-sets} により、
$h_U$ について補題を示せば十分である。ここで $U$ は $X_\etale$ の
準コンパクトかつ準分離な対象である。$\Mor(h_U, \mathcal{G}) = \mathcal{G}(U)$ であることを
思い出せば、結果は Sites の補題 \ref{sites-lemma-directed-colimits-sections} から従う。

\medskip\noindent
アーベル層または加群の層の場合も、補題
\ref{etale-cohomology-lemma-category-constructible-abelian} および
\ref{etale-cohomology-lemma-category-constructible-modules} を用いて同様に従う。
アーベル層の場合には、
$\Mor(j_{U!}\underline{\mathbf{Z}/n\mathbf{Z}}, \mathcal{G})$ が
$n$-ねじれ元全体、すなわち $\mathcal{G}(U)$ のそのような元に等しいことも付け加えておく。
\end{proof}

\begin{lemma}
\label{etale-cohomology-lemma-finite-pushforward-constructible}
$f : X \to Y$ をスキームの有限かつ有限表示な射とし、$\Lambda$ を Noether 環とする。
$\mathcal{F}$ が集合、アーベル群、または $\Lambda$-加群の構成可能層であって
$X_\etale$ 上にあるならば、$f_*\mathcal{F}$ も構成可能である。
\end{lemma}

\begin{proof}
補題 \ref{etale-cohomology-lemma-constructible-local} により、$X$ と $Y$ がアフィンの場合に
証明すれば十分である。補題
\ref{etale-cohomology-lemma-finite-pushforward-commutes-with-base-change} および
\ref{etale-cohomology-lemma-check-constructible} により、$X$ 上全射な任意のアフィンスキームへ
基底変換してよい。補題 \ref{etale-cohomology-lemma-decompose-quasi-finite-morphism} により、
有限エタール射の場合に帰着する（実際、肥厚はエタール・トポス、さらには
小エタール・サイトの同値を導く。定理
\ref{etale-cohomology-theorem-topological-invariance} を参照されたい）。有限エタールの場合は補題
\ref{etale-cohomology-lemma-pushforward-constructible} である。
\end{proof}

\begin{lemma}
\label{etale-cohomology-lemma-category-is-colimit}
$X = \lim_{i \in I} X_i$ を、遷移射がアフィンであるスキームの有向系の極限とする。
$X_i$ は、すべての $i \in I$ に対して準コンパクトかつ準分離であると仮定する。
\begin{enumerate}
\item $X_\etale$ 上の集合の構成可能層の圏は、$(X_i)_\etale$ 上の
集合の構成可能層の圏の余極限である。
\item $X_\etale$ 上の構成可能アーベル層の圏は、$(X_i)_\etale$ 上の
構成可能アーベル層の圏の余極限である。
\item $\Lambda$ を Noether 環とする。$\Lambda$-加群の構成可能層の圏で
$X_\etale$ 上にあるものは、$\Lambda$-加群の構成可能層の圏で
$(X_i)_\etale$ 上にあるものの余極限である。
\end{enumerate}
\end{lemma}

\begin{proof}
(1) を証明する。射影を $f_i : X \to X_i$ と書く。証明は
「忠実」「充満忠実」「本質的全射」に対応する三つの部分からなる。

\medskip\noindent
忠実性。$0 \in I$ を選び、$\mathcal{F}_0$、$\mathcal{G}_0$ を $X_0$ 上の
構成可能層とする。写像 $a, b : \mathcal{F}_0 \to \mathcal{G}_0$ があり、
$f_0^{-1}a = f_0^{-1}b$ を満たすと仮定する。集合 $E \subset X_0$ を、
点 $x \in X_0$ で $a_{\overline{x}} = b_{\overline{x}}$ を満たすものの集合とする。補題
\ref{etale-cohomology-lemma-support-constructible} により部分集合
$E \subset X_0$ は構成可能である。仮定により $X \to X_0$ の像は $E$ に含まれる。
Limits の補題 \ref{limits-lemma-limit-contained-in-constructible} により、
ある $i \geq 0$ が存在し、$X_i \to X_0$ の像が $E$ に含まれる。
したがって $f_{i0}^{-1}a = f_{i0}^{-1}b$ である。

\medskip\noindent
充満忠実性。$0 \in I$ を選び、$\mathcal{F}_0$、$\mathcal{G}_0$ を $X_0$ 上の
構成可能層とする。写像
$a : f_0^{-1}\mathcal{F}_0 \to f_0^{-1}\mathcal{G}_0$ があると仮定する。
ある $i$ と写像
$a_i : f_{i0}^{-1}\mathcal{F}_0 \to f_{i0}^{-1}\mathcal{G}_0$ が存在し、
その引き戻しが $a$ となると主張する。ここで引き戻しは $X$ 上で取る。補題
\ref{etale-cohomology-lemma-category-constructible-sets} により、$\mathcal{F}_0$ を
$(X_0)_\etale$ の準コンパクトかつ準分離な対象によって表される層の
有限余積で置き換えられる。したがって、そのような対象 $U_0 \to X_0$ が
$(X_0)_\etale$ にあるとき、次を示せばよい：
$$
f_0^{-1}\mathcal{G}(U) = \colim_{i \geq 0} f_{i0}^{-1}\mathcal{G}(U_i)
$$
ここで $U = X \times_{X_0} U_0$ かつ $U_i = X_i \times_{X_0} U_0$ である。
これは定理 \ref{etale-cohomology-theorem-colimit} の特別な場合である。

\medskip\noindent
本質的全射性。任意の構成可能層 $\mathcal{F}$ が $X$ 上にあるとき、ある構成可能層について
$f_i^{-1}\mathcal{F}$ と同型であることを示さなければならない。ここで
$\mathcal{F}_i$ は $X_i$ 上の層である。補題
\ref{etale-cohomology-lemma-category-constructible-sets} を適用し、前二段落の結果を用いると、
$h_U$ についてこれを証明すれば十分である。ここで $U$ は $X_\etale$ の
準コンパクトかつ準分離な対象である。この場合、$U$ が $X_i$ 上エタールな
準コンパクトかつ準分離なスキームの基底変換であり、そのような $i$ が存在することを
示せばよい。
これは Limits の補題 \ref{limits-lemma-descend-finite-presentation} および
\ref{limits-lemma-descend-etale} から従う。

\medskip\noindent
(3) を証明する。議論は集合の層の場合とほとんど同じであるが、補題
\ref{etale-cohomology-lemma-category-constructible-modules} を用いる。これは補題
\ref{etale-cohomology-lemma-category-constructible-sets} の代わりである。詳細は省略する。
(2) は (3) から従う。実際、準コンパクトスキーム上の任意の構成可能アーベル層は、
$\mathbf{Z}/n\mathbf{Z}$-加群の構成可能層であり、そのような $n$ が存在する。
\end{proof}

\begin{lemma}
\label{etale-cohomology-lemma-category-loc-cst-is-colimit}
$X = \lim_{i \in I} X_i$ を、遷移射がアフィンであるスキームの有向系の極限とする。
$X_i$ は、すべての $i \in I$ に対して準コンパクトかつ準分離であると仮定する。
\begin{enumerate}
\item $X_\etale$ 上の有限局所定数層の圏は、$(X_i)_\etale$ 上の
有限局所定数層の圏の余極限である。
\item $X_\etale$ 上の有限局所定数アーベル層の圏は、$(X_i)_\etale$ 上の
有限局所定数アーベル層の圏の余極限である。
\item $\Lambda$ を Noether 環とする。有限型局所定数 $\Lambda$-加群層の圏で
$X_\etale$ 上にあるものは、有限型局所定数 $\Lambda$-加群層の圏で
$(X_i)_\etale$ 上にあるものの余極限である。
\end{enumerate}
\end{lemma}

\begin{proof}
補題 \ref{etale-cohomology-lemma-category-is-colimit} により、各場合の関手は充満忠実である。
同じ補題により、(1) の証明を終えるために示すべきことは次の一点だけである。
構成可能層 $\mathcal{F}_i$ が $X_i$ 上にあり、その引き戻し $\mathcal{F}$ が
$X$ 上で有限局所定数であるとする。このとき、ある $i' \geq i$ が存在し、
引き戻し $\mathcal{F}_{i'}$ は $\mathcal{F}_i$ を $X_{i'}$ へ引き戻したもので、
有限局所定数となる。仮定により、エタール被覆
$\mathcal{U} = \{U_j \to X\}_{j \in J}$ が存在し、
$\mathcal{F}|_{U_j} \cong \underline{S_j}$ となる。ここで $S_j$ は有限集合である。
$U_j$ は、すべての $j \in J$ に対してアフィンであると仮定してよい。
$X$ は準コンパクトなので、$J$ は有限と
仮定してよい。補題 \ref{etale-cohomology-lemma-colimit-affine-sites} により、$i' \geq i$ と
エタール被覆 $\mathcal{U}_{i'} = \{U_{i', j} \to X_{i'}\}_{j \in J}$ を取れる。
その $X$ への基底変換は $\mathcal{U}$ である。このとき
$\mathcal{F}_{i'}|_{U_{i', j}}$ と $\underline{S_j}$ は $(U_{i', j})_\etale$ 上の
構成可能層であり、$U_j$ への引き戻しは同型である。したがって $i'$ を大きくすれば、
$\mathcal{F}_{i'}|_{U_{i', j}}$ と $\underline{S_j}$ は同型になる。
よって $\mathcal{F}_{i'}$ は有限局所定数である。(2) と (3) の証明もまったく同じである。
\end{proof}

\begin{lemma}
\label{etale-cohomology-lemma-irreducible-subsheaf-constant-zero}
$X$ を一般点 $\eta$ をもつ既約スキームとする。
\begin{enumerate}
\item $S' \subset S$ を集合の包含とする。
$\underline{S'} \subset \mathcal{G} \subset \underline{S}$ が $\Sh(X_\etale)$ において成り立ち、かつ
$S' = \mathcal{G}_{\overline{\eta}}$ ならば、$\mathcal{G} = \underline{S'}$ である。
\item $A' \subset A$ をアーベル群の包含とする。
$\underline{A'} \subset \mathcal{G} \subset \underline{A}$ が $\textit{Ab}(X_\etale)$ において成り立ち、かつ
$A' = \mathcal{G}_{\overline{\eta}}$ ならば、$\mathcal{G} = \underline{A'}$ である。
\item $M' \subset M$ を環 $\Lambda$ 上の加群の包含とする。
$\underline{M'} \subset \mathcal{G} \subset \underline{M}$ が
$\textit{Mod}(X_\etale, \underline{\Lambda})$ において成り立ち、かつ
$M' = \mathcal{G}_{\overline{\eta}}$ ならば、$\mathcal{G} = \underline{M'}$ である。
\end{enumerate}
\end{lemma}

\begin{proof}
実際、任意のエタール射 $U \to X$ で $U \not = \emptyset$ であるものの像には点
$\eta$ が含まれる。
\end{proof}

\begin{lemma}
\label{etale-cohomology-lemma-push-constant-sheaf-from-open}
$X$ を函数体 $K$ をもつ整正規スキームとし、$E$ を集合とする。
\begin{enumerate}
\item $g : \Spec(K) \to X$ を一般点の包含とする。このとき
$g_*\underline{E} = \underline{E}$ である。
\item $j : U \to X$ を空でない開集合の包含とする。このとき
$j_*\underline{E} = \underline{E}$ である。
\end{enumerate}
\end{lemma}

\begin{proof}
(1) を証明する。点 $x \in X$ を取り、厳密 Hensel 化
$\mathcal{O}^{sh}_{X, \overline{x}}$ を $\mathcal{O}_{X, x}$ に対して取る。More on Algebra の補題
\ref{more-algebra-lemma-henselization-normal} により
$\mathcal{O}^{sh}_{X, \overline{x}}$ は正規整域である。したがって
$\Spec(K) \times_X \Spec(\mathcal{O}^{sh}_{X, \overline{x}})$ は既約である。
ゆえに茎 $(g_*\underline{E})_{\overline{x}}$ は $E$ に等しい。
定理 \ref{etale-cohomology-theorem-higher-direct-images} を参照されたい。

\medskip\noindent
(2) を証明する。$g$ は $j$ を経由するので写像
$j_*\underline{E} \to g_*\underline{E}$ がある。この写像は単射である。
実際、任意のスキーム $V$ で $X$ 上エタールなものに対して
$\Spec(K) \times_X V$ は $U \times_X V$ において稠密である。一方、写像
$\underline{E} \to j_*\underline{E}$ もあるので結論を得る。
\end{proof}

\begin{lemma}
\label{etale-cohomology-lemma-zero-in-generic-point}
$X$ を準コンパクトかつ準分離なスキームとし、$\eta \in X$ を $X$ の
ある既約成分の一般点とする。
\begin{enumerate}
\item $\mathcal{F}$ を $X_\etale$ 上のねじれアーベル層とし、その茎
$\mathcal{F}_{\overline{\eta}}$ が零であるとする。このとき
$\mathcal{F} = \colim \mathcal{F}_i$ は構成可能アーベル層
$\mathcal{F}_i$ のフィルター余極限であって、各 $i$ に対し
$\mathcal{F}_i$ の台は $\eta$ を含まない閉部分スキームに含まれる。
\item $\Lambda$ を Noether 環とし、$\mathcal{F}$ を $\Lambda$-加群の層で
$X_\etale$ 上にあるものとする。その茎 $\mathcal{F}_{\overline{\eta}}$ が零であると仮定する。
このとき $\mathcal{F} = \colim \mathcal{F}_i$ は $\Lambda$-加群の構成可能層
$\mathcal{F}_i$ のフィルター余極限であって、各 $i$ に対し
$\mathcal{F}_i$ の台は $\eta$ を含まない閉部分スキームに含まれる。
\end{enumerate}
\end{lemma}

\begin{proof}
(1) を証明する。補題 \ref{etale-cohomology-lemma-torsion-colimit-constructible} により、
$\mathcal{F} = \colim_{i \in I} \mathcal{F}_i$ と書ける。ここで
$\mathcal{F}_i$ は構成可能アーベル層である。$i \in I$ を選ぶ。仮定により
$\mathcal{F}|_\eta$ は零なので、ある $i'(i) \geq i$ が存在し、
$\mathcal{F}_i|_\eta \to \mathcal{F}_{i'(i)}|_\eta$ が零となる。
補題 \ref{etale-cohomology-lemma-colimit-constructible} を参照されたい。
このとき $\mathcal{G}_i = \Im(\mathcal{F}_i \to \mathcal{F}_{i'(i)})$ は
構成可能アーベル層であり（補題 \ref{etale-cohomology-lemma-constructible-abelian}）、
$\eta$ におけるその茎は零である。したがって、その台 $E_i$ は
$\mathcal{G}_i$ の台であり、$X$ の構成可能部分集合であって $\eta$ を含まない。
$\eta$ は $X$ の
ある既約成分の一般点なので、Topology の補題
\ref{topology-lemma-generic-point-in-constructible} により
$\eta \not \in Z_i = \overline{E_i}$ である。新しい有向集合 $I'$ を、集合 $I$ に
$i_1$ が $i_2$ 以上であることを $i_1 \geq i'(i_2)$ と同値とする順序を入れて、
定める。このとき層 $\mathcal{G}_i$ は $I'$ 上の系をなし、
その余極限は $\mathcal{F}$ である。これで証明は完了する。

\medskip\noindent
(2) の証明もまったく同じなので省略する。
\end{proof}








\section{Noether スキーム上の構成可能層}
\label{etale-cohomology-section-noetherian-constructible}

\noindent
$X$ が Noether スキームならば、任意の局所閉集合は構成可能な局所閉集合である
（Topology の補題 \ref{topology-lemma-constructible-Noetherian-space}）。
したがってアーベル層 $\mathcal{F}$ が $X_\etale$ 上で構成可能であることと、
有限分割 $X = \coprod X_i$ が存在して各部分上で
$\mathcal{F}|_{X_i}$ が有限局所定数となることは同値である。
（慣例により、位相空間の分割の各部分は局所閉である。Topology の節
\ref{topology-section-stratifications} を参照されたい。）
言い換えれば、定義 \ref{etale-cohomology-definition-constructible} では「構成可能」という形容詞を
省略してよい。この節では、Noether スキーム上の構成可能層の圏がもつ追加の性質を
いくつか列挙する。

\begin{proposition}
\label{etale-cohomology-proposition-constructible-over-noetherian}
$X$ を Noether スキームとし、$\Lambda$ を Noether 環とする。
\begin{enumerate}
\item 集合の構成可能層の任意の部分層または商層は構成可能である。
\item $X_\etale$ 上の構成可能アーベル層の圏は
$\textit{Ab}(X_\etale)$ の（強い）Serre 部分圏である。特に、
$X_\etale$ 上の構成可能アーベル層の任意の部分層および商層は構成可能である。
\item $\Lambda$-加群の構成可能層の圏は $X_\etale$ 上で
$\textit{Mod}(X_\etale, \Lambda)$ の（強い）Serre 部分圏である。特に、
任意の部分加群および商加群は $\Lambda$-加群の構成可能層であり、$X_\etale$ 上にある。
\end{enumerate}
\end{proposition}

\begin{proof}
 (1) を証明する。$\mathcal{G} \subset \mathcal{F}$ とし、$\mathcal{F}$ は
$X_\etale$ 上の集合の構成可能層とする。$\eta \in X$ を $X$ の既約成分の
一般点とする。Noether 帰納法により、$U$ を $\eta$ の開近傍として
$\mathcal{G}|_U$ が局所定数となるものを見つければ十分である。そのために $X$ を
$\eta$ のエタール近傍で置き換えてよい。したがって $\mathcal{F}$ は定数層で $X$ は
既約であると仮定してよい。

\medskip\noindent
$\mathcal{F} = \underline{S}$ と書く。ここで $S$ は有限集合である。
このとき $S' = \mathcal{G}_{\overline{\eta}} \subset S$ であり、
$S' = \{s_1, \ldots, s_t\}$ と書く。$(U, \overline{u})$ を
$\overline{\eta}$ のエタール近傍とし、切断
$\sigma_1, \ldots, \sigma_t \in \mathcal{G}(U)$ を取り、$s_i$ が
$\mathcal{G}_{\overline{\eta}} \subset S$ に写るものとする。
$\sigma_i$ は $s_i \in S' \subset S = \Gamma(X, \mathcal{F})$ に写るので、
$\sigma_i$ の $U \times_X U$ への二つの引き戻しは $\mathcal{G}$ の切断として同じである。
$\mathcal{G}$ の層条件により、$\sigma_i$ は
$\mathcal{G}$ の切断から来る $\Im(U \to X)$ という $X$ の開集合上にある。$X$ を縮小すれば
$\underline{S'} \subset \mathcal{G} \subset \underline{S}$ と仮定できる。
すると補題 \ref{etale-cohomology-lemma-irreducible-subsheaf-constant-zero} により
$\underline{S'} = \mathcal{G}$ である。

\medskip\noindent
次に $\mathcal{F} \to \mathcal{Q}$ を全射とし、$\mathcal{F}$ は
$X_\etale$ 上の集合の構成可能層とする。このとき
$\mathcal{G} = \mathcal{F} \times_\mathcal{Q} \mathcal{F}$ と置く。
証明の第一部により、$\mathcal{G}$ は $\mathcal{F} \times \mathcal{F}$ の部分層として
構成可能である。これから $\mathcal{Q}$ も構成可能であることが従う。
補題 \ref{etale-cohomology-lemma-constructible-abelian} を参照されたい。

\medskip\noindent
(3) を証明する。加群の構成可能層が弱い Serre 部分圏をなすことは既に分かっている。
補題 \ref{etale-cohomology-lemma-constructible-abelian} を参照されたい。したがって部分加群について
主張を示せば十分である。

\medskip\noindent
$\mathcal{G} \subset \mathcal{F}$ を部分加群とする。これは $\Lambda$-加群の構成可能層で
$X_\etale$ 上にある。$\eta \in X$ を $X$ の既約成分の一般点とする。
Noether 帰納法により、$U$ を $\eta$ の開近傍として $\mathcal{G}|_U$ が局所定数となる
ものを見つければ十分である。そのために $X$ を $\eta$ のエタール近傍で置き換えてよい。
したがって $\mathcal{F}$ は定数層で $X$ は既約であると仮定してよい。

\medskip\noindent
$\mathcal{F} = \underline{M}$ と書く。ここで $\Lambda$-加群 $M$ は有限である。
このとき $M' = \mathcal{G}_{\overline{\eta}} \subset M$ である。
有限個の元 $s_1, \ldots, s_t$ を取り、これらが $M'$ を $\Lambda$-加群として生成するとする。
（$\Lambda$ が Noether 環で $M$ が有限なので、これは可能である。）
$(U, \overline{u})$ を $\overline{\eta}$ のエタール近傍とし、切断
$\sigma_1, \ldots, \sigma_t \in \mathcal{G}(U)$ を取り、$s_i$ が
$\mathcal{G}_{\overline{\eta}} \subset M$ に写るものとする。
$\sigma_i$ は $s_i \in M' \subset M = \Gamma(X, \mathcal{F})$ に写るので、
$\sigma_i$ の $U \times_X U$ への二つの引き戻しは $\mathcal{G}$ の切断として同じである。
$\mathcal{G}$ の層条件により、$\sigma_i$ は
$\mathcal{G}$ の切断から来る $\Im(U \to X)$ という $X$ の開集合上にある。$X$ を縮小すれば
$\underline{M'} \subset \mathcal{G} \subset \underline{M}$ と仮定できる。
すると補題 \ref{etale-cohomology-lemma-irreducible-subsheaf-constant-zero} により
$\underline{M'} = \mathcal{G}$ である。

\medskip\noindent
(2) を証明する。(3) から通常の方法で従う。詳細は省略する。
\end{proof}

\noindent
次の補題は、$X$ 上の構成可能層のアーベル圏のすべての対象が
「Noether 的」である、すなわち部分対象について昇鎖条件を満たすことを示す。

\begin{lemma}
\label{etale-cohomology-lemma-constructible-over-noetherian-noetherian}
$X$ を Noether スキームとし、$\Lambda$ を Noether 環とする。次の包含を考える。
$$
\mathcal{F}_1 \subset \mathcal{F}_2 \subset \mathcal{F}_3 \subset \ldots
\subset \mathcal{F}
$$
集合、アーベル群、または $\Lambda$-加群の層の圏におけるものとする。
$\mathcal{F}$ が構成可能ならば、ある $n$ に対して
$\mathcal{F}_n = \mathcal{F}_{n + 1} = \mathcal{F}_{n + 2} = \ldots$ となる。
\end{lemma}

\begin{proof}
命題 \ref{etale-cohomology-proposition-constructible-over-noetherian} により
$\mathcal{F}_i$ と $\colim \mathcal{F}_i$ は構成可能である。
したがって補題 \ref{etale-cohomology-lemma-colimit-constructible} から従う。
\end{proof}

\begin{lemma}
\label{etale-cohomology-lemma-constructible-maps-into-constant}
$X$ を Noether スキームとする。
\begin{enumerate}
\item $\mathcal{F}$ を $X_\etale$ 上の集合の構成可能層とする。層の単射
$$
\mathcal{F} \longrightarrow
\prod\nolimits_{i = 1, \ldots, n} f_{i, *}\underline{E_i}
$$
が存在する。ここで $f_i : Y_i \to X$ は有限射で、$E_i$ は有限集合である。
\item $\mathcal{F}$ を $X_\etale$ 上の構成可能アーベル層とする。アーベル層の単射
$$
\mathcal{F} \longrightarrow
\bigoplus\nolimits_{i = 1, \ldots, n} f_{i, *}\underline{M_i}
$$
が存在する。ここで $f_i : Y_i \to X$ は有限射で、$M_i$ は有限アーベル群である。
\item $\Lambda$ を Noether 環とする。$\mathcal{F}$ を
$\Lambda$-加群の構成可能層とし、$X_\etale$ 上の加群層の単射
$$
\mathcal{F} \longrightarrow
\bigoplus\nolimits_{i = 1, \ldots, n} f_{i, *}\underline{M_i}
$$
が存在する。ここで $f_i : Y_i \to X$ は有限射で、$M_i$ は有限 $\Lambda$-加群である。
\end{enumerate}
さらに、各 $Y_i$ は既約かつ被約で、既約かつ被約な閉部分スキーム
$Z_i \subset X$ へ写り、$Y_i \to Z_i$ は $Z_i$ の空でない開集合上で
有限エタールとなるものと仮定してよい。
\end{lemma}

\begin{proof}
 (1) を証明する。補題 \ref{etale-cohomology-lemma-constructible-over-noetherian-noetherian} により
$\mathcal{F}$ の部分層について昇鎖条件が成り立つので、各点 $x \in X$ に対して、
写像 $\varphi : \mathcal{F} \to f_*\underline{E}$ を見つければ十分である。ここで
$f : Y \to X$ は有限射、$E$ は有限集合であり、茎
$\varphi_{\overline{x}} : \mathcal{F}_{\overline{x}} \to
(f_*S)_{\overline{x}}$ は単射となるものとする。
（補題 \ref{etale-cohomology-lemma-constructible-quasi-compact-quasi-separated} のように $X$ を分割し、
各部分の各既約成分に対して $Y_i \to X$ を構成すれば、この議論は避けられる。）
$Z \subset X$ を点の閉包 $\overline{\{x\}}$ に誘導される被約スキーム構造
（Schemes の定義 \ref{schemes-definition-reduced-induced-scheme}）とする。$\mathcal{F}$ は構成可能なので、
有限分離拡大 $K/\kappa(x)$ が存在し、$\mathcal{F}|_{\Spec(K)}$ は値 $E$ の定数層となる
（$E$ は有限集合）。$Y \to Z$ を $Z$ の $\Spec(K)$ における正規化とする。
Morphisms の補題 \ref{morphisms-lemma-normal-normalization} により、$Y$ は正規整スキームである。
$K/\kappa(x)$ は有限拡大なので、$K$ は $Y$ の函数体である。包含を
$g : \Spec(K) \to Y$ と書く。写像 $\mathcal{F}|_{\Spec(K)} \to \underline{E}$ は、
補題 \ref{etale-cohomology-lemma-push-constant-sheaf-from-open} により写像
$\mathcal{F}|_Y \to g_*\underline{E} = \underline{E}$ に随伴する。
これはさらに写像 $\varphi : \mathcal{F} \to f_*\underline{E}$ に随伴する。
$\varphi$ の幾何点 $\overline{x}$ における茎が単射であることに注意する。
$\overline{y} \in Y$ を取り、これを $\overline{x}$ の Y 上の持ち上げとすれば、可換図式
$$
\xymatrix{
\mathcal{F}_{\overline{x}} \ar@{=}[r] \ar[d] &
(\mathcal{F}|_Y)_{\overline{y}} \ar@{=}[d] \\
(f_*\underline{E})_{\overline{x}} \ar[r] &
\underline{E}_{\overline{y}}
}
$$
が単射性を示す。しかしまだ終わっていない。射 $f : Y \to Z$ は整であるが、一般には
有限ではない\footnote{例えば $X$ が体上有限型の Nagata スキームならば、$Y \to Z$ は有限である。}。

\medskip\noindent
前段落の最後に述べた問題を解決するため、$Y = \lim_{i \in I} Y_i$ と書く。
ここで $Y_i$ は既約かつ整で $Z$ 上有限である。すなわち、$f_*\mathcal{O}_Y$ を
$\mathcal{O}_Z$-環の層とみなし、Properties の補題
\ref{properties-lemma-integral-algebra-directed-colimit-finite} を適用してから
関手 $\underline{\Spec}_Z$ を適用する。補題 \ref{etale-cohomology-lemma-relative-colimit} により
$f_*\underline{E} = \colim f_{i, *}\underline{E}$ である。補題
\ref{etale-cohomology-lemma-constructible-is-compact} により、写像
$\mathcal{F} \to f_*\underline{E}$ は
$f_{i, *}\underline{E}$ を経由し、ある $i$ が存在する。$Y_i \to Z$ は整スキーム間の有限射であり、
この射が誘導する函数体拡大は有限分離なので、射は $Z$ の空でない開集合上で有限
エタールである（Algebra の補題 \ref{algebra-lemma-smooth-at-generic-point} を用いる。
詳細は省略する）。これで (1) の証明が完了する。

\medskip\noindent
(2) と (3) の証明は (1) の証明と同一である。
\end{proof}

\noindent
次の補題では、非常に一般的な主張を Noether の場合へ帰着する標準的な技巧を用いる。

\begin{lemma}
\label{etale-cohomology-lemma-constructible-maps-into-constant-general}
\begin{reference}
\cite[Exposee IX, Proposition 2.14]{SGA4}
\end{reference}
$X$ を準コンパクトかつ準分離なスキームとする。
\begin{enumerate}
\item $\mathcal{F}$ を $X_\etale$ 上の集合の構成可能層とする。層の単射
$$
\mathcal{F} \longrightarrow
\prod\nolimits_{i = 1, \ldots, n} f_{i, *}\underline{E_i}
$$
が存在する。ここで $f_i : Y_i \to X$ は有限かつ有限表示な射で、$E_i$ は有限集合である。
\item $\mathcal{F}$ を $X_\etale$ 上の構成可能アーベル層とする。アーベル層の単射
$$
\mathcal{F} \longrightarrow
\bigoplus\nolimits_{i = 1, \ldots, n} f_{i, *}\underline{M_i}
$$
が存在する。ここで $f_i : Y_i \to X$ は有限かつ有限表示な射で、$M_i$ は有限アーベル群である。
\item $\Lambda$ を Noether 環とする。$\mathcal{F}$ を
$\Lambda$-加群の構成可能層とし、$X_\etale$ 上の加群層の単射
$$
\mathcal{F} \longrightarrow
\bigoplus\nolimits_{i = 1, \ldots, n} f_{i, *}\underline{M_i}
$$
が存在する。ここで $f_i : Y_i \to X$ は有限かつ有限表示な射で、$M_i$ は有限 $\Lambda$-加群である。
\end{enumerate}
\end{lemma}

\begin{proof}
この補題を絶対 Noether 近似により Noether の場合へ帰着する。すなわち Limits の命題
\ref{limits-proposition-approximate} により、$X = \lim_{t \in T} X_t$ と書ける。
各 $X_t$ は $\Spec(\mathbf{Z})$ 上有限型で、遷移射はアフィンである。
補題 \ref{etale-cohomology-lemma-category-is-colimit} により、（集合、アーベル群、または
$\Lambda$-加群の）構成可能層の圏は $X_\etale$ 上で、$X_t$ に対応する圏の余極限である。
したがって構成可能層 $\mathcal{F}$ は同様な構成可能層 $\mathcal{F}_t$ の引き戻しであり、
それは $X_t$ 上である（ある $t$ について）。そこで Noether の場合
（補題 \ref{etale-cohomology-lemma-constructible-maps-into-constant}）を適用して単射
$$
\mathcal{F}_t \longrightarrow
\prod\nolimits_{i = 1, \ldots, n} f_{i, *}\underline{E_i}
\quad\text{or}\quad
\mathcal{F}_t \longrightarrow
\bigoplus\nolimits_{i = 1, \ldots, n} f_{i, *}\underline{M_i}
$$
を得る。これは $X_t$ 上で有限射 $f_i : Y_i \to X_t$ に対して成り立つ。
$X_t$ は Noether なので、射 $f_i$ は有限表示である。引き戻しは完全であり、また
$f_{i, *}$ の構成は基底変換と可換する（補題
\ref{etale-cohomology-lemma-finite-pushforward-commutes-with-base-change}）ので、結論を得る。
\end{proof}

\begin{lemma}
\label{etale-cohomology-lemma-support-in-subset}
$X$ を Noether スキームとし、$E \subset X$ を特殊化で閉じている部分集合とする。
\begin{enumerate}
\item $\mathcal{F}$ を $X_\etale$ 上のねじれアーベル層で、その台が $E$ に含まれるものとする。
このとき $\mathcal{F} = \colim \mathcal{F}_i$ は構成可能アーベル層 $\mathcal{F}_i$ の
フィルター余極限であり、各 $i$ に対して $\mathcal{F}_i$ の台は $E$ に含まれる閉集合に
含まれる。
\item $\Lambda$ を Noether 環とし、$\mathcal{F}$ を $\Lambda$-加群の層で $X_\etale$ 上にあり、
その台が $E$ に含まれるものとする。このとき $\mathcal{F} = \colim \mathcal{F}_i$ は
$\Lambda$-加群の構成可能層 $\mathcal{F}_i$ のフィルター余極限であり、各 $i$ に対して
$\mathcal{F}_i$ の台は $E$ に含まれる閉集合に含まれる。
\end{enumerate}
\end{lemma}

\begin{proof}
 (1) を証明する。補題 \ref{etale-cohomology-lemma-torsion-colimit-constructible} により
$\mathcal{F} = \colim_{i \in I} \mathcal{F}_i$ と書け、$\mathcal{F}_i$ は構成可能アーベル層である。
命題 \ref{etale-cohomology-proposition-constructible-over-noetherian} により、その像
$\mathcal{F}'_i \subset \mathcal{F}$ は、写像 $\mathcal{F}_i \to \mathcal{F}$ により構成可能である。
したがって $\mathcal{F} = \colim \mathcal{F}'_i$ であり、$\mathcal{F}'_i$ の台は $E$ に含まれる。
$\mathcal{F}'_i$ の台は（構成可能層の定義により）構成可能なので、その閉包も $E$ に含まれる。
例えば Topology の補題
\ref{topology-lemma-constructible-stable-specialization-closed} を参照されたい。

\medskip\noindent
(2) の証明もまったく同じなので省略する。
\end{proof}








\section{特殊化とエタール層}
\label{etale-cohomology-section-specialization}

\noindent
トポロジー的な見方：$X$ を位相空間とし、$x' \leadsto x$ を $X$ における点の特殊化とする。
$x$ の任意の開近傍は $x'$ を含む。したがって任意の層 $\mathcal{F}$ を $X$ 上にとると
{\it 特殊化写像}
$$
sp : \mathcal{F}_x \longrightarrow \mathcal{F}_{x'}
$$
という茎の写像であり、組 $(U, s)$ の
$\mathcal{F}_x$ における同値類を、組 $(U, s)$ の
$\mathcal{F}_{x'}$ における同値類へ送る。組の同値類による茎の記述については
層、節 \ref{sheaves-section-stalks} を参照せよ。この写像が $\mathcal{F}$ に関して関手的であること、すなわち $sp$
が関手の変換であることは明らかである。

\medskip\noindent
エタール位相の層についてもこの構成をまねることができる（
\cite[Exposee VIII, 7.7, page 397]{SGA4} を参照）。そこでスキーム $S$、幾何点 $\overline{s}$ で $S$ 上のもの、および
幾何点 $\overline{t}$ が $\Spec(\mathcal{O}^{sh}_{S, \overline{s}})$ 上にあるとする。
任意の層 $\mathcal{F}$ を $S_\etale$ 上にとり、次の
{\it 特殊化写像}
$$
sp : \mathcal{F}_{\overline{s}} \longrightarrow \mathcal{F}_{\overline{t}}
$$
ここでは言葉を少し乱用している。$\mathcal{F}_{\overline{t}}$ と書く代わりに $\mathcal{F}_{p(\overline{t})}$ と書くべきであり、ここで $p : \Spec(\mathcal{O}^{sh}_{S, \overline{s}}) \to S$ は標準射である。次を思い出そう。
$$
\mathcal{F}_{\overline{s}} = \colim_{(U, \overline{u})} \mathcal{F}(U)
$$
ここで余極限はすべてのエタール近傍 $(U, \overline{u})$ にわたるものであり、$(S, \overline{s})$ の近傍についての記述は節 \ref{etale-cohomology-section-stalks} を参照せよ。$\mathcal{O}^{sh}_{S, \overline{s}}$ は構造層の茎なので、各エタール近傍 $(U, \overline{u})$ が $(S, \overline{s})$ のものであるとき、標準写像 $\mathcal{O}_{U, u} \to \mathcal{O}^{sh}_{S, \overline{s}}$ が得られる。
したがって一意的な分解を得る。
$$
\Spec(\mathcal{O}^{sh}_{S, \overline{s}}) \to U \to S
$$
$\overline{v}$ を $\overline{t}$ の $U$ における像とすると、$(U, \overline{v})$ は $(S, \overline{t})$ のエタール近傍である。この構成は $(S, \overline{s})$ のエタール近傍の圏から $(S, \overline{t})$ のエタール近傍の圏への関手を定める。したがって写像
$sp : \mathcal{F}_{\overline{s}} \to \mathcal{F}_{\overline{t}}$
を、$(U, \overline{u}, \sigma)$ の同値類（ただし $\sigma \in \mathcal{F}(U)$）を $(U, \overline{v}, \sigma)$ の同値類へ送ることによって定義できる。

\medskip\noindent
次に $K \in D(S_\etale)$ とする。上記の $\overline{s}$ および $\overline{t}$ に対して {\it 特殊化写像}
$$
sp : K_{\overline{s}} \longrightarrow K_{\overline{t}}
\quad\text{ において }\quad D(\textit{Ab})
$$
すなわち、$K$ がアーベル層の複体 $\mathcal{F}^\bullet$ によって表されるなら、写像
$$
K_{\overline{s}} = \mathcal{F}^\bullet_{\overline{s}}
\longrightarrow
\mathcal{F}^\bullet_{\overline{t}} = K_{\overline{t}}
$$
は上で構成した層の特殊化写像を各次数ごとに与える。これは茎関手の完全性（すなわち複体の茎をとる操作が導来圏上で適切に定義されること）により、$K$ を表す複体の選択によらない。

\medskip\noindent
この構成は層 $\mathcal{F}$ に関して $S_\etale$ 上で明らかに関手的である。茎関手をトポスの射
$\overline{s}, \overline{t} : \textit{Sets} \to \Sh(S_\etale)$ とみなせば、$sp$ を次の $2$-射
$$
\xymatrix{
\textit{Sets}
\rrtwocell^{\overline{t}}_{\overline{s}}{\ sp}
&
&
\Sh(S_\etale)
}
$$
とみなすことができる。

\begin{remark}[特殊化写像の別の記述]
\label{etale-cohomology-remark-alternative-sp}
上記の $S$、$\overline{s}$、$\overline{t}$ をとる。
特殊化写像の別の記述には次を用いる。
$$
\mathcal{F}_{\overline{s}} =
\Gamma(\Spec(\mathcal{O}^{sh}_{S, \overline{s}}), p^{-1}\mathcal{F})
\quad\text{ かつ }\quad
\mathcal{F}_{\overline{t}} = \Gamma(\overline{t},
\overline{t}^{-1}p^{-1}\mathcal{F})
$$
最初の等号は $\text{id}_S : S \to S$ に定理 \ref{etale-cohomology-theorem-higher-direct-images} を適用したものから従い、二つ目の等号は補題 \ref{etale-cohomology-lemma-stalk-pullback} から従う。したがって $sp$ を次の写像とみなせる。
$$
sp :
\mathcal{F}_{\overline{s}} =
\Gamma(\Spec(\mathcal{O}^{sh}_{S, \overline{s}}), p^{-1}\mathcal{F})
\xrightarrow{\text{ による引き戻し }\overline{t}}
\Gamma(\overline{t}, \overline{t}^{-1}p^{-1}\mathcal{F}) =
\mathcal{F}_{\overline{t}}
$$
\end{remark}

\begin{remark}[特殊化写像のさらに別の記述]
\label{etale-cohomology-remark-another-sp}
上記の $S$、$\overline{s}$、$\overline{t}$ をとる。
別の方法は、一意的な射を用いることである。
$$
c :
\Spec(\mathcal{O}^{sh}_{S, \overline{t}})
\longrightarrow
\Spec(\mathcal{O}^{sh}_{S, \overline{s}})
$$
これは $S$ 上で、与えられた射 $\overline{t} \to \Spec(\mathcal{O}^{sh}_{S, \overline{s}})$ および射 $\overline{t} \to \Spec(\mathcal{O}^{sh}_{t, \overline{t}})$ と両立する。表示した矢印の存在と一意性は、$\mathcal{O}_{S, s}$、$\mathcal{O}^{sh}_{S, \overline{t}}$、および $\mathcal{O}^{sh}_{S, \overline{s}} \to \kappa(\overline{t})$ に代数、補題 \ref{algebra-lemma-map-into-henselian-colimit} を適用すれば従う。よって
$$
sp :
\mathcal{F}_{\overline{s}} =
\Gamma(\Spec(\mathcal{O}^{sh}_{S, \overline{s}}), \mathcal{F})
\xrightarrow{\text{ による引き戻し }c}
\Gamma(\Spec(\mathcal{O}^{sh}_{S, \overline{t}}), \mathcal{F}) =
\mathcal{F}_{\overline{t}}
$$
（記号については明らかな慣習に従う）。実際、この手続きは対象 $K$ で $D(S_\etale)$ に属するものに対しても成り立つ。$K$ の特殊化写像は次の写像に等しい。
$$
sp :
K_{\overline{s}} =
R\Gamma(\Spec(\mathcal{O}^{sh}_{S, \overline{s}}), K)
\xrightarrow{\text{ による引き戻し }c}
R\Gamma(\Spec(\mathcal{O}^{sh}_{S, \overline{t}}), K) =
K_{\overline{t}}
$$
等号が正当なのは、狭義ヘンゼルスキーム
$\Spec(\mathcal{O}^{sh}_{S, \overline{s}})$ および
$\Spec(\mathcal{O}^{sh}_{S, \overline{t}})$ 上で大域切断をとる操作が完全であり（また $\overline{s}$ および $\overline{t}$ で茎をとる操作と同じであり）、したがって $K$ が非有界であり得ることによる微妙な問題が生じないためである。
\end{remark}

\begin{remark}[特殊化の持ち上げ]
\label{etale-cohomology-remark-can-lift}
$S$ をスキームとし、$t \leadsto s$ を $S$ 上の点の特殊化とする。幾何点 $\overline{t}$ および $\overline{s}$ を選び、それらが $t$ および $s$ の上にあるとする。$t$ は
$\Spec(\mathcal{O}_{S, s})$ の点に対応し、これはスキーム、補題 \ref{schemes-lemma-specialize-points} による。
さらに $\mathcal{O}_{S, s} \to \mathcal{O}^{sh}_{S, \overline{s}}$ は忠実平坦なので、
$t' \in \Spec(\mathcal{O}^{sh}_{S, \overline{s}})$ で $t$ へ写る点を見つけられる。
$\Spec(\mathcal{O}^{sh}_{S, \overline{s}})$ は $S$ 上エタールなスキームの極限なので、$\kappa(t')/\kappa(t)$ は分離代数拡大である（もちろん通常は有限ではない）。
$\kappa(\overline{t})$ は代数閉なので、埋め込み $\kappa(t') \to \kappa(\overline{t})$ を $\kappa(t)$ の拡大として選べる。この選択により次の可換図式
$$
\xymatrix{
\overline{t} \ar[d] \ar[r] &
\Spec(\mathcal{O}^{sh}_{S, \overline{s}}) \ar[d] &
\overline{s} \ar[l] \ar[d] \\
t \ar[r] & S & s \ar[l]
}
$$
が点と幾何点の図式として得られる。したがって $t \leadsto s$ ならば、$\overline{t}$ を $S$ の $\overline{s}$ における狭義ヘンゼル化の幾何点へ常に「持ち上げ」、上記の特殊化写像を得られる。
\end{remark}

\begin{lemma}
\label{etale-cohomology-lemma-specialization-map-pullback}
スキームの射 $g : S' \to S$ をとる。層 $\mathcal{F}$ を $S_\etale$ 上にとる。幾何点 $\overline{s}'$ を $S'$ にとり、幾何点 $\overline{t}'$ を $\Spec(\mathcal{O}^{sh}_{S', \overline{s}'})$ 上にとる。$\overline{s} = g(\overline{s}')$ および $\overline{t} = h(\overline{t}')$ とおく。ここで $h : \Spec(\mathcal{O}^{sh}_{S', \overline{s}'}) \to
\Spec(\mathcal{O}^{sh}_{S, \overline{s}})$ は標準射である。任意の層 $\mathcal{F}$ を $S_\etale$ 上にとったときの特殊化写像
$$
sp :
(g^{-1}\mathcal{F})_{\overline{s}'}
\longrightarrow
(g^{-1}\mathcal{F})_{\overline{t}'}
$$
は特殊化写像
$sp : \mathcal{F}_{\overline{s}} \to \mathcal{F}_{\overline{t}}$
に等しい。これは補題 \ref{etale-cohomology-lemma-stalk-pullback} の同一視
$(g^{-1}\mathcal{F})_{\overline{s}'} = \mathcal{F}_{\overline{s}}$ および
$(g^{-1}\mathcal{F})_{\overline{t}'} = \mathcal{F}_{\overline{t}}$
を介して理解する。
\end{lemma}

\begin{proof}
省略する。
\end{proof}

\begin{lemma}
\label{etale-cohomology-lemma-characterize-locally-constant}
$S$ を、$S$ の各準コンパクト開集合が有限個の既約成分をもつスキームとする（例えば $S$ の基礎位相空間が Noether である場合、または $S$ が局所的に Noether である場合）。
集合の層 $\mathcal{F}$ を $S_\etale$ 上にとる。次は同値である。
\begin{enumerate}
\item $\mathcal{F}$ は有限局所定数である。
\item $\mathcal{F}$ のすべての茎が有限集合であり、すべての特殊化写像
$sp : \mathcal{F}_{\overline{s}} \to \mathcal{F}_{\overline{t}}$
は全単射である。
\end{enumerate}
\end{lemma}

\begin{proof}
(2) を仮定する。幾何点 $\overline{s}$ を $S$ の上で、$s \in S$ の上にとる。(1) を示すには、$(U, \overline{u})$ を $(S, \overline{s})$ のエタール近傍として選び、$\mathcal{F}|_U$ が定数となるものを見つければよい。$S$ はアフィンであると仮定してよい。

\medskip\noindent
$\mathcal{F}_{\overline{s}}$ は有限なので、$(U, \overline{u})$、$n \geq 0$、および相異なる元
$\sigma_1, \ldots, \sigma_n \in \mathcal{F}(U)$ を選ぶ。
$\{\sigma_1, \ldots, \sigma_n\} \subset \mathcal{F}(U)$ が
$\mathcal{F}_{\overline{s}}$ へ、写像 $\mathcal{F}(U) \to \mathcal{F}_{\overline{s}}$ によって
全単射に写るようにする。写像を考える。
$$
\varphi : \underline{\{1, \ldots, n\}} \longrightarrow \mathcal{F}|_U
$$
$U_\etale$ 上で $\sigma_1, \ldots, \sigma_n$ により定める。この写像は構成により $\overline{u}$ における茎上で全単射である。次の部分集合を考える。
$$
E = \{u' \in U \mid \varphi_{\overline{u}'}\text{ は全単射}\} \subset U
$$
ここで $\overline{u}'$ は $U$ の幾何点で $u'$ の上にある任意のものとする（この条件は注意 \ref{etale-cohomology-remark-map-stalks} により選択によらない）。像 $u \in U$ は $\overline{u}$ の像であり、$E$ に属する。$\mathcal{F}$ の特殊化写像に関する仮定、注意 \ref{etale-cohomology-remark-can-lift}、および補題 \ref{etale-cohomology-lemma-specialization-map-pullback} により、$E$ は $U$ の位相空間で特殊化および一般化に関して閉じていることが分かる。

\medskip\noindent
さらに $U$ を縮小して $U$ もアフィンであると仮定してよい。降下、補題 \ref{descent-lemma-locally-finite-nr-irred-local-fppf} により、$U$ は有限個の既約成分をもつ。$u$ を通らない既約成分を除けば、$U$ のすべての既約成分が $u$ を通ると仮定できる。$U$ はソーバー位相空間なので $E = U$ が従い、定理 \ref{etale-cohomology-theorem-exactness-stalks} により $\varphi$ は同型である。したがって (1) が従う。

\medskip\noindent
(1) から (2) が従うことの証明は省略する。
\end{proof}

\begin{lemma}
\label{etale-cohomology-lemma-characterize-locally-constant-module}
$S$ を、$S$ の各準コンパクト開集合が有限個の既約成分をもつスキームとする（例えば $S$ の基礎位相空間が Noether である場合、または $S$ が局所的に Noether である場合）。
Noether 環 $\Lambda$ をとる。
層 $\mathcal{F}$ を $\Lambda$-加群として $S_\etale$ 上にとる。次は同値である。
\begin{enumerate}
\item $\mathcal{F}$ は有限型の局所定数 $\Lambda$-加群の層である。
\item $\mathcal{F}$ のすべての茎が有限 $\Lambda$-加群であり、すべての特殊化写像
$sp : \mathcal{F}_{\overline{s}} \to \mathcal{F}_{\overline{t}}$
は全単射である。
\end{enumerate}
\end{lemma}

\begin{proof}
この補題の証明は補題 \ref{etale-cohomology-lemma-characterize-locally-constant} の証明と同じである。(2) を仮定する。幾何点 $\overline{s}$ を $S$ の上で、$s \in S$ の上にとる。(1) を示すには、$(U, \overline{u})$ を $(S, \overline{s})$ のエタール近傍として選び、$\mathcal{F}|_U$ が定数となるものを見つければよい。$S$ はアフィンであると仮定してよい。

\medskip\noindent
$M = \mathcal{F}_{\overline{s}}$ は有限 $\Lambda$-加群であり、$\Lambda$ は Noether なので、次の表示を選べる。
$$
\Lambda^{\oplus m} \xrightarrow{A} \Lambda^{\oplus n} \to M \to 0
$$
行列 $A = (a_{ji})$ の係数が $\Lambda$ にあるものとして、$(U, \overline{u})$ と元 $\sigma_1, \ldots, \sigma_n \in \mathcal{F}(U)$ を選び、$\sum a_{ji}\sigma_i = 0$ が $\mathcal{F}(U)$ で成り立ち、$\sigma_i$ の $\mathcal{F}_{\overline{s}} = M$ における像が、$\Lambda^n$ の標準基底元の像（上の $M$ の表示）となるようにする。写像を考える。
$$
\varphi : \underline{M} \longrightarrow \mathcal{F}|_U
$$
$U_\etale$ 上で $\sigma_1, \ldots, \sigma_n$ により定める。この写像は構成により $\overline{u}$ における茎上で全単射である。次の部分集合を考える。
$$
E = \{u' \in U \mid \varphi_{\overline{u}'}\text{ は全単射}\} \subset U
$$
ここで $\overline{u}'$ は $U$ の幾何点で $u'$ の上にある任意のものとする（この条件は注意 \ref{etale-cohomology-remark-map-stalks} により選択によらない）。像 $u \in U$ は $\overline{u}$ の像であり、$E$ に属する。$\mathcal{F}$ の特殊化写像に関する仮定、注意 \ref{etale-cohomology-remark-can-lift}、および補題 \ref{etale-cohomology-lemma-specialization-map-pullback} により、$E$ は $U$ の位相空間で特殊化および一般化に関して閉じていることが分かる。

\medskip\noindent
さらに $U$ を縮小して $U$ もアフィンであると仮定してよい。降下、補題 \ref{descent-lemma-locally-finite-nr-irred-local-fppf} により、$U$ は有限個の既約成分をもつ。$u$ を通らない既約成分を除けば、$U$ のすべての既約成分が $u$ を通ると仮定できる。$U$ はソーバー位相空間なので $E = U$ が従い、定理 \ref{etale-cohomology-theorem-exactness-stalks} により $\varphi$ は同型である。したがって (1) が従う。

\medskip\noindent
(1) から (2) が従うことの証明は省略する。
\end{proof}

\begin{lemma}
\label{etale-cohomology-lemma-specialization-map-pushforward}
$f : X \to S$ を準コンパクトかつ準分離なスキームの射とする。
$K \in D^+(X_\etale)$ とする。幾何点 $\overline{s}$ を $S$ にとり、$\overline{t}$ を $\Spec(\mathcal{O}^{sh}_{S, \overline{s}})$ の幾何点としてとる。次の可換図式を得る。
$$
\xymatrix{
(Rf_*K)_{\overline{s}} \ar[r]_{sp} \ar@{=}[d] &
(Rf_*K)_{\overline{t}} \ar@{=}[d] \\
R\Gamma(X \times_S \Spec(\mathcal{O}^{sh}_{S, \overline{s}}), K)
\ar[r] &
R\Gamma(X \times_S \Spec(\mathcal{O}^{sh}_{S, \overline{t}}), K)
}
$$
ここで下の水平矢印は射 $\text{id}_X \times c$ による引き戻しから生じる。ここで
$c : \Spec(\mathcal{O}^{sh}_{S, \overline{t}}) \to
\Spec(\mathcal{O}^{sh}_{S, \overline{s}})$
は注意 \ref{etale-cohomology-remark-another-sp} で導入した射である。垂直矢印は定理 \ref{etale-cohomology-theorem-higher-direct-images} によって与えられる。
\end{lemma}

\begin{proof}
これは注意 \ref{etale-cohomology-remark-another-sp} における $sp$ の記述から直ちに従う。
\end{proof}

\begin{remark}
\label{etale-cohomology-remark-specialization-map-and-fibres}
$f : X \to S$ をスキームの射とする。$K \in D(X_\etale)$ とする。
幾何点 $\overline{s}$ を $S$ にとり、$\overline{t}$ を $\Spec(\mathcal{O}^{sh}_{S, \overline{s}})$ の幾何点としてとる。
$c$ を注意 \ref{etale-cohomology-remark-another-sp} のものとする。常に次の可換図式を作ることができる。
$$
\xymatrix{
(Rf_*K)_{\overline{s}} \ar[r] \ar[d]_{sp} &
R\Gamma(X \times_S \Spec(\mathcal{O}_{S, \overline{s}}^{sh}), K)
\ar[r] \ar[d]_{(\text{id}_X \times c)^{-1}} &
R\Gamma(X_{\overline{s}}, K) \\
(Rf_*K)_{\overline{t}} \ar[r] &
R\Gamma(X \times_S \Spec(\mathcal{O}_{S, \overline{t}}^{sh}), K) \ar[r] &
R\Gamma(X_{\overline{t}}, K)
}
$$
ここで水平矢印は注意 \ref{etale-cohomology-remark-stalk-fibre} のものである。一般には、この図式に適合するような、ファイバー上の $K$ のコホモロジーの間の右側の垂直写像は存在しない。これは補題 \ref{etale-cohomology-lemma-specialization-map-pushforward} の場合でさえ同様である。
\end{remark}















\section{構成可能コホモロジーをもつ複体}
\label{etale-cohomology-section-Dc}

\noindent
$\Lambda$ を環とする。$D(X_\etale, \Lambda)$ を
$\Lambda$-加群の層の $X_\etale$ 上の導来圏と呼ぶ。
$D^b(X_\etale, \Lambda)$（それぞれ $D^+$、$D^-$）を
$D(X_\etale, \Lambda)$ における有界（それぞれ上に、下に）有界な複体の
充満部分圏とする。

\begin{definition}
\label{etale-cohomology-definition-c}
$X$ をスキームとし、$\Lambda$ を Noether 環とする。
{\it $D_c(X_\etale, \Lambda)$} を、$D(X_\etale, \Lambda)$ のうち
コホモロジー層が $\Lambda$-加群の構成可能層である複体の充満部分圏とする。
\end{definition}

\noindent
これは補題 \ref{etale-cohomology-lemma-constructible-abelian} および
Derived Categories, 節 \ref{derived-section-triangulated-sub} により
意味をもつ。したがって $D_c(X_\etale, \Lambda)$ は
$D(X_\etale, \Lambda)$ の厳密充満かつ飽和した三角部分圏である。

\begin{lemma}
\label{etale-cohomology-lemma-restrict-and-shriek-from-etale-c}
$\Lambda$ を Noether 環とする。$j : U \to X$ がスキームの
エタール射ならば、
\begin{enumerate}
\item $K|_U \in D_c(U_\etale, \Lambda)$ である。ただし
$K \in D_c(X_\etale, \Lambda)$ である。
\item $j_!M \in D_c(X_\etale, \Lambda)$ である。ただし
$M \in D_c(U_\etale, \Lambda)$ かつ射 $j$ は準コンパクトかつ準分離である。
\end{enumerate}
\end{lemma}

\begin{proof}
第一の主張は明らかである。第二は $j_!$ が完全であり、
補題 \ref{etale-cohomology-lemma-jshriek-constructible} による。
\end{proof}

\begin{lemma}
\label{etale-cohomology-lemma-pullback-c}
$\Lambda$ を Noether 環とする。$f : X \to Y$ をスキームの射とする。
$K \in D_c(Y_\etale, \Lambda)$ ならば
$Lf^*K \in D_c(X_\etale, \Lambda)$ である。
\end{lemma}

\begin{proof}
これは $f^{-1} = f^*$ が完全であり、
補題 \ref{etale-cohomology-lemma-pullback-constructible} による。
\end{proof}

\begin{lemma}
\label{etale-cohomology-lemma-one-constructible}
$X$ を準コンパクトかつ準分離なスキームとする。
$\Lambda$ を Noether 環とする。$K \in D(X_\etale, \Lambda)$ および
$b \in \mathbf{Z}$ をとり、$H^b(K)$ が構成可能であるとする。
すると、層 $\mathcal{F}$ であって、
$j_{U!}\underline{\Lambda}$ の有限直和であり、$U \in \Ob(X_\etale)$ が
アフィンであるものと、写像 $\mathcal{F}[-b] \to K$ であって
$D(X_\etale, \Lambda)$ において $\mathcal{F} \to H^b(K)$ を
全射にするものが存在する。
\end{lemma}

\begin{proof}
$K$ を $\mathcal{K}^\bullet$ という $\Lambda$-加群の層の複体で表す。
次の全射を考える。
$$
\Ker(\mathcal{K}^b \to \mathcal{K}^{b + 1})
\longrightarrow
H^b(K)
$$
Modules on Sites, 補題
\ref{sites-modules-lemma-module-quotient-direct-sum} により、次の全射を選べる。
$\bigoplus_{i \in I} j_{U_i!} \underline{\Lambda} \to
\Ker(\mathcal{K}^b \to \mathcal{K}^{b + 1})$
ここで $U_i$ はアフィンである。$I' \subset I$ が有限のとき、
$\mathcal{H}_{I'} \subset H^b(K)$ を
$\bigoplus_{i \in I'} j_{U_i!} \underline{\Lambda}$ の像とする。
補題 \ref{etale-cohomology-lemma-colimit-constructible} により、
$\mathcal{H}_{I'} = H^b(K)$ となる有限な $I' \subset I$ が存在する。
補題は
$\mathcal{F} = \bigoplus_{i \in I'} j_{U_i!} \underline{\Lambda}$
ととれば従う。
\end{proof}

\begin{lemma}
\label{etale-cohomology-lemma-bounded-above-c}
$X$ を準コンパクトかつ準分離なスキームとする。
$\Lambda$ を Noether 環とする。$K \in D^-(X_\etale, \Lambda)$ とする。
次は同値である。
\begin{enumerate}
\item $K$ は $D_c(X_\etale, \Lambda)$ に属する。
\item $K$ は上に有界な複体で表され、その項は
$j_{U!}\underline{\Lambda}$ の有限直和であり、
$U \in \Ob(X_\etale)$ はアフィンである。
\item $K$ は上に有界な $\Lambda$-加群の平坦な構成可能層の複体で表される。
\end{enumerate}
\end{lemma}

\begin{proof}
(2) が (3) を、また (3) が (1) を導くことは明らかである。
$K$ が $D_c^-(X_\etale, \Lambda)$ に属すると仮定する。
$H^i(K) = 0$ が $i > b$ であるとする。$a$ に関する帰納法により、
複体 $\mathcal{F}^a \to \ldots \to \mathcal{F}^b$ を構成する。
各 $\mathcal{F}^i$ は $j_{U!}\underline{\Lambda}$ の有限直和で、
$U \in \Ob(X_\etale)$ はアフィンであり、さらに写像
$\mathcal{F}^\bullet \to K$ が存在して、
$H^i(\mathcal{F}^\bullet) \to H^i(K)$ は $i > a$ で同型、
$H^a(\mathcal{F}^\bullet) \to H^a(K)$ は全射となるようにする。
$a = b$ の場合は補題 \ref{etale-cohomology-lemma-one-constructible} により構成できる。
このデータが与えられたとき、次の完全三角を選ぶ。
$$
\mathcal{F}^\bullet \to K \to L \to \mathcal{F}^\bullet[1]
$$
すると $H^i(L) = 0$ が $i \geq a$ で成り立つ。補題
\ref{etale-cohomology-lemma-one-constructible} に従って
$\mathcal{F}^{a - 1}[-a +1] \to L$ を選ぶ。この合成
$\mathcal{F}^{a - 1}[-a +1] \to L \to \mathcal{F}^\bullet$ は
$\mathcal{F}^{a - 1} \to \mathcal{F}^a$ という写像に対応し、
$\mathcal{F}^a \to \mathcal{F}^{a + 1}$ との合成はゼロである。
TR4 により次の写像を得る。
$$
(\mathcal{F}^{a - 1} \to \ldots \to \mathcal{F}^b) \to K
$$
$D(X_\etale, \Lambda)$ においてである。これで帰納段階と補題の証明が
完了する。
\end{proof}

\begin{lemma}
\label{etale-cohomology-lemma-tensor-c}
$X$ をスキームとし、$\Lambda$ を Noether 環とする。
$K, L \in D_c^-(X_\etale, \Lambda)$ とする。このとき
$K \otimes_\Lambda^\mathbf{L} L$ は $D_c^-(X_\etale, \Lambda)$ に属する。
\end{lemma}

\begin{proof}
これは補題 \ref{etale-cohomology-lemma-bounded-above-c} および
\ref{etale-cohomology-lemma-tensor-product-constructible} による。
\end{proof}




\section{構成可能コホモロジーをもつ Tor 有限複体}
\label{etale-cohomology-section-ctf}

\noindent
$X$ をスキームとし、$\Lambda$ を Noether 環とする。
節 \ref{etale-cohomology-section-Dc} で定義した導来圏 $D_c(X_\etale, \Lambda)$ の
よく用いる部分圏として、（局所的に）有限な Tor 次元をもつ対象からなる
充満部分圏がある。形式的な定義は次のとおりである。

\begin{definition}
\label{etale-cohomology-definition-ctf}
$X$ をスキームとする。$\Lambda$ を Noether 環とする。
{\it $D_{ctf}(X_\etale, \Lambda)$} を、$D_c(X_\etale, \Lambda)$ のうち
局所的に有限な Tor 次元をもつ対象からなる充満部分圏とする。
\end{definition}

\noindent
これは $D_c(X_\etale, \Lambda)$ および $D(X_\etale, \Lambda)$ の
厳密充満かつ飽和した三角部分圏である。規約により、Cohomology on Sites,
定義 \ref{sites-cohomology-definition-tor-amplitude} を参照すると、
$$
D_{ctf}(X_\etale, \Lambda) \subset D^b_c(X_\etale, \Lambda)
\subset D^b(X_\etale, \Lambda)
$$
$X$ が準コンパクトならば成り立つ。$D_{ctf}(X_\etale, \Lambda)$ の対象を
理解するよい方法は補題 \ref{etale-cohomology-lemma-when-ctf} に与えられている。

\begin{remark}
\label{etale-cohomology-remark-different}
$D_{ctf}(X_\etale, \Lambda)$ の導来圏の対象は、ある意味で
$D(\mathcal{O}_X)$ の完全対象よりよい大域的性質をもつ。
すなわち、$\mathcal{O}_X$-加群の複体が局所的には有限な局所自由
$\mathcal{O}_X$-加群の有限複体と準同型でありながら、局所自由
$\mathcal{O}_X$-加群の有界複体と大域的には準同型でないことがありうる。
次の補題は、Noether スキーム上の $D_{ctf}$ ではこれが起こらないことを示す。
\end{remark}

\begin{lemma}
\label{etale-cohomology-lemma-when-ctf}
$\Lambda$ を Noether 環とする。$X$ を準コンパクトかつ準分離なスキームとする。
$K \in D(X_\etale, \Lambda)$ とする。次は同値である。
\begin{enumerate}
\item $K \in D_{ctf}(X_\etale, \Lambda)$ であり、
\item $K$ は $\Lambda$-加群の構成可能平坦層の有限複体で表される。
\end{enumerate}
実際、$K$ の Tor 振幅が $[a, b]$ にあるなら、$K$ は
$\mathcal{F}^a \to \ldots \to \mathcal{F}^b$ という複体で表せ、
$\mathcal{F}^p$ は $\Lambda$-加群の構成可能平坦層である。
\end{lemma}

\begin{proof}
構成可能な $\Lambda$-加群の平坦層の有限複体が有限 Tor 次元をもつことは明らかである。
また、それが $D_c(X_\etale, \Lambda)$ の対象であることも明らかである。
したがって (2) は (1) を導く。

\medskip\noindent
(1) を仮定する。$a, b \in \mathbf{Z}$ を選び、
$H^i(K \otimes_\Lambda^\mathbf{L} \mathcal{G}) = 0$ が
$i \not \in [a, b]$ のとき、すべての $\Lambda$-加群の層 $\mathcal{G}$ に対して
成り立つようにする。
最後の主張が成り立つことを $b - a$ に関する帰納法で示す。
$a = b$ ならば、$K = H^a(K)[-a]$ は構成可能平坦層であり、主張が従う。
次に $b > a$ と仮定する。$K$ を複体 $\mathcal{K}^\bullet$ で表す。
これは $\Lambda$-加群の層からなる。次の全射を考える。
$$
\Ker(\mathcal{K}^b \to \mathcal{K}^{b + 1})
\longrightarrow
H^b(K)
$$
補題 \ref{etale-cohomology-lemma-category-constructible-modules} により、$U_i$ を有限個選ぶ。
これらは $X$ 上エタールなアフィンスキームであり、全射
$\bigoplus j_{U_i!}\underline{\Lambda}_{U_i} \to H^b(K)$ を得られる。
$U_i$ を標準エタール被覆 $\{U_{ij} \to U_i\}$ で置き換えると、
この全射は写像
$\mathcal{F} = \bigoplus j_{U_i!}\underline{\Lambda}_{U_i} \to
\Ker(\mathcal{K}^b \to \mathcal{K}^{b + 1})$ に持ち上がると仮定できる。
この写像は次の完全三角を定める。
$$
\mathcal{F}[-b] \to K \to L \to \mathcal{F}[-b + 1]
$$
$D(X_\etale, \Lambda)$ においてである。$D_{ctf}(X_\etale, \Lambda)$ は
三角部分圏なので、$L$ もその中にある。実際、$L$ の Tor 振幅は
$[a, b - 1]$ である。これは $\mathcal{F}$ が $H^b(K)$ へ全射するためである
（詳細は省略する）。帰納法の仮定により、
$\mathcal{F}^a \to \ldots \to \mathcal{F}^{b - 1}$ という有限複体を見つけられる。
これは $\Lambda$-加群の平坦構成可能層からなる。$L$ は
$L \to \mathcal{F}[-b + 1]$ に対応し、その写像は
$\mathcal{F}^b \to \mathcal{F}$ であって
$\mathcal{F}^{b - 1} \to \mathcal{F}^b$ の像を消すものに対応する。
したがって公理 TR3 により、$K$ は次の複体で表される。
$$
\mathcal{F}^a \to \ldots \to \mathcal{F}^{b - 1} \to \mathcal{F}^b
$$
これで証明が完了する。
\end{proof}

\begin{remark}
\label{etale-cohomology-remark-projective-each-degree}
$\Lambda$ を Noether 環とし、$X$ をスキームとする。
構成可能平坦加群の有界複体 $K^\bullet$（各項は $\Lambda$-加群）を
$X_\etale$ 上で考えると、各茎 $K^p_{\overline{x}}$ は有限射影 $\Lambda$-加群である。
したがって複体の茎は $\Lambda$-加群の完全複体である。
\end{remark}

\begin{lemma}
\label{etale-cohomology-lemma-restrict-and-shriek-from-etale-ctf}
$\Lambda$ を Noether 環とする。$j : U \to X$ をスキームのエタール射とすると、
\begin{enumerate}
\item $K|_U \in D_{ctf}(U_\etale, \Lambda)$ ならば
$K \in D_{ctf}(X_\etale, \Lambda)$ であり、
\item $j_!M \in D_{ctf}(X_\etale, \Lambda)$ ならば
$M \in D_{ctf}(U_\etale, \Lambda)$ であり、射 $j$ は準コンパクトかつ準分離である。
\end{enumerate}
\end{lemma}

\begin{proof}
この補題を証明する最も簡単な方法は、$X$ がアフィンの場合に帰着し、
補題 \ref{etale-cohomology-lemma-when-ctf} を適用して、平坦な構成可能層の有限複体
（$\Lambda$-加群の層からなる）に関する主張へ翻訳することである。
その場合の結果は補題 \ref{etale-cohomology-lemma-jshriek-constructible} から従う。
\end{proof}

\begin{lemma}
\label{etale-cohomology-lemma-pullback-ctf}
$\Lambda$ を Noether 環とする。$f : X \to Y$ をスキームの射とする。
$K \in D_{ctf}(Y_\etale, \Lambda)$ ならば
$Lf^*K \in D_{ctf}(X_\etale, \Lambda)$ である。
\end{lemma}

\begin{proof}
補題 \ref{etale-cohomology-lemma-when-ctf} を適用すると、これは $\Lambda$-加群の
平坦構成可能層の有限複体に関する問題へ帰着する。
このとき $f^{-1} = f^*$ は完全であるから、主張は補題
\ref{etale-cohomology-lemma-pullback-constructible} から従う。
\end{proof}

\begin{lemma}
\label{etale-cohomology-lemma-connected-ctf-locally-constant}
$X$ を連結スキームとし、$\Lambda$ を Noether 環とする。
$K \in D_{ctf}(X_\etale, \Lambda)$ のコホモロジー層が局所定数であるとする。
このとき、有限射影 $\Lambda$-加群の有限複体 $M^\bullet$ とエタール被覆
$\{U_i \to X\}$ が存在し、$K|_{U_i} \cong \underline{M^\bullet}|_{U_i}$ が
$D(U_{i, \etale}, \Lambda)$ において成り立つ。
\end{lemma}

\begin{proof}
エタール被覆 $\{U_i \to X\}$ を選び、$K|_{U_i}$ が定数となるようにする。
すなわち、$K|_{U_i} \cong \underline{M_i^\bullet}_{U_i}$ とする。ただし
有限 $\Lambda$-加群からなる有限複体 $M_i^\bullet$ である。
Cohomology on Sites, 補題 \ref{sites-cohomology-lemma-locally-constant} を参照。
$U_i \times_X U_j$ は、$M_i^\bullet$ が $M_j^\bullet$ と
$D(\Lambda)$ において同型でないなら空であることに注意する。
各 $\Lambda$-加群の複体 $M^\bullet$ に対し、
$I_{M^\bullet} =
\{i \in I \mid M_i^\bullet \cong M^\bullet\text{（}D(\Lambda)\text{ において）}\}$
とおく。エタール射は開写像なので、
$U_{M^\bullet} = \bigcup_{i \in I_{M^\bullet}} \Im(U_i \to X)$
は $X$ の開部分集合である。すると
$X = \coprod U_{M^\bullet}$ は $X$ の互いに素な開被覆となる。
$X$ は連結なので、空でない $U_{M^\bullet}$ は一つだけである。
$K$ は $D_{ctf}(X_\etale, \Lambda)$ に属するから、$M^\bullet$ は
$\Lambda$-加群の完全複体である。More on Algebra, 補題
\ref{more-algebra-lemma-perfect} を参照。したがって $M^\bullet$ は
有限射影 $\Lambda$-加群の有限複体であると仮定できる。
\end{proof}








\section{捩れ層}
\label{etale-cohomology-section-torsion}

\noindent
これは捩れアーベル層とそのエタール・コホモロジーについての短い節である。
$\mathcal{C}$ をサイトとする。Cohomology on Sites, 補題
\ref{sites-cohomology-lemma-torsion} で、コホモロジー層が捩れ層である
$D(\mathcal{C})$ の任意の対象は、すべての項が捩れである複体で表せることを
示した。

\begin{lemma}
\label{etale-cohomology-lemma-torsion-cohomology}
 $X$ を準コンパクトかつ準分離なスキームとする。
\begin{enumerate}
\item $\mathcal{F}$ が $X_\etale$ 上の捩れアーベル層ならば、
$H^n_\etale(X, \mathcal{F})$ はすべての $n$ について捩れアーベル群である。
\item $K$ が $D^+(X_\etale)$ に属し捩れコホモロジー層をもつならば、
$H^n_\etale(X, K)$ はすべての $n$ について捩れアーベル群である。
\end{enumerate}
\end{lemma}

\begin{proof}
 (1) を示すため、$\mathcal{F} = \bigcup \mathcal{F}[n]$ と書く。ただし
$\mathcal{F}[d]$ は $d$-捩れ部分層である。補題 \ref{etale-cohomology-lemma-colimit} により
$H^n_\etale(X, \mathcal{F}) = \colim H^n_\etale(X, \mathcal{F}[d])$
を得る。$H^n_\etale(X, \mathcal{F}[d])$ は $d$ で消去されるので、これで
(1) が示される。

\medskip\noindent
 (2) を示すには、$E_2^{p, q} = H^p_\etale(X, H^q(K))$ から
$H^n_\etale(X, K)$ へ収束するスペクトル系列
(Derived Categories, 補題 \ref{derived-lemma-two-ss-complex-functor})
と、層についての結果を用いればよい。
\end{proof}

\begin{lemma}
\label{etale-cohomology-lemma-torsion-direct-image}
 $f : X \to Y$ を準コンパクトかつ準分離なスキームの射とする。
\begin{enumerate}
\item $\mathcal{F}$ が $X_\etale$ 上の捩れアーベル層ならば、
$R^nf_*\mathcal{F}$ は $Y_\etale$ 上の捩れアーベル層であり、すべての $n$ について成り立つ。
\item $K$ が $D^+(X_\etale)$ に属し捩れコホモロジー層をもつならば、
$Rf_*K$ は $D^+(Y_\etale)$ の対象であり、そのコホモロジー層は捩れ
アーベル層である。
\end{enumerate}
\end{lemma}

\begin{proof}
 (1) の証明。$R^nf_*\mathcal{F}$ は、前層
$V \mapsto H^n_\etale(X \times_Y V, \mathcal{F})$ に付随する
$Y_\etale$ 上の層である。
Cohomology on Sites, 補題 \ref{sites-cohomology-lemma-higher-direct-images}
を参照。$V$ をアフィンに選ぶと、$X \times_Y V$ は $f$ が準コンパクトかつ
準分離なので準コンパクトかつ準分離であり、したがって補題
\ref{etale-cohomology-lemma-torsion-cohomology} を適用すると
$H^n_\etale(X \times_Y V, \mathcal{F})$ が捩れであることが分かる。

\medskip\noindent
 (2) の証明。$R^nf_*K$ は、前層
$V \mapsto H^n_\etale(X \times_Y V, K)$ に付随する
$Y_\etale$ 上の層である。
Cohomology on Sites, 補題
\ref{sites-cohomology-lemma-unbounded-describe-higher-direct-images} を参照。
$V$ をアフィンに選ぶと、$X \times_Y V$ は $f$ が準コンパクトかつ準分離なので
準コンパクトかつ準分離であり、したがって補題 \ref{etale-cohomology-lemma-torsion-cohomology}
を適用すると $H^n_\etale(X \times_Y V, K)$ が捩れであることが分かる。
\end{proof}









\section{閉部分スキームを台とするコホモロジー}
\label{etale-cohomology-section-cohomology-support}

\noindent
$X$ をスキームとし、$Z \subset X$ を閉部分スキームとする。
$\mathcal{F}$ を $X_\etale$ 上のアーベル層とする。次を定める。
$$
\Gamma_Z(X, \mathcal{F}) =
\{s \in \mathcal{F}(X) \mid \text{Supp}(s) \subset Z\}
$$
これを $Z$ に台をもつ切断と呼ぶ（定義 \ref{etale-cohomology-definition-support}）。
これは左完全関手だが、一般には完全ではない。したがって導来関手
$$
R\Gamma_Z(X, -) : D(X_\etale) \longrightarrow D(\textit{Ab})
$$
および $Z$ を台とするコホモロジー群を
$H^q_Z(X, \mathcal{F}) = R^q\Gamma_Z(X, \mathcal{F})$.

\medskip\noindent
$\mathcal{I}$ を $X_\etale$ 上の入射アーベル層とする。
$U = X \setminus Z$ とおく。このとき制限写像
$\mathcal{I}(X) \to \mathcal{I}(U)$ は全射であり
（Cohomology on Sites, 補題
\ref{sites-cohomology-lemma-restriction-along-monomorphism-surjective}）、
その核は $\Gamma_Z(X, \mathcal{I})$ である。したがって直ちに、
$K \in D(X_\etale)$ に対して完全三角
$$
R\Gamma_Z(X, K) \to R\Gamma(X, K) \to R\Gamma(U, K) \to R\Gamma_Z(X, K)[1]
$$
が $D(\textit{Ab})$ において存在する。これから長完全コホモロジー系列
$$
\ldots \to H^i_Z(X, K) \to H^i(X, K) \to H^i(U, K) \to
H^{i + 1}_Z(X, K) \to \ldots
$$
を、任意の $K$ が $D(X_\etale)$ に属する場合に得る。

\medskip\noindent
$\mathcal{F}$ を $X_\etale$ 上のアーベル層とする。$Z$ に台をもつ切断の
部分層を考えることができる。これを $\mathcal{H}_Z(\mathcal{F})$ と記し、
次の規則で定める。
$$
\mathcal{H}_Z(\mathcal{F})(U) =
\{s \in \mathcal{F}(U) \mid \text{Supp}(s) \subset U \times_X Z\}
$$
ここでは定義 \ref{etale-cohomology-definition-support} による切断の台を用いた。
命題 \ref{etale-cohomology-proposition-closed-immersion-pushforward} の同値を用いると、
$\mathcal{H}_Z(\mathcal{F})$ を $Z_\etale$ 上のアーベル層とみなせる。
したがって関手
$$
\textit{Ab}(X_\etale) \longrightarrow \textit{Ab}(Z_\etale),\quad
\mathcal{F} \longmapsto \mathcal{H}_Z(\mathcal{F})
$$
を得る。これは左完全だが、一般には完全ではない。

\begin{lemma}
\label{etale-cohomology-lemma-sections-with-support-acyclic}
$i : Z \to X$ をスキームの閉埋め込みとする。
$\mathcal{I}$ を $X_\etale$ 上の入射アーベル層とする。
このとき $\mathcal{H}_Z(\mathcal{I})$ は $Z_\etale$ 上の入射アーベル層である。
\end{lemma}

\begin{proof}
任意のアーベル層 $\mathcal{G}$ を $Z_\etale$ 上で考えると
$$
\Hom_Z(\mathcal{G}, \mathcal{H}_Z(\mathcal{F})) =
\Hom_X(i_*\mathcal{G}, \mathcal{F})
$$
となる。実際、$i_*\mathcal{G}$ の任意の切断は $Z$ に台をもつ。
$i_*$ は完全であり（節 \ref{etale-cohomology-section-closed-immersions}）、かつ
$\mathcal{I}$ は $X_\etale$ 上で入射なので、
$\mathcal{H}_Z(\mathcal{I})$ は $Z_\etale$ 上で入射であると結論できる。
\end{proof}

\noindent
次を導来関手と記す。
$$
R\mathcal{H}_Z : D(X_\etale) \longrightarrow D(Z_\etale)
$$
さらに
$\mathcal{H}^q_Z(\mathcal{F}) = R^q\mathcal{H}_Z(\mathcal{F})$ とおき、
$\mathcal{H}^0_Z(\mathcal{F}) = \mathcal{H}_Z(\mathcal{F})$.
とおく。上の補題により Grothendieck スペクトル系列
$$
E_2^{p, q} = H^p(Z, \mathcal{H}^q_Z(\mathcal{F}))
\Rightarrow H^{p + q}_Z(X, \mathcal{F})
$$

\begin{lemma}
\label{etale-cohomology-lemma-cohomology-with-support-sheaf-on-support}
$i : Z \to X$ をスキームの閉埋め込みとする。
$\mathcal{G}$ を $Z_\etale$ 上の入射アーベル層とする。
このとき $\mathcal{H}^p_Z(i_*\mathcal{G}) = 0$ は $p > 0$ で成り立つ。
\end{lemma}

\begin{proof}
これは関手 $i_*$ が完全であり、入射アーベル層を入射アーベル層へ写すためである
（Cohomology on Sites, 補題
\ref{sites-cohomology-lemma-pushforward-injective-flat}）。
\end{proof}

\begin{lemma}
\label{etale-cohomology-lemma-cohomology-with-support-triangle}
$i : Z \to X$ をスキームの閉埋め込みとする。
$j : U \to X$ を $Z$ の補集合の包含とする。
$\mathcal{F}$ を $X_\etale$ 上のアーベル層とする。
このとき完全三角
$$
i_*R\mathcal{H}_Z(\mathcal{F}) \to \mathcal{F} \to Rj_*(\mathcal{F}|_U) \to
i_*R\mathcal{H}_Z(\mathcal{F})[1]
$$
が $D(X_\etale)$ において存在する。これは完全系列
$$
0 \to i_*\mathcal{H}_Z(\mathcal{F}) \to \mathcal{F} \to
j_*(\mathcal{F}|_U) \to i_*\mathcal{H}^1_Z(\mathcal{F}) \to 0
$$
および同型
$R^pj_*(\mathcal{F}|_U) \cong i_*\mathcal{H}^{p + 1}_Z(\mathcal{F})$
を $p \geq 1$ に対して与える。
\end{lemma}

\begin{proof}
完全三角を得るため、入射分解
$\mathcal{F} \to \mathcal{I}^\bullet$ を選ぶ。すると複体の短完全系列
$$
0 \to
i_*\mathcal{H}_Z(\mathcal{I}^\bullet) \to \mathcal{I}^\bullet
\to j_*(\mathcal{I}^\bullet|_U) \to 0
$$
を上の議論から得る。よって完全三角は
Derived Categories, 節 \ref{derived-section-canonical-delta-functor} から従う。
\end{proof}

\noindent
$X$ をスキームとし、$Z \subset X$ を閉部分スキームとする。
$D_Z(X_\etale)$ を、$D(X_\etale)$ の厳密充満な飽和三角部分圏であって、
コホモロジー層が $Z$ に台をもつ複体からなるものと記す。
$D_Z(X_\etale)$ は $X$ の基礎となる閉集合だけに依存することに注意する。

\begin{lemma}
\label{etale-cohomology-lemma-complexes-with-support-on-closed}
$i : Z \to X$ をスキームの閉埋め込みとする。
写像 $Ri_{small, *} = i_{small, *} : D(Z_\etale) \to D(X_\etale)$ は、
準逆が次である同値 $D(Z_\etale) \to D_Z(X_\etale)$ を誘導する。
$$
i_{small}^{-1}|_{D_Z(X_\etale)} = R\mathcal{H}_Z|_{D_Z(X_\etale)}
$$
\end{lemma}

\begin{proof}
$i_{small}^{-1}$ と $i_{small, *}$ は完全関手の随伴対であり、
$i_{small}^{-1}i_{small, *}$ はアーベル層上の恒等関手と同型であることを思い出す。
命題 \ref{etale-cohomology-proposition-closed-immersion-pushforward} と補題
\ref{etale-cohomology-lemma-stalk-pullback} を参照。したがって
$i_{small, *} : D(Z_\etale) \to D_Z(X_\etale)$ は充満忠実であり、
$i_{small}^{-1}$ は左逆を定める。他方、$K$ が $D_Z(X_\etale)$ の対象であるとし、
随伴写像
$K \to i_{small, *}i_{small}^{-1}K$.
を考える。$i_{small, *}$ と $i_{small}^{-1}$ の完全性により、これはコホモロジー層上の
随伴写像 $H^n(K) \to i_{small, *}i_{small}^{-1}H^n(K)$ を誘導する。
これらのコホモロジー層は $Z$ に台をもつので、随伴写像は同型である。
したがって $D(Z_\etale) \to D_Z(X_\etale)$ は同値である。

\medskip\noindent
証明を終えるには $R\mathcal{H}_Z(K) = i_{small}^{-1}K$ を示せばよい。
ただし $K$ は $D_Z(X_\etale)$ の対象とする。このため、いま示した
$K = i_{small, *}i_{small}^{-1}K$ を用いる。そこで K-入射代表
$\mathcal{I}^\bullet$ を $i_{small}^{-1}K$ に対して選ぶ。
$i_{small, *}$ は完全関手 $i_{small}^{-1}$ の右随伴なので、複体
$i_{small, *}\mathcal{I}^\bullet$ は K-入射である
（Derived Categories, 補題 \ref{derived-lemma-adjoint-preserve-K-injectives}）。
$R\mathcal{H}_Z(K)$ は
$\mathcal{H}_Z(i_{small, *}\mathcal{I}^\bullet) = \mathcal{I}^\bullet$ によって
計算されるので、望む結論を得る。
\end{proof}

\begin{lemma}
\label{etale-cohomology-lemma-cohomology-with-support-quasi-coherent}
$X$ をスキームとし、$Z \subset X$ を閉部分スキームとする。
$\mathcal{F}$ を準連接 $\mathcal{O}_X$-加群とし、対応する準連接層を
$\mathcal{F}^a$ と記す。これは $X$ の小エタールサイト上の層である
（命題 \ref{etale-cohomology-proposition-quasi-coherent-sheaf-fpqc}）。すると
\begin{enumerate}
\item $H^q_Z(X, \mathcal{F})$ は $H^q_Z(X_\etale, \mathcal{F}^a)$ と一致する。
\item $Z$ の補集合が $X$ においてレトロコンパクトならば、
$i_*\mathcal{H}^q_Z(\mathcal{F}^a)$ は $\mathcal{O}_X$-加群の準連接層であり、
$(i_*\mathcal{H}^q_Z(\mathcal{F}))^a$ に等しい。
\end{enumerate}
\end{lemma}

\begin{proof}
$j : U \to X$ を $Z$ の補集合の包含とする。コホモロジー群に関する (1) は、
台をもつコホモロジーの長完全系列と、次の一致
$H^q(X_\etale, \mathcal{F}^a) = H^q(X, \mathcal{F})$ および
$H^q(U_\etale, \mathcal{F}^a) = H^q(U, \mathcal{F})$ から従う。定理
\ref{etale-cohomology-theorem-zariski-fpqc-quasi-coherent} を参照。
$j : U \to X$ が準コンパクト射、すなわち $U \subset X$ がレトロコンパクトならば、
$R^qj_*$ は準連接層を準連接層へ写す
（Cohomology of Schemes, 補題
\ref{coherent-lemma-quasi-coherence-higher-direct-images}）。
さらに、エタールサイト上で対応する層を取る操作と可換である
（Descent, 補題 \ref{descent-lemma-higher-direct-images-small-etale}）。
補題 \ref{etale-cohomology-lemma-cohomology-with-support-triangle} を適用すれば結論を得る。
\end{proof}







\section{局所環が狭義 Hensel 環であるスキーム}
\label{etale-cohomology-section-strictly-henselian-local-rings}

\noindent
この節では、局所環が狭義 Hensel 環であるスキームのエタール・コホモロジーに
関するいくつかの結果を集める。例えば、次は補題
\ref{etale-cohomology-lemma-vanishing-etale-cohomology-strictly-henselian} の興味深い一般化である。

\begin{lemma}
\label{etale-cohomology-lemma-local-rings-strictly-henselian}
$S$ をすべての局所環が狭義 Hensel 環であるスキームとする。
このとき任意のアーベル層 $\mathcal{F}$ を $S_\etale$ 上で考えると
$H^i(S_\etale, \mathcal{F}) = H^i(S_{Zar}, \mathcal{F})$.
\end{lemma}

\begin{proof}
$\epsilon : S_\etale \to S_{Zar}$ を包含関手が与えるサイトの射とする。
Zariski 層 $R^p\epsilon_*\mathcal{F}$ は、前層
$U \mapsto H^p_\etale(U, \mathcal{F})$ に付随する層である。
したがって $x \in X$ における茎は
$\colim H^p_\etale(U, \mathcal{F}) =
H^p_\etale(\Spec(\mathcal{O}_{X, x}), \mathcal{G}_x)$
である。ここで $\mathcal{G}_x$ は $\mathcal{F}$ の
$\Spec(\mathcal{O}_{X, x})$ への引き戻しを表す。補題
\ref{etale-cohomology-lemma-directed-colimit-cohomology} を参照。
したがって $R^p\epsilon_*\mathcal{F}$ の高次直接像は、補題
\ref{etale-cohomology-lemma-vanishing-etale-cohomology-strictly-henselian} により零である。
Leray スペクトル系列から結論が従う。
\end{proof}

\begin{lemma}
\label{etale-cohomology-lemma-gabber-criterion}
$R$ をすべての局所環が狭義 Hensel 環である環とする。
$\mathcal{F}$ を $\Spec(R)_\etale$ 上の層とする。
すべての $f, g \in R$ に対して次の写像の核が
$$
H^1_\etale(D(f + g), \mathcal{F})
\longrightarrow
H^1_\etale(D(f(f + g)), \mathcal{F}) \oplus
H^1_\etale(D(g(f + g)), \mathcal{F})
$$
零であると仮定する。このとき $H^q_\etale(\Spec(R), \mathcal{F}) = 0$ は
$q > 0$ で成り立つ。
\end{lemma}

\begin{proof}
補題 \ref{etale-cohomology-lemma-local-rings-strictly-henselian} により、$\mathcal{F}$ の
エタール・コホモロジーは $\Spec(R)$ の任意の開部分上で Zariski
コホモロジーと一致する。$i$ に関する帰納法で、$h \in R$ に対して
$H^q_\etale(D(h), \mathcal{F}) = 0$ が $1 \leq q \leq i$ で成り立つことを示す。
初期の場合 $i = 0$ は自明である。$i \geq 1$ と仮定する。

\medskip\noindent
$\xi \in H^q_\etale(D(h), \mathcal{F})$ とする。ただし
$1 \leq q \leq i$ かつ $h \in R$ である。$q < i$ なら帰納法により
終わるので、$q = i$ と仮定する。$R$ を $R_h$ で置き換えることにより、
$\xi \in H^i_\etale(\Spec(R), \mathcal{F})$ と仮定してよい。詳細の一部は省く。
$I \subset R$ を、元 $f \in R$ であって $\xi|_{D(f)} = 0$ となるものの集合とする。
$\xi$ は Zariski 局所的に自明なので、各素イデアル $\mathfrak p$（$R$ の）に対し、
$f \in I$ であって $f \not \in \mathfrak p$ となるものが存在する。したがって $I$ が
イデアルであることを示せば $1 \in I$ となり、証明は終わる。
$f \in I$ と $r \in R$ から $rf \in I$ が従うことは明らかである。
そこで $f, g \in I$ と仮定し、$f + g \in I$ を示す。
$q = i = 1$ ならば、これはまさに補題の仮定である。したがって $i = 1$ の
場合が従う。$q = i > 1$ の場合には
$$
D(f + g) = D(f(f + g)) \cup D(g(f + g))
$$
に注意する。Mayer--Vietoris（Cohomology, 補題
\ref{cohomology-lemma-mayer-vietoris}）により、次の完全系列を得る。
これは $\Spec(R)$ の開部分スキーム上のエタール・コホモロジーが Zariski
コホモロジーに等しいため適用できる。
$$
\xymatrix{
H^{i - 1}_\etale(D(fg(f + g)), \mathcal{F}) \ar[d] \\
H^i_\etale(D(f + g), \mathcal{F}) \ar[d] \\
H^i_\etale(D(f(f + g)), \mathcal{F}) \oplus
H^i_\etale(D(g(f + g)), \mathcal{F})
}
$$
最初の群は帰納法により零なので、結果が従う。
\end{proof}

\begin{lemma}
\label{etale-cohomology-lemma-affine-only-closed-points}
$S$ をアフィンスキームであって、(1) すべての点が閉点であり、
(2) すべての剰余体が分離代数閉であるものとする。このとき
任意のアーベル層 $\mathcal{F}$ を $S_\etale$ 上で考えると
$H^i(S_\etale, \mathcal{F}) = 0$ は $i > 0$ で成り立つ。
\end{lemma}

\begin{proof}
条件 (1) により $S$ の基礎位相空間は副有限である。Algebra, 補題
\ref{algebra-lemma-ring-with-only-minimal-primes} を参照。
したがって位相空間 $S$ 上のアーベル層の高次コホモロジー群、すなわち
Zariski コホモロジーは自明である。Cohomology, 補題
\ref{cohomology-lemma-vanishing-for-profinite} を参照。
局所環は Algebra, 補題
\ref{algebra-lemma-local-dimension-zero-henselian} により狭義 Hensel 環である。
よって補題 \ref{etale-cohomology-lemma-local-rings-strictly-henselian} により、$S$ の
エタール・コホモロジーは Zariski コホモロジーで計算され、証明が終わる。
\end{proof}

\noindent
絶対整閉環のスペクトルは、すべての局所環が狭義 Hensel 環であるスキームの
一例である。More on Algebra, 補題
\ref{more-algebra-lemma-absolutely-integrally-closed-strictly-henselian} を参照。
分離閉な分数体をもつ正規整域は、次の補題で述べるように、さらに強い性質をもつ。

\begin{lemma}
\label{etale-cohomology-lemma-normal-scheme-with-alg-closed-function-field}
$X$ を分離閉な函数体をもつ整正規スキームとする。
\begin{enumerate}
\item 分離エタール射 $U \to X$ は開埋め込みの直和である。
\item $X$ のすべての局所環は狭義 Hensel 環である。
\end{enumerate}
\end{lemma}

\begin{proof}
$R$ を、その分数体が分離代数閉である正規整域とする。
$R \to A$ をエタールな環写像とする。このとき $A \otimes_R K$ は
$K$-代数として有限積 $\prod_{i = 1, \ldots, n} K$、すなわち $K$ のコピーの積である。
$e_i$（$i = 1, \ldots, n$）を $A \otimes_R K$ の対応する冪等元とする。
$A$ は正規なので（Algebra, 補題 \ref{algebra-lemma-normal-goes-up}）、
冪等元 $e_i$ は $A$ に属する
（Algebra, 補題 \ref{algebra-lemma-normal-ring-integrally-closed}）。
したがって $A = \prod Ae_i$ であり、$A \otimes_R K = K$ と仮定してよい。
$A \subset A \otimes_R K = K$ なので（$R \to A$ の平坦性と
$R \subset K$ による）、$A$ は整域である。同じ議論により
$A \otimes_R A \subset (A \otimes_R A) \otimes_R K = K$.
写像 $A \otimes_R A \to A$ は全射であるだけでなく単射でもある。したがって
$R \to A$ は、Morphisms, 補題
\ref{morphisms-lemma-universally-injective} および \'Etale Morphisms, 定理
\ref{etale-theorem-etale-radicial-open} により開埋め込みを定める。

\medskip\noindent
$f : U \to X$ を分離エタール射とする。$\eta \in X$ を生成点とし、
$f^{-1}(\{\eta\}) = \{\xi_i\}_{i \in I}$ とおく。前段落の結果は次を示す。
$U' \subset U$ を、その $X$ における像があるアフィン開に含まれる任意の
アフィン開とすると、$U' = \coprod_{i \in I} U'_i$ である。ここで $U'_i$ は
$U'$ の点であって $\xi_i$ の特殊化であるものの集合である。さらに射 $U'_i \to X$ は
開埋め込みである。したがって $U_i = \overline{\{\xi_i\}}$ は $U$ の開かつ閉な
部分スキームであり、$U_i \to X$ は源上局所的に同型である。
Morphisms, 補題 \ref{morphisms-lemma-distinct-local-rings} により、
$U_i \to X$ が分離であることから $U_i \to X$ は単射である。
ゆえに $U_i \to X$ は開埋め込みであり、(1) が成り立つ。

\medskip\noindent
(2) は (1) と、$\mathcal{O}_{X, x}$ の狭義 Hensel 化を $\overline{x}$ における
局所環として記述すること（$X$ のエタールサイト上で。補題
\ref{etale-cohomology-lemma-describe-etale-local-ring}）から従う。直接証明することもできる。
Fundamental Groups, 補題
\ref{pione-lemma-normal-local-domain-separablly-closed-fraction-field} を参照。
\end{proof}

\begin{lemma}
\label{etale-cohomology-lemma-Rf-star-zero-normal-with-alg-closed-function-field}
$f : X \to Y$ をスキームの射とし、$X$ は分離閉な函数体をもつ整正規スキームとする。
このとき $R^qf_*\underline{M} = 0$ は $q > 0$ および任意のアーベル群 $M$ に対して
成り立つ。
\end{lemma}

\begin{proof}
$R^qf_*\underline{M}$ は、前層
$V \mapsto H^q_\etale(V \times_Y X, M)$ に付随する $Y_\etale$ 上の層であることを思い出す。
補題 \ref{etale-cohomology-lemma-higher-direct-images} を参照。$V$ がアフィンならば、
$V \times_Y X \to X$ は分離かつエタールである。したがって
$V \times_Y X = \coprod U_i$ は開部分スキーム $U_i$ の直和であり、各々は $X$ に含まれる。
補題 \ref{etale-cohomology-lemma-normal-scheme-with-alg-closed-function-field} を参照。
補題 \ref{etale-cohomology-lemma-local-rings-strictly-henselian} により、
$H^q_\etale(U_i, M)$ は $H^q_{Zar}(U_i, M)$ に等しい。これは Cohomology,
補題 \ref{cohomology-lemma-irreducible-constant-cohomology-zero} により消える。
\end{proof}

\begin{lemma}
\label{etale-cohomology-lemma-closed-of-affine-normal-scheme-with-alg-closed-function-field}
$X$ を分離閉な函数体をもつアフィン整正規スキームとする。
$Z \subset X$ を閉部分スキームとする。$V \to Z$ をエタール射とし、$V$ は
アフィンであるとする。このとき $V$ は $Z$ の開部分スキームの有限直和である。
$V \to Z$ が全射かつ有限エタールならば、$V \to Z$ は切断をもつ。
\end{lemma}

\begin{proof}
Algebra, 補題 \ref{algebra-lemma-lift-etale} により、$V$ をアフィンスキーム
$U$ へ持ち上げられ、これは $X$ 上エタールである。補題
\ref{etale-cohomology-lemma-normal-scheme-with-alg-closed-function-field} を $U \to X$ に
適用すると、最初の主張を得る。

\medskip\noindent
最後の主張は最初の主張から従う。
$V = \coprod_{i = 1, \ldots, n} V_i$ を開かつ閉な部分スキームへの有限分解とし、
$V_i \to Z$ は開埋め込みであるとする。$V \to Z$ は有限なので、
$V_i \to Z$ は閉でもある。$U_i \subset Z$ を像とする。このとき開かつ閉な
部分スキームへの分解
$$
Z =
\coprod\nolimits_{(A, B)}
\bigcap\nolimits_{i \in A} U_i \cap
\bigcap\nolimits_{i \in B} U_i^c
$$
を得る。ここで直和は $\{1, \ldots, n\} = A \amalg B$ であって、$A$ が少なくとも
一つの元をもつもの全体にわたる。各層は一つの $U_i$ に含まれ、切断が得られる。
\end{proof}

\begin{lemma}
\label{etale-cohomology-lemma-gabber-for-h1-absolutely-algebraically-closed}
$X$ を分離閉な函数体をもつ正規整アフィンスキームとする。
$Z \subset X$ を閉部分スキームとする。任意の有限アーベル群 $M$ に対して
$H^1_\etale(Z, \underline{M}) = 0$ である。
\end{lemma}

\begin{proof}
Cohomology on Sites, 補題 \ref{sites-cohomology-lemma-torsors-h1} により、
$H^1_\etale(Z, \underline{M})$ の元は $\underline{M}$-トーサー $\mathcal{F}$ に対応し、
これは $Z_\etale$ 上にある。このようなトーサーは明らかに有限局所定数層である。
したがって $\mathcal{F}$ はスキーム $V$ により表現可能であり、これは $Z$ 上有限エタールである。
補題 \ref{etale-cohomology-lemma-characterize-finite-locally-constant} を参照。トーサーは局所自明なので、
$V \to Z$ はもちろん全射である。補題
\ref{etale-cohomology-lemma-closed-of-affine-normal-scheme-with-alg-closed-function-field} により
$V \to Z$ は切断をもつので、証明が終わる。
\end{proof}

\begin{lemma}
\label{etale-cohomology-lemma-gabber-for-absolutely-algebraically-closed}
$X$ を分離閉な函数体をもつ正規整アフィンスキームとする。
$Z \subset X$ を閉部分スキームとする。任意の有限アーベル群 $M$ に対して
$H^q_\etale(Z, \underline{M}) = 0$ が $q \geq 1$ で成り立つ。
\end{lemma}

\begin{proof}
$X = \Spec(R)$ および $Z = \Spec(R')$ と書き、環の全射 $R \to R'$ を得る。
補題 \ref{etale-cohomology-lemma-normal-scheme-with-alg-closed-function-field} と Algebra, 補題
\ref{algebra-lemma-quotient-strict-henselization} により、$R'$ のすべての局所環は
狭義 Hensel 環である。さらに、任意の $f' \in R'$ に対し全射
$R_f \to R'_{f'}$ が存在することが分かる。ここで $f \in R$ は $f'$ の持ち上げである。
$R_f$ は分離閉な分数体をもつ正規整域なので、補題
\ref{etale-cohomology-lemma-gabber-for-h1-absolutely-algebraically-closed} により
$H^1_\etale(D(f'), \underline{M}) = 0$ である。したがって補題
\ref{etale-cohomology-lemma-gabber-criterion} を $Z = \Spec(R')$ に適用して結論を得る。
\end{proof}

\begin{lemma}
\label{etale-cohomology-lemma-integral-cover-trivial-cohomology}
$X$ をアフィンスキームとする。
\begin{enumerate}
\item 整全射 $X' \to X$ が存在し、すべての閉部分スキーム $Z' \subset X'$、
すべての有限アーベル群 $M$、およびすべての $q \geq 1$ に対して
$H^q_\etale(Z', \underline{M}) = 0$ となる。
\item 任意の閉部分スキーム $Z \subset X$、有限アーベル群 $M$、
$q \geq 1$、および $\xi \in H^q_\etale(Z, \underline{M})$ に対して、有限表示の
有限全射 $X' \to X$ が存在し、$\xi$ の $H^q_\etale(X' \times_X Z, \underline{M})$
への引き戻しは零となる。
\end{enumerate}
\end{lemma}

\begin{proof}
$X = \Spec(A)$ と書く。$A = \mathbf{Z}[x_i]/J$ と書き、ここで $J$ はあるイデアルとする。
$R$ を、$\mathbf{Z}[x_i]$ の分数体の代数閉包に
おける $\mathbf{Z}[x_i]$ の整閉包とする。$A' = R/JR$ とおき、
$X' = \Spec(A')$ とする。補題
\ref{etale-cohomology-lemma-gabber-for-absolutely-algebraically-closed} により、これは (1) の例を与える。

\medskip\noindent
(2) の証明。$X' \to X$ を上で見つけた整全射とする。もちろん $\xi$ は
$H^q_\etale(X' \times_X Z, \underline{M})$ で零に写る。$X'$ を極限
$X' = \lim X'_i$ と書ける。各項は $X$ 上有限かつ有限表示なスキームである。
現在のアフィンの場合には
これは容易だが、より一般的な Limits, 補題
\ref{limits-lemma-integral-limit-finite-and-finite-presentation} の特殊な場合である。
補題 \ref{etale-cohomology-lemma-directed-colimit-cohomology} により、$\xi$ は
$H^q_\etale(X'_i \times_X Z, \underline{M})$ で零に写る。これは十分大きいある
$i$ について成り立つ。
\end{proof}








\section{絶対整閉環の消失}
\label{etale-cohomology-section-aic-vanishing}

\noindent
環 $R$ が絶対整閉であるとは、$R$ 上のすべてのモニック多項式が
$R$ に根をもつことをいう（More on Algebra, 定義
\ref{more-algebra-definition-absolutely-integrally-closed}).
この節では、有限捩れ群を係数とする $\Spec(R)$ のエタール・コホモロジーが
正の次数で消えること（命題 \ref{etale-cohomology-proposition-aic-vanishing}）を示す。
これは先の補題
\ref{etale-cohomology-lemma-gabber-for-absolutely-algebraically-closed} を少し強めるものである。
読者にはこの節を飛ばすことを勧める。

\begin{lemma}
\label{etale-cohomology-lemma-find-extension}
$A$ を環とする。$a, b \in A$ が
$aA + bA = A$ を満たし、かつ $a \bmod bA$ が 1 の冪根であるとする。
このとき単生成拡大 $A \subset B$ と、$y \in B$ が存在し、
$u = a - by$ は単元となる。
\end{lemma}

\begin{proof}
ある $a^n \equiv 1 \bmod bA$ とする。特に、$a^i$ はすべての
$b^mA$ を法として、$i, m \geq 1$ に対して単元である。
次を満たす $a_1, \ldots, a_n \in A$ が存在すると主張する。
$$
1 = a^n + a_1 a^{n - 1}b + a_2 a^{n - 2}b^2 + \ldots + a_n b^n
$$
実際、$1 - a^n \in bA$ なので、$a_1 \in A$ を取って
$1 - a^n - a_1 a^{n - 1} b \in b^2A$ とできる。これは
$a^{n - 1}$ が $bA$ を法として単元であることによる。次に
$a_2 \in A$ を取って
$1 - a^n - a_1 a^{n - 1} b - a_2 a^{n - 2} b^2 \in b^3A$
とできる。同様に続ける。最終的に $a_1, \ldots, a_{n - 1} \in A$ を取って
$1 - (a^n + a_1 a^{n - 1}b + a_2 a^{n - 2}b^2 +
\ldots + a_{n - 1} ab^{n - 1}) \in b^nA$ とできる。したがって、表示した等式を
満たす $a_n \in A$ を取ることができる。

\medskip\noindent
上の $a_1, \ldots, a_n$ に対し、次のようにおけばよい。
$$
B = A[y]/(y^n + a_1 y^{n - 1} + a_2 y^{n - 2} + \ldots + a_n)
$$
実際、$\mathfrak q \subset B$ を $\mathfrak p \subset A$ の上にある
素イデアルとして取る。矛盾を得るため $u = a - by$ が
$\mathfrak q$ に属すると仮定する。もし $b \in \mathfrak p$ ならば
$a \not \in \mathfrak p$ であり、$aA + bA = A$ だから、したがって
$u$ は $\mathfrak q$ に属さない。よって $b \not \in \mathfrak p$、
すなわち $b \not \in \mathfrak q$ としてよい。これは
$y \bmod \mathfrak q$ が $a/b \bmod \mathfrak q$ に等しいことを意味する。
すると
$$
0 = y^n + a_1 y^{n - 1} + a_2 y^{n - 2} + \ldots + a_n =
b^{-n}(a^n + a_1 a^{n - 1}b + a_2 a^{n - 2}b^2 + \ldots + a_nb^n) =
b^{-n}
$$
となり矛盾である。これで証明が終わる。
\end{proof}

\noindent
証明を説明するため、いくつかの群スキームを導入する。
素数 $\ell$ を固定し、次のようにおく。
$$
A = \mathbf{Z}[\zeta] =
\mathbf{Z}[x]/(x^{\ell - 1} + x^{\ell - 2} + \ldots + 1)
$$
言い換えれば、$A$ は $\mathbf{Z}$ の単生成拡大であり、原始
$\ell$ 乗根 $\zeta$ によって生成される。次をおく。
$$
\pi = \zeta - 1
$$
計算（省略）によれば、$\ell$ は $\pi^{\ell - 1}$ で割り切れる。この主張は
$A$ におけるものである。
$A$ 上の最初の群スキームを
$$
G = \Spec(A[s, \frac{1}{\pi s + 1}])
$$
群法則は次の余乗法で与える。
$$
\mu :
A[s, \frac{1}{\pi s + 1}]
\longrightarrow
A[s, \frac{1}{\pi s + 1}] \otimes_A A[s, \frac{1}{\pi s + 1}],\quad
s \longmapsto \pi s \otimes s + s \otimes 1 + 1 \otimes s
$$
この選択により
$$
\mu(\pi s + 1) = (\pi s + 1) \otimes (\pi s + 1)
$$
となるので、確かに表示したような $A$-代数写像を得る。
これが実際に群法則を定めることの検証は省略する。$A$ 上の二つ目の
群スキームを
$$
H = \Spec(A[t, \frac{1}{\pi^\ell t + 1}])
$$
群法則は次の余乗法で与える。
$$
\mu : A[t, \frac{1}{\pi^\ell t + 1}]
\longrightarrow
A[t, \frac{1}{\pi^\ell t + 1}] \otimes_A A[t, \frac{1}{\pi^\ell t + 1}],\quad
t \longmapsto \pi^\ell t \otimes t + t \otimes 1 + 1 \otimes t
$$
先ほどと同じ検証により、これは群法則を定める。次に多項式
$$
\Phi(s) = \frac{(\pi s + 1)^\ell - 1}{\pi^\ell}
$$
$A[s]$ に属し、次数 $\ell$ で $s$ に関してモニックである。実際、
$s^i$ の $0 < i < \ell$ における係数は ${\ell \choose i}\pi^{i - \ell}$
に等しく、$\pi^{\ell - 1}$ が $\ell$ を $A$ で割り切るので、これは
$A$ の元である。次の環写像を得る。
$$
A[t, \frac{1}{\pi^\ell t + 1}]
\longrightarrow
A[s, \frac{1}{\pi s + 1}],\quad
t \longmapsto \Phi(s)
$$
これは余乗法と両立することを容易に確認できる。したがって群スキームの射
$$
f : G \to H
$$
次の補題は特に、この射が忠実平坦であることを示す（実際には有限エタール
全射であることが分かる）。

\begin{lemma}
\label{etale-cohomology-lemma-monogenic-one}
次が成り立つ。
$$
A[s, \frac{1}{\pi s + 1}] =
\left(A[t, \frac{1}{\pi^\ell t + 1}]\right)[s]/(\Phi(s) - t)
$$
特に、$G$ のホップ代数は $H$ のホップ代数の単生成拡大である。
\end{lemma}

\begin{proof}
これは上の議論と $\Phi(s)$ の形から従う。特に $\Phi(s) = t$ を用いると
$\frac{1}{\pi^\ell t + 1}$ は $\frac{1}{(\pi s + 1)^\ell}$ となることに
注意する。
\end{proof}

\noindent
次に $f$ の核を計算する。$H$ の原点は $t = 0$ で $H$ に与えられ、
したがって $f$ の核は $\Phi(s) = 0$ で与えられる。ここで、次式で与えられる
$A$ 値点 $\sigma_0, \ldots, \sigma_{\ell - 1}$ を $G$ 上で考える。
$$
\sigma_i :
s = \frac{\zeta^i - 1}{\pi} = \frac{\zeta^i - 1}{\zeta - 1} =
\zeta^{i - 1} + \zeta^{i - 2} + \ldots + 1,\quad
i = 0, 1, \ldots, \ell - 1
$$
これらは明らかに $\Ker(f)$ に含まれる。さらに、これらは
$G \to \Spec(A)$ の {\bf すべての}ファイバーで互いに異なる。また
$\sigma_i +_G \sigma_j = \sigma_{i + j \bmod \ell}$ であることも計算できる。
したがって群スキームの閉埋め込み
$$
\underline{\mathbf{Z}/\ell \mathbf{Z}}_A \longrightarrow \Ker(f)
$$
これは $i$ を $\sigma_i$ に送る。ところが構成により $\Ker(f)$ は
$\Spec(A)$ 上次数 $\ell$ の有限平坦である。よってこの写像は同型である。
以上から、$A$ 上の群スキームの短完全系列
\begin{equation}
\label{etale-cohomology-equation-ses}
0 \to
\underline{\mathbf{Z}/\ell \mathbf{Z}}_A
\to G
\to H
\to 0
\end{equation}
を得る。

\begin{lemma}
\label{etale-cohomology-lemma-lift-points-H-to-G}
$R$ を絶対整閉な $A$ 代数とする。このとき
$G(R) \to H(R)$ は全射である。
\end{lemma}

\begin{proof}
$h \in H(R)$ に対応する $A$ 代数写像
$A[t, \frac{1}{\pi^\ell t + 1}] \to R$ で $t$ を $a \in A$ に送るものを取る。
$\Phi(s)$ はモニックなので、$b \in A$ で $\Phi(b) = a$ を満たすものを取れる。
補題 \ref{etale-cohomology-lemma-monogenic-one} により $s$ を $b$ に送ると、
$A$ 代数写像 $A[s, \frac{1}{\pi s + 1}] \to R$ を得る。
これは上の $A[t, \frac{1}{\pi^\ell t + 1}] \to R$ と両立する一意な写像である。
これは $g \in G(R)$ に対応し、$h \in H(R)$ に写る元である。
\end{proof}

\begin{lemma}
\label{etale-cohomology-lemma-interpolate}
$R$ を絶対整閉な $A$ 代数とする。$I, J \subset R$ を
$I + J = R$ を満たすイデアルとする。
$g \in G(R)$ で $g \bmod I = \sigma_0$ かつ
$g \bmod J = \sigma_1$ となるものが存在する。
\end{lemma}

\begin{proof}
$x \in I$ を $x \equiv 1 \bmod J$ となるように取る。
$I$ を $xR$ に、$J$ を $(x - 1)R$ に置き換えてよい。
すると、次を満たす $s \in R$ を探すことになる。
\begin{enumerate}
\item $1 + \pi s$ は単元である、
\item $s \equiv 0 \bmod xR$、かつ
\item $s \equiv 1 \bmod (x - 1)R$.
\end{enumerate}
後二つの条件は、$s = x + x(x - 1)y$ となるある $y \in R$ が
存在することを意味する。第一の条件は
$1 + \pi s = 1 + \pi x + \pi x (x - 1) y$ が $R$ の単元となること
である。しかし $1 + \pi x$ と $\pi x (x - 1)$ は
$R$ の単位イデアルを生成し、$1 + \pi x$ は $\ell$ 乗すると $1$ となり、
$\pi x (x - 1)$ を法とすることに注意する
\footnote{$1 + \pi x$ は $1$ に $\pi$ を法として合同し、$1$ に $x$ を法として合同し、
さらに $1 + \pi = \zeta$ に $x - 1$ を法として合同する。また
$(\pi) \cap (x) \cap (x - 1) = (\pi x (x - 1))$ が $A[x]$ で成り立つ。}。
したがって補題 \ref{etale-cohomology-lemma-find-extension} と $R$ が絶対整閉であることから
結論を得る。
\end{proof}

\begin{proposition}
\label{etale-cohomology-proposition-aic-vanishing}
$R$ を絶対整閉環とする。$M$ を有限アーベル群とする。このとき
$H^i_\etale(\Spec(R), \underline{M}) = 0$ が $i > 0$ で成り立つ。
\end{proposition}

\begin{proof}
任意の有限アーベル群は、素数位数の巡回群を商とする有限フィルトレーションを
もつので、$M = \mathbf{Z}/\ell\mathbf{Z}$ で $\ell$ が素数である場合に
帰着してよい。

\medskip\noindent
すべての $R$ の局所環は狭義ヘンゼル環である
More on Algebra, 補題
\ref{more-algebra-lemma-absolutely-integrally-closed-strictly-henselian} を参照。
さらに、$R$ の任意の局所化も絶対整閉である（More on Algebra, 補題
\ref{more-algebra-lemma-absolutely-integrally-closed-quotient-localization}）。
したがって補題 \ref{etale-cohomology-lemma-gabber-criterion} により、次の写像の核が
$$
H^1_\etale(D(f + g), \mathbf{Z}/\ell\mathbf{Z})
\longrightarrow
H^1_\etale(D(f(f + g)), \mathbf{Z}/\ell\mathbf{Z}) \oplus
H^1_\etale(D(g(f + g)), \mathbf{Z}/\ell\mathbf{Z})
$$
任意の $f, g \in R$ に対して零であることを示せば十分である。
$R$ を $R_{f + g}$ で置き換えると、次の主張に帰着する：
$\xi \in H^1_\etale(\Spec(R), \mathbf{Z}/\ell\mathbf{Z})$ と
アフィン開被覆 $\Spec(R) = U \cup V$ が与えられ、
$\xi|_U$ と $\xi|_V$ が自明ならば $\xi = 0$ である。

\medskip\noindent
$A = \mathbf{Z}[\zeta]$ を上のように取る。
$\mathbf{Z} \subset A$ は単生成なので、環写像 $A \to R$ を取れる。
以下、$R$ を $A$ 代数とみなし、$\Spec(R)$ を $\Spec(A)$ 上の
スキームとみなす。短完全系列 (\ref{etale-cohomology-equation-ses}) を $\Spec(R)$
に基底変換し、エタール・コホモロジーを取ると
$$
G(R) \to H(R) \to
H^1_\etale(\Spec(R), \mathbf{Z}/\ell\mathbf{Z}) \to
H^1_\etale(\Spec(R), G)
$$
証明の残りではこのことを用いる。

\medskip\noindent
$\tau \in \Gamma(U \cap V, \mathbf{Z}/\ell\mathbf{Z})$ を、Mayer--Vietoris
系列（補題 \ref{etale-cohomology-lemma-mayer-vietoris}）における境界が $\xi$ を与える
切断とする。$i = 0, 1, \ldots, \ell - 1$ に対し、
$A_i \subset U \cap V$ を、$\tau$ の値が $i \bmod \ell$ となる
開かつ閉部分集合とする。したがって有限な互いに素な分解
$$
U \cap V = A_0 \amalg \ldots \amalg A_{\ell - 1}
$$
となり、$\tau$ は各 $A_i$ 上で定数である。$i = 0, 1, \ldots, \ell - 1$
に対し、$\tau_i \in H^0(U \cap V, \mathbf{Z}/\ell\mathbf{Z})$ を、
$1$ に等しい元（$A_i$ 上）であり、かつ $0$ に等しい元
（$A_j$ 上、$j \not = i$）として定める。
すると $\tau$ は $\tau_i$ の倍数の和である
\footnote{計算誤差を除けば $\tau = \sum i \tau_i$ である。}。
したがって $\tau_i$ に対応するコホモロジー類が自明であることを示せばよい。
これは $\tau$ が二つの異なる値、すなわち $1$ と $0$ だけを取る場合に帰着する。

\medskip\noindent
$\tau$ が値 $1$ と $0$ だけを取ると仮定する。次のように書く。
$$
U \cap V = A \amalg B
$$
ここで $A$ は $\tau = 0$ となる軌跡、$B$ は $\tau = 1$ となる軌跡である。
したがって $A$ と $B$ は互いに素な閉部分集合である。
$\overline{A}$ と $\overline{B}$ を、$A$ と $B$ の
$\Spec(R)$ における閉包と書く。するといわゆる「バナナ」の形が現れる。すなわち
次のようになる。
$$
\overline{A} \cap \overline{B} = Z_1 \amalg Z_2
$$
ここで $Z_1 \subset U$ および $Z_2 \subset V$ は互いに素な閉部分集合である。
$T_1 = \Spec(R) \setminus V$、$T_2 = \Spec(R) \setminus U$ とおく。
$Z_1 \subset T_1 \subset U$、$Z_2 \subset T_2 \subset V$、および
$T_1 \cap T_2 = \emptyset$ に注意する。位相空間として次のように書ける。
$$
\Spec(R) = \overline{A} \cup \overline{B} \cup T_1 \cup T_2
$$
これを可視化するために図を描くとよい。
$\xi$ が零であることを示すため、命題
\ref{etale-cohomology-proposition-topological-invariance} により $R$ をその還元に取り替えてよい。
以下では $A$、$\overline{A}$、$B$、$\overline{B}$、$T_1$、$T_2$ を
$\Spec(R)$ の被約閉部分スキームとみなす。次に $Z_1$ と $Z_2$ のスキーム構造として
次を用いる。
$$
Z_1 = \overline{A} \cap (\overline{B} \cup T_1)
\quad\text{かつ}\quad
Z_2 = \overline{A} \cap (\overline{B} \cup T_2)
$$
(Morphisms, 定義
\ref{morphisms-definition-scheme-theoretic-intersection-union} にある
スキーム論的な和と交差である)。

\medskip\noindent
$X$ を $G$-トーサーとし、$\Spec(R)$ 上で $\xi$ の
$H^1(\Spec(R), G)$ における像に対応するものとする。$X$ が自明ならば、$\xi$ は
上のコホモロジー完全系列における元 $h \in H(R)$ から来る。
しかし補題 \ref{etale-cohomology-lemma-lift-points-H-to-G} により、元 $h$ は
$G(R)$ の元へ持ち上がるので、望みどおり $\xi = 0$ となる。
したがって、$X$ が自明であることを示せばよい。

\medskip\noindent
思い出しておくと、埋め込み $\mathbf{Z}/\ell \mathbf{Z} \to G(R)$ は
$i \bmod \ell$ を $\sigma_i \in G(R)$ に送る。$\overline{A}$ は絶対整閉環
のスペクトル（すなわち $R$ の商）であることに注意する。補題
\ref{etale-cohomology-lemma-interpolate} により、次を満たす $g \in G(\overline{A})$ を取れる：
$g|_{\overline{A} \cap Z_1} = \sigma_0$ かつ
$g|_{\overline{A} \cap Z_2} = \sigma_1$（スキーム論的な意味で）。
そこで次を定める。
\begin{enumerate}
\item $g_1 \in G(U)$ は、$g$ を $\overline{A} \cap U$ 上で、
$\sigma_0$ を $\overline{B} \cap U$ 上で、さらに $\sigma_0$ を $T_1$ 上で取るもの；
\item $g_2 \in G(V)$ は、$g$ を $\overline{A} \cap V$ 上で、
$\sigma_1$ を $\overline{B} \cap V$ 上で、さらに $\sigma_1$ を $T_2$ 上で取るもの。
\end{enumerate}
具体的には、(1) の $g_1$ を得るため、$\sigma_0$ を
$\Omega = (\overline{B} \cup T_1) \cap U$ 上で、$g$ の
$\Omega' = \overline{A} \cap U$ 上への制限と貼り合わせる。
$U = \Omega \cup \Omega'$（スキーム論的な意味で）であることに注意する。
$U$ は被約であり、$\Omega \cap \Omega' = Z_1$（スキーム論的な意味で）
となるのは $Z_1$ の選び方による。したがって Morphisms, 補題
\ref{morphisms-lemma-scheme-theoretic-union} により、$U$ は
$\Omega$ と $\Omega'$ の $Z_1$ に沿った押し出しである。ゆえに $g_1$ が得られる。
(2) の $g_2$ の存在も同様である。すると
$$
\tau = g_2|_{A \cup B} - g_1|_{A \cup B} \quad(\text{群律における加法})
$$
となり、$X$ が自明であることが分かる。これで証明が終わる。
\end{proof}




















\section{固有基底変換のアフィン版}
\label{etale-cohomology-section-gabber-affine-proper}

\noindent
この節ではオフェル・ガバーの結果を論じる。文献
\cite{gabber-affine-proper} を参照されたい。この結果はローラント・フーバー
によっても証明されている（\cite{Huber-henselian}）。Gabber の証明に
必要な作業の一部は、すでに節 \ref{etale-cohomology-section-strictly-henselian-local-rings}
で行った。

\begin{lemma}
\label{etale-cohomology-lemma-efface-cohomology-on-closed-by-finite-cover}
$X$ をアフィン・スキームとする。$\mathcal{F}$ を $X_\etale$ 上の
捩れアーベル層とし、$Z \subset X$ を閉部分スキームとする。
$\xi \in H^q_\etale(Z, \mathcal{F}|_Z)$ とし、ある $q > 0$ に対して
考える。このとき単射 $\mathcal{F} \to \mathcal{F}'$
（$X_\etale$ 上の捩れアーベル層の間のもの）であって、$\xi$ の
$H^q_\etale(Z, \mathcal{F}'|_Z)$ における像が零となるものが存在する。
\end{lemma}

\begin{proof}
補題 \ref{etale-cohomology-lemma-torsion-colimit-constructible} と \ref{etale-cohomology-lemma-colimit}
により、写像 $\mathcal{G} \to \mathcal{F}$ であって、
$\mathcal{G}$ が構成可能アーベル層であり、$\xi$ が元 $\zeta$
（$H^q_\etale(Z, \mathcal{G}|_Z)$ の元）から来るものを取れる。
単射 $\mathcal{G} \to \mathcal{G}'$
（$X_\etale$ 上の捩れアーベル層の間のもの）であって、$\zeta$ の
$H^q_\etale(Z, \mathcal{G}'|_Z)$ における像が零となるものを取れたとする。
このとき $\mathcal{F}'$ として押し出し
$$
\mathcal{F}' = \mathcal{G}' \amalg_{\mathcal{G}} \mathcal{F}
$$
を取れば補題の結論が従う。（$Z$ への制限は完全なので有限極限および
有限余極限と可換であり、さらに順像の左随伴であるため任意の余極限とも
可換であることに注意する。）したがって $\mathcal{F}$ は構成可能であると
仮定してよい。

\medskip\noindent
$\mathcal{F}$ が構成可能であると仮定する。補題
\ref{etale-cohomology-lemma-constructible-maps-into-constant-general} により、
$\mathcal{F}$ が $f_*\underline{M}$ の形である場合を証明すれば十分である。
ここで $M$ は有限アーベル群、$f : Y \to X$ は有限表示の有限射である
（このような層も補題 \ref{etale-cohomology-lemma-finite-pushforward-constructible}
により構成可能だが、これは用いない）。
$f_*$ の形成は任意の基底変換と可換なので
（補題 \ref{etale-cohomology-lemma-finite-pushforward-commutes-with-base-change}）、
$f_*\underline{M}$ の $Z$ への制限は、$\underline{M}$ の
$Y \times_X Z \to Z$ による順像に等しい。Leray スペクトル系列
（命題 \ref{etale-cohomology-proposition-leray}）と高次順像の消失
（命題 \ref{etale-cohomology-proposition-finite-higher-direct-image-zero}）から
$$
H^q_\etale(Z, f_*\underline{M}|_Z) = H^q_\etale(Y \times_X Z, \underline{M}).
$$
を得る。補題 \ref{etale-cohomology-lemma-integral-cover-trivial-cohomology} により、
有限表示の有限全射 $Y' \to Y$ であって、$\xi$ が
$H^q(Y' \times_X Z, \underline{M})$ で零に写るものを取れる。
$f' : Y' \to X$ を合成 $Y' \to Y \to X$ と書く。このとき写像
$$
f_*\underline{M} \longrightarrow f'_*\underline{M}
$$
は単射であると主張する。上で述べたことから、これで証明は終わる。
所望の単射性は茎で調べればよい。実際、
$\overline{x} : \Spec(k) \to X$ を幾何学的点とすると
$$
(f_*\underline{M})_{\overline{x}} =
\bigoplus\nolimits_{f(\overline{y}) = \overline{x}} M
$$
である（命題 \ref{etale-cohomology-proposition-finite-higher-direct-image-zero}）。
もう一方の層についても同様である。$Y' \to Y$ は有限全射なので、
$\overline{x}$ を持ち上げる幾何学的点の間に誘導される写像も全射であり、
結論が従う。
\end{proof}

\noindent
上の補題は Gabber の結果における高次コホモロジー群を扱う。
大域切断を扱うためには次の補題を用いる。

\begin{lemma}
\label{etale-cohomology-lemma-gabber-h0}
$X$ を準コンパクトかつ準分離なスキームとし、
$i : Z \to X$ を閉埋め込みとする。次を仮定する。
\begin{enumerate}
\item 任意の層 $\mathcal{F}$（$X_{Zar}$ 上）に対し、写像
$\Gamma(X, \mathcal{F}) \to \Gamma(Z, i^{-1}\mathcal{F})$
は全単射である。
\item 任意の有限射 $X' \to X$ に対し、$Z \times_X X' \to X'$
について仮定 (1) が成り立つ。
\end{enumerate}
このとき任意の層 $\mathcal{F}$（$X_\etale$ 上）に対し
$\Gamma(X, \mathcal{F}) = \Gamma(Z, i^{-1}_{small}\mathcal{F})$.
\end{lemma}

\begin{proof}
$\mathcal{F}$ を $X_\etale$ 上の層とする。標準的な（基底変換）写像
$$
i^{-1}(\mathcal{F}|_{X_{Zar}})
\longrightarrow
(i_{small}^{-1}\mathcal{F})|_{Z_{Zar}}
$$
が $Z_{Zar}$ 上の層の間にある。この写像が単射であることを茎で示す。
左辺の $z \in Z$
における茎は、$\mathcal{F}|_{X_{Zar}}$ の $z$ における茎である。
右辺の茎は、すべての基本エタール近傍 $(U, u) \to (X, z)$ であって、
$U \times_X Z \to Z$ が $z$ の近傍上で切断をもつものにわたる余極限である。
エタール射は開写像なので、$U \to X$ の像を $U_0$ と書けば、
これは $z$ の $X$ における開近傍である。写像
$\mathcal{F}(U_0) \to \mathcal{F}(U)$ は、$\mathcal{F}$ の層条件を
エタール被覆 $U \to U_0$ に適用することにより単射である。
すべての $U$ と $U_0$ にわたり余極限を取れば、茎上の単射性を得る。

\medskip\noindent
これと仮定 (1) から、写像
$\Gamma(X, \mathcal{F}) \to \Gamma(Z, i^{-1}_{small}\mathcal{F})$
は単射である。(2) により、すべての $X'$（$X$ 上有限）についても
同じことが成り立つ。

\medskip\noindent
全射性を示す前に記法を導入する。任意の対象 $U$（$X_\etale$ のもの）と
$s \in \mathcal{F}(U)$ に対し、$i^{-1}s \in
(i_{small}^{-1}\mathcal{F})(U \times_X Z)$ で引き戻しを表す
（より正確には、これは Sites, 補題 \ref{sites-lemma-recover} の写像と、
前層の引き戻しから層の引き戻しへの写像を合成したものによる $s$ の像である）。
この構成は $X_\etale$ の射による引き戻し、および（有限）射
$X' \to X$ による引き戻しと可換である。詳細の明確化は読者に委ねる。

\medskip\noindent
$t \in \Gamma(Z, i^{-1}_{small}\mathcal{F})$ とする。
$i^{-1}_{small}\mathcal{F}$ の構成により、エタール被覆
$\{V_j \to Z\}_{j \in J}$、エタール射 $U_j \to X$、切断
$s_j \in \mathcal{F}(U_j)$、および射 $V_j \to U_j$（$X$ 上）であって、
$t|_{V_j}$ が $i^{-1}s_j$ の $V_j \to U_j \times_X Z$ による
引き戻しとなるものが存在する。$V_j \to U_j \times_X Z$ はエタールなので、
その像は開部分スキームである。$U_j \times_X Z \subset U_j$ は閉なので、
$U_j$ を開部分スキームで置き換えた後、$V_j \to U_j \times_X Z$ は
全射であると仮定してよい。このとき $\{V_j \to U_j \times_X Z\}$ は
エタール被覆であり、$t|_{U_j \times_X Z} = i^{-1}s_j$ と結論できる。
空でない任意の閉部分スキーム $T \subset X$ は $Z$ と交わる。
これは例えば層 $(T \to X)_*\underline{\mathbf{Z}}$ に仮定 (1) を
適用すれば分かる。
したがって $\coprod U_j \to X$ は全射である。More on Morphisms, 補題
\ref{more-morphisms-lemma-there-is-a-scheme-integral-over}
により、有限全射 $X' \to X$ であって、$X' \to X$ がザリスキー局所的に
$\coprod U_j \to X$ を経由するものを取れる。
$X' = \bigcup_{\alpha \in A} W_\alpha$ を開被覆とし、
$j : A \to J$ を写像として、射
$g_\alpha : W_\alpha \to U_{j(\alpha)}$（$X$ 上）が存在するとする。
$\mathcal{F}'$ で $\mathcal{F}$ の $X'_\etale$ への引き戻しを表し、
$s'_\alpha \in \mathcal{F}'(W_\alpha)$ で $s_{j(\alpha)}$ の
$g_\alpha$ による引き戻しを表す。
$Z' = X' \times_X Z$ とおき、$i' : Z' \to X'$ で $i$ の基底変換を表す。
このとき $(i'_{small})^{-1}\mathcal{F}' =
(Z' \to Z)_{small}^{-1}i_{small}^{-1}\mathcal{F}$ である。
この同一視の下で、$t' \in
\Gamma(Z', (i')^{-1}_{small}\mathcal{F}')$
により $t$ の $Z'$ への引き戻しを表す。構成により
$t'|_{W_\alpha \times_X Z}$ は
$(i')^{-1}s'_\alpha$ に等しい。したがって $t'$ は部分層
$$
(i')^{-1}(\mathcal{F}'|_{X'_{Zar}})
\subset
(i'_{small})^{-1}\mathcal{F}')|_{Z'_{Zar}}
$$
の切断である。仮定 (2) により、切断 $t'$ は切断
$s' \in \Gamma(X', \mathcal{F}'|_{X'_{Zar}}) = \Gamma(X', \mathcal{F}')$
から来る。第二段落で示した単射性により、$s'$ の
$X' \times_X X'$ への二つの引き戻しは等しい
（実際、$t'$ の $Z' \times_Z Z'$ への引き戻しについてそうである）。
ゆえに注記 \ref{etale-cohomology-remark-cohomological-descent-finite} により、
$s'$ は $\mathcal{F}$ の $X$ 上の切断から来る。
\end{proof}

\begin{lemma}
\label{etale-cohomology-lemma-connected-topological}
$Z \subset X$ を位相空間 $X$ の閉部分集合とする。次を仮定する。
\begin{enumerate}
\item $X$ はスペクトル空間である
（Topology, 定義 \ref{topology-definition-spectral-space}）。
\item $x \in X$ に対し、交わり $Z \cap \overline{\{x\}}$ は
連結である（特に空でない）。
\end{enumerate}
$Z = Z_1 \amalg Z_2$ であり、$Z_i$ が $Z$ で閉であるならば、
分解 $X = X_1 \amalg X_2$ であって、$X_i$ が $X$ で閉かつ
$Z_i = Z \cap X_i$ となるものが存在する。
\end{lemma}

\begin{proof}
$Z_i$ は準コンパクトである。したがって点の集合 $W_i$
（$Z_i$ へ特殊化する点からなる）は、Topology, 補題
\ref{topology-lemma-make-spectral-space} により構成可能位相で閉である。
仮定 (2) から $X = W_1 \amalg W_2$ となる。
$x \in \overline{W_1}$ とする。Topology, 補題
\ref{topology-lemma-constructible-stable-specialization-closed} の (1) により、
特殊化 $x_1 \leadsto x$ であって $x_1 \in W_1$ となるものが存在する。
したがって $\overline{\{x\}} \subset \overline{\{x_1\}}$ であり、
$x \in W_1$ となる。よって $X_i = W_i$ とおけばよい。
\end{proof}

\begin{lemma}
\label{etale-cohomology-lemma-h0-topological}
$Z \subset X$ を位相空間 $X$ の閉部分集合とする。次を仮定する。
\begin{enumerate}
\item $X$ はスペクトル空間である
（Topology, 定義 \ref{topology-definition-spectral-space}）。
\item $x \in X$ に対し、交わり $Z \cap \overline{\{x\}}$ は
連結である（特に空でない）。
\end{enumerate}
このとき任意の層 $\mathcal{F}$（$X$ 上）に対し
$\Gamma(X, \mathcal{F}) = \Gamma(Z, \mathcal{F}|_Z)$.
\end{lemma}

\begin{proof}
$x \leadsto x'$ が点の特殊化ならば、標準写像
$\mathcal{F}_{x'} \to \mathcal{F}_x$ がある。この写像は開集合上の切断と
両立し、$\mathcal{F}$ に関して関手的である。$X$ の各点は
$Z$ の点へ特殊化するので、
$\Gamma(X, \mathcal{F}) \to \Gamma(Z, \mathcal{F}|_Z)$ は単射である。
難しいのは全射性である。

\medskip\noindent
$\mathcal{B}$ で $X$ のすべての準コンパクト開集合の集合を表す。
$\mathcal{F}$ をフィルター余極限
$\mathcal{F} = \colim \mathcal{F}_i$ と書く。ここで各 $\mathcal{F}_i$ は
Modules, 式 (\ref{modules-equation-towards-constructible-sets}) の形である。
Modules, 補題 \ref{modules-lemma-filtered-colimit-constructibles} を参照されたい。
$\mathcal{F}|_Z = \colim \mathcal{F}_i|_Z$ である。実際、$Z$ への制限は
左随伴である（Categories, 補題 \ref{categories-lemma-adjoint-exact}
および Sheaves, 補題 \ref{sheaves-lemma-f-map}）。Sheaves, 補題
\ref{sheaves-lemma-directed-colimits-sections} により、
関手 $\Gamma(X, -)$ と $\Gamma(Z, -)$ はフィルター余極限と可換である。
したがって層 $\mathcal{F}$ は Modules, 式
(\ref{modules-equation-towards-constructible-sets}) の形であると仮定してよい。

\medskip\noindent
埋め込み $\mathcal{F} \subset \mathcal{G}$ があるとする。このとき
$$
\Gamma(X, \mathcal{F}) =
\Gamma(Z, \mathcal{F}|_Z) \cap \Gamma(X, \mathcal{G})
$$
であり、交わりは $\Gamma(Z, \mathcal{G}|_Z)$ の中で取る。
これは証明の最初の注意から従う。実際、$\mathcal{G}$ の大域切断が
$\mathcal{F}$ に属するかどうかは茎で判定でき、$X$ の各点は $Z$ の
点へ特殊化する。

\medskip\noindent
Modules, 補題 \ref{modules-lemma-constructible-in-constant} により、単射
$\mathcal{F} \to \prod (Z_i \to X)_*\underline{S_i}$ が存在する。
ここで積は有限であり、$Z_i \subset X$ は閉、$S_i$ は有限である。
したがって層 $(Z_i \to X)_*\underline{S_i}$ について全射性を
示せば十分である。次に注意する。
$$
\Gamma(X, (Z_i \to X)_*\underline{S_i}) = \Gamma(Z_i, \underline{S_i})
\quad\text{かつ}\quad
\Gamma(X, (Z_i \to X)_*\underline{S_i}|_Z) =
\Gamma(Z \cap Z_i, \underline{S_i})
$$
さらに条件 (1) と (2) は $Z_i$ に継承される。(2) については明らかであり、
(1) については Topology, 補題 \ref{topology-lemma-spectral-sub} から従う。
ゆえに（有限）定数層の場合に補題を示せば十分である。この場合は補題
\ref{etale-cohomology-lemma-connected-topological} の言い換えであり、証明が終わる。
\end{proof}

\begin{example}
\label{etale-cohomology-example-quinard}
$X$ がスペクトル空間でなければ、補題 \ref{etale-cohomology-lemma-h0-topological} は
偽である。次が反例である。$Y$ を $T_1$ 位相空間とし、
$y \in Y$ を開点でない点とする。$X = Y \amalg \{ x \}$ に、
閉集合が $\emptyset$、$\{y\}$、およびすべての
$F \amalg \{ x \}$ となる位相を入れる。ただし $F$ は $Y$ の
閉部分集合である。このとき $Z = \{x, y\}$ は $X$ の閉部分集合であり、
補題 \ref{etale-cohomology-lemma-h0-topological} の仮定 (2) を満たす。しかし $X$ は
連結である一方、$Z$ は連結でない。したがって値 $\{0, 1\}$ を取る
$X$ 上の定数層について、補題の結論は成り立たない。
\end{example}

\begin{lemma}
\label{etale-cohomology-lemma-h0-henselian-pair}
$(A, I)$ をヘンゼル対とする。$X = \Spec(A)$ および
$Z = \Spec(A/I)$ とおく。任意の層 $\mathcal{F}$（$X_\etale$ 上）に対し
$\Gamma(X, \mathcal{F}) = \Gamma(Z, \mathcal{F}|_Z)$ が成り立つ。
\end{lemma}

\begin{proof}
任意の環のスペクトルはスペクトル空間であることを思い出す
（Algebra, 補題 \ref{algebra-lemma-spec-spectral}）。More on Algebra, 補題
\ref{more-algebra-lemma-irreducible-henselian-pair-connected}
により、$\overline{\{x\}} \cap Z$ はすべての $x \in X$ に対し連結である。
補題 \ref{etale-cohomology-lemma-h0-topological} により、この主張は $X_{Zar}$ 上の層について
成り立つ。任意の有限射 $X' \to X$ に対し、
$X' = \Spec(A')$ かつ $Z \times_X X' = \Spec(A'/IA')$ であり、
$(A', IA')$ はヘンゼル対である。More on Algebra, 補題
\ref{more-algebra-lemma-integral-over-henselian-pair}
を参照されたい。したがって $(X')_{Zar}$ 上の層についても同じ主張を得る。
これで補題 \ref{etale-cohomology-lemma-gabber-h0} を適用して結論できる。
\end{proof}

\noindent
最後に Gabber の定理を述べ、証明する。

\begin{theorem}[Gabber]
\label{etale-cohomology-theorem-gabber}
$(A, I)$ をヘンゼル対とする。$X = \Spec(A)$ および
$Z = \Spec(A/I)$ とおく。任意の捩れアーベル層 $\mathcal{F}$
（$X_\etale$ 上）に対し
$H^q_\etale(X, \mathcal{F}) = H^q_\etale(Z, \mathcal{F}|_Z)$ が成り立つ。
\end{theorem}

\begin{proof}
補題 \ref{etale-cohomology-lemma-h0-henselian-pair} により、結果は $q = 0$ について成り立つ。
$q \geq 1$ とし、すべての次数 $< q$ について結果が示されていると仮定する。
$\mathcal{F}$ を捩れアーベル層とする。捩れアーベル層の単射
$\mathcal{F} \to \mathcal{F}'$
（後で選ぶ）を取り、その余核を $\mathcal{Q}$ とする。すると短完全系列
$$
0 \to \mathcal{F} \to \mathcal{F}' \to \mathcal{Q} \to 0
$$
を得る。これは $X_\etale$ 上の捩れアーベル層の系列である。
さらに $X$ と $Z$ 上のコホモロジー長完全系列の間の写像を与え、
その一部は次の形である。
$$
\xymatrix{
H^{q - 1}_\etale(X, \mathcal{F}') \ar[d] \ar[r] &
H^{q - 1}_\etale(X, \mathcal{Q}) \ar[d] \ar[r] &
H^q_\etale(X, \mathcal{F}) \ar[d] \ar[r] &
H^q_\etale(X, \mathcal{F}') \ar[d] \\
H^{q - 1}_\etale(Z, \mathcal{F}'|_Z) \ar[r] &
H^{q - 1}_\etale(Z, \mathcal{Q}|_Z) \ar[r] &
H^q_\etale(Z, \mathcal{F}|_Z) \ar[r] &
H^q_\etale(Z, \mathcal{F}'|_Z)
}
$$
行が完全であるアーベル群のこの可換図式を用いて証明を終える。

\medskip\noindent
$\mathcal{F}$ に対する単射性を示す。$\xi$ を
$H^q_\etale(X, \mathcal{F})$ の非零元とする。補題
\ref{etale-cohomology-lemma-efface-cohomology-on-closed-by-finite-cover} を $Z = X$（！）
として適用すれば、$\mathcal{F} \subset \mathcal{F}'$ であって、
$\xi$ が右で零に写るものを取れる。このとき $\xi$ は
$H^{q - 1}_\etale(X, \mathcal{Q})$ の元の像であり、$q - 1$ に対する
全単射性から、$\xi$ は $H^q_\etale(Z, \mathcal{F}|_Z)$ で零に写らない。

\medskip\noindent
$\mathcal{F}$ に対する全射性を示す。$\xi$ を
$H^q_\etale(Z, \mathcal{F}|_Z)$ の元とする。補題
\ref{etale-cohomology-lemma-efface-cohomology-on-closed-by-finite-cover} を $Z = Z$ として
適用すれば、$\mathcal{F} \subset \mathcal{F}'$ であって、$\xi$ が
右で零に写るものを取れる。このとき $\xi$ は
$H^{q - 1}_\etale(Z, \mathcal{Q}|_Z)$ の元の像であり、$q - 1$ に対する
全単射性から、$\xi$ は垂直写像の像に属する。
\end{proof}

\begin{lemma}
\label{etale-cohomology-lemma-vanishing-restriction-injective}
$X$ を対角射がアフィンなスキームで、$n + 1$ 個のアフィン開集合で
被覆できるものとする。$Z \subset X$ を閉部分スキームとする。
$\mathcal{A}$ を $X_\etale$ 上の捩れ環の層とし、
$\mathcal{I}$ を $\mathcal{A}$ 加群の単射的層（$X_\etale$ 上）とする。
このとき $H^q_\etale(Z, \mathcal{I}|_Z) = 0$ である（$q > n$）。
\end{lemma}

\begin{proof}
$n$ に関する帰納法で示す。$n = 0$ ならば $X$ はアフィンである。
$X = \Spec(A)$ および $Z = \Spec(A/I)$ と書く。$A^h$ を、エタール
$A$ 代数 $B$ であって $A/I \to B/IB$ が同型となるものの
フィルター余極限とする。
このとき $(A^h, IA^h)$ はヘンゼル対であり、$A/I = A^h/IA^h$ である。
More on Algebra, 補題 \ref{more-algebra-lemma-henselization} とその証明を
参照されたい。$X^h = \Spec(A^h)$ とおく。定理 \ref{etale-cohomology-theorem-gabber} により
$$
H^q_\etale(Z, \mathcal{I}|_Z) = H^q_\etale(X^h, \mathcal{I}|_{X^h})
$$
である。定理 \ref{etale-cohomology-theorem-colimit} により
$$
H^q_\etale(X^h, \mathcal{I}|_{X^h}) =
\colim_{A \to B} H^q_\etale(\Spec(B), \mathcal{I}|_{\Spec(B)})
$$
である。ここで余極限は上のような $A$ 代数 $B$ にわたる。
射 $\Spec(B) \to \Spec(A)$ はエタールなので、制限
$\mathcal{I}|_{\Spec(B)}$ は $\mathcal{A}|_{\Spec(B)}$ 加群の
単射的層である（Cohomology on Sites, 補題
\ref{sites-cohomology-lemma-cohomology-of-open}）。したがって右辺の
コホモロジー群は零であり、この場合の結論を得る。

\medskip\noindent
帰納段階には Mayer--Vietoris を用いる。すなわち $X = U \cup V$ とし、
$U$ は $n$ 個のアフィン開集合の和、$V$ はアフィンであるとする。
$X$ の対角射がアフィンであることから、$U \cap V$ は $n$ 個の
アフィン開集合の和である。Mayer--Vietoris は完全系列
$$
H^{q - 1}_\etale(U \cap V \cap Z, \mathcal{I}|_Z) \to
H^q_\etale(Z, \mathcal{I}|_Z) \to
H^q_\etale(U \cap Z, \mathcal{I}|_Z) \oplus
H^q_\etale(V \cap Z, \mathcal{I}|_Z)
$$
を与える。帰納法の仮定により $q > n$ で所望の消失を得る。
\end{proof}





\section{曲線上の捩れ層のコホモロジー}
\label{etale-cohomology-section-vanishing-torsion}

\noindent
本節の目的は、曲線上の捩れ層のコホモロジーに関する基本的な有限性および
消滅の結果を証明することである。定理
\ref{etale-cohomology-theorem-vanishing-affine-curves} を参照されたい。
節 \ref{etale-cohomology-section-vanishing-torsion-coefficients} では、
ネーター環上の捩れ加群の構成可能層を論じる。

\begin{situation}
\label{etale-cohomology-situation-what-to-prove}
ここで $k$ は代数閉体、$X$ は次元 $\leq 1$ の $k$ 上の分離有限型
スキーム、そして $\mathcal{F}$ は $X_\etale$ 上の捩れアーベル層である。
\end{situation}

\noindent
状況 \ref{etale-cohomology-situation-what-to-prove} において、次の諸命題を証明したい。
\begin{enumerate}
\item
\label{etale-cohomology-item-vanishing}
$H^q_\etale(X, \mathcal{F}) = 0$ が $q > 2$ に対して成り立つ。
\item
\label{etale-cohomology-item-vanishing-affine}
$H^q_\etale(X, \mathcal{F}) = 0$ が $q > 1$ に対して成り立つ。ただし
$X$ はアフィンであるとする。
\item
\label{etale-cohomology-item-vanishing-p-p}
$H^q_\etale(X, \mathcal{F}) = 0$ が $q > 1$ に対して成り立つ。ただし
$p = \text{char}(k) > 0$ かつ $\mathcal{F}$ は $p$ 冪捩れであるとする。
\item
\label{etale-cohomology-item-finite-prime-to-p}
$H^q_\etale(X, \mathcal{F})$ は有限である。ただし $\mathcal{F}$ は
構成可能で、$\text{char}(k)$ と互いに素な捩れ層であるとする。
\item
\label{etale-cohomology-item-finite-proper}
$H^q_\etale(X, \mathcal{F})$ は有限である。ただし $X$ は固有で
$\mathcal{F}$ は構成可能であるとする。
\item
\label{etale-cohomology-item-base-change-prime-to-p}
$H^q_\etale(X, \mathcal{F}) \to
H^q_\etale(X_{k'}, \mathcal{F}|_{X_{k'}})$ は、代数閉体の任意の拡大
$k'/k$ に対して同型である。ただし $\mathcal{F}$ は
$\text{char}(k)$ と互いに素な捩れ層であるとする。
\item
\label{etale-cohomology-item-base-change-proper}
$H^q_\etale(X, \mathcal{F}) \to
H^q_\etale(X_{k'}, \mathcal{F}|_{X_{k'}})$ は、代数閉体の任意の拡大
$k'/k$ に対して同型である。ただし $X$ は固有であるとする。
\item
\label{etale-cohomology-item-surjective}
$H^2_\etale(X, \mathcal{F}) \to H^2_\etale(U, \mathcal{F})$
は、任意の開部分 $U \subset X$ に対して全射である。
\end{enumerate}
任意の状況 \ref{etale-cohomology-situation-what-to-prove} が与えられたとき、
その状況に適用される上記の命題がすべて真であることを
「命題 (\ref{etale-cohomology-item-vanishing})--(\ref{etale-cohomology-item-surjective}) が成り立つ」
ということにする。証明は、定数係数コホモロジーの計算から従う
次の帰結から始める。

\begin{lemma}
\label{etale-cohomology-lemma-constant-smooth-statements}
状況 \ref{etale-cohomology-situation-what-to-prove} において、$X$ が滑らかで
$\mathcal{F} = \underline{\mathbf{Z}/\ell\mathbf{Z}}$ であると仮定する。
ここで $\ell$ はある素数である。このとき $\mathcal{F}$ に対して
命題 (\ref{etale-cohomology-item-vanishing})--(\ref{etale-cohomology-item-surjective}) が成り立つ。
\end{lemma}

\begin{proof}
$X$ は滑らかなので、$X$ は滑らかな曲線と点との有限非交和である。
点（$k$ のスペクトル）上の任意の層の高次エタール・コホモロジーは
自明である。例えば補題 \ref{etale-cohomology-lemma-compare-cohomology-point}
または \ref{etale-cohomology-lemma-affine-only-closed-points} による。
したがって $X$ は滑らかな曲線であると仮定してよい。

\medskip\noindent
場合 I：$\ell$ が $k$ の標数と異なる場合。
これは補題 \ref{etale-cohomology-lemma-cohomology-smooth-projective-curve}
（射影的な場合）および補題
\ref{etale-cohomology-lemma-vanishing-cohomology-mu-smooth-curve}
（アフィンの場合）から従う。コホモロジーと代数閉基礎体の拡大に関する
命題 (\ref{etale-cohomology-item-base-change-prime-to-p}) は、種数 $g$ と「欠点」の個数
$r$ が $k$ から $k'$ へ移っても変化しないことから従う。
命題 (\ref{etale-cohomology-item-surjective}) については、$H^2_\etale(U, \mathcal{F})$
が $U \not = X$ となるや否や零になる。実際、このとき $U$ はアフィンである
（Varieties、補題 \ref{varieties-lemma-proper-minus-point} および
\ref{varieties-lemma-curve-affine-projective}）。

\medskip\noindent
場合 II：$\ell$ が $k$ の標数に等しい場合。
消滅は補題 \ref{etale-cohomology-lemma-vanishing-variety-char-p-p} による。
命題 (\ref{etale-cohomology-item-finite-proper}) および
(\ref{etale-cohomology-item-base-change-proper}) は補題
\ref{etale-cohomology-lemma-finiteness-proper-variety-char-p-p} から従う。
\end{proof}

\begin{remark}[代数閉基礎体の拡大に対する不変性]
\label{etale-cohomology-remark-invariance}
$k$ を標数 $p > 0$ の代数閉体とする。
節 \ref{etale-cohomology-section-artin-schreier} で、完全列
$$
k[x] \to k[x] \to
H^1_\etale(\mathbf{A}^1_k, \mathbf{Z}/p\mathbf{Z}) \to 0
$$
が存在することを見た。ここで最初の矢印は $f(x)$ を $f^p - f$ に写す。
余核の代表元の集合は、多項式
$$
\sum\nolimits_{p \not | n} \lambda_n x^n
$$
（ただし $\lambda_n \in k$）からなる。（$k$ が代数閉でなければ、
さらにいくつかの定数も加えなければならない。）特に $k'/k$ が
代数閉体の拡大であるとき、写像
$$
H^1_\etale(\mathbf{A}^1_k, \mathbf{Z}/p\mathbf{Z})
\to
H^1_\etale(\mathbf{A}^1_{k'}, \mathbf{Z}/p\mathbf{Z})
$$
は一般には同型でない。特に写像
$\pi_1(\mathbf{A}^1_{k'}) \to \pi_1(\mathbf{A}^1_k)$
も、エタール基本群の間で（ここに将来の参照を挿入する）同型ではない。
したがって、アフィン直線のエタール・ホモトピー型は代数閉基礎体に依存する。
上の補題 \ref{etale-cohomology-lemma-constant-smooth-statements} から、これは標数 $p$ で
$p$ 冪捩れ係数を用いた場合にのみ起こる現象であることが分かる。
\end{remark}

\begin{lemma}
\label{etale-cohomology-lemma-ses-statements}
$k$ を代数閉体とする。$X$ を $k$ 上の次元 $\leq 1$ の分離有限型
スキームとする。
$0 \to \mathcal{F}_1 \to \mathcal{F} \to \mathcal{F}_2 \to 0$
を $X$ 上の捩れアーベル層の短完全列とする。
$\mathcal{F}_1$ と $\mathcal{F}_2$ に対して命題
(\ref{etale-cohomology-item-vanishing})--(\ref{etale-cohomology-item-surjective}) が成り立つならば、
$\mathcal{F}$ に対しても成り立つ。
\end{lemma}

\begin{proof}
これは大部分、定義とコホモロジー長完全列から直ちに従う。また
$\mathcal{F}$ が構成可能（それぞれ $k$ の標数と互いに素な捩れ）
であることと、$\mathcal{F}_1$ と $\mathcal{F}_2$ がともに構成可能
（それぞれ $k$ の標数と互いに素な捩れ）であることとは同値である。
命題 \ref{etale-cohomology-proposition-constructible-over-noetherian} を参照されたい。
若干の詳細は省略する。
\end{proof}

\begin{lemma}
\label{etale-cohomology-lemma-finite-pushforward-statements}
$k$ を代数閉体とする。$f : X \to Y$ を、$k$ 上の次元 $\leq 1$ の
分離有限型スキームの間の有限射とする。$\mathcal{F}$ を $X$ 上の
捩れアーベル層とする。$\mathcal{F}$ に対して命題
(\ref{etale-cohomology-item-vanishing})--(\ref{etale-cohomology-item-surjective}) が成り立つならば、
$f_*\mathcal{F}$ に対しても成り立つ。
\end{lemma}

\begin{proof}
実際、$H^q_\etale(X, \mathcal{F}) = H^q_\etale(Y, f_*\mathcal{F})$ である。
これは $R^qf_*$ が $q > 0$ に対して消滅すること
（命題 \ref{etale-cohomology-proposition-finite-higher-direct-image-zero}）と
ルレイ・スペクトル系列（Cohomology on Sites、補題
\ref{sites-cohomology-lemma-apply-Leray}）による。
(\ref{etale-cohomology-item-surjective}) には、$f_*$ の形成が任意の基底変換と可換であること
（補題 \ref{etale-cohomology-lemma-finite-pushforward-commutes-with-base-change}）を用いる。
\end{proof}

\begin{lemma}
\label{etale-cohomology-lemma-restrict-to-open}
状況 \ref{etale-cohomology-situation-what-to-prove} において $\mathcal{F}$ が構成可能であると
仮定する。$j : X' \to X$ を稠密な開部分スキームの包含とする。
このとき $\mathcal{F}$ に対して命題
(\ref{etale-cohomology-item-vanishing})--(\ref{etale-cohomology-item-surjective}) が成り立つことと、
$j_!j^{-1}\mathcal{F}$ に対して成り立つこととは同値である。
\end{lemma}

\begin{proof}
$X'$ は稠密なので、$Z = X \setminus X'$ は次元 $0$ であり、
したがって有限集合 $Z = \{x_1, \ldots, x_n\}$ である。その各点は
$k$ 有理点である。補題
\ref{etale-cohomology-lemma-ses-associated-to-open} の短完全列
$$
0 \to j_!j^{-1}\mathcal{F} \to \mathcal{F} \to i_*i^{-1}\mathcal{F} \to 0
$$
を考える。$H^q_\etale(X, i_*i^{-1}\mathcal{F}) =
H^q_\etale(Z, i^*\mathcal{F})$ であることに注意する。実際、
$i : Z \to X$ は閉埋め込みなので有限であり、したがって命題
\ref{etale-cohomology-proposition-finite-higher-direct-image-zero} により $R^qi_*$ は
$q > 0$ に対して消滅する。ゆえにルレイ・スペクトル系列
（Cohomology on Sites、補題 \ref{sites-cohomology-lemma-apply-Leray}）
からこの等式が従う。$Z$ は代数閉体のスペクトルの非交和なので、
$H^q_\etale(Z, i^*\mathcal{F}) = 0$ が $q > 0$ に対して成り立ち、かつ
$$
H^0_\etale(Z, i^{-1}\mathcal{F}) =
\bigoplus\nolimits_{i = 1, \ldots, n} \mathcal{F}_{x_i}
$$
である。各 $\mathcal{F}_{x_i}$ は、$\mathcal{F}$ が構成可能であるという
仮定により有限だから、この群も有限である。
コホモロジー長完全列は完全列
$$
0 \to
H^0_\etale(X, j_!j^{-1}\mathcal{F}) \to
H^0_\etale(X, \mathcal{F}) \to
H^0_\etale(Z, i^{-1}\mathcal{F}) \to
H^1_\etale(X, j_!j^{-1}\mathcal{F}) \to
H^1_\etale(X, \mathcal{F}) \to 0
$$
および同型 $H^q_\etale(X, j_!j^{-1}\mathcal{F}) \to
H^q_\etale(X, \mathcal{F})$（$q > 1$）を与える。

\medskip\noindent
ここから、各命題 (\ref{etale-cohomology-item-vanishing})--(\ref{etale-cohomology-item-surjective}) が
$\mathcal{F}$ に対して成り立つことと $j_!j^{-1}\mathcal{F}$ に対して
成り立つこととの同値性は容易に従う。読者の便宜のため若干補足する：
(a) $\mathcal{F}$ が $k$ の標数と互いに素な捩れ層ならば、
$j_!j^{-1}\mathcal{F}$ もそうである；(b) 層
$j_!j^{-1}\mathcal{F}$ は構成可能である；(c)
$H^0_\etale(Z, i^{-1}\mathcal{F}) =
H^0_\etale(Z_{k'}, i^{-1}\mathcal{F}|_{Z_{k'}})$ である；そして
(d) $U \subset X$ が開ならば $U' = U \cap X'$ は $U$ で稠密である。
\end{proof}

\begin{lemma}
\label{etale-cohomology-lemma-even-easier}
状況 \ref{etale-cohomology-situation-what-to-prove} において $X$ が滑らかであると仮定する。
$j : U \to X$ を開埋め込み、$\ell$ を素数とし、
$\mathcal{F} = j_!\underline{\mathbf{Z}/\ell\mathbf{Z}}$ とおく。
このとき $\mathcal{F}$ に対して命題
(\ref{etale-cohomology-item-vanishing})--(\ref{etale-cohomology-item-surjective}) が成り立つ。
\end{lemma}

\begin{proof}
$X$ は滑らかなので滑らかな曲線の非交和であり、したがって $X$ は曲線
（すなわち既約）であると仮定してよい。すると $U = \emptyset$ であって
証明すべきことがないか、または $U \subset X$ が稠密である。後者の場合、
補題 \ref{etale-cohomology-lemma-constant-smooth-statements} および
\ref{etale-cohomology-lemma-restrict-to-open} から従う。
\end{proof}

\begin{lemma}
\label{etale-cohomology-lemma-somewhat-easier}
状況 \ref{etale-cohomology-situation-what-to-prove} において $X$ が被約であると仮定する。
$j : U \to X$ を開埋め込み、$\ell$ を素数とし、
$\mathcal{F} = j_!\underline{\mathbf{Z}/\ell\mathbf{Z}}$ とおく。
このとき $\mathcal{F}$ に対して命題
(\ref{etale-cohomology-item-vanishing})--(\ref{etale-cohomology-item-surjective}) が成り立つ。
\end{lemma}

\begin{proof}
補題 \ref{etale-cohomology-lemma-even-easier} との違いは、ここでは $X$ の滑らかさを仮定
しない点にある。$\nu : X^\nu \to X$ を正規化射とする。このとき $\nu$ は
有限であり（Varieties、補題
\ref{varieties-lemma-normalization-locally-algebraic}）、$X^\nu$ は滑らかである
（Varieties、補題 \ref{varieties-lemma-regular-point-on-curve}）。
$j^\nu : U^\nu \to X^\nu$ を $U$ の逆像とする。補題
\ref{etale-cohomology-lemma-even-easier} により
$j^\nu_!\underline{\mathbf{Z}/\ell\mathbf{Z}}$ に対して結果が成り立つ。
補題 \ref{etale-cohomology-lemma-finite-pushforward-statements} により
$\nu_*j^\nu_!\underline{\mathbf{Z}/\ell\mathbf{Z}}$ に対しても結果が成り立つ。
一般には $\nu_*j^\nu_!\underline{\mathbf{Z}/\ell\mathbf{Z}}$ が
$j_!\underline{\mathbf{Z}/\ell\mathbf{Z}}$ に等しいとは限らないが、
次のようにこれを回避できる。$X$ は被約なので、射
$\nu : X^\nu \to X$ はある稠密開集合 $j' : X' \to X$ 上で同型である
（Varieties、補題 \ref{varieties-lemma-normalization-locally-algebraic}）。
この開集合上では
$$
(j')^{-1}(\nu_*j^\nu_!\underline{\mathbf{Z}/\ell\mathbf{Z}})
=
(j')^{-1}(j_!\underline{\mathbf{Z}/\ell\mathbf{Z}})
$$
と一致する。上の層と $j' : X' \to X$ に対して補題
\ref{etale-cohomology-lemma-restrict-to-open} を二度適用すれば結論を得る。
\end{proof}

\begin{lemma}
\label{etale-cohomology-lemma-vanishing-easier}
状況 \ref{etale-cohomology-situation-what-to-prove} において $X$ が被約であると仮定する。
$j : U \to X$ を開埋め込みとし、$U$ は連結であるとする。$\ell$ を
素数とする。$\mathcal{G}$ を $\mathbf{F}_\ell$ ベクトル空間からなる
有限局所定数層とし、$U$ 上で考える。$\mathcal{F} = j_!\mathcal{G}$
とおく。このとき $\mathcal{F}$ に対して命題
(\ref{etale-cohomology-item-vanishing})--(\ref{etale-cohomology-item-surjective}) が成り立つ。
\end{lemma}

\begin{proof}
補題 \ref{etale-cohomology-lemma-pullback-filtered} にあるように、$f : V \to U$ を
$\ell$ と互いに素な次数の有限エタール射とする。節
\ref{etale-cohomology-section-trace-method} の議論は写像
$$
\mathcal{G} \to f_*f^{-1}\mathcal{G} \to \mathcal{G}
$$
を与え、その合成は同型である。したがって
$\mathcal{F} = j_!f_*f^{-1}\mathcal{G}$ の場合に補題を証明すれば十分である。
ザリスキーの主定理（More on Morphisms、補題
\ref{more-morphisms-lemma-quasi-finite-separated-pass-through-finite})
により、可換図式
$$
\xymatrix{
V \ar[r]_{j'} \ar[d]_f & Y \ar[d]^{\overline{f}} \\
U \ar[r]^j & X
}
$$
を選べる。ここで $\overline{f} : Y \to X$ は有限であり、$j'$ は稠密像を
もつ開埋め込みである。$Y$ をその被約化で置き換えてよい（この置換で
$V$ は変わらない。実際 $V$ は $U$ 上エタールなので被約である）。
$f$ は有限で $V$ は $Y$ で稠密だから $V = U \times_X Y$ である。
補題 \ref{etale-cohomology-lemma-compatible-shriek-push-finite} により
$$
j_!f_*f^{-1}\mathcal{G} = \overline{f}_*j'_!f^{-1}\mathcal{G}
$$
である。補題 \ref{etale-cohomology-lemma-finite-pushforward-statements} により
$j'_!f^{-1}\mathcal{G}$ を考えれば十分である。
補題 \ref{etale-cohomology-lemma-pullback-filtered} が与えるフィルトレーションの存在、
$j'_!$ が完全であること、および補題 \ref{etale-cohomology-lemma-ses-statements} により、
$\mathcal{F} = j'_!\underline{\mathbf{Z}/\ell\mathbf{Z}}$ の場合へ帰着する。
これは補題 \ref{etale-cohomology-lemma-somewhat-easier} である。
\end{proof}


%10.20.09

\begin{theorem}
\label{etale-cohomology-theorem-vanishing-affine-curves}
$k$ を代数閉体、$X$ を次元 $\leq 1$ の $k$ 上の分離有限型スキーム、
$\mathcal{F}$ を $X_\etale$ 上の捩れアーベル層とする。このとき
\begin{enumerate}
\item
$H^q_\etale(X, \mathcal{F}) = 0$ が $q > 2$ に対して成り立つ。
\item
$H^q_\etale(X, \mathcal{F}) = 0$ が $q > 1$ に対して成り立つ。ただし
$X$ はアフィンであるとする。
\item
$H^q_\etale(X, \mathcal{F}) = 0$ が $q > 1$ に対して成り立つ。ただし
$p = \text{char}(k) > 0$ かつ $\mathcal{F}$ は $p$ 冪捩れであるとする。
\item
$H^q_\etale(X, \mathcal{F})$ は有限である。ただし $\mathcal{F}$ は
構成可能で、$\text{char}(k)$ と互いに素な捩れ層であるとする。
\item
$H^q_\etale(X, \mathcal{F})$ は有限である。ただし $X$ は固有で
$\mathcal{F}$ は構成可能であるとする。
\item
$H^q_\etale(X, \mathcal{F}) \to
H^q_\etale(X_{k'}, \mathcal{F}|_{X_{k'}})$ は、代数閉体の任意の拡大
$k'/k$ に対して同型である。ただし $\mathcal{F}$ は
$\text{char}(k)$ と互いに素な捩れ層であるとする。
\item
$H^q_\etale(X, \mathcal{F}) \to
H^q_\etale(X_{k'}, \mathcal{F}|_{X_{k'}})$ は、代数閉体の任意の拡大
$k'/k$ に対して同型である。ただし $X$ は固有であるとする。
\item
$H^2_\etale(X, \mathcal{F}) \to H^2_\etale(U, \mathcal{F})$
は、任意の開部分 $U \subset X$ に対して全射である。
\end{enumerate}
\end{theorem}

\begin{proof}
定理の主張は、状況 \ref{etale-cohomology-situation-what-to-prove} において命題
(\ref{etale-cohomology-item-vanishing})--(\ref{etale-cohomology-item-surjective}) が成り立つということである。
まず $X$ をその被約化で置き換える。これは命題
\ref{etale-cohomology-proposition-topological-invariance} により許される。
補題 \ref{etale-cohomology-lemma-torsion-colimit-constructible} により、$\mathcal{F}$ を
構成可能アーベル層のフィルター余極限として書ける。コホモロジーを取る
操作は余極限と可換である。補題 \ref{etale-cohomology-lemma-colimit} を参照されたい。
さらに $X_{k'} \to X$ による引き戻しは左随伴なので余極限と可換である。
したがって構成可能層について諸命題を証明すれば十分である。

\medskip\noindent
この段落では補題 \ref{etale-cohomology-lemma-ses-statements} を以後断らずに用いる。
$\mathcal{F} = \mathcal{F}_1 \oplus \ldots \oplus \mathcal{F}_r$
と書く。ここで $\mathcal{F}_i$ はある素数 $\ell_i$ に対して
$\ell_i$ 準素である。$\ell^n$ が $\mathcal{F}$ を零化すると仮定してよい。
ここで $\ell$ はある素数である。完全列
$$
0 \to \mathcal{F}[\ell] \to \mathcal{F} \to \mathcal{F}/\mathcal{F}[\ell] \to 0.
$$
を考える。これより $\mathcal{F}$ が $\ell$ 捩れであると仮定すれば十分と
分かる。すなわち $\mathcal{F}$ は $\mathbf{F}_\ell$ ベクトル空間の
構成可能層である。ここで $\ell$ はある素数である。

\medskip\noindent
定義により、ある稠密開集合 $U \subset X$ が存在して、$\mathcal{F}|_U$ は
$\mathbf{F}_\ell$ ベクトル空間の有限局所定数層である。
$\dim(X) \leq 1$ なので、$U$ を縮小した後、
$U = U_1 \amalg \ldots \amalg U_n$ が既約スキームの非交和であると仮定できる
（$\geq 2$ 個の成分の交わりに含まれる $U$ の閉点を除くだけでよい）。
補題 \ref{etale-cohomology-lemma-restrict-to-open} により
$\mathcal{F} = j_!\mathcal{G}$ の場合に帰着する。ここで $\mathcal{G}$ は
$\mathbf{F}_\ell$ ベクトル空間の有限局所定数層であり、$U$ 上で考える。

\medskip\noindent
$U = U_1 \amalg \ldots \amalg U_n$ を各 $U_i$ が既約となるように
選んだので、
$$
j_!\mathcal{G} =
j_{1!}(\mathcal{G}|_{U_1}) \oplus \ldots \oplus
j_{n!}(\mathcal{G}|_{U_n})
$$
である。ここで $j_i : U_i \to X$ は包含射である。
$j_{i!}(\mathcal{G}|_{U_i})$ の場合は補題
\ref{etale-cohomology-lemma-vanishing-easier} で扱われている。
\end{proof}

\begin{theorem}
\label{etale-cohomology-theorem-vanishing-curves}
$X$ を次元 $1$ の、代数閉体 $k$ 上の有限型スキームとする。
$\mathcal{F}$ を $X_\etale$ 上の捩れ層とする。このとき
$$
H_\etale^q(X, \mathcal{F}) = 0, \quad \forall q \geq 3.
$$
$X$ がアフィンならば、さらに $H_\etale^2(X, \mathcal{F}) = 0$ である。
\end{theorem}

\begin{proof}
$X$ が分離ならば、これはより精密な定理
\ref{etale-cohomology-theorem-vanishing-affine-curves} から直ちに従う。
$X$ が非分離ならば、アフィン開被覆
$X = X_1 \cup \ldots \cup X_n$ を選ぶ。$n$ に関する帰納法により、
$U = X_1 \cup \ldots \cup X_{n - 1}$ 上では消滅が成り立つと仮定してよい。
するとマイヤー–ヴィートリス（補題 \ref{etale-cohomology-lemma-mayer-vietoris}）は
$$
H^2_\etale(U, \mathcal{F}) \oplus H^2_\etale(X_n, \mathcal{F}) \to
H^2_\etale(U \cap X_n, \mathcal{F}) \to
H^3_\etale(X, \mathcal{F}) \to 0
$$
を与える。しかし $U \cap X_n$ はアフィン・スキームの開部分であり、
次元に関する仮定によりそれ自身もアフィンなので、群
$H^2_\etale(U \cap X_n, \mathcal{F})$ は定理
\ref{etale-cohomology-theorem-vanishing-affine-curves} により消滅する。
\end{proof}

\begin{lemma}
\label{etale-cohomology-lemma-base-change-dim-1-separably-closed}
$k'/k$ を分離閉体の拡大とする。$X$ を $k$ 上の次元 $\leq 1$ の
固有スキームとする。$\mathcal{F}$ を $X$ 上の捩れアーベル層とする。
このとき写像 $H^q_\etale(X, \mathcal{F}) \to
H^q_\etale(X_{k'}, \mathcal{F}|_{X_{k'}})$ は $q \geq 0$ に対して同型である。
\end{lemma}

\begin{proof}
代数閉体の場合は定理 \ref{etale-cohomology-theorem-vanishing-affine-curves} で既に見た。
補題の仮定を満たす $k \subset k'$ が与えられたとき、可換図式
$$
\xymatrix{
k' \ar[r] & \overline{k}' \\
k \ar[u] \ar[r] & \overline{k} \ar[u]
}
$$
を選べる。ここで $k \subset \overline{k}$ と
$k' \subset \overline{k}'$ は代数閉包である。$k$ と $k'$ は分離閉なので、
体拡大
$\overline{k}/k$ および $\overline{k}'/k'$
は代数的かつ純非分離である。この場合、射
$X_{\overline{k}} \to X$ および $X_{\overline{k}'} \to X_{k'}$
は普遍同相射である。したがって $\mathcal{F}$ のコホモロジーは
$X_{\overline{k}}$ 上で計算でき、$\mathcal{F}|_{X_{k'}}$ のコホモロジーは
$X_{\overline{k}'}$ 上で計算できる。命題
\ref{etale-cohomology-proposition-topological-invariance} を参照されたい。
ゆえに一般の場合は代数閉体の場合から従う。
\end{proof}






\section{曲線上の捩れ加群のコホモロジー}
\label{etale-cohomology-section-vanishing-torsion-coefficients}

\noindent
本節では、ネーター環上の捩れ加群の構成可能層について、節
\ref{etale-cohomology-section-vanishing-torsion} の議論を繰り返す。最も興味深い段階から
始める。

\begin{lemma}
\label{etale-cohomology-lemma-constant-statements-coefficients}
$\Lambda$ をネーター環、$M$ を有限 $\Lambda$-加群で整数 $n > 0$ により
零化されるもの、$k$ を代数閉体、$X$ を次元 $\leq 1$ の $k$ 上の
分離有限型スキームとする。このとき
\begin{enumerate}
\item $H^q_\etale(X, \underline{M})$ は有限 $\Lambda$-加群である。ただし
$n$ は $\text{char}(k)$ と互いに素であるとする。
\item $H^q_\etale(X, \underline{M})$ は有限 $\Lambda$-加群である。ただし
$X$ は固有であるとする。
\end{enumerate}
\end{lemma}

\begin{proof}
$n = \ell n'$ と書け、$\ell$ が素数ならば、短完全列
$0 \to M[\ell] \to M \to M' \to 0$ を得る。この列は有限
$\Lambda$-加群からなり、$M'$ は $n'$ で零化される。
これから対応する定数層の短完全列が得られ、さらにコホモロジー加群の完全列
$$
H^q_\etale(X, \underline{M[n]}) \to
H^q_\etale(X, \underline{M}) \to
H^q_\etale(X, \underline{M'})
$$
が生じる。したがって $M$ が素数で零化される場合に結果を示せれば、
$n$ に関する帰納法により結論を得る。

\medskip\noindent
$\ell$ を $\ell$ が $M$ を零化するような素数とする。このとき
$\Lambda$ を $\mathbf{F}_\ell$-代数 $\Lambda/\ell \Lambda$ で置き換えられる。
実際、$\mathcal{F}$ を $\Lambda$-加群の層として取ったコホモロジーは、
$\mathcal{F}$ を $\Lambda/\ell \Lambda$-加群の層として取った
コホモロジーと同じである。例えば Cohomology on Sites、補題
\ref{sites-cohomology-lemma-cohomology-modules-abelian-agree}.

\medskip\noindent
$\ell$ を素数とし、$\ell$ が $M$ と $\Lambda$ を零化すると仮定する。
$M$ が有限自由 $\Lambda$-加群である場合に帰着しよう。そのため短完全列
$$
0 \to N \to \Lambda^{\oplus m} \to M \to 0
$$
を選ぶ。これは完全列
$$
H^q_\etale(X, \underline{\Lambda^{\oplus m}}) \to
H^q_\etale(X, \underline{M}) \to
H^{q + 1}_\etale(X, \underline{N})
$$
を定める。$q$ に関する降下帰納法により、$M$ に対する結果は
$\Lambda^{\oplus m}$ に対する結果から従う。ここで、定理
\ref{etale-cohomology-theorem-vanishing-affine-curves} によりコホモロジー群が次数 $> 2$
で消滅することを用いる。

\medskip\noindent
$\ell$ を素数とし、$\ell$ が $\Lambda$ を零化すると仮定する。
残るのは、コホモロジー群 $H^q_\etale(X, \underline{\Lambda})$ が
有限 $\Lambda$-加群であることを示すことである。これを示すため一つの
技巧を用いる。「正しい」議論は、後で証明する係数定理を用いる。
$\Lambda = \bigoplus_{i \in I} \mathbf{F}_\ell e_i$ という基底を選び、
$e_0 = 1$ となるような $0 \in I$ を取る。
この基底の選択は層の同型
$$
\underline{\Lambda} = \bigoplus \underline{\mathbf{F}_\ell} e_i
$$
を $X_\etale$ 上で定める。したがって
$$
H^q_\etale(X, \underline{\Lambda}) =
H^q_\etale(X, \bigoplus \underline{\mathbf{F}_\ell} e_i) =
\bigoplus H^q_\etale(X, \underline{\mathbf{F}_\ell})e_i
$$
である。実際、定理 \ref{etale-cohomology-theorem-colimit}（または補題
\ref{etale-cohomology-lemma-colimit}、または補題
\ref{etale-cohomology-lemma-direct-sum-bounded-below-cohomology}）により、$X$ 上の
コホモロジーを取る操作は直和と可換である。既に
$H^q_\etale(X, \underline{\mathbf{F}_\ell})$ が有限次元の
$\mathbf{F}_\ell$-ベクトル空間であることを知っている
（定理 \ref{etale-cohomology-theorem-vanishing-affine-curves}）。ゆえに
$H^q_\etale(X, \underline{\Lambda})$ は同じ階数の自由
$\Lambda$-加群である。実際、
基底 $\xi_1, \ldots, \xi_m$ が
$H^q_\etale(X, \underline{\mathbf{F}_\ell})$ に対して与えられれば、
$\xi_1 e_0, \ldots, \xi_m e_0$ は $\Lambda$-基底を
$H^q_\etale(X, \underline{\Lambda})$ に対してなす。
\end{proof}

\begin{lemma}
\label{etale-cohomology-lemma-finite-pushforward-coefficients}
$\Lambda$ をネーター環、$k$ を代数閉体とする。$f : X \to Y$ を
$k$ 上の次元 $\leq 1$ の分離有限型スキームの間の有限射とし、
$\mathcal{F}$ を $\Lambda$-加群の層として $X_\etale$ 上で考える。
$H^q_\etale(X, \mathcal{F})$ が有限 $\Lambda$-加群ならば、
$H^q_\etale(Y, f_*\mathcal{F})$ もそうである。
\end{lemma}

\begin{proof}
実際、$H^q_\etale(X, \mathcal{F}) = H^q_\etale(Y, f_*\mathcal{F})$ である。
これは $R^qf_*$ が $q > 0$ に対して消滅すること
（命題 \ref{etale-cohomology-proposition-finite-higher-direct-image-zero}）と
ルレイ・スペクトル系列（Cohomology on Sites、補題
\ref{sites-cohomology-lemma-apply-Leray}）による。
\end{proof}

\begin{lemma}
\label{etale-cohomology-lemma-restrict-to-open-coefficients}
$\Lambda$ をネーター環、$k$ を代数閉体、$X$ を $k$ 上の次元
$\leq 1$ の分離有限型スキームとする。$\mathcal{F}$ を構成可能な
$\Lambda$-加群の層とし、$X_\etale$ 上で考える。$j : X' \to X$ を稠密な開部分
スキームの包含とする。このとき $H^q_\etale(X, \mathcal{F})$ が有限
$\Lambda$-加群であることと、$H^q_\etale(X, j_!j^{-1}\mathcal{F})$ が有限
$\Lambda$-加群であることとは同値である。
\end{lemma}

\begin{proof}
$X'$ は稠密なので、$Z = X \setminus X'$ は次元 $0$ であり、
したがって有限集合 $Z = \{x_1, \ldots, x_n\}$ である。その各点は
$k$ 有理点である。補題 \ref{etale-cohomology-lemma-ses-associated-to-open} の短完全列
$$
0 \to j_!j^{-1}\mathcal{F} \to \mathcal{F} \to i_*i^{-1}\mathcal{F} \to 0
$$
を考える。$H^q_\etale(X, i_*i^{-1}\mathcal{F}) =
H^q_\etale(Z, i^*\mathcal{F})$ であることに注意する。実際、
$i : Z \to X$ は閉埋め込みなので有限であり、したがって命題
\ref{etale-cohomology-proposition-finite-higher-direct-image-zero} により $R^qi_*$ は
$q > 0$ に対して消滅する。ゆえにルレイ・スペクトル系列
（Cohomology on Sites、補題 \ref{sites-cohomology-lemma-apply-Leray}）
からこの等式が従う。$Z$ は代数閉体のスペクトルの非交和なので、
$H^q_\etale(Z, i^*\mathcal{F}) = 0$ が $q > 0$ に対して成り立ち、かつ
$$
H^0_\etale(Z, i^{-1}\mathcal{F}) =
\bigoplus\nolimits_{i = 1, \ldots, n} \mathcal{F}_{x_i}
$$
は有限 $\Lambda$-加群である。$\mathcal{F}_{x_i}$ が有限であるのは、
$\mathcal{F}$ が構成可能な $\Lambda$-加群の層であるという仮定による。
コホモロジー長完全列は完全列
$$
0 \to
H^0_\etale(X, j_!j^{-1}\mathcal{F}) \to
H^0_\etale(X, \mathcal{F}) \to
H^0_\etale(Z, i^{-1}\mathcal{F}) \to
H^1_\etale(X, j_!j^{-1}\mathcal{F}) \to
H^1_\etale(X, \mathcal{F}) \to 0
$$
および同型 $H^0_\etale(X, j_!j^{-1}\mathcal{F}) \to
H^0_\etale(X, \mathcal{F})$（$q > 1$）を与える。これから補題は容易に従う。
\end{proof}

\begin{lemma}
\label{etale-cohomology-lemma-somewhat-easier-coefficients}
$\Lambda$ をネーター環、$M$ を有限 $\Lambda$-加群で整数 $n > 0$ により
零化されるもの、$k$ を代数閉体、$X$ を次元 $\leq 1$ の $k$ 上の
分離有限型スキームとし、$j : U \to X$ を開埋め込みとする。このとき
\begin{enumerate}
\item $H^q_\etale(X, j_!\underline{M})$ は有限 $\Lambda$-加群である。ただし
$n$ は $\text{char}(k)$ と互いに素であるとする。
\item $H^q_\etale(X, j_!\underline{M})$ は有限 $\Lambda$-加群である。ただし
$X$ は固有であるとする。
\end{enumerate}
\end{lemma}

\begin{proof}
$\dim(X) \leq 1$ なので、開集合 $V \subset X$ であって $U$ と交わらず、
$X' = U \cup V$ が $X$ の稠密開集合となるものが存在する（詳細は省略する）。
$j' : X' \to X$ を包含射とすれば、$j_!\underline{M}$ は
$j'_!\underline{M}$ の直和因子である。したがって $U$ が $X$ の稠密開集合
である場合に補題を証明すれば十分である。この場合は補題
\ref{etale-cohomology-lemma-restrict-to-open-coefficients} および
\ref{etale-cohomology-lemma-constant-statements-coefficients} から従う。
\end{proof}

\begin{lemma}
\label{etale-cohomology-lemma-ses-statements-coefficients}
$\Lambda$ をネーター環、$k$ を代数閉体、$X$ を $k$ 上の次元
$\leq 1$ の分離有限型スキームとする。また
$0 \to \mathcal{F}_1 \to \mathcal{F} \to \mathcal{F}_2 \to 0$ を
$\Lambda$-加群の層の短完全列とし、$X_\etale$ 上で考える。
$H^q_\etale(X, \mathcal{F}_i)$（$i = 1, 2$）が有限
$\Lambda$-加群ならば、$H^q_\etale(X, \mathcal{F})$ も有限
$\Lambda$-加群である。
\end{lemma}

\begin{proof}
コホモロジー長完全列から直ちに従う。
\end{proof}

\begin{lemma}
\label{etale-cohomology-lemma-vanishing-easier-coefficients}
$\Lambda$ をネーター環、$k$ を代数閉体、$X$ を次元 $\leq 1$ の
$k$ 上の分離有限型スキームとする。$j : U \to X$ を開埋め込みとし、
$U$ は連結であるとする。$\ell$ を素数、$n > 0$ とし、$\mathcal{G}$ を
有限型かつ局所定数な $\Lambda$-加群の層として $U_\etale$ 上で考え、
$\ell^n$ で零化されるものとする。このとき
\begin{enumerate}
\item $H^q_\etale(X, j_!\mathcal{G})$ は有限 $\Lambda$-加群である。ただし
$\ell$ は $\text{char}(k)$ と互いに素であるとする。
\item $H^q_\etale(X, j_!\mathcal{G})$ は有限 $\Lambda$-加群である。ただし
$X$ は固有であるとする。
\end{enumerate}
\end{lemma}

\begin{proof}
$f : V \to U$ を、補題 \ref{etale-cohomology-lemma-pullback-filtered-modules} にあるような
$\ell$ と互いに素な次数の有限エタール射とする。節
\ref{etale-cohomology-section-trace-method} の議論は写像
$$
\mathcal{G} \to f_*f^{-1}\mathcal{G} \to \mathcal{G}
$$
を与え、その合成は同型である。したがって
$H^q_\etale(X, j_!f_*f^{-1}\mathcal{G})$ の有限性を証明すれば十分である。
ザリスキーの主定理（More on Morphisms、補題
\ref{more-morphisms-lemma-quasi-finite-separated-pass-through-finite})
により、可換図式
$$
\xymatrix{
V \ar[r]_{j'} \ar[d]_f & Y \ar[d]^{\overline{f}} \\
U \ar[r]^j & X
}
$$
を選べる。ここで $\overline{f} : Y \to X$ は有限であり、$j'$ は稠密像を
もつ開埋め込みである。$f$ は有限で $V$ は $Y$ で稠密だから
$V = U \times_X Y$ である。補題
\ref{etale-cohomology-lemma-compatible-shriek-push-finite} により
$$
j_!f_*f^{-1}\mathcal{G} = \overline{f}_*j'_!f^{-1}\mathcal{G}
$$
である。補題 \ref{etale-cohomology-lemma-finite-pushforward-coefficients} により
$j'_!f^{-1}\mathcal{G}$ を考えれば十分である。補題
\ref{etale-cohomology-lemma-pullback-filtered-modules} が与えるフィルトレーションの存在、
$j'_!$ が完全であること、および補題
\ref{etale-cohomology-lemma-ses-statements-coefficients} により、
$\mathcal{F} = j'_!\underline{M}$ の場合へ帰着する。ここで
$\Lambda$-加群 $M$ は有限である。これは補題
\ref{etale-cohomology-lemma-somewhat-easier-coefficients} である。
\end{proof}

\begin{theorem}
\label{etale-cohomology-theorem-vanishing-affine-curves-coefficients}
$\Lambda$ をネーター環、$k$ を代数閉体、$X$ を次元 $\leq 1$ の
$k$ 上の分離有限型スキームとする。また $\mathcal{F}$ を構成可能な
$\Lambda$-加群の層とし、$X_\etale$ 上で考え、捩れ層であるものとする。
このとき
\begin{enumerate}
\item
\label{etale-cohomology-item-finite-prime-to-p-coefficients}
$H^q_\etale(X, \mathcal{F})$ は有限 $\Lambda$-加群である。ただし
$\mathcal{F}$ は $\text{char}(k)$ と互いに素な捩れ層であるとする。
\item
\label{etale-cohomology-item-finite-proper-coefficients}
$H^q_\etale(X, \mathcal{F})$ は有限 $\Lambda$-加群である。ただし
$X$ は固有であるとする。
\end{enumerate}
\end{theorem}

\begin{proof}
以後断らずに用いる。
$\mathcal{F} = \mathcal{F}_1 \oplus \ldots \oplus \mathcal{F}_r$ と書く。
ここで $\mathcal{F}_i$ は $\ell_i^{n_i}$ で零化され、$\ell_i$ はある素数、
$n_i > 0$ は整数である。補題
\ref{etale-cohomology-lemma-ses-statements-coefficients} により $\mathcal{F}_i$ について
定理を証明すれば十分である。したがって $\ell^n$ が $\mathcal{F}$ を
零化すると仮定する。ここで $\ell$ はある素数、$n > 0$ は整数である。

\medskip\noindent
$\mathcal{F}$ は $\Lambda$-加群の層として構成可能なので、稠密開集合
$U \subset X$ が存在し、$\mathcal{F}|_U$ は有限型かつ局所定数な
$\Lambda$-加群の層である。$\dim(X) \leq 1$ なので、$U$ を縮小した後、
$U = U_1 \amalg \ldots \amalg U_n$ が既約スキームの非交和であると仮定できる
（$\geq 2$ 個の成分の交わりに含まれる $U$ の閉点を除くだけでよい）。
補題 \ref{etale-cohomology-lemma-restrict-to-open-coefficients} により
$\mathcal{F} = j_!\mathcal{G}$ の場合へ帰着する。ここで $\mathcal{G}$ は
$\Lambda$-加群の有限型局所定数層であり、$U$ 上で考える（また
$\ell^n$ で零化される）。

\medskip\noindent
$U = U_1 \amalg \ldots \amalg U_n$ を各 $U_i$ が既約となるように
選んだので、
$$
j_!\mathcal{G} =
j_{1!}(\mathcal{G}|_{U_1}) \oplus \ldots \oplus
j_{n!}(\mathcal{G}|_{U_n})
$$
である。ここで $j_i : U_i \to X$ は包含射である。
$j_{i!}(\mathcal{G}|_{U_i})$ の場合は補題
\ref{etale-cohomology-lemma-vanishing-easier-coefficients} で扱われている。
\end{proof}








\section{固有スキームの第一コホモロジー}
\label{etale-cohomology-section-finite-etale-over-proper}

\noindent
Fundamental Groups、節 \ref{pione-section-finite-etale-over-proper} では、
$R^1f_*\underline{G}$ を取る操作が、$f : X \to Y$ が固有射で
$G$ が有限群（可換とは限らない）ならば、ある意味で基底変換と可換であることを
見た。本節では、これらの結果から有用な帰結を導く。

\begin{lemma}
\label{etale-cohomology-lemma-proper-over-henselian-and-h1}
$A$ をヘンゼル局所環とする。$X$ を $A$ 上の固有スキームとし、その閉ファイバーを
$X_0$ とする。$M$ を有限アーベル群とする。このとき
$H^1_\etale(X, \underline{M}) = H^1_\etale(X_0, \underline{M})$ である。
\end{lemma}

\begin{proof}
Cohomology on Sites、補題 \ref{sites-cohomology-lemma-torsors-h1} により、
$H^1_\etale(X, \underline{M})$ の元は、$\underline{M}$-トーサー
$\mathcal{F}$ に対応し、これは $X_\etale$ 上で考えられる。
このようなトーサーは明らかに
有限局所定数層である。したがって補題
\ref{etale-cohomology-lemma-characterize-finite-locally-constant} により、$\mathcal{F}$ は
スキーム $V$ により表現され、これは $X$ 上有限エタールである。逆に、
スキーム $V$ が $X$ 上有限エタールであり、$M$-作用によって
$M$-トーサーとして $X$ 上に置かれるならば、コホモロジー類が得られる。$X_0$ 上の
コホモロジー類と $X_0$ 上有限エタールなトーサーとの間にも同じ翻訳が成り立つ。
したがって補題は Fundamental Groups、補題
\ref{pione-lemma-finite-etale-on-proper-over-henselian}
の圏同値から従う。
\end{proof}

\noindent
次の技術的補題は固有基底変換定理の証明における主要な材料である。
議論は、$A$ 上の任意の固有スキームで特殊ファイバーの次元が
$\leq 1$ であるものに
対して逐語的に成り立つ。しかし実際には結論は固有基底変換定理の帰結となり、
その証明に必要なのは次の特別な形だけである。

\begin{lemma}
\label{etale-cohomology-lemma-efface-cohomology-on-fibre-by-finite-cover}
$A$ をヘンゼル局所環とする。$X = \mathbf{P}^1_A$ とする。
$X_0 \subset X$ を閉ファイバーとする。$\ell$ を素数とする。
$\mathcal{I}$ を $\mathbf{Z}/\ell\mathbf{Z}$-加群の入射層とし、
$X_\etale$ 上で考える。このとき $H^q_\etale(X_0, \mathcal{I}|_{X_0}) = 0$ が
$q > 0$ に対して成り立つ。
\end{lemma}

\begin{proof}
$X$ は $2$ 個のアフィン開集合で被覆できる分離スキームである。したがって
$q > 1$ の場合、これは固有基底変換定理のガバーによるアフィン版から従う。
補題 \ref{etale-cohomology-lemma-vanishing-restriction-injective} を参照されたい。
ゆえに $q = 1$ と仮定してよい。
$\xi \in H^1_\etale(X_0, \mathcal{I}|_{X_0})$ とする。目標は $\xi$ が
$0$ であることを示すことである。補題
\ref{etale-cohomology-lemma-torsion-colimit-constructible} および \ref{etale-cohomology-lemma-colimit} により、
写像 $\mathcal{F} \to \mathcal{I}$ を、$\mathcal{F}$ が構成可能な
$\mathbf{Z}/\ell\mathbf{Z}$-加群の層であり、$\xi$ がある元 $\zeta$ から
来るように選べる。この元は
$H^1_\etale(X_0, \mathcal{F}|_{X_0})$ に属する。層の入射写像
$\mathcal{F} \to \mathcal{F}'$ が与えられたとする。ここでこれらは
$\mathbf{Z}/\ell\mathbf{Z}$-加群の層であり、$X_\etale$ 上で考える。
$\mathcal{I}$ は入射的なので、与えられた写像
$\mathcal{F} \to \mathcal{I}$ を写像 $\mathcal{F}' \to \mathcal{I}$ へ延長できる。
この状況では $\mathcal{F}$ を $\mathcal{F}'$ で置き換え、$\zeta$ を
$\zeta$ の像で
置き換えてよい。この像は $H^1_\etale(X_0, \mathcal{F}'|_{X_0})$ に属する。
また $\mathcal{F} = \mathcal{F}_1 \oplus \mathcal{F}_2$ が直和ならば、
$\mathcal{F}$ を $\mathcal{F}_i$ で置き換え、$\zeta$ を $\zeta$ の像で置き換えてよい。
この像は $H^1_\etale(X_0, \mathcal{F}_i|_{X_0})$ に属する。

\medskip\noindent
補題 \ref{etale-cohomology-lemma-constructible-maps-into-constant-general} と上の注意により、
$\mathcal{F}$ は $f_*\underline{M}$ の形であると仮定してよい。ここで $M$ は
有限 $\mathbf{Z}/\ell\mathbf{Z}$-加群であり、$f : Y \to X$ は有限表示な
有限射である（補題 \ref{etale-cohomology-lemma-finite-pushforward-constructible} により
このような層も構成可能だが、ここでは用いない）。
$f_*$ の形成は任意の基底変換と可換なので（補題
\ref{etale-cohomology-lemma-finite-pushforward-commutes-with-base-change}）、
$f_*\underline{M}$ の $X_0$ への制限は、$\underline{M}$ の、特殊ファイバーの
誘導射 $Y_0 \to X_0$ による前進と等しい。ルレイ・スペクトル系列
（命題 \ref{etale-cohomology-proposition-leray}）および高次順像の消滅
（命題 \ref{etale-cohomology-proposition-finite-higher-direct-image-zero}）により
$$
H^1_\etale(X_0, f_*\underline{M}|_{X_0}) = H^1_\etale(Y_0, \underline{M}).
$$
$Y \to \Spec(A)$ は固有なので、補題
\ref{etale-cohomology-lemma-proper-over-henselian-and-h1} により
$H^1_\etale(Y_0, \underline{M})$ は $H^1_\etale(Y, \underline{M})$ に等しい。
したがってコホモロジー類 $\zeta$ はコホモロジー類
$$
\tilde\zeta \in H^1_\etale(Y, \underline{M}) = H^1_\etale(X, f_*\underline{M})
$$
へ持ち上がる。しかし $\tilde \zeta$ は
$H^1_\etale(X, \mathcal{I})$ において零へ写る。これは $\mathcal{I}$ が
入射的だからである。また次の図式の可換性により、
$$
\xymatrix{
H^1_\etale(X, f_*\underline{M}) \ar[r] \ar[d] &
H^1_\etale(X, \mathcal{I}) \ar[d] \\
H^1_\etale(X_0, (f_*\underline{M})|_{X_0}) \ar[r] &
H^1_\etale(X_0, \mathcal{I}|_{X_0})
}
$$
像 $\xi$（$\zeta$ の像）も零であると結論する。
\end{proof}











\section{基底変換の準備}
\label{etale-cohomology-section-base-change-preliminaries}

\noindent
滑らかな基底変換定理または固有基底変換定理のいずれかに関心がある読者は、
対応する節まで直接読み飛ばしてよい。
本節および続く数節では、スキームの可換図式
$$
\xymatrix{
X \ar[d]_f & Y \ar[l]^h \ar[d]^e \\
S & T \ar[l]_g
}
$$
を考える。通常、この図式はカルテシアン、すなわち
$Y = X \times_S T$ であると仮定する。上の可換図式から、小エタール・サイトの可換図式
$$
\xymatrix{
X_\etale \ar[d]_{f_{small}} & Y_\etale \ar[d]^{e_{small}} \ar[l]^{h_{small}} \\
S_\etale & T_\etale \ar[l]_{g_{small}}
}
$$
が得られる。次の記法を用いることにする。
$$
f^{-1} = f_{small}^{-1}, \quad
g_* = g_{small, *}, \quad
e^{-1} = e_{small}^{-1}, \text{かつ}\quad
h_* = h_{small, *}.
$$
Sites の節 \ref{sites-section-pullback} により、基底変換写像（または引戻し写像）
$$
f^{-1}g_*\mathcal{F}
\longrightarrow
h_*e^{-1}\mathcal{F}
$$
を得る。これは層 $\mathcal{F}$ が $T_\etale$ 上にある場合のものである。
$\mathcal{F}$ が $T_\etale$ 上のアーベル層ならば、導来基底変換写像
$$
f^{-1}Rg_*\mathcal{F}
\longrightarrow
Rh_*e^{-1}\mathcal{F}
$$
を得る。Cohomology on Sites の補題
\ref{sites-cohomology-lemma-base-change-map-flat-case} を参照されたい。
最後に、$K$ が $D(T_\etale)$ の任意の対象ならば、基底変換写像
$$
f^{-1}Rg_*K
\longrightarrow
Rh_*e^{-1}K
$$
がある。Cohomology on Sites の注意
\ref{sites-cohomology-remark-base-change} を参照されたい。

\begin{lemma}
\label{etale-cohomology-lemma-base-change-local}
スキームのカルテシアン図式
$$
\xymatrix{
X \ar[d]_f & Y \ar[l]^h \ar[d]^e \\
S & T \ar[l]_g
}
$$
を考える。$\{U_i \to X\}$ を、$U_i \to S$ が
$U_i \to V_i \to S$ と分解し、$V_i \to S$ がエタールとなるようなエタール被覆とし、
カルテシアン図式
$$
\xymatrix{
U_i \ar[d]_{f_i} & U_i \times_X Y \ar[l]^{h_i} \ar[d]^{e_i} \\
V_i & V_i \times_S T \ar[l]_{g_i}
}
$$
を考える。$\mathcal{F}$ を $T_\etale$ 上の層とし、$K$ を $D(T_\etale)$ の対象とする。
$K_i = K|_{V_i \times_S T}$ および
$\mathcal{F}_i = \mathcal{F}|_{V_i \times_S T}$ とおく。
\begin{enumerate}
\item $f_i^{-1}g_{i, *}\mathcal{F}_i = h_{i, *}e_i^{-1}\mathcal{F}_i$
がすべての $i$ に対して成り立つならば、$f^{-1}g_*\mathcal{F} = h_*e^{-1}\mathcal{F}$ である。
\item $f_i^{-1}Rg_{i, *}K_i = Rh_{i, *}e_i^{-1}K_i$
がすべての $i$ に対して成り立つならば、$f^{-1}Rg_*K = Rh_*e^{-1}K$ である。
\item $\mathcal{F}$ がアーベル層であり、
$f_i^{-1}R^qg_{i, *}\mathcal{F}_i = R^qh_{i, *}e_i^{-1}\mathcal{F}_i$
がすべての $i$ に対して成り立つならば、
$f^{-1}R^qg_*\mathcal{F} = R^qh_*e^{-1}\mathcal{F}$ である。
\end{enumerate}
\end{lemma}

\begin{proof}
(1) を証明する。まず
$$
(f^{-1}g_*\mathcal{F})|_{U_i} =
f_i^{-1}(g_*\mathcal{F}|_{V_i}) =
f_i^{-1}g_{i, *}\mathcal{F}_i
$$
であることに注意する。第一の等号は $U_i \to X \to S$ が
$U_i \to V_i \to S$ に等しいことから、第二の等号は
Sites の補題 \ref{sites-lemma-localize-morphism-strong} により
$g_*\mathcal{F}|_{V_i} = g_{i, *}\mathcal{F}_i$ であることから従う。
同様に
$$
(h_*e^{-1}\mathcal{F})|_{U_i} =
h_{i, *}(e^{-1}\mathcal{F}|_{U_i \times_X Y}) =
h_{i, *}e_i^{-1}\mathcal{F}_i
$$
を得る。したがって、基底変換写像
$f_i^{-1}g_{i, *}\mathcal{F}_i \to h_{i, *}e_i^{-1}\mathcal{F}_i$
がすべての $i$ に対して同型ならば、基底変換写像
$f^{-1}g_*\mathcal{F} \to h_*e^{-1}\mathcal{F}$
の $U_i$ 上への制限はすべての $i$ に対して同型である。
$\{U_i \to X\}$ はエタール被覆なので、この写像自身が同型であると結論する。

\medskip\noindent
残り二つの主張については、Sites の補題
\ref{sites-lemma-localize-morphism-strong} の代わりに
Cohomology on Sites の補題
\ref{sites-cohomology-lemma-restrict-direct-image-open} を用いればよい。
\end{proof}

\begin{lemma}
\label{etale-cohomology-lemma-base-change-compose}
スキームのカルテシアン図式の塔
$$
\xymatrix{
W \ar[d]_i & Z \ar[l]^j \ar[d]^k \\
X \ar[d]_f & Y \ar[l]^h \ar[d]^e \\
S & T \ar[l]_g
}
$$
を考える。$K$ を $D(T_\etale)$ の対象とする。写像
$$
f^{-1}Rg_*K \to Rh_*e^{-1}K
\quad\text{および}\quad
i^{-1}Rh_*e^{-1}K \to Rj_*k^{-1}e^{-1}K
$$
が同型ならば、
$(f \circ i)^{-1}Rg_*K \to Rj_*(e \circ k)^{-1}K$
は同型である。
同様に、$\mathcal{F}$ が $T_\etale$ 上のアーベル層であり、写像
$$
f^{-1}R^qg_*\mathcal{F} \to R^qh_*e^{-1}\mathcal{F}
\quad\text{および}\quad
i^{-1}R^qh_*e^{-1}\mathcal{F} \to R^qj_*k^{-1}e^{-1}\mathcal{F}
$$
が同型ならば、
$(f \circ i)^{-1}R^qg_*\mathcal{F} \to R^qj_*(e \circ k)^{-1}\mathcal{F}$
は同型である。
\end{lemma}

\begin{proof}
これらの基底変換写像の合成が外側の長方形に対する基底変換写像であることを確認すれば、
主張は形式的に従う。Cohomology on Sites の注意
\ref{sites-cohomology-remark-compose-base-change-horizontal} を参照されたい。
\end{proof}

\begin{lemma}
\label{etale-cohomology-lemma-base-change-Rf-star-colim}
$I$ を有向集合とする。アフィンな遷移射をもつスキームのカルテシアン図式の逆系
$$
\xymatrix{
X_i \ar[d]_{f_i} & Y_i \ar[l]^{h_i} \ar[d]^{e_i} \\
S_i & T_i \ar[l]_{g_i}
}
$$
を考え、$g_i$ は準コンパクトかつ準分離であるとする。
$X = \lim X_i$、$S = \lim S_i$、$T = \lim T_i$ および $Y = \lim Y_i$ とおくと、
カルテシアン図式
$$
\xymatrix{
X \ar[d]_f & Y \ar[l]^h \ar[d]^e \\
S & T \ar[l]_g
}
$$
を得る。$(\mathcal{F}_i, \varphi_{i'i})$ を定義
\ref{etale-cohomology-definition-inverse-system-sheaves} にいう $(T_i)$ 上の層の系とする。
$\mathcal{F} = \colim p_i^{-1}\mathcal{F}_i$ を $T$ 上で定める。ここで
$p_i : T \to T_i$ は射影である。このとき次が成り立つ。
\begin{enumerate}
\item $f_i^{-1}g_{i, *}\mathcal{F}_i = h_{i, *}e_i^{-1}\mathcal{F}_i$
がすべての $i$ に対して成り立つならば、
$f^{-1}g_*\mathcal{F} = h_*e^{-1}\mathcal{F}$ である。
\item $\mathcal{F}_i$ がすべての $i$ に対してアーベル層であり、
$f_i^{-1}R^qg_{i, *}\mathcal{F}_i = R^qh_{i, *}e_i^{-1}\mathcal{F}_i$
がすべての $i$ に対して成り立つならば、
$f^{-1}R^qg_*\mathcal{F} = R^qh_*e^{-1}\mathcal{F}$ である。
\end{enumerate}
\end{lemma}

\begin{proof}
(2) を証明し、(1) の証明は省略する。層の引戻しは左随伴なので余極限と可換することを、
以下では断りなく用いる。$h_i$ は $g_i$ の基底変換であるから、準コンパクトかつ準分離である。
$q_i : Y \to Y_i$ を射影と書くと、
$e^{-1}\mathcal{F} = \colim e^{-1}p_i^{-1}\mathcal{F}_i =
\colim q_i^{-1}e_i^{-1}\mathcal{F}_i$ である。補題
\ref{etale-cohomology-lemma-relative-colimit-general} により
$$
R^qh_*e^{-1}\mathcal{F} = \colim r_i^{-1}R^qh_{i, *}e_i^{-1}\mathcal{F}_i
$$
を得る。ここで $r_i : X \to X_i$ は射影である。同様に
$$
f^{-1}Rg_*\mathcal{F} =
f^{-1}\colim s_i^{-1}R^qg_{i, *}\mathcal{F}_i =
\colim r_i^{-1}f_i^{-1}R^qg_{i, *}\mathcal{F}_i
$$
である。ここで $s_i : S \to S_i$ は射影である。これで補題が従う。
\end{proof}

\begin{lemma}
\label{etale-cohomology-lemma-base-change-Rf-star-colim-complexes}
$I$、$X_i$、$Y_i$、$S_i$、$T_i$、$f_i$、$h_i$、$e_i$、$g_i$、
$X$、$Y$、$S$、$T$、$f$、$h$、$e$、$g$ は補題
\ref{etale-cohomology-lemma-base-change-Rf-star-colim} の主張におけるものとする。
$0 \in I$ とし、$K_0 \in D^+(T_{0, \etale})$ とする。
$i \in I$、$i \geq 0$ に対して、$K_i$ を $K_0$ の $T_i$ への引戻しと書く。
$K$ を $K_0$ の $T$ への引戻しと書く。
$f_i^{-1}Rg_{i, *}K_i = Rh_{i, *}e_i^{-1}K_i$
がすべての $i \geq 0$ に対して成り立つならば、$f^{-1}Rg_*K = Rh_*e^{-1}K$ である。
\end{lemma}

\begin{proof}
基底変換写像
$f^{-1}Rg_*K \to Rh_*e^{-1}K$ がコホモロジー層上の同型を誘導することを示せば十分である。
言い換えると、$f^{-1}R^pg_*K \to R^ph_*e^{-1}K$ がすべての
$p \in \mathbf{Z}$ に対して同型であることを示さなければならない。その仮定は
$f_i^{-1}R^pg_{i, *}K_i \to R^ph_{i, *}e_i^{-1}K_i$ がすべての
$i \geq 0$ および $p \in \mathbf{Z}$ に対して同型であるということである。
ここで補題 \ref{etale-cohomology-lemma-base-change-Rf-star-colim} の証明において、
補題 \ref{etale-cohomology-lemma-relative-colimit-general} の代わりに補題
\ref{etale-cohomology-lemma-relative-colimit-general-complexes} を参照すれば、まったく同じ議論ができる。
\end{proof}

\begin{lemma}
\label{etale-cohomology-lemma-base-change-f-star-general-stalks}
スキームのカルテシアン図式
$$
\xymatrix{
X \ar[d]_f & Y \ar[l]^h \ar[d]^e \\
S & T \ar[l]_g
}
$$
を考え、$g : T \to S$ は準コンパクトかつ準分離であるとする。
$\mathcal{F}$ を $T_\etale$ 上のアーベル層とし、$q \geq 0$ とする。
次の条件は同値である。
\begin{enumerate}
\item 任意の幾何学的点 $\overline{x}$ が $X$ 上にあり、その像が
$\overline{s} = f(\overline{x})$ であるとき
$$
H^q(\Spec(\mathcal{O}^{sh}_{X, \overline{x}}) \times_S T, \mathcal{F})
=
H^q(\Spec(\mathcal{O}^{sh}_{S, \overline{s}}) \times_S T, \mathcal{F})
$$
が成り立つ。
\item $f^{-1}R^qg_*\mathcal{F} \to R^qh_*e^{-1}\mathcal{F}$ は同型である。
\end{enumerate}
\end{lemma}

\begin{proof}
$Y = X \times_S T$ であるから、
$\Spec(\mathcal{O}^{sh}_{X, \overline{x}}) \times_X Y =
\Spec(\mathcal{O}^{sh}_{X, \overline{x}}) \times_S T$ である。したがって定理
\ref{etale-cohomology-theorem-higher-direct-images}（および補題
\ref{etale-cohomology-lemma-stalk-pullback}）により、(1) の写像は (2) の写像の
$\overline{x}$ における茎の写像である。ゆえに定理
\ref{etale-cohomology-theorem-exactness-stalks} から結果が従う。
\end{proof}

\begin{lemma}
\label{etale-cohomology-lemma-check-stalks-better}
$f : X \to S$ をスキームの射とする。
$\overline{x}$ を $X$ の幾何学的点とし、その像 $\overline{s}$ を $S$ における点とする。
$\Spec(K) \to \Spec(\mathcal{O}^{sh}_{S, \overline{s}})$ を射とし、
$K$ は分離閉体であるとする。$\mathcal{F}$ を
$\Spec(K)_\etale$ 上のアーベル層とし、$q \geq 0$ とする。次の条件は同値である。
\begin{enumerate}
\item
$H^q(\Spec(\mathcal{O}^{sh}_{X, \overline{x}}) \times_S \Spec(K), \mathcal{F}) =
H^q(\Spec(\mathcal{O}^{sh}_{S, \overline{s}}) \times_S \Spec(K), \mathcal{F})$
\item
$H^q(\Spec(\mathcal{O}^{sh}_{X, \overline{x}})
\times_{\Spec(\mathcal{O}^{sh}_{S, \overline{s}})} \Spec(K), \mathcal{F}) =
H^q(\Spec(K), \mathcal{F})$
\end{enumerate}
\end{lemma}

\begin{proof}
$\Spec(K) \times_S \Spec(\mathcal{O}^{sh}_{S, \overline{s}})$ は、$K$ 上の
エタール代数のフィルター付き余極限のスペクトルであることに注意する。
$K$ は分離閉なので、各エタール $K$-代数は $K$ のコピーの有限積である。したがって
$$
\Spec(K) \times_S \Spec(\mathcal{O}^{sh}_{S, \overline{s}}) =
\lim_{i \in I} \coprod\nolimits_{a \in A_i} \Spec(K)
$$
と余フィルター付き極限として書ける。ここで各項は $\Spec(K)$ のコピーの非交和であり、
有限集合 $A_i$ で添字づけられている。$A_i$ は、写像
$\Spec(K) \to \Spec(\mathcal{O}^{sh}_{S, \overline{s}})$ が与えられているため空でない。
したがって
\begin{align*}
\Spec(\mathcal{O}^{sh}_{X, \overline{x}}) \times_S \Spec(K)
& =
\Spec(\mathcal{O}^{sh}_{X, \overline{x}})
\times_{\Spec(\mathcal{O}^{sh}_{S, \overline{s}})}
\left(
\Spec(\mathcal{O}^{sh}_{S, \overline{s}})
\times_S \Spec(K)\right) \\
& =
\lim_{i \in I} \coprod\nolimits_{a \in A_i}
\Spec(\mathcal{O}^{sh}_{X, \overline{x}})
\times_{\Spec(\mathcal{O}^{sh}_{S, \overline{s}})} \Spec(K)
\end{align*}
この状況ではコホモロジーを取る操作はスキームの極限と可換するので
（定理 \ref{etale-cohomology-theorem-colimit}）、結論が従う。
\end{proof}













\section{順像の基底変換}
\label{etale-cohomology-section-base-change-f-star}

\noindent
本節は準備的な内容なので、初読時には飛ばしてよい。
本節では、どのような射 $f : X \to S$ に対し、任意のカルテシアン図式
についてすべての（集合の）層上で $f^{-1}g_* = h_*e^{-1}$ が成り立つかを論じる。
$$
\xymatrix{
X \ar[d]_f & Y \ar[l]^h \ar[d]^e \\
S & T \ar[l]_g
}
$$
ここで $g$ は準コンパクトかつ準分離である。

\begin{lemma}
\label{etale-cohomology-lemma-base-change-f-star-general}
スキームのカルテシアン図式
$$
\xymatrix{
X \ar[d]_f & Y \ar[l]^h \ar[d]^e \\
S & T \ar[l]_g
}
$$
を考える。$f$ は平坦であり、任意の対象 $U$ が $X_\etale$ に属するとき、
その対象は被覆 $\{U_i \to U\}$ をもち、$U_i \to S$ は
$U_i \to V_i \to S$ と分解し、$V_i \to S$ はエタール、$U_i \to V_i$ は
準コンパクトで幾何学的に連結なファイバーをもつと仮定する。
このとき、任意の集合の層 $\mathcal{F}$ が $T_\etale$ 上にあれば
$f^{-1}g_*\mathcal{F} = h_*e^{-1}\mathcal{F}$ である。
\end{lemma}

\begin{proof}
$U \to X$ をエタール射とし、$U \to S$ は $U \to V \to S$ と分解し、
$V \to S$ はエタール、$U \to V$ は準コンパクトで幾何学的に連結なファイバーを
もつとする。$U \to V$ は平坦であることに注意する
（More on Flatness の補題 \ref{flat-lemma-etale-flat-up-down}）。次を主張する。
\begin{align*}
f^{-1}g_*\mathcal{F}(U)
& = g_*\mathcal{F}(V) \\
& = \mathcal{F}(V \times_S T) \\
& = e^{-1}\mathcal{F}(U \times_X Y) \\
& = h_*e^{-1}\mathcal{F}(U)
\end{align*}
実際、$U$ を $X_\etale$ の対象、$V$ を $S_\etale$ の対象とみなすと、
第一の等号は補題 \ref{etale-cohomology-lemma-sections-upstairs} から従う\footnote{厳密には、
$f^{-1}g_*\mathcal{F}$ の $U_\etale$ への制限が、$U \to V$ による
$g_*\mathcal{F}$ の $V_\etale$ への制限の引戻しであることも用いている。
Sites の補題 \ref{sites-lemma-localize-morphism-strong} を参照されたい。}。
$V \times_S T$ を $T_\etale$ の対象とみなせば、第二の等号は $g_*$ の定義から従う。
$U \times_X Y = U \times_S T$（なぜなら $Y = X \times_S T$）であることに注意する。
したがって $U \times_X Y \to V \times_S T$ は、$U \to V$ の基底変換として
幾何学的に連結なファイバーをもつ。$U \times_X Y$ を $Y_\etale$ の対象とみなすと、
第三の等号は先ほどと同様に補題 \ref{etale-cohomology-lemma-sections-upstairs} から従う。
最後に、第四の等号は $h_*$ の定義から従う。

\medskip\noindent
仮定により、$X_\etale$ の任意の対象は前段落の議論を適用できるエタール被覆をもつ。
したがって補題が成り立つ。
\end{proof}

\begin{lemma}
\label{etale-cohomology-lemma-fppf-reduced-fibres-base-change-f-star}
スキームのカルテシアン図式
$$
\xymatrix{
X \ar[d]_f & Y \ar[l]^h \ar[d]^e \\
S & T \ar[l]_g
}
$$
を考える。$f$ は平坦かつ局所有限表示であり、幾何学的に被約なファイバーをもつとする。
このとき $f^{-1}g_*\mathcal{F} = h_*e^{-1}\mathcal{F}$ が、
任意の層 $\mathcal{F}$ が $T_\etale$ 上にある場合に成り立つ。
\end{lemma}

\begin{proof}
補題 \ref{etale-cohomology-lemma-base-change-f-star-general} と More on Morphisms の補題
\ref{more-morphisms-lemma-cover-by-geometrically-connected} を組み合わせればよい。
\end{proof}

\begin{lemma}
\label{etale-cohomology-lemma-base-change-f-star-field}
スキームのカルテシアン図式
$$
\xymatrix{
X \ar[d]_f & Y \ar[l]^h \ar[d]^e \\
S & T \ar[l]_g
}
$$
を考える。$S$ は分離閉体のスペクトルであると仮定する。
このとき $f^{-1}g_*\mathcal{F} = h_*e^{-1}\mathcal{F}$ が、
任意の層 $\mathcal{F}$ が $T_\etale$ 上にある場合に成り立つ。
\end{lemma}

\begin{proof}
$X$ 上局所的に議論してよい。したがって $X$ はアフィンであると仮定できる。
このとき $X$ は $S$ 上有限型のアフィン・スキームの余フィルター付き極限として書ける。
補題 \ref{etale-cohomology-lemma-base-change-Rf-star-colim} により、$X$ は $S$ 上有限型であると仮定してよい。
すると補題 \ref{etale-cohomology-lemma-base-change-f-star-general} が適用できる。実際、分離閉体上有限型の
任意のスキームは、連結かつ幾何学的に連結なスキームの有限非交和である
（Varieties の補題
\ref{varieties-lemma-separably-closed-field-connected-components} を参照）。
\end{proof}

\begin{lemma}
\label{etale-cohomology-lemma-base-change-f-star-valuation}
スキームのカルテシアン図式
$$
\xymatrix{
X \ar[d]_f & Y \ar[l]^h \ar[d]^e \\
S & T \ar[l]_g
}
$$
を考える。次を仮定する。
\begin{enumerate}
\item $f$ は平坦かつ開である。
\item $S$ の剰余体は分離的代数閉である。
\item $U \to X$ をエタール射とし $U$ がアフィンならば、$U$ は $X$ の開部分スキームの
有限非交和として書ける（たとえば $X$ が分離閉な関数体をもつ正規整スキームである場合）。
\item ファイバー $X_s$（$f$ に関するもの）の任意の空でない開集合は連結である
（たとえば $X_s$ が既約または空である場合）。
\end{enumerate}
このとき、任意の集合の層 $\mathcal{F}$ が $T_\etale$ 上にあれば
$f^{-1}g_*\mathcal{F} = h_*e^{-1}\mathcal{F}$ である。
\end{lemma}

\begin{proof}
省略する。ヒント：仮定から補題 \ref{etale-cohomology-lemma-base-change-f-star-general} の条件がほとんど自明に従う。
(3) の括弧内の例は補題
\ref{etale-cohomology-lemma-normal-scheme-with-alg-closed-function-field} から従う。
\end{proof}

\noindent
次の補題は本来ここに置くべきものではないが、ほかに適当な場所も見当たらない。

\begin{lemma}
\label{etale-cohomology-lemma-fppf-reduced-fibres-pullback-products}
$f : X \to S$ を、平坦かつ局所有限表示で幾何学的に被約なファイバーをもつスキームの射とする。
このとき $f^{-1} : \Sh(S_\etale) \to \Sh(X_\etale)$ は積と可換する。
\end{lemma}

\begin{proof}
$I$ を集合とし、$\mathcal{G}_i$ を $S_\etale$ 上の層とする（$i \in I$）。
$U \to X$ をエタール射とし、$U \to S$ は $U \to V \to S$ と分解し、
$V \to S$ はエタール、$U \to V$ は有限表示平坦で幾何学的に連結なファイバーをもつとする。
このとき
\begin{align*}
f^{-1}(\prod \mathcal{G}_i)(U)
& =
(\prod \mathcal{G}_i)(V) \\
& =
\prod \mathcal{G}_i(V) \\
& =
\prod f^{-1}\mathcal{G}_i(U) \\
& =
(\prod f^{-1}\mathcal{G}_i)(U)
\end{align*}
である。第一および第三の等号では補題 \ref{etale-cohomology-lemma-sections-upstairs} を用いた
（さらに、$f^{-1}\mathcal{G}$ の $U_\etale$ への制限が、$U \to V$ による
$\mathcal{G}$ の $V_\etale$ への制限の引戻しであることも用いている。
Sites の補題 \ref{sites-lemma-localize-morphism-strong} を参照）。
More on Morphisms の補題
\ref{more-morphisms-lemma-cover-by-geometrically-connected} により、
任意の対象 $U$ が $X_\etale$ に属するとき、その対象はエタール被覆
$\{U_i \to U\}$ をもち、前段落の議論を $U_i$ に適用できる。これで補題が従う。
\end{proof}

\begin{lemma}
\label{etale-cohomology-lemma-base-change-f-star}
$f : X \to S$ を平坦なスキームの射とし、任意の幾何学的点 $\overline{x}$ が $X$ 上にあるとき、写像
$$
\mathcal{O}_{S, f(\overline{x})}^{sh}
\longrightarrow
\mathcal{O}_{X, \overline{x}}^{sh}
$$
が幾何学的に連結なファイバーをもつとする。このとき、スキームの任意のカルテシアン図式
$$
\xymatrix{
X \ar[d]_f & Y \ar[l]^h \ar[d]^e \\
S & T \ar[l]_g
}
$$
において $g$ が準コンパクトかつ準分離ならば、
$f^{-1}g_*\mathcal{F} = h_*e^{-1}\mathcal{F}$
が、任意の集合の層 $\mathcal{F}$ が $T_\etale$ 上にある場合に成り立つ。
\end{lemma}

\begin{proof}
茎上で等号を確認すれば十分である。定理
\ref{etale-cohomology-theorem-exactness-stalks} を参照されたい。定理
\ref{etale-cohomology-theorem-higher-direct-images} により
$$
(h_*e^{-1}\mathcal{F})_{\overline{x}} =
\Gamma(\Spec(\mathcal{O}_{X, \overline{x}}^{sh}) \times_X Y, e^{-1}\mathcal{F})
$$
であり、同様に
$$
(f^{-1}g_*^{-1}\mathcal{F})_{\overline{x}} =
(g_*^{-1}\mathcal{F})_{f(\overline{x})} =
\Gamma(\Spec(\mathcal{O}_{S, f(\overline{x})}^{sh}) \times_S T, \mathcal{F})
$$
これらの集合は、射
$$
\Spec(\mathcal{O}_{X, \overline{x}}^{sh}) \times_X Y
\longrightarrow
\Spec(\mathcal{O}_{S, f(\overline{x})}^{sh}) \times_S T
$$
に補題 \ref{etale-cohomology-lemma-sections-upstairs} を適用することにより等しい。この射は
$\Spec(\mathcal{O}_{X, \overline{x}}^{sh}) \to
\Spec(\mathcal{O}_{S, f(\overline{x})}^{sh})$
の基底変換である。実際、$Y = X \times_S T$ である。
\end{proof}







\section{高次順像の基底変換}
\label{etale-cohomology-section-base-change}

\noindent
本節は高次順像についての節 \ref{etale-cohomology-section-base-change-f-star} の類似である。
本節は準備的な内容なので、初読時には飛ばしてよい。

\begin{remark}
\label{etale-cohomology-remark-base-change-holds}
$f : X \to S$ をスキームの射とし、$n$ を整数とする。
$BC(f, n, q_0)$ が真であるとは、任意の可換図式
$$
\xymatrix{
X \ar[d]_f & X' \ar[l] \ar[d]_{f'} & Y \ar[l]^h \ar[d]^e \\
S & S' \ar[l] & T \ar[l]_g
}
$$
において $X' = X \times_S S'$ および $Y = X' \times_{S'} T$ であり、
$g$ が準コンパクトかつ準分離で、任意のアーベル層 $\mathcal{F}$ が
$T_\etale$ 上にあり $n$ により零化されるとき、その基底変換写像
$$
(f')^{-1}R^qg_*\mathcal{F}
\longrightarrow
R^qh_*e^{-1}\mathcal{F}
$$
が $q \leq q_0$ に対して同型であることをいう。
\end{remark}

\begin{lemma}
\label{etale-cohomology-lemma-base-change-q-injective}
$f : X \to S$ および $n$ は注意 \ref{etale-cohomology-remark-base-change-holds} のとおりとし、
ある $q \geq 1$ に対して $BC(f, n, q - 1)$ が成り立つと仮定する。このとき任意の可換図式
$$
\xymatrix{
X \ar[d]_f & X' \ar[l] \ar[d]_{f'} & Y \ar[l]^h \ar[d]^e \\
S & S' \ar[l] & T \ar[l]_g
}
$$
において $X' = X \times_S S'$ および $Y = X' \times_{S'} T$ であり、
$g$ が準コンパクトかつ準分離で、任意のアーベル層 $\mathcal{F}$ が
$T_\etale$ 上にあり $n$ により零化されるとき、次が成り立つ。
\begin{enumerate}
\item 基底変換写像
$(f')^{-1}R^qg_*\mathcal{F}\to R^qh_*e^{-1}\mathcal{F}$
は単射である。
\item $\mathcal{F} \subset \mathcal{G}$ であり、$\mathcal{G}$ が
$T_\etale$ 上で $n$ により零化されるならば、
$$
\Coker\left(
(f')^{-1}R^qg_*\mathcal{F}\to R^qh_*e^{-1}\mathcal{F}
\right)
\subset
\Coker\left(
(f')^{-1}R^qg_*\mathcal{G}\to R^qh_*e^{-1}\mathcal{G}
\right)
$$
\item (2) において層 $\mathcal{G}$ が $\mathbf{Z}/n\mathbf{Z}$-加群の単射層ならば、
$$
\Coker\left((f')^{-1}R^qg_*\mathcal{F}\to R^qh_*e^{-1}\mathcal{F} \right)
\subset R^qh_*e^{-1}\mathcal{G}
$$
\end{enumerate}
\end{lemma}

\begin{proof}
短完全列
$0 \to \mathcal{F} \to \mathcal{I} \to \mathcal{Q} \to 0$
を選び、$\mathcal{I}$ は $\mathbf{Z}/n\mathbf{Z}$-加群の単射層であるとする。
誘導される図式
$$
\xymatrix{
(f')^{-1}R^{q - 1}g_*\mathcal{I} \ar[d]_{\cong} \ar[r] &
(f')^{-1}R^{q - 1}g_*\mathcal{Q} \ar[d]_{\cong} \ar[r] &
(f')^{-1}R^qg_*\mathcal{F} \ar[d] \ar[r] &
0 \ar[d] \\
R^{q - 1}h_*e^{-1}\mathcal{I} \ar[r] &
R^{q - 1}h_*e^{-1}\mathcal{Q} \ar[r] &
R^qh_*e^{-1}\mathcal{F} \ar[r] &
R^qh_*e^{-1}\mathcal{I}
}
$$
を考える。各行は完全である。$\mathcal{I}$ が単射的なので右上隅は零である。
左側の二つの垂直矢印は $BC(f, n, q - 1)$ により同型である。したがって (1) が従う。
上の図式はさらに
$$
\Coker\left(
(f')^{-1}R^qg_*\mathcal{F}\to R^qh_*e^{-1}\mathcal{F}
\right)
\subset
R^qh_*e^{-1}\mathcal{I}
$$
も示すので、(3) が成り立つ。(2) を証明するには
$\mathcal{F} \subset \mathcal{G} \subset \mathcal{I}$ と選べばよい。
\end{proof}

\begin{lemma}
\label{etale-cohomology-lemma-base-change-q-integral-top}
$f : X \to S$ および $n$ は注意 \ref{etale-cohomology-remark-base-change-holds} のとおりとし、
ある $q \geq 1$ に対して $BC(f, n, q - 1)$ が成り立つと仮定する。可換図式
$$
\vcenter{
\xymatrix{
X \ar[d]_f &
X' \ar[d]_{f'} \ar[l] &
Y \ar[l]^h \ar[d]^e &
Y' \ar[l]^{\pi'} \ar[d]^{e'} \\
S &
S' \ar[l] &
T \ar[l]_g &
T' \ar[l]_\pi
}
}
\quad\text{および}\quad
\vcenter{
\xymatrix{
X' \ar[d]_{f'} & & Y' \ar[ll]^{h' = h \circ \pi'} \ar[d]^{e'} \\
S' & & T' \ar[ll]_{g' = g \circ \pi}
}
}
$$
を考える。すべての正方形はカルテシアンであり、$g$ は準コンパクトかつ準分離、
$\pi$ は整かつ全射であるとする。$\mathcal{F}$ を $T_\etale$ 上のアーベル層で
$n$ により零化されるものとし、$\mathcal{F}' = \pi^{-1}\mathcal{F}$ とおく。
基底変換写像
$$
(f')^{-1}R^qg'_*\mathcal{F}' \longrightarrow R^qh'_*(e')^{-1}\mathcal{F}'
$$
が同型ならば、基底変換写像
$(f')^{-1}R^qg_*\mathcal{F} \to R^qh_*e^{-1}\mathcal{F}$
も同型である。
\end{lemma}

\begin{proof}
$\mathcal{F} \to \pi_*\pi^{-1}\mathcal{F}'$ は $\pi$ が全射なので単射である
（茎上で確認せよ）。したがって補題 \ref{etale-cohomology-lemma-base-change-q-injective} により、
基底変換写像
$$
(f')^{-1}R^qg_*\pi_*\mathcal{F}'
\longrightarrow
R^qh_*e^{-1}\pi_*\mathcal{F}'
$$
が同型であることを示せば十分である。これは仮定から従う。実際、
$R^qg_*\pi_*\mathcal{F}' = R^qg'_*\mathcal{F}'$、
$e^{-1}\pi_*\mathcal{F}' =\pi'_*(e')^{-1}\mathcal{F}'$、および
$R^qh_*\pi'_*(e')^{-1}\mathcal{F}' = R^qh'_*(e')^{-1}\mathcal{F}'$ である。
これらは補題
\ref{etale-cohomology-lemma-integral-pushforward-commutes-with-base-change}、補題
\ref{etale-cohomology-lemma-what-integral}、および相対 Leray スペクトル系列
（Cohomology on Sites の補題 \ref{sites-cohomology-lemma-relative-Leray}）から従う。
\end{proof}

\begin{lemma}
\label{etale-cohomology-lemma-base-change-q-integral-bottom}
$f : X \to S$ および $n$ は注意 \ref{etale-cohomology-remark-base-change-holds} のとおりとし、
ある $q \geq 1$ に対して $BC(f, n, q - 1)$ が成り立つと仮定する。可換図式
$$
\vcenter{
\xymatrix{
X \ar[d]_f &
X' \ar[d]_{f'} \ar[l] &
X'' \ar[l]^{\pi'} \ar[d]_{f''} &
Y \ar[l]^{h'} \ar[d]^e \\
S &
S' \ar[l] &
S'' \ar[l]_\pi &
T \ar[l]_{g'}
}
}
\quad\text{および}\quad
\vcenter{
\xymatrix{
X' \ar[d]_{f'} & & Y \ar[ll]^{h = h' \circ \pi'} \ar[d]^e \\
S' & & T \ar[ll]_{g = g' \circ \pi}
}
}
$$
を考える。すべての正方形はカルテシアンであり、$g'$ は準コンパクトかつ準分離、
$\pi$ は整であるとする。$\mathcal{F}$ を $T_\etale$ 上のアーベル層で
$n$ により零化されるものとする。基底変換写像
$$
(f')^{-1}R^qg_*\mathcal{F} \longrightarrow R^qh_*e^{-1}\mathcal{F}
$$
が同型ならば、基底変換写像
$(f'')^{-1}R^qg'_*\mathcal{F} \to R^qh'_*e^{-1}\mathcal{F}$
も同型である。
\end{lemma}

\begin{proof}
写像 $\pi$ と $\pi'$ は整なので $R\pi_* = \pi_*$ および
$R\pi'_* = \pi'_*$ である。補題 \ref{etale-cohomology-lemma-what-integral} を参照されたい。
また $(f')^{-1}\pi_* = \pi'_*(f'')^{-1}$ である。したがって
$\pi'_*(f'')^{-1}R^qg'_*\mathcal{F} = (f')^{-1}R^qg_*\mathcal{F}$
および
$\pi'_*R^qh'_*e^{-1}\mathcal{F} = R^qh_*e^{-1}\mathcal{F}$.
ゆえに仮定は、問題の写像に関手 $\pi'_*$ を適用すると同型になることを意味する。
したがって補題 \ref{etale-cohomology-lemma-what-integral} により、その写像自身も同型である。
\end{proof}

\begin{lemma}
\label{etale-cohomology-lemma-formal-argument}
$T$ を準コンパクトかつ準分離なスキームとする。
$P$ を $T$ 上の準コンパクトかつ準分離なスキームに関する性質とする。次を仮定する。
\begin{enumerate}
\item $T'' \to T'$ が $T$ 上の準コンパクトかつ準分離なスキームの厚化ならば、
$P(T'')$ であることと $P(T')$ であることは同値である。
\item $T' = \lim T_i$ が、$T$ 上の準コンパクトかつ準分離なスキームからなる、
アフィンな遷移射をもつ逆系の極限であり、性質 $P(T_i)$ がすべての $i$ に対して成り立つならば、
$P(T')$ が成り立つ。
\item $Z \subset T'$ が閉部分スキームで、その準コンパクトな補集合が
$V \subset T'$ であり、$P(T')$ が成り立つならば、$P(V)$ または $P(Z)$ が成り立つ。
\end{enumerate}
このとき $P(T)$ ならば、$P(\Spec(K))$ が成り立つような射
$\Spec(K) \to T$ が存在する。ここで $K$ は体である。
\end{lemma}

\begin{proof}
集合 $\mathfrak T$ を考える。その元は、閉部分スキーム $T' \subset T$ であって
$P(T')$ を満たすものである。
仮定 (2) によりこの集合は極小元をもち、それを $T'$ とする。仮定 (1) により $T'$ は被約である。
$\eta \in T'$ を $T'$ の既約成分の一般点とする。このとき $\eta = \Spec(K)$ であり、
$K$ はある体である。さらに $\eta = \lim V$ であり、極限はアフィン開部分スキーム
$V \subset T'$ で $\eta$ を含むもの全体にわたる。
仮定 (3) と $T'$ の極小性により、性質 $P(V)$ はこれらすべての $V$ に対して成り立つ。
したがって (2) により $P(\eta)$ であり、証明が完了する。
\end{proof}

\begin{lemma}
\label{etale-cohomology-lemma-base-change-does-not-hold-pre}
$f : X \to S$ および $n$ は注意 \ref{etale-cohomology-remark-base-change-holds} のとおりとし、
ある $q \geq 1$ に対して $BC(f, n, q - 1)$ は真であるが
$BC(f, n, q)$ は真でないと仮定する。このとき可換図式
$$
\xymatrix{
X \ar[d]_f & X' \ar[d]_{f'} \ar[l] & Y \ar[l]^h \ar[d]^e \\
S & S' \ar[l] & \Spec(K) \ar[l]_g
}
$$
が存在する。ここで $X' = X \times_S S'$、$Y = X' \times_{S'} \Spec(K)$ であり、
$K$ は体で、$\mathcal{F}$ は $\Spec(K)$ 上のアーベル層で $n$ により零化され、
$(f')^{-1}R^qg_*\mathcal{F} \to R^qh_*e^{-1}\mathcal{F}$
は同型でない。
\end{lemma}

\begin{proof}
可換図式
$$
\xymatrix{
X \ar[d]_f & X' \ar[l] \ar[d]_{f'} & Y \ar[l]^h \ar[d]^e \\
S & S' \ar[l] & T \ar[l]_g
}
$$
を選ぶ。ここで $X' = X \times_S S'$ および $Y = X' \times_{S'} T$ であり、
$g$ は準コンパクトかつ準分離で、$\mathcal{F}$ は $T_\etale$ 上のアーベル層で
$n$ により零化され、基底変換写像
$(f')^{-1}R^qg_*\mathcal{F} \to R^qh_*e^{-1}\mathcal{F}$
は同型でない。もちろん $S'$ を $S'$ のアフィン開集合で置き換えてよく、実際そうする。
すると $T$ は準コンパクトかつ準分離になる。補題
\ref{etale-cohomology-lemma-base-change-q-injective} により
$(f')^{-1}R^qg_*\mathcal{F} \to R^qh_*e^{-1}\mathcal{F}$
は単射である。幾何学的点 $\overline{x}$ を $X'$ 上に選び、元 $\xi$ を
$(R^qh_*q^{-1}\mathcal{F})_{\overline{x}}$ から選ぶ。ただしこれは写像
$((f')^{-1}R^qg_*\mathcal{F})_{\overline{x}} \to
(R^qh_*e^{-1}\mathcal{F})_{\overline{x}}$ の像に属さないものとする。

\medskip\noindent
射 $\pi : T' \to T$ を考え、$T'$ は準コンパクトかつ準分離であるとし、
$\mathcal{F}' = \pi^{-1}\mathcal{F}$ と書く。
$\pi' : Y' = Y \times_T T' \to Y$ を $\pi$ の基底変換と書き、
$e' : Y' \to T'$ を $e$ の基底変換と書く。図式は
$$
\vcenter{
\xymatrix{
X' \ar[d]_{f'} & Y \ar[l]^h \ar[d]^e & Y' \ar[l]^{\pi'} \ar[d]^{e'} \\
S' & T \ar[l]_g & T' \ar[l]_\pi
}
}
\quad\text{および}\quad
\vcenter{
\xymatrix{
X' \ar[d]_{f'} & & Y' \ar[ll]^{h' = h \circ \pi'} \ar[d]^{e'} \\
S' & & T' \ar[ll]_{g' = g \circ \pi}
}
}
$$
である。引戻し写像を用いると、アーベル層の標準的可換図式
$$
\xymatrix{
(f')^{-1}R^qg_*\mathcal{F} \ar[r] \ar[d] &
(f')^{-1}R^qg'_*\mathcal{F}' \ar[d] \\
R^qh_*e^{-1}\mathcal{F} \ar[r] &
R^qh'_*(e')^{-1}\mathcal{F}'
}
$$
を $X'$ 上に得る。$P(T')$ を次の性質とする。
\begin{itemize}
\item $\xi'$ を $\xi$ の $(Rh'_*(e')^{-1}\mathcal{F}')_{\overline{x}}$ における像とする。
これは写像
$(f^{-1}R^qg'_*\mathcal{F}')_{\overline{x}} \to
(R^qh'_*(e')^{-1}\mathcal{F}')_{\overline{x}}$ の像には属さない。
\end{itemize}
補題 \ref{etale-cohomology-lemma-formal-argument} の仮定 (1)、(2)、(3) が $P$ に対して成り立つと主張する。
これにより本補題が証明される。

\medskip\noindent
補題 \ref{etale-cohomology-lemma-formal-argument} の条件 (1) は $P$ に対して成り立つ。
実際、スキームとその厚化のエタール位相は同じである。命題
\ref{etale-cohomology-proposition-topological-invariance} を参照されたい。

\medskip\noindent
$I$ を有向集合とし、$T_i$ を $I$ 上の逆系とする。この逆系は、$T$ 上の
準コンパクトかつ準分離なスキームからなり、アフィンな遷移射をもつとする。
$T' = \lim T_i$ とおく。$\mathcal{F}'$ と $\mathcal{F}_i$ をそれぞれ
$\mathcal{F}$ の $T'$ と $T_i$ への引戻しと書く。図式
$$
\vcenter{
\xymatrix{
X \ar[d]_{f'} & Y \ar[l]^h \ar[d]^e & Y_i \ar[l]^{\pi_i'} \ar[d]^{e_i} \\
S & T \ar[l]_g & T_i \ar[l]_{\pi_i}
}
}
\quad\text{および}\quad
\vcenter{
\xymatrix{
X \ar[d]_{f'} & & Y_i \ar[ll]^{h_i = h \circ \pi_i'} \ar[d]^{e_i} \\
S & & T_i \ar[ll]_{g_i = g \circ \pi_i}
}
}
$$
を前段落と同様に考える。$\mathcal{F}'$ は $T'$ 上で、$\mathcal{F}_i$ の
$T'$ への引戻しの余極限であり、$(e')^{-1}\mathcal{F}'$ は
$e_i^{-1}\mathcal{F}_i$ の $Y'$ への引戻しの余極限であることは明らかである。
補題 \ref{etale-cohomology-lemma-relative-colimit-general} により
$$
R^qh'_*(e')^{-1}\mathcal{F}' = \colim R^qh_{i, *}e_i^{-1}\mathcal{F}_i
\quad\text{および}\quad
(f')^{-1}R^qg'_*\mathcal{F}' = \colim (f')^{-1}R^qg_{i, *}\mathcal{F}_i
$$
したがって、性質 $P(T_i)$ がすべての $i$ に対して真ならば $P(T')$ が成り立つ。
ゆえに補題 \ref{etale-cohomology-lemma-formal-argument} の条件 (2) は $P$ に対して成り立つ。

\medskip\noindent
最も興味深いのは補題 \ref{etale-cohomology-lemma-formal-argument} の条件 (3) である。
$T'$ を $T$ 上の準コンパクトかつ準分離なスキームとし、$P(T')$ が真であると仮定する。
$Z \subset T'$ を閉部分スキームとし、その補集合 $V \subset T'$ は準コンパクトであるとする。
図式
$$
\xymatrix{
Y' \times_{T'} Z \ar[d]_{e_Z} \ar[r]_{i'} &
Y' \ar[d]_{e'} &
Y' \times_{T'} V \ar[l]^{j'} \ar[d]^{e_V} \\
Z \ar[r]^i &
T' &
V \ar[l]_j
}
$$
を考える。単射 $j^{-1}\mathcal{F}' \to \mathcal{J}$ を選ぶ。ここで $\mathcal{J}$ は
$\mathbf{Z}/n\mathbf{Z}$-加群の単射層であり、$V$ 上にある。茎を見ると、写像
$$
\mathcal{F}' \to \mathcal{G} = j_*\mathcal{J} \oplus i_*i^{-1}\mathcal{F}'
$$
が単射であることが分かる。したがって $\xi'$ は次の非零元へ写る。
\begin{align*}
& \Coker\left(
((f')^{-1}R^qg'_*\mathcal{G})_{\overline{x}}
\to
(R^qh'_*(e')^{-1}\mathcal{G})_{\overline{x}}
\right) = \\
&
\Coker\left(
((f')^{-1}R^qg'_*j_*\mathcal{J})_{\overline{x}}
\to
(R^qh'_*(e')^{-1}j_*\mathcal{J})_{\overline{x}}
\right) \oplus \\
& \Coker\left(
((f')^{-1}R^qg'_*i_*i^{-1}\mathcal{F}')_{\overline{x}}
\to
(R^qh'_*(e')^{-1}i_*i^{-1}\mathcal{F}')_{\overline{x}}
\right)
\end{align*}
これは補題 \ref{etale-cohomology-lemma-base-change-q-injective} の (2) による。
$\xi'$ の第二直和因子への像が零でないならば、
$$
(f')^{-1}R^qg'_*i_*i^{-1}\mathcal{F}' =
(f')^{-1}R^q(g' \circ i)_*i^{-1}\mathcal{F}'
$$
を用いる（命題 \ref{etale-cohomology-proposition-finite-higher-direct-image-zero} により
$Ri_* = i_*$ だからである）。また
$$
R^qh'_*(e')^{-1}i_*i^{-1}\mathcal{F} =
R^qh'_*i'_*e_Z^{-1}i^{-1}\mathcal{F} =
R^q(h' \circ i')_*e_Z^{-1}i^{-1}\mathcal{F}'
$$
を用いる（第一の等号は補題
\ref{etale-cohomology-lemma-finite-pushforward-commutes-with-base-change} により、第二の等号は命題
\ref{etale-cohomology-proposition-finite-higher-direct-image-zero} により $Ri'_* = i'_*$ であることによる）。
これにより $P(Z)$ を得る。
最後に、$\xi'$ の第一直和因子への像が零でないと仮定する。このとき
$$
(e')^{-1}j_*\mathcal{J} = j'_*e_V^{-1}\mathcal{J}
\quad\text{および}\quad
R^aj'_*e_V^{-1}\mathcal{J} = 0, \quad a = 1, \ldots, q - 1
$$
である。これは $BC(f, n, q - 1)$ を図式
$$
\xymatrix{
X \ar[d]_f & Y' \ar[l] \ar[d]_{e'} & Y \ar[l]^{j'} \ar[d]^{e_V} \\
S & T' \ar[l] & V \ar[l]_j
}
$$
に適用し、$\mathcal{J}$ が単射的であることから従う。
$h' \circ j'$ に対する相対 Leray スペクトル系列
（Cohomology on Sites の補題 \ref{sites-cohomology-lemma-relative-Leray}）により
$$
R^qh'_*(e')^{-1}j_*\mathcal{J} =
R^qh'_*j'_*e_V^{-1}\mathcal{J}
\longrightarrow
R^q(h' \circ j')_* e_V^{-1}\mathcal{J}
$$
は単射である。したがって $\xi$ は
$(R^q(h' \circ j')_* e_V^{-1}\mathcal{J})_{\overline{x}}$ の非零元へ写る。
単射 $j^{-1}\mathcal{F}' \to \mathcal{J}$ に補題
\ref{etale-cohomology-lemma-base-change-q-injective} の (3) を適用すると、$P(V)$ が成り立つと結論できる。
\end{proof}

\begin{lemma}
\label{etale-cohomology-lemma-base-change-does-not-hold}
$f : X \to S$ および $n$ は注意 \ref{etale-cohomology-remark-base-change-holds} のとおりとし、
ある $q \geq 1$ に対して $BC(f, n, q - 1)$ は真であるが
$BC(f, n, q)$ は真でないと仮定する。このとき可換図式
$$
\xymatrix{
X \ar[d]_f & X' \ar[d] \ar[l] & Y \ar[l]^h \ar[d] \\
S & S' \ar[l] & \Spec(K) \ar[l]
}
$$
が存在し、両方の正方形はカルテシアンで、次が成り立つ。
\begin{enumerate}
\item $S'$ はアフィン、整、正規で、代数閉な関数体をもつ。
\item $K$ は代数閉で、$\Spec(K) \to S'$ は優勢である
（言い換えると、$K$ は $S'$ の関数体の拡大である）。
\end{enumerate}
さらに、ある整数 $d | n$ が存在し、$R^qh_*(\mathbf{Z}/d\mathbf{Z})$ は非零である。
\end{lemma}

\noindent
逆に、補題における $R^qh_*(\mathbf{Z}/d\mathbf{Z})$ の非消滅は
$BC(f, n, q)$ が真でないことを導く。実際、補題
\ref{etale-cohomology-lemma-Rf-star-zero-normal-with-alg-closed-function-field} は
$R^q(\Spec(K) \to S')_*\mathbf{Z}/d\mathbf{Z} = 0$ を示す。

\begin{proof}
まず補題 \ref{etale-cohomology-lemma-base-change-does-not-hold-pre} にあるような図式と
$\mathcal{F}$ を選ぶ。$S'$ はアフィンであると仮定してよく、実際そうする
（これは明らかだが、疑わしければ補題の証明を参照されたい）。補題
\ref{etale-cohomology-lemma-base-change-q-integral-top} により、$K$ は代数閉であると仮定してよい。
すると $\mathcal{F}$ は $\mathbf{Z}/n\mathbf{Z}$-加群に対応する。そのような加群は、
$\mathbf{Z}/d\mathbf{Z}$ のコピーの直和であり、ここで $d | n$ は種々に動く。したがって
$\mathcal{F}$ は値 $\mathbf{Z}/d\mathbf{Z}$ をもつ定数層であると仮定してよい。
補題 \ref{etale-cohomology-lemma-base-change-q-integral-bottom} により、$S'$ を
$S'$ の $\Spec(K)$ における正規化で置き換えてよい。これで証明が完了する。
\end{proof}










\section{滑らかな基底変換}
\label{etale-cohomology-section-smooth-base-change}

\noindent
本節では，滑らかな基底変換定理を証明する．

\begin{lemma}
\label{etale-cohomology-lemma-smooth-base-change-fields}
$K/k$ を体の拡大とする．$X$ を $k$ 上の滑らかなアフィン曲線で，
有理点 $x \in X(k)$ をもつものとする．$\mathcal{F}$ を $\Spec(K)$ 上の
アーベル層で，整数 $n$ によって零化され，この整数が $k$ において
可逆であるものとする．
$q > 0$ とし，
$$
\xi \in H^q(X_K, (X_K \to \Spec(K))^{-1}\mathcal{F})
$$
とする．次を満たすものが存在する：
\begin{enumerate}
\item 有限拡大 $K'/K$ および $k'/k$ で，$k' \subset K'$ であるもの；
\item 有限 \'etale Galois 被覆 $Z \to X_{k'}$ とその群 $G$．
\end{enumerate}
ここで $G$ の位数は $n$ のある冪を割り，$Z \to X_{k'}$ は
$x_{k'}$ 上で分裂し，$\xi$ は
$H^q(Z_{K'}, (Z_{K'} \to \Spec(K))^{-1}\mathcal{F})$ において消える．
\end{lemma}

\begin{proof}
$q > 1$ のとき，$\xi$ は
$H^q(X_{\overline{K}}, (X_{\overline{K}} \to \Spec(K))^{-1}\mathcal{F})$
において消えることが分かっている（定理
\ref{etale-cohomology-theorem-vanishing-affine-curves}）．補題
\ref{etale-cohomology-lemma-directed-colimit-cohomology} により，これは，有限拡大 $K'/K$ が
存在して，$\xi$ が
$H^q(X_{K'}, (X_{K'} \to \Spec(K))^{-1}\mathcal{F})$ において消えることを
意味する．したがってこの場合は
$k' = k$ および $Z = X$ と取ればよい．

\medskip\noindent
$q = 1$ と仮定する．$\mathcal{F}$ は離散加群 $M$ に対応し，この加群は連続な
$\text{Gal}_K$-作用をもつことを思い出そう；補題
\ref{etale-cohomology-lemma-equivalence-abelian-sheaves-point} を参照されたい．
$M$ は $n$-捩れなので，有限な $\text{Gal}_K$-安定部分群の合併である．
したがって補題 \ref{etale-cohomology-lemma-colimit} により，$M$ が $n$ で零化される
有限アーベル群である場合に帰着する．$K$ を有限拡大で置き換えれば，
$\text{Gal}_K$ の $M$ への作用が自明であると仮定してよい．ゆえに，
$\mathcal{F} = \underline{M}$ は有限アーベル群 $M$ を値とし，この群が
$n$ で零化される定数層であると仮定してよい．

\medskip\noindent
$M$ は巡回群の直和として書ける．Galois 群の位数が $k$ において
可逆である任意の二つの有限 \'etale Galois 被覆は，Galois 群の位数が
$k$ において可逆である第三の被覆によって支配される
（Fundamental Groups，節 \ref{pione-section-finite-etale-under-galois}）．
したがって，$M = \mathbf{Z}/d\mathbf{Z}$ かつ $d | n$ である場合に
補題を証明すれば十分である．

\medskip\noindent
$M = \mathbf{Z}/d\mathbf{Z}$ かつ $d | n$ と仮定する．この場合，
$\overline{\xi} = \xi|_{X_{\overline{K}}}$ は
$$
H^1(X_{\overline{k}}, \mathbf{Z}/d\mathbf{Z}) =
H^1(X_{\overline{K}}, \mathbf{Z}/d\mathbf{Z})
$$
の元である；定理 \ref{etale-cohomology-theorem-vanishing-affine-curves} を参照されたい．
この群は $\mathbf{Z}/d\mathbf{Z}$-トーサーを分類する；Cohomology on Sites，
補題 \ref{sites-cohomology-lemma-torsors-h1} を参照されたい．
$\overline{\xi}$ に対応するトーサー（$X_{\overline{k}, \etale}$ 上の層とみなす）は，
有限 \'etale 射 $T \to X_{\overline{k}}$ とその上の
$\mathbf{Z}/d\mathbf{Z}$ の作用を与え，この作用は $T$ の
$x_{\overline{k}}$ 上のファイバーに推移的に作用する；補題
\ref{etale-cohomology-lemma-characterize-finite-locally-constant} を参照されたい．
連結成分 $T' \subset T$ を選ぶ（$\overline{\xi}$ の位数が $d$ ならば，
$T$ はすでに連結である）．すると $T' \to X_{\overline{k}}$ は，
Galois 群が部分群 $G \subset \mathbf{Z}/d\mathbf{Z}$ である有限 \'etale
Galois 被覆である（細部は省略する）．さらに元 $\overline{\xi}$ は写像
$H^1(X_{\overline{k}}, \mathbf{Z}/d\mathbf{Z}) \to
H^1(T', \mathbf{Z}/d\mathbf{Z})$ の下で零に写る．これは $T$ の
定義的性質の一つだからである．

\medskip\noindent
次に極限論法を用いると，有限拡大 $k'/k$ で $\overline{k}$ に含まれ，
$T' \to X_{\overline{k}}$ が有限 \'etale Galois 被覆
$Z \to X_{k'}$ に降下し，その群が $G$ であるものを選べる．Limits の補題
\ref{limits-lemma-descend-finite-presentation},
\ref{limits-lemma-descend-finite-finite-presentation},
\ref{limits-lemma-descend-etale} を参照されたい．$k'$ を拡大すれば，
$Z$ が $x_{k'}$ 上で分裂すると仮定してよい．構成により，$\xi$ の
$H^1(Z_{\overline{K}}, \mathbf{Z}/d\mathbf{Z})$ における像は零である．
したがって補題 \ref{etale-cohomology-lemma-directed-colimit-cohomology} により，有限部分拡大
$\overline{K}/K'/K$ で $k'$ を含み，$\xi$ が
$H^1(Z_{K'}, \mathbf{Z}/d\mathbf{Z})$ において消えるものを見つけられる．
これで証明が完了する．
\end{proof}

\begin{theorem}[滑らかな基底変換]
\label{etale-cohomology-theorem-smooth-base-change}
スキームのファイバー積図式
$$
\xymatrix{
X \ar[d]_f & Y \ar[l]^h \ar[d]^e \\
S & T \ar[l]_g
}
$$
を考える．$f$ は滑らかで，$g$ は準コンパクトかつ準分離的であるとする．
このとき
$$
f^{-1}R^qg_*\mathcal{F} = R^qh_*e^{-1}\mathcal{F}
$$
が，任意の $q$，および任意のアーベル層 $\mathcal{F}$ で $T_\etale$ 上のもののうち，
幾何学点におけるすべての茎が $S$ 上可逆な位数の捩れであるものに対して
成り立つ．
\end{theorem}

\begin{proof}[滑らかな基底変換の第一証明]
この証明は非常に長いが，以下に与える第二証明よりも直接的である
（用いる一般論が少ない）．

\medskip\noindent
定理は $X_\etale$ 上局所的である．より正確には，\'etale 射
$U \to X$ があり，$U \to S$ が $U \to V \to S$ と分解し，
$V \to S$ が \'etale であると仮定する．このときファイバー積正方形
$$
\xymatrix{
U \ar[d]_{f'} & U \times_X Y \ar[l]^{h'} \ar[d]^{e'} \\
V & V \times_S T \ar[l]_{g'}
}
$$
を考えられる．$\mathcal{F}' = \mathcal{F}|_{V \times_S T}$ とおけば，
$f^{-1}R^qg_*\mathcal{F}|_U = (f')^{-1}R^qg'_*\mathcal{F}'$ および
$R^qh_*e^{-1}\mathcal{F}|_U = R^qh'_*(e')^{-1}\mathcal{F}'$ である
（これは局所化とサイトの射との両立性から従う；Sites の補題
\ref{sites-lemma-localize-morphism-strong} および Cohomology on Sites の補題
\ref{sites-cohomology-lemma-restrict-direct-image-open} を参照されたい）．
したがって，$X$ の \'etale 被覆をなす $U \to X$ と上のような分解
$U \to V \to S$ であって，$f'$, $h'$, $g'$, $e'$ をもつ図式に対して
定理が成り立つものを構成すれば十分である．

\medskip\noindent
滑らかな射の局所構造により（Morphisms，補題
\ref{morphisms-lemma-smooth-etale-over-affine-space} を参照），
$X$ と $S$ はアフィンであり，$X \to S$ は \'etale 射
$X \to \mathbf{A}^d_S$ を経由して分解すると仮定してよい．
ファイバー積図式の塔
$$
\xymatrix{
W \ar[d]_i & Z \ar[l]^j \ar[d]^k \\
X \ar[d]_f & Y \ar[l]^h \ar[d]^e \\
S & T \ar[l]_g
}
$$
があり，下側と上側の正方形に対して定理が成り立つならば，外側の
長方形に対しても定理が成り立つ；これは形式的である．$X \to S$ を合成
$$
X \to \mathbf{A}^{d - 1}_S \to \mathbf{A}^{d - 2}_S \to \ldots \to
\mathbf{A}^1_S \to S
$$
として書けば，$X$ と $S$ がアフィンで，$X \to S$ の相対次元が $1$
である場合に定理を証明すれば十分であると分かる．

\medskip\noindent
$n \geq 1$ が $S$ 上可逆であるごとに，$\mathcal{F}[n]$ を，$\mathcal{F}$ の
$n$ で零化される切断からなる部分層とする．すると
$\mathcal{F} = \colim \mathcal{F}[n]$ である．これは $\mathcal{F}$ の茎に
関する仮定による．関手 $e^{-1}$ と $f^{-1}$ は
左随伴なので余極限と可換である．関手 $R^qh_*$ と $R^qg_*$ は補題
\ref{etale-cohomology-lemma-relative-colimit} によりフィルター余極限と可換である．
したがって $\mathcal{F}[n]$ に対して定理を証明すれば十分である．
以下では整数 $n$ を固定し，$\mathbf{Z}/n\mathbf{Z}$-加群の層を扱い，
$S$ は $\Spec(\mathbf{Z}[1/n])$ 上のスキームであると仮定する．

\medskip\noindent
次に，$T$ がアフィンである場合に帰着する．$g$ は準コンパクトかつ
準分離的であり，$S$ はアフィンなので，スキーム $T$ は準コンパクトかつ
準分離的である．したがって Cohomology of Schemes の補題
\ref{coherent-lemma-induction-principle} の帰納原理を用いられる．ゆえに，
$T = W \cup W'$ が開被覆であり，次の各正方形に対して定理が成り立つならば，
$$
\xymatrix{
X \ar[d] & e^{-1}(W) \ar[l]^i \ar[d] \\
S & W \ar[l]_a
}
\quad
\xymatrix{
X \ar[d] & e^{-1}(W') \ar[l]^j \ar[d] \\
S & W' \ar[l]_b
}
\quad
\xymatrix{
X \ar[d] & e^{-1}(W \cap W') \ar[l]^-k \ar[d] \\
S & W \cap W' \ar[l]_c
}
$$
もとの図式に対しても定理が成り立つことを示せば十分である．このことを見るため，
図式
$$
\xymatrix{
f^{-1}R^{q - 1}c_*\mathcal{F}|_{W \cap W'} \ar[d]^{\cong} \ar[r] &
f^{-1}R^qg_*\mathcal{F} \ar[d] \ar[r] &
f^{-1}R^qa_*\mathcal{F}|_W \oplus
f^{-1}R^qb_*\mathcal{F}|_{W'} \ar[d]_{\cong} \\
R^qk_*e^{-1}\mathcal{F}|_{e^{-1}(W \cap W')} \ar[r] &
R^qh_*e^{-1}\mathcal{F} \ar[r] &
R^qi_*e^{-1}\mathcal{F}|_{e^{-1}(W)} \oplus
R^qj_*e^{-1}\mathcal{F}|_{e^{-1}(W')}
}
$$
を考える．その各行は補題 \ref{etale-cohomology-lemma-relative-mayer-vietoris} の長完全列である．
したがって $5$-補題から所望の結論が得られる．

\medskip\noindent
まとめると，$S$, $X$, $T$, $Y$ はアフィン，$\mathcal{F}$ は $n$-捩れ，
$X \to S$ は相対次元 $1$ の滑らかな射，かつ $S$ は
$\mathbf{Z}[1/n]$ 上のスキームであると仮定してよい．$q$ に関する帰納法で
定理を証明する．初期の場合 $q = 0$ は補題
\ref{etale-cohomology-lemma-fppf-reduced-fibres-base-change-f-star} で扱われている．
$q > 0$ とし，より小さいすべての次数で定理が成り立つと仮定する．短完全列
$0 \to \mathcal{F} \to \mathcal{I} \to \mathcal{Q} \to 0$
で，$\mathcal{I}$ が $\mathbf{Z}/n\mathbf{Z}$-加群の単射的層であるものを選ぶ．
誘導される図式
$$
\xymatrix{
f^{-1}R^{q - 1}g_*\mathcal{I} \ar[d]_{\cong} \ar[r] &
f^{-1}R^{q - 1}g_*\mathcal{Q} \ar[d]_{\cong} \ar[r] &
f^{-1}R^qg_*\mathcal{F} \ar[d] \ar[r] &
0 \ar[d] \\
R^{q - 1}h_*e^{-1}\mathcal{I} \ar[r] &
R^{q - 1}h_*e^{-1}\mathcal{Q} \ar[r] &
R^qh_*e^{-1}\mathcal{F} \ar[r] &
R^qh_*e^{-1}\mathcal{I}
}
$$
を考える．各行は完全である．$\mathcal{I}$ が単射的なので，右上隅は零である．
帰納法の仮定により，左側の二つの垂直射は同型である．したがって
$R^qh_*e^{-1}\mathcal{I} = 0$ を証明すれば十分である．

\medskip\noindent
$S = \Spec(A)$ および $T = \Spec(B)$ と書き，射 $T \to S$ は環準同型
$A \to B$ で与えられるとする．$A \to B = \colim_{i \in I} (A_i \to B_i)$
を，$\mathbf{Z}[1/n]$ 上有限型な環の準同型のフィルター余極限として書ける
（Algebra，補題 \ref{algebra-lemma-limit-no-condition} を参照）．$i \in I$ に対し
$S_i = \Spec(A_i)$ および $T_i = \Spec(B_i)$ とおく．十分大きい $i$ に対し，
滑らかな射 $X_i \to S_i$ で，相対次元が $1$ であり，
$X = X_i \times_{S_i} S$ となるものを見つけられる；Limits の補題
\ref{limits-lemma-descend-finite-presentation},
\ref{limits-lemma-descend-smooth},
\ref{limits-lemma-descend-dimension-d} を参照されたい．
$Y_i = X_i \times_{S_i} T_i$ とおくと，正方形
$$
\xymatrix{
X_i \ar[d]_{f_i} & Y_i \ar[l]^{h_i} \ar[d]^{e_i} \\
S_i & T_i \ar[l]_{g_i}
}
$$
を得る．$\mathcal{I}_i = (T \to T_i)_*\mathcal{I}$ は
$\mathbf{Z}/n\mathbf{Z}$-加群の $T_i$ 上の単射的層である；Cohomology on Sites の補題
\ref{sites-cohomology-lemma-pushforward-injective-flat} を参照されたい．補題
\ref{etale-cohomology-lemma-linus-hamann} により
$\mathcal{I} = \colim (T \to T_i)^{-1}\mathcal{I}_i$ である．$e$ により引き戻すと
$e^{-1}\mathcal{I} = \colim (Y \to Y_i)^{-1}e_i^{-1}\mathcal{I}_i$ を得る．
$Y_i \to X_i$ という射の系で，極限が $Y \to X$ であるものに補題
\ref{etale-cohomology-lemma-relative-colimit-general} を適用すると，
$$
R^qh_*e^{-1}\mathcal{I} =
\colim (X \to X_i)^{-1} R^qh_{i, *} e_i^{-1}\mathcal{I}_i
$$
を得る．これにより，$T$ と $S$ が $\mathbf{Z}[1/n]$ 上有限型な
アフィンスキームである場合に帰着する．

\medskip\noindent
まとめると，$q \geq 1$ は次数 $< q$ で定理が成り立つ整数，$S$ と $T$ は
$\mathbf{Z}[1/n]$ 上有限型なアフィンスキーム，$X \to S$ は相対次元 $1$ の
滑らかな射で $X$ はアフィン，そして $\mathcal{I}$ は
$\mathbf{Z}/n\mathbf{Z}$-加群の単射的層であると仮定してよく，示すべきことは
$R^qh_*e^{-1}\mathcal{I} = 0$ である．これを $\dim(T)$ に関する帰納法で示す．

\medskip\noindent
初期の場合は $T = \emptyset$，すなわち $\dim(T) < 0$ である．
これを好まなければ，$\dim(T) = 0$ の場合を初期の場合としてもよい．
この場合 $T \to S$ は有限である（実際，終域が Jacobson なので
$T \to \Spec(\mathbf{Z}[1/n])$ さえ有限である；細部は省略する）．
したがって $h$ も有限であり，その高次順像は消える
（以下の参考文献を参照されたい）．

\medskip\noindent
$\dim(T) = d \geq 0$ と仮定し，$T$ の次元がより小さいすべての状況で
結果を知っているとする．アフィンな $U$ で $X$ 上 \'etale なものと，
切断 $\xi$，すなわち $R^qh_*q^{-1}\mathcal{I}$ の $U$ 上の切断を選ぶ．
$\xi$ が零であることを
示さなければならない．もちろん，$X$ を $U$ で（それに応じて $Y$ を
$U \times_X Y$ で）置き換え，$\xi \in H^0(X, R^qh_*e^{-1}\mathcal{I})$
と仮定してよい．さらに，$R^qh_*e^{-1}\mathcal{I}$ は層なので，$\xi$ が
$X$ 上局所的に零であることを示せば十分である．したがって $X$ を
ある \'etale 被覆の各要素で置き換えてよい．特に，補題
\ref{etale-cohomology-lemma-higher-direct-images} を用いると，$\xi$ はある元
$\tilde \xi \in H^q(Y, e^{-1}\mathcal{I})$ の像であると仮定してよい．
$\tilde \xi$ を用いていえば，我々の課題は，$\tilde \xi$ が
$H^q(U_i \times_X Y, e^{-1}\mathcal{I})$ において，ある \'etale 被覆
$\{U_i \to X\}$ に対して消えることを示すことである．

\medskip\noindent
More on Morphisms の補題 \ref{more-morphisms-lemma-cover-smooth-by-special}
により，$X \to S$ が $X \to V \to S$ と分解し，$V \to S$ は \'etale，
$X \to V$ は相対次元 $1$ のアフィンスキームの滑らかな射で，切断をもち，
幾何学的に連結なファイバーをもつと仮定してよい．例えば More on Algebra の補題
\ref{more-algebra-lemma-dimension-etale-extension} により，
$\dim(V \times_S T) \leq \dim(T) = d$ であることに注意する．したがって，
$S$ を $V$ で，$T$ を $V \times_S T$ で置き換えてよい（証明の第一段落の
議論とまったく同様である）．ゆえに $X \to S$ は相対次元 $1$ の滑らかな射で，
幾何学的に連結なファイバーをもち，切断 $\sigma : S \to X$ をもつと
仮定してよい．

\medskip\noindent
$\pi : T' \to T$ を有限全射とする．以下では
$\dim(T') \leq \dim(T) = d$ を用いる；Algebra，補題
\ref{algebra-lemma-integral-dim-up} を参照されたい．単射
$\pi^{-1}\mathcal{I} \to \mathcal{I}'$ を，$\mathbf{Z}/n\mathbf{Z}$-加群の
単射的層へのものとして選ぶ．
すると $\mathcal{I} \to \pi_*\mathcal{I}'$ は単射であり，したがって分裂をもつ
（$\mathcal{I}$ は $\mathbf{Z}/n\mathbf{Z}$-加群の単射的層だからである）．
$\pi' : Y' = Y \times_T T' \to Y$ を $\pi$ の基底変換，
$e' : Y' \to T'$ を $e$ の基底変換と書く．図式は
$$
\xymatrix{
X \ar[d]_f & Y \ar[l]^h \ar[d]^e & Y' \ar[l]^{\pi'} \ar[d]^{e'} \\
S & T \ar[l]_g & T' \ar[l]_\pi
}
$$
である．命題 \ref{etale-cohomology-proposition-finite-higher-direct-image-zero} と補題
\ref{etale-cohomology-lemma-finite-pushforward-commutes-with-base-change} により
$R\pi'_*(e')^{-1}\mathcal{I}' = e^{-1}\pi_*\mathcal{I}'$ である．したがって
Leray スペクトル系列（Cohomology on Sites，補題
\ref{sites-cohomology-lemma-Leray}）により
$$
H^q(Y', (e')^{-1}\mathcal{I}') =
H^q(Y, e^{-1}\pi_*\mathcal{I}') \supset H^q(Y, e^{-1}\mathcal{I})
$$
であり，これは任意の \'etale 射 $U \to X$ による基底変換後にも成り立つ．
したがって $T$ を $T'$ で，$\mathcal{I}$ を $\mathcal{I}'$ で，
$\tilde \xi$ をその $H^q(Y', (e')^{-1}\mathcal{I}')$ における像で
置き換えてよい．

\medskip\noindent
$T \to S' \to S$ という分解があり，$\pi : S' \to S$ が有限であるとする．
$X' = S' \times_S X$ とおくと，誘導される図式
$$
\xymatrix{
X \ar[d]_f & X' \ar[l]^{\pi'} \ar[d]^{f'} & Y \ar[l]^{h'} \ar[d]^e \\
S & S' \ar[l]_\pi & T \ar[l]_g
}
$$
を考えられる．$\pi'$ の高次順像は消えるので，Leray スペクトル系列により
$R^qh_*e^{-1}\mathcal{I} = \pi'_*R^qh'_*e^{-1}\mathcal{I}$ である．したがって
$H^0(X, R^qh_*e^{-1}\mathcal{I}) = H^0(X', R^qh'_*e^{-1}\mathcal{I})$ である．
ゆえに $\xi$ が零であることと，対応する $R^qh'_*e^{-1}\mathcal{I}$ の
切断が零であることは同値である\footnote{この段階は別の仕方でも分かる．
実際，ある \'etale 被覆 $\{U_i \to X\}$ が存在し，$\tilde \xi$ が
$H^q(U_i \times_X Y, e^{-1}\mathcal{I})$ において消えることを示す必要がある．
一方，ある \'etale 被覆 $\{U'_j \to X'\}$ が存在し，$\tilde \xi$ が
$H^q(U'_i \times_{X'} Y, e^{-1}\mathcal{I})$ において消えることを証明できたとする．
このとき $X' \to X$ に対する性質 (B)（補題 \ref{etale-cohomology-lemma-finite-B}）により，
ある \'etale 被覆 $\{U_i \to X\}$ が存在して，$U_i \times_X X'$ は
$X'$ 上のスキームの非交和となり，その各成分は $U'_j$ を経由するような
$j$ をもつ．したがって所望のように，$\tilde \xi$ は
$H^q(U_i \times_X Y, e^{-1}\mathcal{I})$ において消える．}．
したがって $S$ を $S'$ で，$X$ を $X'$ で置き換えてよい．
$\sigma : S \to X$ は基底変換により $\sigma' : S' \to X'$ となるので，
この置換後にも $X \to S$ は切断 $\sigma$ と幾何学的に連結なファイバーを
もつことに注意する．

\medskip\noindent
$S$ と $T$ が Nagata スキームであることを用いる；Algebra，命題
\ref{algebra-proposition-ubiquity-nagata} を参照されたい．これにより種々の
正規化が有限になる；Morphisms の補題
\ref{morphisms-lemma-nagata-normalization-finite} および
\ref{morphisms-lemma-nagata-normalization} を参照されたい．特に，まず $T$ を
その正規化で置き換え，次に $S$ を，$S$ の $T$ における正規化で置き換えてよい．
すると $T \to S$ は整正規スキームの支配的射の非交和となる；Morphisms，補題
\ref{morphisms-lemma-normal-normalization} を参照されたい．明らかに連結成分ごとに
議論できるから，$T \to S$ は整正規スキームの支配的射であると仮定してよい．

\medskip\noindent
$s \in S$ と $t \in T$ を生成点とする．補題
\ref{etale-cohomology-lemma-smooth-base-change-fields} により，有限体拡大 $K/\kappa(t)$ および
$k/\kappa(s)$ で $k$ が $K$ に含まれるものと，有限 \'etale Galois 被覆
$Z \to X_k$ が存在する．その Galois 群 $G$ の位数は $n$ のある冪を割り，
$\sigma(\Spec(k))$ 上で分裂し，$\tilde \xi$ は
$H^q(Z_K, e^{-1}\mathcal{I}|_{Z_K})$ において零に写る．
$T' \to T$ を $T$ の $\Spec(K)$ における正規化，$S' \to S$ を $S$ の
$\Spec(k)$ における正規化とする．すると可換図式
$$
\xymatrix{
S' \ar[d] & T' \ar[l] \ar[d] \\
S & T \ar[l]
}
$$
を得る．その垂直射は有限である．上で述べた議論により，$S$ と $T$ を
$S'$ と $T'$ で置き換える（それに応じて $X$ を $X \times_S S'$ で，
$Y$ を $Y \times_T T'$ で置き換える）．この置換後，有限 \'etale Galois 被覆
$Z \to X_s$ で，$X \to S$ の生成ファイバーを被覆し，Galois 群 $G$ の位数が
$n$ のある冪を割り，$\sigma(s)$ 上で分裂し，$\tilde \xi$ が
$H^q(Z_t, (Z_t \to Y)^{-1}e^{-1}\mathcal{I})$ において零に写るものを得る．
ここで $Z_t = Z \times_S t = Z \times_s t = Z \times_{X_s} Y_t$ である．
$n$ は $S$ 上可逆なので，Fundamental Groups の補題
\ref{pione-lemma-extend-covering} により，有限 \'etale 射 $U \to X$ で，
$X_s$ への制限が $Z$ であるものを見つけられる．

\medskip\noindent
ここで $X$ を $U$ で，$Y$ を $U \times_X Y$ で置き換える．この置換後には
$X \to S$ のファイバーが幾何学的に連結でなくなるかもしれない
（切断は依然として存在するが，これは用いない）．その代わり，置換後には
$\tilde \xi$ が $H^q(Y_t, e^{-1}\mathcal{I})$ において零に写る．すなわち，
$\tilde \xi$ の $Y \to T$ の生成ファイバーへの制限は零である．

\medskip\noindent
$t$ は $T$ の函数体のスペクトルであること，すなわちスキームとして $t$ は
$T$ の空でないアフィン開部分スキームの極限であることを思い出そう．補題
\ref{etale-cohomology-lemma-directed-colimit-cohomology} により，空でない開部分スキーム
$V \subset T$ が存在して，$\tilde \xi$ は
$H^q(Y \times_T V, e^{-1}\mathcal{I}|_{Y \times_T V})$ において零に写る．

\medskip\noindent
$Z = T \setminus V$ とおく．図式
$$
\xymatrix{
Y \times_T Z \ar[d]_{e_Z} \ar[r]_{i'} &
Y \ar[d]_e &
Y \times_T V \ar[l]^{j'} \ar[d]^{e_V} \\
Z \ar[r]^i &
T &
V \ar[l]_j
}
$$
を考える．単射 $i^{-1}\mathcal{I} \to \mathcal{I}'$ を，
$\mathbf{Z}/n\mathbf{Z}$-加群の $Z$ 上の単射的層へのものとして選ぶ．
茎を調べると，写像
$$
\mathcal{I} \to j_*\mathcal{I}|_V \oplus i_*\mathcal{I}'
$$
は単射であり，$\mathcal{I}$ が $\mathbf{Z}/n\mathbf{Z}$-加群の単射的層なので
分裂することが分かる．したがって，$\tilde \xi$ が
$$
H^q(Y, e^{-1}j_*\mathcal{I}|_V) \oplus
H^q(Y, e^{-1}i_*\mathcal{I}')
$$
において，少なくとも $X$ をある \'etale 被覆の各要素で置き換えた後に
零に写ることを示せば十分である．次に注意する：
$$
e^{-1}j_*\mathcal{I}|_V = j'_*e_V^{-1}\mathcal{I}|_V,\quad
e^{-1}i_*\mathcal{I}' = i'_*e_Z^{-1}\mathcal{I}'
$$
$q$ に関する帰納法の仮定により，
$$
R^aj'_*e_V^{-1}\mathcal{I}|_V = 0, \quad a = 1, \ldots, q - 1
$$
である．$j'$ に対する Leray スペクトル系列と上の消滅から，写像
$$
H^q(Y, j'_*(e_V^{-1}\mathcal{I}|_V))
\longrightarrow
H^q(Y \times_T V, e_V^{-1}\mathcal{I}_V) =
H^q(Y \times_T V, e^{-1}\mathcal{I}|_{Y \times_T V})
$$
は単射である．したがって，上の第一直和成分における $\tilde \xi$ の像は
零である．実際，$\tilde \xi$ は
$H^q(Y \times_T V, e^{-1}\mathcal{I}|_{Y \times_T V})$ において消えることが
分かっている．$\dim(Z) < \dim(T) = d$ だから，帰納法により，$\tilde \xi$ の
第二直和成分
$$
H^q(Y, e^{-1}i_*\mathcal{I}') =
H^q(Y, i'_*e_Z^{-1}\mathcal{I}') =
H^q(Y \times_T Z, e_Z^{-1}\mathcal{I}')
$$
における像は，$X$ を適切な \'etale 被覆の各要素で
置き換えた後に消える．これで滑らかな基底変換定理の証明が完了する．
\end{proof}

\begin{proof}[滑らかな基底変換の第二証明]
この証明は，より長い第一証明と同じである．いくつかの補題で用いた論法を
別に取り出してあるために短くなっているにすぎない．

\medskip\noindent
$q = 0$ の場合は補題 \ref{etale-cohomology-lemma-fppf-reduced-fibres-base-change-f-star}
である．したがって $q > 0$ とし，より小さいすべての次数で結果が
成り立つと仮定してよい．

\medskip\noindent
$n \geq 1$ が $S$ 上可逆であるごとに，$\mathcal{F}[n]$ を，$\mathcal{F}$ の
$n$ で零化される切断からなる部分層とする．すると
$\mathcal{F} = \colim \mathcal{F}[n]$ である．これは $\mathcal{F}$ の茎に
関する仮定による．関手 $e^{-1}$ と $f^{-1}$ は
左随伴なので余極限と可換する．関手 $R^qh_*$ と $R^qg_*$ は補題
\ref{etale-cohomology-lemma-relative-colimit} によりフィルター余極限と可換する．
したがって $\mathcal{F}[n]$ に対して定理を証明すれば十分である．以下では
整数 $n$ で $S$ 上可逆なものを固定し，$\mathbf{Z}/n\mathbf{Z}$-加群の層を扱う．

\medskip\noindent
補題 \ref{etale-cohomology-lemma-base-change-local} により，問題は $X$ と $S$ 上 \'etale
局所的である．滑らかな射の局所構造により（Morphisms，補題
\ref{morphisms-lemma-smooth-etale-over-affine-space} を参照），$X$ と $S$ は
アフィンであり，$X \to S$ は \'etale 射 $X \to \mathbf{A}^d_S$ を経由して
分解すると仮定してよい．$X \to S$ を合成
$$
X \to \mathbf{A}^{d - 1}_S \to \mathbf{A}^{d - 2}_S \to \ldots \to
\mathbf{A}^1_S \to S
$$
として書くと，補題 \ref{etale-cohomology-lemma-base-change-compose} により，$X$ と $S$ が
アフィンで，$X \to S$ の相対次元が $1$ である場合に定理を証明すれば
十分であると分かる．

\medskip\noindent
補題 \ref{etale-cohomology-lemma-base-change-does-not-hold} により，
$R^qh_*\mathbf{Z}/d\mathbf{Z} = 0$ を $d | n$ に対して示せば十分である．
ここで考えるのはファイバー積図式
$$
\xymatrix{
X \ar[d] & Y \ar[d] \ar[l]^h \\
S & \Spec(K) \ar[l]
}
$$
であり，$X \to S$ はアフィンかつ相対次元 $1$ の滑らかな射，$S$ は
正規整域 $A$ のスペクトルで，その分数体 $L$ は代数閉であり，$K/L$ は
代数閉体の拡大である．

\medskip\noindent
$R^qh_*\mathbf{Z}/d\mathbf{Z}$ は前層
$$
U \longmapsto
H^q(U \times_X Y, \mathbf{Z}/d\mathbf{Z}) =
H^q(U \times_S \Spec(K), \mathbf{Z}/d\mathbf{Z})
$$
に付随する $X_\etale$ 上の層であることを思い出そう（補題
\ref{etale-cohomology-lemma-higher-direct-images}）．したがって，$U$ と
$\xi \in H^q(U \times_S \Spec(K), \mathbf{Z}/d\mathbf{Z})$ が与えられたとき，
\'etale 被覆 $\{U_i \to U\}$ が存在して，$\xi$ が
$H^q(U_i \times_S \Spec(K), \mathbf{Z}/d\mathbf{Z})$ において消えることを
示せば十分である．

\medskip\noindent
もちろん $U$ はアフィンとしてよい．すると $U \times_S \Spec(K)$ は
$K$ 上の（滑らかな）アフィン曲線であるから，定理
\ref{etale-cohomology-theorem-vanishing-affine-curves} により $q > 1$ では消滅する．

\medskip\noindent
最後に $q = 1$ の場合を扱う．More on Morphisms の補題
\ref{more-morphisms-lemma-cover-smooth-by-special} と同様に，$U$ をある
\'etale 被覆の各要素で置き換えてよい．すると $U \to S$ は
$U \to V \to S$ と分解する．ここで $U \to V$ は幾何学的に連結な
ファイバーをもち，$U$, $V$ はアフィン，$V \to S$ は \'etale であり，
切断 $\sigma : V \to U$ が存在する．補題
\ref{etale-cohomology-lemma-normal-scheme-with-alg-closed-function-field} により，$V$ は $S$ の
（アフィン）開部分スキームの（有限）非交和と同型である．明らかに $S$ を
その一つで置き換え，$X$ を $U$ の対応する成分で置き換えてよい．したがって
$X \to S$ は幾何学的に連結なファイバーと切断 $\sigma$ をもち，
$\xi \in H^1(Y, \mathbf{Z}/d\mathbf{Z})$ であると仮定してよい．
$K$ と $L$ は代数閉なので，
$$
H^1(X_L, \mathbf{Z}/d\mathbf{Z}) =
H^1(Y, \mathbf{Z}/d\mathbf{Z})
$$
である；補題 \ref{etale-cohomology-lemma-base-change-dim-1-separably-closed} を参照されたい．
したがって，有限 \'etale Galois 被覆 $Z \to X_L$ で，Galois 群が
$G \subset \mathbf{Z}/d\mathbf{Z}$ であり，$\xi$ を零化するものが存在する．これは補題
\ref{etale-cohomology-lemma-smooth-base-change-fields} の主張または証明を見てもよいし，
$\xi$ が $\mathbf{Z}/d\mathbf{Z}$-トーサーで $X_L$ 上のものに対応することを
直接用いてもよい．最後に Fundamental Groups の補題
\ref{pione-lemma-extend-covering-general} により，（必然的に全射な）有限
\'etale 射 $X' \to X$ で，その $X_L$ への制限が $Z \to X_L$ であるものを得る．
$\xi$ は $X'_K$ において消えるから，これで証明が完了する．
\end{proof}

\noindent
実際にしばしば用いられるのは，滑らかな基底変換定理の次の直接の帰結である．

\begin{lemma}
\label{etale-cohomology-lemma-smooth-base-change-general}
$S$ をスキームとする．$S' = \lim S_i$ を，スキーム $S_i$ で $S$ 上滑らかなものの
有向逆極限で，遷移射がアフィンであるものとする．$f : X \to S$ を
準コンパクトかつ準分離的とし，ファイバー正方形
$$
\xymatrix{
X' \ar[d]_{f'} \ar[r]_{g'} & X \ar[d]^f \\
S' \ar[r]^g & S
}
$$
を作る．このとき
$$
g^{-1}Rf_*E = R(f')_*(g')^{-1}E
$$
が，任意の $E \in D^+(X_\etale)$ で，そのコホモロジー層 $H^q(E)$ の茎が
$S$ 上可逆な位数の捩れであるものに対して成り立つ．
\end{lemma}

\begin{proof}
スペクトル系列
$$
E_2^{p, q} = R^pf_*H^q(E)
\quad\text{かつ}\quad
{E'}_2^{p, q} = R^pf'_*H^q((g')^{-1}E) = R^pf'_*(g')^{-1}H^q(E)
$$
を考える．これらは $R^nf_*E$ および $R^nf'_*(g')^{-1}E$ に収束する．
これらのスペクトル系列は Derived Categories の補題
\ref{derived-lemma-two-ss-complex-functor} で構成されている．滑らかな基底変換定理
（定理 \ref{etale-cohomology-theorem-smooth-base-change}）と補題
\ref{etale-cohomology-lemma-base-change-Rf-star-colim} を組み合わせると，
$$
g^{-1}R^pf_*H^q(E) = R^p(f')_*(g')^{-1}H^q(E)
$$
を得る．以上をすべて組み合わせれば補題を得る．
\end{proof}












\section{滑らかな基底変換の応用}
\label{etale-cohomology-section-applications-smooth-base-change}

\noindent
本節では，滑らかな基底変換定理からほぼ直ちに従ういくつかの帰結を論じる．

\begin{lemma}
\label{etale-cohomology-lemma-base-change-field-extension}
$L/K$ を体の拡大とする．$g : T \to S$ を $K$ 上のスキームの準コンパクト
かつ準分離的な射とする．$g_L : T_L \to S_L$ を $g$ の $\Spec(L)$ への
基底変換と書く．$E \in D^+(T_\etale)$ のコホモロジー層の茎は，$K$ において
可逆な位数の捩れであるとする．$E_L$ を $E$ の $(T_L)_\etale$ への引戻しとする．
このとき $Rg_{L, *}E_L$ は $Rg_*E$ の $S_L$ への引戻しである．
\end{lemma}

\begin{proof}
$L/K$ が分離的ならば，$L$ は滑らかな $K$-代数のフィルター余極限である；
Algebra，補題 \ref{algebra-lemma-colimit-syntomic} を参照されたい．したがって
この場合の補題は補題 \ref{etale-cohomology-lemma-smooth-base-change-general} から直ちに従う．
一般の場合，$K'$ と $L'$ を $K$ と $L$ の完全閉包とする
（Algebra，定義 \ref{algebra-definition-perfection}）．すると
$\Spec(K') \to \Spec(K)$ と $\Spec(L') \to \Spec(L)$ は普遍同相写像である．
実際，$K'/K$ と $L'/L$ は純非分離的である（Algebra，補題
\ref{algebra-lemma-p-ring-map} を参照）．したがって，
$(T_{K'})_\etale = T_\etale$,
$(S_{K'})_\etale = S_\etale$,
$(T_{L'})_\etale = (T_L)\etale$，および
$(S_{L'})_\etale = (S_L)_\etale$ である．これは \'etale コホモロジーの
位相的不変性による；命題 \ref{etale-cohomology-proposition-topological-invariance} を
参照されたい．これにより，補題は分離的な体拡大 $L'/K'$ の場合に帰着する
（完全体の定義による；Algebra，定義
\ref{algebra-definition-perfect} を参照されたい）．
\end{proof}

\begin{lemma}
\label{etale-cohomology-lemma-smooth-base-change-separably-closed}
$K/k$ を分離閉体の拡大とする．$X$ を $k$ 上の準コンパクトかつ
準分離的なスキームとする．$E \in D^+(X_\etale)$ のコホモロジー層の茎は，
$k$ において可逆な位数の捩れであるとする．このとき
\begin{enumerate}
\item 写像 $H^q_\etale(X, E) \to H^q_\etale(X_K, E|_{X_K})$ は同型である；
\item $E \to R(X_K \to X)_*E|_{X_K}$ は同型である．
\end{enumerate}
\end{lemma}

\begin{proof}
(1) の証明．まず $\overline{k}$ と $\overline{K}$ を $k$ と $K$ の
代数閉包とする．射 $\Spec(\overline{k}) \to \Spec(k)$ と
$\Spec(\overline{K}) \to \Spec(K)$ は普遍同相写像である．実際，
$\overline{k}/k$ と $\overline{K}/K$ は純非分離的である
（Algebra，補題 \ref{algebra-lemma-p-ring-map} を参照）．したがって，
$H^q_\etale(X, \mathcal{F}) =
H^q_\etale(X_{\overline{k}}, \mathcal{F}_{X_{\overline{k}}})$
である．これは \'etale コホモロジーの位相的不変性による；命題
\ref{etale-cohomology-proposition-topological-invariance} を参照されたい．$X_K$ と
$X_{\overline{K}}$ についても同様である．したがって $k$ と $K$ は代数閉と
仮定してよい．この場合 $K$ は滑らかな $k$-代数の極限である；Algebra，補題
\ref{algebra-lemma-colimit-syntomic} を参照されたい．したがって本補題は，
定理 \ref{etale-cohomology-theorem-smooth-base-change} を補題
\ref{etale-cohomology-lemma-smooth-base-change-general} の形に言い換えたものの特殊な場合である．

\medskip\noindent
(2) の証明．任意の準コンパクトかつ準分離的な $U$ で $X_\etale$ に属するものに対し，
上の議論から，写像 $E \to R(X_K \to X)_*E|_{X_K}$ の制限は
コホモロジー上の同型を与える．$X_\etale$ の各対象はこのような $U$ による
\'etale 被覆をもつので，所望の主張が従う．
\end{proof}

\begin{lemma}
\label{etale-cohomology-lemma-base-change-does-not-hold-post}
注意 \ref{etale-cohomology-remark-base-change-holds} における $f : X \to S$ と $n$ に対し，
$n$ は $S$ 上可逆であり，ある $q \geq 1$ について $BC(f, n, q - 1)$ は
真だが $BC(f, n, q)$ は真でないと仮定する．このとき可換図式
$$
\xymatrix{
X \ar[d]_f & X' \ar[d] \ar[l] & Y \ar[l]^h \ar[d] \\
S & S' \ar[l] & \Spec(K) \ar[l]
}
$$
が存在し，両方の正方形はファイバー積正方形である．ここで $S'$ は
アフィン，整かつ正規で，その函数体 $K$ は代数閉であり，さらに整数
$d | n$ で $R^qh_*(\mathbf{Z}/d\mathbf{Z})$ が零でないものが存在する．
\end{lemma}

\begin{proof}
まず補題 \ref{etale-cohomology-lemma-base-change-does-not-hold} における図式と
$\mathcal{F}$ を選ぶ．$S'$ はアフィンであると仮定する（これは明らかだが，
疑わしければ同補題の証明を参照されたい）．$K'$ を $S'$ の函数体とし，
$Y' = X' \times_{S'} \Spec(K')$ とおくと，図式
$$
\xymatrix{
X \ar[d]_f &
X' \ar[d] \ar[l] &
Y' \ar[l]^{h'} \ar[d] &
Y \ar[l] \ar[d] \\
S &
S' \ar[l] &
\Spec(K') \ar[l] &
\Spec(K) \ar[l]
}
$$
を得る．補題 \ref{etale-cohomology-lemma-smooth-base-change-separably-closed} により，全順像
$R(Y \to Y')_*\mathbf{Z}/d\mathbf{Z}$ は $\mathbf{Z}/d\mathbf{Z}$ と
$D(Y'_\etale)$ において同型である；ここで $n$ が $S$ 上可逆であることを
用いた．したがって相対 Leray スペクトル系列により
$Rh'_*\mathbf{Z}/d\mathbf{Z} = Rh_*\mathbf{Z}/d\mathbf{Z}$ である．
これで証明が完了する．
\end{proof}










\section{固有基底変換定理}
\label{etale-cohomology-section-proper-base-change}

\noindent
本節では固有基底変換定理を述べ，証明する．我々の方法は概ね
\cite[XII, Theorem 5.1]{SGA4} の証明に従い，Gabber の（アフィンの場合の）
着想を用いて議論を少し簡略化する．

\begin{lemma}
\label{etale-cohomology-lemma-zariski-h0-proper-over-henselian-pair}
$(A, I)$ を Hensel 対とする．$f : X \to \Spec(A)$ をスキームの固有射とする．
$Z = X \times_{\Spec(A)} \Spec(A/I)$ とおく．任意の層 $\mathcal{F}$ を
$X$ に付随する位相空間上にとると，
$\Gamma(X, \mathcal{F}) = \Gamma(Z, \mathcal{F}|_Z)$ である．
\end{lemma}

\begin{proof}
これを証明するため補題 \ref{etale-cohomology-lemma-h0-topological} を用いる．まず Properties の補題
\ref{properties-lemma-quasi-compact-quasi-separated-spectral} により，$X$ の
台位相空間はスペクトル空間である．$Y \subset X$ を既約閉部分スキームとする．
証明を終えるには，$Y \cap Z = Y \times_{\Spec(A)} \Spec(A/I)$ が連結であることを
示せばよい．$X$ を $Y$ で置き換えると，$X$ は既約であり，$Z$ が連結であることを
示せばよいと仮定できる．$X \to \Spec(B) \to \Spec(A)$ を $f$ の Stein 分解とする
（More on Morphisms，定理
\ref{more-morphisms-theorem-stein-factorization-general}）．すると $A \to B$ は
整であり，$(B, IB)$ は Hensel 対である（More on Algebra，補題
\ref{more-algebra-lemma-integral-over-henselian-pair}）．したがって
$X \to \Spec(A)$ のファイバーは幾何学的に連結であると仮定してよい．
他方，$T \subset \Spec(A)$ を $f$ の像とすると，これは $X$ が $A$ 上
固有なので既約かつ閉である．
ゆえに More on Algebra の補題
\ref{more-algebra-lemma-irreducible-henselian-pair-connected} により
$T \cap V(I)$ は連結である．いま
$Y \times_{\Spec(A)} \Spec(A/I) \to T \cap V(I)$ は，連結なファイバーをもつ
全射な閉写像である．したがって結果は Topology の補題
\ref{topology-lemma-connected-fibres-quotient-topology-connected-components} から従う．
\end{proof}

\begin{lemma}
\label{etale-cohomology-lemma-h0-proper-over-henselian-pair}
$(A, I)$ を Hensel 対とする．$f : X \to \Spec(A)$ をスキームの固有射とする．
$i : Z \to X$ を $X \times_{\Spec(A)} \Spec(A/I)$ の $X$ への閉埋込みとする．
任意の層 $\mathcal{F}$ を $X_\etale$ 上にとると，
$\Gamma(X, \mathcal{F}) = \Gamma(Z, i_{small}^{-1}\mathcal{F})$ である．
\end{lemma}

\begin{proof}
これは補題 \ref{etale-cohomology-lemma-gabber-h0} と
\ref{etale-cohomology-lemma-zariski-h0-proper-over-henselian-pair}，および $X$ 上有限な任意の
スキームが $\Spec(A)$ 上固有であるという事実から従う．
\end{proof}

\begin{lemma}
\label{etale-cohomology-lemma-h0-proper-over-henselian-local}
$A$ を Hensel 局所環とする．$f : X \to \Spec(A)$ をスキームの固有射とする．
$X_0 \subset X$ を $f$ の閉点上のファイバーとする．任意の層
$\mathcal{F}$ を $X_\etale$ 上にとると，$\Gamma(X, \mathcal{F}) =
\Gamma(X_0, \mathcal{F}|_{X_0})$ である．
\end{lemma}

\begin{proof}
これは補題 \ref{etale-cohomology-lemma-h0-proper-over-henselian-pair} の特殊な場合である．
\end{proof}

\noindent
$f : X \to S$ をスキームの射とする．
$\overline{s} : \Spec(k) \to S$ を幾何学点とする．$f$ の
$\overline{s}$ におけるファイバーとは，スキーム $X_{\overline{s}} =
\Spec(k) \times_{\overline{s}, S} X$ を $\Spec(k)$ 上のスキームとみなしたものである．
$\mathcal{F}$ が $X_\etale$ 上の層ならば，
$\mathcal{F}_{\overline{s}} = p_{small}^{-1}\mathcal{F}$ によって，
$\mathcal{F}$ の $(X_{\overline{s}})_\etale$ への引戻しを表す．以下では集合
$$
\Gamma(X_{\overline{s}}, \mathcal{F}_{\overline{s}})
$$
を考える．$s \in S$ を $\overline{s}$ の像とする．$\kappa(s)^{sep}$ を，定義
\ref{etale-cohomology-definition-algebraic-geometric-point} と同様に，$\kappa(s)$ の $k$ における
分離代数閉包とする．補題 \ref{etale-cohomology-lemma-sections-base-field-extension} により，
引戻しは全単射
$$
\Gamma(X_{\kappa(s)^{sep}},
p_{sep}^{-1} \mathcal{F})
\longrightarrow
\Gamma(X_{\overline{s}}, \mathcal{F}_{\overline{s}})
$$
を定める．ここで $p_{sep} : X_{\kappa(s)^{sep}} =
\Spec(\kappa(s)^{sep}) \times_S X \to X$ は射影である．

\begin{lemma}
\label{etale-cohomology-lemma-proper-pushforward-stalk}
$f : X \to S$ をスキームの固有射とする．$\overline{s} \to S$ を幾何学点とする．
任意の層 $\mathcal{F}$ を $X_\etale$ 上にとると，標準写像
$$
(f_*\mathcal{F})_{\overline{s}} \longrightarrow
\Gamma(X_{\overline{s}}, \mathcal{F}_{\overline{s}})
$$
は全単射である．
\end{lemma}

\begin{proof}
定理 \ref{etale-cohomology-theorem-higher-direct-images}（集合の層に対するもの）により，
$$
(f_*\mathcal{F})_{\overline{s}} =
\Gamma(X \times_S \Spec(\mathcal{O}_{S, \overline{s}}^{sh}),
p_{small}^{-1}\mathcal{F})
$$
である．ここで $p : X \times_S \Spec(\mathcal{O}_{S, \overline{s}}^{sh}) \to X$
は射影である．狭義 Hensel 局所環 $\mathcal{O}_{S, \overline{s}}^{sh}$ の剰余体は
$\kappa(s)^{sep}$ なので，上の議論と補題
\ref{etale-cohomology-lemma-h0-proper-over-henselian-local} から本補題が従う．
\end{proof}

\begin{lemma}
\label{etale-cohomology-lemma-proper-base-change-f-star}
$f : X \to Y$ をスキームの固有射とする．$g : Y' \to Y$ をスキームの射とする．
$X' = Y' \times_Y X$ とおき，射影を $f' : X' \to Y'$ および
$g' : X' \to X$ とする．$\mathcal{F}$ を $X_\etale$ 上の任意の層とする．
このとき $g^{-1}f_*\mathcal{F} = f'_*(g')^{-1}\mathcal{F}$ である．
\end{lemma}

\begin{proof}
標準写像 $g^{-1}f_*\mathcal{F} \to f'_*(g')^{-1}\mathcal{F}$ がある．
すなわち，これは写像
$$
f_*\mathcal{F} \longrightarrow
g_*f'_*(g')^{-1}\mathcal{F} = f_*g'_*(g')^{-1}\mathcal{F}
$$
に随伴する．この写像は，$f_*$ を標準写像
$\mathcal{F} \to g'_*(g')^{-1}\mathcal{F}$ に適用したものである．
この写像が同型であることを確認するには，茎上で何が起こるかを計算すればよい．
$y' : \Spec(k) \to Y'$ を幾何学点とし，$y$ をその $Y$ における像とする．補題
\ref{etale-cohomology-lemma-proper-pushforward-stalk} により，二つの茎はそれぞれ
$\Gamma(X'_{y'}, \mathcal{F}_{y'})$ と $\Gamma(X_y, \mathcal{F}_y)$ である．
ここで層 $\mathcal{F}_y$ と $\mathcal{F}_{y'}$ は，$\mathcal{F}$ の，射影
$X_y \to X$ と $X'_{y'} \to X$ による引戻しである．スキームとして
$X_y = X'_{y'}$ なので，茎は一致する．この同型が我々の写像と両立することの
確認は省略する．
\end{proof}

\noindent
ここから固有基底変換定理の議論を始める．そのために記法を導入する．
スキームの射の可換図式
\begin{equation}
\label{etale-cohomology-equation-base-change-diagram}
\vcenter{
\xymatrix{
X' \ar[r]_{g'} \ar[d]_{f'} & X \ar[d]^f \\
Y' \ar[r]^g & Y
}
}
\end{equation}
を考える．するとサイトの可換図式
$$
\xymatrix{
X'_\etale \ar[r]_{g'_{small}} \ar[d]_{f'_{small}} &
X_\etale \ar[d]^{f_{small}} \\
Y'_\etale \ar[r]^{g_{small}} &
Y_\etale
}
$$
を得る．任意の対象 $E$ を $D(X_\etale)$ にとると，標準的な基底変換写像
\begin{equation}
\label{etale-cohomology-equation-base-change}
g_{small}^{-1}Rf_{small, *}E \longrightarrow Rf'_{small, *}(g'_{small})^{-1}E
\end{equation}
を $D(Y'_\etale)$ において得る．Cohomology on Sites の注意
\ref{sites-cohomology-remark-base-change} を参照されたい．そこでは環の層として
定数層 $\mathbf{Z}$ を用いる．この公式では通常，下付き添字 ${}_{small}$ を省略する．
例えば，$E = \mathcal{F}[0]$ で，$\mathcal{F}$ が $X_\etale$ 上のアーベル層ならば，
基底変換写像は写像
\begin{equation}
\label{etale-cohomology-equation-base-change-sheaf}
g^{-1}Rf_*\mathcal{F} \longrightarrow Rf'_*(g')^{-1}\mathcal{F}
\end{equation}
であり，$D(Y'_\etale)$ に属する．

\medskip\noindent
写像 (\ref{etale-cohomology-equation-base-change}) は，上で与えた一般性のもとでは同型に
なりえない．目標は，図式 (\ref{etale-cohomology-equation-base-change-diagram}) がファイバー積図式，
$f : X \to Y$ が固有，$E$ のコホモロジー層が捩れで，$E$ が下に有界ならば，
これが同型であることを示すことである．この問題を調べるため次の用語を導入する．
{\it $f : X \to Y$ に対してコホモロジーは基底変換と可換する}とは，
(\ref{etale-cohomology-equation-base-change-sheaf}) が，図式
(\ref{etale-cohomology-equation-base-change-diagram}) で $X' = Y' \times_Y X$ であるものすべてと，
任意の捩れアーベル層 $\mathcal{F}$ に対して同型であることをいう．

\begin{lemma}
\label{etale-cohomology-lemma-proper-base-change-in-terms-of-injectives}
$f : X \to Y$ をスキームの固有射とする．次は同値である：
\begin{enumerate}
\item $f$ に対してコホモロジーは基底変換と可換する（上を参照）；
\item 任意の素数 $\ell$，任意の $\mathbf{Z}/\ell\mathbf{Z}$-加群の
単射的層 $\mathcal{I}$ を $X_\etale$ 上にとり，さらに図式
(\ref{etale-cohomology-equation-base-change-diagram}) で $X' = Y' \times_Y X$ である任意のものに対し，
層 $R^qf'_*(g')^{-1}\mathcal{I}$ は $q > 0$ ならば零である．
\end{enumerate}
\end{lemma}

\begin{proof}
(1) が (2) を含意することは明らかである．逆に (2) を仮定し，$\mathcal{F}$ を
$X_\etale$ 上の捩れアーベル層とする．$Y' \to Y$ をスキームの射とし，
$X' = Y' \times_Y X$ とおく．図式 (\ref{etale-cohomology-equation-base-change-diagram}) と同様，
射影を $g' : X' \to X$ および $f' : X' \to Y'$ とする．層の写像
$$
g^{-1}R^qf_*\mathcal{F} \longrightarrow R^qf'_*(g')^{-1}\mathcal{F}
$$
がすべての $q \geq 0$ に対して同型であることを示したい．

\medskip\noindent
各 $n \geq 1$ に対し，$\mathcal{F}[n]$ を，$\mathcal{F}$ の切断のうち
$n$ で零化されるものからなる部分層とする．このとき
$\mathcal{F} = \colim \mathcal{F}[n]$ である．
関手 $g^{-1}$ と $(g')^{-1}$ は（左随伴なので）任意の余極限と可換する．
$f$ または $f'$ に沿う高次順像は，補題 \ref{etale-cohomology-lemma-relative-colimit} により
フィルター余極限と可換する．ゆえに
$$
g^{-1}R^qf_*\mathcal{F} = \colim g^{-1}R^qf_*\mathcal{F}[n]
\quad\text{かつ}\quad
R^qf'_*(g')^{-1}\mathcal{F} =
\colim R^qf'_*(g')^{-1}\mathcal{F}[n]
$$
を得る．したがって，$\mathcal{F}$ が正整数 $n$ で零化される場合に結果を
証明すれば十分である．

\medskip\noindent
$n = \ell n'$ で，$\ell$ が素数であるとすると，短完全列
$$
0 \to \mathcal{F}[\ell] \to \mathcal{F} \to
\mathcal{F}/\mathcal{F}[\ell] \to 0
$$
を得る．$\mathcal{F}/\mathcal{F}[\ell]$ は $n'$ で零化されることに注意する．
さらに，$\mathcal{F}[\ell]$ と $\mathcal{F}/\mathcal{F}[\ell]$ の双方に対して
結果が成り立つならば，高次順像の長完全列（および $5$-補題）により結果が
成り立つ．このようにして，$\mathcal{F}$ が素数 $\ell$ で零化される場合に帰着する．

\medskip\noindent
$\mathcal{F}$ は素数 $\ell$ で零化されると仮定する．単射分解
$\mathcal{F} \to \mathcal{I}^\bullet$ を
$D(X_\etale, \mathbf{Z}/\ell\mathbf{Z})$ において選ぶ．仮定 (2) と Leray の
非輪状性補題（Derived Categories，補題
\ref{derived-lemma-leray-acyclicity}）を適用すると，
$$
f'_*(g')^{-1}\mathcal{I}^\bullet
$$
は $Rf'_*(g')^{-1}\mathcal{F}$ を計算する．最後に補題
\ref{etale-cohomology-lemma-proper-base-change-f-star} を適用すればよい．
\end{proof}

\begin{lemma}
\label{etale-cohomology-lemma-sandwich}
$f : X \to Y$ と $g : Y \to Z$ をスキームの固有射とする．次を仮定する：
\begin{enumerate}
\item $f$ に対してコホモロジーは基底変換と可換する；
\item $g \circ f$ に対してコホモロジーは基底変換と可換する；
\item $f$ は全射である．
\end{enumerate}
このとき $g$ に対してコホモロジーは基底変換と可換する．
\end{lemma}

\begin{proof}
以下，補題 \ref{etale-cohomology-lemma-proper-base-change-in-terms-of-injectives} の同値性を
断りなく用いる．$\ell$ を素数とする．$\mathcal{I}$ を
$\mathbf{Z}/\ell\mathbf{Z}$-加群の単射的層として $Y_\etale$ 上にとる．層の単射
$f^{-1}\mathcal{I} \to \mathcal{J}$ を選び，$\mathcal{J}$ は
$\mathbf{Z}/\ell\mathbf{Z}$-加群の単射的層として $X_\etale$ 上にあるとする．$f$ は全射なので，
写像 $\mathcal{I} \to f_*\mathcal{J}$ は単射である（幾何学点における茎を見よ）．
$\mathcal{I}$ は単射的だから，$\mathcal{I}$ は $f_*\mathcal{J}$ の直和因子である．
したがって $f_*\mathcal{J}$ に対する所望の消滅を証明すれば十分である．

\medskip\noindent
スキームの射 $Z' \to Z$ を取り，$Y' = Z' \times_Z Y$ および
$X' = Z' \times_Z X = Y' \times_ Y X$ とおく．射影を
$a : X' \to X$, $b : Y' \to Y$, $c : Z' \to Z$ と書く．
$f' : X' \to Y'$ および $g' : Y' \to Z'$ についても同様である．補題
\ref{etale-cohomology-lemma-proper-base-change-f-star} により
$b^{-1}f_*\mathcal{J} = f'_*a^{-1}\mathcal{J}$.
他方，$R^qf'_*a^{-1}\mathcal{J}$ と
$R^q(g' \circ f')_*a^{-1}\mathcal{J}$ は $q > 0$ に対して零であることが
分かっている．スペクトル系列（Cohomology on Sites，補題
\ref{sites-cohomology-lemma-relative-Leray}）
$$
R^pg'_* R^qf'_* a^{-1}\mathcal{J} \Rightarrow
R^{p + q}(g' \circ f')_* a^{-1}\mathcal{J}
$$
を用いると，
$ R^pg'_*(b^{-1}f_*\mathcal{J}) = R^pg'_*(f'_*a^{-1}\mathcal{J}) = 0$
を $p > 0$ に対して得る．これは所望の結論である．
\end{proof}

\begin{lemma}
\label{etale-cohomology-lemma-composition}
$f : X \to Y$ と $g : Y \to Z$ をスキームの固有射とする．次を仮定する：
\begin{enumerate}
\item $f$ に対してコホモロジーは基底変換と可換する；
\item $g$ に対してコホモロジーは基底変換と可換する．
\end{enumerate}
このとき $g \circ f$ に対してコホモロジーは基底変換と可換する．
\end{lemma}

\begin{proof}
以下，補題 \ref{etale-cohomology-lemma-proper-base-change-in-terms-of-injectives} の同値性を
断りなく用いる．$\ell$ を素数とする．$\mathcal{I}$ を
$\mathbf{Z}/\ell\mathbf{Z}$-加群の単射的層として $X_\etale$ 上にとる．すると
$f_*\mathcal{I}$ は $\mathbf{Z}/\ell\mathbf{Z}$-加群の単射的層として
$Y_\etale$ 上にある（Cohomology on Sites，補題
\ref{sites-cohomology-lemma-pushforward-injective-flat}）．結果はここから形式的に
従うが，以下で詳述する．

\medskip\noindent
スキームの射 $Z' \to Z$ を取り，$Y' = Z' \times_Z Y$ および
$X' = Z' \times_Z X = Y' \times_ Y X$ とおく．射影を
$a : X' \to X$, $b : Y' \to Y$, $c : Z' \to Z$ と書く．
$f' : X' \to Y'$ および $g' : Y' \to Z'$ についても同様である．補題
\ref{etale-cohomology-lemma-proper-base-change-f-star} により
$b^{-1}f_*\mathcal{I} = f'_*a^{-1}\mathcal{I}$.
他方，$R^qf'_*a^{-1}\mathcal{I}$ と
$R^q(g')_*b^{-1}f_*\mathcal{I}$ は $q > 0$ に対して零であることが
分かっている．スペクトル系列（Cohomology on Sites，補題
\ref{sites-cohomology-lemma-relative-Leray}）
$$
R^pg'_* R^qf'_* a^{-1}\mathcal{I} \Rightarrow
R^{p + q}(g' \circ f')_* a^{-1}\mathcal{I}
$$
を用いると，所望のとおり
$R^p(g' \circ f')_*a^{-1}\mathcal{I} = 0$ が $p > 0$ に対して成り立つ．
\end{proof}

\begin{lemma}
\label{etale-cohomology-lemma-finite}
\begin{slogan}
\'Etale コホモロジーの固有基底変換は有限射に対して成り立つ．
\end{slogan}
$f : X \to Y$ をスキームの有限射とする．このとき $f$ に対して
コホモロジーは基底変換と可換する．
\end{lemma}

\begin{proof}
有限射は固有であることに注意する；Morphisms，補題
\ref{morphisms-lemma-finite-proper} を参照されたい．さらに有限射の基底変換は
有限である；Morphisms，補題 \ref{morphisms-lemma-base-change-finite} を
参照されたい．したがって結果は補題
\ref{etale-cohomology-lemma-proper-base-change-in-terms-of-injectives} と命題
\ref{etale-cohomology-proposition-finite-higher-direct-image-zero} を組み合わせると従う．
\end{proof}

\begin{lemma}
\label{etale-cohomology-lemma-reduce-to-P1}
スキームの任意の固有射に対してコホモロジーが基底変換と可換することを
証明するには，射 $\mathbf{P}^1_S \to S$ について，任意のスキーム $S$ に対し
それが成り立つことを証明すれば十分である．
\end{lemma}

\begin{proof}
$f : X \to Y$ をスキームの固有射とする．$Y = \bigcup Y_i$ をアフィン開被覆とし，
$X_i = f^{-1}(Y_i)$ とおく．$X_i \to Y_i$ に対してコホモロジーが基底変換と
可換することを証明できれば，$f$ に対してもコホモロジーは基底変換と可換する．
実際，高次順像の形成は基底上の Zariski（さらには \'etale）局所化と可換する；
補題 \ref{etale-cohomology-lemma-higher-direct-images} を参照されたい．したがって $Y$ は
アフィンであると仮定してよい．

\medskip\noindent
$Y$ をアフィンスキームとし，$X \to Y$ を固有射とする．Chow の補題により
可換図式
$$
\xymatrix{
X \ar[rd] & X' \ar[d] \ar[l]^\pi \ar[r] & \mathbf{P}^n_Y \ar[dl] \\
& Y &
}
$$
が存在する．ここで $X' \to \mathbf{P}^n_Y$ は埋込み，かつ
$\pi : X' \to X$ は固有全射である；Limits，補題
\ref{limits-lemma-chow-finite-type} を参照されたい．$X \to Y$ が固有なので，
$X' \to Y$ は固有である（Morphisms，補題
\ref{morphisms-lemma-composition-proper}）．したがって
$X' \to \mathbf{P}^n_Y$ は閉埋込みである（Morphisms，補題
\ref{morphisms-lemma-image-proper-scheme-closed}）．ゆえに
$X' \to X \times_Y \mathbf{P}^n_Y = \mathbf{P}^n_X$ は閉埋込みである
（像が閉である埋込みだからである）．

\medskip\noindent
補題 \ref{etale-cohomology-lemma-sandwich} により，$\pi$ および $X' \to Y$ に対して
コホモロジーが基底変換と可換することを証明すれば十分である．これらの射は
いずれも閉埋込みと射影 $\mathbf{P}^n_S \to S$（ある $S$ に対するもの）の
合成に分解する．補題 \ref{etale-cohomology-lemma-finite} により結果は閉埋込みに対して成り立つ
（閉埋込みは有限だからである）．補題 \ref{etale-cohomology-lemma-composition} により，射影
$\mathbf{P}^n_S \to S$ に対して結果を証明すれば十分である．

\medskip\noindent
各 $n \geq 1$ に対し，有限全射
$$
\mathbf{P}^1_S \times_S \ldots \times_S \mathbf{P}^1_S
\longrightarrow
\mathbf{P}^n_S
$$
が存在し，座標上では
$$
((x_1 : y_1), (x_2 : y_2), \ldots, (x_n : y_n))
\longmapsto
(F_0 : \ldots : F_n)
$$
で与えられる．ここで $F_0, \ldots, F_n$ は変数 $x_1, \ldots, y_n$ の
整数係数多項式で，
$$
\prod (x_i t + y_i) = F_0 t^n + F_1 t^{n - 1} + \ldots + F_n
$$
を満たす．補題 \ref{etale-cohomology-lemma-sandwich}, \ref{etale-cohomology-lemma-finite},
\ref{etale-cohomology-lemma-composition} をもう一度適用すると，本補題が成り立つと分かる．
\end{proof}

\begin{theorem}[固有基底変換]
\label{etale-cohomology-theorem-proper-base-change}
$f : X \to Y$ をスキームの固有射とする．$g : Y' \to Y$ をスキームの射とする．
$X' = Y' \times_Y X$ とおき，ファイバー積図式
$$
\xymatrix{
X' \ar[r]_{g'} \ar[d]_{f'} & X \ar[d]^f \\
Y' \ar[r]^g & Y
}
$$
を考える．$\mathcal{F}$ を $X_\etale$ 上の捩れアーベル層とする．
このとき基底変換写像
$$
g^{-1}Rf_*\mathcal{F} \longrightarrow Rf'_*(g')^{-1}\mathcal{F}
$$
は同型である．
\end{theorem}

\begin{proof}
上で導入した用語では，これはスキームの任意の固有射に対して
コホモロジーが基底変換と可換することを意味する．補題
\ref{etale-cohomology-lemma-reduce-to-P1} により，射 $\mathbf{P}^1_S \to S$ について，任意の
スキーム $S$ に対しコホモロジーが基底変換と可換することを証明すれば十分である．

\medskip\noindent
$S$ を閉点 $s$ をもつ狭義 Hensel 局所環のスペクトルとする．
$X = \mathbf{P}^1_S$ および $X_0 = X_s = \mathbf{P}^1_s$ とおく．
$\mathcal{F}$ を $\mathbf{Z}/\ell\mathbf{Z}$-加群の層として $X_\etale$ 上にとる．
証明の鍵は
$$
H^q_\etale(X, \mathcal{F}) = H^q_\etale(X_0, \mathcal{F}|_{X_0}).
$$
であることである．実際，$\mathcal{F} \to \mathcal{I}^\bullet$ を
$\mathbf{Z}/\ell\mathbf{Z}$-加群の単射的層による分解として選ぶ．すると
$\mathcal{I}^\bullet|_{X_0}$ は $\mathcal{F}|_{X_0}$ の，右
$H^0_\etale(X_0, -)$-非輪状対象による分解である；補題
\ref{etale-cohomology-lemma-efface-cohomology-on-fibre-by-finite-cover} を参照されたい．
Leray の非輪状性補題によれば，右辺は複体
$H^0_\etale(X_0, \mathcal{I}^\bullet|_{X_0})$ により計算される．これは補題
\ref{etale-cohomology-lemma-h0-proper-over-henselian-local} により
$H^0_\etale(X, \mathcal{I}^\bullet)$ に等しい．この複体は左辺を計算する．

\medskip\noindent
$S$ を一般のものとし，$\mathcal{F}$ を
$\mathbf{Z}/\ell\mathbf{Z}$-加群の層として $X_\etale$ 上にとる．
$\overline{s} : \Spec(k) \to S$ を $S$ の幾何学点で，$s \in S$ 上にあるものとする．
このとき
$$
(R^qf_*\mathcal{F})_{\overline{s}} =
H^q_\etale(\mathbf{P}^1_{\mathcal{O}_{S, \overline{s}}^{sh}},
\mathcal{F}|_{\mathbf{P}^1_{\mathcal{O}_{S, \overline{s}}^{sh}}}) =
H^q_\etale(\mathbf{P}^1_{\kappa(s)^{sep}},
\mathcal{F}|_{\mathbf{P}^1_{\kappa(s)^{sep}}})
$$
である．ここで $\kappa(s)^{sep}$ は
$\mathcal{O}_{S, \overline{s}}^{sh}$ の剰余体，すなわち $\kappa(s)$ の
$k$ における分離代数閉包である．第一の等号は定理
\ref{etale-cohomology-theorem-higher-direct-images} により，第二の等号は前段落の表示式により成り立つ．

\medskip\noindent
最後に，スキームの任意の射 $g : T \to S$ を考える．ここで
$S$ と $\mathcal{F}$ は上のとおりとする．射影を $f' : \mathbf{P}^1_T \to T$ とおき，
$g' : \mathbf{P}^1_T \to \mathbf{P}^1_S$ を $g$ が誘導する射とする．
基底変換写像
$$
g^{-1}R^qf_*\mathcal{F}
\longrightarrow
R^qf'_*(g')^{-1}\mathcal{F}
$$
を考える．$\overline{t}$ を $T$ の幾何学点で，その像が
$\overline{s} = g(\overline{t})$ であるものとする．上の議論により，
$\overline{t}$ における茎上の写像は
$$
H^q_\etale(\mathbf{P}^1_{\kappa(s)^{sep}},
\mathcal{F}|_{\mathbf{P}^1_{\kappa(s)^{sep}}})
\longrightarrow
H^q_\etale(\mathbf{P}^1_{\kappa(t)^{sep}},
\mathcal{F}|_{\mathbf{P}^1_{\kappa(t)^{sep}}})
$$
である．$\kappa(s)^{sep} \subset \kappa(t)^{sep}$ なので，この写像は補題
\ref{etale-cohomology-lemma-base-change-dim-1-separably-closed} により同型である．

\medskip\noindent
これにより，$\mathbf{P}^1_S \to S$ と $\mathbf{Z}/\ell\mathbf{Z}$-加群の層に
対してコホモロジーが基底変換と可換することが証明された．特に，
$\mathbf{Z}/\ell\mathbf{Z}$-加群の単射的層に対しては，任意の基底変換の
高次順像が零である．言い換えると，補題
\ref{etale-cohomology-lemma-proper-base-change-in-terms-of-injectives} の条件 (2) が成り立ち，
証明は完了する．
\end{proof}

\begin{lemma}
\label{etale-cohomology-lemma-proper-base-change}
$f : X \to Y$ をスキームの固有射とする．$g : Y' \to Y$ をスキームの射とする．
$X' = Y' \times_Y X$ とおき，射影を $f' : X' \to Y'$ および
$g' : X' \to X$ と書く．$E \in D^+(X_\etale)$ のコホモロジー層は捩れであるとする．
このとき基底変換写像 (\ref{etale-cohomology-equation-base-change})
$g^{-1}Rf_*E \to Rf'_*(g')^{-1}E$ は同型である．
\end{lemma}

\begin{proof}
これは，スペクトル系列
$$
E_2^{p, q} = R^pf_*H^q(E)
\quad\text{かつ}\quad
{E'}_2^{p, q} = R^pf'_*(g')^{-1}H^q(E)
$$
を用いれば，固有基底変換定理（定理 \ref{etale-cohomology-theorem-proper-base-change}）の
簡単な帰結である．これらは $R^nf_*E$ および
$R^nf'_*(g')^{-1}E$ に収束する．スペクトル系列は Derived Categories の補題
\ref{derived-lemma-two-ss-complex-functor} で構成されている．細部は省略する．
\end{proof}

\begin{lemma}
\label{etale-cohomology-lemma-proper-base-change-stalk}
$f : X \to Y$ をスキームの固有射とする．$\overline{y} \to Y$ を幾何学点とする．
\begin{enumerate}
\item 捩れアーベル層 $\mathcal{F}$ を $X_\etale$ 上にとると，
$(R^nf_*\mathcal{F})_{\overline{y}} =
H^n_\etale(X_{\overline{y}}, \mathcal{F}_{\overline{y}})$ である．
\item $E \in D^+(X_\etale)$ のコホモロジー層が捩れならば，
$(R^nf_*E)_{\overline{y}} =
H^n_\etale(X_{\overline{y}}, E|_{X_{\overline{y}}})$ である．
\end{enumerate}
\end{lemma}

\begin{proof}
主張において，$\mathcal{F}_{\overline{y}}$ は $\mathcal{F}$ のスキーム論的
ファイバー $X_{\overline{y}} = \overline{y} \times_Y X$ への引戻しを表す．
$\overline{y} \to Y$ による引戻しは $\mathcal{F}$ の茎を与えるので，補題の
第一の主張は定理 \ref{etale-cohomology-theorem-proper-base-change} の特殊な場合である．
第二の主張は補題 \ref{etale-cohomology-lemma-proper-base-change} の特殊な場合である．
\end{proof}






\section{固有基底変換の応用}
\label{etale-cohomology-section-applications-proper-base-change}

\noindent
本節では，固有基底変換定理から多少なりとも直ちに従う帰結をいくつか論じる．

\begin{lemma}
\label{etale-cohomology-lemma-base-change-separably-closed}
$K/k$ を分離閉体の拡大とする．$X$ を $k$ 上の固有スキームとする．
$\mathcal{F}$ を $X_\etale$ 上の捩れアーベル層とする．このとき写像
$H^q_\etale(X, \mathcal{F}) \to
H^q_\etale(X_K, \mathcal{F}|_{X_K})$ は $q \geq 0$ に対して同型である．
\end{lemma}

\begin{proof}
茎を調べれば，これは定理 \ref{etale-cohomology-theorem-proper-base-change} の特殊な場合であると分かる．
\end{proof}

\begin{lemma}
\label{etale-cohomology-lemma-cohomological-dimension-proper}
$f : X \to Y$ を，すべてのファイバーの次元が $\leq n$ であるスキームの
固有射とする．このとき任意の捩れアーベル層 $\mathcal{F}$ を $X_\etale$ 上にとると，
$R^qf_*\mathcal{F} = 0$ が $q > 2n$ に対して成り立つ．
\end{lemma}

\begin{proof}
すべての固有射について，$n$ に関する帰納法でこれを証明する．

\medskip\noindent
$n = 0$ ならば，$f$ は有限射であり（More on Morphisms，補題
\ref{more-morphisms-lemma-characterize-finite}），結果は命題
\ref{etale-cohomology-proposition-finite-higher-direct-image-zero} により成り立つ．

\medskip\noindent
$n > 0$ とする．補題 \ref{etale-cohomology-lemma-proper-base-change-stalk} を用いると，
$H^i_\etale(X, \mathcal{F}) = 0$ を，$i > 2n$ で，$X$ が
$\dim(X) \leq n$ を満たす代数閉体 $k$ 上の固有スキームであり，
$\mathcal{F}$ が $X$ 上の捩れアーベル層である場合に証明すれば十分だと分かる．

\medskip\noindent
$n = 1$ ならば，これは定理 \ref{etale-cohomology-theorem-vanishing-curves} から従う．
$n > 1$ と仮定する．命題 \ref{etale-cohomology-proposition-topological-invariance} により，
$X$ をその被約化で置き換えてよい．$\nu : X^\nu \to X$ を正規化とする．
これは全射な双有理有限射であり（Varieties，補題
\ref{varieties-lemma-normalization-locally-algebraic} を参照），したがって稠密開集合
$U \subset X$ 上で同型である（Morphisms，補題
\ref{morphisms-lemma-birational-birational}）．すると
$c : \mathcal{F} \to \nu_*\nu^{-1}\mathcal{F}$ は（$\nu$ が全射なので）単射であり，
$U$ 上で同型である．$i : Z \to X$ を $U$ の補集合の包含と書く．
$U$ は $X$ に稠密なので，$\dim(Z) < \dim(X) = n$ である．命題
\ref{etale-cohomology-proposition-closed-immersion-pushforward} により，
$\Coker(c) = i_*\mathcal{G}$ となるある捩れアーベル層 $\mathcal{G}$ が
$Z_\etale$ 上に存在する．
このとき $H^q_\etale(X, \Coker(c)) = H^q_\etale(Z, \mathcal{G})$ であり
（命題 \ref{etale-cohomology-proposition-finite-higher-direct-image-zero} と Leray スペクトル系列による），
帰納法の仮定から $c$ の余核は次数 $\leq 2(n - 1)$ にのみコホモロジーをもつ．
したがって $\nu_*\nu^{-1}\mathcal{F}$ に対して結果を証明すれば十分である．
$\nu$ は有限なので，これは $H^i_\etale(X^\nu, \nu^{-1}\mathcal{F})$ が
$i > 2n$ に対して零であることに帰着する．
この場合を次段落で扱う．

\medskip\noindent
$X$ は $k$ 上の次元 $n$ の整正規固有スキームであると仮定する．
非定数有理関数 $f$ を $X$ 上に選ぶ．グラフ
$X' \subset X \times \mathbf{P}^1_k$ は $f$ に対する図式
$$
X \xleftarrow{b} X' \xrightarrow{f} \mathbf{P}^1_k
$$
に入る．$b$ はある開部分スキーム $U \subset X$ 上で同型であり，その補集合は
閉部分スキーム $Z \subset X$ で，余次元 $\geq 2$ であることに注意する．実際，
$U$ は $f$ の定義域であり，余次元 $1$ の点を $X$ の中ですべて含む；Morphisms の補題
\ref{morphisms-lemma-rational-map-from-reduced-to-separated} および
\ref{morphisms-lemma-extend-across}（正規性に対する Serre の判定法，Properties の補題
\ref{properties-lemma-criterion-normal} と組み合わせる）を参照されたい．さらに，
$b$ のファイバーは次元 $\leq 1$ である（$\mathbf{P}^1$ の閉部分スキームだからである）．
したがって帰納法により，$R^ib_*b^{-1}\mathcal{F}$ が非零となりうるのは
$i \in \{0, 1, 2\}$ の場合だけである．特別三角形
$$
\mathcal{F} \to Rb_*b^{-1}\mathcal{F} \to Q \to \mathcal{F}[1]
$$
を選ぶ．前と同様に $\mathcal{F} \to b_*b^{-1}\mathcal{F}$ が単射であることと，
今述べたことから，$Q$ の非零なコホモロジー層は次数 $0, 1, 2$ にのみ現れ，
$Z$ 上に台をもつと分かる．さらに補題 \ref{etale-cohomology-lemma-torsion-direct-image} により，
これらのコホモロジー層は捩れである．帰納法と Derived Categories の補題
\ref{derived-lemma-two-ss-complex-functor} のスペクトル系列により，
$H^i(X, Q)$ は $i > 2 + 2\dim(Z) \leq 2 + 2(n - 2) = 2n - 2$ に対して
零である．したがって $H^i(X', b^{-1}\mathcal{F}) = 0$ を $i > 2n$ に対して
証明すれば十分である．ここで射
$$
f : X' \to \mathbf{P}^1_k
$$
を用いる．そのファイバーの次元は $< n$ である．したがって帰納法により，
$R^if_*b^{-1}\mathcal{F} = 0$ が $i > 2(n - 1)$ に対して成り立つ．Leray スペクトル系列
$$
H^i(\mathbf{P}^1_k, R^jf_*b^{-1}\mathcal{F})
\Rightarrow
H^{i + j}(X', b^{-1}\mathcal{F})
$$
と $\dim(\mathbf{P}^1_k) = 1$ であることから結論が従う．
\end{proof}

\noindent
法 $n$ 係数を用いる場合，非有界複体に対しても固有基底変換を行える．

\begin{lemma}
\label{etale-cohomology-lemma-proper-base-change-mod-n}
$f : X \to Y$ をスキームの固有射とする．$g : Y' \to Y$ をスキームの射とする．
$X' = Y' \times_Y X$ とおき，射影を $f' : X' \to Y'$ および
$g' : X' \to X$ と書く．$n \geq 1$ を整数とする．
$E \in D(X_\etale, \mathbf{Z}/n\mathbf{Z})$ とする．このとき基底変換写像
(\ref{etale-cohomology-equation-base-change})，すなわち
$g^{-1}Rf_*E \to Rf'_*(g')^{-1}E$ は同型である．
\end{lemma}

\begin{proof}
$Y$ と $Y'$ が準コンパクトである場合にこれを証明すれば十分である．Morphisms の補題
\ref{morphisms-lemma-morphism-finite-type-bounded-dimension} により，射
$f : X \to Y$ と $f' : X' \to Y'$ のファイバーの次元は有界である．したがって補題
\ref{etale-cohomology-lemma-cohomological-dimension-proper} から，
$$
f_* :
\textit{Mod}(X_\etale, \mathbf{Z}/n\mathbf{Z})
\longrightarrow
\textit{Mod}(Y_\etale, \mathbf{Z}/n\mathbf{Z})
$$
および
$$
f'_* :
\textit{Mod}(X'_\etale, \mathbf{Z}/n\mathbf{Z})
\longrightarrow
\textit{Mod}(Y'_\etale, \mathbf{Z}/n\mathbf{Z})
$$
は Derived Categories の補題 \ref{derived-lemma-unbounded-right-derived} の意味で
有限コホモロジー次元をもつ．K-単射的複体 $\mathcal{I}^\bullet$ を
$\mathbf{Z}/n\mathbf{Z}$-加群の複体として選ぶ．その各項 $\mathcal{I}^n$ は
$\mathbf{Z}/n\mathbf{Z}$-加群の単射的層であり，この複体は $E$ を表すものとする．
Injectives の定理
\ref{injectives-theorem-K-injective-embedding-grothendieck} を参照されたい．通常の
固有基底変換定理により，$R^qf'_*(g')^{-1}\mathcal{I}^n = 0$ が
$q > 0$ に対して成り立つ；定理
\ref{etale-cohomology-theorem-proper-base-change} を参照されたい．したがって Derived Categories の補題
\ref{derived-lemma-unbounded-right-derived} により，$Rf'_*(g')^{-1}E$ は複体
$f'_*(g')^{-1}\mathcal{I}^\bullet$ で計算できる．通常の固有基底変換定理をもう一度
適用すると，これは所望の $g^{-1}f_*\mathcal{I}^\bullet$ に等しい．
\end{proof}

\begin{lemma}
\label{etale-cohomology-lemma-pull-out-constant}
$X$ を準コンパクトかつ準分離なスキームとする．
$E \in D^+(X_\etale)$ および $K \in D^+(\mathbf{Z})$ とする．このとき
$$
R\Gamma(X, E \otimes_\mathbf{Z}^\mathbf{L} \underline{K}) =
R\Gamma(X, E) \otimes_\mathbf{Z}^\mathbf{L} K
$$
\end{lemma}

\begin{proof}
$H^i(E) = 0$ が $i \geq a$ に対して成り立ち，$H^j(K) = 0$ が
$j \geq b$ に対して成り立つとする．$K$ は下に有界な複体 $K^\bullet$ で表せ，
その各項は捩れなし $\mathbf{Z}$-加群である．（K-平坦複体 $L^\bullet$ を選んで
$K$ を表し，
$K^\bullet = \tau_{\geq b - 1}L^\bullet$ とすればよい．これは $\mathbf{Z}$ の
大域次元が $1$ だから可能である．More on Algebra，補題
\ref{more-algebra-lemma-last-one-flat} を参照されたい．）$E$ は下に有界な複体
$\mathcal{E}^\bullet$ で表せる．このとき
$E \otimes_\mathbf{Z}^\mathbf{L} \underline{K}$ は
$$
\text{Tot}(\mathcal{E}^\bullet \otimes_\mathbf{Z} \underline{K}^\bullet)
$$
で表される．特別三角形
$$
\sigma_{\geq -b + n + 1}K^\bullet
\to K^\bullet \to
\sigma_{\leq -b + n}K^\bullet
$$
と自明な消滅
$$
H^n(X,
\text{Tot}(\mathcal{E}^\bullet \otimes_\mathbf{Z}
\sigma_{\geq -a + n + 1}\underline{K}^\bullet) = 0
$$
および
$$
H^n(R\Gamma(X, E)
\otimes_\mathbf{Z}^\mathbf{L} \sigma_{\geq -a + n + 1}K^\bullet) = 0
$$
を用いると，$K^\bullet$ が平坦 $\mathbf{Z}$-加群からなる有界複体である場合に
帰着する．この議論を繰り返すと，$K^\bullet$ がある次数に置かれた一つの平坦
$\mathbf{Z}$-加群に等しい場合に帰着する．次に $\mathcal{E}^\bullet$ の素朴切断を
用いると，全く同様に，$\mathcal{E}^\bullet$ がある次数に置かれた一つの
アーベル層である場合に帰着する．したがって
$$
H^n(X, \mathcal{E} \otimes_\mathbf{Z} \underline{M}) =
H^n(X, \mathcal{E}) \otimes_\mathbf{Z} M
$$
を，$M$ が平坦 $\mathbf{Z}$-加群で，$\mathcal{E}$ が $X$ 上のアーベル層である場合に
示せば十分である．この場合，$M$ を有限自由 $\mathbf{Z}$-加群のフィルター余極限として
書く（ラザールの定理；Algebra，定理 \ref{algebra-theorem-lazard} を参照）．定理
\ref{etale-cohomology-theorem-colimit} により，有限自由 $\mathbf{Z}$-加群 $M$ の場合に帰着し，
この場合は結果は自明である．
\end{proof}

\begin{lemma}
\label{etale-cohomology-lemma-projection-formula-proper}
$f : X \to Y$ をスキームの固有射とする．
$E \in D^+(X_\etale)$ のコホモロジー層は捩れであるとする．
$K \in D^+(Y_\etale)$ とする．このとき
$$
Rf_*E \otimes_\mathbf{Z}^\mathbf{L} K =
Rf_*(E \otimes_\mathbf{Z}^\mathbf{L} f^{-1}K)
$$
$D^+(Y_\etale)$ において成り立つ．
\end{lemma}

\begin{proof}
Cohomology on Sites，節 \ref{sites-cohomology-section-projection-formula} により，
左辺から右辺への標準写像がある．この等式を茎上で確認する．導来テンソル積の計算は
引戻しと可換することを思い出そう．Cohomology on Sites，補題
\ref{sites-cohomology-lemma-pullback-tensor-product} を参照されたい．したがって
$$
(E \otimes_\mathbf{Z}^\mathbf{L} f^{-1}K)_{\overline{x}}
=
E_{\overline{x}} \otimes_\mathbf{Z}^\mathbf{L} K_{\overline{y}}
$$
である．ここで $\overline{y}$ は $\overline{x}$ の $Y$ における像である．
$\mathbf{Z}$ の大域次元は $1$ なので，この複体の次数 $< - 1 + a + b$ の
コホモロジーは，$H^i(E) = 0$ が $i \geq a$ に対して成り立ち，
$H^j(K) = 0$ が $j \geq b$ に対して成り立つならば零である．さらに $H^i(E)$ は
各 $i$ について
捩れアーベル層なので，複体 $E \otimes_\mathbf{Z}^\mathbf{L} K$ のコホモロジー層も
同じ性質をもつ．実際，
$$
(E \otimes_\mathbf{Z}^\mathbf{L} f^{-1}K)
\otimes_{\mathbf{Z}}^\mathbf{L} \mathbf{Q} =
(E \otimes_\mathbf{Z}^\mathbf{L} \mathbf{Q})
\otimes_{\mathbf{Q}}^\mathbf{L}
(f^{-1}K \otimes_{\mathbf{Z}}^\mathbf{L} \mathbf{Q})
$$
であり，これは導来圏において零である．以上から，補題
\ref{etale-cohomology-lemma-proper-base-change-stalk} を両辺に適用でき，
$$
R\Gamma(X_{\overline{y}},
E|_{X_{\overline{y}}} \otimes_\mathbf{Z}^\mathbf{L}
(X_{\overline{y}} \to \overline{y})^{-1}K_{\overline{y}}) =
R\Gamma(X_{\overline{y}},
E|_{X_{\overline{y}}}) \otimes_\mathbf{Z}^\mathbf{L} K_{\overline{y}}
$$
を示せば十分であると分かる．これは補題 \ref{etale-cohomology-lemma-pull-out-constant} で示されている．
\end{proof}








\section{局所非輪状性}
\label{etale-cohomology-section-local-acyclicity}

\noindent
本節では，滑らかな射の局所非輪状性を滑らかな基底変換定理から導く．
SGA 4 または SGA 4.5 では，著者らはまず滑らかな射に対する局所非輪状性の一形態を
証明し，そこから滑らかな基底変換定理を導いている．

\medskip\noindent
Deligne \cite[Definition 2.12, page 242]{SGA4.5} による局所非輪状性の定式化を用いる．
$f : X \to S$ をスキームの射とする．$\overline{x}$ を $X$ の幾何学点で，
$\overline{s} = f(\overline{x})$ をその $S$ における像とする．
$\overline{t}$ を $\Spec(\mathcal{O}^{sh}_{S, \overline{s}})$ の幾何学点とする．
可換図式
$$
\xymatrix{
F_{\overline{x}, \overline{t}} =
\overline{t} \times_{\Spec(\mathcal{O}^{sh}_{S, \overline{s}})}
\Spec(\mathcal{O}^{sh}_{X, \overline{x}}) \ar[r] \ar[d] &
\Spec(\mathcal{O}^{sh}_{X, \overline{x}}) \ar[r] \ar[d] &
X \ar[d] \\
\overline{t} \ar[r] &
\Spec(\mathcal{O}^{sh}_{S, \overline{s}}) \ar[r] &
S
}
$$
を得る．スキーム $F_{\overline{x}, \overline{t}}$ を
{\it $f$ の $\overline{x}$ における消滅サイクル多様体}という．
$K$ を $D(X_\etale)$ の対象とする．スキームの任意の射 $g : Y\to X$ に対し，
$R\Gamma(Y, K)$ と書いて $R\Gamma(Y_\etale, g_{small}^{-1}K)$ を表す．
$\mathcal{O}^{sh}_{X, \overline{x}}$ は狭義 Hensel 環なので，
$K_{\overline{x}} = R\Gamma(\Spec(\mathcal{O}^{sh}_{X, \overline{x}}), K)$ である．
したがって標準写像
\begin{equation}
\label{etale-cohomology-equation-alpha-K}
\alpha_{K, \overline{x}, \overline{t}} :
K_{\overline{x}}
\longrightarrow
R\Gamma(F_{\overline{x}, \overline{t}}, K)
\end{equation}
を，$F_{\overline{x}, \overline{t}} \to
\Spec(\mathcal{O}^{sh}_{X, \overline{x}})$ に沿ってコホモロジーを引き戻すことで得る．

\begin{definition}
\label{etale-cohomology-definition-locally-acyclic}
\begin{reference}
\cite[Definition 2.12, page 242]{SGA4.5} および
\cite[Definition (1.3), page 54]{SGA4.5}
\end{reference}
$f : X \to S$ をスキームの射とする．$K$ を $D(X_\etale)$ の対象とする．
\begin{enumerate}
\item
$\overline{x}$ を $X$ の幾何学点で，像が $\overline{s} = f(\overline{x})$ であるものとする．
$f$ が {\it $\overline{x}$ で $K$ に関して局所非輪状}であるとは，すべての幾何学点
$\overline{t}$ を $\Spec(\mathcal{O}^{sh}_{S, \overline{s}})$ にとったとき，
写像 (\ref{etale-cohomology-equation-alpha-K}) が同型であることをいう\footnote{
$\overline{t}$ が $\Spec(\mathcal{O}^{sh}_{S, \overline{s}})$ の代数的幾何学点であるとは
仮定しない．補題 \ref{etale-cohomology-lemma-smooth-base-change-separably-closed} を用いて，しばしば
この場合に帰着できる．}．
\item $f$ が {\it $K$ に関して局所非輪状}であるとは，$f$ が
$\overline{x}$ で $K$ に関して局所非輪状であることが，すべての幾何学点
$\overline{x}$ を $X$ にとったときに成り立つことをいう．
\item $f$ が {\it $K$ に関して普遍的局所非輪状}であるとは，スキームの任意の射
$S' \to S$ に対して，基底変換 $f' : X' \to S'$ が $K$ の $X'$ への引戻しに
関して局所非輪状であることをいう．
\item $f$ が {\it 局所非輪状}であるとは，すべての幾何学点 $\overline{x}$ を
$X$ にとり，任意の整数 $n$ が $\kappa(\overline{x})$ の標数と互いに素であるとき，
射 $f$ が $\overline{x}$ で値 $\mathbf{Z}/n\mathbf{Z}$ の定数層に関して
局所非輪状であることをいう．
\item $f$ が {\it 普遍的局所非輪状}であるとは，スキームの任意の射
$S' \to S$ に対して基底変換 $f' : X' \to S'$ が局所非輪状であることをいう．
\end{enumerate}
\end{definition}

\noindent
$M$ をアーベル群とする．このとき，$f : X \to S$ の，定数層
$\underline{M}$ に関する局所非輪状性は，
$$
H^q(F_{\overline{x}, \overline{t}}, \underline{M}) =
\left\{
\begin{matrix}
M & \text{ならば} & q = 0 \\
0 & \text{ならば} & q \not = 0
\end{matrix}
\right.
$$
が，任意の幾何学点 $\overline{x}$ を $X$ にとり，任意の幾何学点
$\overline{t}$ を $\Spec(\mathcal{O}^{sh}_{S, f(\overline{x})})$ にとったときに
成り立つという条件に帰着する．このように，局所非輪状であることは，幾何学的ファイバー
$F_{\overline{x}, \overline{t}}$，すなわち $\overline{x}$ と $\overline{s}$ における
狭義 Hensel 化の間の写像のファイバーの高次コホモロジー群が消滅することに対応する．

\begin{proposition}
\label{etale-cohomology-proposition-smooth-locally-acyclic}
$f : X \to S$ をスキームの滑らかな射とする．このとき $f$ は普遍的局所非輪状である．
\end{proposition}

\begin{proof}
滑らかな射の基底変換は滑らかなので，滑らかな射が局所非輪状であることを示せば
十分である．$\overline{x}$ を $X$ の幾何学点で，像が
$\overline{s} = f(\overline{x})$ であるものとする．$\overline{t}$ を
$\Spec(\mathcal{O}^{sh}_{S, f(\overline{x})})$ の幾何学点とする．環準同型
$\mathcal{O}^{sh}_{S, \overline{s}} \to \mathcal{O}^{sh}_{X, \overline{x}}$ の性質を
証明しようとしているので（定義 \ref{etale-cohomology-definition-locally-acyclic} に続く議論を参照），
$f : X \to S$ を基底変換
$X \times_S \Spec(\mathcal{O}^{sh}_{S, \overline{s}}) \to
\Spec(\mathcal{O}^{sh}_{S, \overline{s}})$ で置き換えてよいし，そうする．したがって
$S$ は狭義 Hensel 局所環のスペクトルで，$\overline{s}$ は $S$ の閉点上にあると
仮定してよいし，そうする．

\medskip\noindent
補題 \ref{etale-cohomology-lemma-base-change-f-star-general-stalks} を図式
$$
\xymatrix{
X \ar[d]_f & X_{\overline{t}} \ar[l]^h \ar[d]^e \\
S & \overline{t} \ar[l]_g
}
$$
および層 $\mathcal{F} = \underline{M}$ に適用する．ここで
$M = \mathbf{Z}/n\mathbf{Z}$，かつ $n$ は $\overline{x}$ の剰余体の標数と互いに素な
ある整数とする．滑らかな基底変換により，写像
$f^{-1}R^qg_*\mathcal{F} \to R^qh_*e^{-1}\mathcal{F}$ は同型である；定理
\ref{etale-cohomology-theorem-smooth-base-change} を参照されたい（捩れに関する仮定は $n$ の選び方により
満たされる）．したがって補題 \ref{etale-cohomology-lemma-base-change-f-star-general-stalks} により，
$$
H^q(F_{\overline{x}, \overline{t}}, \underline{M}) =
H^q(\Spec(\mathcal{O}^{sh}_{X, \overline{x}}) \times_S \overline{t},
\underline{M})
=
H^q(\Spec(\mathcal{O}^{sh}_{S, \overline{s}}) \times_S \overline{t},
\underline{M}) =
H^q(\overline{t}, \underline{M})
$$
の中央の等号を得る．外側の二つの等号には
$S = \Spec(\mathcal{O}^{sh}_{S, \overline{s}})$ を用いる．
$\overline{t}$ は分離閉体のスペクトルなので，
$$
H^q(F_{\overline{x}, \overline{t}}, \underline{M}) =
\left\{
\begin{matrix}
M & \text{ならば} & q = 0 \\
0 & \text{ならば} & q \not = 0
\end{matrix}
\right.
$$
を得る．これは示すべきことであった（定義
\ref{etale-cohomology-definition-locally-acyclic} に続く議論を参照）．
\end{proof}

\begin{lemma}
\label{etale-cohomology-lemma-locally-acyclic-locally-constant}
$f : X \to S$ をスキームの射とする．$\mathcal{F}$ を $X_\etale$ 上の局所定数
アーベル層とし，すべての幾何学点 $\overline{x}$ を $X$ にとったときにアーベル群
$\mathcal{F}_{\overline{x}}$ は捩れ群で，そのすべての元の位数が $\overline{x}$ の
剰余体の標数と互いに素であるとする．$f$ が局所非輪状ならば，$f$ は
$\mathcal{F}$ に関して局所非輪状である．
\end{lemma}

\begin{proof}
実際，$\overline{x}$ を $X$ の幾何学点とする．$\mathcal{F}$ は局所定数なので，
$\mathcal{F}$ の $\Spec(\mathcal{O}^{sh}_{X, \overline{x}})$ への制限は，定数層
$\underline{M}$ と同型である．ここで $M = \mathcal{F}_{\overline{x}}$ とする．仮定により，
$M = \colim M_i$ と書ける．ここでこれは，有限アーベル群 $M_i$ のフィルター余極限で，
各群の位数は $\overline{x}$ の剰余体の標数と互いに素である．
$\overline{t}$ を $\Spec(\mathcal{O}^{sh}_{S, f(\overline{x})})$ の幾何学点とする．
$F_{\overline{x}, \overline{t}}$ はアフィンなので，
$$
H^q(F_{\overline{x}, \overline{t}}, \underline{M}) =
\colim H^q(F_{\overline{x}, \overline{t}}, \underline{M_i})
$$
が補題 \ref{etale-cohomology-lemma-colimit} により成り立つ．各 $i$ に対して
$M_i = \bigoplus \mathbf{Z}/n_{i, j}\mathbf{Z}$ と有限直和に書ける．ここで整数
$n_{i, j}$ は $\overline{x}$ の剰余体の標数と互いに素である．$f$ は局所非輪状なので，
$$
H^q(F_{\overline{x}, \overline{t}},
\underline{\mathbf{Z}/n_{i, j}\mathbf{Z}}) =
\left\{
\begin{matrix}
\mathbf{Z}/n_{i, j}\mathbf{Z} & \text{ならば} & q = 0 \\
0 & \text{ならば} & q \not = 0
\end{matrix}
\right.
$$
である．定義 \ref{etale-cohomology-definition-locally-acyclic} に続く議論を参照されたい．
直和と余極限をとると，
$$
H^q(F_{\overline{x}, \overline{t}}, \underline{M}) =
\left\{
\begin{matrix}
M & \text{ならば} & q = 0 \\
0 & \text{ならば} & q \not = 0
\end{matrix}
\right.
$$
を得て，証明が完了する．
\end{proof}

\begin{lemma}
\label{etale-cohomology-lemma-locally-acyclic-quasi-finite-base-change}
図式
$$
\xymatrix{
X' \ar[r]_{g'} \ar[d]_{f'} & X \ar[d]^f \\
S' \ar[r]^g & S
}
$$
をスキームのファイバー積図式とする．$K$ を $D(X_\etale)$ の対象とする．
$\overline{x}'$ を $X'$ の幾何学点とし，$\overline{x}$ をその $X$ における像とする．
次を仮定する：
\begin{enumerate}
\item $f$ は $\overline{x}$ で $K$ に関して局所非輪状である；
\item $g$ は局所準有限である，または $S' = \lim S_i$ は $S$ 上局所準有限な
スキームのアフィン遷移射をもつ有向逆極限である，または $g : S' \to S$ は整である．
\end{enumerate}
このとき $f'$ は $\overline{x}'$ で $(g')^{-1}K$ に関して局所非輪状である．
\end{lemma}

\begin{proof}
$\overline{s}'$ および $\overline{s}$ を，それぞれ $\overline{x}'$ および
$\overline{x}$ の $S'$ および $S$ における像とする．$\overline{t}'$ を
$\Spec(\mathcal{O}^{sh}_{S', \overline{s}'})$ の幾何学点とし，$\overline{t}$ を
$\Spec(\mathcal{O}^{sh}_{S, \overline{s}})$ におけるその像とする．Algebra の補題
\ref{algebra-lemma-base-change-strict-henselization-quasi-finite} と $g$ に関する仮定により，
$$
\mathcal{O}^{sh}_{X, \overline{x}}
\otimes_{\mathcal{O}^{sh}_{S, \overline{s}}}
\mathcal{O}^{sh}_{S', \overline{s}'}
\longrightarrow
\mathcal{O}^{sh}_{X', \overline{x}'}
$$
は同型である．我々の規約により $\kappa(\overline{t}) = \kappa(\overline{t}')$ なので，
$$
F_{\overline{x}', \overline{t}'} =
\Spec\left(
\mathcal{O}^{sh}_{X', \overline{x}'}
\otimes_{\mathcal{O}^{sh}_{S', \overline{s}'}}
\kappa(\overline{t}')\right) =
\Spec\left(
\mathcal{O}^{sh}_{X, \overline{x}}
\otimes_{\mathcal{O}^{sh}_{S, \overline{s}}}
\kappa(\overline{t})\right) =
F_{\overline{x}, \overline{t}}
$$
を得る．言い換えると，$f'$ の $\overline{x}'$ における消滅サイクル多様体は，
$f$ の $\overline{x}$ における消滅サイクル多様体の例である．本補題はこれと定義から
直ちに従う．
\end{proof}










\section{余特殊化写像}
\label{etale-cohomology-section-cospecialization}

\noindent
$f : X \to S$ をスキームの射とする．$\overline{x}$ を $X$ の幾何学点で，
$\overline{s} = f(\overline{x})$ をその $S$ における像とする．$\overline{t}$ を
$\Spec(\mathcal{O}^{sh}_{S, \overline{s}})$ の幾何学点とする．
$K \in D(X_\etale)$ とする．スキームの任意の射 $g : Y \to X$ に対し，
$K|_Y$ と書いて $g_{small}^{-1}K$ を表し，$R\Gamma(Y, K)$ と書いて
$R\Gamma(Y_\etale, g_{small}^{-1}K)$ を表す．
次を仮定すると，
\begin{enumerate}
\item $K$ は下に有界，すなわち $K \in D^+(X_\etale)$ である；
\item $f$ は $K$ に関して局所非輪状である．
\end{enumerate}
{\it 余特殊化写像}
$$
cosp :
R\Gamma(X_{\overline{t}}, K)
\longrightarrow
R\Gamma(X_{\overline{s}}, K)
$$
が存在すると主張する．これは節 \ref{etale-cohomology-section-specialization}，とりわけ注意
\ref{etale-cohomology-remark-specialization-map-and-fibres} で考えた特殊化写像と密接に関係する．

\medskip\noindent
この写像を構成するため，射
$$
X_{\overline{t}}
\xrightarrow{h}
X \times_S \Spec(\mathcal{O}^{sh}_{S, \overline{s}})
\xleftarrow{i}
X_{\overline{s}}
$$
を考える．$h^{-1}$ と $Rh_*$ の随伴の単位は写像
$$
\beta_{K, \overline{s}, \overline{t}} :
K|_{X \times_S \Spec(\mathcal{O}^{sh}_{S, \overline{s}})}
\longrightarrow
Rh_*(K|_{X_{\overline{t}}})
$$
を $D((X \times_S \Spec(\mathcal{O}^{sh}_{S, \overline{s}}))_\etale)$ において与える．
下の補題 \ref{etale-cohomology-lemma-beta-is-isomorphism} により，上の仮定のもとで引戻し
$i^{-1}\beta_{K, \overline{s}, \overline{t}}$ は同型である．したがって余特殊化写像を
合成
\begin{align*}
R\Gamma(X_{\overline{t}}, K)
& =
R\Gamma(X \times_S \Spec(\mathcal{O}^{sh}_{S, \overline{s}}),
Rh_*(K|_{X_{\overline{t}}})) \\
&
\xrightarrow{i^{-1}}
R\Gamma(X_{\overline{s}}, i^{-1}Rh_*(K|_{X_{\overline{t}}})) \\
&
\xrightarrow{(i^{-1}\beta_{K, \overline{s}, \overline{t}})^{-1}}
R\Gamma(X_{\overline{s}},
i^{-1}(K|_{X \times_S \Spec(\mathcal{O}^{sh}_{S, \overline{s}})})) \\
& =
R\Gamma(X_{\overline{s}}, K)
\end{align*}
として定義できる．

\begin{lemma}
\label{etale-cohomology-lemma-beta-is-isomorphism}
写像 $i^{-1}\beta_{K, \overline{s}, \overline{t}}$ は同型である．
\end{lemma}

\begin{proof}
写像 $h$, $i$, $\beta_{K, \overline{s}, \overline{t}}$ の構成は，$X$ と $K$ の
$\Spec(\mathcal{O}^{sh}_{S, \overline{s}})$ への基底変換だけに依存する．したがって
$S$ は閉点 $\overline{s}$ をもつ狭義 Hensel スキームであると仮定してよいし，そうする．
$f$ の $K$ に関する局所非輪状性はこの基底変換で保たれることに注意する（例えば補題
\ref{etale-cohomology-lemma-locally-acyclic-quasi-finite-base-change} によるか，この非常に特殊な場合に
狭義 Hensel 環を直接比較すればよい）．

\medskip\noindent
$\overline{x}$ を $X_{\overline{s}}$ の幾何学点とする．同じことだが，
$\overline{x}$ を $f$ による像が $\overline{s}$ である幾何学点とする．
$i^{-1}\beta_{K, \overline{s}, \overline{t}}$ の $\overline{x}$ における茎を計算しよう．
まず，
$$
(i^{-1}\beta_{K, \overline{s}, \overline{t}})_{\overline{x}} =
(\beta_{K, \overline{s}, \overline{t}})_{\overline{x}}
$$
である．引戻しは茎を保つからである；補題 \ref{etale-cohomology-lemma-stalk-pullback} を参照されたい．
$S = \Spec(\mathcal{O}^{sh}_{S, \overline{s}})$ という状況なので，
$h : X_{\overline{t}} \to X$ は
$X_{\overline{t}} \times_X \Spec(\mathcal{O}^{sh}_{X, \overline{x}})
= F_{\overline{x}, \overline{t}}$ を満たす．したがって
$$
(\beta_{K, \overline{s}, \overline{t}})_{\overline{x}} :
K_{\overline{x}}
\longrightarrow
Rh_*(K|_{X_{\overline{t}}})_{\overline{x}} =
R\Gamma(F_{\overline{x}, \overline{t}}, K)
$$
を満たすと分かる．ここで等号は定理 \ref{etale-cohomology-theorem-higher-direct-images} による．
したがって写像 $(\beta_{K, \overline{s}, \overline{t}})_{\overline{x}}$ は，定義
\ref{etale-cohomology-definition-locally-acyclic} で用いた写像
$\alpha_{K, \overline{x}, \overline{t}}$ にほかならない．写像が同型かどうかは茎で
確認できるので，定理 \ref{etale-cohomology-theorem-exactness-stalks} により結果が従う．
\end{proof}

\noindent
余特殊化写像が存在するとき，それは特殊化写像の逆写像となることを目指す．

\begin{lemma}
\label{etale-cohomology-lemma-specialization-cospecialization}
上の状況で，さらに $f$ が準コンパクトかつ準分離ならば，図式
$$
\xymatrix{
(Rf_*K)_{\overline{s}} \ar[r] \ar[d]_{sp} &
R\Gamma(X_{\overline{s}}, K) \\
(Rf_*K)_{\overline{t}} \ar[r] &
R\Gamma(X_{\overline{t}}, K) \ar[u]_{cosp}
}
$$
は可換である．
\end{lemma}

\begin{proof}
補題 \ref{etale-cohomology-lemma-beta-is-isomorphism} の証明と同様，$S$ を
$\Spec(\mathcal{O}^{sh}_{S, \overline{s}})$ で置き換えてよい．すると写像は
$h : X_{\overline{t}} \to X$, $i : X_{\overline{s}} \to X$, および
$\beta_{K, \overline{s}, \overline{t}} : K \to Rh_*(K|_{X_{\overline{t}}})$ に簡約される．
定理 \ref{etale-cohomology-theorem-higher-direct-images} により
$(Rf_*K)_{\overline{s}} = R\Gamma(X, K)$ であることを用いると，$sp$ と基底変換写像
$(Rf_*K)_{\overline{t}} \to R\Gamma(X_{\overline{t}}, K)$ の合成は，$h$ に沿う
コホモロジーの引戻しにほかならない．これは写像
$$
R\Gamma(X, K) \xrightarrow{\beta_{K, \overline{s}, \overline{t}}}
R\Gamma(X, Rh_*(K|_{X_{\overline{t}}})) =
R\Gamma(X_{\overline{t}}, K)
$$
と同じである．いま写像 $cosp$ は，まずこの表示式の $=$ を逆向きに用い，次に $i$ に
沿って引き戻し，最後に $i^{-1}\beta_{K, \overline{s}, \overline{t}}$ の逆を適用する．
したがって所望の可換性を得る．
\end{proof}

\begin{lemma}
\label{etale-cohomology-lemma-sp-isom-proper-torsion-loc-ac}
$f : X \to S$ をスキームの射とする．$K \in D(X_\etale)$ とする．次を仮定する：
\begin{enumerate}
\item $K$ は下に有界，すなわち $K \in D^+(X_\etale)$ である；
\item $f$ は $K$ に関して局所非輪状である；
\item $f$ は固有である；
\item $K$ のコホモロジー層は捩れである．
\end{enumerate}
このとき，すべての幾何学点 $\overline{s}$ を $S$ にとり，すべての幾何学点
$\overline{t}$ を $\Spec(\mathcal{O}^{sh}_{S, \overline{s}})$ にとると，
特殊化写像
$sp : (Rf_*K)_{\overline{s}} \to (Rf_*K)_{\overline{t}}$
および余特殊化写像
$cosp : R\Gamma(X_{\overline{t}}, K) \to R\Gamma(X_{\overline{s}}, K)$
はともに同型である．
\end{lemma}

\begin{proof}
固有基底変換定理（補題 \ref{etale-cohomology-lemma-proper-base-change-stalk} の形）により，
$(Rf_*K)_{\overline{s}} = R\Gamma(X_{\overline{s}}, K)$ であり，$\overline{t}$ についても
同様である．「正しい」証明は，補題 \ref{etale-cohomology-lemma-specialization-cospecialization} の議論が，
この場合に $sp$ と $cosp$ が互いに逆な同型であることを示すと確認するものであろう．
代わりに，$cosp$ が同型であることを直接示す．上の議論から，$cosp$ が同型であることと，
$i$ による引戻し
$$
R\Gamma(X \times_S \Spec(\mathcal{O}^{sh}_{S, \overline{s}}),
Rh_*(K|_{X_{\overline{t}}}))
\longrightarrow
R\Gamma(X_{\overline{s}}, i^{-1}Rh_*(K|_{X_{\overline{t}}}))
$$
が $D^+(\textit{Ab})$ において同型であることは同値である．これは，固有射
$f' : X \times_S \Spec(\mathcal{O}^{sh}_{S, \overline{s}}) \to
\Spec(\mathcal{O}^{sh}_{S, \overline{s}})$
に，射 $\overline{s} \to \Spec(\mathcal{O}^{sh}_{S, \overline{s}})$ と複体
$K' = Rh_*(K|_{X_{\overline{t}}})$ に関する固有基底変換定理を適用すれば従う．複体
$K'$ は補題 \ref{etale-cohomology-lemma-torsion-direct-image} により下に有界で，捩れコホモロジー層をもつ．
$\Spec(\mathcal{O}^{sh}_{S, \overline{s}})$ は狭義 Hensel で，$\overline{s}$ は閉点上に
あるので，表示した矢印の始域は $(Rf'_*K')_{\overline{s}}$ に，終域は
$R\Gamma(X_{\overline{s}}, K')$ に等しく，既に用いた補題
\ref{etale-cohomology-lemma-proper-base-change-stalk} により表示写像は同型である．したがって補題
\ref{etale-cohomology-lemma-specialization-cospecialization} の図式の四本の矢印のうち三本が同型であり，
結論が従う．
\end{proof}

\begin{lemma}
\label{etale-cohomology-lemma-sp-isom-proper-loc-cst-torsion}
$f : X \to S$ をスキームの射とする．$\mathcal{F}$ を $X_\etale$ 上の
アーベル層とする．次を仮定する：
\begin{enumerate}
\item $f$ は滑らかかつ固有である；
\item $\mathcal{F}$ は局所定数である；
\item $\mathcal{F}_{\overline{x}}$ は捩れ群で，そのすべての元の位数が
$\overline{x}$ の剰余標数と互いに素である．これはすべての幾何学点
$\overline{x}$ を $X$ にとったときに成り立つ．
\end{enumerate}
このとき，すべての幾何学点 $\overline{s}$ を $S$ にとり，すべての幾何学点
$\overline{t}$ を $\Spec(\mathcal{O}^{sh}_{S, \overline{s}})$ にとると，特殊化写像
$sp : (Rf_*\mathcal{F})_{\overline{s}} \to (Rf_*\mathcal{F})_{\overline{t}}$
は同型である．
\end{lemma}

\begin{proof}
これは補題 \ref{etale-cohomology-lemma-sp-isom-proper-torsion-loc-ac} および
\ref{etale-cohomology-lemma-locally-acyclic-locally-constant} と，命題
\ref{etale-cohomology-proposition-smooth-locally-acyclic} から従う．
\end{proof}

















\section{コホモロジー次元}
\label{etale-cohomology-section-cd}

\noindent
節 \ref{etale-cohomology-section-vanishing-torsion} の結果と，一見無害な固有基底変換定理の一回の適用
（命題 \ref{etale-cohomology-proposition-cd-affine} の証明中）を用いて，スキームおよび体の
コホモロジー次元に対するいくつかの上界を導ける．

\begin{definition}
\label{etale-cohomology-definition-cd}
$X$ を準コンパクトかつ準分離なスキームとする．
{\it $X$ のコホモロジー次元}とは，最小の元
$$
\text{cd}(X) \in \{0, 1, 2, \ldots\} \cup \{\infty\}
$$
で，任意の捩れアーベル層 $\mathcal{F}$ を $X_\etale$ 上にとったときに
$H^i_\etale(X, \mathcal{F}) = 0$ が $i > \text{cd}(X)$ に対して成り立つものをいう．
$X = \Spec(A)$ ならば，これを $A$ のコホモロジー次元と呼ぶこともある．
\end{definition}

\noindent
スキームの標数が $p$ ならば，$p$-冪捩れ層のコホモロジーの消滅について，しばしば
より鋭い上界が得られる．これは別の箇所で扱う（将来の参照をここに挿入する）．

\begin{lemma}
\label{etale-cohomology-lemma-cd-limit}
$X = \lim X_i$ を，アフィン遷移射をもつ準コンパクトかつ準分離なスキームの系の
有向極限とする．このとき $\text{cd}(X) \leq \max \text{cd}(X_i)$ である．
\end{lemma}

\begin{proof}
$f_i : X \to X_i$ を射影と書く．$\mathcal{F}$ を $X_\etale$ 上の捩れアーベル層とする．
補題 \ref{etale-cohomology-lemma-linus-hamann} により
$\mathcal{F} = \lim f_i^{-1}f_{i, *}\mathcal{F}$ である．したがって定理
\ref{etale-cohomology-theorem-colimit} により $H^q_\etale(X, \mathcal{F}) =
\colim H^q_\etale(X_i, f_{i, *}\mathcal{F})$ であり，本補題が従う．
\end{proof}

\begin{lemma}
\label{etale-cohomology-lemma-cd-curve-over-field}
$K$ を体とする．$X$ を $1$ 次元アフィンスキームで $K$ 上有限型のものとする．このとき
$\text{cd}(X) \leq 1 + \text{cd}(K)$ である．
\end{lemma}

\begin{proof}
$\mathcal{F}$ を $X_\etale$ 上の捩れアーベル層とする．射
$f : X \to \Spec(K)$ に対する Leray スペクトル系列を考えると，
$$
E_2^{p, q} = H^p(\Spec(K), R^qf_*\mathcal{F})
$$
を得て，これは $H^{p + q}_\etale(X, \mathcal{F})$ に収束する．
$R^qf_*\mathcal{F}$ の幾何学点 $\Spec(\overline{K}) \to \Spec(K)$ における茎は，
$\mathcal{F}$ の $X_{\overline{K}}$ への引戻しのコホモロジーである．したがって定理
\ref{etale-cohomology-theorem-vanishing-affine-curves} により，次数 $\geq 2$ で消滅する．
\end{proof}

\begin{lemma}
\label{etale-cohomology-lemma-cd-field-extension}
$L/K$ を体の拡大とする．このとき
$\text{cd}(L) \leq \text{cd}(K) + \text{trdeg}_K(L)$ である．
\end{lemma}

\begin{proof}
$\text{trdeg}_K(L) = \infty$ ならば明らかである．そうでなければ，拡大の列
$L= L_r/L_{r - 1}/ \ldots /L_1/L_0 = K$ で，
$\text{trdeg}_{L_i}(L_{i + 1}) = 1$ かつ $r = \text{trdeg}_K(L)$ となるものをとれる．
したがって $r = 1$ の場合に本補題を証明すれば十分である．この場合，$L = \colim A_i$
と，有限型 $K$-部分代数のフィルター余極限に書ける．補題 \ref{etale-cohomology-lemma-cd-limit} により，
$\text{cd}(A_i) \leq 1 + \text{cd}(K)$ を証明すれば十分であり，これは補題
\ref{etale-cohomology-lemma-cd-curve-over-field} から従う．
\end{proof}

\begin{lemma}
\label{etale-cohomology-lemma-strictly-henselian}
$K$ を体とする．$X$ を $K$ 上有限型のスキームとする．$x \in X$ とする．
$a = \text{trdeg}_K(\kappa(x))$ および $d = \dim_x(X)$ とおく．このとき写像
$$
K(t_1, \ldots, t_a)^{sep} \longrightarrow \mathcal{O}_{X, x}^{sh}
$$
で，次を満たすものが存在する：
\begin{enumerate}
\item $\mathcal{O}_{X, x}^{sh}$ の剰余体は
$K(t_1, \ldots, t_a)^{sep}$ の純非分離拡大である；
\item $\mathcal{O}_{X, x}^{sh}$ は，有限型 $K(t_1, \ldots, t_a)^{sep}$-代数で
次元 $\leq d - a$ のもののフィルター余極限である．
\end{enumerate}
\end{lemma}

\begin{proof}
$X$ はアフィンであると仮定してよい．Noether 正規化により，必要なら再び $X$ を
縮小した後，有限射 $\pi : X \to \mathbf{A}^d_K$ を選べる；Algebra，補題
\ref{algebra-lemma-Noether-normalization-at-point} を参照されたい．$\kappa(x)$ は
$\pi(x)$ の剰余体の有限拡大なので，この剰余体も超越次数 $a$ を $K$ 上にもつ．
したがって有限射 $\pi' : \mathbf{A}^d_K \to \mathbf{A}^d_K$ で，
$\pi'(\pi(x))$ が，線形部分空間 $\mathbf{A}^a_K \subset \mathbf{A}^d_K$ で，
最後の $d - a$ 個の座標を零とおいて得られるものの一般点に対応するようにできる．ゆえに合成
$$
X \xrightarrow{\pi' \circ \pi} \mathbf{A}^d_K \xrightarrow{p} \mathbf{A}^a_K
$$
$\pi' \circ \pi$ と射影 $p$，すなわち最初の $a$ 個の座標への射影との合成は，$x$ を一般点
$\eta \in \mathbf{A}^a_K$ に写す．誘導される 'etale 局所環の写像
$$
K(t_1, \ldots, t_a)^{sep} =
\mathcal{O}_{\mathbf{A}^a_k, \eta}^{sh}
\longrightarrow \mathcal{O}_{X, x}^{sh}
$$
は (1) を満たす．実際，$\mathcal{O}_{X, x}^{sh}$ の剰余体が分離閉体
$K(t_1, \ldots, t_a)^{sep}$ の代数拡大であることは明らかである．他方，
$X = \Spec(B)$ ならば，$\mathcal{O}_{X, x}^{sh} = \colim B_j$ は 'etale
$B$-代数 $B_j$ のフィルター余極限である．$B_j$ は
$K[t_1, \ldots, t_d]$ 上準有限であることに注意する．これは $B$ が
$K[t_1, \ldots, t_d]$ 上有限だからである．同様に
$K(t_1, \ldots, t_a)^{sep} = \colim A_i$ と，'etale
$K[t_1, \ldots, t_a]$-代数のフィルター余極限に書ける．各 $i$ に対し，ある $j$ が存在して
$A_i \to K(t_1, \ldots, t_a)^{sep} \to \mathcal{O}_{X, x}^{sh}$
は写像 $\psi_{i, j} : A_i \to B_j$ を経由する．このとき $B_j$ は
$A_i[t_{a + 1}, \ldots, t_d]$ 上準有限である．したがって
$$
B_{i, j} = B_j \otimes_{\psi_{i, j}, A_i} K(t_1, \ldots, t_a)^{sep}
$$
は次元 $\leq d - a$ である．なぜなら
$K(t_1, \ldots, t_a)^{sep}[t_{a + 1}, \ldots, t_d]$ 上準有限だからである．
$\mathcal{O}_{X, x}^{sh}$ はフィルター余極限\footnote{$R$ を環とする．
$A = \colim_{i \in I} A_i$ を有限表示
$R$-代数のフィルター余極限とする．$B = \colim_{j \in J} B_j$ を $R$-代数の
フィルター余極限とする．$A \to B$ を $R$-代数準同型とする．すべての
$i \in I$ に対し，ある $j \in J$ と $R$-代数準同型
$\psi_{i, j} : A_i \to B_j$ が存在すると仮定する．
$(i', j', \psi_{i', j'}) \geq (i, j, \psi_{i, j})$ とは，$i' \geq i$，
$j' \geq j$ で，$\psi_{i, j}$ と $\psi_{i', j'}$ が両立することとする．このとき
三つ組の集まりは有向集合をなし，$B = \colim B_j \otimes_{\psi_{i, j} A_i} A$ である．}
として代数 $B_{i, j}$ から得られるので，これで (2) の証明が完了する．
いくつかの細部を省略する．
\end{proof}

\begin{proposition}
\label{etale-cohomology-proposition-cd-affine}
$K$ を体とする．$X$ を $K$ 上有限型のアフィンスキームとする．このとき
$\text{cd}(X) \leq \dim(X) + \text{cd}(K)$ である．
\end{proposition}

\begin{proof}
$\dim(X)$ に関する帰納法で証明する．$\mathcal{F}$ を $X_\etale$ 上の
捩れアーベル層とする．

\medskip\noindent
$\dim(X) = 0$ の場合．このとき構造射 $f : X \to \Spec(K)$ は有限である．
したがって $R^if_*\mathcal{F} = 0$ が $i > 0$ に対して成り立つ；命題
\ref{etale-cohomology-proposition-finite-higher-direct-image-zero} を参照されたい．ゆえに
$H^i_\etale(X, \mathcal{F}) = H^i_\etale(\Spec(K), f_*\mathcal{F})$ である．これは
$f$ に対する Leray スペクトル系列（Cohomology on Sites，補題
\ref{sites-cohomology-lemma-Leray}）によるので，結果は明らかである．

\medskip\noindent
$\dim(X) = 1$ の場合．これは補題 \ref{etale-cohomology-lemma-cd-curve-over-field} である．

\medskip\noindent
$d = \dim(X) > 1$ とし，命題は体上の次元 $< d$ の有限型アフィンスキームに対して
成り立つと仮定する．Noether 正規化により，例えば Varieties の補題
\ref{varieties-lemma-noether-normalization-affine} を参照すると，有限射
$f : X \to \mathbf{A}^d_K$ が存在する．命題
\ref{etale-cohomology-proposition-finite-higher-direct-image-zero} により，$R^if_*\mathcal{F} = 0$ が
$i > 0$ に対して成り立つことを思い出そう．$f$ に対する Leray スペクトル系列
（Cohomology on Sites，補題 \ref{sites-cohomology-lemma-Leray}）により，
$\pi_*\mathcal{F}$ に対して $\mathbf{A}^d_K$ 上で結果を証明すれば十分である．

\medskip\noindent
挿話 I．$j : X \to Y$ を，滑らかな $d$ 次元多様体で $K$ 上のもの
（必ずしもアフィンでない）の
開埋込みで，補集合が有効 Cartier 因子 $D$ の台であるものとする．層
$R^qj_*\mathcal{F}$ は $q > 0$ ならば $D$ 上に台をもつ．
$(R^qj_*\mathcal{F})_{\overline{y}} = 0$ が
$a = \text{trdeg}_K(\kappa(y)) > d - q$ に対して成り立つと主張する．実際，定理
\ref{etale-cohomology-theorem-higher-direct-images} により
$$
(R^qj_*\mathcal{F})_{\overline{y}} =
H^q(\Spec(\mathcal{O}_{Y, y}^{sh}) \times_Y X, \mathcal{F})
$$
局所方程式 $f \in \mathfrak m_y \subset \mathcal{O}_{Y, y}$ を $D$ に対して選ぶ．このとき
$$
\Spec(\mathcal{O}_{Y, y}^{sh}) \times_Y X =
\Spec(\mathcal{O}_{Y, y}^{sh}[1/f])
$$
補題 \ref{etale-cohomology-lemma-strictly-henselian} を用いると埋込み
$$
K(t_1, \ldots, t_a)^{sep}(x) =
K(t_1, \ldots, t_a)^{sep}[x]_{(x)}[1/x]
\longrightarrow
\mathcal{O}_{Y, y}^{sh}[1/f]
$$
を得る．$K$ 上で見た $\mathcal{O}_{Y, y}^{sh}$ の分数体の超越次数は $d$ なので，
$\mathcal{O}_{Y, y}^{sh}[1/f]$ は $(d - a - 1)$ 次元有限型代数のフィルター余極限で，
その係数体は $K(t_1, \ldots, t_a)^{sep}(x)$ である．この体自身の
コホモロジー次元は補題 \ref{etale-cohomology-lemma-cd-field-extension} により $1$ である．したがって
帰納法の仮定と補題 \ref{etale-cohomology-lemma-cd-limit} から所望の消滅を得る．

\medskip\noindent
挿話 II．$Z$ を $K$ 上の次元 $d - 1$ の滑らかな多様体とする．$E_a \subset Z$ を，
点 $z \in Z$ で $\text{trdeg}_K(\kappa(z)) \leq a$ を満たすものの集合とする．$E_a$ は
特殊化について閉じていることに注意する；Varieties，補題
\ref{varieties-lemma-dimension-locally-algebraic} を参照されたい．$\mathcal{G}$ を $Z$ 上の
捩れアーベル層で，その台が $E_a$ に含まれるものとする．このとき
$H^b_\etale(Z, \mathcal{G}) = 0$ が $b > a + \text{cd}(K)$ に対して成り立つと主張する．
実際，$\mathcal{G} = \colim \mathcal{G}_i$ と書ける．ここで $\mathcal{G}_i$ は，
閉部分スキーム $Z_i$ で $E_a$ に含まれるものの上に台をもつ捩れアーベル層である；補題
\ref{etale-cohomology-lemma-support-in-subset} を参照されたい．すると帰納法の仮定により
$\mathcal{G}_i$ に対する所望の消滅が従う\footnote{ここではまず命題
\ref{etale-cohomology-proposition-closed-immersion-pushforward} を用いて $\mathcal{G}_i$ を $Z_i$ 上の層の
順像として書き，帰納法の仮定からこの層の $Z_i$ 上での消滅を得て，最後に
$Z_i \to Z$ に対する Leray スペクトル系列から $\mathcal{G}_i$ の消滅を得る．}．
最後に定理 \ref{etale-cohomology-theorem-colimit} により結論する．

\medskip\noindent
可換図式
$$
\xymatrix{
\mathbf{A}^d_K \ar[rd]_f \ar[rr]_-j & &
\mathbf{P}^1_K \times_K \mathbf{A}^{d - 1}_K \ar[ld]^g \\
& \mathbf{A}^{d - 1}_K
}
$$
を考える．$j$ は滑らかな $d$ 次元多様体の開埋込みで，その補集合は有効 Cartier 因子
$D$ であることに注意する．したがって挿話 I で得た結果を使える．相対 Leray
スペクトル系列
$$
E_2^{p, q} = R^pg_*R^qj_*\mathcal{F} \Rightarrow R^{p + q}f_*\mathcal{F}
$$
を調べる．$R^qj_*\mathcal{F}$ は $q > 0$ ならば $D$ 上に台をもち，
$g|_D : D \to \mathbf{A}^{d - 1}_K$ は同型なので，
$R^pg_*R^qj_*\mathcal{F} = 0$ が $p > 0$ かつ $q > 0$ に対して成り立つ．さらに
$R^qj_*\mathcal{F} = 0$ が $q > d$ に対して成り立つ．他方，$g$ は相対次元 $1$ の
固有射である．したがって補題 \ref{etale-cohomology-lemma-cohomological-dimension-proper} により，
$R^pg_*j_*\mathcal{F} = 0$ が $p > 2$ に対して成り立つ．よってスペクトル系列の
$E_2$-頁は次の形をもつ：
$$
\begin{matrix}
g_*R^dj_*\mathcal{F} & 0 & 0 \\
\ldots & \ldots & \ldots \\
g_*R^2j_*\mathcal{F} & 0 & 0 \\
g_*R^1j_*\mathcal{F} & 0 & 0 \\
g_*j_*\mathcal{F} & R^1g_*j_*\mathcal{F} & R^2g_*j_*\mathcal{F}
\end{matrix}
$$
$R^qf_*\mathcal{F} = g_*R^qj_*\mathcal{F}$ が $q > 2$ に対して成り立つと結論する．
挿話 I により，$R^qf_*\mathcal{F}$ の台は $q > 2$ ならば
$\mathbf{A}^{d - 1}_K$ の点のうち，その剰余体の超越次数が $\leq d - q$ であるものの
集合に含まれる．挿話 II により
$$
H^p(\mathbf{A}^{d - 1}_K, R^qf_*\mathcal{F}) = 0
\text{ ただし }p > d - q + \text{cd}(K)\text{ かつ }q > 2
$$
である．他方，定理 \ref{etale-cohomology-theorem-higher-direct-images} により
$R^2f_*\mathcal{F}_{\overline{\eta}} =
H^2(\mathbf{A}^1_{\overline{\eta}}, \mathcal{F}) = 0$ である（次元 $1$ の場合による消滅）．
ここで $\eta$ は $\mathbf{A}^{d - 1}_K$ の一般点である．したがって挿話 II を再び用いると
$$
H^p(\mathbf{A}^{d - 1}_K, R^2f_*\mathcal{F}) = 0
\text{ ただし }p > d - 2 + \text{cd}(K)
$$
を得る．最後に
$$
H^p(\mathbf{A}^{d - 1}_K, R^qf_*\mathcal{F}) = 0
\text{ ただし }p > d - 1 + \text{cd}(K)\text{ かつ }q = 0, 1
$$
が帰納法の仮定により成り立つ．以上を Leray スペクトル系列
$H^p(\mathbf{A}^{d - 1}_K, R^qf_*\mathcal{F}) \Rightarrow
H^{p + q}(\mathbf{A}^d_K, \mathcal{F})$ と組み合わせて結論を得る．
\end{proof}

\begin{lemma}
\label{etale-cohomology-lemma-interlude-II}
$K$ を体とする．$X$ を $K$ 上有限型のアフィンスキームとする．$E_a \subset X$ を，
点 $x \in X$ で $\text{trdeg}_K(\kappa(x)) \leq a$ を満たすものの集合とする．
$\mathcal{F}$ を $X_\etale$ 上の捩れアーベル層で，その台が $E_a$ に含まれるものとする．
このとき $H^b_\etale(X, \mathcal{F}) = 0$ が $b > a + \text{cd}(K)$ に対して成り立つ．
\end{lemma}

\begin{proof}
$\mathcal{F} = \colim \mathcal{F}_i$ と書ける．ここで $\mathcal{F}_i$ は，閉部分スキーム
$Z_i$ で $E_a$ に含まれるものの上に台をもつ捩れアーベル層である；補題
\ref{etale-cohomology-lemma-support-in-subset} を参照されたい．すると命題 \ref{etale-cohomology-proposition-cd-affine} は
$\mathcal{F}_i$ に対する所望の消滅を与える．細部を省略する．ヒント：まず命題
\ref{etale-cohomology-proposition-closed-immersion-pushforward} を用いて $\mathcal{F}_i$ を $Z_i$ 上の層の
順像として書き，この層の $Z_i$ 上での消滅を用い，さらに $Z_i \to Z$ に対する
Leray スペクトル系列を用いて $\mathcal{F}_i$ の消滅を得る．最後に定理
\ref{etale-cohomology-theorem-colimit} により結論する．
\end{proof}

\begin{lemma}
\label{etale-cohomology-lemma-interlude-I}
$f : X \to Y$ を体 $K$ 上有限型のスキームのアフィン射とする．$E_a(X)$ を，
点 $x \in X$ で $\text{trdeg}_K(\kappa(x)) \leq a$ を満たすものの集合とする．
$\mathcal{F}$ を $X_\etale$ 上の捩れアーベル層で，その台が $E_a$ に含まれるものとする．
このとき $R^qf_*\mathcal{F}$ の台は $E_{a - q}(Y)$ に含まれる．
\end{lemma}

\begin{proof}
問題は $Y$ 上局所的なので，$Y$ はアフィンであると仮定してよい．すると $X$ も
アフィンであり，図式
$$
\xymatrix{
X \ar[d]_f \ar[r]_i & \mathbf{A}^{n + m}_K \ar[d]^{\text{pr}} \\
Y \ar[r]^j & \mathbf{A}^n_K
}
$$
を選べる．ここで水平矢印は閉埋込み，右側の垂直矢印は射影である（細部は省略）．
$j_*R^qf_*\mathcal{F} = R^q\text{pr}_*i_*\mathcal{F}$ である．これは $i$ と $j$ の
高次順像が消滅するからである；命題
\ref{etale-cohomology-proposition-finite-higher-direct-image-zero} を参照されたい．さらに，命題中の
$j_*$ の茎の記述から，$j_*R^qf_*\mathcal{F}$ に対して消滅を証明すれば十分である．
したがって $f$ は射影 $\text{pr} : \mathbf{A}^{n + m}_K \to \mathbf{A}^n_K$ であり，
$\mathcal{F}$ は $\mathbf{A}^{n + m}_K$ 上の捩れアーベル層で，本補題の主張中の仮定を
満たすと仮定してよい．

\medskip\noindent
$y$ を $\mathbf{A}^n_K$ の点とする．定理 \ref{etale-cohomology-theorem-higher-direct-images} により
$$
(R^q\text{pr}_*\mathcal{F})_{\overline{y}} =
H^q(\mathbf{A}^{n + m}_K \times_{A^n_K} \Spec(\mathcal{O}_{Y, y}^{sh}),
\mathcal{F}) =
H^q(\mathbf{A}^m_{\mathcal{O}_{Y, y}^{sh}}, \mathcal{F})
$$
$b = \text{trdeg}_K(\kappa(y))$ とする．補題 \ref{etale-cohomology-lemma-strictly-henselian} から埋込み
$$
L = K(t_1, \ldots, t_b)^{sep} \longrightarrow \mathcal{O}_{Y, y}^{sh}
$$
$\mathcal{O}_{Y, y}^{sh} = \colim B_i$ と書く．これは有限型 $L$-部分代数
$B_i \subset \mathcal{O}_{Y, y}^{sh}$ で，正則関数の環
$K[T_1, \ldots, T_n]$，すなわち $\mathbf{A}^n_K$ 上の正則関数環を含むものの
フィルター余極限である．このとき
$$
\mathbf{A}^m_{\mathcal{O}_{Y, y}^{sh}} =
\lim \mathbf{A}^m_{B_i}
$$
$z \in \mathbf{A}^m_{B_i}$ が $\mathcal{F}$ の台に属する点ならば，像 $x$，すなわち
$z$ の $\mathbf{A}^{m + n}_K$ における像は，
$\text{trdeg}_K(\kappa(x)) \leq a$ を満たす．これは本補題の $\mathcal{F}$ に関する
仮定による．$\mathcal{O}_{Y, y}^{sh}$ は
$K[T_1, \ldots, T_n]$ 上の \'etale 代数のフィルター余極限であり，かつ
$B_i \subset \mathcal{O}_{Y, y}^{sh}$ なので，$\kappa(z)/\kappa(x)$ は代数的である
（いくつかの細部を省略する）．したがって
$\text{trdeg}_K(\kappa(z)) \leq a$，ゆえに
$\text{trdeg}_L(\kappa(z)) \leq a - b$ である．補題
\ref{etale-cohomology-lemma-interlude-II} により
$$
H^q(\mathbf{A}^m_{B_i}, \mathcal{F}) = 0\text{，ただし }q > a - b
$$
したがって定理 \ref{etale-cohomology-theorem-colimit} により，所望のとおり
$(Rf_*\mathcal{F})_{\overline{y}} = 0$ が $q > a - b$ に対して成り立つ．
\end{proof}








\section{有限コホモロジー次元}
\label{etale-cohomology-section-finite-cd}

\noindent
節 \ref{etale-cohomology-section-cd} で始めた議論を続ける．

\begin{definition}
\label{etale-cohomology-definition-cd-f}
$f : X \to Y$ をスキームの準コンパクトかつ準分離な射とする．
{\it $f$ のコホモロジー次元}とは，最小の元
$$
\text{cd}(f) \in \{0, 1, 2, \ldots\} \cup \{\infty\}
$$
で，任意の捩れアーベル層 $\mathcal{F}$ を $X_\etale$ 上にとったときに
$R^if_*\mathcal{F} = 0$ が $i > \text{cd}(f)$ ならば成り立つものをいう．
\end{definition}

\begin{lemma}
\label{etale-cohomology-lemma-finite-cd}
$K$ を体とする．
\begin{enumerate}
\item $f : X \to Y$ が $K$ 上有限型のスキームの射ならば，
$\text{cd}(f) < \infty$ である．
\item $\text{cd}(K) < \infty$ ならば，$\text{cd}(X) < \infty$ は，
任意の有限型スキーム $X$ が $K$ 上定義されているとき成り立つ．
\end{enumerate}
\end{lemma}

\begin{proof}
(1) の証明．$Y$ はアフィンであると仮定してよい．これを証明するため，
スキームのコホモロジー，補題 \ref{coherent-lemma-induction-principle} の
帰納原理を用いる．$X$ もアフィンならば，結果は補題
\ref{etale-cohomology-lemma-interlude-I} により成り立つ．したがって，$X = U \cup V$ であり，
$U \to Y$，$V \to Y$，および $U \cap V \to Y$ に対して結果が成り立つならば，
$f$ に対しても成り立つことを示せば十分である．これは相対 Mayer--Vietoris 完全列から従う；
補題 \ref{etale-cohomology-lemma-relative-mayer-vietoris} を参照されたい．

\medskip\noindent
(2) の証明．これを証明するため，スキームのコホモロジー，補題
\ref{coherent-lemma-induction-principle} の帰納原理を用いる．$X$ がアフィンならば，
結果は命題 \ref{etale-cohomology-proposition-cd-affine} により成り立つ．したがって，
$X = U \cup V$ であり，$U$，$V$，および $U \cap V $ に対して結果が成り立つならば，
$X$ に対しても成り立つことを示せば十分である．これは Mayer--Vietoris 完全列から従う；
補題 \ref{etale-cohomology-lemma-mayer-vietoris} を参照されたい．
\end{proof}

\begin{lemma}
\label{etale-cohomology-lemma-finite-cd-mod-n-direct-sums}
コホモロジーと直和．$n \geq 1$ を整数とする．
\begin{enumerate}
\item $f : X \to Y$ を，$\text{cd}(f) < \infty$ を満たすスキームの
準コンパクトかつ準分離な射とする．このとき関手
$$
Rf_* :
D(X_\etale, \mathbf{Z}/n\mathbf{Z})
\longrightarrow
D(Y_\etale, \mathbf{Z}/n\mathbf{Z})
$$
は直和と可換である．
\item $X$ を，$\text{cd}(X) < \infty$ を満たす準コンパクトかつ準分離な
スキームとする．このとき関手
$$
R\Gamma(X, -) :
D(X_\etale, \mathbf{Z}/n\mathbf{Z})
\longrightarrow
D(\mathbf{Z}/n\mathbf{Z})
$$
は直和と可換である．
\end{enumerate}
\end{lemma}

\begin{proof}
(1) の証明．$\text{cd}(f) < \infty$ なので，
$$
f_* :
\textit{Mod}(X_\etale, \mathbf{Z}/n\mathbf{Z})
\longrightarrow
\textit{Mod}(Y_\etale, \mathbf{Z}/n\mathbf{Z})
$$
は導来圏，補題 \ref{derived-lemma-unbounded-right-derived} の意味で
有限コホモロジー次元をもつ．$I$ を集合とし，各 $i \in I$ に対して
$E_i$ を $D(X_\etale, \mathbf{Z}/n\mathbf{Z})$ の対象とする．
$\mathcal{I}_i^\bullet$ を $\mathbf{Z}/n\mathbf{Z}$-加群の K-単射的複体で，
各項 $\mathcal{I}_i^n$ が $\mathbf{Z}/n\mathbf{Z}$-加群の単射的層であり，
$E_i$ を表すものとして選ぶ．
単射的対象，定理 \ref{injectives-theorem-K-injective-embedding-grothendieck}
を参照されたい．このとき $\bigoplus E_i$ は複体
$\bigoplus \mathcal{I}_i^\bullet$（各項ごとの直和）で表される；単射的対象，補題
\ref{injectives-lemma-derived-products} を参照されたい．補題
\ref{etale-cohomology-lemma-relative-colimit} により
$$
R^qf_*(\bigoplus \mathcal{I}_i^n) =
\bigoplus R^qf_*(\mathcal{I}_i^n) = 0
$$
が $q > 0$ および任意の $n$ に対して成り立つ．したがって導来圏，補題
\ref{derived-lemma-unbounded-right-derived} により，$Rf_*(\bigoplus E_i)$ を複体
$$
f_*(\bigoplus \mathcal{I}_i^\bullet) =
\bigoplus f_*(\mathcal{I}_i^\bullet)
$$
で計算できる（等号は再び補題 \ref{etale-cohomology-lemma-relative-colimit} による）．すでに用いた
単射的対象，補題 \ref{injectives-lemma-derived-products} により，この複体は
$\bigoplus Rf_*E_i$ を表す．

\medskip\noindent
(2) の証明は (1) の証明と同一なので省略する．
\end{proof}

\begin{lemma}
\label{etale-cohomology-lemma-proper-mod-n-direct-sums}
$f : X \to Y$ をスキームの固有射とする．$n \geq 1$ を整数とする．このとき関手
$$
Rf_* :
D(X_\etale, \mathbf{Z}/n\mathbf{Z})
\longrightarrow
D(Y_\etale, \mathbf{Z}/n\mathbf{Z})
$$
は直和と可換である．
\end{lemma}

\begin{proof}
$Y$ が準コンパクトな場合に証明すれば十分である．射，補題
\ref{morphisms-lemma-morphism-finite-type-bounded-dimension} により，
$f : X \to Y$ のファイバーの次元は有界である．したがって補題
\ref{etale-cohomology-lemma-cohomological-dimension-proper} により $\text{cd}(f) < \infty$ である．ゆえに
結果は補題 \ref{etale-cohomology-lemma-finite-cd-mod-n-direct-sums} から従う．
\end{proof}

\begin{lemma}
\label{etale-cohomology-lemma-pull-out-constant-mod-n}
$X$ を，$\text{cd}(X) < \infty$ を満たす準コンパクトかつ準分離なスキームとする．
$\Lambda$ を捩れ環とする．$E \in D(X_\etale, \Lambda)$ および
$K \in D(\Lambda)$ とする．このとき
$$
R\Gamma(X, E \otimes_\Lambda^\mathbf{L} \underline{K}) =
R\Gamma(X, E) \otimes_\Lambda^\mathbf{L} K
$$
\end{lemma}

\begin{proof}
サイト上のコホモロジー，節 \ref{sites-cohomology-section-projection-formula} により，
右辺から左辺への標準的な写像がある．$T(K)$ を，$K \in D(\Lambda)$ に対して
補題の主張が成り立つという性質とする．結論を得るため，代数学詳論，注意
\ref{more-algebra-remark-P-resolution} の条件 (1)，(2)，(3) が $T$ に対して
成り立つことを確かめる．性質 (1) は等式の両辺が直和と可換であることから従う；
補題 \ref{etale-cohomology-lemma-finite-cd-mod-n-direct-sums} を参照されたい．性質 (2) は，
三角圏の間の完全関手を比較しており，導来圏，補題
\ref{derived-lemma-third-isomorphism-triangle} を用いられることから従う．
性質 (3) は，$K = \Lambda[k]$ のとき補題が任意のシフト $k \in \mathbf{Z}$ に対して
成り立つという主張であり，これは明らかである．
\end{proof}

\begin{lemma}
\label{etale-cohomology-lemma-projection-formula-proper-mod-n}
$f : X \to Y$ をスキームの固有射とする．$\Lambda$ を捩れ環とする．
$E \in D(X_\etale, \Lambda)$ および $K \in D(Y_\etale, \Lambda)$ とする．このとき
$$
Rf_*E \otimes_\Lambda^\mathbf{L} K =
Rf_*(E \otimes_\Lambda^\mathbf{L} f^{-1}K)
$$
が $D(Y_\etale, \Lambda)$ において成り立つ．
\end{lemma}

\begin{proof}
サイト上のコホモロジー，節 \ref{sites-cohomology-section-projection-formula} により，
左辺から右辺への標準的な写像がある．$\overline{y}$ における茎上で等式を確かめる．
固有基底変換（$Y' = \overline{y}$ とした補題
\ref{etale-cohomology-lemma-proper-base-change-mod-n} の形）により，$Y$ が代数閉体のスペクトルである場合に
帰着する．これは補題 \ref{etale-cohomology-lemma-pull-out-constant-mod-n} で示されている．そこでは補題
\ref{etale-cohomology-lemma-cohomological-dimension-proper} による $\text{cd}(X) < \infty$ を用いる．
\end{proof}






\section{\'etale コホモロジーにおける K\"unneth 公式}
\label{etale-cohomology-section-kunneth}

\noindent
まず，一方の因子が固有である場合の K\"unneth 公式を証明する．次に，
この公式を用いて開埋込みに対する基底変換の性質を証明する．これから，
（滑らかな基底変換に似た）「体のスペクトルへ向かう射による基底変換」が得られる．
最後に，これを用いてより一般的な K\"unneth 公式を得る．

\begin{remark}
\label{etale-cohomology-remark-define-kunneth-map}
スキームの圏における Cartesian 図式
$$
\xymatrix{
X \times_S Y \ar[d]_p \ar[r]_q \ar[rd]_c & Y \ar[d]^g \\
X \ar[r]^f & S
}
$$
$\Lambda$ を環とし，$E \in D(X_\etale, \Lambda)$ および
$K \in D(Y_\etale, \Lambda)$ とする．このとき標準的な写像
$$
Rf_*E \otimes_\Lambda^\mathbf{L} Rg_*K
\longrightarrow
Rc_*(p^{-1}E \otimes_\Lambda^\mathbf{L} q^{-1}K)
$$
がある．例えば，標準的な写像 $Rf_*E \to Rc_*p^{-1}E$ と
$Rg_*K \to Rc_*q^{-1}K$，およびサイト上のコホモロジー，注意
\ref{sites-cohomology-remark-cup-product} で定義された相対カップ積を用いて
これを定義できる．あるいは，写像
$$
c^{-1}(Rf_*E \otimes_\Lambda^\mathbf{L} Rg_*K)
=
p^{-1}f^{-1}Rf_*E \otimes_\Lambda^\mathbf{L} q^{-1} g^{-1}Rg_*K
\to
p^{-1}E \otimes_\Lambda^\mathbf{L} q^{-1}K
$$
の随伴を用いてもよい．この写像には随伴写像 $f^{-1}Rf_*E \to E$ および
$g^{-1}Rg_*K \to K$ を用いる．
\end{remark}


\begin{lemma}
\label{etale-cohomology-lemma-kunneth-one-proper}
$k$ を分離閉体とする．$X$ を $k$ 上の固有スキームとする．
$Y$ を $k$ 上の準コンパクトかつ準分離なスキームとする．
\begin{enumerate}
\item $E \in D^+(X_\etale)$ のコホモロジー層が捩れ層であり，
$K \in D^+(Y_\etale)$ ならば，
$$
R\Gamma(X \times_{\Spec(k)} Y,
\text{pr}_1^{-1}E
\otimes_\mathbf{Z}^\mathbf{L}
\text{pr}_2^{-1}K
)
=
R\Gamma(X, E)
\otimes_\mathbf{Z}^\mathbf{L}
R\Gamma(Y, K)
$$
\item $n \geq 1$ が整数で，$Y$ が $k$ 上有限型であり，
$E \in D(X_\etale, \mathbf{Z}/n\mathbf{Z})$ および
$K \in D(Y_\etale, \mathbf{Z}/n\mathbf{Z})$ ならば，
$$
R\Gamma(X \times_{\Spec(k)} Y,
\text{pr}_1^{-1}E
\otimes_{\mathbf{Z}/n\mathbf{Z}}^\mathbf{L}
\text{pr}_2^{-1}K
)
=
R\Gamma(X, E)
\otimes_{\mathbf{Z}/n\mathbf{Z}}^\mathbf{L}
R\Gamma(Y, K)
$$
\end{enumerate}
\end{lemma}

\begin{proof}
(1) の証明．補題 \ref{etale-cohomology-lemma-projection-formula-proper} により
$$
R\text{pr}_{2, *}(
\text{pr}_1^{-1}E
\otimes_\mathbf{Z}^\mathbf{L}
\text{pr}_2^{-1}K) =
R\text{pr}_{2, *}(\text{pr}_1^{-1}E)
\otimes_\mathbf{Z}^\mathbf{L}
K
$$
である．固有基底変換（補題 \ref{etale-cohomology-lemma-proper-base-change} の形）により，これは
$$
\underline{R\Gamma(X, E)}
\otimes_\mathbf{Z}^\mathbf{L}
K
$$
という $D(Y_\etale)$ の対象に等しい．この対象に $R\Gamma(Y, -)$ を適用すると，
$\text{pr}_2$ に対する Leray スペクトル系列により (1) の等式の左辺が得られる．
したがって補題 \ref{etale-cohomology-lemma-pull-out-constant} から結論が従う．

\medskip\noindent
(2) の証明．補題 \ref{etale-cohomology-lemma-projection-formula-proper-mod-n}，
\ref{etale-cohomology-lemma-proper-base-change-mod-n}，\ref{etale-cohomology-lemma-pull-out-constant-mod-n}，および
補題 \ref{etale-cohomology-lemma-finite-cd} による $\text{cd}(Y) < \infty$ を用いる点を除けば，
(1) の証明とまったく同じである．
\end{proof}

\begin{lemma}
\label{etale-cohomology-lemma-supported-in-closed-points}
$K$ を分離閉体とする．$X$ を $K$ 上有限型のスキームとする．
$\mathcal{F}$ を $X_\etale$ 上のアーベル層で，その台が $X$ の閉点全体に
含まれるものとする．このとき $H^q(X, \mathcal{F}) = 0$ が $q > 0$ に対して
成り立ち，$\mathcal{F}$ は大域切断で生成される．
\end{lemma}

\begin{proof}
（$\mathcal{F}$ が捩れ層ならば，消滅は補題 \ref{etale-cohomology-lemma-interlude-II} から直ちに従う．）
補題 \ref{etale-cohomology-lemma-support-in-subset} により，$\mathcal{F}$ は構成可能層
$\mathcal{F}_i$ のフィルター余極限として書ける．これらは $\mathbf{Z}$-加群の層で，その台 $Z_i \subset X$ が
閉点の有限集合となるものとして書ける．命題
\ref{etale-cohomology-proposition-closed-immersion-pushforward} により，そのような層は
$(Z_i \to X)_*\mathcal{G}_i$ の形である．ここで $\mathcal{G}_i$ は $Z_i$ 上の層である．
$K$ は分離閉なので，スキーム $Z_i$ は分離閉体のスペクトルの有限非交和である．
$H^q(Z_i, \mathcal{G}_i) = H^q(X, \mathcal{F}_i)$ であることを思い出そう．これは
$Z_i \to X$ に対する Leray スペクトル系列と，この射の高次順像の消滅
（命題 \ref{etale-cohomology-proposition-finite-higher-direct-image-zero}）による．
補題 \ref{etale-cohomology-lemma-equivalence-abelian-sheaves-point} および
\ref{etale-cohomology-lemma-compare-cohomology-point} により，$H^q(Z_i, \mathcal{G}_i)$ は
$q > 0$ に対して零であり，$H^0(Z_i, \mathcal{G}_i)$ は $\mathcal{G}_i$ を生成する．
したがって $H^q(X, \mathcal{F}_i)$ は $q > 0$ に対して消滅し，
$\mathcal{F}_i$ は大域切断で生成される．定理 \ref{etale-cohomology-theorem-colimit} により，
$H^q(X, \mathcal{F}) = 0$ が $q > 0$ に対して成り立つ．大域切断で生成される
アーベル群の層のフィルター余極限も大域切断で生成されるので，これで証明は完了する
（詳細は省略する）．
\end{proof}

\begin{lemma}
\label{etale-cohomology-lemma-vanishing-closed-points}
$K$ を分離閉体とする．$X$ を $K$ 上有限型のスキームとする．
$Q \in D(X_\etale)$ とする．$Q_{\overline{x}}$ が非零であるのは，
$x$ が $X$ の閉点である場合に限ると仮定する．このとき
$$
Q = 0 \Leftrightarrow H^i(X, Q) = 0 \text{（任意の）}i
$$
\end{lemma}

\begin{proof}
左から右への含意は自明である．したがって逆の含意を証明すればよい．

\medskip\noindent
$Q$ が下に有界であると仮定する；ほとんどすべての応用にはこの場合で十分である．
$Q$ が零でなければ，最小の $i$ で，コホモロジー層 $H^i(Q)$ が非零となるものを考えられる．
補題 \ref{etale-cohomology-lemma-supported-in-closed-points} により
$H^i(X, Q) = H^0(X, H^i(Q)) \not = 0$ であり，結論を得る．

\medskip\noindent
一般の場合．$\mathcal{B} \subset \Ob(X_\etale)$ を準コンパクトな対象全体とする．
補題 \ref{etale-cohomology-lemma-supported-in-closed-points} により，サイト上のコホモロジー，補題
\ref{sites-cohomology-lemma-cohomology-over-U-trivial} の仮定が満たされる．したがって
$H^q(U, Q) = H^0(U, H^q(Q))$ がすべての $U \in \mathcal{B}$ に対して成り立つ．
特に，$U = X$ に対して成り立つ．ゆえに補題
\ref{etale-cohomology-lemma-supported-in-closed-points} から結論が従う．実際，$Q$ が
$D(X_\etale)$ において零であることと，$H^q(Q)$ がすべての $q$ に対して
零であることとは同値である．
\end{proof}

\begin{lemma}
\label{etale-cohomology-lemma-kunneth-localize-on-X}
$K$ を体とする．$j : U \to X$ を $K$ 上有限型のスキームの開埋込みとする．
$Y$ を $K$ 上有限型のスキームとする．図式
$$
\xymatrix{
Y \times_{\Spec(K)} X \ar[d]_q  &
Y \times_{\Spec(K)} U \ar[l]^h \ar[d]^p \\
X & U \ar[l]_j
}
$$
を考える．基底変換写像 $q^{-1}Rj_*\mathcal{F} \to Rh_*p^{-1}\mathcal{F}$ は，
$\mathcal{F}$ が $U_\etale$ 上のアーベル層で，その茎が $K$ で可逆な位数の
捩れ群であるならば，同型である．
\end{lemma}

\begin{proof}
$\mathcal{F} = \colim \mathcal{F}[n]$ と書く．ここで余極限は $K$ で可逆な整数の
乗法系上でとる．この状況ではコホモロジーはフィルター余極限と可換なので
（正確な参照については補題 \ref{etale-cohomology-lemma-base-change-Rf-star-colim} を参照），
$\mathcal{F}[n]$ に対して補題を証明すれば十分である．したがって，$\mathcal{F}$ は
$\mathbf{Z}/n\mathbf{Z}$-加群の層で，ある $n$ が $K$ で可逆であるものと仮定してよい
（これは証明の最後に用いる）．証明中，$X \times_K Y$ を $\Spec(K)$ 上のファイバー積の
略記として用いる．$\dim(X) + \dim(Y)$ に関する帰納法で補題を証明する．
$\dim(X) \leq 0$ ならば，$U$ は $X$ の開かつ閉な部分スキームなので補題は自明である．
点 $z \in X \times_K Y$ を選ぶ．$\overline{z}$ における茎が同型であることを示す．

\medskip\noindent
$z \mapsto x \in X$ であり，$\text{trdeg}_K(\kappa(x)) > 0$ と仮定する．
$X' = \Spec(\mathcal{O}_{X, x}^{sh})$ とおき，$U' \subset X'$ で $U$ の逆像を表す．
基底変換
$$
\xymatrix{
Y \times_K X' \ar[d]_{q'}  &
Y \times_K U' \ar[l]^{h'} \ar[d]^{p'} \\
X' & U' \ar[l]_{j'}
}
$$
を考える．これは元の図式を $X' \to X$ により基底変換したものである．
$X' \to X$ は \'etale 射のフィルター余極限であることに注意する．補題
\ref{etale-cohomology-lemma-smooth-base-change-general} の形の滑らかな基底変換により，写像
$q^{-1}Rj_*\mathcal{F} \to Rh_*p^{-1}\mathcal{F}$ の $X'$，したがって
$Y \times_K X'$ への引戻しは写像
$(q')^{-1}Rj'_*\mathcal{F}' \to Rj'_*(p')^{-1}\mathcal{F}'$ である．ここで
$\mathcal{F}'$ は $\mathcal{F}$ の $U'$ への引戻しである．
（この段階では，本質的に自明な結果である \'etale 基底変換を用いれば十分である．）
したがって，$(q')^{-1}Rj'_*\mathcal{F}' \to Rj'_*(p')^{-1}\mathcal{F}'$ が同型であることを
示せば，元の写像が $\overline{z}$ における茎上で同型であることが従う．補題
\ref{etale-cohomology-lemma-strictly-henselian} により，分離閉体 $L/K$ で，
$X' = \lim X_i$ と書け，$X_i$ が $L$ 上有限型のアフィンスキームであり，
$\dim(X_i) < \dim(X)$ となるものが存在する．十分大きな $i$ に対して，
開部分 $U_i \subset X_i$ が存在し，その制限は $U'$ であって，これは $X'$ 上の開部分である．
図式
$$
\vcenter{
\xymatrix{
Y \times_K X_i \ar[d]_{q_i}  &
Y \times_K U_i \ar[l]^{h_i} \ar[d]^{p_i} \\
X_i & U_i \ar[l]_{j_i}
}
}
\quad\text{すなわち}\quad
\vcenter{
\xymatrix{
Y_L \times_L X_i \ar[d]_{q_i}  &
Y_L \times_L U_i \ar[l]^{h_i} \ar[d]^{p_i} \\
X_i & U_i \ar[l]_{j_i}
}
}
$$
に，体 $L$ 上で，また $\mathcal{F}$ のこれらの図式への引戻しに対して帰納法の仮定を
適用できる．補題 \ref{etale-cohomology-lemma-base-change-Rf-star-colim} により，写像
$(q')^{-1}Rj'_*\mathcal{F}' \to Rj'_*(p')^{-1}\mathcal{F}$ は同型である．

\medskip\noindent
$z \mapsto y \in Y$ であり，$\text{trdeg}_K(\kappa(y)) > 0$ と仮定する．
$Y' = \Spec(\mathcal{O}_{X, x}^{sh})$ とする．補題
\ref{etale-cohomology-lemma-strictly-henselian} により，分離閉体 $L/K$ で，
$Y' = \lim Y_i$ と書け，$Y_i$ が $L$ 上有限型のアフィンスキームであり，
$\dim(Y_i) < \dim(Y)$ となるものが存在する．特に $Y'$ は $L$ 上のスキームである．
添字 $L$ で，$K$ 上のスキームから $L$ 上のスキームへの基底変換を表す．
可換図式
$$
\vcenter{
\xymatrix{
Y' \times_K X \ar[d]_f  &
Y' \times_K U \ar[l]^{h'} \ar[d]^{f'} \\
Y \times_K X \ar[d]_q  &
Y \times_K U \ar[l]^h \ar[d]^p \\
X & U \ar[l]_j
}
}
\quad\text{および}\quad
\vcenter{
\xymatrix{
Y' \times_L X_L \ar[d]_{q'}  &
Y' \times_L U_L \ar[l]^{h'} \ar[d]^{p'} \\
X_L \ar[d] &
U_L \ar[l]^{j_L} \ar[d] \\
X & U \ar[l]_j
}
}
$$
を考える．左側と右側で，最上段と最下段がそれぞれ同じであることに注意する．
滑らかな基底変換により
$f^{-1}Rh_*p^{-1}\mathcal{F} = Rh'_*(f')^{-1}p^{-1}\mathcal{F}$
である（前段落と同様）．$\Spec(L) \to \Spec(K)$ に対する滑らかな基底変換
（補題 \ref{etale-cohomology-lemma-base-change-field-extension}）により，
$Rj_{L, *}\mathcal{F}_L$ は $Rj_*\mathcal{F}$ の $X_L$ への引戻しである．
これら二つの観察を合わせると，元の写像が $\overline{z}$ における茎上で同型であることを
証明するには，右側の図式の上の正方形に対する基底変換写像が同型であることを
証明すれば十分である\footnote{ここでは，サイト上のコホモロジー，注意
\ref{sites-cohomology-remark-compose-base-change} で説明されているように，
基底変換写像の「垂直合成」も基底変換写像であることを用いる．}．
そこで $Y' = \lim Y_i$ を用い，前段落とまったく同様に議論すると，
帰納法の仮定により $Y' \times_K X$ 上の写像が同型になることがわかる．

\medskip\noindent
したがって，$\dim(X) + \dim(Y)$ が最小となる反例があれば，写像
$q^{-1}Rj_*\mathcal{F} \to Rh_*p^{-1}\mathcal{F}$ が茎上で同型でないのは
$X \times_K Y$ の閉点においてのみである（実際，点 $z$ が $X \times_K Y$ の
閉点であることと，$z$ の $X$ における像と $Y$ における像がともに閉点であることとは
同値である）．同型であることは局所的に証明すれば十分なので，$X$ と $Y$ はアフィンであると
仮定してよい．このとき，開稠密埋込み $Y \to Y'$ で，$Y'$ が射影的であるものを選べる．
（閉埋込み $Y \to \mathbf{A}^n_K$ を選び，$Y'$ を $Y$ の $\mathbf{P}^n_K$ における
スキーム論的閉包とすればよい．）このとき $\dim(Y') = \dim(Y)$ なので，
$Y$ が $K$ 上射影的である「最小の」反例を得る．次の段落で，これは起こりえないことを示す．

\medskip\noindent
補題の主張にあるような図式で，$q^{-1}Rj_*\mathcal{F} \to Rh_*p^{-1}\mathcal{F}$ が
$X \times_K Y$ のすべての非閉点において同型であり，かつ $Y$ が射影的であるものを考える．
この写像の $(X \times_K Y)_{K^{sep}}$ への制限は，補題の図式を $K^{sep}$ へ
基底変換したものに対する対応する写像である．したがって $K$ は分離的代数閉であると
仮定してよく，以下そう仮定する．完全三角
$$
q^{-1}Rj_*\mathcal{F} \to Rh_*p^{-1}\mathcal{F} \to Q \to
(q^{-1}Rj_*\mathcal{F})[1]
$$
を $D((X \times_K Y)_\etale)$ において選ぶ．$Q$ は閉点に台をもつので，
$H^i(X \times_K Y, Q) = 0$ がすべての $i$ に対して成り立つことを示せば十分である；
補題 \ref{etale-cohomology-lemma-vanishing-closed-points} を参照されたい．したがって，
$q^{-1}Rj_*\mathcal{F} \to Rh_*p^{-1}\mathcal{F}$ がコホモロジー上に同型を誘導することを
示せば十分である．$\mathcal{F}$ は $n$ により零化され，この整数は $K$ で可逆であることを思い出そう．
補題 \ref{etale-cohomology-lemma-kunneth-one-proper} の K\"unneth 公式により
\begin{align*}
R\Gamma(X \times_K Y, q^{-1}Rj_*\mathcal{F})
&=
R\Gamma(X, Rj_*\mathcal{F})
\otimes_{\mathbf{Z}/n\mathbf{Z}}^\mathbf{L}
R\Gamma(Y, \mathbf{Z}/n\mathbf{Z}) \\
& =
R\Gamma(U, \mathcal{F})
\otimes_{\mathbf{Z}/n\mathbf{Z}}^\mathbf{L}
R\Gamma(Y, \mathbf{Z}/n\mathbf{Z})
\end{align*}
であり，また
$$
R\Gamma(X \times_K Y, Rh_*p^{-1}\mathcal{F}) =
R\Gamma(U \times_K Y, p^{-1}\mathcal{F}) =
R\Gamma(U, \mathcal{F})
\otimes_{\mathbf{Z}/n\mathbf{Z}}^\mathbf{L}
R\Gamma(Y, \mathbf{Z}/n\mathbf{Z})
$$
これで証明は完了する．
\end{proof}

\begin{lemma}
\label{etale-cohomology-lemma-punctual-base-change}
$K$ を体とする．任意の可換図式
$$
\xymatrix{
X \ar[d] & X' \ar[l] \ar[d]_{f'} & Y \ar[l]^h \ar[d]^e \\
\Spec(K) & S' \ar[l] & T \ar[l]_g
}
$$
で，$K$ 上のスキームからなり，$X' = X \times_{\Spec(K)} S'$ および
$Y = X' \times_{S'} T$ が成り立ち，$g$ が準コンパクトかつ準分離であるもの，ならびに
任意のアーベル層 $\mathcal{F}$ を $T_\etale$ 上にとり，その茎が $K$ で可逆な位数の
捩れ群であるものに対して，基底変換写像
$$
(f')^{-1}Rg_*\mathcal{F}
\longrightarrow
Rh_*e^{-1}\mathcal{F}
$$
は同型である．
\end{lemma}

\begin{proof}
問題は $X$ 上局所的なので，$X$ はアフィンであると仮定してよい．極限，補題
\ref{limits-lemma-relative-approximation} により，$X = \lim X_i$ を，遷移射がアフィンで，
$X_i$ が $K$ 上有限型のスキームである余フィルター極限として書ける．
$X'_i =  X_i \times_{\Spec(K)} S'$ および $Y_i = X'_i \times_{S'} T$ と書く．
補題 \ref{etale-cohomology-lemma-base-change-Rf-star-colim} により，頂点が $X_i, Y_i, S_i, T_i$ である
正方形に対して主張を証明すれば十分である．したがって $X$ は $K$ 上有限型であると
仮定してよい．同様に，$\mathcal{F} = \colim \mathcal{F}[n]$ と書ける．ここで余極限は
$K$ で可逆な整数の乗法系上でとる．上で用いた同じ補題により，$\mathcal{F}$ が
$\mathbf{Z}/n\mathbf{Z}$-加群の層で，ある $n$ が $K$ で可逆である場合に帰着する．

\medskip\noindent
$K$ をその代数閉包 $\overline{K}$ で置き換えてよい．実際，補題
\ref{etale-cohomology-lemma-base-change-field-extension} により，順像の形成は $\overline{K}$ への
基底変換と可換である（$g$ と $h$ のいずれにも適用できる）．また，
$X'_{\overline{K}}$ への制限後に一致を証明すれば十分である．次に，\'etale
コホモロジーの位相的不変性があるので，$X$ をその被約化で置き換えてよい；命題
\ref{etale-cohomology-proposition-topological-invariance} を参照されたい．この置換後，射
$X \to \Spec(K)$ は平坦かつ有限表示で，幾何学的に被約なファイバーをもち，
任意の基底変換，特に $X' \to S'$ についても同じことが成り立つ．したがって
$(f')^{-1}g_*\mathcal{F} \to Rh_*e^{-1}\mathcal{F}$ は補題
\ref{etale-cohomology-lemma-fppf-reduced-fibres-base-change-f-star} により同型である．

\medskip\noindent
ここで補題 \ref{etale-cohomology-lemma-base-change-does-not-hold-post} を適用すると，次を証明すれば
十分であることがわかる：可換図式
$$
\xymatrix{
X \ar[d]_f & X' \ar[d] \ar[l] & Y \ar[l]^h \ar[d] \\
\Spec(K) & S' \ar[l] & \Spec(L) \ar[l]
}
$$
で，両方の正方形が Cartesian であり，$S'$ がアフィン，整，正規で，
その関数体 $K$ が代数閉であるものが与えられたとき，
$R^qh_*(\mathbf{Z}/d\mathbf{Z})$ は $q > 0$ および $d | n$ に対して零である．
この消滅は，写像
$$
(f')^{-1}R^q(\Spec(L) \to S')_*\mathbf{Z}/d\mathbf{Z}
\longrightarrow
R^qh_*\mathbf{Z}/d\mathbf{Z}
$$
が同型であるという主張と同値であることに注意する．実際，左辺は，例えば補題
\ref{etale-cohomology-lemma-Rf-star-zero-normal-with-alg-closed-function-field} により零である．

\medskip\noindent
$S' = \Spec(B)$ と書き，$L$ を $B$ の分数体とする．
$B = \bigcup_{i \in I} B_i$ と，有限型 $K$-部分代数 $B_i$ の和集合に書く．
$J$ を組 $(i, g)$ で，$i \in I$ かつ $g \in B_i$ が非零であるもの全体とし，
$(i', g') \geq (i, g)$ であることを，$i' \geq i$ かつ $g$ の像が
$(B_{i'})_{g'}$ の可逆元であることと定める．このとき
$L = \colim_{(i, g) \in J} (B_i)_g$ である．$j = (i, g) \in J$ に対して
$S_j = \Spec(B_i)$ および $U_j = \Spec((B_i)_g)$ とおく．このとき
$$
\vcenter{
\xymatrix{
X' \ar[d] & Y \ar[l]^h \ar[d] \\
S' & \Spec(L) \ar[l]
}
}
\quad\text{は次の図式の余極限である}\quad
\vcenter{
\xymatrix{
X \times_K S_j \ar[d] & X \times_K U_j \ar[l]^{h_j} \ar[d] \\
S_j & U_j \ar[l]
}
}
$$
したがって補題 \ref{etale-cohomology-lemma-base-change-Rf-star-colim} を適用すると，右側の図式において
基底変換が成り立つことを証明すれば十分であるとわかる．これは補題
\ref{etale-cohomology-lemma-kunneth-localize-on-X} で証明した内容である．
\end{proof}

\begin{lemma}
\label{etale-cohomology-lemma-punctual-base-change-upgrade}
$K$ を体とする．$n \geq 1$ は $K$ で可逆であるとする．可換図式
$$
\xymatrix{
X \ar[d] & X' \ar[l]^p \ar[d]_{f'} & Y \ar[l]^h \ar[d]^e \\
\Spec(K) & S' \ar[l] & T \ar[l]_g
}
$$
で，$X' = X \times_{\Spec(K)} S'$ および $Y = X' \times_{S'} T$ が成り立ち，
$g$ が準コンパクトかつ準分離であるものを考える．標準的な写像
$$
p^{-1}E \otimes_{\mathbf{Z}/n\mathbf{Z}}^\mathbf{L} (f')^{-1}Rg_*F
\longrightarrow
Rh_*(h^{-1}p^{-1}E \otimes_{\mathbf{Z}/n\mathbf{Z}}^\mathbf{L} e^{-1}F)
$$
は，$E$ が $D^+(X_\etale, \mathbf{Z}/n\mathbf{Z})$ の対象で，Tor 振幅が
$[a, \infty]$ に含まれるような $a \in \mathbf{Z}$ が存在し，かつ
$F$ が $D^+(T_\etale, \mathbf{Z}/n\mathbf{Z})$ の対象ならば同型である．
\end{lemma}

\begin{proof}
この補題は補題 \ref{etale-cohomology-lemma-punctual-base-change} を導来圏の対象へ一般化したものである．
補題 \ref{etale-cohomology-lemma-punctual-base-change} ではスキーム $X$ が $K$ 上任意であるため，
本補題の主張が成り立つ．読者には，この骨の折れる証明を飛ばすことを強く勧める
（あるいは最後の段落だけを読まれたい）．

\medskip\noindent
$E$ は，下に有界な K-平坦複体 $\mathcal{E}^\bullet$ で，各項が平坦な
$\mathbf{Z}/n\mathbf{Z}$-加群であるものにより表せる．サイト上のコホモロジー，補題
\ref{sites-cohomology-lemma-bounded-below-tor-amplitude} を参照されたい．
整数 $b$ を，$H^i(F) = 0$ が $i < b$ に対して成り立つように選ぶ．
十分大きな整数 $N$ を選び，愚切断の短完全列
$$
0 \to \sigma_{\geq N + 1}\mathcal{E}^\bullet \to
\mathcal{E}^\bullet \to
\sigma_{\leq N}\mathcal{E}^\bullet \to 0
$$
を考える．これは完全三角 $E'' \to E \to E' \to E''[1]$ を
$D(X_\etale, \mathbf{Z}/n\mathbf{Z})$ において与える．$F$ を固定すると，補題の主張にある
矢印の両辺は $E$ に関する完全関手である．
$$
p^{-1}E'' \otimes_{\mathbf{Z}/n\mathbf{Z}}^\mathbf{L} (f')^{-1}Rg_*F
\quad\text{および}\quad
Rh_*(h^{-1}p^{-1}E'' \otimes_{\mathbf{Z}/n\mathbf{Z}}^\mathbf{L} e^{-1}F)
$$
はいずれも次数 $\geq N + b$ にあることに注意する．したがって，対象 $E'$ に対して
補題を証明できれば，補題は次数 $\leq N + b$ で成り立ち，結論を得る．細部の一部は
省略する．ゆえに，$E$ は平坦な $\mathbf{Z}/n\mathbf{Z}$-加群の有界複体で表されると
仮定してよい．同じ種類の議論をもう一度行うと，$E$ は一つの平坦な
$\mathbf{Z}/n\mathbf{Z}$-加群 $\mathcal{E}$ で与えられると仮定してよい．

\medskip\noindent
次に，変数 $F$ に同じ議論を用いると，$F$ が一つの
$\mathbf{Z}/n\mathbf{Z}$-加群の層 $\mathcal{F}$ で与えられる場合に帰着する．
$\mathcal{F}$ は整数 $m | n$ で零化されるとする．$\ell$ が $m$ を割る素数で
$m > \ell$ ならば，短完全列
$0 \to \mathcal{F}[\ell] \to \mathcal{F} \to
\mathcal{F}/\mathcal{F}[\ell] \to 0$ を考え，より小さい $m$ に帰着できる．
最終的に，$\mathcal{F}$ が素数 $\ell$ により零化され，この素数が $n$ を割る場合に帰着する．
この場合，
$$
p^{-1}\mathcal{E}
\otimes_{\mathbf{Z}/n\mathbf{Z}}^\mathbf{L}
(f')^{-1}Rg_*\mathcal{F}
=
p^{-1}(\mathcal{E}/\ell \mathcal{E})
\otimes_{\mathbf{F}_\ell}^\mathbf{L}
(f')^{-1}Rg_*\mathcal{F}
$$
$\mathcal{E}$ の平坦性により，この等式が成り立つことに注意する．もう一方の項についても
同様である．したがって，以下で論じる $\mathbf{F}_\ell$-ベクトル空間の層の場合に帰着する．

\medskip\noindent
$\ell$ は $K$ で可逆な素数であると仮定する．$\mathcal{E}$，$\mathcal{F}$ は，
$\mathbf{F}_\ell$-ベクトル空間の層で，それぞれ $X_\etale$ および $T_\etale$ 上にあると仮定する．
写像
$$
p^{-1}\mathcal{E} \otimes_{\mathbf{F}_\ell} (f')^{-1}R^qg_*\mathcal{F}
\longrightarrow
R^qh_*(h^{-1}p^{-1}\mathcal{E} \otimes_{\mathbf{F}_\ell} e^{-1}\mathcal{F})
$$
がすべての $q \geq 0$ に対して同型であることを示したい．この問題は $X$ 上局所的なので，
$X$ はアフィンであると仮定してよい．$\mathcal{E}$ は，$\mathbf{F}_\ell$-ベクトル空間の
構成可能層を $X_\etale$ 上でとったもののフィルター余極限として書ける；補題
\ref{etale-cohomology-lemma-torsion-colimit-constructible} を参照されたい．テンソル積はフィルター余極限と
可換であり，高次順像も同様なので（補題 \ref{etale-cohomology-lemma-relative-colimit}），
$\mathcal{E}$ は $\mathbf{F}_\ell$-ベクトル空間の構成可能層で，$X_\etale$ 上にあると
仮定してよい．このとき整数 $m$ と，有限かつ有限表示な射
$\pi_i : X_i \to X$，$i = 1, \ldots, m$ を，単射
$$
\mathcal{E} \to
\bigoplus\nolimits_{i = 1, \ldots, m}
\pi_{i, *}\mathbf{F}_\ell
$$
が存在するように選べる．補題 \ref{etale-cohomology-lemma-constructible-maps-into-constant-general} を
参照されたい．この直和も構成可能層であることに注意する（補題
\ref{etale-cohomology-lemma-finite-pushforward-constructible}）．したがって余核も構成可能である
（補題 \ref{etale-cohomology-lemma-constructible-abelian}）．次元移動，すなわち $q$ に関する帰納法を
$\mathbf{F}_\ell$-ベクトル空間の構成可能層の圏を $X_\etale$ 上で考え，そこで行うことにより，
層 $\pi_{i, *}\mathbf{F}_\ell$ に対して結果を証明すれば十分である
（詳細は省略する；ヒント：まず $q = 0$ のときの単射性を，すべての構成可能な
$\mathcal{E}$ に対して証明せよ）．この場合を証明するため，補題の図式を
$$
\xymatrix{
X_i \ar[d]^{\pi_i} &
X'_i \ar[l]^{p_i} \ar[d]^{\pi'_i} &
Y_i \ar[l]^{h_i} \ar[d]^{\rho_i} \\
X \ar[d] & X' \ar[l]^p \ar[d]_{f'} & Y \ar[l]^h \ar[d]^e \\
\Spec(K) & S' \ar[l] & T \ar[l]_g
}
$$
へ拡張し，すべての正方形を Cartesian とする．以下の等式では，
$R\pi_{i, *} = \pi_{i, *}$，ならびに $\pi'_i$，$\rho_i$ に対する同様の等式，
Leray スペクトル系列，補題 \ref{etale-cohomology-lemma-finite-pushforward-commutes-with-base-change}，および
$\pi_i$，$\pi'_i$，$\rho_i$ に対する補題
\ref{etale-cohomology-lemma-projection-formula-proper-mod-n}（有限射に対してはほとんど自明である）を用いる．
すると
\begin{align*}
p^{-1}\pi_{i, *}\mathbf{F}_\ell
\otimes_{\mathbf{F}_\ell} (f')^{-1}R^qg_*\mathcal{F}
& =
\pi'_{i, *}\mathbf{F}_\ell
\otimes_{\mathbf{F}_\ell} (f')^{-1}R^qg_*\mathcal{F} \\
& =
\pi'_{i, *}((\pi'_i)^{-1} (f')^{-1}R^qg_*\mathcal{F})
\end{align*}
同様に，
\begin{align*}
R^qh_*(h^{-1}p^{-1} \pi_{i, *}\mathbf{F}_\ell
\otimes_{\mathbf{F}_\ell} e^{-1}\mathcal{F})
& =
R^qh_*(\rho_{i, *}\mathbf{F}_\ell
\otimes_{\mathbf{F}_\ell} e^{-1}\mathcal{F}) \\
& =
R^qh_*(\rho_i^{-1}e^{-1}\mathcal{F}) \\
& =
\pi'_{i, *}R^qh_{i, *} \rho_i^{-1}e^{-1}\mathcal{F})
\end{align*}
ここで
$R^qh_{i, *} \rho_i^{-1}e^{-1}\mathcal{F} = 
(\pi'_i)^{-1} (f')^{-1}R^qg_*\mathcal{F}$ が補題
\ref{etale-cohomology-lemma-punctual-base-change} により成り立つので，結論を得る．
\end{proof}

\begin{lemma}
\label{etale-cohomology-lemma-punctual-base-change-upgrade-unbounded}
$K$ を体とする．$n \geq 1$ は $K$ で可逆であるとする．可換図式
$$
\xymatrix{
X \ar[d] & X' \ar[l]^p \ar[d]_{f'} & Y \ar[l]^h \ar[d]^e \\
\Spec(K) & S' \ar[l] & T \ar[l]_g
}
$$
で，$K$ 上有限型のスキームからなり，$X' = X \times_{\Spec(K)} S'$ および
$Y = X' \times_{S'} T$ が成り立つものを考える．標準的な写像
$$
p^{-1}E \otimes_{\mathbf{Z}/n\mathbf{Z}}^\mathbf{L} (f')^{-1}Rg_*F
\longrightarrow
Rh_*(h^{-1}p^{-1}E \otimes_{\mathbf{Z}/n\mathbf{Z}}^\mathbf{L} e^{-1}F)
$$
は，$E$ が $D(X_\etale, \mathbf{Z}/n\mathbf{Z})$ の対象で，
$F$ が $D(T_\etale, \mathbf{Z}/n\mathbf{Z})$ の対象ならば同型である．
\end{lemma}

\begin{proof}
各関手が直和と可換であることを用いて，補題
\ref{etale-cohomology-lemma-punctual-base-change-upgrade} に帰着する．読者には証明を飛ばすことを勧める．
導来テンソル積は直和と可換することを思い出そう．（導来）引戻しも直和と可換する．
$Rh_*$ と $Rg_*$ も直和と可換する；補題 \ref{etale-cohomology-lemma-finite-cd} および
\ref{etale-cohomology-lemma-finite-cd-mod-n-direct-sums} を参照されたい（ここでスキームが $K$ 上
有限型であることを用いる）．

\medskip\noindent
証明を終えるため，次のように論じる．まず
$E = \text{hocolim} \tau_{\leq N} E$ と書く．各関手は直和と可換するので，ホモトピー余極限とも
可換する．したがって，上に有界な $E$ に対して補題を証明すれば十分である．
$F$ についても同様に，$F$ は上に有界であると仮定してよい．このとき $E$ は，
上に有界な複体 $\mathcal{E}^\bullet$ で，その各項が $\mathbf{Z}/n\mathbf{Z}$-加群の層であるものにより表せる．すると
$$
\mathcal{E}^\bullet = \colim \sigma_{\geq -N}\mathcal{E}^\bullet
$$
である（愚切断）．したがって，$\mathcal{E}^\bullet$ は
$\mathbf{Z}/n\mathbf{Z}$-加群の層の有界複体であると仮定してよい．
$F$ に対しては，平坦な（！）$\mathbf{Z}/n\mathbf{Z}$-加群の層からなる上に有界な複体を
選ぶ．すると，$F$ が平坦な $\mathbf{Z}/n\mathbf{Z}$-加群の層の有界複体で表される場合に
帰着する．ここで補題 \ref{etale-cohomology-lemma-punctual-base-change-upgrade} を適用して結論を得る．
\end{proof}

\begin{lemma}
\label{etale-cohomology-lemma-kunneth}
$k$ を分離閉体とする．$X$ と $Y$ を $k$ 上有限型のスキームとする．
$n \geq 1$ を $k$ で可逆な整数とする．このとき
$E \in D(X_\etale, \mathbf{Z}/n\mathbf{Z})$ および
$K \in D(Y_\etale, \mathbf{Z}/n\mathbf{Z})$ に対して
$$
R\Gamma(X \times_{\Spec(k)} Y,
\text{pr}_1^{-1}E
\otimes_{\mathbf{Z}/n\mathbf{Z}}^\mathbf{L}
\text{pr}_2^{-1}K
)
=
R\Gamma(X, E)
\otimes_{\mathbf{Z}/n\mathbf{Z}}^\mathbf{L}
R\Gamma(Y, K)
$$
\end{lemma}

\begin{proof}
補題 \ref{etale-cohomology-lemma-punctual-base-change-upgrade-unbounded} により
$$
R\text{pr}_{1, *}(
\text{pr}_1^{-1}E
\otimes_{\mathbf{Z}/n\mathbf{Z}}^\mathbf{L}
\text{pr}_2^{-1}K) =
E \otimes_{\mathbf{Z}/n\mathbf{Z}}^\mathbf{L}
\underline{R\Gamma(Y, K)}
$$
である．補題 \ref{etale-cohomology-lemma-pull-out-constant-mod-n} から結論が従う．この補題を
用いられるのは，補題 \ref{etale-cohomology-lemma-finite-cd} により $\text{cd}(X) < \infty$ だからである．
\end{proof}














\section{カオス位相と Zariski 位相の比較}
\label{etale-cohomology-section-compare-chaotic-Zariski}

\noindent
アフィンスキームの構造層を構成するとき，まずアフィン開部分上の値を構成し，
次にそれをすべての開部分へ拡張する．同様の構成は，スキーム $X$ 上のアーベル群の
複体を構成する際にしばしば有用である．$X_{affine, Zar}$ は，$X$ のアフィン開部分の
圏に標準 Zariski 被覆による位相を入れたものを表すことを思い出そう；位相，定義
\ref{topologies-definition-big-small-Zariski} を参照されたい．$X_{affine, Zar}$ のトポスは
$X$ の小 Zariski トポスである；位相，補題
\ref{topologies-lemma-alternative-zariski} を参照されたい．本節では，
$X_{affine}$ で，同じ台となる圏にカオス位相，すなわち層と前層が一致する位相を入れたものを
表す．サイトの射
$$
\epsilon : X_{affine, Zar} \longrightarrow X_{affine}
$$
を得る．これはサイト上のコホモロジー，節
\ref{sites-cohomology-section-compare} におけるものと同様である．

\begin{lemma}
\label{etale-cohomology-lemma-check-zar}
上の状況で，$K$ を $D^+(X_{affine})$ の対象とする．このとき $K$ が（充満忠実な）関手
$R\epsilon_* ; D(X_{affine, Zar}) \to D(X_{affine})$ の本質的像に入るための必要十分条件は，
次の二条件が成り立つことである．
\begin{enumerate}
\item $R\Gamma(\emptyset, K)$ が $D(\textit{Ab})$ において零である．
\item $U = V \cup W$ で，$U, V, W \subset X$ がアフィン開部分，
$V, W \subset U$ が標準開部分であるとき（代数学，定義
\ref{algebra-definition-Zariski-topology}），サイト上のコホモロジー，補題
\ref{sites-cohomology-lemma-c-square} の写像 $c^K_{U, V, W, V \cap W}$ は擬同型である．
\end{enumerate}
\end{lemma}

\begin{proof}
（関手 $R\epsilon_*$ は，サイト上のコホモロジー，節
\ref{sites-cohomology-section-compare} の議論により充満忠実である．）
空集合に関する小さな難点を除けば，これは非常に一般的なサイト上のコホモロジー，補題
\ref{sites-cohomology-lemma-descent-squares} から従う．その仮定はスキーム，補題
\ref{schemes-lemma-sheaf-on-affines} およびサイト上のコホモロジー，補題
\ref{sites-cohomology-lemma-descent-squares-helper} により満たされる．

\medskip\noindent
この難点を避けるため，$X_{affine, almost-chaotic}$ で，空対象 $\emptyset$ が二つの被覆，
すなわち $\{\emptyset \to \emptyset\}$ と空被覆をもつサイトを表す
（議論についてはサイト，例 \ref{sites-example-site-topological} を参照されたい）．
このときサイトの射
$$
X_{affine, Zar} \to X_{affine, almost-chaotic} \to X_{affine}
$$
がある．上の議論は第一の矢印に対して適用できる．そこで，$K$ を $D^+(X_{affine})$ の
対象としたとき，その対象が（充満忠実な）関手
$D(X_{affine, almost-chaotic}) \to D(X_{affine})$ の本質的像に入ることと，
$R\Gamma(\emptyset, K)$ が $D(\textit{Ab})$ において零であることとが同値であることは，
読者に委ねる．
\end{proof}









\section{大トポスと小トポスの比較}
\label{etale-cohomology-section-compare}

\noindent
$S$ をスキームとする．位相，補題 \ref{topologies-lemma-at-the-bottom-etale} では，
比較射 $\pi_S : (\Sch/S)_\etale \to S_\etale$ および
$i_S : \Sh(S_\etale) \to \Sh((\Sch/S)_\etale)$ を導入した．これらは
$\pi_S \circ i_S = \text{id}$ および $\pi_{S, *} = i_S^{-1}$ を満たす．より一般に，
$f : T \to S$ が $(\Sch/S)_\etale$ の対象ならば，射
$i_f : \Sh(T_\etale) \to \Sh((\Sch/S)_\etale)$ で
$f_{small} = \pi_S \circ i_f$ を満たすものがある；位相，補題
\ref{topologies-lemma-put-in-T-etale} および
\ref{topologies-lemma-morphism-big-small-etale} を参照されたい．降下，注意
\ref{descent-remark-change-topologies-ringed} では，これらを環付きサイトの射
$$
\pi_S : ((\Sch/S)_\etale, \mathcal{O}) \to (S_\etale, \mathcal{O}_S)
$$
および環付きトポスの射
$$
i_S : (\Sh(S_\etale), \mathcal{O}_S) \to (\Sh((\Sch/S)_\etale), \mathcal{O})
$$
ならびに
$$
i_f : (\Sh(T_\etale), \mathcal{O}_T) \to (\Sh((\Sch/S)_\etale, \mathcal{O}))
$$
へ拡張した．制限 $i_S^{-1} = \pi_{S, *}$（位相，定義
\ref{topologies-definition-restriction-small-etale} を参照）は，$\mathcal{O}$ を
$\mathcal{O}_S$ に写すことに注意する．同様に，$i_f^{-1}$ は $\mathcal{O}$ を
$\mathcal{O}_T$ に写す．降下，注意 \ref{descent-remark-change-topologies-ringed} を参照されたい．
したがって，$i_S^*\mathcal{F} = i_S^{-1}\mathcal{F}$ および
$i_f^*\mathcal{F} = i_f^{-1}\mathcal{F}$ は，任意の $\mathcal{O}$-加群
$\mathcal{F}$ を $(\Sch/S)_\etale$ 上にとるとき成り立つ．特に $i_S^*$ と $i_f^*$ は完全関手である．
関手 $i_S^*$ はしばしば $\mathcal{F} \mapsto \mathcal{F}|_{S_\etale}$ と表される
（これは位相，定義 \ref{topologies-definition-restriction-small-etale} の記法と衝突しない）．

\begin{lemma}
\label{etale-cohomology-lemma-compare-injectives}
$S$ をスキームとする．$T$ を $(\Sch/S)_\etale$ の対象とする．
\begin{enumerate}
\item $\mathcal{I}$ が $\textit{Ab}((\Sch/S)_\etale)$ において単射的ならば，
\begin{enumerate}
\item $i_f^{-1}\mathcal{I}$ は $\textit{Ab}(T_\etale)$ において単射的である．
\item $\mathcal{I}|_{S_\etale}$ は $\textit{Ab}(S_\etale)$ において単射的である．
\end{enumerate}
\item $\mathcal{I}^\bullet$ が $\textit{Ab}((\Sch/S)_\etale)$ における K-単射的複体ならば，
\begin{enumerate}
\item $i_f^{-1}\mathcal{I}^\bullet$ は $\textit{Ab}(T_\etale)$ における K-単射的複体である．
\item $\mathcal{I}^\bullet|_{S_\etale}$ は $\textit{Ab}(S_\etale)$ における K-単射的複体である．
\end{enumerate}
\end{enumerate}
加群に対する対応する主張は成り立たない．
\end{lemma}

\begin{proof}
(1)(b) と (2)(b) は，制限関手 $\pi_{S, *} = i_S^{-1}$ が完全関手
$\pi_S^{-1}$ の右随伴であることから形式的に従う；ホモロジー，補題
\ref{homology-lemma-adjoint-preserve-injectives} および導来圏，補題
\ref{derived-lemma-adjoint-preserve-K-injectives} を参照されたい．

\medskip\noindent
(1)(a) と (2)(a) は二通りに示せる．第一の証明：$i_f^{-1}$ が完全関手
$i_{f, !}$ の右随伴であることを用いる．この関手は，集合の層に対しては位相，補題
\ref{topologies-lemma-put-in-T-etale} で，アーベル層に対してはサイト上の加群，補題
\ref{sites-modules-lemma-g-shriek-adjoint} で構成されている．サイト上の加群，補題
\ref{sites-modules-lemma-exactness-lower-shriek} では，これが完全であることが示されている．
第二の証明：位相，補題 \ref{topologies-lemma-morphism-big-small-etale} で示されている
$i_f = i_T \circ f_{big}$ を用いる．$f_{big}$ は局所化なので，これによる引戻しは
単射的対象と K-単射的対象を保つ；サイト上のコホモロジー，補題
\ref{sites-cohomology-lemma-cohomology-of-open} および
\ref{sites-cohomology-lemma-restrict-K-injective-to-open} を参照されたい．そこで，すでに証明した
(1)(b) と (2)(b) を関手 $i_T^{-1}$ に適用して結論を得る．

\medskip\noindent
$S = \Spec(\mathbf{Z})$ とし，写像 $2 : \mathcal{O}_S \to \mathcal{O}_S$ を考える．
これは $\mathcal{O}_S$-加群の単射で，$S_\etale$ 上にある．しかし引戻し
$\pi_S^*(2) : \mathcal{O} \to \mathcal{O}$ は単射でない．これは
$\Spec(\mathbf{F}_2)$ 上で評価すればわかる．ここで，単射
$\alpha : \mathcal{O} \to \mathcal{I}$ を，単射的な $\mathcal{O}$-加群
$\mathcal{I}$ で $(\Sch/S)_\etale$ 上にあるものへの単射として選ぶ．図式
$$
\xymatrix{
\mathcal{O}_S \ar[d]_2 \ar[rr]_{\alpha|_{S_\etale}} & &
\mathcal{I}|_{S_\etale} \\
\mathcal{O}_S \ar@{..>}[rru]
}
$$
を考える．点線の矢印は $\mathcal{O}_S$-加群の圏には存在しえない．実際，存在すれば
（随伴により）単射 $\alpha$ が単射でない写像 $\pi_S^*(2)$ を経由することになり，
これは不可能である．したがって $\mathcal{I}|_{S_\etale}$ は単射的な
$\mathcal{O}_S$-加群ではない．
\end{proof}

\noindent
$f : T \to S$ をスキームの射とする．位相，補題
\ref{topologies-lemma-morphism-big-small-etale} (3) の可換図式から，環付きサイトの可換図式
$$
\xymatrix{
(T_\etale, \mathcal{O}_T) \ar[d]_{f_{small}} &
((\Sch/T)_\etale, \mathcal{O}) \ar[d]^{f_{big}} \ar[l]^{\pi_T} \\
(S_\etale, \mathcal{O}_S) &
((\Sch/S)_\etale, \mathcal{O}) \ar[l]_{\pi_S}
}
$$
が得られる．これは $f_{small}^\sharp$，$f_{big}^\sharp$，$\pi_S^\sharp$，および
$\pi_T^\sharp$ の定義を書き下せば容易にわかる．特に，
\begin{equation}
\label{etale-cohomology-equation-compare-big-small}
(f_{big, *}\mathcal{F})|_{S_\etale} =
f_{small, *}(\mathcal{F}|_{T_\etale})
\end{equation}
は，任意の層 $\mathcal{F}$ を $(\Sch/T)_\etale$ 上にとるとき成り立つ．さらに
$\mathcal{F}$ が $\mathcal{O}$-加群の層ならば，(\ref{etale-cohomology-equation-compare-big-small}) は
$\mathcal{O}_S$-加群の同型で，$S_\etale$ 上にある．

\begin{lemma}
\label{etale-cohomology-lemma-compare-higher-direct-image}
$f : T \to S$ をスキームの射とする．
\begin{enumerate}
\item $K$ が $D((\Sch/T)_\etale)$ の対象ならば，
$
(Rf_{big, *}K)|_{S_\etale} = Rf_{small, *}(K|_{T_\etale})
$
が $D(S_\etale)$ において成り立つ．
\item $K$ が $D((\Sch/T)_\etale, \mathcal{O})$ の対象ならば，
$
(Rf_{big, *}K)|_{S_\etale} = Rf_{small, *}(K|_{T_\etale})
$
が $D(\textit{Mod}(S_\etale, \mathcal{O}_S))$ において成り立つ．
\end{enumerate}
より一般に，$g : S' \to S$ を $(\Sch/S)_\etale$ の対象とする．ファイバー積
$$
\xymatrix{
T' \ar[r]_{g'} \ar[d]_{f'} & T \ar[d]^f \\
S' \ar[r]^g & S
}
$$
を考える．このとき
\begin{enumerate}
\item[(3)] $K$ が $D((\Sch/T)_\etale)$ の対象ならば，
$i_g^{-1}(Rf_{big, *}K) = Rf'_{small, *}(i_{g'}^{-1}K)$
が $D(S'_\etale)$ において成り立つ．
\item[(4)] $K$ が $D((\Sch/T)_\etale, \mathcal{O})$ の対象ならば，
$i_g^*(Rf_{big, *}K) = Rf'_{small, *}(i_{g'}^*K)$
が $D(\textit{Mod}(S'_\etale, \mathcal{O}_{S'}))$ において成り立つ．
\item[(5)] $K$ が $D((\Sch/T)_\etale)$ の対象ならば，
$g_{big}^{-1}(Rf_{big, *}K) = Rf'_{big, *}((g'_{big})^{-1}K)$
が $D((\Sch/S')_\etale)$ において成り立つ．
\item[(6)] $K$ が $D((\Sch/T)_\etale, \mathcal{O})$ の対象ならば，
$g_{big}^*(Rf_{big, *}K) = Rf'_{big, *}((g'_{big})^*K)$
が $D(\textit{Mod}(S'_\etale, \mathcal{O}_{S'}))$ において成り立つ．
\end{enumerate}
\end{lemma}

\begin{proof}
(1) は，$K$ を表すアーベル層の K-単射的複体を選び，補題
\ref{etale-cohomology-lemma-compare-injectives} と (\ref{etale-cohomology-equation-compare-big-small}) を用いれば従う．

\medskip\noindent
(3) は，$K$ を表すアーベル層の K-単射的複体を選び，補題
\ref{etale-cohomology-lemma-compare-injectives} と位相，補題
\ref{topologies-lemma-morphism-big-small-cartesian-diagram-etale} を用いれば従う．

\medskip\noindent
(5) はサイト上のコホモロジー，補題
\ref{sites-cohomology-lemma-localize-cartesian-square} である．

\medskip\noindent
(6) はサイト上のコホモロジー，補題
\ref{sites-cohomology-lemma-localize-cartesian-square-modules} である．

\medskip\noindent
(2) は次のように証明できる．上で，環付きサイトの射として
$\pi_S \circ f_{big} = f_{small} \circ \pi_T$ であることを見た．したがって
$R\pi_{S, *} \circ Rf_{big, *} = Rf_{small, *} \circ R\pi_{T, *}$
をサイト上のコホモロジー，補題
\ref{sites-cohomology-lemma-derived-pushforward-composition} により得る．
制限関手 $\pi_{S, *}$ と $\pi_{T, *}$ は完全なので，結論が従う．

\medskip\noindent
(4) は (6) と，$f' : T' \to S'$ に適用した (2) から従う．
\end{proof}

\noindent
$S$ をスキームとし，$\mathcal{H}$ を $(\Sch/S)_\etale$ 上のアーベル層とする．
$H^n_\etale(U, \mathcal{H})$ は，$\mathcal{H}$ の対象 $U$ 上のコホモロジーを表し，
この対象は $(\Sch/S)_\etale$ に属することを思い出そう．

\begin{lemma}
\label{etale-cohomology-lemma-compare-cohomology}
$f : T \to S$ をスキームの射とする．このとき
\begin{enumerate}
\item $K$ が $D(S_\etale)$ の対象ならば，
$H^n_\etale(S, \pi_S^{-1}K) = H^n(S_\etale, K)$.
\item $K$ が $D(S_\etale, \mathcal{O}_S)$ の対象ならば，
$H^n_\etale(S, L\pi_S^*K) = H^n(S_\etale, K)$.
\item $K$ が $D(S_\etale)$ の対象ならば，
$H^n_\etale(T, \pi_S^{-1}K) = H^n(T_\etale, f_{small}^{-1}K)$.
\item $K$ が $D(S_\etale, \mathcal{O}_S)$ の対象ならば，
$H^n_\etale(T, L\pi_S^*K) = H^n(T_\etale, Lf_{small}^*K)$.
\item $M$ が $D((\Sch/S)_\etale)$ の対象ならば，
$H^n_\etale(T, M) = H^n(T_\etale, i_f^{-1}M)$.
\item $M$ が $D((\Sch/S)_\etale, \mathcal{O})$ の対象ならば，
$H^n_\etale(T, M) = H^n(T_\etale, i_f^*M)$.
\end{enumerate}
\end{lemma}

\begin{proof}
(5) を証明するには，$M$ をアーベル層の K-単射的複体で表し，補題
\ref{etale-cohomology-lemma-compare-injectives} を適用して定義を書き下せばよい．(3) は，
$i_f^{-1}\pi_S^{-1} = f_{small}^{-1}$ なのでこれから従う．(1) は (3) の特別な場合である．

\medskip\noindent
(6) は非常に一般的なサイト上のコホモロジー，補題
\ref{sites-cohomology-lemma-pullback-same-cohomology} から従う．次いで (4) は
$Lf_{small}^* = i_f^* \circ L\pi_S^*$ なので従う．(2) は (4) の特別な場合である．
\end{proof}

\begin{lemma}
\label{etale-cohomology-lemma-cohomological-descent-etale}
$S$ をスキームとする．$K \in D(S_\etale)$ に対して，写像
$$
K \longrightarrow R\pi_{S, *}\pi_S^{-1}K
$$
は同型である．
\end{lemma}

\begin{proof}
実際，$\pi_S^{-1}$ と $\pi_{S, *} = i_S^{-1}$ はいずれも完全関手であり，合成
$\pi_{S, *} \circ \pi_S^{-1}$ は恒等関手である．
\end{proof}

\begin{lemma}
\label{etale-cohomology-lemma-compare-higher-direct-image-proper}
$f : T \to S$ をスキームの固有射とする．このとき
\begin{enumerate}
\item $\pi_S^{-1} \circ f_{small, *} = f_{big, *} \circ \pi_T^{-1}$
が関手 $\Sh(T_\etale) \to \Sh((\Sch/S)_\etale)$ の等式として成り立つ．
\item $\pi_S^{-1}Rf_{small, *}K = Rf_{big, *}\pi_T^{-1}K$
が，$K$ を $D^+(T_\etale)$ の対象で，コホモロジー層が捩れ層であるものとするとき成り立つ．
\item $\pi_S^{-1}Rf_{small, *}K = Rf_{big, *}\pi_T^{-1}K$
が，$K$ を $D(T_\etale, \mathbf{Z}/n\mathbf{Z})$ の対象とするとき成り立つ．
\item $\pi_S^{-1}Rf_{small, *}K = Rf_{big, *}\pi_T^{-1}K$
が，すべての $K$ を $D(T_\etale)$ にとるとき，$f$ が有限ならば成り立つ．
\end{enumerate}
\end{lemma}

\begin{proof}
(1) の証明．$\mathcal{F}$ を $T_\etale$ 上の層とする．
$g : S' \to S$ を $(\Sch/S)_\etale$ の対象とする．ファイバー積
$$
\xymatrix{
T' \ar[r]_{f'} \ar[d]_{g'} & S' \ar[d]^g \\
T \ar[r]^f & S
}
$$
を考える．このとき
$$
(f_{big, *}\pi_T^{-1}\mathcal{F})(S') =
(\pi_T^{-1}\mathcal{F})(T') =
((g'_{small})^{-1}\mathcal{F})(T')  =
(f'_{small, *}(g'_{small})^{-1}\mathcal{F})(S')
$$
であり，第二の等号は補題 \ref{etale-cohomology-lemma-describe-pullback} による．一方，
$$
(\pi_S^{-1}f_{small, *}\mathcal{F})(S') =
(g_{small}^{-1}f_{small, *}\mathcal{F})(S')
$$
であり，これも補題 \ref{etale-cohomology-lemma-describe-pullback} による．したがって集合の層に対する
固有基底変換（補題 \ref{etale-cohomology-lemma-proper-base-change-f-star}）により，二つの集合は標準的に
同型である．この同型は制限写像と両立し，同型
$\pi_S^{-1}f_{small, *}\mathcal{F} = f_{big, *}\pi_T^{-1}\mathcal{F}$.
を定める．したがって関手の同型
$\pi_S^{-1} \circ f_{small, *} = f_{big, *} \circ \pi_T^{-1}$.

\medskip\noindent
(2) の証明．標準的な基底変換写像
$\pi_S^{-1}Rf_{small, *}K \to Rf_{big, *}\pi_T^{-1}K$
が，任意の $K$ を $D(T_\etale)$ にとるとき存在する；サイト上のコホモロジー，注意
\ref{sites-cohomology-remark-base-change} を参照されたい．これが同型であることを証明するには，
基底変換写像の $i_g : \Sh(S'_\etale) \to \Sh((\Sch/S)_\etale)$ による引戻しが，
任意の対象 $g : S' \to S$ を $(\Sch/S)_\etale$ にとるとき同型であることを示せば十分である．
$T', g', f'$ は前段落と同じものとする．基底変換写像の引戻しは
\begin{align*}
g_{small}^{-1}Rf_{small, *}K
& =
i_g^{-1}\pi_S^{-1}Rf_{small, *}K \\
& \to
i_g^{-1}Rf_{big, *}\pi_T^{-1}K \\
& =
Rf'_{small, *}(i_{g'}^{-1}\pi_T^{-1}K) \\
& =
Rf'_{small, *}((g'_{small})^{-1}K)
\end{align*}
である．ここでは $\pi_S \circ i_g = g_{small}$，
$\pi_T \circ i_{g'} = g'_{small}$，および補題
\ref{etale-cohomology-lemma-compare-higher-direct-image} を用いた．$K$ が下に有界で，$K$ のコホモロジー層が
捩れ層ならば，この写像は固有基底変換定理（補題 \ref{etale-cohomology-lemma-proper-base-change}）により同型である．

\medskip\noindent
(3) の証明は，補題 \ref{etale-cohomology-lemma-proper-base-change-mod-n} を補題
\ref{etale-cohomology-lemma-proper-base-change} の代わりに用いる点を除けば，(2) の証明と同じである．

\medskip\noindent
(4) の証明．$f$ が有限ならば，関手 $f_{small, *}$ と $f_{big, *}$ は完全である．
$f_{small}$ については，これは命題 \ref{etale-cohomology-proposition-finite-higher-direct-image-zero} から従う．
任意の基底変換 $f'$ は $f$ の基底変換であり，これも有限なので，補題
\ref{etale-cohomology-lemma-compare-higher-direct-image} (3) により $f_{big, *}$ も完全である
（高次導来関手が零だからである）．したがってこの場合は (1) から従う．
\end{proof}




\section{fppf 位相と \'etale 位相の比較}
\label{etale-cohomology-section-fppf-etale}

\noindent
本節のひな型は，サイト上のコホモロジー，節
\ref{sites-cohomology-section-cohomology-LC} で論じられた，局所コンパクト位相空間上の
通常の位相と qc 位相の比較である．まず，位相，節
\ref{topologies-section-change-topologies} および
\ref{topologies-section-etale} の内容をいくつか復習する．

\medskip\noindent
$S$ をスキームとし，$(\Sch/S)_{fppf}$ を fppf サイトとする．同じ台となる圏には
第二の位相，すなわち \'etale 位相があり，したがって第二のサイト
$(\Sch/S)_\etale$ がある．恒等関手
$(\Sch/S)_\etale \to (\Sch/S)_{fppf}$ は連続であり，サイトの射
$$
\epsilon_S : (\Sch/S)_{fppf} \longrightarrow (\Sch/S)_\etale
$$
を定める．サイト上のコホモロジー，節 \ref{sites-cohomology-section-compare} を
参照されたい．$\epsilon_{S, *}$ は台となる前層上の恒等関手であり，
$\epsilon_S^{-1}$ は \'etale 層にその fppf 層化を対応させることに注意されたい．
$S_\etale$ を小 \'etale サイトとする．サイトの射
$$
\pi_S : (\Sch/S)_\etale \longrightarrow S_\etale
$$
があり，これは連続関手 $S_\etale \to (\Sch/S)_\etale$，$U \mapsto U$ により与えられる．
実際，$S_\etale$ はファイバー積と終対象をもち，上の関手はこれらと可換するので，
サイト，命題 \ref{sites-proposition-get-morphism} が適用できる．

\begin{lemma}
\label{etale-cohomology-lemma-describe-pullback-pi-fppf}
上の記法を用いる．$\mathcal{F}$ を $S_\etale$ 上の層とする．規則
$$
(\Sch/S)_{fppf} \longrightarrow \textit{Sets},\quad
(f : X \to S) \longmapsto \Gamma(X, f_{small}^{-1}\mathcal{F})
$$
は層であり，したがってもちろん $(\Sch/S)_\etale$ 上の層でもある．実際，この層は
$\pi_S^{-1}\mathcal{F}$ に等しく，これは $(\Sch/S)_\etale$ 上の等式である；また
$\epsilon_S^{-1}\pi_S^{-1}\mathcal{F}$ に等しく，これは $(\Sch/S)_{fppf}$ 上の等式である．
\end{lemma}

\begin{proof}
\'etale 位相に関する主張は補題 \ref{etale-cohomology-lemma-describe-pullback} の内容である．証明を
終えるには，$\pi_S^{-1}\mathcal{F}$ が fppf 位相に対する層であることを示せば十分である．
これも補題 \ref{etale-cohomology-lemma-describe-pullback} で示されている．
\end{proof}

\noindent
補題 \ref{etale-cohomology-lemma-describe-pullback-pi-fppf} の状況では，$\epsilon_S$ と $\pi_S$ の合成と
上の等式により，サイトの射
$$
a_S : (\Sch/S)_{fppf} \longrightarrow S_\etale
$$

\begin{lemma}
\label{etale-cohomology-lemma-push-pull-fppf-etale}
上の記法を用いる．$f : X \to Y$ を $(\Sch/S)_{fppf}$ の射とする．
このときトポスの可換図式
$$
\xymatrix{
\Sh((\Sch/X)_{fppf}) \ar[rr]_{f_{big, fppf}} \ar[d]_{\epsilon_X} & &
\Sh((\Sch/Y)_{fppf}) \ar[d]^{\epsilon_Y} \\
\Sh((\Sch/X)_\etale) \ar[rr]^{f_{big, \etale}} & &
\Sh((\Sch/Y)_\etale)
}
$$
および
$$
\xymatrix{
\Sh((\Sch/X)_{fppf}) \ar[rr]_{f_{big, fppf}} \ar[d]_{a_X} & &
\Sh((\Sch/Y)_{fppf}) \ar[d]^{a_Y} \\
\Sh(X_\etale) \ar[rr]^{f_{small}} & &
\Sh(Y_\etale)
}
$$
があり，$a_X = \pi_X \circ \epsilon_X$ および $a_Y = \pi_X \circ \epsilon_X$ を満たす．
\end{lemma}

\begin{proof}
図式の可換性は位相，節 \ref{topologies-section-change-topologies} の議論から従う．
\end{proof}

\begin{lemma}
\label{etale-cohomology-lemma-proper-push-pull-fppf-etale}
補題 \ref{etale-cohomology-lemma-push-pull-fppf-etale} において $f$ が固有ならば，
$a_Y^{-1} \circ f_{small, *} = f_{big, fppf, *} \circ a_X^{-1}$.
\end{lemma}

\begin{proof}
補題 \ref{etale-cohomology-lemma-compare-higher-direct-image-proper} (1) の証明を繰り返しても証明できるが，
ここではこの結果から導く．$\epsilon_{Y, *}$ は台となる前層上の恒等関手なので，
同型を反映する．補題 \ref{etale-cohomology-lemma-describe-pullback-pi-fppf} の記述により
$\epsilon_{Y, *} \circ a_Y^{-1} = \pi_Y^{-1}$ であり，$X$ についても同様である．
標準的な写像
$a_Y^{-1}f_{small, *}\mathcal{F} \to f_{big, fppf, *}a_X^{-1}\mathcal{F}$ が同型であることを
示すには，
\begin{align*}
\pi_Y^{-1}f_{small, *}\mathcal{F}
& =
\epsilon_{Y, *}a_Y^{-1}f_{small, *}\mathcal{F} \\
& \to 
\epsilon_{Y, *}f_{big, fppf, *}a_X^{-1}\mathcal{F} \\
& =
f_{big, \etale, *} \epsilon_{X, *}a_X^{-1}\mathcal{F} \\
& =
f_{big, \etale, *}\pi_X^{-1}\mathcal{F}
\end{align*}
が同型であることを示せば十分である．これは補題
\ref{etale-cohomology-lemma-compare-higher-direct-image-proper} (1) である．
\end{proof}

\begin{lemma}
\label{etale-cohomology-lemma-descent-sheaf-fppf-etale}
補題 \ref{etale-cohomology-lemma-push-pull-fppf-etale} において，$f$ が平坦，有限表示局所的，
かつ全射であると仮定する．このとき関手
$$
\Sh(Y_\etale) \longrightarrow
\left\{
(\mathcal{G}, \mathcal{H}, \alpha)
\middle|
\begin{matrix}
\mathcal{G} \in \Sh(X_\etale),\ \mathcal{H} \in \Sh((\Sch/Y)_{fppf}), \\
\alpha : a_X^{-1}\mathcal{G} \to f_{big, fppf}^{-1}\mathcal{H}
\text{ は同型}
\end{matrix}
\right\}
$$
$\mathcal{F}$ を
$(f_{small}^{-1}\mathcal{F}, a_Y^{-1}\mathcal{F}, can)$ に送る関手は圏同値である．
\end{lemma}

\begin{proof}
関手 $a_X^{-1}$ は充満忠実である（補題
\ref{etale-cohomology-lemma-describe-pullback-pi-fppf} により
$a_{X, *}a_X^{-1} = \text{id}$ だからである）．したがって忘却関手
$(\mathcal{G}, \mathcal{H}, \alpha) \mapsto \mathcal{H}$ は，三つ組の圏を
$\Sh((\Sch/Y)_{fppf})$ の充満部分圏と同一視する．さらに関手
$a_Y^{-1}$ も充満忠実なので，補題の関手もまた充満忠実である．

\medskip\noindent
\'etale 被覆 $\{Y_i \to Y\}$ が与えられているとする．
$f_i : X_i \to Y_i$ を $f$ の基底変換とし，
$f_{ij} = f_i \times f_j : X_i \times_X X_j  \to Y_i \times_Y Y_j$
と書く．補題が $f_i$ と $f_{ij}$ に対してすべての $i, j$ について成り立つならば，
$f$ に対しても成り立つ，と主張する．実際，与えられた \'etale 被覆は，4つのサイト
$Y_\etale, X_\etale, (\Sch/Y)_{fppf}, (\Sch/X)_{fppf}$ のそれぞれにおいて
終対象の \'etale 被覆を定める．したがって4つの場合のそれぞれで，層の圏はこの被覆に
関する貼合せデータの圏と圏同値である（Sites，補題
\ref{sites-lemma-mapping-property-glue}）．そこで圏からなる大きな可換図式を用いれば
主張の証明は完了する．詳細は省略する．この主張により，$Y$ 上 \'etale 局所的に
議論してよい．

\medskip\noindent
$\{X \to Y\}$ は fppf 被覆であることに注意する．$Y$ 上 \'etale 局所的に
議論することにより，射 $s : X' \to X$ が存在し，その合成
$f' = f \circ s : X' \to Y$ が全射な有限局所自由射であると仮定してよい．
More on Morphisms，補題
\ref{more-morphisms-lemma-dominate-fppf-etale-locally} を参照せよ．
補題が $f'$ に対して成り立つならば $f$ に対しても成り立つ，と主張する．実際，
三つ組 $(\mathcal{G}, \mathcal{H}, \alpha)$ が $f$ に対して与えられたとき，$s$ で
引き戻せば，三つ組
$(s_{small}^{-1}\mathcal{G}, \mathcal{H}, s_{big, fppf}^{-1}\alpha)$
を $f'$ に対して得る．この三つ組の解から，層 $\mathcal{F}$ を $Y_\etale$ 上に得て，
$a_Y^{-1}\mathcal{F} = \mathcal{H}$ となる．証明の第1段落により，これは元の三つ組が
本質的像に属することを意味する．したがって次の段落で述べる場合に帰着する．

\medskip\noindent
$f$ が全射な有限局所自由射であると仮定する．三つ組
$(\mathcal{G}, \mathcal{H}, \alpha)$ をとる．この場合，三つ組
$$
(\mathcal{G}_1, \mathcal{H}_1, \alpha_1) =
(f_{small}^{-1}f_{small, *}\mathcal{G},
f_{big, fppf, *}f_{big, fppf}^{-1}\mathcal{H}, \alpha_1)
$$
を考える．ここで $\alpha_1$ は次の同一視から得られる：
\begin{align*}
a_X^{-1}f_{small}^{-1}f_{small, *}\mathcal{G}
& =
f_{big, fppf}^{-1}a_Y^{-1}f_{small, *}\mathcal{G} \\
& =
f_{big, fppf}^{-1}f_{big, fppf, *}a_X^{-1}\mathcal{G} \\
& \to
f_{big, fppf}^{-1}f_{big, fppf, *}f_{big, fppf}^{-1}\mathcal{H}
\end{align*}
ここで第3の等号は補題 \ref{etale-cohomology-lemma-proper-push-pull-fppf-etale} であり，矢印は
$\alpha$ によって与えられる．この三つ組は考えている関手の像に属する．実際，
$\mathcal{F}_1 = f_{small, *}\mathcal{F}$ が解である（これを見るには再び補題
\ref{etale-cohomology-lemma-proper-push-pull-fppf-etale} を用いる；詳細は省略する）．三つ組の標準的な写像
$$
(\mathcal{G}, \mathcal{H}, \alpha)
\to
(\mathcal{G}_1, \mathcal{H}_1, \alpha_1)
$$
があり，第2成分では単位射
$\text{id} \to f_{big, fppf, *}f_{big, fppf}^{-1}$ を用いる（証明の第1段落により，
第2成分上の射を指定すれば十分である）．$\{f : X \to Y\}$ は fppf 被覆なので，写像
$\mathcal{H} \to \mathcal{H}_1$ は単射である（詳細は省略する）．
$$
\mathcal{G}_2 = \mathcal{G}_1 \amalg_\mathcal{G} \mathcal{G}_1\quad
\mathcal{H}_2 = \mathcal{H}_1 \amalg_\mathcal{H} \mathcal{H}_1
$$
とおき，$\alpha_2$ を誘導される同型とする（引戻し関手は完全なので，これは意味をもつ）．
このとき $\mathcal{H}$ は2つの写像 $\mathcal{H}_1 \to \mathcal{H}_2$ の等化子である．
三つ組 $(\mathcal{G}_2, \mathcal{H}_2, \alpha_2)$ に対して上の議論を繰り返すと，
三つ組の単射な射
$$
(\mathcal{G}_2, \mathcal{H}_2, \alpha_2)
\to
(\mathcal{G}_3, \mathcal{H}_3, \alpha_3)
$$
で，後者の三つ組が考えている関手の像に属するものを得る．それが
$\mathcal{F}_3$ に対応し，これは $\Sh(Y_\etale)$ の対象であるとする．充満忠実性により2つの写像
$\mathcal{F}_1 \to \mathcal{F}_3$ が得られるので，$\mathcal{F}$ をこれら2写像の
等化子とすることができる．現れる引戻し関手の完全性により，望むとおり
$a_Y^{-1}\mathcal{F} = \mathcal{H}$ を得る．
\end{proof}

\begin{lemma}
\label{etale-cohomology-lemma-compare-fppf-etale}
比較射
$\epsilon : (\Sch/S)_{fppf} \to (\Sch/S)_\etale$ を考える．
$\mathcal{P}$ をスキームの有限射全体の類とする．$X$ を $(\Sch/S)_\etale$ の対象とする．
$\mathcal{A}'_X \subset \textit{Ab}((\Sch/X)_\etale)$ を，形
$\pi_X^{-1}\mathcal{F}$ の層からなる充満部分圏とする．ここで $\mathcal{F}$ は
$\textit{Ab}(X_\etale)$ の対象である．このとき Cohomology on Sites の性質
(\ref{sites-cohomology-item-base-change-P}),
(\ref{sites-cohomology-item-restriction-A}),
(\ref{sites-cohomology-item-A-sheaf}),
(\ref{sites-cohomology-item-A-and-P})，および
(\ref{sites-cohomology-item-refine-tau-by-P})
が，Cohomology on Sites，状況 \ref{sites-cohomology-situation-compare} において成り立つ．
\end{lemma}

\begin{proof}
まず，Homology，補題 \ref{homology-lemma-characterize-weak-serre-subcategory} の
条件 (1)，(2)，(3)，(4) を確認して，
$\mathcal{A}'_X \subset \textit{Ab}((\Sch/X)_\etale)$ が弱 Serre 部分圏であることを示す．
たとえば補題 \ref{etale-cohomology-lemma-cohomological-descent-etale} により $\pi_X^{-1}$ は完全かつ
充満忠実なので，(1)，(2)，(3) は直ちに従う．
$0 \to \pi_X^{-1}\mathcal{F} \to \mathcal{G} \to \pi_X^{-1}\mathcal{F}' \to 0$
が $\textit{Ab}((\Sch/X)_\etale)$ の短完全列ならば，補題
\ref{etale-cohomology-lemma-cohomological-descent-etale} により
$0 \to \mathcal{F} \to \pi_{X, *}\mathcal{G} \to \mathcal{F}' \to 0$
は完全である．したがって
$\mathcal{G} = \pi_X^{-1}\pi_{X, *}\mathcal{G}$ は $\mathcal{A}'_X$ に属し，
最後の条件も確認できる．

\medskip\noindent
Cohomology on Sites の性質 (\ref{sites-cohomology-item-base-change-P}) は，
スキームのファイバー積が存在すること，およびスキームの有限射の基底変換が有限射であることから
従う．Morphisms，補題 \ref{morphisms-lemma-base-change-finite} を参照せよ．

\medskip\noindent
Cohomology on Sites の性質 (\ref{sites-cohomology-item-restriction-A}) は，
Topologies，補題 \ref{topologies-lemma-morphism-big-small-etale} の可換図式 (3) から従う．

\medskip\noindent
Cohomology on Sites の性質 (\ref{sites-cohomology-item-A-sheaf}) は補題
\ref{etale-cohomology-lemma-describe-pullback-pi-fppf} である．

\medskip\noindent
Cohomology on Sites の性質 (\ref{sites-cohomology-item-A-and-P}) は補題
\ref{etale-cohomology-lemma-compare-higher-direct-image-proper} (4) により成り立つ．

\medskip\noindent
Cohomology on Sites の性質 (\ref{sites-cohomology-item-refine-tau-by-P}) は
More on Morphisms，補題
\ref{more-morphisms-lemma-dominate-fppf-etale-locally} から従う．
\end{proof}

\begin{lemma}
\label{etale-cohomology-lemma-V-C-all-n-etale-fppf}
上の記号を用いる．
\begin{enumerate}
\item $X \in \Ob((\Sch/S)_{fppf})$ とアーベル層 $\mathcal{F}$ をとり，後者は
$X_\etale$ 上の層とする．このとき
$\epsilon_{X, *}a_X^{-1}\mathcal{F} = \pi_X^{-1}\mathcal{F}$
かつ $R^i\epsilon_{X, *}(a_X^{-1}\mathcal{F}) = 0$ が $i > 0$ に対して成り立つ．
\item 有限射 $f : X \to Y$ を $(\Sch/S)_{fppf}$ でとり，アーベル層
$\mathcal{F}$ を $X$ 上にとる．このとき
$a_Y^{-1}(R^if_{small, *}\mathcal{F}) =
R^if_{big, fppf, *}(a_X^{-1}\mathcal{F})$
がすべての $i$ について成り立つ．
\item スキーム $X$ と $K$ をとり，後者は $D^+(X_\etale)$ に属するとする．写像
$\pi_X^{-1}K \to R\epsilon_{X, *}(a_X^{-1}K)$ は同型である．
\item スキームの有限射 $f : X \to Y$ と $K$ をとり，後者は
$D^+(X_\etale)$ に属するとする．このとき
$a_Y^{-1}(Rf_{small, *}K) = Rf_{big, fppf, *}(a_X^{-1}K)$.
\item スキームの固有射 $f : X \to Y$ と $K$ をとり，後者は捩れコホモロジー層をもつ
$D^+(X_\etale)$ の対象とする．このとき
$a_Y^{-1}(Rf_{small, *}K) = Rf_{big, fppf, *}(a_X^{-1}K)$.
\end{enumerate}
\end{lemma}

\begin{proof}
補題 \ref{etale-cohomology-lemma-compare-fppf-etale} により，Cohomology on Sites，節
\ref{sites-cohomology-section-compare-general} の補題はすべて現在の状況に適用できる．
結果をこの状況に読み替えるため，Cohomology on Sites，補題
\ref{sites-cohomology-lemma-A} の圏 $\mathcal{A}_X$ が
$a_X^{-1} : \textit{Ab}(X_\etale) \to \textit{Ab}((\Sch/X)_{fppf})$
の本質的像であることに注意する．

\medskip\noindent
(1) は $(V_n)$ がすべての $n$ について成り立つことと同値であり，これは Cohomology on Sites，補題
\ref{sites-cohomology-lemma-V-C-all-n-general} により成り立つ．

\medskip\noindent
(2) は Cohomology on Sites，補題
\ref{sites-cohomology-lemma-V-implies-C-general} の結論に $\epsilon_Y^{-1}$ を
適用することから従う．

\medskip\noindent
(3) は Cohomology on Sites，補題
\ref{sites-cohomology-lemma-V-C-all-n-general} (1) から従う．実際，
$\pi_X^{-1}K$ は $D^+_{\mathcal{A}'_X}((\Sch/X)_\etale)$ に属し，
$a_X^{-1} = \epsilon_X^{-1} \circ a_X^{-1}$ である．

\medskip\noindent
(4) も同じ理由により Cohomology on Sites，補題
\ref{sites-cohomology-lemma-V-C-all-n-general} (2) から従う．

\medskip\noindent
(5)．次の等式を用いる：
\begin{align*}
R\epsilon_{Y, *}Rf_{big, fppf, *}a_X^{-1}K
& =
Rf_{big, \etale, *}R\epsilon_{X, *}a_X^{-1}K \\
& =
Rf_{big, \etale, *}\pi_X^{-1}K \\
& =
\pi_Y^{-1}Rf_{small, *}K \\
& =
R\epsilon_{Y, *} a_Y^{-1}Rf_{small, *}K
\end{align*}
第1の等号は補題 \ref{etale-cohomology-lemma-push-pull-fppf-etale} の可換図式と Cohomology on Sites，
補題 \ref{sites-cohomology-lemma-derived-pushforward-composition} による．
第2の等号は (3)，第3の等号は補題
\ref{etale-cohomology-lemma-compare-higher-direct-image-proper} (2)，第4の等号は再び (3) である．
したがって基底変換写像
$a_Y^{-1}(Rf_{small, *}K) \to Rf_{big, fppf, *}(a_X^{-1}K)$
は同型
$$
R\epsilon_{Y, *}a_Y^{-1}Rf_{small, *}K \to
R\epsilon_{Y, *}Rf_{big, fppf, *}a_X^{-1}K
$$
を誘導する．次の注意により証明が完了する：写像
$\alpha : a_Y^{-1}L \to M$ をとり，$L$ は $D^+(Y_\etale)$ に，$M$ は
$D^+((\Sch/Y)_{fppf})$ に属するとする．$R\epsilon_{Y, *}\alpha$ が同型ならば，
その写像も同型である．実際，$i$ について帰納的に $H^i(\alpha)$ が同型であることを
示す．これは十分小さいすべての $i$ に対して成り立つ．$i \leq i_0$ に対して
成り立つとすると，(1) により $R^j\epsilon_{Y, *}H^i(M) = 0$ が
$j > 0$ かつ $i \leq i_0$ に対して成り立つ．実際，この範囲では
$H^i(M) = a_Y^{-1}H^i(L)$ である．したがってスペクトル系列の議論により
$\epsilon_{Y, *}H^{i_0 + 1}(M) = H^{i_0 + 1}(R\epsilon_{Y, *}M)$ である．ゆえに
$\epsilon_{Y, *}H^{i_0 + 1}(M) = \pi_Y^{-1}H^{i_0 + 1}(L) =
\epsilon_{Y, *}a_Y^{-1}H^{i_0 + 1}(L)$ である．したがって望むとおり
$H^{i_0 + 1}(\alpha)$ は同型である（$\epsilon_{Y, *}$ は台となる前層上の恒等関手なので，
同型を反映する）．
\end{proof}

\begin{lemma}
\label{etale-cohomology-lemma-cohomological-descent-etale-fppf}
スキーム $X$ をとる．$K \in D^+(X_\etale)$ に対し，写像
$$
K \longrightarrow Ra_{X, *}a_X^{-1}K
$$
は同型である．ここで $a_X : \Sh((\Sch/X)_{fppf}) \to \Sh(X_\etale)$ は上の射である．
\end{lemma}

\begin{proof}
まず主張を，$K$ が1つのアーベル層で与えられる場合に帰着させる．実際，$K$ を
下に有界な複体 $\mathcal{F}^\bullet$ で表す．層の場合から
$\mathcal{F}^n = a_{X, *} a_X^{-1} \mathcal{F}^n$
であり，層 $R^qa_{X, *}a_X^{-1}\mathcal{F}^n$ は $q > 0$ に対して零である．
$a_X^{-1}\mathcal{F}^\bullet$ と関手 $a_{X, *}$ に Leray の非輪状性補題
（Derived Categories，補題 \ref{derived-lemma-leray-acyclicity}）を適用すれば結論を得る．
以下では $K = \mathcal{F}$ と仮定する．

\medskip\noindent
補題 \ref{etale-cohomology-lemma-describe-pullback-pi-fppf} により
$a_{X, *}a_X^{-1}\mathcal{F} = \mathcal{F}$ である．したがって
$R^qa_{X, *}a_X^{-1}\mathcal{F} = 0$ が $q > 0$ に対して成り立つことを示せば十分である．これには
$a_X = \epsilon_X \circ \pi_X$ と Leray スペクトル系列（Cohomology on Sites，補題
\ref{sites-cohomology-lemma-relative-Leray}）を用いることができる．補題
\ref{etale-cohomology-lemma-V-C-all-n-etale-fppf} により，
$R^i\epsilon_{X, *}(a_X^{-1}\mathcal{F}) = 0$ が $i > 0$ に対して成り立ち，かつ
$\epsilon_{X, *}a_X^{-1}\mathcal{F} = \pi_X^{-1}\mathcal{F}$ である．補題
\ref{etale-cohomology-lemma-cohomological-descent-etale} により，
$R^j\pi_{X, *}(\pi_X^{-1}\mathcal{F}) = 0$ が $j > 0$ に対して成り立つ．これで証明は完了する．
\end{proof}

\begin{lemma}
\label{etale-cohomology-lemma-compare-cohomology-etale-fppf}
スキーム $X$ と上の射
$a_X : \Sh((\Sch/X)_{fppf}) \to \Sh(X_\etale)$ に対して，次が成り立つ：
\begin{enumerate}
\item $H^q(X_\etale, \mathcal{F}) = H^q_{fppf}(X, a_X^{-1}\mathcal{F})$
がアーベル層 $\mathcal{F}$，すなわち $X_\etale$ 上の層に対して成り立つ．
\item $H^q(X_\etale, K) = H^q_{fppf}(X, a_X^{-1}K)$ が
$K \in D^+(X_\etale)$ に対して成り立つ．
\end{enumerate}
たとえば，$A$ がアーベル群ならば
$H^q_\etale(X, \underline{A}) = H^q_{fppf}(X, \underline{A})$.
\end{lemma}

\begin{proof}
これは補題 \ref{etale-cohomology-lemma-cohomological-descent-etale-fppf} と Cohomology on Sites，
注意 \ref{sites-cohomology-remark-before-Leray} により従う．
\end{proof}







\section{fppf 位相と \'etale 位相の比較：加群}
\label{etale-cohomology-section-fppf-etale-modules}

\noindent
節 \ref{etale-cohomology-section-fppf-etale} の議論を続けるが，本節では加群の層に対して何が起こるかを
簡潔に論じる．

\medskip\noindent
$S$ をスキームとする．サイトの射 $\epsilon_S$，$\pi_S$，およびその合成 $a_S$ は，
節 \ref{etale-cohomology-section-fppf-etale} で導入したものであり，環付きサイトの射へ自然に拡張される．
第1の射を
$$
\epsilon_S :
((\Sch/S)_{fppf}, \mathcal{O})
\longrightarrow
((\Sch/S)_\etale, \mathcal{O})
$$
と書く．実際これらの構造層は台となる前層が同じなので，同じ記号を用いてよいことに
注意する．第2の射は
$$
\pi_S :
((\Sch/S)_\etale, \mathcal{O})
\longrightarrow
(S_\etale, \mathcal{O}_S)
$$
である．第3の射は
$$
a_S :
((\Sch/S)_{fppf}, \mathcal{O})
\longrightarrow
(S_\etale, \mathcal{O}_S)
$$
である．スキーム $S$ 上の準連接加群の圏が
$(S_\etale, \mathcal{O}_S)$ 上の準連接加群の圏と同じであることは既に知っている．
Descent，命題 \ref{descent-proposition-equivalence-quasi-coherent} を参照せよ．
ここでは \'etale コホモロジーと fppf コホモロジーの比較を述べたいので，本節の残りでは
準連接層を小 \'etale サイト上のものとして考える．これらのサイト上の準連接加群について
既知の事実を復習しよう．

\begin{lemma}
\label{etale-cohomology-lemma-review-quasi-coherent}
$S$ をスキームとする．$\mathcal{F}$ を準連接 $\mathcal{O}_S$-加群とし，
$S_\etale$ 上にあるものとする．
\begin{enumerate}
\item 規則
$$
\mathcal{F}^a : (\Sch/S)_\etale \longrightarrow \textit{Ab},\quad
(f : T \to S) \longmapsto \Gamma(T, f_{small}^*\mathcal{F})
$$
は fppf 被覆，したがって特に \'etale 被覆に対する層条件を満たす．
\item $\mathcal{F}^a = \pi_S^*\mathcal{F}$ が $(\Sch/S)_\etale$ 上で成り立つ．
\item $\mathcal{F}^a = a_S^*\mathcal{F}$ が $(\Sch/S)_{fppf}$ 上で成り立つ．
\item 規則 $\mathcal{F} \mapsto \mathcal{F}^a$ は，準連接 $\mathcal{O}_S$-加群と
$((\Sch/S)_\etale, \mathcal{O})$ 上の準連接加群との圏同値を定める．
\item 規則 $\mathcal{F} \mapsto \mathcal{F}^a$ は，準連接 $\mathcal{O}_S$-加群と
$((\Sch/S)_{fppf}, \mathcal{O})$ 上の準連接加群との圏同値を定める．
\item $\epsilon_{S, *}a_S^*\mathcal{F} = \pi_S^*\mathcal{F}$ かつ
$a_{S, *}a_S^*\mathcal{F} = \mathcal{F}$ である．
\item $R^i\epsilon_{S, *}(a_S^*\mathcal{F}) = 0$ かつ
$R^ia_{S, *}(a_S^*\mathcal{F}) = 0$ が $i > 0$ に対して成り立つ．
\end{enumerate}
\end{lemma}

\begin{proof}
読者には，Descent の諸節 \ref{descent-section-quasi-coherent-sheaves}，
\ref{descent-section-quasi-coherent-cohomology}，および
\ref{descent-section-quasi-coherent-sheaves-bis} の内容に基づいて，これらの結果の証明を
自ら見いだすことを勧める．

\medskip\noindent
まず，この補題の記号 $\mathcal{F}^a$ が，節 \ref{etale-cohomology-section-quasi-coherent} および
\ref{etale-cohomology-section-cohomology-quasi-coherent}，ならびに Descent，節
\ref{descent-section-quasi-coherent-sheaves} における以前の用法と
整合する理由を説明する．すなわち Descent，命題
\ref{descent-proposition-equivalence-quasi-coherent} により，準連接加群
$\mathcal{F}_0$ がスキーム $S$ 上（換言すれば小 Zariski サイト上）に存在し，
$\mathcal{F}$ は規則
$$
\mathcal{F}_0^a : (\Sch/S)_\etale \longrightarrow \textit{Ab},\quad
(f : T \to S) \longmapsto \Gamma(T, f^*\mathcal{F})
$$
を部分圏 $S_\etale \subset (\Sch/S)_\etale$ に制限したものである．ここで $f^*$ は
スキーム上の加群の層の通常の引戻しを表す．$\mathcal{F}_0^a$ は環付きサイトの射
$$
((\Sch/S)_\etale, \mathcal{O}) \longrightarrow (S_{Zar}, \mathcal{O}_{S_{Zar}})
$$
による引戻しであるから，Descent，注意
\ref{descent-remark-change-topologies-ringed-sites} により，引戻しの合成から直ちに
$\mathcal{F}^a = \mathcal{F}_0^a$ が従う．これは Descent，補題
\ref{descent-lemma-sheaf-condition-holds} により，fpqc 被覆に対してすら層条件を証明する
（命題 \ref{etale-cohomology-proposition-quasi-coherent-sheaf-fpqc} も参照せよ）．次いで (2) と (3) は再び
Descent，注意 \ref{descent-remark-change-topologies-ringed-sites} により，(4) と (5) は
Descent，命題 \ref{descent-proposition-equivalence-quasi-coherent} により従う
（メタ結果である定理 \ref{etale-cohomology-theorem-quasi-coherent} も参照せよ）．

\medskip\noindent
(6) は層 $\mathcal{F}^a = \pi_S^*\mathcal{F} = a_S^*\mathcal{F}$ の記述から直ちに従う．

\medskip\noindent
任意のアーベル層 $\mathcal{H}$ を $(\Sch/S)_{fppf}$ 上にとる．高次順像
$R^p\epsilon_{S, *}\mathcal{H}$ は，前層 $U \mapsto H^p_{fppf}(U, \mathcal{H})$ に
付随する $(\Sch/S)_\etale$ 上の層である．Cohomology on Sites，補題
\ref{sites-cohomology-lemma-higher-direct-images} を参照せよ．したがって
$R^p\epsilon_{S, *}a_S^*\mathcal{F} = R^p\epsilon_{S, *}\mathcal{F}^a = 0$
を $p > 0$ に対して証明するには，任意のスキーム $U$ が $S$ 上にあり，\'etale 被覆
$\{U_i \to U\}_{i \in I}$ をもち，$H^p_{fppf}(U_i, \mathcal{F}^a) = 0$ が
$p > 0$ に対して成り立つことを示せば十分である．アフィンスキームによる開被覆をとれば，
必要な消滅は通常のコホモロジーとの比較（Descent，命題
\ref{descent-proposition-same-cohomology-quasi-coherent} または定理
\ref{etale-cohomology-theorem-zariski-fpqc-quasi-coherent}）と，Cohomology of Schemes，補題
\ref{coherent-lemma-quasi-coherent-affine-cohomology-zero} によるアフィンスキーム上の
準連接層のコホモロジーの消滅から従う．

\medskip\noindent
$R^pa_{S, *}a_S^{-1}\mathcal{F} = R^pa_{S, *}\mathcal{F}^a = 0$ を
$p > 0$ に対して示すには，全く同様に論じる．これで証明は完了する．
\end{proof}

\begin{lemma}
\label{etale-cohomology-lemma-cohomological-descent-etale-fppf-modules}
$S$ をスキームとする．$\mathcal{F}$ を準連接 $\mathcal{O}_S$-加群とし，
$S_\etale$ 上にとる．このとき写像
$$
\pi_S^*\mathcal{F} \longrightarrow R\epsilon_{S, *}(a_S^*\mathcal{F})
\quad\text{および}\quad
\mathcal{F} \longrightarrow Ra_{S, *}(a_S^*\mathcal{F})
$$
は同型である．ここで
$a_S : \Sh((\Sch/S)_{fppf}) \to \Sh(S_\etale)$ は上の射である．
\end{lemma}

\begin{proof}
これは補題 \ref{etale-cohomology-lemma-review-quasi-coherent} の (6) と (7) から直ちに従う．
\end{proof}

\begin{lemma}
\label{etale-cohomology-lemma-cohomological-descent-complex-modules}
$S = \Spec(A)$ をアフィンスキームとする．$M^\bullet$ を $A$-加群の複体とする．
複体 $\mathcal{F}^\bullet$ を考える．これは $\mathcal{O}$-加群の前層からなり，
$(\textit{Aff}/S)_{fppf}$ 上で規則
$$
(U/S) = (\Spec(B)/\Spec(A)) \longmapsto M^\bullet \otimes_A B
$$
によって与えられる．これは加群の複体であり，標準的な写像
$$
M^\bullet \longrightarrow
R\Gamma((\textit{Aff}/S)_{fppf}, \mathcal{F}^\bullet)
$$
は擬同型である．
\end{lemma}

\begin{proof}
各 $\mathcal{F}^n$ は加群の層である．実際，これは加群
$\mathcal{G}^n = (\widetilde{M}^n)^a$（補題
\ref{etale-cohomology-lemma-review-quasi-coherent} のもの）を
$(\textit{Aff}/S)_{fppf} \subset (\Sch/S)_{fppf}$ に制限したものと一致する．
この包含は環付きトポスの圏同値を定めるので（Topologies，補題
\ref{topologies-lemma-affine-big-site-fppf}），
$$
R\Gamma((\textit{Aff}/S)_{fppf}, \mathcal{F}^\bullet) =
R\Gamma((\Sch/S)_{fppf}, \mathcal{G}^\bullet)
$$
を得る．$M^\bullet = R\Gamma(S, \widetilde{M}^\bullet)$ であることに注意する．これは，
たとえば Derived Categories of Schemes，補題
\ref{perfect-lemma-affine-compare-bounded} による．したがって比較写像
$$
R\Gamma(S, \widetilde{M}^\bullet)
\longrightarrow
R\Gamma((\Sch/S)_{fppf}, (\widetilde{M}^\bullet)^a)
$$
が同型であることを示そうとしている．$M^\bullet$ が下に有界ならば，これは Descent，命題
\ref{descent-proposition-same-cohomology-quasi-coherent}
および Derived Categories，補題 \ref{derived-lemma-two-ss-complex-functor} の
第1スペクトル系列により成り立つ．一般の場合には
$M^\bullet = \lim M_n^\bullet$ と書き，
$M_n^\bullet = \tau_{\geq -n}M^\bullet$ とする．すると系 $M_n^p$ は最終的に
$M^p$ を値とする定値系になる．次を主張する：
$$
(\widetilde{M}^\bullet)^a = R\lim (\widetilde{M}_n^\bullet)^a
$$
実際，左辺から右辺への自然な写像が，任意のアフィン対象 $U = \Spec(B)$，すなわち
$(\Sch/S)_{fppf}$ の対象の上のコホモロジーに同型を誘導することを示せば十分である．
$i \in \mathbf{Z}$ と $n > |i|$ に対して
$$
H^i(U, (\widetilde{M}_n^\bullet)^a) =
H^i(\tau_{\geq -n}M^\bullet \otimes_A B) =
H^i(M^\bullet \otimes_A B)
$$
である．第1の等号は上で扱った下に有界な場合から従う．したがって Cohomology on Sites，
補題 \ref{sites-cohomology-lemma-RGamma-commutes-with-Rlim} により主張が成り立つ．
最後に
\begin{align*}
R\Gamma((\Sch/S)_{fppf}, (\widetilde{M}^\bullet)^a)
& =
R\Gamma((\Sch/S)_{fppf}, R\lim (\widetilde{M}_n^\bullet)^a) \\
& =
R\lim R\Gamma((\Sch/S)_{fppf}, (\widetilde{M}_n^\bullet)^a) \\
& =
R\lim M_n^\bullet \\
& =
M^\bullet
\end{align*}
を望むとおり得る．第2の等号は $R\lim$ が $R\Gamma$ と可換することから従う．
Cohomology on Sites，補題
\ref{sites-cohomology-lemma-RGamma-commutes-with-Rlim} を参照せよ．
\end{proof}










\section{ph 位相と \'etale 位相の比較}
\label{etale-cohomology-section-ph-etale}

\noindent
本節のひな型は，Cohomology on Sites，節
\ref{sites-cohomology-section-cohomology-LC} で論じられた，局所コンパクト位相空間上の
通常の位相と qc 位相の比較である．まず，Topologies の諸節
\ref{topologies-section-change-topologies} および
\ref{topologies-section-etale} の内容をいくつか復習する．

\medskip\noindent
$S$ をスキームとし，$(\Sch/S)_{ph}$ を ph サイトとする．同じ台となる圏には
第二の位相，すなわち \'etale 位相があり，したがって第二のサイト
$(\Sch/S)_\etale$ がある．恒等関手
$(\Sch/S)_\etale \to (\Sch/S)_{ph}$ は連続であり（More on Morphisms，補題
\ref{more-morphisms-lemma-fppf-ph} および Topologies，補題
\ref{topologies-lemma-zariski-etale-smooth-syntomic-fppf} による），サイトの射
$$
\epsilon_S : (\Sch/S)_{ph} \longrightarrow (\Sch/S)_\etale
$$
を定める．Cohomology on Sites，節 \ref{sites-cohomology-section-compare} を参照せよ．
$\epsilon_{S, *}$ は台となる前層上の恒等関手であり，$\epsilon_S^{-1}$ は
\'etale 層にその ph 層化を対応させることに注意する．$S_\etale$ を小 \'etale サイトとする．
サイトの射
$$
\pi_S : (\Sch/S)_\etale \longrightarrow S_\etale
$$
があり，これは連続関手 $S_\etale \to (\Sch/S)_\etale$，$U \mapsto U$ により与えられる．
実際，$S_\etale$ はファイバー積と終対象をもち，上の関手はこれらと可換するので，
Sites，命題 \ref{sites-proposition-get-morphism} が適用できる．

\begin{lemma}
\label{etale-cohomology-lemma-describe-pullback-pi-ph}
上の記法を用いる．$\mathcal{F}$ を $S_\etale$ 上の層とする．規則
$$
(\Sch/S)_{ph} \longrightarrow \textit{Sets},\quad
(f : X \to S) \longmapsto \Gamma(X, f_{small}^{-1}\mathcal{F})
$$
は層であり，したがってもちろん $(\Sch/S)_\etale$ 上の層でもある．実際，この層は
$\pi_S^{-1}\mathcal{F}$ に等しく，これは $(\Sch/S)_\etale$ 上の等式である；また
$\epsilon_S^{-1}\pi_S^{-1}\mathcal{F}$ に等しく，これは $(\Sch/S)_{ph}$ 上の等式である．
\end{lemma}

\begin{proof}
\'etale 位相に関する主張は補題 \ref{etale-cohomology-lemma-describe-pullback} の内容である．証明を
終えるには，$\pi_S^{-1}\mathcal{F}$ が ph 位相に対する層であることを示せば十分である．
Topologies，補題 \ref{topologies-lemma-characterize-sheaf} により，全射な固有射
$V \to U$ が $S$ 上のスキームの射として与えられたとき，次が等化子図式であることを
示せば十分である：
$$
\xymatrix{
(\pi_S^{-1}\mathcal{F})(U) \ar[r] &
(\pi_S^{-1}\mathcal{F})(V) \ar@<1ex>[r] \ar@<-1ex>[r] &
(\pi_S^{-1}\mathcal{F})(V \times_U V)
}
$$
$\mathcal{G} = \pi_S^{-1}\mathcal{F}|_{U_\etale}$ とおく．可換図式
$$
\xymatrix{
V \times_U V \ar[r] \ar[rd]_g \ar[d] & V \ar[d]^f \\
V \ar[r]^f & U
}
$$
を考える．次が成り立つ：
$$
(\pi_S^{-1}\mathcal{F})(V) = \Gamma(V, f^{-1}\mathcal{G}) =
\Gamma(U, f_*f^{-1}\mathcal{G})
$$
ここで $f_*$ と $f^{-1}$ は小 \'etale サイト間の関手性を表すために用いている．次に，
$$
(\pi_S^{-1}\mathcal{F})(V \times_U V) =
\Gamma(V \times_U V, g^{-1}\mathcal{G}) =
\Gamma(U, g_*g^{-1}\mathcal{G})
$$
である．等化子図式の2つの写像は，2つの写像
$$
f_*f^{-1}\mathcal{G} \longrightarrow g_*g^{-1}\mathcal{G}
$$
から得られる．したがって $\mathcal{G}$ がこれら2つの層の写像の等化子であることを
証明すれば十分である．$\overline{u}$ を $U$ の幾何点とし，
$\Omega = \mathcal{G}_{\overline{u}}$ とおく．補題
\ref{etale-cohomology-lemma-proper-pushforward-stalk} により $\overline{u}$ で茎をとると，2つの写像
$$
H^0(V_{\overline{u}}, \underline{\Omega}) \longrightarrow
H^0((V \times_U V)_{\overline{u}}, \underline{\Omega}) =
H^0(V_{\overline{u}} \times_{\overline{u}} V_{\overline{u}},
\underline{\Omega})
$$
を得る．ここで $\underline{\Omega}$ は値 $\Omega$ の定値層を表す．もちろん，これらの
写像は射影による引戻しである．そこで，第1因子への射影による引戻しから来る切断は
第1射影のファイバー上で定値であり，また第1因子への射影による引戻しから来る切断は
第1射影のファイバー上で定値であることは明らかである．これらの引戻し写像の像の共通部分に
属する切断は $V_{\overline{u}} \times_{\overline{u}} V_{\overline{u}}$ 全体で定値であり，
すなわち望むとおり $\Omega$ の元から来る．
\end{proof}

\noindent
補題 \ref{etale-cohomology-lemma-describe-pullback-pi-ph} の状況では，$\epsilon_S$ と $\pi_S$ の合成と
上の等式により，サイトの射
$$
a_S : (\Sch/S)_{ph} \longrightarrow S_\etale
$$

\begin{lemma}
\label{etale-cohomology-lemma-push-pull-ph-etale}
上の記法を用いる．$f : X \to Y$ を $(\Sch/S)_{ph}$ の射とする．
このときトポスの可換図式
$$
\xymatrix{
\Sh((\Sch/X)_{ph}) \ar[rr]_{f_{big, ph}} \ar[d]_{\epsilon_X} & &
\Sh((\Sch/Y)_{ph}) \ar[d]^{\epsilon_Y} \\
\Sh((\Sch/X)_\etale) \ar[rr]^{f_{big, \etale}} & &
\Sh((\Sch/Y)_\etale)
}
$$
および
$$
\xymatrix{
\Sh((\Sch/X)_{ph}) \ar[rr]_{f_{big, ph}} \ar[d]_{a_X} & &
\Sh((\Sch/Y)_{ph}) \ar[d]^{a_Y} \\
\Sh(X_\etale) \ar[rr]^{f_{small}} & &
\Sh(Y_\etale)
}
$$
があり，$a_X = \pi_X \circ \epsilon_X$ および $a_Y = \pi_X \circ \epsilon_X$ を満たす．
\end{lemma}

\begin{proof}
図式の可換性は Topologies，節 \ref{topologies-section-change-topologies} の議論から従う．
\end{proof}

\begin{lemma}
\label{etale-cohomology-lemma-proper-push-pull-ph-etale}
補題 \ref{etale-cohomology-lemma-push-pull-ph-etale} において $f$ が固有ならば，
$a_Y^{-1} \circ f_{small, *} = f_{big, ph, *} \circ a_X^{-1}$.
\end{lemma}

\begin{proof}
補題 \ref{etale-cohomology-lemma-compare-higher-direct-image-proper} (1) の証明を繰り返しても証明できるが，
ここではこの結果から導く．$\epsilon_{Y, *}$ は台となる前層上の恒等関手なので，
同型を反映する．補題 \ref{etale-cohomology-lemma-describe-pullback-pi-ph} の記述により
$\epsilon_{Y, *} \circ a_Y^{-1} = \pi_Y^{-1}$ であり，$X$ についても同様である．
標準的な写像
$a_Y^{-1}f_{small, *}\mathcal{F} \to f_{big, ph, *}a_X^{-1}\mathcal{F}$ が同型であることを
示すには，
\begin{align*}
\pi_Y^{-1}f_{small, *}\mathcal{F}
& =
\epsilon_{Y, *}a_Y^{-1}f_{small, *}\mathcal{F} \\
& \to 
\epsilon_{Y, *}f_{big, ph, *}a_X^{-1}\mathcal{F} \\
& =
f_{big, \etale, *} \epsilon_{X, *}a_X^{-1}\mathcal{F} \\
& =
f_{big, \etale, *}\pi_X^{-1}\mathcal{F}
\end{align*}
が同型であることを示せば十分である．これは補題
\ref{etale-cohomology-lemma-compare-higher-direct-image-proper} (1) である．
\end{proof}

\begin{lemma}
\label{etale-cohomology-lemma-compare-ph-etale}
比較射 $\epsilon : (\Sch/S)_{ph} \to (\Sch/S)_\etale$ を考える．
$\mathcal{P}$ をスキームの固有射全体の類とする．$X$ を $(\Sch/S)_\etale$ の対象とする．
$\mathcal{A}'_X \subset \textit{Ab}((\Sch/X)_\etale)$ を，形
$\pi_X^{-1}\mathcal{F}$ の層からなる充満部分圏とする．ここで $\mathcal{F}$ は
$X_\etale$ 上の捩れアーベル層である．このとき Cohomology on Sites の性質
(\ref{sites-cohomology-item-base-change-P}),
(\ref{sites-cohomology-item-restriction-A}),
(\ref{sites-cohomology-item-A-sheaf}),
(\ref{sites-cohomology-item-A-and-P})，および
(\ref{sites-cohomology-item-refine-tau-by-P})
が，Cohomology on Sites，状況 \ref{sites-cohomology-situation-compare} において成り立つ．
\end{lemma}

\begin{proof}
まず，Homology，補題 \ref{homology-lemma-characterize-weak-serre-subcategory} の
条件 (1)，(2)，(3)，(4) を確認して，
$\mathcal{A}'_X \subset \textit{Ab}((\Sch/X)_\etale)$ が弱 Serre 部分圏であることを示す．
たとえば補題 \ref{etale-cohomology-lemma-cohomological-descent-etale} により $\pi_X^{-1}$ は完全かつ
充満忠実なので，(1)，(2)，(3) は直ちに従う．
$0 \to \pi_X^{-1}\mathcal{F} \to \mathcal{G} \to \pi_X^{-1}\mathcal{F}' \to 0$
が $\textit{Ab}((\Sch/X)_\etale)$ の短完全列ならば，補題
\ref{etale-cohomology-lemma-cohomological-descent-etale} により
$0 \to \mathcal{F} \to \pi_{X, *}\mathcal{G} \to \mathcal{F}' \to 0$
は完全である．特に $\pi_{X, *}\mathcal{G}$ は $X_\etale$ 上の捩れアーベル層である．
したがって $\mathcal{G} = \pi_X^{-1}\pi_{X, *}\mathcal{G}$ は $\mathcal{A}'_X$ に属し，
最後の条件も確認できる．

\medskip\noindent
Cohomology on Sites の性質 (\ref{sites-cohomology-item-base-change-P}) は，
スキームのファイバー積が存在すること，およびスキームの固有射の基底変換が固有射であることから
従う．Morphisms，補題 \ref{morphisms-lemma-base-change-proper} を参照せよ．

\medskip\noindent
Cohomology on Sites の性質 (\ref{sites-cohomology-item-restriction-A}) は，
Topologies，補題 \ref{topologies-lemma-morphism-big-small-etale} の可換図式 (3) から従う．

\medskip\noindent
Cohomology on Sites の性質 (\ref{sites-cohomology-item-A-sheaf}) は補題
\ref{etale-cohomology-lemma-describe-pullback-pi-ph} である．

\medskip\noindent
Cohomology on Sites の性質 (\ref{sites-cohomology-item-A-and-P}) は補題
\ref{etale-cohomology-lemma-compare-higher-direct-image-proper} (2)，および
$R^if_{small}\mathcal{F}$ が捩れであるという事実から従う．ここで $\mathcal{F}$ は
$X_\etale$ 上の捩れアーベル層である．補題 \ref{etale-cohomology-lemma-torsion-direct-image} を参照せよ．

\medskip\noindent
Cohomology on Sites の性質 (\ref{sites-cohomology-item-refine-tau-by-P}) は，
More on Morphisms，補題 \ref{more-morphisms-lemma-dominate-fppf-etale-locally} と，
有限射は固有であり，全射な固有射は ph 被覆を定めるという事実を合わせれば従う．
Topologies，補題 \ref{topologies-lemma-surjective-proper-ph} を参照せよ．
\end{proof}

\begin{lemma}
\label{etale-cohomology-lemma-V-C-all-n-etale-ph}
上の記号を用いる．
\begin{enumerate}
\item $X \in \Ob((\Sch/S)_{ph})$ と捩れアーベル層 $\mathcal{F}$ をとり，後者は
$X_\etale$ 上の層とする．このとき
$\epsilon_{X, *}a_X^{-1}\mathcal{F} = \pi_X^{-1}\mathcal{F}$
かつ $R^i\epsilon_{X, *}(a_X^{-1}\mathcal{F}) = 0$ が $i > 0$ に対して成り立つ．
\item 固有射 $f : X \to Y$ を $(\Sch/S)_{ph}$ でとり，捩れアーベル層
$\mathcal{F}$ を $X$ 上にとる．このとき
$a_Y^{-1}(R^if_{small, *}\mathcal{F}) =
R^if_{big, ph, *}(a_X^{-1}\mathcal{F})$
がすべての $i$ について成り立つ．
\item スキーム $X$ と $K$ をとり，後者は捩れコホモロジー層をもつ
$D^+(X_\etale)$ の対象とする．このとき写像
$\pi_X^{-1}K \to R\epsilon_{X, *}(a_X^{-1}K)$ は同型である．
\item スキームの固有射 $f : X \to Y$ と $K$ をとり，後者は捩れコホモロジー層をもつ
$D^+(X_\etale)$ の対象とする．このとき
$a_Y^{-1}(Rf_{small, *}K) = Rf_{big, ph, *}(a_X^{-1}K)$.
\end{enumerate}
\end{lemma}

\begin{proof}
補題 \ref{etale-cohomology-lemma-compare-ph-etale} により，Cohomology on Sites，節
\ref{sites-cohomology-section-compare-general} の補題はすべて現在の状況に適用できる．
結果を読み替えるため，Cohomology on Sites，補題
\ref{sites-cohomology-lemma-A} の圏 $\mathcal{A}_X$ は，
$\textit{Ab}((\Sch/X)_{ph})$ のうち，形 $a_X^{-1}\mathcal{F}$ の層からなる充満部分圏で
あることに注意する．ここで $\mathcal{F}$ は $X_\etale$ 上の捩れアーベル層である．

\medskip\noindent
(1) は $(V_n)$ がすべての $n$ について成り立つことと同値であり，これは
Cohomology on Sites，補題 \ref{sites-cohomology-lemma-V-C-all-n-general} により成り立つ．

\medskip\noindent
(2) は Cohomology on Sites，補題
\ref{sites-cohomology-lemma-V-implies-C-general} の結論に $\epsilon_Y^{-1}$ を
適用することから従う．

\medskip\noindent
(3) は Cohomology on Sites，補題
\ref{sites-cohomology-lemma-V-C-all-n-general} (1) から従う．実際，
$\pi_X^{-1}K$ は $D^+_{\mathcal{A}'_X}((\Sch/X)_\etale)$ に属し，
$a_X^{-1} = \epsilon_X^{-1} \circ a_X^{-1}$ である．

\medskip\noindent
(4) も同じ理由により Cohomology on Sites，補題
\ref{sites-cohomology-lemma-V-C-all-n-general} (2) から従う．
\end{proof}

\begin{lemma}
\label{etale-cohomology-lemma-cohomological-descent-etale-ph}
$X$ をスキームとする．捩れコホモロジー層をもつ $K \in D^+(X_\etale)$ に対し，写像
$$
K \longrightarrow Ra_{X, *}a_X^{-1}K
$$
は同型である．ここで $a_X : \Sh((\Sch/X)_{ph}) \to \Sh(X_\etale)$ は上の射である．
\end{lemma}

\begin{proof}
まず主張を，$K$ が1つのアーベル層で与えられる場合に帰着させる．実際，$K$ を
捩れアーベル層からなる下に有界な複体 $\mathcal{F}^\bullet$ で表す．これは
Cohomology on Sites，補題 \ref{sites-cohomology-lemma-torsion} により可能である．
層の場合から
$\mathcal{F}^n = a_{X, *} a_X^{-1} \mathcal{F}^n$
であり，層 $R^qa_{X, *}a_X^{-1}\mathcal{F}^n$ は $q > 0$ に対して零である．
$a_X^{-1}\mathcal{F}^\bullet$ と関手 $a_{X, *}$ に Leray の非輪状性補題
（Derived Categories，補題 \ref{derived-lemma-leray-acyclicity}）を適用すれば結論を得る．
以下では $K = \mathcal{F}$ と仮定し，$\mathcal{F}$ は捩れアーベル層とする．

\medskip\noindent
補題 \ref{etale-cohomology-lemma-describe-pullback-pi-ph} により
$a_{X, *}a_X^{-1}\mathcal{F} = \mathcal{F}$ である．したがって
$R^qa_{X, *}a_X^{-1}\mathcal{F} = 0$ が $q > 0$ に対して成り立つことを示せば十分である．
これには $a_X = \epsilon_X \circ \pi_X$ と Leray スペクトル系列（Cohomology on Sites，
補題 \ref{sites-cohomology-lemma-relative-Leray}）を用いることができる．補題
\ref{etale-cohomology-lemma-V-C-all-n-etale-ph} により，
$R^i\epsilon_{X, *}(a_X^{-1}\mathcal{F}) = 0$ が $i > 0$ に対して成り立ち，かつ
$\epsilon_{X, *}a_X^{-1}\mathcal{F} = \pi_X^{-1}\mathcal{F}$ である．補題
\ref{etale-cohomology-lemma-cohomological-descent-etale} により，
$R^j\pi_{X, *}(\pi_X^{-1}\mathcal{F}) = 0$ が $j > 0$ に対して成り立つ．これで証明は完了する．
\end{proof}

\begin{lemma}
\label{etale-cohomology-lemma-compare-cohomology-etale-ph}
スキーム $X$ と上の射 $a_X : \Sh((\Sch/X)_{ph}) \to \Sh(X_\etale)$ に対して，
次が成り立つ：
\begin{enumerate}
\item $H^q(X_\etale, \mathcal{F}) = H^q_{ph}(X, a_X^{-1}\mathcal{F})$
が捩れアーベル層 $\mathcal{F}$，すなわち $X_\etale$ 上の層に対して成り立つ．
\item $H^q(X_\etale, K) = H^q_{ph}(X, a_X^{-1}K)$
が，捩れコホモロジー層をもつ $K \in D^+(X_\etale)$ に対して成り立つ．
\end{enumerate}
たとえば，$A$ が捩れアーベル群ならば
$H^q_\etale(X, \underline{A}) = H^q_{ph}(X, \underline{A})$.
\end{lemma}

\begin{proof}
これは補題 \ref{etale-cohomology-lemma-cohomological-descent-etale-ph} と Cohomology on Sites，
注意 \ref{sites-cohomology-remark-before-Leray} により従う．
\end{proof}












\section{h 位相と \'etale 位相の比較}
\label{etale-cohomology-section-h-etale}

\noindent
本節のひな型は，Cohomology on Sites，節
\ref{sites-cohomology-section-cohomology-LC} で論じられた，局所コンパクト位相空間上の
通常の位相と qc 位相の比較である．さらに本節は，ph 位相と \'etale 位相を比較する節と
ほとんど逐語的に同じである．まず，Topologies の諸節
\ref{topologies-section-change-topologies} および
\ref{topologies-section-etale}，ならびに More on Flatness，節
\ref{flat-section-h} の内容をいくつか復習する．

\medskip\noindent
$S$ をスキームとし，$(\Sch/S)_h$ を h サイトとする．同じ台となる圏には第二の位相，
すなわち \'etale 位相があり，したがって第二のサイト $(\Sch/S)_\etale$ がある．恒等関手
$(\Sch/S)_\etale \to (\Sch/S)_h$ は連続であり（More on Flatness，補題
\ref{flat-lemma-zariski-h} および Topologies，補題
\ref{topologies-lemma-zariski-etale-smooth-syntomic-fppf} による），サイトの射
$$
\epsilon_S : (\Sch/S)_h \longrightarrow (\Sch/S)_\etale
$$
を定める．Cohomology on Sites，節 \ref{sites-cohomology-section-compare} を参照せよ．
$\epsilon_{S, *}$ は台となる前層上の恒等関手であり，$\epsilon_S^{-1}$ は
\'etale 層にその h 層化を対応させることに注意する．$S_\etale$ を小 \'etale サイトとする．
サイトの射
$$
\pi_S : (\Sch/S)_\etale \longrightarrow S_\etale
$$
があり，これは連続関手 $S_\etale \to (\Sch/S)_\etale$，$U \mapsto U$ により与えられる．
実際，$S_\etale$ はファイバー積と終対象をもち，上の関手はこれらと可換するので，
Sites，命題 \ref{sites-proposition-get-morphism} が適用できる．

\begin{lemma}
\label{etale-cohomology-lemma-describe-pullback-pi-h}
上の記法を用いる．$\mathcal{F}$ を $S_\etale$ 上の層とする．規則
$$
(\Sch/S)_h \longrightarrow \textit{Sets},\quad
(f : X \to S) \longmapsto \Gamma(X, f_{small}^{-1}\mathcal{F})
$$
は層であり，したがってもちろん $(\Sch/S)_\etale$ 上の層でもある．実際，この層は
$\pi_S^{-1}\mathcal{F}$ に等しく，これは $(\Sch/S)_\etale$ 上の等式である；また
$\epsilon_S^{-1}\pi_S^{-1}\mathcal{F}$ に等しく，これは $(\Sch/S)_h$ 上の等式である．
\end{lemma}

\begin{proof}
\'etale 位相に関する主張は補題 \ref{etale-cohomology-lemma-describe-pullback} の内容である．証明を
終えるには，$\pi_S^{-1}\mathcal{F}$ が h 位相に対する層であることを示せば十分である．
しかし補題 \ref{etale-cohomology-lemma-describe-pullback-pi-ph} では，
$\pi_S^{-1}\mathcal{F}$ がより強い ph 位相に対してすら層であることを示した．
\end{proof}

\noindent
補題 \ref{etale-cohomology-lemma-describe-pullback-pi-h} の状況では，$\epsilon_S$ と $\pi_S$ の合成と
上の等式により，サイトの射
$$
a_S : (\Sch/S)_h \longrightarrow S_\etale
$$

\begin{lemma}
\label{etale-cohomology-lemma-push-pull-h-etale}
上の記法を用いる．$f : X \to Y$ を $(\Sch/S)_h$ の射とする．
このときトポスの可換図式
$$
\xymatrix{
\Sh((\Sch/X)_h) \ar[rr]_{f_{big, h}} \ar[d]_{\epsilon_X} & &
\Sh((\Sch/Y)_h) \ar[d]^{\epsilon_Y} \\
\Sh((\Sch/X)_\etale) \ar[rr]^{f_{big, \etale}} & &
\Sh((\Sch/Y)_\etale)
}
$$
および
$$
\xymatrix{
\Sh((\Sch/X)_h) \ar[rr]_{f_{big, h}} \ar[d]_{a_X} & &
\Sh((\Sch/Y)_h) \ar[d]^{a_Y} \\
\Sh(X_\etale) \ar[rr]^{f_{small}} & &
\Sh(Y_\etale)
}
$$
があり，$a_X = \pi_X \circ \epsilon_X$ および $a_Y = \pi_X \circ \epsilon_X$ を満たす．
\end{lemma}

\begin{proof}
図式の可換性は Topologies，節 \ref{topologies-section-change-topologies} の議論と
同様に従う．
\end{proof}

\begin{lemma}
\label{etale-cohomology-lemma-proper-push-pull-h-etale}
補題 \ref{etale-cohomology-lemma-push-pull-h-etale} において $f$ が固有ならば，
$a_Y^{-1} \circ f_{small, *} = f_{big, h, *} \circ a_X^{-1}$.
\end{lemma}

\begin{proof}
補題 \ref{etale-cohomology-lemma-compare-higher-direct-image-proper} (1) の証明を繰り返しても証明できるが，
ここではこの結果から導く．$\epsilon_{Y, *}$ は台となる前層上の恒等関手なので，
同型を反映する．補題 \ref{etale-cohomology-lemma-describe-pullback-pi-h} の記述により
$\epsilon_{Y, *} \circ a_Y^{-1} = \pi_Y^{-1}$ であり，$X$ についても同様である．
標準的な写像
$a_Y^{-1}f_{small, *}\mathcal{F} \to f_{big, h, *}a_X^{-1}\mathcal{F}$ が同型であることを
示すには，
\begin{align*}
\pi_Y^{-1}f_{small, *}\mathcal{F}
& =
\epsilon_{Y, *}a_Y^{-1}f_{small, *}\mathcal{F} \\
& \to 
\epsilon_{Y, *}f_{big, h, *}a_X^{-1}\mathcal{F} \\
& =
f_{big, \etale, *} \epsilon_{X, *}a_X^{-1}\mathcal{F} \\
& =
f_{big, \etale, *}\pi_X^{-1}\mathcal{F}
\end{align*}
が同型であることを示せば十分である．これは補題
\ref{etale-cohomology-lemma-compare-higher-direct-image-proper} (1) である．
\end{proof}

\begin{lemma}
\label{etale-cohomology-lemma-compare-h-etale}
比較射 $\epsilon : (\Sch/S)_h \to (\Sch/S)_\etale$ を考える．
$\mathcal{P}$ を固有射全体の類とする．$X$ を $(\Sch/S)_\etale$ の対象とする．
$\mathcal{A}'_X \subset \textit{Ab}((\Sch/X)_\etale)$ を，形
$\pi_X^{-1}\mathcal{F}$ の層からなる充満部分圏とする．ここで $\mathcal{F}$ は
$X_\etale$ 上の捩れアーベル層である．このとき Cohomology on Sites の性質
(\ref{sites-cohomology-item-base-change-P}),
(\ref{sites-cohomology-item-restriction-A}),
(\ref{sites-cohomology-item-A-sheaf}),
(\ref{sites-cohomology-item-A-and-P})，および
(\ref{sites-cohomology-item-refine-tau-by-P})
が，Cohomology on Sites，状況 \ref{sites-cohomology-situation-compare} において成り立つ．
\end{lemma}

\begin{proof}
まず，Homology，補題 \ref{homology-lemma-characterize-weak-serre-subcategory} の
条件 (1)，(2)，(3)，(4) を確認して，
$\mathcal{A}'_X \subset \textit{Ab}((\Sch/X)_\etale)$ が弱 Serre 部分圏であることを示す．
たとえば補題 \ref{etale-cohomology-lemma-cohomological-descent-etale} により $\pi_X^{-1}$ は完全かつ
充満忠実なので，(1)，(2)，(3) は直ちに従う．
$0 \to \pi_X^{-1}\mathcal{F} \to \mathcal{G} \to \pi_X^{-1}\mathcal{F}' \to 0$
が $\textit{Ab}((\Sch/X)_\etale)$ の短完全列ならば，補題
\ref{etale-cohomology-lemma-cohomological-descent-etale} により
$0 \to \mathcal{F} \to \pi_{X, *}\mathcal{G} \to \mathcal{F}' \to 0$
は完全である．特に $\pi_{X, *}\mathcal{G}$ は $X_\etale$ 上の捩れアーベル層である．
したがって $\mathcal{G} = \pi_X^{-1}\pi_{X, *}\mathcal{G}$ は $\mathcal{A}'_X$ に属し，
最後の条件も確認できる．

\medskip\noindent
Cohomology on Sites の性質 (\ref{sites-cohomology-item-base-change-P}) は，
スキームのファイバー積が存在すること，スキームの固有射の基底変換が固有射であること
（Morphisms，補題 \ref{morphisms-lemma-base-change-proper}），および有限表示射の基底変換が
有限表示射であること（Morphisms，補題
\ref{morphisms-lemma-base-change-finite-presentation}）から従う．

\medskip\noindent
Cohomology on Sites の性質 (\ref{sites-cohomology-item-restriction-A}) は，
Topologies，補題 \ref{topologies-lemma-morphism-big-small-etale} の可換図式 (3) から従う．

\medskip\noindent
Cohomology on Sites の性質 (\ref{sites-cohomology-item-A-sheaf}) は補題
\ref{etale-cohomology-lemma-describe-pullback-pi-h} である．

\medskip\noindent
Cohomology on Sites の性質 (\ref{sites-cohomology-item-A-and-P}) は補題
\ref{etale-cohomology-lemma-compare-higher-direct-image-proper} (2)，および
$R^if_{small}\mathcal{F}$ が捩れであるという事実から従う．ここで $\mathcal{F}$ は
$X_\etale$ 上の捩れアーベル層である．補題 \ref{etale-cohomology-lemma-torsion-direct-image} を参照せよ．

\medskip\noindent
Cohomology on Sites の性質 (\ref{sites-cohomology-item-refine-tau-by-P}) は，
More on Morphisms，補題 \ref{more-morphisms-lemma-dominate-fppf-etale-locally} と，
全射な有限局所自由射は全射，固有，かつ有限表示であり，したがって More on Flatness，
補題 \ref{flat-lemma-surjective-proper-finite-presentation-h} により h 被覆を定めるという
事実を合わせれば従う．
\end{proof}

\begin{lemma}
\label{etale-cohomology-lemma-V-C-all-n-etale-h}
上の記号を用いる．
\begin{enumerate}
\item $X \in \Ob((\Sch/S)_{h})$ と捩れアーベル層 $\mathcal{F}$ をとり，後者は
$X_\etale$ 上の層とする．このとき
$\epsilon_{X, *}a_X^{-1}\mathcal{F} = \pi_X^{-1}\mathcal{F}$
かつ $R^i\epsilon_{X, *}(a_X^{-1}\mathcal{F}) = 0$ が $i > 0$ に対して成り立つ．
\item 固有射 $f : X \to Y$ を $(\Sch/S)_h$ でとり，捩れアーベル層
$\mathcal{F}$ を $X$ 上にとる．このとき
$a_Y^{-1}(R^if_{small, *}\mathcal{F}) =
R^if_{big, h, *}(a_X^{-1}\mathcal{F})$
がすべての $i$ について成り立つ．
\item スキーム $X$ と $K$ をとり，後者は捩れコホモロジー層をもつ
$D^+(X_\etale)$ の対象とする．このとき写像
$\pi_X^{-1}K \to R\epsilon_{X, *}(a_X^{-1}K)$ は同型である．
\item スキームの固有射 $f : X \to Y$ と $K$ をとり，後者は捩れコホモロジー層をもつ
$D^+(X_\etale)$ の対象とする．このとき
$a_Y^{-1}(Rf_{small, *}K) = Rf_{big, h, *}(a_X^{-1}K)$.
\end{enumerate}
\end{lemma}

\begin{proof}
補題 \ref{etale-cohomology-lemma-compare-h-etale} により，Cohomology on Sites，節
\ref{sites-cohomology-section-compare-general} の補題はすべて現在の状況に適用できる．
結果を読み替えるため，Cohomology on Sites，補題
\ref{sites-cohomology-lemma-A} の圏 $\mathcal{A}_X$ は，
$\textit{Ab}((\Sch/X)_h)$ のうち，形 $a_X^{-1}\mathcal{F}$ の層からなる充満部分圏で
あることに注意する．ここで $\mathcal{F}$ は $X_\etale$ 上の捩れアーベル層である．

\medskip\noindent
(1) は $(V_n)$ がすべての $n$ について成り立つことと同値であり，これは
Cohomology on Sites，補題 \ref{sites-cohomology-lemma-V-C-all-n-general} により成り立つ．

\medskip\noindent
(2) は Cohomology on Sites，補題
\ref{sites-cohomology-lemma-V-implies-C-general} の結論に $\epsilon_Y^{-1}$ を
適用することから従う．

\medskip\noindent
(3) は Cohomology on Sites，補題
\ref{sites-cohomology-lemma-V-C-all-n-general} (1) から従う．実際，
$\pi_X^{-1}K$ は $D^+_{\mathcal{A}'_X}((\Sch/X)_\etale)$ に属し，
$a_X^{-1} = \epsilon_X^{-1} \circ a_X^{-1}$ である．

\medskip\noindent
(4) も同じ理由により Cohomology on Sites，補題
\ref{sites-cohomology-lemma-V-C-all-n-general} (2) から従う．
\end{proof}

\begin{lemma}
\label{etale-cohomology-lemma-cohomological-descent-etale-h}
$X$ をスキームとする．捩れコホモロジー層をもつ $K \in D^+(X_\etale)$ に対し，写像
$$
K \longrightarrow Ra_{X, *}a_X^{-1}K
$$
は同型である．ここで $a_X : \Sh((\Sch/X)_h) \to \Sh(X_\etale)$ は上の射である．
\end{lemma}

\begin{proof}
まず主張を，$K$ が1つのアーベル層で与えられる場合に帰着させる．実際，$K$ を
捩れアーベル層からなる下に有界な複体 $\mathcal{F}^\bullet$ で表す．これは
Cohomology on Sites，補題 \ref{sites-cohomology-lemma-torsion} により可能である．
層の場合から
$\mathcal{F}^n = a_{X, *} a_X^{-1} \mathcal{F}^n$
であり，層 $R^qa_{X, *}a_X^{-1}\mathcal{F}^n$ は $q > 0$ に対して零である．
$a_X^{-1}\mathcal{F}^\bullet$ と関手 $a_{X, *}$ に Leray の非輪状性補題
（Derived Categories，補題 \ref{derived-lemma-leray-acyclicity}）を適用すれば結論を得る．
以下では $K = \mathcal{F}$ と仮定し，$\mathcal{F}$ は捩れアーベル層とする．

\medskip\noindent
補題 \ref{etale-cohomology-lemma-describe-pullback-pi-h} により
$a_{X, *}a_X^{-1}\mathcal{F} = \mathcal{F}$ である．したがって
$R^qa_{X, *}a_X^{-1}\mathcal{F} = 0$ が $q > 0$ に対して成り立つことを示せば十分である．
これには $a_X = \epsilon_X \circ \pi_X$ と Leray スペクトル系列（Cohomology on Sites，
補題 \ref{sites-cohomology-lemma-relative-Leray}）を用いることができる．補題
\ref{etale-cohomology-lemma-V-C-all-n-etale-h} により，
$R^i\epsilon_{X, *}(a_X^{-1}\mathcal{F}) = 0$ が $i > 0$ に対して成り立ち，かつ
$\epsilon_{X, *}a_X^{-1}\mathcal{F} = \pi_X^{-1}\mathcal{F}$ である．補題
\ref{etale-cohomology-lemma-cohomological-descent-etale} により，
$R^j\pi_{X, *}(\pi_X^{-1}\mathcal{F}) = 0$ が $j > 0$ に対して成り立つ．これで証明は完了する．
\end{proof}

\begin{lemma}
\label{etale-cohomology-lemma-compare-cohomology-etale-h}
スキーム $X$ と上の射 $a_X : \Sh((\Sch/X)_h) \to \Sh(X_\etale)$ に対して，
次が成り立つ：
\begin{enumerate}
\item $H^q(X_\etale, \mathcal{F}) = H^q_h(X, a_X^{-1}\mathcal{F})$
が捩れアーベル層 $\mathcal{F}$，すなわち $X_\etale$ 上の層に対して成り立つ．
\item $H^q(X_\etale, K) = H^q_h(X, a_X^{-1}K)$
が，捩れコホモロジー層をもつ $K \in D^+(X_\etale)$ に対して成り立つ．
\end{enumerate}
たとえば，$A$ が捩れアーベル群ならば
$H^q_\etale(X, \underline{A}) = H^q_h(X, \underline{A})$.
\end{lemma}

\begin{proof}
これは補題 \ref{etale-cohomology-lemma-cohomological-descent-etale-h} と Cohomology on Sites，
注意 \ref{sites-cohomology-remark-before-Leray} により従う．
\end{proof}











\section{\'etale 層の降下}
\label{etale-cohomology-section-glueing-etale}

\noindent
\'etale 層が fppf 位相および h 位相で「貼り合わさる」こと，ならびに関連する結果を
証明する．既に次の関連結果を示した：
\begin{enumerate}
\item 補題 \ref{etale-cohomology-lemma-describe-pullback} によれば，スキーム $S$ の小 \'etale サイト上の
層は，$S$ の大 \'etale サイト上の層で，fpqc 被覆（したがって特に Zariski，\'etale，
平滑，syntomic，および fppf 被覆）に対する層条件を満たすものを定める．
\item 補題 \ref{etale-cohomology-lemma-describe-pullback-pi-fppf} は前項を fppf 位相について言い換えた
ものである．
\item 補題 \ref{etale-cohomology-lemma-describe-pullback-pi-ph} は ph 位相について同じことを証明する．
\item 補題 \ref{etale-cohomology-lemma-describe-pullback-pi-h} は h 位相について同じことを証明する．
\item 補題 \ref{etale-cohomology-lemma-descent-sheaf-fppf-etale} は \'etale 層に対する fppf 降下の一形態である．
\item 注意 \ref{etale-cohomology-remark-cohomological-descent-finite} によれば，有限全射に対して \'etale 層の
降下が成り立つ（以下でこれを明確化し，強める）．
\end{enumerate}
単体的空間の章では，これについてさらにいくつかの結果を証明する．たとえば
Simplicial Spaces の諸節 \ref{spaces-simplicial-section-fppf-hypercovering} および
\ref{spaces-simplicial-section-proper-hypercovering-spaces} を参照せよ．

\medskip\noindent
結果を便利に表すため，いくつか記号が必要である．
$\mathcal{U} = \{f_i : X_i \to X\}$ を，共通の終域をもつスキームの射の族とする．
$\mathcal{U}$ に関する \'etale 層の {\it 降下データ}とは，族
$((\mathcal{F}_i)_{i \in I}, (\varphi_{ij})_{i, j \in I})$ で，次を満たすものである：
\begin{enumerate}
\item $\mathcal{F}_i$ は $\Sh(X_{i, \etale})$ に属する．
\item $\varphi_{ij} :
\text{pr}_{0, small}^{-1} \mathcal{F}_i
\longrightarrow
\text{pr}_{1, small}^{-1} \mathcal{F}_j$
$\Sh((X_i \times_X X_j)_\etale)$ における同型である．
\end{enumerate}
さらに {\it コサイクル条件}が成り立つ，すなわち図式
$$
\xymatrix{
\text{pr}_{0, small}^{-1}\mathcal{F}_i
\ar[dr]_{\text{pr}_{02, small}^{-1}\varphi_{ik}}
\ar[rr]^{\text{pr}_{01, small}^{-1}\varphi_{ij}} & &
\text{pr}_{1, small}^{-1}\mathcal{F}_j
\ar[dl]^{\text{pr}_{12, small}^{-1}\varphi_{jk}} \\
& \text{pr}_{2, small}^{-1}\mathcal{F}_k
}
$$
は $\Sh((X_i \times_X X_j \times_X X_k)_\etale)$ において可換である．
{\it 降下データの射}の明らかな概念があり，降下データの圏を得る．降下データ
$((\mathcal{F}_i)_{i \in I}, (\varphi_{ij})_{i, j \in I})$
が {\it 有効}であるとは，$\mathcal{F}$ が $\Sh(X_\etale)$ に存在し，同型
$\varphi_i : f_{i, small}^{-1} \mathcal{F} \to \mathcal{F}_i$
が $\Sh(X_{i, \etale})$ に存在し，$\varphi_{ij}$ と両立する，すなわち
$$
\varphi_{ij} =
\text{pr}_{1, small}^{-1} (\varphi_j) \circ
\text{pr}_{0, small}^{-1} (\varphi_i^{-1})
$$
ことをいう．同じことを別の仕方で述べよう．対象 $\mathcal{F}$ が
$\Sh(X_\etale)$ に与えられると，
{\it 標準降下データ} $(f_{i, small}^{-1}\mathcal{F}_i, c_{ij})$ を得る．
ここで $c_{ij}$ は標準同型
$$
c_{ij} : \text{pr}_{0, small}^{-1} f_{i, small}^{-1}\mathcal{F}
\longrightarrow
\text{pr}_{1, small}^{-1} f_{j, small}^{-1}\mathcal{F}
$$
である．降下データ
$((\mathcal{F}_i)_{i \in I}, (\varphi_{ij})_{i, j \in I})$
が有効であることと，ある $\mathcal{F}$ に付随する標準降下データと同型であることは，
その対象が $\Sh(X_\etale)$ に属するとき同値である．

\medskip\noindent
族がただ1つの射 $\{X \to Y\}$ からなる場合，降下データを対
$(\mathcal{F}, \varphi)$ とみなす．ここで $\mathcal{F}$ は $\Sh(X_\etale)$ の対象であり，
$\varphi$ は同型
$$
\text{pr}_{0, small}^{-1} \mathcal{F}
\longrightarrow
\text{pr}_{1, small}^{-1} \mathcal{F}
$$
であって，$\Sh((X \times_Y X)_\etale)$ に属し，コサイクル条件を満たす，すなわち
$$
\xymatrix{
\text{pr}_{0, small}^{-1}\mathcal{F}
\ar[dr]_{\text{pr}_{02, small}^{-1}\varphi}
\ar[rr]^{\text{pr}_{01, small}^{-1}\varphi} & &
\text{pr}_{1, small}^{-1}\mathcal{F}
\ar[dl]^{\text{pr}_{12, small}^{-1}\varphi} \\
& \text{pr}_{2, small}^{-1}\mathcal{F}
}
$$
が $\Sh((X \times_Y X \times_Y X)_\etale)$ において可換である．降下データの射と
有効性の概念は先ほどと全く同じように定める．

\medskip\noindent
まず，全射な整射に対する有効降下を証明する．

\begin{lemma}
\label{etale-cohomology-lemma-glue-etale-sheaf-section}
$f : X \to Y$ を切断をもつスキームの射とする．このとき関手
$$
\Sh(Y_\etale)
\longrightarrow
\text{\'etale 層の降下データ（対象族：}\{X \to Y\}\text{）}
$$
$\mathcal{G}$ を標準降下データに送る関手は，その対象が $\Sh(Y_\etale)$ に
属するとき圏同値である．
\end{lemma}

\begin{proof}
これは形式的であり，引戻し関手の関手性のみに依存する．詳細は省略する．ヒント：
$s : Y \to X$ が切断ならば，$(\mathcal{F}, \varphi)$ を
$s_{small}^{-1}\mathcal{F}$ に送る関手が擬逆である．
\end{proof}

\begin{lemma}
\label{etale-cohomology-lemma-glue-etale-sheaf-integral-surjective}
$f : X \to Y$ をスキームの全射な整射とする．関手
$$
\Sh(Y_\etale)
\longrightarrow
\text{\'etale 層の降下データ（対象族：}\{X \to Y\}\text{）}
$$
は圏同値である．
\end{lemma}

\begin{proof}
この証明では，引戻し関手と順像関手の添字 ${}_{small}$ を省略する．
$X_1 = X \times_Y X$ と書き，$f_1 : X_1 \to Y$ を射
$f \circ \text{pr}_0 = f \circ \text{pr}_1$ とする．
$(\mathcal{F}, \varphi)$ を $\{X \to Y\}$ に関する降下データとする．
$\mathcal{F}_1 = \text{pr}_0^{-1}\mathcal{F}$ とおく．$\varphi$ は同型
$\mathcal{F}_1 \to \text{pr}_1^{-1}\mathcal{F}$ を定めるとみなせる．降下データ
$(\mathcal{F}, \varphi)$ を層
$$
\mathcal{G} =
\text{等化子}\left(
\xymatrix{
f_*\mathcal{F}
\ar@<1ex>[r] \ar@<-1ex>[r] &
f_{1, *}\mathcal{F}_1
}
\right)
$$
に送る規則が，補題の主張にある関手の擬逆であると主張する．2本の矢印のうち第1は写像
$$
f_*\mathcal{F} \to
f_*\text{pr}_{0, *}\text{pr}_0^{-1}\mathcal{F} =
f_{1, *}\mathcal{F}_1
$$
から来て，第2の矢印は写像
$$
f_*\mathcal{F} \to
f_* \text{pr}_{1, *}\text{pr}_1^{-1}\mathcal{F}
\xleftarrow{\varphi}
f_* \text{pr}_{0, *} \text{pr}_0^{-1}\mathcal{F} =
f_{1, *}\mathcal{F}_1
$$
から来る．ここで左向きの矢印は可逆である．これが機能することを証明するには，標準写像
$f^{-1}\mathcal{G} \to \mathcal{F}$ が同型であることを示さなければならない；詳細は
省略する．これを証明するには，任意の射の族 $\Spec(k) \to Y$ で引き戻した後に確認すれば
十分である．ここで $k$ は代数閉体である．実際，対応する基底変換
$X_k \to X$ は共同全射であり，$X_\etale$ 上の層の写像が同型かどうかは幾何点での
茎を見れば確認できる．定理 \ref{etale-cohomology-theorem-exactness-stalks} を参照せよ．補題
\ref{etale-cohomology-lemma-integral-pushforward-commutes-with-base-change} により，$\mathcal{G}$ を降下データ
$(\mathcal{F}, \varphi)$ から構成する操作は任意の基底変換と可換する．
したがって $Y$ は代数閉体のスペクトルであると仮定してよい（基底変換は射 $f$ の性質を
保つことに注意する．Morphisms の補題 \ref{morphisms-lemma-base-change-surjective} および
\ref{morphisms-lemma-base-change-finite} を参照せよ）．この場合，射 $X \to Y$ は切断をもち，
補題 \ref{etale-cohomology-lemma-glue-etale-sheaf-section} により関手が圏同値であることが分かる．一方，
関手が圏同値であることと上の構成が擬逆であることが同値であることは読者が示せる；
詳細は省略する．これで証明は完了する．
\end{proof}

\begin{lemma}
\label{etale-cohomology-lemma-glue-etale-sheaf-proper-surjective}
$f : X \to Y$ をスキームの全射な固有射とする．関手
$$
\Sh(Y_\etale)
\longrightarrow
\text{\'etale 層の降下データ（対象族：}\{X \to Y\}\text{）}
$$
は圏同値である．
\end{lemma}

\begin{proof}
補題 \ref{etale-cohomology-lemma-glue-etale-sheaf-integral-surjective} で与えたものと全く同じ証明が使える．
ただし補題 \ref{etale-cohomology-lemma-integral-pushforward-commutes-with-base-change} の使用を，補題
\ref{etale-cohomology-lemma-proper-base-change-f-star} の使用に置き換える．
\end{proof}

\begin{lemma}
\label{etale-cohomology-lemma-glue-etale-sheaf-check-after-base-change}
$f : X \to Y$ をスキームの射とする．$Z \to Y$ をスキームの全射な整射または
全射な固有射とする．関手
$$
\Sh(Z_\etale)
\longrightarrow
\text{\'etale 層の降下データ（対象族：}\{X \times_Y Z \to Z\}\text{）}
$$
および
$$
\Sh((Z \times_Y Z)_\etale)
\longrightarrow
\text{\'etale 層の降下データ（対象族：}
\{X \times_Y (Z \times_Y Z) \to Z \times_Y Z\}\text{）}
$$
が圏同値ならば，
$$
\Sh(Y_\etale)
\longrightarrow
\text{\'etale 層の降下データ（対象族：}\{X \to Y\}\text{）}
$$
は圏同値である．
\end{lemma}

\begin{proof}
これは定義と補題 \ref{etale-cohomology-lemma-glue-etale-sheaf-integral-surjective} および
\ref{etale-cohomology-lemma-glue-etale-sheaf-proper-surjective} の形式的帰結である．詳細は省略する．
\end{proof}

\begin{lemma}
\label{etale-cohomology-lemma-glue-etale-sheaf-fppf-cover}
$f : X \to Y$ を，全射，平坦，かつ有限表示局所的なスキームの射とする．関手
$$
\Sh(Y_\etale)
\longrightarrow
\text{\'etale 層の降下データ（対象族：}\{X \to Y\}\text{）}
$$
は圏同値である．
\end{lemma}

\begin{proof}
補題 \ref{etale-cohomology-lemma-glue-etale-sheaf-integral-surjective} の証明と全く同様に，
$(\mathcal{F}, \varphi)$ を
$$
\mathcal{G} =
\text{等化子}\left(
\xymatrix{
f_*\mathcal{F}
\ar@<1ex>[r] \ar@<-1ex>[r] &
f_{1, *}\mathcal{F}_1
}
\right)
$$
に送る関手が擬逆を与えると主張する．これを証明するには
$f^{-1}\mathcal{G} \to \mathcal{F}$ が同型であることを示せば十分である．これは局所的に
確認できるので，$Y$ はアフィンであると仮定する．すると有限全射 $Z \to Y$ と開被覆
$Z = \bigcup W_i$ で，各 $W_i \to Y$ が $X$ を経由するものを見いだせる．
More on Morphisms，補題 \ref{more-morphisms-lemma-dominate-fppf-finite} を参照せよ．
補題 \ref{etale-cohomology-lemma-glue-etale-sheaf-check-after-base-change} を適用すると，$Y$ を $Z$ および
$Z \times_Y Z$ で置き換え，$f$ をその基底変換で置き換えた後に補題を証明すれば十分である．
したがって $f$ は Zariski 局所的に切断をもつと仮定してよい．もちろん，問題が $Y$ 上
局所的であることを用いれば，切断をもつ場合，すなわち補題
\ref{etale-cohomology-lemma-glue-etale-sheaf-section} に帰着する．
\end{proof}

\begin{lemma}
\label{etale-cohomology-lemma-glue-etale-sheaf-fppf}
$\{f_i : X_i \to X\}$ をスキームの fppf 被覆とする．関手
$$
\Sh(X_\etale)
\longrightarrow
\text{\'etale 層の降下データ（対象族：}\{f_i : X_i \to X\}\text{）}
$$
は圏同値である．
\end{lemma}

\begin{proof}
射 $f : \coprod X_i \to X$ に補題 \ref{etale-cohomology-lemma-glue-etale-sheaf-fppf-cover} を適用する．
すると形式的な議論により，$f$ に関する降下データは被覆に関する降下データと同じもので
あることが分かる．Descent，補題 \ref{descent-lemma-family-is-one} と比較せよ．詳細は省略する．
\end{proof}

\begin{lemma}
\label{etale-cohomology-lemma-glue-etale-sheaf-modification}
$f : X' \to X$ をスキームの固有射とする．$i : Z \to X$ を閉埋め込みとし，
$E = Z \times_X X'$ とおく．図式
$$
\xymatrix{
E \ar[d]_g \ar[r]_j & X' \ar[d]^f \\
Z \ar[r]^i & X
}
$$
を考える．$f$ が $X \setminus Z$ 上で同型ならば，関手
$$
\Sh(X_\etale)
\longrightarrow
\Sh(X'_\etale) \times_{\Sh(E_\etale)} \Sh(Z_\etale)
$$
は圏同値である．
\end{lemma}

\begin{proof}
Categories，例 \ref{categories-example-2-fibre-product-categories} で構成した
$2$-ファイバー積圏を用いる．この関手は $\mathcal{F}$ を三つ組
$(f^{-1}\mathcal{F}, i^{-1}\mathcal{F}, c)$ に送る．ここで
$c : j^{-1}f^{-1}\mathcal{F} \to g^{-1}i^{-1}\mathcal{F}$
は標準同型である．擬逆関手を構成しよう．
$(\mathcal{F}', \mathcal{G}, \alpha)$ を矢印の右辺の対象とする．
同型
$$
i^{-1}f_*\mathcal{F}' = g_*j^{-1}\mathcal{F}'
\xrightarrow{g_*\alpha}
g_*g^{-1}\mathcal{G}
$$
を得る．最初の等式は補題 \ref{etale-cohomology-lemma-proper-base-change-f-star} である．
これを用いて射
$i_*\mathcal{G} \to i_*g_*g^{-1}\mathcal{G}$
および $f'_*\mathcal{F}' \to i_*g_*g^{-1}\mathcal{G}$ を得る．
$$
\mathcal{F} = f_*\mathcal{F}' \times_{i_*g_*g^{-1}\mathcal{G}} i_*\mathcal{G}
$$
とおく．$\mathcal{F}$ は矢印の左辺の対象であり，その右辺における像は出発点の
三つ組と同型であると主張する．$\mathcal{F}$ の，幾何学的点 $\overline{x}$（$X$ の点）
における茎を計算しよう．

\medskip\noindent
$\overline{x}$ が $Z$ に属さないなら，一方では $\overline{x}$ は一意な幾何学的点
$\overline{x}'$（$X'$ の点）から生じ，
$\mathcal{F}'_{\overline{x}'} = (f_*\mathcal{F}')_{\overline{x}}$
である．他方，$(i_*\mathcal{G})_{\overline{x}}$ および
$(i_*g_*g^{-1}\mathcal{G})_{\overline{x}}$ は一点集合である．したがって
$\mathcal{F}_{\overline{x}}$ は $\mathcal{F}'_{\overline{x}'}$ に等しい．

\medskip\noindent
$\overline{x}$ が $Z$ に属する，すなわち $\overline{x}$ が幾何学的点 $\overline{z}$
（$Z$ の点）の像であるなら，
$(i_*\mathcal{G})_{\overline{x}} = \mathcal{G}_{\overline{z}}$
および
$$
(i_*g_*g^{-1}\mathcal{G})_{\overline{x}} =
(g_*g^{-1}\mathcal{G})_{\overline{z}} =
\Gamma(E_{\overline{z}}, g^{-1}\mathcal{G}|_{E_{\overline{z}}})
$$
を得る（ここでは上で用いた順像に対する固有基底変換を使った）．同様に
$$
(f_*\mathcal{F}')_{\overline{x}} =
\Gamma(X'_{\overline{x}}, \mathcal{F}'|_{X'_{\overline{x}}})
$$
を得る．同一視 $E_{\overline{z}} = X'_{\overline{x}}$ があり，$\alpha$ は層
$\mathcal{F}'|_{X'_{\overline{x}}}$ と $g^{-1}\mathcal{G}|_{E_{\overline{z}}}$ の間の
同型を定めるので，この場合には
$$
\mathcal{F}_{\overline{x}} = \mathcal{G}_{\overline{z}}
$$
を得る．

\medskip\noindent
証明を終えるため，標準射
$i^{-1}\mathcal{F} \to \mathcal{G}$ および $f^{-1}\mathcal{F} \to \mathcal{F}'$
があり，これらは $\alpha$ と両立し，茎上では上で得た同型を与えることを確認する．
これらの射の詳しい構成は省略する．
\end{proof}

\begin{lemma}
\label{etale-cohomology-lemma-h-descent-etale-sheaves}
$S$ をスキームとする．亜群にファイバー化された圏
$$
p : \mathcal{S} \longrightarrow (\Sch/S)_h
$$
であって，$U$ 上のファイバー圏が圏 $\Sh(U_\etale)$，すなわち $U$ の小 \'etale
サイト上の層の圏であるものは，亜群におけるスタックである．
\end{lemma}

\begin{proof}
補題を証明するため，More on Flatness，補題
\ref{flat-lemma-refine-check-h-stack} の条件 (1)，(2)，(3) を確認する．

\medskip\noindent
層には貼り合わせがあり（また Zariski 被覆は \'etale 被覆である）ので，条件 (1) は成り立つ．
Sites，補題 \ref{sites-lemma-glue-sheaves} を参照せよ．

\medskip\noindent
条件 (2) を見るため，$f : X \to Y$ は $S$ 上有限表示の全射，平坦かつ固有な射で，
$Y$ はアフィンであると仮定する．すると補題
\ref{etale-cohomology-lemma-glue-etale-sheaf-fppf-cover} または補題
\ref{etale-cohomology-lemma-glue-etale-sheaf-proper-surjective} により，$\{X \to Y\}$ に対する降下が成り立つ．

\medskip\noindent
条件 (3) は，より一般的な補題 \ref{etale-cohomology-lemma-glue-etale-sheaf-modification} から直ちに従う．
\end{proof}



















\section{ブローアップ正方形と \'etale コホモロジー}
\label{etale-cohomology-section-blow-up-square}

\noindent
ブローアップ正方形は More on Flatness，節 \ref{flat-section-blow-up-ph} で導入した．
固有基底変換定理を用いると，ブローアップ正方形について Mayer--Vietoris 型の
結果が成り立つことが分かる．

\begin{lemma}
\label{etale-cohomology-lemma-blow-up-square-cohomological-descent}
$X$ をスキームとし，$Z \subset X$ を有限型準連接イデアルによって切り出される
閉部分スキームとする．対応するブローアップ正方形
$$
\xymatrix{
E \ar[d]_\pi \ar[r]_j & X' \ar[d]^b \\
Z \ar[r]^i & X
}
$$
を考える．捩れコホモロジー層をもつ $K \in D^+(X_\etale)$ に対して，
識別三角形
$$
K \to Ri_*(K|_Z) \oplus Rb_*(K|_{X'}) \to Rc_*(K|_E) \to K[1]
$$
が $D(X_\etale)$ に存在する．ここで $c = i \circ \pi = b \circ j$ である．
\end{lemma}

\begin{proof}
記号 $K|_{X'}$ は $b_{small}^{-1}K$ を表す．アーベル層の下に有界な複体
$\mathcal{F}^\bullet$ で $K$ を表現するものを選ぶ．$i_*(\mathcal{F}^\bullet|_Z)$ は
$Ri_*(K|_Z)$ を表現することに注意せよ．実際，$i_*$ は完全である
（命題 \ref{etale-cohomology-proposition-finite-higher-direct-image-zero}）．擬同型
$b_{small}^{-1}\mathcal{F}^\bullet \to \mathcal{I}^\bullet$
を選ぶ．ここで $\mathcal{I}^\bullet$ は $X'_\etale$ 上の入射的アーベル層からなる
下に有界な複体である．この射は射 $\mathcal{F}^\bullet \to b_*(\mathcal{I}^\bullet)$
に随伴し，$b_*(\mathcal{I}^\bullet)$ は $Rb_*(K|_{X'})$ を表現する．補題
\ref{etale-cohomology-lemma-proper-base-change-f-star} により
$\pi_*(\mathcal{I}^\bullet|_E) = (b_*\mathcal{I}^\bullet)|_Z$ であり，補題
\ref{etale-cohomology-lemma-proper-base-change} によりこの複体は $R\pi_*(K|_E)$ を表現する．
したがって射
$$
Ri_*(K|_Z) \oplus Rb_*(K|_{X'}) \to Rc_*(K|_E)
$$
は，下に有界な複体の全射
$$
i_*(\mathcal{F}^\bullet|_Z) \oplus
b_*(\mathcal{I}^\bullet)
\to
i_*\left(b_*(\mathcal{I}^\bullet)|_Z\right)
$$
によって表現される．所望の識別三角形を得るには，標準射
$\mathcal{F}^\bullet \to i_*(\mathcal{F}^\bullet|_Z) \oplus
b_*(\mathcal{I}^\bullet)$
が上に表示した複体の射の核へ擬同型に写ることを示せば十分である（すなわち，
複体の短完全列は導来圏における識別三角形を定める．Derived Categories，節
\ref{derived-section-canonical-delta-functor} を参照せよ）．これは幾何学的点
$\overline{x}$（$X$ の点）における茎上で確認できる．$\overline{x}$ が $Z$ に属さない
なら，$X' \to X$ は $\overline{x}$ のある開近傍上で同型である．この場合，対応する
幾何学的点を $\overline{x}'$（$X'$ の点）で表すと，示すべきことは
$$
\mathcal{F}^\bullet_{\overline{x}} \to \mathcal{I}^\bullet_{\overline{x}'}
$$
が擬同型であることである．これは $\mathcal{I}^\bullet$ の選び方から従う．
$\overline{x}$ が $Z$ に属するなら，
$b_(\mathcal{I}^\bullet)_{\overline{x}} \to 
i_*\left(b_*(\mathcal{I}^\bullet)|_Z\right)_{\overline{x}}$
はアーベル群の複体の同型である．したがって核は，所望どおり
$i_*(\mathcal{F}^\bullet|_Z)_{\overline{x}} =
\mathcal{F}^\bullet_{\overline{x}}$ に等しい．
\end{proof}

\begin{lemma}
\label{etale-cohomology-lemma-blow-up-square-etale-cohomology}
$X$ をスキームとし，$K \in D^+(X_\etale)$ は捩れコホモロジー層をもつとする．
$Z \subset X$ を有限型準連接イデアルによって切り出される閉部分スキームとする．
対応するブローアップ正方形
$$
\xymatrix{
E \ar[d] \ar[r] & X' \ar[d]^b \\
Z \ar[r] & X
}
$$
を考える．このとき標準的な長完全列
$$
H^p_\etale(X, K) \to
H^p_\etale(X', K|_{X'}) \oplus
H^p_\etale(Z, K|_Z) \to
H^p_\etale(E, K|_E) \to
H^{p + 1}_\etale(X, K)
$$
が存在する．
\end{lemma}

\begin{proof}[第1の証明]
これは補題 \ref{etale-cohomology-lemma-blow-up-square-cohomological-descent} と等式
$$
R\Gamma(X, Rb_*(K|_{X'})) = R\Gamma(X', K|_{X'})
$$
（Cohomology on Sites，節 \ref{sites-cohomology-section-leray} を参照せよ），
および他の項についての同様の等式から直ちに従う．
\end{proof}

\begin{proof}[第2の証明]
補題 \ref{etale-cohomology-lemma-compare-cohomology-etale-ph} により，これらのコホモロジー群は
$X, X', E, Z$ の，サイト $(\Sch/X)_{ph}$ 上のあるアーベル層の複体を係数とする
コホモロジーである（具体的には導来圏の対象 $a_X^{-1}K$ である．上の補題
\ref{etale-cohomology-lemma-describe-pullback-pi-ph} を参照し，$K|_{X'} = b_{small}^{-1}K$ を想起せよ）．
More on Flatness，補題 \ref{flat-lemma-blow-up-square-ph} により，表現可能前層の図式の
ph 層化は余 Cartesian である．したがって，サイト $(\Sch/X)_{ph}$ と補題の可換図式に
きわめて一般的な Cohomology on Sites，補題
\ref{sites-cohomology-lemma-square-triangle} を適用すれば主張が従う．
\end{proof}

\begin{lemma}
\label{etale-cohomology-lemma-blow-up-square-equivalence}
$X$ をスキームとし，$Z \subset X$ を有限型準連接イデアルによって切り出される
閉部分スキームとする．対応するブローアップ正方形
$$
\xymatrix{
E \ar[d]_\pi \ar[r]_j & X' \ar[d]^b \\
Z \ar[r]^i & X
}
$$
を考える．次が与えられていると仮定する：
\begin{enumerate}
\item $K'$，すなわち $D^+(X'_\etale)$ の捩れコホモロジー層をもつ対象，
\item $L$，すなわち $D^+(Z_\etale)$ の捩れコホモロジー層をもつ対象，
\item 同型 $\gamma : K'|_E \to L|_E$．
\end{enumerate}
このとき対象 $K$（$D^+(X_\etale)$ の対象）と同型 $f : K|_{X'} \to K'$，
$g : K|_Z \to L$ で，$\gamma = g|_E \circ f^{-1}|_E$ を満たすものが存在する．
さらに，次が与えられているとする：
\begin{enumerate}
\item $M$，すなわち $D^+(X_\etale)$ の捩れコホモロジー層をもつ対象，
\item 射 $\alpha : K' \to M|_{X'}$（$D(X'_\etale)$ における射），
\item 射 $\beta : L \to M|_Z$（$D(Z_\etale)$ における射）．
\end{enumerate}
これらが
$$
\alpha|_E  = \beta|_E \circ \gamma.
$$
を満たすなら，射 $M \to K$ が $D(X_\etale)$ に存在し，その $X'$ への制限は
$a \circ f$，$Z$ への制限は $b \circ g$ である．
\end{lemma}

\begin{proof}
$K$ が存在するなら，補題 \ref{etale-cohomology-lemma-blow-up-square-cohomological-descent} はそれが入る
識別三角形を与える．そこで単に識別三角形
$$
K \to Ri_*(L) \oplus Rb_*(K') \to Rc_*(L|_E) \to K[1]
$$
を選ぶ．ここで $c = i \circ \pi = b \circ j$ である．射
$Ri_*(L) \to Rc_*(L|_E)$ は，$Ri_*$ を随伴射 $E \to R\pi_*(L|_E)$ に適用したものである．
射 $Rb_*(K') \to Rc_*(L|_E)$ は標準射
$Rb_*(K') \to Rc_*(K'|_E)) = R$ と $Rc_*(\gamma)$ の合成である．補題の主張における
射 $g$ および $f$ はこれらの射の随伴である．この識別三角形を $Z$ に制限すると，
基底変換定理（補題 \ref{etale-cohomology-lemma-proper-base-change}）により射
$Rb_*(K) \to Rc_*(L|_E)$ は同型となるので，射 $g : K|_Z \to L$ は同型である．
識別三角形を見ると，$f : K|_{X'} \to K'$ は $X' \setminus E = X \setminus Z$ 上で
同型であることが分かる．さらに，構成により $\gamma \circ f|_E = g|_E$ である．
$\gamma$ と $g$ は同型なので，$f$ は $E$ の幾何学的点における茎上でも同型を誘導する．
したがって $f$ は同型である．

\medskip\noindent
最後の主張については，$K'$ を $K|_{X'}$ で，$L$ を $K|_Z$ で，$\gamma$ を標準的な
同一視で置き換えてよい．$\alpha$ と $\beta$ は可換正方形
$$
\xymatrix{
K \ar[r] \ar@{..>}[d] &
Ri_*(K|_Z) \oplus Rb_*(K|_{X'}) \ar[r] \ar[d]_{\beta \oplus \alpha} &
Rc_*(K|_E) \ar[r] \ar[d]_{\alpha|_E} &
K[1] \ar@{..>}[d] \\
M \ar[r] &
Ri_*(M|_Z) \oplus Rb_*(M|_{X'}) \ar[r] &
Rc_*(M|_E) \ar[r] &
M[1]
}
$$
を誘導する．したがって導来圏の公理により，識別三角形の射を与える点線の矢印を得る．
\end{proof}














\section{概ブローアップ正方形と h 位相}
\label{etale-cohomology-section-blow-up-h}

\noindent
本節では More on Flatness，節 \ref{flat-section-blow-up-h} の議論を続ける．
読者の便宜のため想起すると，概ブローアップ正方形とは可換図式
\begin{equation}
\label{etale-cohomology-equation-almost-blow-up-square}
\vcenter{
\xymatrix{
E \ar[d] \ar[r] & X' \ar[d]^b \\
Z \ar[r] & X
}
}
\end{equation}
であって，次の条件を満たすスキームの図式である：
\begin{enumerate}
\item $Z \to X$ は有限表示の閉埋め込みである，
\item $E = b^{-1}(Z)$ は $X'$ の局所主閉部分スキームである，
\item $b$ は固有かつ有限表示である，
\item $X'' \subset X'$ を，$\mathcal{O}_{X'}$ の切断で $E$ に台をもつものからなる準連接
イデアルによって切り出される閉部分スキームとすると（Properties，補題
\ref{properties-lemma-sections-supported-on-closed-subset}），これは $X$ の $Z$ における
ブローアップである．
\end{enumerate}
したがって射 $b$ は同型 $X' \setminus E \to X \setminus Z$ を誘導する．

\medskip\noindent
スキームの大 fppf サイト $(\Sch/S)_{fppf}$（ここで $S$ はスキーム）の導来圏の
対象が「h 層である」ための判定条件を与える．同じ基礎圏にはもう1つの位相，すなわち
h 位相がある（More on Flatness，節 \ref{flat-section-h}）．fppf 被覆は h 被覆であることを
想起せよ（More on Flatness，補題 \ref{flat-lemma-zariski-h}）．したがって射
$$
\epsilon : (\Sch/S)_h \longrightarrow (\Sch/S)_{fppf}
$$
を考えられる．Cohomology on Sites，節 \ref{sites-cohomology-section-compare} を参照せよ．
特に充満忠実関手
$$
R\epsilon_* : D((\Sch/S)_h) \longrightarrow D((\Sch/S)_{fppf})
$$
がある．この関手の本質像は何であろうか．

\begin{lemma}
\label{etale-cohomology-lemma-blow-up-square-h}
上の記号のもとで，$K$ が $R\epsilon_*$ の本質像に属するなら，Cohomology on Sites，
補題 \ref{sites-cohomology-lemma-c-square} の射 $c^K_{X, Z, X', E}$ は擬同型である．
\end{lemma}

\begin{proof}
${}^\#$ で h 位相における層化を表す．More on Flatness，補題
\ref{flat-lemma-blow-up-square-h} で
$h_X^\# = h_Z^\# \amalg_{h_E^\#} h_{X'}^\#$ を示した．他方，射
$h_E^\# \to h_{X'}^\#$ は，$E \to X'$ が単射なので単射である．したがって本補題は
Cohomology on Sites，補題
\ref{sites-cohomology-lemma-descent-squares-helper} の特別な場合である（この補題自体は
Cohomology on Sites，補題 \ref{sites-cohomology-lemma-square-triangle} の形式的帰結である）．
\end{proof}

\begin{proposition}
\label{etale-cohomology-proposition-check-h}
$K$ を $D^+((\Sch/S)_{fppf})$ の対象とする．このとき $K$ が
$R\epsilon_* : D((\Sch/S)_h) \to D((\Sch/S)_{fppf})$
の本質像に属するための必要十分条件は，$c^K_{X, X', Z, E}$ が，$(\Sch/S)_h$ における，
$X$ がアフィンである任意の概ブローアップ正方形
（\ref{etale-cohomology-equation-almost-blow-up-square}）について擬同型であることである．
\end{proposition}

\begin{proof}
Cohomology on Sites，補題 \ref{sites-cohomology-lemma-descent-squares} を適用して証明する．
その仮定は補題 \ref{etale-cohomology-lemma-blow-up-square-h} と More on Flatness，命題
\ref{flat-proposition-check-h} により成り立つ．
\end{proof}

\begin{lemma}
\label{etale-cohomology-lemma-refine-check-h}
$K$ を $D^+((\Sch/S)_{fppf})$ の対象とする．このとき $K$ が
$R\epsilon_* : D((\Sch/S)_h) \to D((\Sch/S)_{fppf})$ の本質像に属するための必要十分条件は，
More on Flatness，例 \ref{flat-example-one-generator} および
\ref{flat-example-two-generators} のような任意の概ブローアップ正方形について
$c^K_{X, X', Z, E}$ が擬同型であることである．
\end{lemma}

\begin{proof}
Cohomology on Sites，補題 \ref{sites-cohomology-lemma-descent-squares} を適用して証明する．
その仮定は補題 \ref{etale-cohomology-lemma-blow-up-square-h} と More on Flatness，補題
\ref{flat-lemma-refine-check-h} により成り立つ．
\end{proof}






\section{h 位相における構造層のコホモロジー}
\label{etale-cohomology-section-cohomology-O-h}

\noindent
$p$ を素数とする．$(\mathcal{C}, \mathcal{O})$ を $p\mathcal{O} = 0$ を満たす環付きサイトとする．
$\colim_F \mathcal{O}$ を，環の層の圏における系
$$
\mathcal{O} \xrightarrow{F} \mathcal{O} \xrightarrow{F}
\mathcal{O} \xrightarrow{F} \ldots
$$
の余極限とする．ここで $F : \mathcal{O} \to \mathcal{O}$，$f \mapsto f^p$ は
Frobenius 自己準同型である．

\begin{lemma}
\label{etale-cohomology-lemma-h-sheaf-colim-F}
$p$ を素数とする．$S$ を $\mathbf{F}_p$ 上のスキームとする．層
$\mathcal{O}^{perf} = \colim_F \mathcal{O}$ を $(\Sch/S)_{fppf}$ 上で考える．
このとき $\mathcal{O}^{perf}$ は
$R\epsilon_* : D((\Sch/S)_h) \to D((\Sch/S)_{fppf})$ の本質像に属する．
\end{lemma}

\begin{proof}
補題 \ref{etale-cohomology-lemma-refine-check-h} の判定条件を用いて証明する．条件を確認する前に，
準コンパクトかつ準分離な対象 $X$（$(\Sch/S)_{fppf}$ の対象）に対して
$$
H^i_{fppf}(X, \mathcal{O}^{perf}) = \colim_F H^i_{fppf}(X, \mathcal{O})
$$
が成り立つことに注意する．Cohomology on Sites，補題
\ref{sites-cohomology-lemma-colim-works-over-collection} を参照せよ．さらに
$H^i_{fppf}(X, \mathcal{O}) = H^i(X, \mathcal{O})$ も用いる．Descent，命題
\ref{descent-proposition-same-cohomology-quasi-coherent} を参照せよ．

\medskip\noindent
$A, f, J$ を More on Flatness，例 \ref{flat-example-one-generator} のとおりとし，
付随する概ブローアップ正方形を考える．$X$，$X'$，$Z$，$E$ はアフィンなので，
$\mathcal{O}$ の高次コホモロジーは消える．したがって
$$
0 \to
\mathcal{O}^{perf}(X) \to
\mathcal{O}^{perf}(X') \oplus \mathcal{O}^{perf}(Z) \to
\mathcal{O}^{perf}(E) \to 0
$$
が短完全列であることだけを確認すればよい．これは More on Flatness，補題
\ref{flat-lemma-h-sheaf-colim-F} の証明で示されている．

\medskip\noindent
$X, X', Z, E$ を More on Flatness，例 \ref{flat-example-two-generators} のとおりとする．
$X$ と $Z$ はアフィンなので，$H^p(X, \mathcal{O}_X) = H^p(Z, \mathcal{O}_X) = 0$ が
$p > 0$ に対して成り立つ．More on Flatness，補題 \ref{flat-lemma-funny-blow-up} により，
$H^p(X', \mathcal{O}_{X'}) = 0$ が $p > 0$ に対して成り立つ．
$E = \mathbf{P}^1_Z$ であり $Z$ はアフィンなので，$H^p(E, \mathcal{O}_E) = 0$ も
$p > 0$ に対して成り立つ．前段落と同様に，
$$
0 \to
\mathcal{O}^{perf}(X) \to
\mathcal{O}^{perf}(X') \oplus \mathcal{O}^{perf}(Z) \to
\mathcal{O}^{perf}(E) \to 0
$$
が短完全列であることの確認に帰着する．これは More on Flatness，補題
\ref{flat-lemma-h-sheaf-colim-F} の証明で示されている．
\end{proof}

\begin{proposition}
\label{etale-cohomology-proposition-h-cohomology-structure-sheaf}
$p$ を素数とする．$S$ を $\mathbf{F}_p$ 上の準コンパクトかつ準分離なスキームとする．
このとき
$$
H^i((\Sch/S)_h, \mathcal{O}^h) =
\colim_F H^i(S, \mathcal{O})
$$
が成り立つ．ここで左辺の $\mathcal{O}^h$ は構造層の h 層化を表す．
\end{proposition}

\begin{proof}
これは補題 \ref{etale-cohomology-lemma-h-sheaf-colim-F} の言い換えにすぎない．
$\mathcal{O}^h = \mathcal{O}^{perf} = \colim_F \mathcal{O}$ を想起せよ．
More on Flatness，補題 \ref{flat-lemma-char-p} を参照せよ．補題
\ref{etale-cohomology-lemma-h-sheaf-colim-F} により，$\mathcal{O}^{perf}$ を
$D((\Sch/S)_{fppf})$ の対象とみなすと，$R\epsilon_*K$ の形である．ここで
$K \in D((\Sch/S)_h)$ である．このとき $K = \epsilon^{-1}\mathcal{O}^{perf}$ であり，
$\mathcal{O}^{perf}$ は h 層なので，これは実際 $\mathcal{O}^{perf}$ に等しい．
Cohomology on Sites，節 \ref{sites-cohomology-section-compare} を参照せよ．したがって
$R\epsilon_*\mathcal{O}^{perf} = \mathcal{O}^{perf}$ である（紛らわしい記号法をご容赦願いたい）．
ゆえに主張は Leray の等式
$$
R\Gamma_h(S, \mathcal{O}^{perf}) =
R\Gamma_{fppf}(S, R\epsilon_*\mathcal{O}^{perf}) =
R\Gamma_{fppf}(S, \mathcal{O}^{perf})
$$
および等式
$$
H^i_{fppf}(S, \mathcal{O}^{perf}) =
H^i_{fppf}(S, \colim_F \mathcal{O}) =
\colim_F H^i_{fppf}(S, \mathcal{O})
$$
から従う．実際，$S$ は準コンパクトかつ準分離である（補題
\ref{etale-cohomology-lemma-h-sheaf-colim-F} の証明を参照せよ）．
\end{proof}
